\section{The common family and its cohomology bound}
\label{sec:family-window}
\hypertarget{source-family-window}{}

The average in \eqref{eq:fourier-compare} uses every retained translate at one scale. Its upper bound must therefore come from arithmetic estimates valid throughout that set of translates. We construct an interval of
scales on which the quadratic-twist estimates and the scalar
cohomology budget hold for the same auxiliary field. We then fix
the ranks and weights before choosing the translate. Construction~\ref{hyp:family-prescribed-realization} forms the signed rows at the admissible word levels. It copies the reference tuple when adjoining the chosen character values. In Section~\ref{sec:word-division} the signed transport indices and the Hermitian-dual index enter one slope identity. Its two estimates must therefore refer to the same lattices.

\subsection{Order of the choices}

Fix a primitive Dirichlet character \(\psi\) of conductor \(N_\psi=f_\psi\) in the notation of Definition~\ref{def:dirichlet-modulus}.
Two choices determine the auxiliary field once \(\psi\) is fixed. The primes \(q,d\) are chosen for an interval; the dyadic index \(e=e(z)\) can vary with the original cohomology scale \(z\) inside it.

\begin{definition}[The original auxiliary field]\label{def:family-original-host}
Put \(M=\operatorname{lcm}(2^{e+2},N_\psi,q,d)\) and identify
\(\operatorname{Gal}(\mathbb Q(\mu_M)/\mathbb Q)\) with
\((\mathbb Z/M\mathbb Z)^\times\) by the action on a primitive root of unity. Write \(U_M=(\mathbb Z/M\mathbb Z)^\times\) as in Definition~\ref{def:dirichlet-modulus} and put
\[
 U_{\psi,e;q,d}=\left\{a\in U_M:
 a\equiv\pm1\pmod{2^{e+2}},\ \psi(a)=1, \left(\frac aq\right)=\left(\frac ad\right)=1\right\}.
\]
The original auxiliary field and its Galois group are
\[
 E_{\psi,e;q,d}=\mathbb Q(\mu_M)^{U_{\psi,e;q,d}},\qquad
 \operatorname{Gal}(E_{\psi,e;q,d}/\mathbb Q)
       \simeq(\mathbb Z/M\mathbb Z)^\times/U_{\psi,e;q,d}.
\]
Write \(E=E_{\psi,e;q,d}\) when these parameters are fixed.

Let \(K_\psi\subset\mathbb Q(\mu_{f_\psi})\) be the field fixed by
\(\ker\psi\), and set \(C_e=\mathbb Q(\mu_{2^{e+2}})^+\).
The four kernel conditions are the stabilizers of the four factors
in the compositum
\begin{equation}\label{eq:host-character-compositum}
 E=K_\psi C_e(\sqrt{-q},\sqrt{-d}),\qquad
 Q_S=\operatorname{lcm}(f_\psi,2^{e+2}).
\end{equation}

Indeed \(q,d\equiv3\pmod4\), so the two quadratic characters
cut out the fields of the displayed negative square roots. The
intersection of the four stabilizers fixes their compositum.
On sufficiently late selected intervals \(q,d\nmid Q_S\) and \(M=Q_Sqd\).
This is the field in \eqref{eq:cohom-host-field}.
\end{definition}

The family uses this same \(E\) for the strict lengths in
Definition~\ref{def:strict-length} and for all translates at
the corresponding rescaled point. Its late intervals satisfy
Propositions~\ref{prop:original-window-budget}
and~\ref{prop:common-host-prime-pair}.
Fix a sufficiently small
\(\epsilon>0\). Each estimate will hold once the family parameters
exceed specified lower bounds. The bounds may depend on
\(\epsilon\) and on the parameters chosen earlier. For estimate \(i\), write \(T_{ij}\) for its lower threshold on the \(j\)-th parameter, with dependencies taken in the order of choice. Intersect the finitely many ranges for \(\epsilon\), and at step \(j\) impose
\[
 t_j>\max_i T_{ij}(\epsilon,t_1,\ldots,t_{j-1}).
\]
The existence theorem in the construction then
supplies a family satisfying all these inequalities with one
choice of weights. In particular the cohomology, rank and lattice
estimates apply to the same family.

The final comparison uses local length and prime-count
bounds, dimension and Hilbert--Schmidt estimates for the
row spaces, irreducibility with controlled multiplicities,
and divisibility and slope estimates for the lattices.
In their final bounds, let \(e_D\) be the exponent of
\(x^\epsilon\) from the divisibility estimate and let \(e_L\)
be the corresponding exponent from the remaining losses.
Section~\ref{subsec:fixed-weight-absorption} uses the values
\(e_D=2\) and \(e_L=0\). The prime pair is chosen first. Its associated field will then be the field of definition of the strict cohomology groups in the divisibility estimate.

\subsection{Choosing one pair of quadratic twists on an interval}
\label{sec:pair-selection}
\hypertarget{source-prime-selection}{}

The quotient of the translate group retains twists by each auxiliary quadratic character and by their product. To discard their contributions uniformly, the same two primes must control all three twists at every argument in the selected interval.

\begin{proposition}[A quadratic pair on the full interval]\label{prop:quadratic-prime-pair}
Fix the modulus \(M\), \(0<\epsilon<1/52\) and \(K\ge1\).
For all sufficiently large \(Y\), put \(w=Y^{\epsilon/5}\)
and \(Z=Y^{400K}\). There are primes \(q\equiv3\pmod 8\) and
\(d\equiv7\pmod{16}\), with
\[
 \tfrac12w<\log q\le w,\qquad
 \tfrac12 w\log^2Y<\log d\le w\log^2Y,
 \qquad d\equiv-1\pmod{q\prod_{2<p\le w}p},
\]
where the product is over primes. For every primitive inducer at
level \(8M\), the twists by the Legendre character at \(q\), at
\(d\), or by their product satisfy
\[
 |G_\chi(y)|\le y^{1/2+\epsilon/2},\qquad
 |L'(1,\chi)/L(1,\chi)|\le 4\log y
       \qquad(Y\le y\le Y^K).
\]

For the twist by the character at \(q\), take \((F,\log U)=(8M,w)\).
For the twists by the character at \(d\) and by the product of the two
characters, take \((F,\log U)=(8Me^w,w\log^2Y)\). Each twist is
zero-free for \(\RePart s>1-1/400\) and \(|\Im s|\le FU\) with its
indicated pair. The lower
threshold for \(Y\) depends only on \(M,K,\epsilon\).
\end{proposition}

\begin{proof}
\proofstep{Count the candidate primes and control correlations.}
Consider a reduced class \(a_0\bmod m_0\), with \(8\mid m_0\).
Each prime dividing \(m_0\) is at most \(w\) or exceeds \(Z\).
The candidates \(\mathcal P\) are the primes in this class in
\((\sqrt U,U]\). In each application
\[
  \tfrac14\log^2Y\log m_0\le\log U,\qquad
  Y^{\epsilon/6}\le\log U\le w\log^2Y .
\]
Put \(B=m_0 Z^2\), so \(\log B\le5\log U/\log^2Y\).
For each modulus \(m\le B\) under consideration, apply the relative
prime number theorem in progressions of Thorner and Zaman
\cite[Corollary 1.4 and Remark 1.5]{tz-progressions}.

For \(\sqrt U\le t\le U\) one has \(t\ge m^{12}\) eventually.
The relative error in the log-weighted count is \(O_A(Y^{-A})\)
for every fixed \(A\), since \(\log t/\log m\gg\log^2Y\) and
every fixed positive power of \(\log t\) dominates \(\log^2Y\).
The possible secondary term may be taken from one primitive real
character \(\xi\) for all these moduli. For an absolute sufficiently small \(c>0\), the Landau--Page theorem \cite[Lemma 7.1]{landau-page} gives at most one primitive real character satisfying both conditions
\[
 \operatorname{cond}\xi\le B,\qquad L(\beta,\xi)=0,\quad \beta>1-c/\log B.
\]
Such a zero is in the
exceptional region for every multiple of the conductor up to \(B\).
A local exception below this threshold changes the main term by
a relative \(O(\exp(-c\log t/\log B))\).

For a reduced class \(a\bmod m\) write \(u_a=\xi(a)\) if the
exception exists and its conductor divides \(m\), and \(u_a=0\)
otherwise. Terms involving \(\beta\) are omitted in the latter
case. Partial summation gives, with the same relative error,
the unweighted count
\[
  \frac{1}{\varphi(m)}\int_{\sqrt U}^U
                   \frac{1-u_a t^{\beta-1}}{\log t}\,dt .
\]
In particular the small-main-term case \(u_a=1\) is allowed.
With \(\lambda=\min(1,(1-\beta)\log U)\), the numerator is
\(\asymp\lambda\) on the interval. Indeed
\[
 1-t^{\beta-1}=1-e^{-(1-\beta)\log t}
 \asymp\min\{1,(1-\beta)\log t\}\asymp\lambda
 \qquad(\sqrt U\le t\le U).
\]

The weighted main term
\((t-t^\beta/\beta)/\varphi(m)\) is \(O(t\lambda/\varphi(m))\)
there, by \(-\log\beta\le2(1-\beta)\) for \(\beta\) near one.
Consequently both the boundary terms and the integral of the
weighted error in partial summation are
\(O_A(Y^{-A}U\lambda/(\varphi(m)\log U))\).
The effective zero bound \cite[(1.2)]{tz-progressions} and the size of the modulus give
\[
 1-\beta\gg B^{-1/2},\qquad B=U^{o(1)}.
\]
Hence the relative loss from the exceptional factor is at most a subpower of \(U\), and \(|\mathcal P|\ge U^{9/10}\) eventually.

For distinct \(w<r,s\le Z\), the candidate primes exceed \(Z^2\),
so quadratic reciprocity identifies \((rs/p)\), in the prescribed
class modulo \(m_0\), with a fixed sign times
\(\eta_{r,s}(a)=(a/r)(a/s)\) on the classes \(a\bmod m_0rs\)
above \(a_0\). The two new prime factors are coprime to \(m_0\).
The primary term in this weighted class sum is proportional to
\[
       \sum_{a_r\in\mathbb F_r^\times}
       \sum_{a_s\in\mathbb F_s^\times}(a_r/r)(a_s/s)=0.
\]

If the exceptional conductor divides \(m_0\), its contribution is
also constant across these classes and vanishes by the same sum.
Otherwise the progression modulo \(m_0\) has its full primary
main term. The exception contributes on the larger modulus only
when its conductor divides that modulus and its local characters
on both new factors equal the two nontrivial quadratic factors.
Thus its prime-to-\(m_0\) conductor must be \(rs\).
Only one unordered pair can have that property. The relative
progression errors give \(O_A(Y^{-A})\) for every other correlation.

For any \(|\kappa(n)|\le1\) write
\[
 S_p(y)=\sum_{n\le Z}\Lambda(n)\min(1,y/n)\kappa(n)(n/p).
\]
Primes \(\le w\) contribute \(O(w)\). The powers of exponent at
least two contribute \(O(\sqrt y)\), since Chebyshev's bound gives
\(\sum_{r^j\le t,\,j\ge2}\log r\ll\sqrt t\).
On the shell \(2^k y<r^j\le2^{k+1}y\), the weight is at most
\(2^{-k}\), so the tail costs at most
\[
 \sum_{k\ge0}2^{-k} O\bigl((2^{k+1}y)^{1/2}\bigr)\ll y^{1/2}.
\] Dyadic summation bounds the diagonal mass and the off-diagonal mass separately:
\[
 \sum_r |a_r|^2\ll y\log(2y),\qquad \sum_r |a_r|\ll y(1+\log Z).
\]
The exceptional pair costs at
most \(|a_r|^2+|a_s|^2\) in the expanded mean square. All other
off-diagonal terms together cost
\(O(Y^{-K-1}y^2(1+\log Z)^2)=O(y)\).

Hence \(\sum_p|S_p(y)|^2\ll |\mathcal P| y\log(2y)\) for every
weight and every \(Y\le y\le Y^K\). For nonintegral
\(y\), \(yS_p'(y)\) is the sum with the weight replaced by
\(\kappa(n)\mathbf 1_{n>y}\). Thus the same estimate gives
\(\sum_p |S_p'(y)|^2\le C y^{-1}\log y\,|\mathcal P|\).
On \([V,b]\subseteq[V,2V]\), absolute continuity and
Cauchy--Schwarz imply
\[
 \sum_p\sup_{y\in[V,b]}|S_p(y)|^2
 \le 2\sum_p|S_p(V)|^2
      +2V\int_V^b\sum_p|S_p'(u)|^2\,du
 \ll V\log(2V)\,|\mathcal P|.
\]
For fixed \(\tau>0\) the exceptional fraction for threshold
\(V^{1/2+\tau}/2\) is \(O(V^{-2\tau}\log(2V))\).
The total exceptional fraction on dyadic intervals up to any endpoint
in \([Y,Y^K]\) is bounded by extending the sum to infinity:
\[
 \sum_{j\ge0}(2^jY)^{-2\tau}\log(2^{j+1}Y)
 =Y^{-2\tau}\sum_{j\ge0}2^{-2j\tau}
            \bigl(\log(2Y)+j\log2\bigr)
 \ll_\tau Y^{-2\tau}\log(2Y).
\] The same conclusion holds for
any fixed finite set of base characters.

\proofstep{Discard the primes giving a zero in the box.}
Let the weight now be any of a fixed number of primitive base
characters of conductor at most \(F\le\sqrt U\), and put \(Q=FU\)
with \(F\ge1\). Each twisted character is primitive and
nonprincipal, of conductor \(fp\), where \(f\le F\) is its base
conductor. This conductor determines \(p\) as its unique prime
factor greater than \(F\). Huxley's density estimate for primitive
characters up to conductor \(Q\), at height \(T=Q\), bounds the
number with a zero of real part \(>1-1/400\) in that box by
\[
 O_\eta\big((Q^2 T)^{12/(5\cdot400)+\eta}\big)
                 \ll Q^{1/4}\le U^{3/8}
\]
on taking \(\eta=1/100\) \cite[(1.1)]{huxley}. Thus the union of
the failed zero-box and maximal-sum tests omits a positive
fraction of the candidate primes, uniformly in the list for a
fixed bound on its size.

\proofstep{Recover the centered sums.}
The mean-square and zero-density estimates have removed exceptional choices of the auxiliary prime. To use a remaining choice in the translate average, we need a bound for the centered sum itself, uniform throughout the original interval.

For a surviving character choose the margin and strip width together:
\[
 \alpha=\epsilon/(1600K),\qquad \nu=1/400-\alpha.
\]
The logarithmic derivative is
\(O_\alpha(\log Q)\) on \(1-\nu\le\sigma\le 2\),
\(|t|\le Q/2\). Indeed, the partial fraction expansion of the
completed primitive \(L\)-function, differenced at \(2+it\),
has an \(O(\log Q)\) remainder after retaining the zeros within
unit vertical distance. The unit-band count is \(O(\log Q)\)
by the argument principle (equivalently by the real part of the
logarithmic derivative at \(2+it\) and Stirling's formula).
The distant bands contribute \(O(\log Q)\) by inverse-square decay. In the unit band the zero-box condition gives
\[
 \sum_{|\Im\rho-t|\le1}\frac{1}{|\sigma+it-\rho|}\ll \frac{\log Q}{\alpha}.
\]


Set \(Z'=\lfloor Z\rfloor+1/2\) and apply truncated Perron to
\(\sum_{n\le Z}\Lambda(n)\chi(n)/n\), on
\(\Re s=1/\log Z'\) at height \(Q/2\). The truncation costs
\(O(Z'\log^2 Z'/Q)\). For \(Z'/2<n<2Z'\) use the
distance of \(Z'\) from integers in
\(1/|\log(Z'/n)|\ll Z'/|Z'-n|\). Away from that range use
absolute convergence and \(-\zeta'/\zeta(1+1/\log Z')
\ll\log Z'\). Moving to \(\Re s=-\nu\) crosses the Perron pole
with residue \(-L'(1,\chi)/L(1,\chi)\).

For \(y\le Y^K\), the contour bounds after multiplying by \(y\) are
\[
 O_\alpha\!\left(y\left[
   \frac{Z'\log^2Z'}{Q}
   +\frac{\log Q}{Q}
   +(Z')^{-\nu}\log^2 Q\right]\right).
\]
The first term is truncation, the second is horizontal and the third
is vertical. The lower bound \(\log Q\ge Y^{\epsilon/6}\) makes
the first two negligible. For the last, \(Z'\asymp Y^{400K}\)
and \(400K\alpha=\epsilon/4\), so its size is
\(O_\epsilon(Y^{\epsilon/4}Y^{\epsilon/2})=
O_\epsilon(Y^{3\epsilon/4})\), using \(\log Q\le Y^{\epsilon/4}\).
By \eqref{eq:centered-sum} this is the error in \(G_\chi(y)+S_p(y)\)
with the primitive base used as the weight.

Set \(\tau=\epsilon/2\) in the maximal-sum test. For large \(Y\)
the sum and the tail together give the asserted bound for
\(|G_\chi(y)|\). Dividing the centered identity by \(y\) gives
\[
 |L'(1,\chi)/L(1,\chi)|\le
 y^{-1/2+\tau}+\sum_{n\le y}\Lambda(n)/n\le4\log y .
\]
Indeed, summing the divisor identity for \(\log n\) and bounding
the integer-part error by Chebyshev's estimate gives
\begin{equation}\label{eq:factorial-mangoldt-bound}
 \log(\lfloor y\rfloor !)
 =\sum_{d\le y}\Lambda(d)\lfloor y/d\rfloor
 =y\sum_{d\le y}\frac{\Lambda(d)}d+O(y).
\end{equation}
Stirling's formula on the left yields the stated \(\log y+O(1)\).

\proofstep{Select \(q\) and then \(d\).}
In the first application \(\log U=w\), \(p=q\equiv3\bmod8\)
and the conductor bound is \(F=8M\). Use all primitive
characters inducing characters at level \(8M\).
In the second application \(\log U=w\log^2Y\). Choose
\(p=d\equiv7\bmod16\) and \(d\equiv-1\) modulo the product
of \(q\) and the odd primes \(\le w\). This is a reduced class
by the Chinese remainder theorem. Its modulus has logarithm
at most \(4w\) by the prime number theorem.

The sole factor
greater than \(w\) is \(q>Z\). Include the first list and every
twist in that list by \(q\), with \(F=8M e^w\le\sqrt U\).
The two applications give
\[
 w/2<\log q\le w,\qquad \tfrac12 w\log^2Y<\log d\le w\log^2Y.
\]
This proves all three twist estimates and zero-box assertions on
one interval. Their thresholds depend only on \(M,K,\epsilon\).
Quadratic reciprocity identifies the two Legendre characters
with those of discriminants \(-q,-d\), including their value at
two.
\end{proof}


\subsection{Two level primes with disjoint exceptional sets}
\hypertarget{source-level-primes}{}

The primes defining the quadratic characters are now fixed. The levels of the forms require two further primes. Both must split as prescribed, and the prime divisors introduced by their level factors must satisfy the separation below.

Fix the primes \(q,d\) of Proposition~\ref{prop:quadratic-prime-pair}.
For \(x\) in their common interval take the modulus
\(N=Q_Sqd\) of \eqref{eq:host-character-compositum}.
The restriction from \(\operatorname{Gal}(K_\psi C_e/\mathbb Q)\)
onto \(\operatorname{Gal}(C_e/\mathbb Q)\) gives
\[
   2^e\le [K_\psi C_e:\mathbb Q]=n,\qquad
   Q_S\le 4 f_\psi 2^e\le 4f_\psi n,\qquad n\asymp x^{2+8\epsilon}.
\] With
\(b=\exp(x^\epsilon)\), \(k\ge1\) the weight parameter and
\(1\le r\le x^5\) the rank, define \(X=(kb)^8\). The preceding
prime size bounds and \(Y\le x\) give
\[
 d\mid N,\quad d\equiv3\pmod4,\quad
 \log N\le3x^{\epsilon/5}\log^2 x,\qquad
 \log(4dN^2)=o(\sqrt{\log X}).
\]
Write \(K_0=\mathbb Q(\sqrt{-d})\). \begin{definition}[Exceptional primes for a type level]\label{def:type-exceptional-primes}
For an integer \(g\ge2\) and a positive rank \(r\), define
\[
 B_r(g)=\{g\}\ \cup\!
 \bigcup_{j=1}^r\{\ell\text{ prime}\mid
                 \ell^{4v_\ell(g^j-1)}>g^j\}.
\]
These are the possible exceptional residual primes at the type
level \(g\).
\end{definition}

\begin{proposition}[Disjoint exceptions at the two type levels]\label{prop:two-type-primes}
For the auxiliary primes \(q,d\) in Proposition~\ref{prop:quadratic-prime-pair}
with \(N=Q_Sqd\), \(X=(kb)^8\) and \(1\le r\le x^5\), there are
at every sufficiently large eligible \(x\) distinct
primes \(g_0,g_1\) and elements \(\pi_i\in\mathcal O_{K_0}\) satisfying
\[
 X<g_i\le2X,\quad N_{K_0/\mathbb Q}(\pi_i)=g_i,\quad \pi_i\equiv1\pmod{N\mathcal O_{K_0}},\quad B_r(g_0)\cap B_r(g_1)=\varnothing.
\] There are also distinct
primes \(p,p',q_0\le N^{C_8}\), with \(p,p'\equiv1\pmod{15N}\)
and \(q_0\equiv1\pmod N\), for an absolute \(C_8\).
\end{proposition}

\begin{proof}
\proofstep{Count degree-one ideals in the ray class.}
Let \(L/K_0\) be the ray class field for \(N\mathcal O_{K_0}\).
Its degree \(h_N\) is at most \(h_{K_0}N^2\) by the ray class
exact sequence. Minkowski's bound gives a representative ideal
of norm \(O(\sqrt d)\) for each ideal class, and the number of
ideals of norm \(j\) is at most the divisor function. Thus
\(h_N\ll N^2\sqrt d\log(2d)\).
In the abelian relative Chebotarev theorem take the base field
\(K_0\) of degree two and absolute discriminant \(d\), take
the whole Galois group as the abelian subgroup, and take the
identity conjugacy class \cite[Theorems 1.4 and 3.3]{tz-chebotarev}.
Each relative Hecke conductor divides \(N\mathcal O_{K_0}\)
and has norm at most \(N^2\). Put \(D=4dN^2\).

For
\(u=X,2X\) the theorem counts unramified prime ideals of norm
\(\le u\) with trivial Frobenius as
\[
 \frac{\operatorname{Li}(u)-\theta\operatorname{Li}(u^\beta)}
 {h_N}\left(1+O\left(
 e^{-c\log u/\log D}+e^{-c\sqrt{\log u}}\right)\right).
\]
Here \(\theta=0\) in the absence of an exceptional zero and
\(\theta=1\) otherwise. The hypothesis \(u\ge D^{C}\) for
a sufficiently large absolute \(C\) holds. Stark's bound in
the same theorem sequence gives \(1-\beta\gg D^{-C'}\) for
some absolute \(C'\). In its notation
\(\lambda_1=(1-\beta)\log Q\gg Q^{-2}\) with \(Q\) polynomial
in \(D\) for degree two. This includes the trivial relative
character.

Set
\[
 \lambda=\min(1,(1-\beta)\log X)
\]
in the exceptional
case and \(\lambda=1\) otherwise. The difference between the
two main numerators is \(\asymp X\lambda/\log X\), by integrating
\((1-\theta t^{\beta-1})/\log t\) on \([X,2X]\). Each numerator
is itself \(O(X\lambda/\log X)\). For \(\beta\ge3/4\) this
follows by integrating \(1/\log t\) from \(u^\beta\) to \(u\),
and when the zero is farther away use \(\lambda=1\) for large
\(X\). Thus the two endpoint errors are \(o(X\lambda/(h_N\log X))\).
The resulting count is \(\gg X/(D^{C''}\log X)\).

Let \(I\) be the ray-class ideal count on \([X,2X]\). After excluding residue degree two, ramification and divisors of the modulus, the number of underlying rational primes is at least
\[
 \tfrac12\bigl(I-O(\sqrt X+\log N)\bigr).
\]
The factor \(1/2\) allows for both degree-one primes over the same rational prime.
Consequently there are at least \(X^{9/10}\) candidate rational
primes for every sufficiently large \(x\), uniformly in these hosts.

For a counted ideal the ray condition gives a generator
\(\pi\equiv1\bmod N\mathcal O_{K_0}\). Integrality follows
from the integrality of its ideal. Write
\(\pi=(A+B\sqrt{-d})/2\) in the basis
\(1,(1+\sqrt{-d})/2\). Then
\[
 A^2+dB^2=4g,\qquad N\mid B,\qquad 2N\mid A-2-B .
\] In particular
\(g\equiv1\bmod N\). For the type primes use the short norm calculation on the interval \((X,2X]\).

\proofstep{Avoid the exceptional set of the first prime.}
The ray class imposes the splitting conditions. The second norm must also avoid the exceptions of the first.

For a given \(j\) there are at most three such primes. Four of
the corresponding primary factors would have product greater
than \(g^j\) and also dividing \(g^j-1\). Thus
\(|B_r(g)|\le3r+1\). Fix \(\ell,j\), and consider all the
integers \(X<g\le2X\) satisfying the indicated strict inequality.
If the set is nonempty, let \(v\) be the least of their
\(v_\ell(g^j-1)\). The minimizer satisfies \(\ell^{4v}>X\). By minimality every counted integer satisfies
\[
 g^j\equiv1\pmod{\ell^v},\qquad \ell^v>X^{1/4}.
\]

There are at most \(2j\) roots, since the unit group is cyclic
for odd \(\ell\) and is a product of a cyclic group and a
group of order at most two for \(\ell=2\).
Each residue class occurs at most \(1+X/\ell^v\le1+X^{3/4}\)
times. For any fixed candidate \(g_1\), the number of integers
with \(B_r(g)\cap B_r(g_1)\ne\varnothing\) is therefore at most
\[
 (3r+1)\big(1+2r^2(1+X^{3/4})\big)<X^{9/10}.
\]
The last inequality follows from \(r\le x^5\) and
\(X\ge\exp(8x^\epsilon)\). A second candidate \(g_0\) has disjoint
exceptional set. It is distinct since \(g_1\in B_r(g_1)\).

\proofstep{Choose the remaining split primes.}
The small primes needed with these two are obtained by applying
the least-prime bound in the class one successively to \(15N\),
\(15Np\), and \(Npp'\). For a modulus \(m\ge2\) the bound
\(\ll m^{C}\) follows also from the relative progression estimate
above at \(t=m^C\), with \(C\) an absolute constant large enough
to make the relative error less than \(1/2\) and the positive
weighted main term nonzero. Finitely many small moduli are
absorbed by enlarging the exponent. The successive moduli ensure distinctness and give
\[
 p,p'\equiv1\pmod{15N},\qquad q_0\equiv1\pmod N,\qquad \max(p,p',q_0)\le N^{C_8}
\]
for an absolute \(C_8\), hence smaller than \(b\) and both level
primes for sufficiently large \(x\).
\end{proof}

\begin{auxlemma}[The two characters used to remove ramification]
\label{lem:split-deletion-characters}
Let \(E\) be the original CM field of
Definition~\ref{def:family-original-host}, let \(F=E^+\), and let
\(p,p'\) be the distinct primes congruent to \(1\bmod 15N\) in
Proposition~\ref{prop:two-type-primes}. For each \(h\in\{p,p'\}\)
there is a finite-order Hecke character \(\delta_h\) of \(E\),
trivial on \(\mathbb A_F^\times\) and at every infinite place.
It is unramified away from the primes above \(h\). For each \(w\mid h\)
its restriction to \(\mathcal O_{E_w}^\times\) factors through the
residue field and satisfies
\[
       \bigl(\delta_{h,w}|_{\mathcal O_{E_w}^\times}\bigr)^a=1
                   \quad\Longleftrightarrow\quad 15\mid a
                   \qquad(a\ge1).
\]
In particular its finite conductor divides \(h\mathcal O_E\).
Take \(\delta=\delta_p\) and \(\delta'=\delta_{p'}\) in
Construction~\ref{hyp:family-prescribed-realization}.
\end{auxlemma}

\begin{proof}
Put \(k_0=\mathbb Q(\sqrt{-d})\subset E\) and let \(c\) denote complex
conjugation. The cyclotomic modulus \(N=Q_Sqd\) for \(E\) divides
\(h-1\), so \(h\) splits completely in \(E\).
Choose one of the two primes \(\mathfrak q\) of \(k_0\) above \(h\).
The set \(S=\{w:w\mid\mathfrak q\}\) and its conjugate \(cS\)
partition the primes of \(E\) above \(h\). Every residue field here
is identified with \(\mathbb F_h\) by its prime subfield.
Since \(\mathbb F_h^\times\) is cyclic and \(15\mid h-1\), choose
a character \(\omega:\mathbb F_h^\times\to\mathbb C^\times\) of order
\(15\). Put \(\mathfrak m=\prod_{w\in S}w\) and prescribe on its residue
units the character
\[
        \theta((z_w)_{w\in S})=\omega\!\left(\prod_{w\in S}z_w\right).
\]

This prescription is compatible with every global unit. Indeed
\(d\equiv7\bmod16\), so \(\mathcal O_{k_0}^\times=\{\pm1\}\). For
\(u\in\mathcal O_E^\times\), reduction of the norm at the completely
split prime gives
\[
       \prod_{w\in S}(u\bmod w)
          =N_{E/k_0}(u)\bmod\mathfrak q\ \in\{\pm1\}.
\]
Both values lie in the kernel of the odd-order character \(\omega\).
In the ray class sequence
\[
 \mathcal O_E^\times \longrightarrow
 \prod_{w\in S}(\mathcal O_E/w)^\times\longrightarrow
 \operatorname{Cl}_{\mathfrak m}(E)\longrightarrow
 \operatorname{Cl}(E)\longrightarrow 1
\]
use the middle map induced by the local unit ideles.
Thus \(\theta\) gives a character on their image.
Extend it to a character \(\lambda\) of the finite abelian ray class
group. Explicitly, when adjoining an element \(b\) of order \(n\)
modulo the current subgroup, choose an \(n\)-th root of the assigned
value of \(b^n\) as the value of \(b\). The resulting ray class
character is trivial at infinity and unramified outside \(S\).

Define \(\delta_h(z)=\lambda(z)/\lambda(cz)\) on the idele classes
of \(E\). Complex conjugation fixes the image of \(\mathbb A_F^\times\),
so the quotient is trivial there. On local residue units this gives
\[
 \delta_{h,w}|_{\mathbb F_h^\times}=
       \begin{cases}\omega,&w\in S,\\ \omega^{-1},&w\in cS.\end{cases}
\] This gives the asserted order at every prime
over \(h\). Away from these primes both factors are unramified.
Both characters have finite image, so their values can be included in
one coefficient field by adjoining finitely many roots of unity.
The two constructions use disjoint sets of ramified primes.
At coefficient characteristic \(\ell\) the choice is \(\delta'\)
if \(\ell=p\), and \(\delta\) otherwise. In either case it is
unramified above \(\ell\) and its inertia image
has order \(15\) above \(p'\) or \(p\), respectively.
\end{proof}


\subsection{Selecting the same top positions for all Frobenius cycles}
\label{subsec:frobenius-selection}

The top indices in Section~\ref{sec:comparison} determine the coefficients of the surviving characters. Their prescribed proportions in the quotient must be reconciled with restrictions along the Frobenius cycles at many primes. Fix a scale \(x\) in a selected interval
of Proposition~\ref{prop:quadratic-prime-pair} and \(0<\epsilon<1/52\). Set
\[
 h=2+8\epsilon,\quad \delta=x^{-\epsilon/8},\quad
 K=\lceil8(h+2)/\epsilon\rceil,\quad r=K+1 .
\]
\begin{paramdefinition}[Group data and admissible cycle windows]\label{hyp:frobenius-top-cycle}
The finite abelian group \(A\) has order \(n\asymp x^h\).
In the arithmetic application \(A_0=\ker\pi_\psi\) is the
subgroup of Definition~\ref{def:translate-scale}; its index
\(o\) is the order of the fixed target character \(\psi\).
The combinatorial argument uses only this index. For each of a finite
set of primes \(p\), fix a step \(a_p\in A\) of order
\(d_p\ge x^{\epsilon/4}\) and an origin on each of its
cycles. Number the positions on that cycle from \(0\) to
\(d_p-1\). In the application the prime list is restricted to a prescribed congruence class and to primes inert in the auxiliary quadratic field. The two size tests on such a prime and its step are
\[
 p<X'=\exp(x^{\epsilon/3}),\qquad \operatorname{ord}(a_p)=d_p\ge x^{\epsilon/4}.
\]
Each prime has
weight \(\omega_p=\log(p)/p\).

Let \(L_2\) be a common upper bound for the short distances
used at these primes. We require \(L_2\le L_b\) for a
constant fixed before \(x\). Put
\[
\begin{gathered}
 u=\lceil1/\delta\rceil,\qquad l=(6r+1)u,\qquad
 M=\lceil(h+3)10^{6r+1}\log x\rceil,\\
 G=(M+5)l,\qquad
 N_1=1+\lceil\max(6eG\delta,(h+3)\log_2 x)\rceil .
\end{gathered}
\]
A full block is one of the consecutive aligned intervals
of length \(l\) starting at position zero. Its \(6r+1\)
cells have length \(u\). A marked block has its first,
third and subsequent odd-numbered cells empty. Each
intervening cell contains exactly one selected position
of even index on the cycle. Call a cycle good if

\begin{enumerate}
\item every cyclic interval of length \(L_2\) contains at
most \(K\) selected positions,
\item among every \(M+1\) successive full blocks, taken
cyclically in the list of full blocks, there is a marked block,
\item every cyclic interval of length \(G\) has at most
\(N_1\) selected positions.
\end{enumerate}

The quantities \(L_b\) and \(\epsilon\) are fixed independently of the scale.
\end{paramdefinition}

\begin{auxlemma}[Lengths of admissible Frobenius cycles]
\label{lem:cycle-geometric-margins}
On the arithmetic list in Definition~\ref{hyp:frobenius-top-cycle}, uniformly in the prime list,
\[
 L_2\le 2(620J/\epsilon+1),\qquad L_2<u,\qquad 3G\le d_p
\]
for large \(x\). In a good cycle each marked block supplies \(3r\) even selections separated by more than \(u\). The gap to the next marked even position is at most \(G\) and contains at most \(N_1\) intermediate selections.
\end{auxlemma}
\begin{proof}
For the arithmetic choice,
\(L_2\) is twice the supremum of
\(\lceil\log_p((kx^{26})^4)\rceil\) over the primes on this list. Every odd prime up to \(Y^{\epsilon/5}\)
splits in \(\mathbb Q(\sqrt{-d})\) by the congruence on
the auxiliary discriminant, and \(2\) splits as well.
Here the prime-selection interval has \(x\le Y^J\)
with \(J\) bounded in this family. Thus \(\log p>(\epsilon/(5J))\log x\) at each relevant prime.

Substituting \(k\le2x^{9/2+10\epsilon}\) in the definition gives, uniformly for large \(x\),
\[
 L_2\le 2\bigl(620J/\epsilon+1\bigr).
\]
These conditions bound both the crowded regions and the gaps.
The deterministic geometry of these tests requires
\[
 L_2<u,\qquad 3G\le d_p.
\]
Both hold eventually, uniformly in the primes of the list, since
\[
 G/d_p=O_\epsilon(x^{-\epsilon/8}\log x)=o(1),\qquad L_2/u=O_\epsilon(x^{-\epsilon/8}).
\]

Every marked block supplies \(3r\) even selected positions
with gaps greater than \(u\) and buffers of at least \(u\)
at its ends. The next marked block starts within \((M+2)l\),
including a possible final partial block. In particular the
gap to the next selected even position is at most \(G\) and
the number of intermediate selected positions is at most
\(N_1\). The first condition bounds by \(r-1\) the number
of possible incoming selected indices in a preceding short
window. Thus the constants are also valid when an origin
falls near the end of the numbered cycle.

\end{proof}

\begin{proposition}[One top set for the selected cycles]\label{prop:top-set-cycles}
For the finite group and cycle data of Definition~\ref{hyp:frobenius-top-cycle},
there is a single set \(T\subset A\) with
\(\delta n/2\le|T|\le2\delta n\) and all of the following
bounds. For any coset \(C\) of \(A_0\) let \(\lambda_C\)
equal \(4/(3o+1)\) for \(C=A_0\) and \(3/(3o+1)\)
otherwise. Then
\begin{align}
  \big||T\cap C|-\lambda_C|T|\big|
    &\le x^{-1+\epsilon}|T|,\label{eq:coset-shares}\\
  \sum_p\omega_p\#\{t\in T:\text{the }p
           \text{-cycle through }t\text{ is bad}\}
    &\le 4C_bx^{-2}\delta n\sum_p\omega_p .\label{eq:bad-cycles}
\end{align}
Here \(C_b\) and the lower threshold for \(x\) depend only
on \(\epsilon,L_b\), the bound on \(o\) and the constants in
\(n\asymp x^h\). There is no dependence on the steps, the
origins or the number of primes.
\end{proposition}

\begin{proof}
\proofstep{Bound the failure probability at one selected position.}
Choose the membership indicators independently with success
probability \(\delta\) on \(A_0\) and \(3\delta/4\) on
its complement. Fix a position and condition on its being
selected. All the other trials retain their initial laws.
In each short window failure requires at least \(K\) further
successes. There are \(d_p\) starting positions. Summing over the possible \(K\)-subsets in each window and using \(d_p\le n\) gives
\[
 d_p\binom{L_2}{K}\delta^K\ll_{L_b,\epsilon}x^{h}x^{-(h+2)}\ll x^{-2}.
\]
For a run of \(M+1\) full
blocks discard the block containing the conditioned position
if present. At least \(M\) untouched blocks remain, on
independent trials. A cell required to be empty has
probability at least \((1-\delta)^u\). The probability for
a cell required to contain exactly one even index is at least
\(\lfloor u/2\rfloor(3\delta/4)(1-\delta)^{u-1}\), by
summing the disjoint singleton events.

For \(\delta\le1/100\)
we have
\[
 (1-\delta)^u\ge e^{-(1+\delta)/(1-\delta)}>0.3
\]
and \(\lfloor u/2\rfloor\delta\ge(1-\delta)/2\).
The \(6r+1\) cell events in an untouched block are independent, each with probability greater than \(1/10\). Thus
\[
 \Pr(\text{block is marked})\ge10^{-(6r+1)}.
\]
The
chance of a failed run anywhere is at most
\(d_p\exp(-10^{-(6r+1)} M)\ll x^{-3}\).
For a length-\(G\) window the conditioned count is bounded
by one plus a binomial variable with parameters \(G,\delta\).
For \(s=N_1-1\le G\) the union bound over \(s\)-subsets gives
\[
 \Pr(\text{at least }s\text{ unconditioned successes})\le \binom{G}{s}\delta^s\le(eG\delta/s)^s\le2^{-s}.
\]
The tail is zero for \(s>G\).
Summing over starts gives \(O(x^{-3})\). Together the
conditional failure probability is at most \(C_bx^{-2}\).

\proofstep{Sum the failure counts with the prime weights.}
Let \(I_t\) be the membership indicator at position \(t\)
and let \(F_{p,t}\) be the failure of its cycle test at \(p\).
Combining the three conditional probabilities gives
\(\Pr(F_{p,t}\mid I_t=1)\le C_bx^{-2}\). Consequently
\[
 \mathbb E\sum_p \omega_p\sum_t I_t\mathbf1_{F_{p,t}}
   \le C_bx^{-2}\sum_p\omega_p\sum_t\Pr(I_t=1)
   \le C_b x^{-2}\delta n\sum_p\omega_p .
\]
Markov's inequality gives \eqref{eq:bad-cycles} with
probability at least \(3/4\), also when this expectation vanishes.

\proofstep{Control the coset counts for the same set.}
The surviving selections satisfy the restrictions along the Frobenius cycles. Their coset counts still have to give the prescribed character weights, so the probability estimate must concern this same selection.

The mean of \(|T|\) lies between \(3\delta n/4\)
and \(\delta n\), and its variance is at most \(\delta n\).
For each \(C\), the signed count
\(|T\cap C|-\lambda_C|T|\) has mean zero, because the
cosets have equal cardinalities and the membership
probabilities have ratio \(4:3\). Its variance is at most
\(\delta n\) by independence. Chebyshev's inequality for
the total size and these \(o\) counts shows that the
probability of violating their bounds, using
\(x^{-1+\epsilon}\delta n/2\) as the coset tolerance, is
\[
 O\!\left((\delta n)^{-1}
             +\frac{o x^{2-2\epsilon}}{\delta n}\right)=o(1).
\]
Intersect the events. This proves both requirements on the
same \(T\) and allows all translates used in the Fourier
projection to be taken afterwards.
\end{proof}

For the
stated prime weights and \(p<X'\),
\(\sum_p\log(p)/p\le\log X'(1+\log X')\).
Thus \eqref{eq:bad-cycles} costs at most
\(|T|x^{-2+\epsilon}\) for sufficiently large \(x\).
Only this weighted count of failed cycles enters the global
estimate, even if an individual prime has several failed
cycles.

\paragraph{The large-weight set inside the identity coset.}
For the larger weights choose the labels \(s\) in
\(\mathbb Z/n_0\mathbb Z\) via \eqref{eq:dyadic-coordinate-orders}.
A label \(s\) denotes the unique member of \(A_0\) whose real
cyclotomic coordinate is \(t_0 s\bmod D\). Write \(T_2\) for the set
of these members. At fixed
\(\epsilon\) and bounded \(o\), let \(2\le L\ll_\epsilon\log x\)
bound the lengths of the arithmetic progressions being tested.

\begin{lemma}[Progression bounds for the second set]\label{lem:second-set-progressions}
Fix \(0<\epsilon<1/52\), the group of
Definition~\ref{hyp:frobenius-top-cycle} with bounded index \(o\),
and an integer \(2\le L\ll_\epsilon\log x\). Put \(n_0=n/o\).
There is an absolute \(C\) and a subset \(T_2\) of the cyclic
kernel, with coordinate labels \(0<s<n_0/(4L)\) in
\eqref{eq:dyadic-coordinate-orders}, such that
\[
 |T_2|\ge \frac{n_0}{L}\exp(-C\sqrt{\log n_0\log L}),
 \qquad x^{-\epsilon}n\le |T_2|\le\frac{n_0}{4L}
\]
for all sufficiently large \(x\). Moreover every translate
of \(T_2\) in \(\mathbb Z/n_0\mathbb Z\) contains at most two
members of any progression of length \(L\) whose step has
order greater than \(4L\). An entire coset of a cyclic
subgroup of order at most \(4L\) meets that translate in at
most one point.
\end{lemma}

\begin{proof}
\proofstep{Encode lattice points on one sphere.}
Put \(X=n_0/(8L)\),
\(t=\lfloor\sqrt{\log X/\log(4L)}\rfloor\),
\(B=\lfloor X^{1/t}\rfloor\) and
\(D=\lfloor B/(4L)\rfloor\). Eventually \(t\ge2\),
\(D\ge B/(8L)\) uniformly in the stated range.

In
\(\{0,\ldots,D-1\}^t\), some sphere centered at the
origin contains at least \(D^t/(1+tD^2)\) points, since
the squared norms are integers. Encode its points \(z\)
by \(s=1+\sum_{a=0}^{t-1}z_aB^a\). This is injective,
and \(s\le B^t\le X\). The size bound follows from
\(B^t\ge X/2^t\), \(\log B\le (\log X)/t\) and
\[
 t\log(16L)+2(\log X)/t+\log(1+t)
                      =O(\sqrt{\log X\log(4L)}).
\]
Since \(L\ll_\epsilon\log x\) and \(n_0\asymp x^h/o\),
the logarithmic density loss is \(o(\log x)\) uniformly.
This proves the second lower bound.

Suppose three distinct points satisfy
\(j(s_2-s_1)=i(s_3-s_1)\) with distinct nonzero integers
\(|i|,|j|\le L\). The contribution at each digit has
absolute value at most \(2L(D-1)<B\), so reduction modulo
\(B\) and induction on digit positions give
\(j(z_2-z_1)=i(z_3-z_1)\) as vectors. In fact, put \(v=z_2-z_1\ne0\). On the line \(z_1+tv\) the common
squared norm condition is
\[
 \|z_1+tv\|^2-\|z_1\|^2
   =t\bigl(2\langle z_1,v\rangle+t\|v\|^2\bigr)=0.
\]
Its roots are \(0\) and \(1\), whereas the third point would give
\(t=j/i\notin\{0,1\}\). This proves
the exclusion for integer equalities.

If the order of the
step is greater than \(4L\), three visits to a translate
within a progression of length \(L\) would give three
distinct integers \(s_a\) in the chosen range and the
same equality modulo \(n_0\) by eliminating the step.
The absolute value of the difference of the two sides is
less than \(n_0\), yielding the forbidden integer equality.
For a step of order \(b\le4L\) two distinct visits would force the incompatible conditions
\[
 n_0\mid b(s_2-s_1),\qquad 0<|b(s_2-s_1)|<n_0.
\]
This proves both intersection
bounds. The construction depends on \(n_0,L\); it works
for every translate and every step simultaneously.

\end{proof}

\paragraph{Short orders and prime powers near the cutoff.}
\begin{auxlemma}[Short orders in the real dyadic quotient]
\label{lem:short-real-dyadic-order}
Let \(e\ge0\), let \(n\) be odd, and suppose its image in
\((\mathbb Z/2^{e+2}\mathbb Z)^\times/\{\pm1\}\) has order \(2^v\).
Then
\[
 2^{e+2-v}\mid n-1\qquad\text{or}\qquad2^{e+2-v}\mid n+1.
\]
\end{auxlemma}

\begin{proof}
Choose \(z=\pm n\) with \(z\equiv1\pmod4\). Its \(2^v\)-th power is
\(\pm1\pmod{2^{e+2}}\), and reduction modulo four forces the plus sign.
The assertion is immediate for \(z=1\). Otherwise every
\(z^{2^j}+1\) has 2-adic valuation one, giving
\[
 e+2\le v_2(z^{2^v}-1)
      =v_2(z-1)+\sum_{j=0}^{v-1}v_2(z^{2^j}+1)=v_2(z-1)+v .
\]
\end{proof}

Two other exceptional sets in the finite calculation admit
bounds without averaging the top positions. The host contains
the real cyclotomic field of conductor \(2^{e+2}\) with
\(2^e\asymp_\psi x^h\). The image of the Frobenius step
at an unramified odd prime \(p\) on this factor is the class
of \(p\bmod 2^{e+2}\) modulo signs. If the order of the
step is below \(x^{\epsilon/4}\), that image has order
\(2^v\le x^{\epsilon/4}\). Lemma~\ref{lem:short-real-dyadic-order} gives
\(p\equiv\pm1\pmod{2^{e+2-v}}\). All the exceptional primes therefore satisfy
\[
 p\equiv\pm1\pmod B,\qquad B=2^{e+2-\lfloor\log_2(x^{\epsilon/4})\rfloor}\gg_\psi x^{h-\epsilon/4}.
\]

Sum over integers \(jB\pm1>1\) and use \(jB-1\ge jB/2\).
For \(X''=xnX'\) this gives
\[
 \sum_{\substack{p<X''\\ {\rm ord}(a_p)<x^{\epsilon/4}}}
              \frac{\log p}{p}
 \ll \frac{\log X''(1+\log X'')}{B}
 \ll_\psi x^{-h+11\epsilon/12}\le x^{-h+\epsilon}
\]
for all sufficiently large scales, also after restricting the sum to the prescribed unramified primes.

There is also an exceptional band when some positive power
of \(p\) is in \([X'(1-\rho),X'(1+\rho)]\) with
\(p_0=k^2x^{26}\le p<X'\). Here the larger width used
for all of the choices satisfies
\[
 0\le\rho\ll
  \frac{(n\delta x^\epsilon+\log x)N_1^2\log a}{k},
 \qquad a=x\exp(x^\epsilon).
\]
The two terms arise from the widths of the weight intervals,
with common factor \(C N_1^2\log(a)/k\) for a fixed constant
\(C\). The rank estimates give
\[
 k\ge x^{9/2+10\epsilon-\epsilon/100},\qquad
 N_1=O_\epsilon(\log x),\qquad n\ll x^h .
\]
Thus \(\rho\le x^{-2}\) for sufficiently large \(x\).

For fixed power \(v\)
the ratio of the greatest and least eligible primes is
at most \(((1+\rho)/(1-\rho))^{1/v}\le1+3\rho\).
If the least prime is \(p_*\), there are at most
\(1+3\rho p_*\) integers in this interval. Each contributes
at most \((\log X')/p_*\). There are at most
\(1+2\log X'\) possible powers since \(p\ge2\).
Thus the band contributes
\[
  \sum_p (\log p)/p\ll (1+\log X')^2(\rho+p_0^{-1})
                    \le x^{-1}
\]
for sufficiently large \(x\), absorbing the implied constant.
Both estimates have thresholds uniform in the particular
host of the fixed target family.

\paragraph{Choosing the top weights after the positions.}
\label{par:top-weight-selection}
The cutting off of exceptional local subspaces involves a
second simultaneous choice. At each \(t\in T\), choose two
integer weights independently and uniformly in the intervals
\[
  [a,a+a/(kx^{13})],\qquad
  [b',b'+2n\delta x^\epsilon\Lambda b'/k],
  \qquad b'=a/X'.
\]
Set the weights to zero off \(T\). Both intervals contain
integers for large \(x\), since \(a,b'\) grow faster
than every fixed power of \(x\) and \(k\) has polynomial
growth. The same independent uniform choices on a single
cycle give exactly the marginal law of the global weights.
For a candidate \(p\), exclude the short orders and the
band above. Write \(N_v\) for the number of selected
positions on a good cycle \(v\). The local interpolation
estimate required here is a nonnegative fraction \(\kappa_v\)
of exceptional directions with
\[
                 \mathbb E(\kappa_v)\le N_v x^{-12}.
\]

For \(p<p_0\) its numerical bound has the form
\(\min(1,C N_v r^2 k x^{-26})\), which suffices since
\(k\ll x^{9/2+10\epsilon}\). For the complementary range \(p\ge p_0\), the local estimate is
\[
 \text{local expectation}\ \le C N_vN_1 x^{-13}.
\] Section~\ref{sec:interpolation} identifies
the tested integral rank and proves the determinant estimate
for varying a weight on a progression.
Section~\ref{sec:edge-transfer} obtains the large-prime mean
on the full edge lattices from the two-letter tensor edge
test. Sections~\ref{sec:segment-window}--\ref{sec:block-pairing} prove
the segment mean, the retained rank and the symbol pairing. Sections~\ref{sec:column-lifts} and \ref{sec:symbol-correction} construct
the required simultaneous lifts at the prescribed order and conductor.
Section~\ref{sec:host-edge-tests} verifies the sizes in both
ranges and extends the tests to the specified place of the host.
Section~\ref{sec:plane-lifts} proves the smaller-prime bound by
counting source vectors and constructing a target covector for
every primitive combination at a single plane.

We apply these bounds to select the weights simultaneously.
Let \(m\) be the length of the lift chain and put
\[
 \Sigma_p=m\sum_{v\ {\rm good}}\kappa_v.
\]

The cycle decomposition bounds its mean, and Markov's inequality bounds the rejected primes:
\[
 \mathbb E\Sigma_p\le m|T|x^{-12},\qquad \Pr(\Sigma_p>x^{-\epsilon})\le x^\epsilon\mathbb E\Sigma_p.
\]
Since
\(\mathbf1_{\Sigma_p>x^{-\epsilon}}\le x^\epsilon\Sigma_p\),
the expectation of
\[
   \sum_{p\ {\rm candidate}} \frac{\log p}{p}
       (2\mathbf1_{\Sigma_p>x^{-\epsilon}}+4\Sigma_p)
\]
is at most
\((2x^\epsilon+4)m|T|x^{-12}\log X'(1+\log X')\).
With \(m\le x^{1/2+\epsilon/2}\) and \(|T|\le 2\delta n\), the size bounds give
\[
 (2x^\epsilon+4)m|T|x^{-12}\log X'(1+\log X')\ll x^{-19/2+11\epsilon}\le x^{-9}
\]
eventually. Some single pair of
weight vectors attains this bound. At a candidate where
\(\Sigma_p\le x^{-\epsilon}\) use every good cycle, and
at the others use none.

By Section~\ref{subsec:column-orbit-length} the local length
error on the retained cycles is bounded by \(2\Sigma_p\max_v S_v\)
with \(S_v\le2mN_v a/p\). Its total logarithmic cost is therefore
at most \(m|T|a x^{-9}\). The missing prime sum at rejected
candidates is controlled by the same choice. The smoothed
positive part at one top in one lift is at most
\[
 2a\sum_{j\ge1,\ j\text{ odd}}p^{-j}
   =\frac{2a}{p(1-p^{-2})}\le\frac{8a}{3p}\le\frac{3a}p .
\] These estimates charge a prime for its actual
weight and hold simultaneously over every subset of the
candidate primes.


\subsection{The cohomology input to the window}

The bounds for the selected positions do not account for division in the integral lattices. A second input is the length of a global Galois cohomology group with vanishing localizations. We need that length over the field already chosen and for all twists below the eventual rank.

\begin{definition}[Strict localization length]\label{def:strict-length}
Let \(E\) be the auxiliary abelian field chosen at the scale in
question, let \(\ell\) be a rational prime, and let \(j\ge1\)
be an integer. The prime \(\ell\) varies independently of the
auxiliary primes \(q,d\). Let \(S_\ell\) consist of the places
of \(E\) above \(\ell\) and the infinite places. Write
\(G_{E,S_\ell}\) for the Galois group of the maximal extension
of \(E\) unramified outside \(S_\ell\).
Use the Tate-twisted modules
\(A=(\mathbb Q_\ell/\mathbb Z_\ell)(-j)\) and
\(T=\mathbb Z_\ell(j+1)\). Write \(H^i(E,A)\) and \(H^i(E,T)\)
for cohomology of \(G_{E,S_\ell}\) with these coefficients.
Write \(\operatorname{len}\) for \(\mathbb Z_\ell\)-length. The strict group and its length over \(\mathbb Z_\ell\) are
\[
 \operatorname{Sha}^1(E,A)=
 \ker\!\left(H^1(G_{E,S_\ell},A)\longrightarrow
        \prod_{v\mid\ell}H^1(E_v,A)\right),\qquad
 V(E,\ell,j)=\operatorname{length}_{\mathbb Z_\ell}\operatorname{Sha}^1(E,A).
\]
\end{definition}
The places at infinity are included in the allowed ramification set \(S_\ell\).
For this abelian host, \(q\equiv3\pmod8\) gives
\(\bigl(\frac{-1}{q}\bigr)=-1\), so complex conjugation acts nontrivially on \(E\). All its archimedean completions are therefore complex. Their degree-one local cohomology is zero, also when \(\ell=2\); the displayed kernel agrees with the localization kernel of Appendix~\ref{sec:cohomology}.

\begin{lemma}[The Poitou--Tate length identity]\label{lem:pt-length-identity}
Use \(A,T,S_\ell\) of Definition~\ref{def:strict-length}.
For odd \(\ell\), or when \(E\) has no real places, Poitou--Tate duality \cite{milne-duality} gives the following identity when the
lengths are finite. In the sum, \(v\) runs over the places of
\(E\) above \(\ell\), and \(E_v\) is the completion at \(v\).
\[
 \operatorname{len}H^2(E,T)+\operatorname{len}H^0(E,A)
 =
 \operatorname{len}\operatorname{Sha}^1(E,A)+
 \sum_{v\mid\ell}\operatorname{len}H^0(E_v,A).
\]
\end{lemma}
\begin{proof}
Write \(D^\vee\) for the Pontryagin dual. The final terms of
Poitou--Tate and local Tate duality give the exact sequence
\[
 0\longrightarrow\operatorname{Sha}^1(E,A)^\vee
 \longrightarrow H^2(E,T)\longrightarrow
 \bigoplus_{v\mid\ell}H^0(E_v,A)^\vee
 \longrightarrow H^0(E,A)^\vee\longrightarrow0 .
\]
In the stated range the infinite places contribute no modified
local terms. Additivity of length and invariance under finite
Pontryagin duality give the identity. When \(\ell\) is odd it
also splits under the two idempotents \((1\pm c)/2\) for complex
conjugation \(c\). The real \(2\)-primary terms are treated separately.
\end{proof}

\begin{paramdefinition}[Cyclotomic tower for strict coinvariants]\label{hyp:strict-tower}
Write \(\mu_r\) for the group of \(r\)-th roots of unity, for
integers \(r\ge1\), and let \(\chi_\ell\) be the \(\ell\)-adic
cyclotomic character. Choose a base field \(L_0\) containing
\(\mu_\ell\) for odd \(\ell\) or \(\mu_4\) for \(\ell=2\).
Let \(L_\infty\) be its cyclotomic \(\mathbb Z_\ell\)-extension,
\(\gamma\) a topological generator of its Galois group, and
\(u=\chi_\ell(\gamma)\). Let \(X_\ell\) be the Galois group over \(L_\infty\) of
the maximal abelian pro-\(\ell\) extension of \(L_\infty\)
unramified away from \(\ell\) in which every place above \(\ell\)
splits completely. For a nonnegative integer \(m\), let \(L_m\)
be the layer of degree \(\ell^m\).
\end{paramdefinition}

\begin{lemma}[Coinvariants for the strict kernel]\label{lem:tower-coinvariants}
For the tower of Definition~\ref{hyp:strict-tower} and any
\(j\ge1\), the Pontryagin dual of the strict kernel over \(L_m\) is
\[
 X_\ell/\bigl((u^j\gamma)^{\ell^m}-1\bigr)X_\ell .
\]
\end{lemma}
\begin{proof}
Write \(\Gamma_m=\langle\gamma^{\ell^m}\rangle_{\rm top}\).
The module \(A\) is trivial on the absolute Galois group over
\(L_\infty\). For the quotient acting on \(A\),
\[
 H^1(\Gamma_m,A)=A/(u^{-j\ell^m}-1)A=0,\qquad
 H^2(\Gamma_m,A)=0 .
\]
The first equality uses divisibility and \(u^{-j\ell^m}\ne1\);
the second uses \(\operatorname{cd}_\ell\Gamma_m=1\).
The same calculation with an open decomposition subgroup gives
the local inflation--restriction isomorphism. A class trivial
at all places above \(\ell\) thus corresponds to a continuous
equivariant homomorphism \(f:X_\ell\to A\). Writing the action
on \(X_\ell\) additively, equivariance is exactly
\[
 f(\gamma^{\ell^m}h)=u^{-j\ell^m}f(h)
 \quad\Longleftrightarrow\quad
 f\bigl(((u^j\gamma)^{\ell^m}-1)h\bigr)=0
 \qquad(h\in X_\ell).
\]
These homomorphisms form the Pontryagin dual of the displayed
compact quotient, after identifying the underlying group of
\(A\) with \(\mathbb Q_\ell/\mathbb Z_\ell\).
\end{proof}

Lemma~\ref{lem:tower-length-centres} uses this presentation to
bound the lengths outside finitely many exceptional twists.

Let \(z\) be the original scale for the cohomology estimates,
let \(C_t=401\), and put \(T_{C_t}(z)=\lfloor z^{C_t}\rfloor\).
This cutoff is the
upper endpoint of the strict range \(1\le j<T_{C_t}(z)\).
The auxiliary field has a \(2\)-power cyclotomic factor.
Its exponent \(e\) satisfies
\(z^{2+8\epsilon}\le2^e\le2z^{2+8\epsilon}\).
\begin{lemma}[Conductor bound at the strict cutoff]
Let \(C_t=401\), \(T_{C_t}(z)=\lfloor z^{C_t}\rfloor\), and take
the auxiliary field \(E\subseteq\mathbb Q(\mu_N)\) with
\[
 N\le2^{e+2}\operatorname{cond}(\psi)qd,\qquad
 2^e\le2z^{2+8\epsilon},\qquad \log(qd)\le z^{\epsilon/4}.
\]
Then \([E:\mathbb Q]\le N\) and
\(\log(NT_{C_t}(z))\le2z^{\epsilon/4}\) for sufficiently large
\(z\) depending on \(\psi,\epsilon,C_t\).
\end{lemma}
\begin{proof}
The fixed field lies in the cyclotomic field of conductor
\(N\le2^{e+2}m_\psi qd\), where \(m_\psi=\operatorname{cond}(\psi)\).
In particular \([E:\mathbb Q]\le\varphi(N)\le N\). The upper
bound on \(2^e\) and the range of \(T_{C_t}\) give
\[
 \log(NT_{C_t}(z))
 \le \log(qd)+\log(8m_\psi)+(2+8\epsilon+C_t)\log z .
\]
The first term is at most \(z^{\epsilon/4}\), and the remaining
terms together are at most \(z^{\epsilon/4}\) past the claimed threshold.
\end{proof}

\begin{proposition}[A bound for strict cohomology at the original scale]
\label{prop:original-window-budget}
Fix the primitive character \(\psi\), \(0<\epsilon<1/52\), and the
dyadic index function \(e=e(z)\) satisfying
\[
                z^{2+8\epsilon}\le2^{e(z)}\le2z^{2+8\epsilon}.
\]
There is \(J\ge1\) such that, for all sufficiently large \(X\) and a
prime pair satisfying Proposition~\ref{prop:common-host-prime-pair} on
\([X^{1/J},X]\), there is an interval
\([W,W^2]\subseteq[X^{1/J},X]\) with the following property.
At every \(z\) in this interval all strict lengths in the range
\(1\le j<T_{C_t}(z)\) are finite, and for every finite set \(P\)
of rational primes
\begin{equation}\label{eq:cohomwin}
 \sum_{\ell\in P}(\log\ell)
       \sum_{1\le j<T_{C_t}(z)} V(E,\ell,j)
       \le \exp(z^\epsilon/2).
\end{equation}
\end{proposition}
\begin{proof}
The window is constructed in Lemma~\ref{lem:exceptional-twist-window}.
The position of the window may depend on the prime pair. Its exponent
\(J\) is fixed before that pair is chosen.
For the bounds entering this estimate, put
\(B=\exp(z^{4\epsilon/5})\) and let \(K_0,C>0\) be fixed
constants. The three contributions are
\begin{center}
\begin{tabular}{@{}lll@{}}
Contribution & Upper bound & Proof \\
\hline
Prime \(2\) & \(B\log 2\) & Prop.~\ref{prop:dyadic-window-total} \\
Odd primes, negative component & \(K_0(NT_{C_t}(z))^4\) &
  \S\ref{subsec:odd-prime-negative} \\
Odd primes, positive component & \((NT_{C_t}(z))^C+CB\) &
  \S\ref{subsec:odd-prime-sum}
\end{tabular}
\end{center}
Since \(\log(NT_{C_t}(z))\le2z^{\epsilon/4}\), the three contributions sum to at most
\[
 B\log2+ K_0(NT_{C_t}(z))^4+(NT_{C_t}(z))^C+CB\ll \exp(Dz^{4\epsilon/5})
\]
for a fixed \(D>0\). The inequality
\(4\epsilon/5<\epsilon\) gives \eqref{eq:cohomwin} for all
sufficiently large \(z\) in the chosen interval. The lower
threshold is uniform in the field, the twist indices and \(P\).
\end{proof}

\begin{lemma}[Rescaling the interval and its strict cutoff]\label{lem:rescale-host-budget}
Let \(c\ge1\) be the constant in the frame order \(n=c2^e\),
set \(\kappa_{\rm sc}=c^{1/(2+8\epsilon)}\), the positive rescaling constant, and take a window of
Proposition~\ref{prop:original-window-budget} with
\(W\ge\max(\kappa_{\rm sc}^3,\kappa_{\rm sc}^{400})\). Set \(Y=\kappa_{\rm sc}W\).
For each \(Y\le x\le Y^{3/2}\), use the field and dyadic index
at \(x'=x/\kappa_{\rm sc}\). Then both scales lie in \([W,W^2]\), and
\[
 x^{2+8\epsilon}\le n\le 2x^{2+8\epsilon},\qquad
 \lfloor x^{400}\rfloor\le \lfloor(x')^{401}\rfloor .
\]
For every finite set of primes the budget in \eqref{eq:cohomwin},
restricted to \(1\le j<\lfloor x^{400}\rfloor\), is at most
\(\sqrt b\), where \(b=\exp(x^\epsilon)\).
\end{lemma}

\begin{proof}
The index \(t_0\) in \eqref{eq:dyadic-coordinate-orders} is the degree
of \(K_\psi\cap C_e\) over \(\mathbb Q\). Indeed restriction
identifies the subgroup fixing \(K_\psi\) with
\(\operatorname{Gal}(C_e/(K_\psi\cap C_e))\). The intersections
increase with \(e\) inside the fixed finite extension \(K_\psi\),
so they stabilize. Thus eventually
\[
           n=oD/t_0=c2^e,\qquad c=o/t_0\ge1 .
\]
For \(\psi=1\) this gives \(c=1\). Lemma~\ref{lem:host-rescaling}
places both scales in \([W,W^2]\).
Multiplying the bounds for \(2^{e(x')}\) by \(c=\kappa_{\rm sc}^{2+8\epsilon}\)
gives the asserted frame order. As \(x'\ge \kappa_{\rm sc}^{400}\),
\((x')^{401}\ge \kappa_{\rm sc}^{400}(x')^{400}=x^{400}\); this includes equality.
All summands of \eqref{eq:cohomwin} are nonnegative, so restricting
its strict cutoff costs nothing, and \(\exp((x')^\epsilon/2)\le\sqrt b\).

\end{proof}

\begin{lemma}[The selected rank lies in the cohomology range]
\label{lem:rank-window}
Fix \(0<\epsilon<1/52\). Let \(x\ge 2^{100}\),
\(\delta_1=x^{-\epsilon/8}\) and \(L=x^\epsilon\). Suppose
the frame order \(n\) and the positive even ranks \(m,k\)
satisfy the family bounds
\[
 x^{2+8\epsilon}\le n\le 2x^{2+8\epsilon},\qquad
 \sqrt{x\delta_1}\le m\le 2\sqrt{x\delta_1},\qquad
 x^{9/2+10\epsilon}\le k\le 2x^{9/2+10\epsilon},\qquad m+3\le k.
\]
Set \(f=2n\) and \(r=k+m-1\). The rank datum of the host
chooses the positive integer \(t\) by the equality
\[
 T_0=\lceil x^{377}\rceil,\qquad
 N_W=8\left\lceil\frac{10r}{8}
           \left(f+\frac{f^2(k-1)^3L^2}{m}\right)\right\rceil,
 \qquad r+t=8fr+4+T_0N_W .
\]
Then \(t\le\lfloor x^{400}\rfloor\).
\end{lemma}

\begin{proof}
Put \(y=x^{1/100}\ge2\). The exponents in the family
bounds give
\[
 f\le4y^{216},\quad k\le2y^{470},\quad m\ge y^{49},\quad
 L^2\le y^4,\quad r\le2k\le4y^{470}.
\]
Here \(2+8\epsilon\le216/100\), \(9/2+10\epsilon\le470/100\),
\((1-\epsilon/8)/2\ge49/100\), and \(2\epsilon\le4/100\).
Thus the parenthesis in \(N_W\) is at most
\[
 4y^{216}+\frac{(4y^{216})^2(2y^{470})^3y^4}{y^{49}}
 \le 132y^{1797}.
\]
The ceiling contributes at most \(8\) to \(N_W\), so
\(N_W\le5288y^{2267}\); also \(T_0\le2y^{37700}\).
It follows that
\[
 t\le128y^{686}+4+10576y^{39967}
   \le10708y^{39967}\le y^{40000}=x^{400},
\]
since \(y^{33}\ge 2^{33}\ge 10708\). The rank \(t\) is integral.
\end{proof}
\begin{construction}
\begin{tabular}{@{}l@{\hspace{1.2em}}l@{}}
\toprule
\textit{Datum fixed in the construction} & \textit{Range of the twist index} \\
\midrule
Original scale \(x'\) & \(1\le a<T_{C_t}(x')\) \\
Rescaled point \(x=\kappa_{\rm sc}x'\) & \(\lfloor x^{400}\rfloor\le \lfloor(x')^{401}\rfloor\le T_{C_t}(x')\) \\
Selected integer \(t\le\lfloor x^{400}\rfloor\) & \(1\le a<t\) \\
\bottomrule
\end{tabular}
\caption{The increasing scale and the restriction to a smaller summation range in Lemmas~\ref{lem:rescale-host-budget} and~\ref{lem:rank-window}. The field in the budget is the original field at \(x'\).}\label{fig:analytic-cutoff}
\end{construction}

\begin{proposition}[A bound for strict cohomology at the selected rank]\label{prop:actual-rank-budget}
On the eventual windows of Lemma~\ref{lem:rescale-host-budget},
the frame at \(x\) may be chosen with the ranks of
Lemma~\ref{lem:rank-window}. With the original field \(E\) at
\(x'=x/\kappa_{\rm sc}\) and the actual chosen rank \(t\), put
\[
 V_\ell=\sum_{a=1}^{t-1}
   \operatorname{len}_{\Z_\ell}\operatorname{Sha}^1
     (E,(\Q_\ell/\Z_\ell)(-a)).
\]
Each \(V_\ell\) is finite and
\(\sum_{\ell\in P}(\log\ell)V_\ell\le\sqrt b\)
for every finite set \(P\) of primes.
\end{proposition}

\begin{proof}
Identify the field at \(x\) using the fixed kernel conditions
for \(\psi,q,d,e(x')\). To extend its order-\(n\) frame to
Lemma~\ref{lem:rank-window}, use \(2\lceil A/2\rceil\) at the
lower endpoints \(A\) of the \(m,k\) intervals. This even
integer lies in \([A,2A]\) for \(A\ge2\), as holds for \(x\ge16\).
The rank tests follow in this order:
\[
 \begin{aligned}
 m\le2x,\quad 2x^2\le f\le4x^3,\quad k\ge x^4
   &\ \Longrightarrow\ m+3\le k,\quad m\le f\le k,\\
 t\ge7r+4 &\ \Longrightarrow\ t\ge r+5,\quad t>3m,\quad
                                 m^2\le k+m(t+1).
 \end{aligned}
\]
The formula for \(t\) gives the second lower bound using \(f\ge1\)
and gives odd parity using odd \(r\) and \(8\mid N_W\). Adjoining
type primes, local choices and lift data preserves the field and
rank. Choose the interval with lower endpoint exceeding \(2^{100}\).
The preceding lemmas give precisely the range needed in \eqref{eq:cohomwin},
\[
             1\le a<t\le\lfloor x^{400}\rfloor
                          \le\lfloor(x')^{401}\rfloor .
\]

At \(2\), let \(K_\psi\) be the fixed
field of the kernel of \(\psi\) in the cyclotomic field of its
conductor. Set \(L_0=K_\psi(i,\sqrt{-q},\sqrt{-d})\) and let
\(L_m=E(i)\) be the layer of degree \(2^m\) over \(L_0\) in
its cyclotomic \(\mathbb Z_2\)-extension. The restriction
estimate \(V(E,2,j)\le 1+V(L_m,2,j)\) holds for each
\(j\ge1\), as proved in Lemma~\ref{lem:dyadic-host}.
At odd primes both eigenspaces for complex conjugation
enter the sum. This proves the bound for the actual field
and rank at every selected point of the interval.
\label{cor:family-window-1}

\end{proof}

\begin{lemma}[Character components of strict cohomology]\label{lem:strict-character-budget}
Use the field \(E\), translate \(x\) and rank \(t\) of
Proposition~\ref{prop:actual-rank-budget}. At each odd prime \(\ell\) and each \(1\le a<t\), the positive component is indexed by
\[
 \{\xi\in\operatorname{Gal}(E/\mathbb Q)^\wedge:\xi(-1)=(-1)^a\}.
\]
The localization kernel over \(\mathbb Q\) indexed by \(\xi\) has coefficient character \(\epsilon_\ell^{-a}\xi^{-1}\).
The character estimates in Sections~\ref{subsec:odd-prime-trace}
and~\ref{subsec:odd-prime-pointwise}, with the exceptional cases
specified in the proof, give the positive part of the budget in
\eqref{eq:cohomwin} on the strict range.
\end{lemma}
\begin{proof}
Write \(x'=x/\kappa_{\rm sc}\) for the scale used to choose the field, and
\(T'=\lfloor(x')^{401}\rfloor\). Let \(a\) be the positive Tate
index, so that the scalar module has twist \(-a\).
The positive sign at an odd prime \(\ell\) is the eigenspace
of complex conjugation with eigenvalue \((-1)^a\) on the
twisted group. Put \(\mathcal N=\operatorname{lcm}(2^{e+2},
\operatorname{cond}\psi)qd\) and let \(H\) be the subgroup
fixing \(E\) inside \(\mathbb Q(\mu_{\mathcal N})\). Write
\(\epsilon_\ell\) for the \(\ell\)-adic cyclotomic character.
Decomposition by the quotient cutting out \(E\) gives
\[
 \xi\colon (\mathbb Z/\mathcal N\mathbb Z)^\times/H\longrightarrow
 \overline{\mathbb Q}_\ell^\times,\qquad
 \xi(-1)=(-1)^a,\qquad
 \text{coefficient character}=\epsilon_\ell^{-a}\xi^{-1}.
\]
Only the field and the twist enter these conditions. They therefore
give the same sum after a change of automorphic type at a level prime.

Fix the primes \(p_K,p_1\) used in the auxiliary construction
for \(\psi\) in Section~\ref{sec:odd-prime-cases}. The first
specifies the imaginary quadratic field \(K_*=\mathbb Q(\sqrt{-p_K})\)
of discriminant \(-p_K\). Write \(h_*\) for its class number.
The finite exceptional set is
\[
 B_\psi=\{3\}\cup\{\ell>2:\ell\text{ prime},\
                \ell\mid \operatorname{ord}(\psi)\operatorname{cond}(\psi)p_Kh_*p_1\}.
\]
For \(\ell\notin B_\psi\) the character quotient has order
prime to \(\ell\). If moreover \(\ell-1\mid2a\), the principal term vanishes. Neither auxiliary prime occurs as a residual prime in this case:
\[
 1\le a<T',\quad \ell-1\mid2a\quad\Longrightarrow\quad
 \ell\le 2a+1\le2T'<\min(q,d).
\]

At \(\ell\in B_\psi\) the character comparison has polynomial error. For \(a\ge2\), the parity test first removes the \(\ell\)-conductor factor and takes the prime-to-\(\ell\) part. In terms of the factors by conductor prime and their prime-to-\(\ell\) components,
\[
 \xi=\prod_{p\mid\operatorname{cond}\xi}\xi_p,\qquad
 \xi_P=\prod_{\substack{p\mid\operatorname{cond}\xi\\p\ne\ell}}(\xi_p)_{\ell'} .
\]
Let \(c_2=c_2(\xi)\) be the exponent of the dyadic conductor. With \(n_0\) as in \eqref{eq:ordinary-dyadic-exponents}, Section~\ref{sec:odd-prime-cases} uses the threshold
\[
 R_1=R_1(\ell,\xi)=16\operatorname{ord}(\psi)
            (\ell\operatorname{cond}\psi)^2\,4^{n_0}.
\]
\begin{center}
\begin{tabular}{@{}p{0.52\linewidth}p{0.43\linewidth}@{}}
Hypothesis & Estimate and proof \\
\hline
\(\ell\notin B_\psi\), \(\ell-1\nmid2a\) &
  Rank-three trace sum over primes and characters (\S\ref{subsec:odd-prime-trace}) \\
\(\ell\notin B_\psi\), \(\ell-1\mid2a\), \(2\mid\operatorname{ord}\xi\), or
\(\ell\in B_\psi\), \(a\ge2\), \(2\mid\operatorname{ord}\xi_P\) &
  Even-order estimate (\S\ref{subsec:odd-prime-pointwise}) \\
\(\ell\notin B_\psi\), \(\ell-1\mid2a\), \(2\nmid\operatorname{ord}\xi\) &
  Rank-four construction (\S\ref{sec:count-three-target}) \\
\(\ell\in B_\psi\), \(a\ge2\), \(2\nmid\operatorname{ord}\xi_P\) &
  Fixed cyclotomic layer (\S\ref{subsec:odd-prime-pointwise}) \\
\(\ell\in B_\psi\), \(a=1\), \(2^{c_2}\ge R_1\) &
  Gaussian estimate (\S\ref{subsec:odd-prime-pointwise}) \\
\(\ell\in B_\psi\), \(a=1\), \(2^{c_2}<R_1\) &
  Bounded-degree height (\S\ref{sec:cohom-height})
\end{tabular}
\end{center}
Descent gives these bounds over the original character fields. \label{cor:family-window-2}
For the fixed layer, conductor divisibility bounds the primitive characters by
\[
 \#\{\xi:\operatorname{cond}\xi\mid\operatorname{cond}\psi\}
 \le\varphi(\operatorname{cond}\psi).
\]
After summation over \(a<T'\), each character charges the reference tower once,
and its coefficient after summing over the exceptional primes
is at most
\[
 \varphi(\operatorname{cond}\psi) \sum_{\ell\in B_\psi}\log\ell.
\]

All the other contributions on the positive component are bounded by a
fixed power of \(\mathcal NT'\) on the counting estimates.
Together with the negative component and the full length at
\(2\), these estimates give the weighted bound \(\sqrt b\).
\end{proof}

\begin{lemma}[The cohomology bound after restriction to the good types]\label{lem:strict-budget-after-cut}
At each retained translate and localization \(i\) under consideration,
let \(S_i\) be the support of Definition~\ref{def:good-projector-support}.
Write \(\operatorname{pr}_{S_i}\) for the rational projection onto the sum of
the eigenspaces indexed by \(S_i\). Use the image lattice
\(\operatorname{pr}_{S_i}(M_i)\) of Proposition~\ref{prop:cut-selector}.
The scalar cohomology
lengths are computed over the same field \(E\) and at the same rank \(t\)
as in Proposition~\ref{prop:actual-rank-budget}.
At that translate the lengths satisfy
\[
 V_\ell\ge0,\qquad \sum_\ell(\log\ell)V_\ell\le\sqrt b.
\]
\end{lemma}
\begin{proof}
The realization at a translate preserves the field \(E\), the
rank \(t\), and the indices of the larger weights.
Changing the algebraic closure transports the coefficient module.
Changing the prime over a place conjugates the restriction map.
Neither operation changes the length of its localization kernel.
The restriction to good types concerns the Hecke eigenspaces.
The coefficient module for the scalar lengths remains
\((\Q_\ell/\Z_\ell)(-a)\) over \(E\).

By Lemmas~\ref{lem:rescale-host-budget} and~\ref{lem:rank-window},
the cutoff at the selected rank satisfies
\[
 \{a\in\mathbb Z:1\le a<t\}\subseteq\{a\in\mathbb Z:1\le a<T'\}
                   \qquad(t\le T').
\]
This includes \(t=T'\).
All the lengths are finite, so they may be summed as ordinary
nonnegative numbers. Proposition~\ref{prop:actual-rank-budget}
bounds every finite subsum of \(\sum_\ell(\log\ell)V_\ell\).
Taking the supremum over finite sets of primes gives the full sum at this translate, with the field and rank still fixed.
In the marked row and column estimates the good multiplicities at
\((j,w)\) enter separately, with weights \(e_w\omega_w\).
\label{cor:family-window-3}
\end{proof}

\subsection{The integral type cut}\label{subsec:selected-cut}
\hypertarget{source-selected-cut}{}

A scalar cohomology length is counted with the multiplicity of the corresponding good character. To determine which characters contribute, we must realize the restriction to good types on the integral Hecke module at each of the two distinguished primes.

\begin{definition}[Type group and theta support at a localization]\label{def:localized-type-support}
Fix a point of a chosen chain, with its two distinguished depth-zero
primes \(g_0,g_1\). These are the primes defining the type conditions.
For each coefficient localization \(i\), let \(O_i\) be the coefficient
DVR, \(M_i\) the integral Hecke module and \(X_i\) its set of generic
Hecke characters. Write \(Z_i\subseteq X_i\) for the characters with
the prescribed theta summand and \(Z_i^+\subseteq Z_i\) for the good
ones, whose remaining type polynomials satisfy the tests of Definition~\ref{hyp:retained-type-tests}. The row group is
\[
 \Gamma=\prod_{k=0}^1\,\prod_{v\in\mathcal P_k}
              \mathrm{GL}_r(\mathbb F_{g_k}),\qquad h=|\Gamma|,
\]
where \(\mathcal P_k\) indexes the type places at \(g_k\) up to the
specified conjugacy. The same \(h\) is used at all localizations of
this point. It may vary along the family.
\end{definition}

\begin{lemma}[Integral central operators]\label{lem:integral-central}
Let \(e\) be a central idempotent of \(\mathbb C[\Gamma]\). For any
embedding of a characteristic-zero coefficient DVR \(O\) into \(\mathbb C\) whose image contains a primitive \(h\)-th root of unity, the induced map on group algebras satisfies
\[
 he\in\operatorname{im}\bigl(Z(O[\Gamma])\longrightarrow Z(\mathbb C[\Gamma])\bigr).
\]
The assertion also holds for a sum of products of central
idempotents.
\end{lemma}

\begin{proof}
For \(\gamma\in\Gamma\), let \(T\) be left multiplication by
\(\gamma^{-1}e\) on the regular representation. Its diagonal entries
on the group basis all equal the coefficient of \(\gamma\) in \(e\).
Since \(e\) is central, \((\gamma^{-1}e)^{h+1}=\gamma^{-1}e\).
The eigenvalues of \(T\) are therefore zero or powers of the
primitive root. Its trace is in the embedded ring \(O\). Taking
this trace for each \(\gamma\) gives the coefficients of \(he\).
Coefficientwise injectivity descends centrality. Products of
commuting central idempotents are central idempotents, and the
scaled lifts may be added in the center. In particular the
clearing factor for a product in the type cut is \(h\).
\end{proof}

We use the point data tied to the integral module. A complex
automorphic point \(P\) over a Hecke character \(a\) has a nonzero
form space \(U_P\) of the required tuple and lift index. Every
\(a\in X_i\) has at least one such point. At each \((k,v)\) the
irreducible finite group types have monic polynomial labels over
\(\mathbb F_{g_k}\). The central idempotent of a label is the sum
of the central idempotents of the types with that label. Its
action on the complex row space is the projection to those
isotypic components, by complete reducibility of finite group
representations in characteristic zero. Acting at every
\((k,v)\) gives the product of these commuting projections. The good
selector \(e_0\) is the sum of the products specified by the good
label tuples.

\begin{definition}[Support of the good type projector]\label{def:good-projector-support}
At a localization of Definition~\ref{def:localized-type-support} with \(Z_i^+\ne\varnothing\), the support of the projector \(e_0\) on the point spaces \(U_P\) is
\[
 S_i=\{a\in X_i: \text{some }P\text{ over }a
                    \text{ and }F\in U_P\text{ satisfy }e_0F\ne0\}.
\]
\end{definition}

\begin{proposition}[The good support of the type selector]
\label{prop:cut-selector}
For the integral center action and the depth-zero dictionary
on these point models, \(Z_i\cap S_i=Z_i^+\). There is an integral Hecke-equivariant endomorphism \(T:M_i\to M_i\). On the full eigenspace \(V_x\) at \(x\) it satisfies
\[
 T|_{V_x}=\begin{cases}h\,\mathrm{id}_{V_x},&x\in Z_i^+,\\0,&x\in X_i\setminus S_i.\end{cases}
\]
At each split type place \(w\) at \(g_k\), every point over
\(S_i\) has wild inertia with characteristic polynomial a power
of \(X-1\), zero monodromy, and each repeated eigenvalue of the tame generator has order
dividing \(g_k-1=\mathrm Nw-1\).
\end{proposition}

\begin{proof}
\proofstep{Check the projector on all points of the theta locus.}
When \(a\in Z_i\), the isobaric partition includes the specified
character block of rank \(t\). Removing it and undoing the common
twist gives the rank-\(r\) datum used for the type dictionary.
To see that the removal is available at any \(P\) over \(a\), start
with a point over \(a\) whose isobaric partition contains this block. The isobaric
partitions of points over \(a\) agree up to permutation. The
character block thus occurs at \(P\) with the same twisting
convention.

The parameter surjectivity in the dictionary gives
a type parameter with the same local data. The form space
\(U_P\) lies in the summand for that parameter. The series support
identity acts there by one for a good parameter and by zero
otherwise. The resulting test on the characteristic-zero fibre is
\[
 a\in Z_i^+\quad\Longleftrightarrow\quad\text{the parameter left after removal is good}.
\]
For \(a\in Z_i^+\), transport a partition with good remaining parameter to \(P\)
by a permutation. Conversely, if the parameter after removal is good,
the partition at \(P\) gives \(a\in Z_i^+\). This gives
\[
 e_0 F=F\quad(a\in Z_i^+),\qquad
 e_0 F=0\quad(a\in Z_i\setminus Z_i^+),\qquad F\in U_P .
\]
\proofstep{Descend from the point spaces to the integral Hecke module.}
The fibre calculation determines where the scalar should act. Its use in estimating lengths depends on realizing that scalar by an endomorphism before tensoring with the fraction field.

The integral action of the center of \(O_i[\Gamma]\)
commutes with the Hecke action on \(M_i\). Its comparison with
the action on the form spaces detects any scalar that acts on
all \(U_P\) over \(a\), as a scalar on the entire generic
\(a\)-eigenspace of \(M_i\). This is the tied eigenspace spanning
condition. The coefficient DVRs in this comparison already
contain the roots of unity in Lemma~\ref{lem:integral-central}.

The two action formulas give
\[
 Z_i\cap S_i=Z_i^+,
\]
with inclusion
of \(Z_i^+\) using nonzero form spaces. Lemma~\ref{lem:integral-central}
and the integral center action give an endomorphism of \(M_i\)
acting by \(h\) on the full eigenspaces over \(Z_i^+\). For
\(a\notin S_i\), the same endomorphism acts by zero on the
full eigenspace since all the form spaces there have zero
image. This is the operator used in the length comparison.

\proofstep{Read the inertia polynomial from one occurring type, then use it throughout the fibre.}
If \(e_0F\ne0\), some good tuple product has nonzero
value on \(F\). Each factor has nonzero value since the
factors commute. Expanding a factor as a sum of individual
type idempotents gives an occurring type with the indicated
label at that place. The depth-zero dictionary attaches to
that label \(P_v(X)\) the inertia polynomial in residue-field
labels
\[
                         P_v(X)(X-u_v)^t .
\]
Here \(P_v\) is squarefree with nonzero constant term and
avoids the two forbidden reciprocal labels from the theta
block. The dictionary gives zero monodromy and
\(u_v\in\mathbb F_{g_k}^{\times}\). These assertions include
the conjugate places, with the inverse labels on conjugation.
All roots are nonzero. A root of multiplicity at least two
must come from the repeated rank-\(t\) block and therefore
belongs to \(\mathbb F_{g_k}^{\times}\). The dictionary attaches
distinct continuous tame inertia characters to distinct
nonzero labels.

For two such characters \(\theta,\theta'\) and a tame generator
\(t\), agreement at \(t\) would give
\[
 (\theta/\theta')(t)=1\ \Longrightarrow\
 \overline{\langle t\rangle}\subseteq\ker(\theta/\theta').
\]
The kernel is closed by continuity and the cyclic subgroup is dense,
so the two characters would coincide. The dictionary is
multiplicative in the residue-field labels, so eigenvalues
of multiplicity at least two have order dividing \(g_k-1\).
At these completely split places \(g_k=\mathrm Nw\). Thus the local parameter satisfies the type condition. Writing \(m\) for the local rank and \(\lambda\) for any repeated tame eigenvalue, the three tests read
\[
 P_{\rm wild}(X)=(X-1)^m,\qquad N=0,\qquad \lambda^{\mathrm Nw-1}=1.
\]

This allows the repeated character block.

This argument used one point \(P\) over \(a\) with \(e_0F\ne0\) for some \(F\in U_P\).
At any other point above \(a\) the equality of the isobaric
partitions gives the same local parameter at each split
type place. Consequently this condition holds at every point of the cut, including points whose form spaces were unnecessary for deciding membership in \(S_i\). The cut comparison
therefore applies to the whole retained Hecke module.
\end{proof}

\subsubsection*{Compatibility with the integral module and rank}

Proposition~\ref{prop:cut-selector} identifies the generic eigensystems on \(S_i\). Their integral structure is the image of the original lattice under the same projection. Write \(V=V_{S_i}\oplus V_{S_i^c}\) for the rational
decomposition and \(M_{S_i}\) for the projected lattice.
Its exact presentation is
\[
  0\longrightarrow M_i\cap V_{S_i^c}
    \longrightarrow M_i\longrightarrow M_{S_i}\longrightarrow0.
\]
A central operator commuting with Hecke preserves each
eigenspace, so it preserves the kernel and descends integrally.
Writing \(p\) for the cut projection and \(V_a\) for the generic eigenspace at a retained character, one has
\[
 p|_{V_a}=\operatorname{id}_{V_a}.
\]
Hence the same point-space tests
detect its scalar action before and after taking the cut.

The row labels, the
inertia polynomials and the attached representations at the
retained points still come from the original place and tuple.
The finite paired-character equations restrict with them.
In particular the restriction uses the original norm choice,
with fewer local characters to test.

The reduced space still realizes the same tame and Hodge--Tate
labels at each retained character. At places over the coefficient prime use local--global compatibility with the specified inertia and monodromy condition at the type primes and the unramified parameter elsewhere. For the de Rham, potentially semistable representation at an unramified type, the criterion is
\[
 \mathrm{WD}(R)|_{I_v}=1,\quad N=0\quad\Longrightarrow\quad R|_{G_{E_v}}\text{ is crystalline}.
\]

At the
type primes the monodromy and inertia calculation is used in
the potentially crystalline comparison. The Hodge--Tate jumps
are those of the algebraic level-zero row. Their multiplicity
one implies that the semisimplified representation has distinct
absolutely irreducible constituents. Every summand of a de Rham representation is de Rham. At a fixed embedding \(\tau\), if an irreducible \(U\) occurred at least twice, any jump \(a\) of \(U\) would give
\[
 \dim\operatorname{gr}^a D_{\rm dR,\tau}(R^{\rm ss})\ge 2\dim\operatorname{gr}^a D_{\rm dR,\tau}(U)\ge2,
\]
contrary to simplicity of the labelled jumps. Here \(R\) denotes the attached representation. Thus the retained eigenspaces carry both the required inertia polynomials and the distinct labelled jumps used to identify their multiplicities.

\subsubsection*{Multiplicities at a twisted coefficient place}

Projection retains entire eigenspaces over the selected characters. We now compare their total multiplicity with the rank of the twisted theta image at a single place.

Fix a lift \(j\), a twisted coefficient place \(w\) and the
twist \(\mu=\mu(w)\) assigned to it. The row is the internal
direct sum \(M_j=\bigoplus_\sigma M_{j,\sigma}\).
Write \(\Theta_\mu\) for the theta map on this row and
\(\Theta_\mu^*\) for its adjoint. On \(M_{j,\sigma}\)
the operator \(\Theta_\mu^*\Theta_\mu\) is a positive
scalar. The twisted zero-space is the image of \(\Theta_\mu\)
and has dimension \(d_j(w)=\dim M_j\).

\begin{proposition}[Multiplicity after the type cut]
\label{prop:cut-multiplicity}
For a lift \(j\) and twisted coefficient place \(w\) on the cut
of Proposition~\ref{prop:cut-selector}, let
\(m_{i,\pi}\) be the dimension of the full good Hecke
eigenspace at a retained localization \(i\). Then
\begin{equation}\label{eq:retained-multiplicity}
 \sum_{i\ {\rm over}\ (j,w)}\;\sum_{\pi\in Z_i^+} m_{i,\pi}
             \le d_j(w)=d_j .
\end{equation}
\end{proposition}

\begin{proof}
\proofstep{Assign each good lower-row label by its integral split polynomials.}
For a localization \(i\) over \((j,w)\), assign a good
row label \(\sigma\) to \(i\) if the normalized polynomials
of its lifted parameter have integral representatives with
the residual Hecke values of \(i\) at every split place
away from the residue characteristic, the original
ramification set and the auxiliary level. Form the lifted
parameter by adding the rank-\(t\) character part and
taking the common normalizing twist. The local Hecke order
has a single residue character, so the test may use any
point of the localization. At surviving localizations it
may use any retained point. In the cut the polynomials,
their residue test and the assigned row label set are
therefore preserved. Projection of the split generic Hecke
algebra to a subset of characters retains their whole
eigenspaces. In particular, if \(V_a\) and \(V_a^{\rm cut}\) denote the original and cut eigenspaces at a good character, then
\[
 V_a^{\rm cut}\simeq V_a,\qquad \dim V_a^{\rm cut}=\dim V_a.
\]


\proofstep{Identify the full good eigenspace by the theta map and its adjoint.}
For a good character \(\pi\) and an isobaric row parameter
\(\Sigma\) for it, let \(M_{j,\Sigma}\) be the sum of row
spaces with that parameter and \(V_{i,\pi}\) the full
Hecke eigenspace. A row label belongs to this sum when its
local reading agrees with that of \(\Sigma\) at each split
place. We use the split local theta comparison
on the tied row and point models. The upward image of each
row type is in its specified point type. The local parameter
formula adds the character part and applies the common twist.
The split components of this base change determine the
Hecke character. The primes excluded from this split comparison have
Dirichlet density zero over \(E\).
For an inert place \(w\) above \(v\),
\[
 \mathrm Nw=(\mathrm Nv)^2.
\]

The Euler product of the lower field at \(2\) converges, so
\[
 \sum_{w\ {\rm inert}}(\mathrm Nw)^{-1}
   \le\sum_{v}(\mathrm Nv)^{-2}<\infty .
\]
The other exclusions form a finite set. In each finite Galois quotient
of \(G_E\), a nonempty conjugacy class has positive
Chebotarev density, so it meets the remaining Frobenius
classes. Local--global compatibility at those primes
and the common twist identify the functions
\[
       g\longmapsto \det(1-X\rho(g))
\]
for the two attached representations on a dense set.
Their continuity gives equality on \(G_E\). At all
unramified places the corresponding standard parameters
therefore agree. Strong multiplicity one for isobaric
representations identifies their constituents with multiplicity.
In the tied model these lists determine the local split
and unramified inert parameters used for the point
comparison.

Conversely, equality of the good split
Hecke readings gives equal split parameters by taking
the roots of the common Hecke polynomials and undoing
the normalization. \label{cor:selected-cut-1}
Thus \(\Theta_\mu M_{j,\Sigma}\subseteq V_{i,\pi}\).
For a row vector \(v\) in a different parameter and
\(u\in V_{i,\pi}\), use
\[
 \langle v,\Theta_\mu^*u\rangle
                    =\langle\Theta_\mu v,u\rangle .
\]
A nonzero value requires an occurring type over \(\pi\)
by orthogonality. Its split Hecke polynomials then agree
with those of \(\pi\), and cancelling the common character
block and twist identifies its row parameter with \(\Sigma\).
Thus \(\Theta_\mu^*V_{i,\pi}\subseteq M_{j,\Sigma}\).
The upward injection and the injection of the adjoint
on the full good eigenspace
(Proposition~\ref{prop:downward-on-good-space}) give
\[
 \dim M_{j,\Sigma}\le\dim V_{i,\pi}\le\dim M_{j,\Sigma},
 \qquad \Theta_\mu M_{j,\Sigma}=V_{i,\pi}.
\]

\proofstep{Identify the assignment with the sum of full good eigenspaces.}
A nonzero good row component assigned to \(i\) lifts
into a point type. Its lifted parameter has the split
polynomials of the residual test. The local point exhaustion
on this model puts the lift in the localization \(i\).
The two type-prime labels of its rank-\(r\) part are those
of the row, so the point is good. Conversely a good point
supplies its own integral Hecke polynomials for the residual
test. Any row of its parameter has good labels, since
base changes over this point have the same split components,
and the character part and twist cancel as above. These
two inclusions give the sum of good row components assigned
to \(i\) as the sum of the \(M_{j,\Sigma}\) for good
characters at \(i\). Put \(M^+\) for this assigned sum. Its image and dimension are
\[
 \Theta_\mu(M^+)=\bigoplus_{\pi\in Z_i^+}V_{i,\pi},\qquad \dim\Theta_\mu(M^+)=\sum_{\pi\in Z_i^+}m_{i,\pi}.
\]
For the upper bound an empty good point set contributes zero.

\proofstep{Separate the residual Hecke characters before summing dimensions.}
Summing over localizations gives a rank bound only if distinct localizations count disjoint lower-row labels.

Distinct localizations over \((j,w)\) have disjoint
assigned row sets. Their coefficient DVRs are identified
compatibly with the complex comparison. The same row
label has the same split parameter at both, and the added
character depends on \(j,\mu(w)\). If a label were
assigned to both localizations, choose a split place \(v\)
good for both and an index \(a\) whose Hecke generator has
different residues, using separation of the residual ideals.
The two lifted polynomials at \(v\) agree over the
identified DVR because their complex images agree.
Their coefficients of \(T^{r+t-a}\) have the same
reduction. Writing \(H_v(T)\) for the lifted polynomial and \(t_{v,a}\) for the Hecke value, the coefficient is
\[
 [T^{r+t-a}]H_v(T)=(-1)^a(\mathrm Nv)^{a(a-1)/2}t_{v,a}.
\]

The multiplier is a unit since \(v\) is away from the
residue characteristic of \(w\). Thus the two Hecke values
have the same reduction, a contradiction. Summing dimensions over the disjoint sets inside \(M_j\)
gives the inequality. The comparison of twisted and base
zero-spaces gives the last equality for this lift.
\label{cor:selected-cut-2}
\end{proof}

\subsubsection{The downward lift on the full good eigenspace}
\label{subsec:selected-point-adjoint}

The preceding multiplicity bound uses the rank of the whole good eigenspace. To compute it by the downward lift, we need injectivity of the adjoint on that eigenspace, including the multiplicity space of each occurring representation.

\begin{established}[Tied level, Weil data and good ranks]\label{hyp:tied-adjoint-level}
Consider the realizations at the fixed lift and its twist
on the selected cut of Proposition~\ref{prop:cut-selector}. Write \(A=V_j\) for the definite
Hermitian space used to realize the point coefficient type and
\(W=W_j\) for the lower space. Their dimensions and stable range are
\[
 (\dim_E A,\dim_E W)=(r+t,r),\qquad 0<r,\qquad t\ge r+5 .
\] The finite bad set \(S\) is the
set below the ramification and auxiliary level of this twist.
The source form has the fixed principal level and Cartan type,
and the splitting characters and Weil representation at both
levels are the ones in the tied realization. In particular
the host splits at the dyadic places in the application of
the first term identity.
\end{established}
\begin{proof}
Proposition~\ref{prop:cut-selector} selects points inside the fixed signed theta realization of Proposition~\ref{prop:signed-theta-data}. Projection to a type or an eigenspace restricts the vector space of functions and leaves its Hermitian spaces, splitting characters, Schwartz data and principal level fixed. The dimension at a lift is \(r+t\), with its lower constituent of dimension \(r\). The selected-rank estimates in Section~\ref{sec:family-window} give \(r>0\) and \(t\ge7r+4\ge r+5\). Finally \(-d\equiv1\pmod 8\) by the choice of \(d\). Hence \(\mathbb Q(\sqrt{-d})\otimes\mathbb Q_2=\mathbb Q_2\times\mathbb Q_2\). Base change to every dyadic completion of the totally real subfield splits the CM quadratic algebra there.
\end{proof}

\begin{proposition}[Injectivity on a full good Hecke eigenspace]
\label{prop:downward-on-good-space}
For the tied models of Lemma~\ref{hyp:tied-adjoint-level}
with \(r>0\) and \(t\ge r+5\), at each selected lift and twist let \(V\) be its full good Hecke eigenspace. The adjoint theta operator satisfies
\[
 \ker(\Theta_\mu^*|_V)=0.
\]
\end{proposition}

\begin{proof}
\proofstep{Read the good spherical system on both kinds of places.}
Let \(f\) be any element of the full Hecke eigenspace at a
good point. Read \(f\) in the vector-valued automorphic model
on \(A\) and take its inverse image \(u\) in the multiplicity
space. This reading is injective. The actions of the finite
Hecke operators on the reading and on the multiplicity space
are intertwined. In particular, if \(u\ne0\), then it has
the required spherical system of eigenvalues away from \(S\).
Here the assertion concerns a joint eigenvector, and the
system of eigenvalues can have arbitrary multiplicity.
The untwisted rank-\(r\) constituent and common normalizing
characters determined by the point give the following
property of the system.

For a good place \(w\) of \(E\)
over \(F\), let \(x_w\) be the unitary character value
of the twisted datum. The system contains the string
\[
 \left\{ x_w\,(\mathrm Nw)^{\,j-(t-1)/2}\mid 0\le j<t\right\}.
\]
Removing this multiset leaves \(r\) entries,
each of absolute value at most \((\mathrm Nw)^{(r-1)/2}\).
The local comparison with base change gives this description
at split and inert good places, including the passage from a
spherical parameter over \(F\) to the entries over \(E\).
The transported good projector acts as the identity on
\(u\). In proving this last assertion, use that the
archimedean signature and both selected finite type labels
of a localization form are good and that the arithmetic
idempotent reads as the same type projector on the
automorphic model. These two transports use the actual
twisted place and the same row index. \label{cor:selected-point-adjoint-1}

\proofstep{Detect a nonzero vector by the doubling integral.}
Suppose \(u\ne0\). At each real label take a vector \(v_\tau\)
of the prescribed lower compact type, and let \(w_0\) belong to
the definite source type. Use their Cartan inclusion
\[
 P=\iota_P((\bigotimes_\tau v_\tau)\otimes w_0)
\]
into the
Fock polynomials. The polarized vacuum test on the chosen
levels is nonzero on this Cartan summand. Pure tensors span it,
so choose \((v_\tau),w_0\) for which \(P\) times the vacuum
gives a nonzero vector \(\Phi\).

Take as \(a\) the scalar
component of the adelic extension of \(u\) paired against
this same \(w_0\). Schur orthogonality in the definite type gives
\[
 \int_{U(A)_\infty}
    \langle\omega(b)\Phi,\Phi\rangle\,r_a(b)\,db
  = c_1c_2\int_{U(A)_\infty}
       \big|\langle w_0,\rho(b)w_0\rangle\big|^2\,db ,
\]
with the conventions for the two coefficients chosen
conjugately. The Weil factor \(c_1\) is nonzero by the
nonzero chosen Cartan polynomial and the norm of its finite
part. The factor \(c_2\) is nonzero because the definite
multiplicity norm of \(u\) is positive. The remaining
continuous matrix coefficient is nonzero at identity.
This proves the real occurrence integral is nonzero and
also shows \(a\ne0\). Let \(p_0\) denote extraction of the chosen scalar component. At each good finite Hecke operator \(T\),
\[
 T a=T p_0(u)=p_0(Tu).
\]
Hence \(a\) has the same spherical parameter. The chosen finite Schwartz readings
are spherical away from \(S\) and principal-congruence
fixed at the level of \(a\). \label{cor:selected-point-adjoint-2}

The occurrence integral supplies the nonzero local factor in
the scalar zeta integral. To locate its pole, in the standard
partial \(L\)-function of \(a\) twist by
the inverse of \(x\). Write \(L_c^S\) for the Euler product on the \(r\) remaining entries after the twist. With \(w\mid S\) denoting all places over \(S\), the incomplete zeta factor is
\[
 \zeta_E^S(s)=\zeta_E(s)\prod_{w\mid S}(1-(\mathrm Nw)^{-s}).
\]
The product for the standard factor is
\[
 L_{\rm std}^S(s,a,x^{-1}) =
   L_c^S(s)\prod_{j=0}^{t-1}
          \zeta_E^S\big(s+(t-1)/2-j\big).
\]
The absolute parameter bound makes the Euler logarithm for
\(L_c^S\) normally convergent on
\(\RePart s>(r+1)/2\).

At \(s_*=(t+1)/2\) the last
zeta argument is \(1\), and all other arguments are integers
at least \(2\). The simple pole of the Dedekind zeta function
has nonzero residue, and none of its finitely many removed
factors vanish at \(1\). The factors with \(j<t-1\) are nonzero at \(s_*\), as is
\(L_c^S(s_*)\) by convergence of its Euler logarithm. Hence
\[
 \operatorname*{Res}_{s=s_*} L_{\rm std}^S(s,a,x^{-1})
  =L_c^S(s_*)\!
      \prod_{j<t-1}\zeta_E^S(t-j)
         \operatorname*{Res}_{u=1}\zeta_E^S(u)\ne0 .
\]
It is a simple pole.

Add \(t\) hyperbolic planes to \(W\) to obtain \(B\), of
rank \(r+2t\). Its doubled Siegel section is built from the
pair of copies of \(\Phi\boxtimes \Phi_H\), where \(\Phi_H\)
is the split Weil vacuum with principal-congruence support.
The complementary ranks in the doubled space and the positive doubling point satisfy
\[
 (r+2t)+r=2(r+t),\qquad s_0=\frac{(r+2t)-(r+t)}2=\frac t2.
\]

For clarity, the dimension parameter of the split group in
the first term identity is its Siegel parabolic rank \(r+t\).
Apply that identity in the complementary-space range,
with the compatible splitting character, the majorant
compact average, and the Ikeda map toward \(W\).
When \(W\) is anisotropic the paired residue is proportional
to the convergent integral of its theta kernel. When \(W\)
is isotropic it is the indicated leading regularized theta
coefficient on \(W\).

In this second case apply the
regularizing operator before integration, using the Witt
decomposition and the nonzero scalar of its auxiliary
Eisenstein series. Write \(\mathcal R\) for the paired residue before integration against the two forms. The implication used in both cases is
\[
 \Theta_{W,\phi}(a)=0\ \ (\text{all }\phi)\quad\Longrightarrow\quad \langle\mathcal R,a\otimes\overline a\rangle=0.
\]
Indeed, on doubled
Schwartz tensors the pulled-back theta kernel separates as
two theta factors, so integrating against the two forms
gives products of the corresponding lifts. The same
vanishing applies to the compact-averaged and regularized
inputs by the seesaw identity, continuity on Schwartz
vectors and the nonzero regularizing scalar.
\cite{gqt-rallis,yamana-first-term}\label{cor:selected-point-adjoint-3}

On the right of absolute convergence the doubling
unfolding instead gives the scalar zeta integral as
\[
 Z(s)=C\,\frac{L_{\rm std}^S(s+1/2,a,x^{-1})}{d_S(s)}.
\]
The denominator \(d_S\) is the unramified spherical
normalizer on the doubled group. For positive real part its
shifted unitary-character factors are absolutely
convergent and nonzero. The local factor \(C\) is independent
of \(s\) for the chosen data. To see this, the hyperbolic
vacuum coefficient vanishes outside the principal
congruence factor in the finite integration at level.

On this support the exponentiated height in the unfolding
is \(1\) and every right-level translate fixes the
coefficients. On the real compact group the hyperbolic
vacuum is fixed, leaving its nonzero squared norm times the
definite occurrence integral above. This yields \(C\ne0\).
Both \(C\) and \(d_S(s_0)\) are nonzero. The unfolding and
meromorphic continuation therefore give
\[
      \operatorname*{Res}_{s=s_0} Z(s)
        =\frac{C}{d_S(s_0)}
                 \operatorname*{Res}_{s=s_*}L_{\rm std}^S(s,a,x^{-1})
        \ne0,\qquad s_*=s_0+\tfrac12 .
\] This contradicts
vanishing of all the lifts and proves the existence of a
Schwartz vector and a scalar reading for which the integral
on \(W\) is nonzero. \label{cor:selected-point-adjoint-4}

\proofstep{Use the finite level generators to test the downward pairing.}
The test vectors used so far range over all Schwartz functions. At the selected level the adjoint uses a specified kernel; the finite level generators connect these two pairing tests.

The preceding conclusion uses all Schwartz data. To pass
to the actual downward operator, suppose its output on
\(u\) is zero. Its good target projection agrees with the
whole normalized level output on a good source by the
local type comparisons at the selected primes. The
output pairs the Cartan tensors against the coefficients
of the target compact types.

For a target argument \(g\) and a source linear functional \(\lambda\), denote the downward pairing with test \(\phi\) by \(B_{g,\lambda}(\phi)\). Vanishing on the generating tensors implies, by linearity,
\[
 B_{g,\lambda}(P\,\phi_{\rm vac})=0\qquad(\text{every allowed polynomial }P).
\]
Here \(\phi_{\rm vac}\) includes the fixed finite vacuum.
This uses the span of the real Fock polynomials under the
compact tensor maps and target translation. The same
vanishing holds for all Schwartz vectors whose finite
Heisenberg translates satisfy the linear relations of
the chosen finite vacuum. The real Schwartz approximation
on this fixed finite module is in the Sobolev norm for
which the downward pairing at each fixed argument is
continuous.

For a source level element \(k\) with character \(\chi_V(k)\),
right translation of the integration variable gives
\[
 I_u(\omega(k,1)\phi,l,g)=\chi_V(k)^{-1}I_u(\phi,l,g).
\]
Project a Schwartz test onto the inverse character by
averaging over the compact source level. On every fixed
finite vector this is a finite sum over cosets of an open
stabilizer after intersecting with the character kernel.
The difference from the projection pairs to zero. On the
inverse-character space let \(G\) be the finite level
quotient and \(k_b\) lifts of its elements \(b\). The resulting right action on the tests is
\[
 \pi(b)=\chi_{\rm flag}(k_b)^{-1}\omega(k_b^{-1},1),\qquad b\in G,
\]
where \(\chi_{\rm flag}\) is the flag part of the level character. The extra generators are the images of
\(\sum_{b\in G}c_b\pi(b)\) of the following types.

At one
of the selected residual places and for a nonzero vector
\(a\), take the subgroup of matrices fixing \(a^\perp\)
pointwise or inducing the identity on the quotient by
the line through \(a\). Use the identity at the other
places. For either choice of subgroup \(H\), write \(\eta\) for its prescribed determinant character or its inverse. The group algebra element acting on the source is
\[
 \sum_{b\in G}c_b\pi(b)=\sum_{b\in H}\eta(b)\pi(b).
\]

The selected good source projector kills these same
character sums with the normalized right action on
every scalar component. The substitution
\(h\mapsto hk_b\) in the theta integral transfers test
averaging to this source averaging, including the factor
\(\chi_{\rm flag}(k_b)^{-1}\). On a finite-dimensional orbit space \(V\), each subgroup \(H_j\) in this construction gives the normalized average
\[
 P_j=\frac1{|H_j|}\sum_{b\in H_j}c_b\pi(b),
\]
where the normalization is taken in characteristic zero. These projections
are self-adjoint in the invariant metric, so
\[
              \left(\sum_j\operatorname{im}P_j\right)^\perp
                            =\bigcap_j\ker P_j .
\]
In particular the self-adjoint projections give the orthogonal decomposition
\[
 V=\biggl(\bigcap_j\ker P_j\biggr)\oplus\biggl(\sum_j\operatorname{im}P_j\biggr).
\]
The good source
projector kills the pairings from the latter.

On the common kernel the local generation comparison
spans the inverse-character tests by target translates of
the fixed finite Heisenberg module above. This comparison
is applied one place at a time, with the spherical
generation at split and inert off-level places and the
prescribed local types at the twist and type primes.
Target translation replaces \(g\) by its right translate.
The normalized output thus tests a spanning family of the
inverse-character quotient and forces every downward
pairing of \(u\) to vanish. The pole test now applies in its full form:
\[
 \bigl(\Theta_{W,\phi}(u)=0\ \text{for every allowed }\phi\bigr)\quad\Longrightarrow\quad u=0.
\]

\label{cor:selected-point-adjoint-5}

\proofstep{Identify the normalized level operator with the actual adjoint.}
Finally identify this normalized level operator with the
point-to-row map \(\Theta_\mu^*\). The theta rule on simple
tensors is the adelic integral in the chosen source and
target coordinates. Both rules preserve level. They agree
after composition with the chosen embedding of the finite
coefficient model. The level comparison gives injectivity
of this restriction on the subspace of level rules, so the
resulting full coefficient tensors agree. Hence zero under
\(\Theta_\mu^*\) gives zero under the normalized level
operator. Applying its injectivity and the source coordinate isomorphism
proves \(f=0\) for every vector of the full eigenspace.
\label{cor:selected-point-adjoint-6}
\end{proof}


\subsubsection*{The character and the two lattices on the retained family}
\label{subsec:marked-localization}

The adjoint comparison supplies the good row at the specified lift. Its contributions to the numerator are local: at each twisted place they use a particular Hecke character and the two integral lattices attached to it.

Fix a member of the cut family at scale \(x\), a lift \(j\), and
a twisted coefficient place \(w\mid\ell\). Let \(i\) be one of
the Hecke localizations at \((j,w)\), with coefficient integers
\(O\), and let \(\pi\) be a good eigencharacter there. In the
following argument the index \(i\) specifies the twist, its side
and its auxiliary level \(Q_i\). Write \(E\) for the auxiliary CM field
and \(v\) for a place of \(E\) above \(\ell\).

Let \(r\) be the rank of the complementary representation and \(t\) the
number of characters in the theta string. For the \(\ell\)-adic cyclotomic
character \(\chi_\ell\) of \(G_E=\operatorname{Gal}(\overline E/E)\), set
\begin{equation}\label{eq:signature-cyclotomic-minimum}
 \nu_\ell=\min_{g\in G_E}v_\ell\!\left(\chi_\ell(g)^{tr}-1\right),
 \qquad v_\ell(\ell)=1.
\end{equation}

\begin{proposition}[The signature and the complement at a marked point]
\label{prop:marked-signature}
Fix the norm choice on the retained family before choosing the
integral realizations at \(i\). The complement of the theta string
at \(\pi\) identifies, with its integral lattice, with one of the
good signatures tested on the prescribed side in
Proposition~\ref{prop:one-retained-character}. For the length \(s_\pi\) of
the annihilating quotient and the inertia readings \(\zeta_v\)
of the chosen twist, the simultaneous estimates are
\[
 s_\pi<\infty,\qquad s_\pi\le\nu_\ell+\ell^{-2},\qquad
 t\sum_{v\in Q_i} v_\ell(\zeta_v-1)\le\ell^{-2}.
\]
The one choice at \(w\) applies to all these presentations.
\end{proposition}

\begin{proof}
\proofstep{Read the normalized characteristic polynomial on the chosen side.}
Proposition~\ref{prop:cut-multiplicity} gives an occurring good
lower row \(\sigma\) for \(\pi\) at lift \(j\). In the coefficient model
take a nonzero common eigenvector \(x_\sigma\) in its good
isotypic space. The descent of an eigenvector here uses only
linear algebra. Take the basis of the small-field form, also a
basis after extension to the row field, and choose a nonzero
coordinate \(u\) of \(x_\sigma\). Over the subfield containing
the eigenvalues and the small-field matrices, choose a linear
functional on the row field with value \(1\) at \(u\).

Apply it
to each coordinate and then apply the coefficient embedding of
the chosen side. Write \(\lambda\) for the chosen functional and apply it coordinatewise. For every small-field Hecke matrix \(T\) and eigenvalue \(a_T\) in the field of linearity,
\[
 T\lambda(x_\sigma)=\lambda(Tx_\sigma)=a_T\lambda(x_\sigma),\qquad \lambda(u)=1.
\]
Applying the coefficient embedding preserves these identities. For the level action, undo the row character on each
group element. Its inverse value is in the eigenvalue field,
so the map also intertwines the sum giving the good projector.

The resulting vector lies in the good level space on that side.
At a split unramified prime \(P\) outside the finite omitted set,
its \(h\)-th Hecke eigenvalue is the embedded row eigenvalue
times \(c(P)^h\), with \(c(P)\ne0\) the normalizing character
factor in the level action. Decompose the good level space at \(j\) on this side into its recorded simultaneous eigenspaces. A nonzero component yields a signature \(s\) with the same normalized eigenvalues for every \(1\le h\le r\).

For corresponding unnormalized eigenvalues \(A_h,B_h\) and coefficient embedding \(\iota\), the cancellation reads
\[
 c(P)^h\iota(A_h)=c(P)^h\iota(B_h)\quad\Longrightarrow\quad A_h=B_h.
\]
For the normalized split roots \(\alpha_1,\ldots,\alpha_r\)
these values are the coefficients in
\[
 \prod_{i=1}^r (T-\alpha_i)
     =\sum_{h=0}^r(-1)^h e_h(\alpha_1,\ldots,\alpha_r) T^{r-h}.
\]
Thus equality of the embedded eigenvalues after dividing by
\(c(P)^h\) gives equality of polynomials, with all multiplicities. Undoing the common reciprocal and
scalar normalization identifies the Frobenius polynomials.
The signature is one of the good ones tested in
Section~\ref{subsec:retained-twist}: the localization at \(w\)
uses a nonzero twist index, and its good type and lower
cuspidal system are precisely those just selected.
\label{cor:marked-localization-1}

\proofstep{Choose integral realizations of all retained characters.}
Take the full semisimple integral Galois realization \(R_x\)
for each generic eigencharacter \(x\) of this localization,
as in Section~\ref{subsec:retained-galois}. At a theta
character the characteristic-zero representation has its
labelled string of characters and a continuous integral
complement \(\rho_x\). Denote the labelled characters by \(\chi_h\). At a good split Frobenius element \(g\), the comparison removes the string by the factorization
\[
 \det(1-XR_x(g))=\det(1-X\rho_x(g))\prod_h(1-X\chi_h(g)).
\]
The attached cuspidal summands realize the complementary factor. Density and semisimplicity identify the rational complement.

To obtain a stable lattice, start with any full lattice
\(L_0\) in the complement. Compactness gives a single integer
\(a\ge 0\) such that, for a coefficient uniformizer \(\pi\),
\[
 L_0\ \subseteq\ L=\sum_{g\in G_E}\rho_x(g)L_0
       \ \subseteq\ \pi^{-a}L_0,\qquad
 \rho_x(h)L=L\quad(h\in G_E).
\]
The upper lattice is Noetherian, so the middle module is a
full lattice. This construction only concerns the chosen complement.
The coefficient choice of Section~\ref{subsec:retained-galois}
splits the constituents and the roots of the marked corner.
At a retained non-theta character only \(R_x\) enters the matrix algebra. The unused complement may be assigned
\[
 \rho_x=\mathbf1^{\oplus r}\qquad(x\ \text{outside the theta locus}).
\]

This constructs one family at all the retained characters.
\label{cor:marked-localization-2}

\proofstep{Identify the two character evaluation kernels.}
For the comparison with \(s\), write \(O_0\) for the signature's
base coefficient ring. The character \(\psi_0\) on the parent
at \(w\) has its values in a common coefficient field. Its
two evaluations give the Galois character at \(i\) and the
character in the twist field, with the recorded side action.
The localization construction extends the embeddings to a
common field. On \(O_0\) it uses the prescribed map into
\(O\) and the inclusion into the twist integers. On the
parent field the two maps compare the same embedding at \(w\).
Concretely, two embeddings inducing \(w\) differ on the
completed algebraic coefficients by a local field embedding.
Extend over the finite coefficient fields into an algebraic
closure and transport the integral embedding by this
conjugacy.

For the local conjugacy \(\sigma\) and the absolute valuation \(v\), one has
\[
 v(\sigma a)=v(a)\ge0\qquad(a\in O_0).
\]
Thus the transported base coefficients are integral. The field used
for the localized row includes these base readings by
construction. Let \(\operatorname{ev}_i\) and \(\operatorname{ev}_s\) be the
two character evaluations on \(O_0[G_E]\), and let \(\iota_i,\iota_s\)
be their embeddings into the common field. For
\(a=\sum_g a_g[g]\) one has
\[
 \iota_i(\operatorname{ev}_i a)
     =\sum_g\iota_i(a_g)\iota_i(\operatorname{ev}_i[g])
     =\sum_g\iota_s(a_g)\iota_s(\operatorname{ev}_s[g])
     =\iota_s(\operatorname{ev}_s a).
\]
Injectivity of the common-field embeddings gives, for \(a\in O_0[G_E]\),
\[
 \operatorname{ev}_i(a)=0\iff\iota_i\operatorname{ev}_i(a)=0\iff\iota_s\operatorname{ev}_s(a)=0\iff\operatorname{ev}_s(a)=0.
\]
Applying this equality to \([g]-[1]\) also identifies
triviality on each decomposition group. The side inversion in taking the Galois
character of the ray class character is undone when the
signature is read back on that side.
\label{cor:marked-localization-3}

\proofstep{Transfer the length and inertia estimates on normalized lattices.}
The comparison of representations becomes a length estimate only after the character, the signature lattice and the chosen side of the place have been identified together.

Apply Lemma~\ref{lem:retained-signature-specialization}
to this \(\rho_\pi\) and this \(s\). The split polynomial identity
identifies the two rational rank-\(r\) representations with the
prescribed twists. After integral scaling the intertwiner \(A\) has nonzero reduction. Absolute residual irreducibility gives
\[
 \ker\bar A=0\quad\Longrightarrow\quad\det A\in O^\times,
\]
since the two lattices have the same rank. The equality of evaluation kernels
identifies the generated orders. Extension from their common
normalization preserves normalized length and the injection
into the signature Selmer group. The ramification set on the
signature side is exactly \(Q_i\). At each such place a marked inertia generator gives a generator
\(\zeta\) of the same finite character image. A replacement by
\(\zeta^a\) with \(a\) prime to its order preserves
\(v_\ell(\zeta-1)\).

Choose positive representatives \(a,b\) with \(ab\equiv1\) modulo
the order of \(\zeta\). The two divisibilities follow from
\[
 \zeta^a-1=(\zeta-1)\sum_{j=0}^{a-1}\zeta^j,\qquad
 \zeta-1=(\zeta^a-1)\sum_{j=0}^{b-1}\zeta^{aj}.
\]
Thus the two elements generate the same ideal. Thus its simultaneous character estimate gives the asserted
bounds on the actual complement.
The character used for this application was fixed with the
norm choice before \(R_x,\rho_x\) were selected. At the cut
the maps on the retained points restrict the original
readings, including \(Q_i\) and the side. This proves the
annihilation estimate for all of these integral presentations
with one choice at \(w\). \label{cor:marked-localization-4}
\end{proof}

At a labelled embedding of the host at \(v\), the archimedean
weights give the following candidate lists. Write \(c=(t-1)/2\), and use the real place \(\tau\)
under this embedding. In the Harish--Chandra parameter at
\(\tau\) put \(s=e_{j-1}\), write \(v_a\) for its \(k-1\)
decreasing offsets and \(b_h\), \(0\le h<m\), for the
increasing list interlaced with the \(e\)'s. On the positive
half the three sets for \(R_x\chi_1^{-1}\) are
\[
 C_+=\{c+s+v_a\}_a,\qquad D_+=\{c+s-b_h\}_h,\qquad
 I=\{0,\ldots,t-1\}.
\]
On the other half reflect the lists by \(a\mapsto t-1-a\),
writing \(C_-,D_-\) for the first two images.

\begin{lemma}[Signed jumps on the row and column]
\label{lem:marked-signed-jumps}
Writing \(\operatorname{HT}_d\) for the multiset of jumps on side \(d\) at a retained eigencharacter, one has
\[
 \operatorname{HT}_d(R_x\chi_1^{-1})=C_d\sqcup D_d\sqcup I,
\]
with all multiplicities one. At \(\pi\) the complement \(W=\rho_\pi\chi_1^{-1}\)
has jumps \(C_d\cup D_d\) and avoids \(0,-1\). At paired split
places the positive column of \(W\) and negative row of
\(W^\vee(1)\) each have \(m-j\) negative jumps per embedding.
\end{lemma}

\begin{proof}
\proofstep{Normalize the level-zero jump by the first character.}
Let \(p\) be a Harish--Chandra parameter
entry, \(d=\pm1\) the sign of the embedding and \(c_\tau\)
the common archimedean shift. The common shift cancels on twisting by \(\chi_1^{-1}\). Subtracting its jump from the level-zero jump gives
\[
 (dp+c_\tau)-(-ds-c+c_\tau)=d(p+s)+c.
\]

For the three pieces of the shifted parameter \(p+s\) the
entries are \(s+v_a\), \(s-b_h\), and the consecutive integers
from \(-c\) to \(c\). This proves the indicated reflection and
multiplicity formula. The cut gives the level-zero jump
identity on every retained eigencharacter, including at a
type prime. Away from the type primes its unramified
Weil--Deligne parameter and zero monodromy give
crystallinity; at a type prime use the potentially crystalline
comparison with the retained tame label.
\label{cor:marked-localization-5}

\proofstep{Separate the consecutive string from the signed complement.}
The inequalities on the selected weights give
\(0\le v_a\le S_c\), \(|s-b_h|\le\eta_c\), and
\(s>\eta_c+t+S_c\). The interlacing margins give the opposite signs on either side of \(j\):
\[
 s-b_h\ge(t+3)/2\ (h<j),\qquad s-b_h\le-(t+3)/2\ (h\ge j).
\]
These follow from the interlacing
margins on either side of \(e_{j-1}\) and monotonicity of
the \(b_h\). The three sets are disjoint and each jump is
simple. In particular removing \(I\) at \(\pi\) leaves
exactly \(C_d\cup D_d\) for \(W=\rho_\pi\chi_1^{-1}\);
this avoids \(0,-1\). The Tate dual \(W^\vee(1)\) at
the \emph{same} embedding sends each remaining \(a\) to
\(-1-a\).

For a complementary shifted entry \(z=p+s\), the positive
column jump and the negative row's dual jump are
\[
     a_+=z+c,\qquad a_{\rm row}=-1-(-z+c)=z-c-1 .
\]
For the \(b_h\)-entries with \(h<j\) one has \(z\ge c+2\);
both jumps are then positive. For \(h\ge j\) one has
\(z\le-c-2\), and both are negative. The outer entries
\(z=s+v_a\) exceed \(c+1\). Thus each side has \(m-j\)
negative jumps per embedding.
Thus each contributes at most the degree times \(m\) to
the free finite cohomology bound, even at the last lift.
\end{proof}

Above the prime cutoff the stronger separation required in
paragraph~\ref{subsec:outside-windows} follows directly from
the parameter inequality
\( |b_h-e_{j-1}|\ge 2b/\ell_*+t/2+3\). On either side
this absolute value is \(|a-c|\) for \(a\in D_d\).
On the complementary block \(C_d\) the place windows instead give
\[
 d(a-c)=s+v_a,\qquad W_0\le s\le W_0+\eta_c+1.
\]

Here \(W_0=b\) on the part with larger weights, \(\omega\) on the part with smaller weights and \(b_v+\lfloor b'/x\rfloor\) at a designated top.
The floor accounts for the extra one at the top. Thus the
exceptional integers tested there use the actual jump sets of
the marked map. These weight inequalities belong to the same
parameter record as \(\ell_*\).
\label{cor:marked-localization-6}

\subsubsection*{The two integral maps and the local cases}

Lemma~\ref{lem:marked-signed-jumps} determines the ordered jumps on the two marked maps. Proposition~\ref{prop:marked-signature} supplies the length on the good complement. The prime range determines which map to use.
Fix the exceptional primes computed from the host and the
three weight windows, and put the selected character and its
two auxiliary twist parameters in the family before choosing
the matrices \(R_x,\rho_x\). Choose the integral period functors
for the completions of this host with their comparison
isomorphisms, as in Sections~\ref{subsec:cut-wach-models}
and~\ref{subsec:split-local-length}.  In the rational matrix product
over the characters put the fibre at \(\pi\) in the direct
sum basis for the string and its complement.

If this changes
basis by \(B\) over the fraction field, it carries every sum
\(\sum_g a_g R(g)\), with coefficients in the localized
Hecke order, to the sum with the same coefficients in the new
generators. Indeed
\[
 B\Big(\sum_g a_gR(g)\Big)B^{-1}
      =\sum_g a_g\,BR(g)B^{-1}.
\]
Applying the identity with \(B^{-1}\) gives the reverse inclusion of
the integral spans even if \(B\) is defined only over the fraction field. This transports the marked
idempotents, the ordered corner roots and the element
perturbing the cyclotomic corner. On the rank-\(r\) block
the integral span maps onto \(M_r(O)\): its image modulo the
maximal ideal is the whole matrix algebra by absolute residual
irreducibility, and the cokernel vanishes by Nakayama.
Section~\ref{subsec:split-selmer} carries this out with the
split Frobenius polynomials of the chosen localization.

Write \(T\) for the stable lattice of \(W\) in this direct sum.

\begin{lemma}[Lattices for the local row and column comparisons]
\label{lem:marked-row-column-lattices}
In the integral presentation the evaluation at the rank-one corner gives a surjection
\[
 e_1\mathcal A\twoheadrightarrow T^*=\operatorname{Hom}(T,O)(1).
\] The rational column
representations have constituents in sums of \(R_x\chi_1^{-1}\).
Their relations give finite classes on the rank-\(r\) target,
an unramified class on the first character line, and zero on
the other character lines.
\end{lemma}

\begin{proof}
\proofstep{Use the full matrix image for the character row.}
Write the rank-one projection as \(v\lambda\) with
\(v\in T\), \(\lambda\in\operatorname{Hom}_O(T,O)\), and
\(\lambda(v)=1\). For any \(\lambda'\in\operatorname{Hom}_O(T,O)\),
choose the matrix \(M=v\lambda'\) in the full matrix image.
The row of this product, evaluated on any \(x\in T\), is determined by
\[
 ((v\lambda)M)x=v\lambda\!\left(v\lambda'(x)\right)
                =v\lambda(v)\lambda'(x)=v\lambda'(x).
\]
Thus the image contains every integral row.
Put \(U_g=\chi_1(g)^{-1}R(g)\). The action on the row functional, including the Tate factor, is
\[
 g\cdot\lambda'=\epsilon(g)\lambda'\circ U_g^{-1}\quad\text{on }T^*.
\]
The source embeds rationally into a sum
of the dual-twisted family by recording all its matrix rows.
\proofstep{Use the direct-block projection for the column lifts.}
The row argument used the entire matrix image. On the conjugate column the finite classes are read from the projection onto the original complementary block.

For the column lifts one uses the left action on matrix
columns and projection to the direct block. Its rational
constituents belong to sums of \(R_x\chi_1^{-1}\).
The column relations give finite classes in the rank-\(r\)
block, an unramified class on the first character line and
zero on the other lines, as shown in
Section~\ref{subsec:split-selmer}.

For the character-row map
the extension with trivial quotient is integral: on
\(Y=\sum_g\mathcal A(U_g-1)E\) use the cocycle \((U_g-1)E\).
The map
\[
 (z,a)\mapsto(z+aE,a)
\]
embeds its rationalization
in the matrix columns and a trivial line. The latter already
occurs in the twisted string. These constructions justify the
passage of the jump bounds to the lifting maps by exactness
on crystalline subquotients. The modulus multiplying an
integral finite class is computed on the actual surjection.
For a finite extension \(O\subset O'\) of coefficient DVRs and a finite-length module \(M\),
\[
 \frac{\operatorname{length}_{O'}(M\otimes_O O')}{e(O')}=\frac{e(O'/O)\operatorname{length}_O(M)}{e(O'/O)e(O)}=\frac{\operatorname{length}_O(M)}{e(O)}.
\]
The normalized colength
returns to this presentation. \label{cor:marked-localization-7}
\end{proof}

Apply Lemma~\ref{lem:marked-row-column-lattices} to the
following cases on this same marked family.

\paragraph{Unramified split and inert places.}
For split places choose the side whose embedding sign makes
the direct column positive, and take the conjugates of these
places for its rows. If \(\mathcal R_\ell\) and \(\mathcal C_\ell\) denote the row and column places, at an inert prime these sets are
\[
 \mathcal R_\ell=\{v:v\mid\ell\},\qquad \mathcal C_\ell=\varnothing.
\]
The unramified Wach estimates apply to these surjections, with the order of Frobenius on the embeddings of \(E\) and the signed jump arrays above.
Sections~\ref{subsec:ordinary-split-length} and~\ref{subsec:inert-rows}
give the resulting split and inert allowances. Above the
cutoff, Paragraph~\ref{subsec:outside-windows} applies to
this same family with the exceptional set already specified.
Proposition~\ref{prop:offwindow-local-vanishing} kills the additional
local allowance outside these exceptions.

\paragraph{The split conditions at two and three.}
At two, and at three when \(3\nmid Q_S\), partition the places into columns and conjugate rows using the archimedean sign. One choice covers every good \(\pi\) at the localization and preserves the matrices, twist, support and character-row bound.
The column and row lifting exponents are proved for these crystalline surjections over
\(\mathbb Q_\ell(\mu_{\operatorname{lcm}(\ell,N)})\) and
returned by corestriction in Section~\ref{subsec:split-local-length}.
The preliminary choice of zero for an unused exponent supplies
no such lifting result. The character-row estimate and its dual use the respective heights
\[
 h=\lfloor2b\rfloor,\qquad h^\vee=h+1,\qquad n_0=4nR.
\]
The extra unit allows the Tate-dual shift.
The depth is at least two at \(2\). At \(3\mid Q_S\) use the
ramified odd-prime estimate of
Section~\ref{subsec:special-prime-lengths}.
This includes the torsion of finite integral cohomology when
its rational rank is zero. \label{cor:marked-localization-8}

\paragraph{The common denominator of the two type projectors.}
At a type prime apply the pair of finite series projectors
for the labels of \(\pi\), multiplied by their common group
denominator. If \(e_\pi\) denotes their product in the rational
finite group algebra and \(a_\pi=h_Ge_\pi\) its integral
multiple by the common group denominator \(h_G\), then
\[
       a_\pi|_{\operatorname{im} e_\pi}=h_G\,\mathrm{id},
        \qquad a_\pi|_{\ker e_\pi}=0 .
\]
Both projectors commute with the Hecke action since the
level actions and the split Hecke correspondences commute.
The remaining support has the two inertia labels of \(\pi\)
at its theta points. Such a point is good by the cut.

The
norm choice identifies the inertia characters on this
support with the labels of the same finite series, including
inversion on the conjugate side. Hence the common tame
lattice comparison of Section~\ref{subsec:special-prime-lengths}
applies on the support of this operator. Use
\[
 \eta=\lfloor\eta_c\rfloor,\ s_T=\lfloor3b\rfloor+2,\ S_T=0
\]
with the opposite signs at the paired places and \(C_T=40\).
The auxiliary prime bound there gives the strict size inequality
once \(b\ge C_T\). The pushout estimate then bounds the
finite-index defect for the rank-\(r\) target and each of
its character targets. Multiplying the resulting element
of the congruence order by that very projector clears its
complement. This gives the type-prime length on each graded
piece. \label{cor:marked-localization-9}

\paragraph{A threshold shared by the local length estimates.}
All these estimates may be used at a single late point.
The unramified row comparison first fixes its absolute
constant \(C_2\ge1\). The type-prime, small-prime, parameter
and exceptional-prime comparisons supply finitely many
further constants. Write \(A_1,\ldots,A_s\) for the further upper bounds after insertion of \(C_2\). Choose
\[
 C\ge\max\{C_2,A_1,\ldots,A_s\}.
\]

For
that bound take the maximum of the thresholds for the
integral setup, the size of the exceptional set, the
special primes, the large-prime vanishing, the ordinary
split and inert estimates, the parameter inequalities,
and the bounds at two and three. Each has already been
obtained for all point indices above its own threshold.
Only now take the point and evaluate its exceptional set
and integral models. An empty good locus has zero projected
lattice. At each other localization the rank-\(r\) term has
the signature bound above and the character terms have
the strict and class-number bounds of
Section~\ref{subsec:division-length}. This gives the local
table used in the weighted return on precisely these cuts.
\label{cor:marked-localization-10}

\subsection{The split local estimate at the conductor primes and at three}
\label{subsec:split-local-length}
\hypertarget{source-small-prime-lengths}{}

At the small split primes, the depth of the cyclotomic level enters the colength of the marked lattice. We keep this depth in the cohomology estimates before normalizing the length. Let \(\ell\) be
the coefficient prime at a retained localization and let
\(O\) be its DVR, with fraction field \(K_\ell\).
Write \(f,m,r,t,R\) for the host degree parameter, the number
of lifts, the rank of the lower block, the rank of the
character block, and the total rank. The parameter range is
\[
 1\le m\le f,\qquad 0\le r\le t,\qquad t,R\ge1.
\]

The CM field has degree \(2f\) over \(\mathbb Q\).
Use the level exponent \(s=\max(1,v_\ell(N))\) for the modulus defining the auxiliary field \(N=Q_Sqd\). At the small primes \(s\) is the
local level exponent in the row comparison. The same is true
at \(3\) when \(3\nmid Q_S\), with level exponent \(1\).

The primes in this comparison are
\[
 \mathcal S_{\rm sp}=\{\ell:\ell\mid Q_S\}\cup\{3\},\qquad
 Q_S=\operatorname{lcm}(2^{e+2},f_\psi).
\]
For \(\ell=2\), one has \(s\ge2\) and \(-d\equiv1\pmod8\).
An odd \(\ell\in\mathcal S_{\rm sp}\) is either 3 or a divisor of
the fixed conductor \(f_\psi\). Eventually every such prime is at
most \(w\) in Proposition~\ref{prop:quadratic-prime-pair}. Hence
\[
     -d\equiv1\pmod\ell\quad(\ell\text{ odd}),\qquad
     \mathbb Q(\sqrt{-d})\otimes\mathbb Q_\ell
          \simeq\mathbb Q_\ell\times\mathbb Q_\ell
                  \quad(\ell\in\mathcal S_{\rm sp}).
\]
Hensel's lemma gives the splitting at odd primes, and the dyadic
unit-square criterion gives it at two. Thus the row and column lie at conjugate places
even when \(K_\psi\) ramifies at \(\ell\). For each such completion
\(L\), put \(N'=N/\ell^{v_\ell(N)}\). Equation~\eqref{eq:host-character-compositum} gives
\[
    L\subseteq F_s=F_0(\mu_{\ell^s}),\qquad
    F_0=\mathbb Q_\ell(\mu_{N'}),\qquad
    F_0/\mathbb Q_\ell\text{ unramified},
\]
since \(\ell\nmid N'\) \cite[Corollary~10.16]{sutherland-local-fields}.
In particular \(s=1\) at three if \(3\nmid Q_S\).

Use the cohomology conditions on these completions. They include the finite-class lifting criterion, the bound on invariants over the cyclotomic subfield and the finite character criterion. The
split row and column constructions below invoke these
criteria on the restricted original lattices. In the
finite-class criterion the target class comes from
integral finite cohomology and the map of crystalline
lattices is surjective. In the character criterion of
Proposition~\ref{prop:cyclotomic-scalar-finite-condition},
the auxiliary extension is crystalline and has a
trivial rank-one quotient with a chosen integral lift
of its generator. \label{cor:small-prime-lengths-1}

\begin{lemma}[Invariants on the character string]\label{lem:small-prime-scalar-invariants}
At a place of \(E\) in the split case,
for \(1\le a<t\) write \(F\) for this completion. The normalized length satisfies
\[
 \operatorname{len}\bigl((K_\ell/O)(-a)^{G_F}\bigr)\le s+v_\ell(a)\le s+\lfloor\log_\ell t\rfloor ,
\]
uniformly in the twist and place.
\end{lemma}

\begin{proof}
For a place of \(E\) over \(\ell\), the prescribed cyclotomic
inertia image contains \(1+\ell^s\). On
\((K_\ell/O)(-a)\) with \(1\le a<t\), invariants lie in the
kernel of multiplication by
\((1+\ell^s)^a-1\). This kernel has normalized length
\[
     v_\ell((1+\ell^s)^a-1)=s+v_\ell(a).
\]
For the equality, the binomial formula gives
\(v_\ell((1+u)^b-1)=v_\ell(u)\) if \(\ell\nmid b\).

For a step of exponent \(\ell\), the first term \(\ell u\)
has valuation \(1+v_\ell(u)\). Each later term has greater
valuation when \(\ell\) is odd and \(v_\ell(u)\ge1\), or
when \(\ell=2\) and \(v_\ell(u)\ge2\).
Iterating through the \(\ell\)-part of \(a\) proves the
formula. The maximum invariant length over the indicated
twists and places is at most
\(s+\lfloor\log_\ell t\rfloor\). In particular the loss at two depends on its actual local depth. We retain that depth in passing from these invariant lengths to the split local conditions.
\end{proof}

\subsubsection*{The split row and column local conditions}

On the two sides of a conjugate split pair, the same good character supplies a column and a Tate-dual row. Their finite local conditions have different roles: the column provides classes to lift, and the row measures the colength of their images.

\begin{lemma}[Finite-class rank and colength on the small split row]\label{lem:small-split-finite-colength}
For every \(\ell\in\mathcal S_{\rm sp}\), with the split host
completions and cyclotomic depth \(s=\max(1,v_\ell(N))\)
of Section~\ref{subsec:split-local-length}, use the paired
places \(u,w\) of the same good character \(\pi\), the rank-\(r\) complement
\(W\) on the column and \(W^\vee(1)\) on the row.
The rational finite dimensions are at most
\([E_u:\mathbb Q_\ell]m\) on both sides. For the marked lattice
surjection, if the finite lift uses \(\ell^{e'}\) and the row
invariants are killed by \(\ell^{c_0(u)}\), the normalized row colength satisfies
\[
 \operatorname{colength}_{\rm row}\le [E_u:\mathbb Q_\ell]m e'+rc_0(u).
\]
\end{lemma}

\begin{proof}
\proofstep{Count the finite dimensions on the two signed jump lists.}
Use the same good Hecke character \(\pi\) on the retained
family. Write \(R_\pi^\circ\) for the full twisted
representation and \(W\) for its rank \(r\) constituent
after removing the direct sum of the \(t\) cyclotomic
characters. Write
\[
 T_2=W^\vee(1),
\]
using the convention
that the Tate twist by \(1\) has filtration jump \(-1\).
The local crystalline data give the following jumps,
with multiplicity one, for \(R_\pi^\circ\) at each
embedding of the completion into a splitting coefficient
field. The cyclotomic block has jumps \(0,\ldots,t-1\).

Of the other \(r\) jumps, \(k-1\) are in an interval of
strictly positive jumps at every embedding of a column
place and in an interval of strictly negative jumps at
every embedding of its conjugate row place. The remaining
\(m\) jumps can have either sign. These sets are disjoint,
and the two sets for \(W\) avoid both \(0\) and \(-1\).
Conjugation of the embedding type gives the switch of
sign on the row places. The local description of \(T_2\)
uses the following two exclusions at the same embedding:
\begin{equation}\label{eq:split-row-zero-jump}
 J(T_2)=\{-a-1:a\in J(W)\},\qquad
 0\notin J(W),\qquad -1\notin J(W)\ \Longrightarrow\ 0\notin J(T_2).
\end{equation}

At a column place \(u\), \(W\) is crystalline and
\(H^0(E_u,W)=0\), since a trivial subrepresentation
would contribute a zero jump. For a crystalline representation
\(V\) over a splitting coefficient field, the
Bloch--Kato local dimension formula is
\[
 \dim H^1_f(E_u,V)=
 \dim H^0(E_u,V)+
 \sum_{\tau:E_u\hookrightarrow K_\ell}
    \dim \bigl(D_{\mathrm{dR},\tau}(V)/
                       \operatorname{Fil}^0 D_{\mathrm{dR},\tau}(V)\bigr).
\]
Here and below one may enlarge the coefficient field to
contain all local embeddings. The filtration quotient
counts the negative jumps. Additivity of the crystalline
filtration under the split decomposition removes the
cyclotomic jumps from the full list. Thus
\(\dim H^1_f(E_u,W)\le [E_u:\mathbb Q_\ell]m\).

On the conjugate row place \(w\), \(T_2\) also has zero
rational invariants. Under the passage to the Tate dual the
\(k-1\) negative integer jumps of \(W\) become
nonnegative jumps of \(T_2\), by \(j\mapsto -j-1\).
Thus at most \(m\) jumps on each row embedding contribute:
\[
 \dim H^1_f(E_w,T_2)
   \le \sum_{\tau:E_w\hookrightarrow K_\ell} m
   =[E_w:\mathbb Q_\ell]m=[E_u:\mathbb Q_\ell]m .
\]
The comparison uses the Tate-dual row and its actual conjugate
column, even though their degrees agree.
\cite[Corollary 3.8.4]{bloch-kato}

\proofstep{Lift finite classes on the actual crystalline lattices and return by corestriction.}
For the integral condition, Proposition~\ref{prop:full-cyclotomic-finite-lift} gives an exponent \(e'\). Its input at a
column place is a surjection of stable lattices whose
rational representations are subquotients of finite
direct sums of the \(R_x^\circ\) in the family. At a row place the constituent family is instead obtained by the operation
\[
 R_x^\circ\longmapsto (R_x^\circ)^\vee(1).
\]

The family members are crystalline at these places,
and the absolute jumps of both families lie in the
bounded range used to choose \(e'\). Crystallinity and
the jump window pass to rational subquotients by
exactness and strictness of the crystalline filtration.
The local field is inside the prescribed cyclotomic
extension at level \(N\). The estimate says that the
cokernel of the map on integral \(H^1_f\) is killed
by \(\ell^{e'}\).

For this use of the Wach lattice estimate, the local
field \(L\) at a row or column is extended to
\[
 L'=\mathbb Q_\ell(\mu_{\operatorname{lcm}(\ell,N)}).
\]

The coefficients contain the corresponding roots of
unity. The extension \(L'/L\) is Galois and has
\(\ell\)-valuation of its degree at most
\(\lfloor\log_\ell N\rfloor+1\). Indeed
\[
 \ell^{v_\ell([L':L])}\le[L':L]\le[L':\mathbb Q_\ell]\le\ell N
 \quad\Longrightarrow\quad
 v_\ell([L':L])\le1+\lfloor\log_\ell N\rfloor .
\]
Restriction of a crystalline representation preserves
its jump window by base change at each embedding.
On \(L'\) the full cyclotomic level is \(s\).
The crystalline lattice comparison there gives the
lifting loss \(32(2hs+1)\) for jumps in \([-h,h]\).
The restriction--corestriction identity on the
finite local groups adds at most
\(\lfloor\log_\ell N\rfloor+1\) to the loss.

For \(c\in H^1_f(L,T)\) choose a source class \(d\)
over \(L'\) with \(f_*(d)=\ell^{32(2hs+1)}{\rm res}(c)\).
Naturality and restriction--corestriction give
\[
 f_*{\rm cor}(d)
   =\ell^{32(2hs+1)}{\rm cor}({\rm res}(c))
   =\ell^{32(2hs+1)}[L':L]c .
\] The prime-to-\(\ell\)
factor of the degree is a unit in the coefficients.
Thus the exponent \(e'\) in the conversion below
suffices. An invariant over \(L\) is already invariant
over \(L'\), so the invariant-killing argument uses
restriction alone.

\proofstep{Add the torsion length in the original coefficient normalization.}
The free rank is known; the integral row still contributes torsion.

Apply this at a row place to a lattice surjection
onto the target lattice \(T\) in \(T_2\) in its tied
matrix basis. The source sits rationally in a sum of the
dual-twisted family, and the target is therefore a
rational subquotient of that sum. The cokernel is a
quotient of \(H^1_f(E_w,T)/\ell^{e'}H^1_f(E_w,T)\).
This integral group is the preimage of the rational
finite part in \(H^1(E_w,T)\); by local cohomology
finiteness it is finitely generated. By the definition as an inverse image under rationalization,
\[
 \operatorname{rank}_O H^1_f(E_w,T)\le\dim_{K_\ell}H^1_f(E_w,T_2).
\]
The torsion lies in the image
of \(H^0(E_w,T_2/T)\) by the cohomology sequence.

If \(\ell^{c_0(u)}\) kills the torsion invariants, they embed into
\[
 \ell^{-c_0(u)}T/T\simeq T/\ell^{c_0(u)}T.
\]
The normalized length of the right side is \(r c_0(u)\).
Dividing lengths by the ramification degree of the
coefficient DVR gives the row colength bound
\[
       [E_u:\mathbb Q_\ell]m e' + r c_0(u).
\]
This computation gives the same bound on any of the
tied extensions of the coefficient field. Apply the degree-zero conclusion of
Lemma~\ref{lem:cyclotomic-invariant-loss} to \(W\) and \(T_2\).
The two exclusions in \eqref{eq:split-row-zero-jump} persist
under restriction to the full cyclotomic field, since a field
extension repeats each labelled list at the embeddings above it. The union of lists over
the embeddings of \(E_w\) has at most
\([E_w:\mathbb Q_\ell]r\) entries, so the bound uses the
degree of the base field and the rank, independently of
the degree of the splitting coefficient field.

The rank and colength comparison uses the crystalline lifts and invariant bounds for the specified local data.
The global lattice-index argument now uses these
local estimates on the same family.
\label{cor:small-prime-lengths-2}
\end{proof}

\paragraph{How the distinct jumps enter the local exponents.}
The two estimates involving the set of jumps use distinct
jump values. For a window \(U\subset [-h,h]\) with
\(|U|\le n\), put \(M=2s+n_0+1\), where \(n_0\) is least
with \(2h\le(\ell-1)\ell^{n_0}\), and put
\(\lambda=\lambda_N(h)=s+\lfloor\log_\ell(2h)\rfloor+
\lfloor\log_\ell N\rfloor+1\). The degree term is
\(D=v_\ell([F_M:F_s])=M-s\) on the full cyclotomic field.
The jump contribution to the local length is the positive integer
\begin{equation}\label{eq:local-jump-loss}
 C(k)=D+\sum_{i=0}^{k-1}\left(2+\max\left(1,
       \sum_{a\in U,\ a\ne i}(M+v_\ell(a-i))\right)\right),
 \qquad 1\le k\le h+1.
\end{equation}
Indeed \(n_0\le\lfloor\log_\ell(2h)\rfloor+1\), and
\(0<|a-i|\le2h\) in the inner sum. Hence
\(C(k)\le 2\lambda+3k(n+1)\lambda\).

For the invariant bound the fine loss is \(C(1)\),
so \(20(n+1)\lambda\) suffices. For a character of
jump \(e\), where \(0\le e<t\le h\), the full-field
specialization uses \(C(e+1)\) plus
\[
 (e+2)s+v_\ell(e)+\left\lfloor\frac{e}{\ell-1}\right\rfloor
       +1+s+v_\ell(e+1),
\]
with the zero term \(v_\ell(e)\) interpreted as \(0\) at
\(e=0\). The restriction--corestriction loss at the smaller
field is at most \(\lfloor\log_\ell N\rfloor+1\).
Together these are bounded by \(40(t+1)(n+1)\lambda\),
since \(e+1\le t\), \(s\le\lambda\) and each valuation
at a positive integer is at most that integer.
 It uses the
fine lifting and invariant losses of the local lattice
comparison. \label{cor:small-prime-lengths-3}

\subsubsection*{The integral Hodge--Tate estimate at full cyclotomic level}
\label{subsec:ramified-ht}
\hypertarget{source-ramified-cyclotomic}{}

The small-prime comparison needs an integral bound for invariants. We first pass to the field containing all roots of unity at the chosen cyclotomic level and keep track of the power of the coefficient prime needed to return to the original lattice.

Let \(F_0/\mathbb Q_\ell\) be finite unramified and
\(F_j=F_0(\mu_{\ell^j})\). The local lattices in this
argument are first restricted to \(F_s\), with \(s\ge1\)
and \(s\ge2\) at two. Include \(F_s\) in the finite
coefficient field. Normalize \(v(\ell)=1\) and let
\(L\subset V\) be the original stable lattice.
Write its jump interval as \([-a,b]\), and choose the depth by
\[
 H=a+b,\qquad n_0=\min\{n\ge0:H\le(\ell-1)\ell^n\},\qquad M=2s+n_0+1.
\]

\begin{proposition}[Integral cyclotomic comparison at the full ramified level]
\label{prop:ramified-integral-wach-comparison}\label{hyp:ramified-integral-wach-comparison}
Put \(S=\mathcal O_{F_0}[[\pi]]\), \(X=(1+\pi)^{\ell^{s-1}}-1\), \(Q=\varphi(X)/X\), and \(\Gamma_s=\operatorname{Gal}(F_\infty/F_s)\), where \(F_\infty=\bigcup_jF_j\). Let \(T\) be a crystalline lattice over \(F_s\) with jumps in \([-a,b]\). Its cyclotomic module \(N=N(T)\) is finite free over \(S\), is functorial in \(T\), and generates the associated integral étale module. For a generator \(\gamma\) of \(\Gamma_s\),
\[
 (\gamma-1)N\subset XN,\qquad
 Q^a\varphi^*N\subset N,\qquad Q^bN\subset\varphi^*N.
\]
Here \(\varphi^*N\) denotes the image of linearized Frobenius.

Inside the common period space the corresponding bounds are
\[
  X^b(A_{\inf}\otimes T)\subset A_{\inf}N
                     \subset X^{-a}(A_{\inf}\otimes T).
\]
At every root of \(Q\), write \(t=\log(1+\pi)\) for a local uniformizer. In a basis \((v_j)_j\) adapted to the filtration, indexed with multiplicity by its jumps, the rational completed lattice is \(\bigoplus_j F'[[t]]t^{-j}v_j\) over a splitting field \(F'\). The module of a kernel is the kernel on integral modules.
\end{proposition}

\begin{proof}
\proofstep{Specialize the Kisin--Ren classification.}
Use the multiplicative Lubin--Tate group over \(L=\mathbb Q_\ell\) with uniformizer \(\ell\). Its \(\ell^s\)-division points are \(\zeta-1\) for \(\zeta^{\ell^s}=1\), and \(K=F_s\) is contained in \(F_0L_\infty\). Thus one may take \(m=s\) in the construction of Kisin and Ren. Their three elements are
\[
 [\ell](\pi)=(1+\pi)^\ell-1,\quad
 P=[\ell^{s-1}](\pi)=X,\quad
 [\ell^s](\pi)/[\ell^{s-1}](\pi)=Q .
\]
Since \(K\otimes_{\mathbb Q_\ell}L=K\), their crystalline condition imposes no further filtration condition. The Lie algebra condition of \(L\)-analyticity is automatic for \(L=\mathbb Q_\ell\). Their integral equivalence therefore applies to the underlying \(\mathbb Z_\ell\)-lattices of all the representations here. We use the covariant module, equivalently the contravariant construction on the dual lattice. \cite[Sections~2.1, 2.4 and Corollary~3.3.8]{kisin-ren-lubin-tate}

Any extra coefficient action comes along by functoriality. For coefficient integers \(O\) containing \(\mathcal O_{F_0}\), the unramified embeddings give
\[
 O\otimes_{\mathbb Z_\ell} S\simeq
   \prod_{\tau:F_0\hookrightarrow\operatorname{Frac}O} O[[\pi]] .
\] On each summand, reduction of \(N\) modulo \(\pi\) is \(\ell\)-torsion-free and hence free over \(O\). Lift a basis. It generates by Nakayama's lemma and is independent by successive reduction modulo \(\pi\), since \(N\) has no \(\pi\)-torsion.

\proofstep{Read the filtration and the two integral bounds.}
The filtered-module construction makes modifications at the roots of \(Q,\varphi Q,\varphi^2 Q,\ldots\). At a root of \(Q\) the defining completed lattice is the degree-zero filtration on \(F'((t))\otimes D_{\rm dR}\); the preceding root reached by Frobenius is unmodified. An adapted basis therefore gives
\[
   N^{\wedge}=\bigoplus_j F'[[t]]t^{-j}v_j,\qquad
   (\varphi^*N)^{\wedge}=\bigoplus_j F'[[t]]v_j .
\]
Here \(Q\) has a simple zero, so these relative exponents give both height bounds after inverting \(\ell\). Integrally, the Frobenius matrix and its inverse already have entries in \(S[1/Q]\). The required products by powers of \(Q\) have no poles at any factor of \(Q\) by the calculation just made. Hence they lie in \(S\): division by the distinguished monic polynomial \(Q^k\) has an integral remainder, and a remainder vanishing after inverting \(\ell\) is zero. \cite[Section~2.2, Lemmas~2.2.1 and~2.2.5]{kisin-ren-lubin-tate}

The lower height bound gives
\[
      \varphi^*(X^a N)=X^aQ^a\varphi^*N\subseteq X^aN .
\] Lemma~3.2.1 of Kisin and Ren evaluates every integral functional of its étale lattice on this module in their integral period ring \(S^{\rm ur}\subset A_{\inf}\). Multiplication by \(X\) is invertible in the étale coefficient ring. Evaluating on a basis of \(T^\vee\) thus gives \(X^aN\subset A_{\inf}\otimes T\). The dual module has jumps in \([-b,a]\); the identical argument yields
\[
     X^bN^\vee\subset A_{\inf}\otimes T^\vee
        \quad\Longrightarrow\quad
     X^b(A_{\inf}\otimes T)\subset A_{\inf}N .
\]
The implication uses the perfect evaluation pairing for a finite free module. \cite[Lemma~3.2.1 and Corollary~3.3.8]{kisin-ren-lubin-tate}

\proofstep{Pass from the parameter \(\pi\) to the full divisor \(X\).}
The classification gives trivial action modulo \(\pi\). Consider any \(z=\zeta-1\) with \(\zeta^{\ell^{s-1}}=1\). The group \(\Gamma_s\) fixes \(z\), and \(Q(z)=\ell\ne0\). Over the fraction field, Frobenius identifies the fibre at \(z\) with the coefficient twist of the fibre at \(\zeta^\ell-1\). Frobenius commutes with \(\gamma\). Iteration from the fibre at zero shows that \(\gamma\) acts trivially on all these fibres. They are the distinct roots of \(X\), so monic division gives
\[
                         (\gamma-1)N\subset XN.
\]

For a surjection \(T_1\to T_2\) with kernel \(T'\), the rational modified lattices are exact. Away from their modifications this follows from exactness of the crystalline modules. At each modification the completed lattices are the filtered modules above, and the induced map is strictly surjective for the filtration. This also proves exactness on \(S[1/\ell]\): any nonunit irreducible there is represented by a distinguished polynomial and is tested in the corresponding local completion.

The integral module for \(T'\) is the intersection of its étale lattice with its rational module, as in the proof of Corollary~3.3.8 of Kisin and Ren. Combining rational and étale exactness yields
\[
                         N(T')=\ker(N(T_1)\longrightarrow N(T_2)).
\]
The integral cokernel \(C\) is killed by a power of \(\ell\) by rational surjectivity. After reduction modulo \(\ell\) it vanishes upon inverting \(\pi\), by étale surjectivity. Hence \(\pi^d C\subset\ell C\) for some \(d\); iteration shows that some power of \(\pi\) kills \(C\). Section~\ref{subsec:ramified-lift} removes the dependence on this preliminary power. \cite[Proposition~2.2.6 and Corollary~3.3.8]{kisin-ren-lubin-tate}
\end{proof}

\begin{lemma}[Finite classes and restricted Herr pairs]\label{lem:ramified-restricted-herr}
For \(T\) as in Proposition~\ref{prop:ramified-integral-wach-comparison} with jumps in \([-h,h]\), put \(N=N(T)\). Every cocycle in
\[
   Q^{-h}N\oplus XN,\qquad
   d(y,z)=(\gamma-1)y-(\varphi-1)z ,
\]
defines an integral class with crystalline rational extension. After inverting \(\ell\), these pairs represent all finite classes, and two such pairs have the same class precisely when their difference is \(((\varphi-1)w,(\gamma-1)w)\) with \(w\in N[1/\ell]\).
\end{lemma}

\begin{proof}
Given a cocycle \((y,z)\), adjoin one generator \(v\) to \(N\). Define its Frobenius and cyclotomic action by
\[
      \varphi(v)=v+y,\qquad\gamma(v)=v+z .
\]
The cocycle equation says exactly that the operators commute. Write \(\Phi_N\) for the Frobenius matrix on \(N\). After inverting \(Q\), the linearized Frobenius has inverse
\[
   \begin{pmatrix}\Phi_N & y\\ 0 & 1\end{pmatrix}^{-1}
     =\begin{pmatrix}\Phi_N^{-1}&-\Phi_N^{-1}y\\0&1\end{pmatrix}.
\] The cyclotomic generator acts trivially modulo \(\pi\) because \(z\in XN\). Its powers extend continuously to \(\Gamma_s\simeq\mathbb Z_\ell\). Indeed, \(\gamma\) acts trivially on \(I/I^2\) and on the module modulo \(I\), where \(I=(\ell,\pi)\). For \(M=N\oplus Sv\) this gives
\[
  (\gamma-1)I^kM\subset I^{k+1}M,\qquad
  \lim_{n\to\infty}\gamma^{\ell^n}=1\quad\text{on each }M/I^kM
\]
by the binomial expansion.

This is an object of the Kisin--Ren category for \(L=\mathbb Q_\ell\); its associated lattice is crystalline after tensoring with \(\mathbb Q_\ell\). In the étale module the two differences of \(v\) are its usual Herr representative. For the new lift \(v+w\) its pair is
\[
    (y,z)+\bigl((\varphi-1)w,(\gamma-1)w\bigr).
\] \cite[Theorem~1.6 and Corollary~3.3.8]{kisin-ren-lubin-tate}

Conversely, a rational finite class is a crystalline extension of the trivial representation by \(T[1/\ell]\). It admits a stable lattice extension after scaling the class: a continuous cocycle on the compact Galois group has bounded image. On the corresponding cyclotomic modules the surjection to the trivial module is rationally onto by Proposition~\ref{prop:ramified-integral-wach-comparison}. The kernel identity and the bounds on the extension module put the differences of any such lift \(v\) in
\[
  (\varphi-1)v\in Q^{-h}N[1/\ell],\qquad
  (\gamma-1)v\in XN[1/\ell].
\]

Indeed, if \(N'\) is the extension module, multiplication by \(Q^h\) reduces the first intersection to the kernel identity. For the second, \(Xn'\) in the kernel implies that the image of \(n'\) in the free trivial quotient is zero. Thus the intersections inside the rational étale module are
\[
 Q^{-h}N'[1/\ell]\cap D_{\rm et}(T[1/\ell])=Q^{-h}N[1/\ell],
 \qquad XN'[1/\ell]\cap D_{\rm et}(T[1/\ell])=XN[1/\ell].
\]
Here \(D_{\rm et}\) denotes the cyclotomic étale module.

An equality of rational extension classes gives an isomorphism of extensions inducing the identity on kernel and quotient. Full faithfulness of the cyclotomic equivalence carries it back to the modules. It changes the chosen lift by a single \(w\in N[1/\ell]\), giving the asserted description of coboundaries. \cite[Corollaries~2.4.4 and~3.3.8]{kisin-ren-lubin-tate}
\end{proof}

At the full level, Proposition~\ref{prop:ramified-integral-wach-comparison} gives the integral Hodge--Tate bounds below. Lemma~\ref{lem:ramified-restricted-herr} will enter when lifting finite classes across a surjection.\label{cor:ramified-cyclotomic-1}

\begin{proposition}[Integral Hodge--Tate coordinates at full cyclotomic level]\label{prop:full-level-ht-coordinates}
With the module of Proposition~\ref{prop:ramified-integral-wach-comparison}, over \(F_s=F_0(\mu_{\ell^s})\), with \(F_0/\mathbb Q_\ell\)
unramified and \(s\ge2\) at two, for each embedding \(\tau\)
there are coordinates with image lattice \(B'\) and semilinear generator matrix \(U\) at level \(M\) such that
\[
 B_r\subseteq B'\subseteq\ell^{-1}B_r,\qquad v((U-1)_{ij})>1/(\ell-1)\quad\text{for all }i,j.
\] A vector of the rational
space whose image belongs to \(B'\) lies in the original
stable lattice, and the graded de Rham comparisons identify
the twist coordinates with the Hodge--Tate jumps at \(\tau\).
\end{proposition}

\begin{proof}
\proofstep{Specialize the two height inclusions in the same Wach basis.}
In the Wach comparison the height sandwich is
\[
 X^b (A_{\inf}\otimes L)\ \subseteq\ A_{\inf}N
    \ \subseteq\ X^{-a}(A_{\inf}\otimes L),
 \qquad X=(1+\pi)^{\ell^{s-1}}-1.
\]
Here the coefficient tensor products and the images
in the common period module are understood.
Use the Wach basis after multiplying its columns
by \(X^a\), so they can be read in the \(A_{\inf}\)
lattice.

Choose compatible primitive roots with \(\zeta_{\ell^0}=1\)
and put \(\epsilon=(1,\zeta_\ell,\zeta_{\ell^2},\ldots)\) and
\(\pi=[\epsilon]-1\). For every \(k\ge0\), Fontaine's map
\(\theta\) and the Witt Frobenius give
\[
  (\theta\circ\varphi^{-k})([\epsilon])
       =(\epsilon^{1/\ell^k})^\sharp=\zeta_{\ell^k}.
\]
The sharp map takes the limit of the \(\ell^j\)-th powers of
the lifted \(j\)-th components. Here every such power equals
\(\zeta_{\ell^k}\). We specialize at the absolute depth
\(d_0=s+n_0\). Thus
\begin{equation}\label{eq:full-level-fontaine-map}
 \pi\longmapsto\zeta_{\ell^{d_0}}-1,\qquad
 X\longmapsto c=\zeta_{\ell^{n_0+1}}-1,\qquad
 v(c)=\frac{1}{(\ell-1)\ell^{n_0}} .
\end{equation}

On \(\mathcal O_{F_0}\) the map acts by inverse unramified
Frobenius to power \(d_0\). It permutes the unramified embeddings.
In an integral \(\mathbb Z_\ell\)-basis the action of the extra coefficient
integers is by endomorphisms of the underlying Wach module.
Their matrix entries lie in \(\mathcal O_{F_0}[[\pi]]\).
Evaluation therefore puts these entries in
\(\mathcal O_{F_{d_0}}\subset\mathcal O_{F_M}\).
The change of basis to the original period lattice itself can
have entries in \(\mathcal O_{\mathbb C_\ell}\).

Taking the Wach basis columns times a
\(\mathbb Z_\ell\)-basis of the coefficient integers gives
a square matrix \(D\) on the underlying \(\mathbb Q_\ell\)
coordinates. The specialized sandwich puts its ball
image between \(c^H B_d\) and \(B_d\), where
\(B_d=\mathcal O_{\mathbb C_\ell}^d\).
Since
\[
 v(\ell)-v(c^H)=1-\frac{H}{(\ell-1)\ell^{n_0}}\ge0,
\]
it follows that \(D\) is invertible and
\[
                       \ell B_d\subseteq DB_d\subseteq B_d.
\]
The action and the coefficient multiplication in this
basis have matrices over \(\mathcal O_{F_M}\).
Indeed their columns are the specialized
Wach-coordinate expressions in the same basis.
The multiplication by \(X^a\) cancels in the action
comparison because \(c\) is fixed at level \(M\).
On the cyclotomic kernel the matrix part of the
action is the identity, by the \'etale--Wach
comparison. \label{cor:ramified-cyclotomic-2}

\proofstep{Increase the depth until the generator is strictly small in the integral model.}
For the generator estimate, start with a
cyclotomic generator at depth \(s\), with action
\(\gamma_0\), and set \(I=(\ell,\pi)\) in the
Wach coordinate ring. On \(S/(\pi)=\mathcal O_{F_0}\) the cyclotomic
generator acts trivially. Its action on \(N/\pi N\) is trivial
by \((\gamma_0-1)N\subset XN\) in
Proposition~\ref{prop:ramified-integral-wach-comparison}, and \(\gamma_0(\pi)/\pi\equiv1\pmod I\).
The Leibniz rule for a semilinear difference gives
\[
   (\gamma_0-1)(I^j N)\subset I^{j+1}N
           \qquad(j\ge 0).
\]
Let \(k=s+n_0+1\). Expand the operator on \(N\):
\[
  \gamma_0^{\ell^k}-1
       =\sum_{j=1}^{\ell^k}\binom{\ell^k}{j}(\gamma_0-1)^j .
\]
For \(j<\ell^{k-1}\), \(v_\ell(j)\le k-2\), so the small-index binomial terms satisfy
\[
 v_\ell\binom{\ell^k}{j}
 =k-v_\ell(j)+v_\ell\binom{\ell^k-1}{j-1}\ge2
 \qquad(1\le j<\ell^{k-1}).
\]

For \(j\ge\ell^{k-1}\), the
operator \((\gamma_0-1)^j\) takes values in
\(I^{\ell^{k-1}}N\). By \eqref{eq:full-level-fontaine-map}, the least valuation of a
generator of the specialized ideal \(I\) is
\(1/((\ell-1)\ell^{s+n_0-1})\). Since \(k-1=s+n_0\),
the valuations for the two ranges are respectively at least
\[
     2,\qquad
     \frac{\ell^{k-1}}{(\ell-1)\ell^{s+n_0-1}}
                =\frac{\ell}{\ell-1}.
\]
In both ranges the strict logarithm threshold holds because
\[
 \min\{2,\ell/(\ell-1)\}>1/(\ell-1).
\]

Consequently \(\gamma=\gamma_0^{\ell^k}\) at level \(M\)
has action matrix \(U\) with this entrywise bound on \(U-1\).
For \(u=\chi(\gamma_0)\) and a nonzero integer \(j\), the generator depth is computed by
\[
 v(u^j-1)=v(u-1)+v_\ell(j).
\]
For \(j\) prime to \(\ell\)
this follows by expanding the binomial and for
multiplication of \(j\) by \(\ell\) the valuation
increases by one. The latter step requires
\(v(u-1)>1/(\ell-1)\), which holds for the admitted
depths. Negative \(j\) give the same valuation by
inversion.

\proofstep{Project by the separability idempotent and choose a maximal minor.}
Fix a field embedding \(\tau:F_s\to L_{\rm coeff}\).
The separability idempotent \(E_\tau\) of
\(F_s\otimes_{\mathbb Q_\ell}\mathbb C_\ell\), acting
through \(\tau\) on the coefficients and the standard
embedding on the scalar factor, selects the
\(\tau\)-component. The conjugated idempotent
\(D^{-1}E_\tau D\) has entries in \(F_M\).
Write
\[
 E'=D^{-1}E_\tau D.
\]
Select a nonzero
maximal-rank minor of largest absolute value.
Multiply the selected rows by the inverse of
that minor's matrix to obtain a row map \(R\).
Its entries are integral by the minor ratio rule.

For the selected column inclusion \(S_0\), we have
\(RS_0=1\), \(RE'=R\), and \(E'=E'S_0R\).
Thus \(S_0\) is an integral section of \(R\).
The induced action on \(P=RD^{-1}\) is computed by
\(R U_g S_0\). This is well defined since \(\ker R=\ker E'\)
is stable and \(S_0\) lifts every coordinate.
It gives integral matrices and preserves the strict
bound for \(U_\gamma-1\).
In these coordinates the image \(B'=P B_d\) satisfies
\[
                         B_r\subseteq B'\subseteq \ell^{-1}B_r.
\]
The first inclusion follows from \(D B_d\subseteq B_d\)
and the integral section \(S_0\), and the second follows
from \(\ell D^{-1} B_d\subseteq B_d\).

To detect integrality on \(V\), extend \(\tau^{-1}\)
to an embedding of the coefficient field in
\(\mathbb C_\ell\). Evaluation in this embedding sends
\(B_d\) into the unit ball of the coefficient
coordinates and factors through \(P\) since it
factors through \(E_\tau\).
On a rational vector these evaluations have the
valuations of its coordinates in an integral
coefficient basis of \(L\). For \(v=\sum_i v_i e_i\) in this basis, the integrality test reads
\[
 Pv\in B'\quad\Longrightarrow\quad v_\ell(v_i)\ge0\ (\text{all }i)\quad\Longrightarrow\quad v\in L.
\]
Finally the graded de Rham
comparison, using the same separability idempotent
on both factors, gives a basis of the coordinate
space of vectors fixed after twists by jumps at \(\tau\).
Indeed adapted filtered vectors give their symbols in
\(\operatorname{gr} B_{\rm dR}\), and the projection
selects the \(\tau\)-part of each graded piece.
These are the Hodge--Tate coordinates of the
specialized lattice. \label{cor:ramified-cyclotomic-3}
\end{proof}

\begin{lemma}[Cohomology and invariant loss at full cyclotomic level]\label{lem:cyclotomic-invariant-loss}
In the cyclotomic comparison of Proposition~\ref{prop:full-level-ht-coordinates}, put
\[
 N_i=\max\left(1,\sum_{a\in J\setminus\{i\}}
                    (M+v_\ell(a-i))\right)
\]
for the indicated distinct jumps \(J\) in the marked comparison.
Let \(B(i)\) be the integral coordinate ball with twisted action,
and keep the projected jump set \(J\) at the chosen embedding.
The degree-one bound holds for every \(i\). For the second bound
take an embedding giving the integrality test of
Proposition~\ref{prop:full-level-ht-coordinates} with \(0\notin J\).
Then, for the original lattice \(L\subset V\),
\[
 \ell^{1+N_i} H^1(G_{F_M},B(i))_{\rm tors}=0,\qquad
 0\notin J\ \Longrightarrow\ \ell^{2+N_0}(V/L)^{G_{F_M}}=0.
\]
These exponents are independent of the coefficient degree.
\end{lemma}

\begin{proof}
\proofstep{Use the integral trace on the cyclotomic kernel.}
Set \(H_\infty=G_{\bigcup_jF_j}\).
For a continuous cocycle of \(H_\infty\) in the scalar
ball \(\mathcal O_{\mathbb C_\ell}\), multiplication by
\(\ell\) gives an integral coboundary. In fact Tate's
trace theorem for a finite extension of the
cyclotomic tower puts every element of positive
valuation in the integral trace image. To approximate
the coboundary, choose an open normal subgroup on
which the cocycle \(f\) has norm at most \(\varepsilon\).
Take an integral \(b\) of trace \(\ell\) in the fixed
field of this subgroup. For a set \(S\) of coset
representatives put \(x=-\sum_{\sigma\in S}\sigma(b)f(\sigma)\).
In the difference \(hx-x\), replace each \(h\sigma\)
by its representative times an element of the subgroup.

If \(h\sigma=\tau n_\sigma\) with \(\tau\in S\), the
cocycle identity gives
\[
 h f(\sigma)=f(\tau)+ \tau f(n_\sigma)-f(h).
\]
Since \(n_\sigma(b)=b\) and the
\(\tau\)'s permute \(S\), we obtain
\[
       hx-x=\ell f(h)-\sum_{\sigma\in S}
                       \tau(b)\tau f(n_\sigma).
\]
The last sum has norm at most \(\varepsilon\) by the
ultrametric inequality and integrality of \(b\), uniformly
in \(h\).
Also \(\|x\|\le\sup_h\|f(h)\|\).
Apply this construction successively to the residual
cocycle divided by \(\ell\), with the error bounds
tending to zero geometrically. The successive corrections
then converge in the ball to the integral cobounding
vector. Apply the construction coordinatewise on
\(H_\infty\), where the matrix action and the
cyclotomic twist are trivial. Tate's fixed-field
comparison gives
\[
 \mathbb C_\ell^{H_\infty}=\widehat{F_\infty}.
\]

\cite[\S3, Propositions 9--10]{tate-pdivisible}

\proofstep{Split off the base by the compatible normalized trace.}
Let \(A_g=\chi(g)^iU_g g\) denote the semilinear
action in twist \(i\). On the completed integers of
\(F_\infty\), split off \(\mathcal O_{F_M}\) by the
normalized trace. At finite level \(j\ge M\) the
basis \(\zeta_{\ell^j}^{a}\) for
\(0\le a<\ell^{j-M}\) is integral over
\(\mathcal O_{F_M}\) by the Eisenstein uniformizer
description, with normalized trace the constant
coordinate. The integral projections are compatible
as \(j\) grows. Write \(Y\) for the kernel of the extended normalized trace. Compatibility gives the integral splitting
\[
 \widehat{\mathcal O}_{F_\infty}=\mathcal O_{F_M}\oplus Y.
\]
Both this module and the base are
stable under \(A_\gamma\).

On \(Y\), \(\gamma-1\) is invertible after inverting
\(\ell\), with inverse taking \((\zeta_\ell-1)Y\)
into \(Y\). To see the uniform bound, at each finite
level the nonzero basis exponents are permuted by
multiplication by \(u=\chi(\gamma)\) modulo
\(\ell^{j-M}\), with multipliers in the base roots of
unity. If \(v_\ell(a)=b_0\), the length of its orbit
is \(d_a=\ell^{\max(0,j-2M-b_0)}\). The return multiplier
is \(\zeta_{\ell^j}^{a(u^{d_a}-1)}\), which is nontrivial
since
\[
 v_\ell(a(u^{d_a}-1))=b_0+M+\max(0,j-2M-b_0)=\max(M+b_0,j-M)<j.
\]
Let \(c_a=\zeta_{\ell^j}^{a(u^{d_a}-1)}\) be this
return multiplier, whose difference from one has valuation
at most \(1/(\ell-1)\).

On the orbit module,
\[
   (\gamma-1)\sum_{i=0}^{d_a-1}\gamma^i
                          =\gamma^{d_a}-1=(c_a-1)\mathrm{id}.
\]
The sum of powers is integral, so the inverse obtained by
dividing by \(c_a-1\) sends \((\zeta_\ell-1)\) times
the orbit module into the orbit module. This proves the
integral bound and the inverses agree on the nested
finite trace-zero spaces. Since they commute with
the base uniformizer, they extend through completion
with the same integral bound. Now
\(u^iU_\gamma-1\) has strictly higher valuation than
\(\zeta_\ell-1\) entrywise. On this summand write \(D=\gamma-1\) and \(E=(A_\gamma-1)-D\).
The denominator of \(D^{-1}\) costs at most the valuation of
\(\zeta_\ell-1\), whereas \(E\) has strictly higher valuation. Hence
\(D^{-1}E\) raises the base-uniformizer filtration and
\[
 (A_\gamma-1)^{-1}=\sum_{n\ge0}(-D^{-1}E)^nD^{-1}
\]
converges in that filtration. In particular
\((A_\gamma-1)^{-1}\ell Y^r\subseteq Y^r\).

The Hodge--Tate basis vectors have coordinates fixed
by \(H_\infty\) and are fixed by \(A_\gamma\)
in their individual twists. By the trace-zero
invertibility just proved, their coordinates are
in the base field \(F_M\). Thus \(U_\gamma\) has
an annihilating polynomial with the distinct roots
\(u^{-a}\) for the set \(J\) of their jumps.
For twist \(i\), let
\[
 Q_i(Z)=\prod_{a\in J\setminus\{i\}}(Z-u^{-a}),\qquad
 N_i=\max\left(1,\sum_{a\in J\setminus\{i\}}
                   (M+v_\ell(a-i))\right).
\]
If \(w\) is integral on the base and
\(\ell^k w=(u^iU_\gamma-1)w'\) there, then
\(Q_i(U_\gamma)w=0\).
Integral polynomial division gives
\[
 Q_i(Z)-Q_i(u^{-i})=(Z-u^{-i})D_i(Z),\qquad D_i\in \mathcal O_{F_M}[Z].
\]
Substitute \(Z=U_\gamma\) and apply to \(w\). Since \(u^i\) is a unit, this gives
\[
 Q_i(u^{-i})w\in(A_\gamma-1)B(i),
\] where \(B(i)\) denotes the integral ball with the twisted action. The valuation of this scalar is the sum in \(N_i\). Together with the trace-zero
bound, this shows that \(\ell^{N_i}\) kills the
torsion in the integral generator coinvariants.

\proofstep{Return to integral cohomology, then kill invariants modulo the original lattice.}
The coinvariant bound is now used to clear denominators of invariant vectors.

For an integral cocycle \(f\) of \(G_{F_M}\) whose
class is \(\ell\)-power torsion, subtract the
integral coboundary from \(\ell f\) on \(H_\infty\).
The result descends to the procyclic quotient, with
values fixed by \(H_\infty\). Its torsion cobounding
vector is fixed there as well, since the cocycle
vanishes on the kernel. The generator calculation
makes the descended value at \(\gamma\) an integral
coboundary after multiplication by \(\ell^{N_i}\).
Equality at the generator extends by the cocycle
law and continuity. Restriction to the kernel and the generator calculation together give
\[
 \ell^{1+N_i}H^1(G_{F_M},B(i))_{\rm tors}=0.
\]

\label{cor:ramified-cyclotomic-4}

For an invariant class \(v\bmod L\), write \(y=Pv\).
The cocycle \(g\mapsto\ell(gy-y)\) in twist zero is
integral in the ball and a sufficiently large
\(\ell\)-power makes it an integral coboundary.
Applying the bound yields an integral vector whose
coboundary is that of \(\ell^{2+N_0}y\).
When the jumps avoid zero, the difference vanishes
by Tate's nonzero-twist invariant theorem over the
finite local field \(F_M\).
Thus
\[
 \ell^{2+N_0}y\in B_r\subseteq B'
\]
and \(\ell^{2+N_0}v\in L\).
This proves the invariant loss used in the row
estimate with no dependence on the coefficient degree.
\cite[Theorem 2]{tate-pdivisible}
\end{proof}


\subsubsection*{Finite-class lifting at full cyclotomic level}
\label{subsec:ramified-lift}

The preceding invariant bound treats a single lattice. Lifting finite cohomology classes across a surjection requires an additional estimate. We make the lift at the same full cyclotomic level before applying corestriction.

\begin{established}[Full cyclotomic field and crystalline height]\label{hyp:full-cyclotomic-height}
Let \(F_0/\mathbb Q_\ell\) be finite unramified, and set
\(F_s=F_0(\mu_{\ell^s})\) for \(s\ge1\), with \(s\ge2\)
at \(\ell=2\). The coefficient field contains \(F_s\),
as in Section~\ref{subsec:ramified-ht}. Let \(T_1\to T_2\)
be a surjection of integral crystalline lattices with
source jumps in \([-h,h]\) at every embedding. Put
\(H=2h\). For every embedding \(\tau\) and every generator \(\gamma\) of the full cyclotomic tower,
\[
 \operatorname{HT}_\tau(\ker f)\cup\operatorname{HT}_\tau(T_2)
 \subseteq[-h,h],\qquad v_\ell(\chi(\gamma)-1)=s,
 \qquad f:T_1\twoheadrightarrow T_2 .
\]
The integral comparison holds for these lattices.
\end{established}
\begin{proof}
The filtered de Rham functor is exact with strict maps on crystalline representations. Thus the graded pieces of the kernel and quotient vanish outside the source window. Since \(F_0\) is unramified and the cyclotomic extension of \(\mathbb Q_\ell\) is totally ramified, restriction identifies the group above \(F_s\) with \(1+\ell^s\mathbb Z_\ell\). In the indicated range, including \(s\ge2\) at two, the logarithm takes this group isomorphically onto \(\ell^s\mathbb Z_\ell\) and preserves valuations. Any generator therefore has precisely the asserted valuation. Proposition~\ref{prop:ramified-integral-wach-comparison} and Lemma~\ref{lem:ramified-restricted-herr} identify the fixed lattices over this full cyclotomic base. In their notation
\[
 X=(1+\pi)^{\ell^{s-1}}-1,\qquad Q=\varphi(X)/X .
\]
\label{cor:full-cyclotomic-lifting-1}
\end{proof}

\begin{proposition}[Finite lifting over the full cyclotomic field]\label{prop:full-cyclotomic-finite-lift}
For the crystalline surjection in Lemma~\ref{hyp:full-cyclotomic-height}, write \(f_*\) for its map on integral finite cohomology. Uniformly in the coefficient degree,
\[
 \ell^{32(2hs+1)}\operatorname{coker}(f_*)=0.
\]
The proof first lifts a representative after a Wach cokernel
loss, then corrects its actual Herr defect in the kernel.
At \(\ell=2\) the depth satisfies \(s\ge2\).
\end{proposition}

\begin{proof}
\proofstep{Control the cokernel of the integral Wach map.}
Write \(A\) for the integral Wach coefficient ring
and \(N_i\) for the Wach lattices. On each coefficient
letter, let \(\mathfrak a\) annihilate the quotient of
\(N_2\) by the image of \(N_1\).
The Wach length comparison gives
\(\pi^c\in\mathfrak a\) for some \(c\).
The height inclusions
\(Q^h\varphi^*N_i\subset N_i\) and
\(Q^hN_i\subset\varphi^*N_i\) give
\[
 Q^H\mathfrak a\subset A\varphi(\mathfrak a),
\]
with the transported ideal on the Frobenius letter.

In fact multiplying by \(Q^h\) on either side of
the original inclusion through the image passes from
the Frobenius target to the original target and
then back to the Frobenius source. The Frobenius
coordinate basis over \(\varphi(A)\) is
\(1,\epsilon,\ldots,\epsilon^{\ell-1}\), with
\(\epsilon=1+\pi\), and \(\varphi(\epsilon)=\epsilon^\ell\).
Thus the Frobenius extension of the annihilator is
the ideal with precisely the original ideal as all
its coefficient components. These statements use
the same Wach map in the étale and height comparisons.

The coefficient extraction can be written on the original and transported letters as
\[
 X^H\sum_j\epsilon^j\varphi(b_j)\in\mathfrak a\quad\Longrightarrow\quad X^H b_j\in\mathfrak a_{\rm prev}\quad(\text{each }j),
\]
where \(\mathfrak a_{\rm prev}\) is the corresponding annihilator on the preceding letter. Indeed multiply by \(Q^H\) before extracting the Frobenius components. Repeatedly perform
this operation on \(X^H(\epsilon-1)^c\).
Write \(f(\epsilon)=(\epsilon-1)^c\), and at each
extraction write
\[
 f(\epsilon)=\sum_{j=0}^{\ell-1}\epsilon^j f_j(\epsilon^\ell).
\]
The leading term lies in \(j=c\bmod\ell\). After
taking that component its coefficient is still one and its
degree becomes \(\lfloor c/\ell\rfloor\).

Iterate on that
component; each extraction preserves the property that
\(X^H\) times the current polynomial belongs to the
annihilator on the corresponding letter. If \(d_j\) denotes the successive degrees, the monic extraction gives
\[
 d_0=c,\qquad d_{j+1}=\lfloor d_j/\ell\rfloor.
\]
Thus the degree becomes zero after finitely many steps and the terminal monic constant is one.
Thus \(X^H\) lies in the annihilator. This step is uniform
even if \(c\) depends on the lattice.

Put \(u=\chi(\gamma)\).
The full-level action on \(N_i/XN_i\) is trivial and \(\gamma(X)/X\equiv u\pmod X\), so
\[
 \gamma|_{X^jN_i/X^{j+1}N_i}=u^j\operatorname{id}.
\]
Apply the polynomials with roots
\(u^0,\ldots,u^{H-1}\) and \(u^H,\ldots,u^{2H-1}\)
to \(X^H y\) for \(y\in N_2\).
The first result lies in \(X^H\operatorname{im}N_1\)
by first lifting \(X^Hy\), then using the first
polynomial on that lift. The other polynomial gives a vector in the deeper filtration step, and the common annihilator gives
\[
 X^{2H}N_2=X^H(X^HN_2)\subseteq X^H\operatorname{im}N_1.
\]

For \(H>0\), interpolate an inverse of the second
polynomial at the roots of the first.
Every denominator divides a product of the
differences between one root \(u^i\) and all the
other roots in \(0\le j<2H\).
Its valuation is at most
\[
 (2H-1)s+v_\ell(i!)+v_\ell((2H-1-i)!)\le
                 E_W:=2Hs+v_\ell((2H)!).
\]
Here the valuation of each nonzero difference is
\(s+v_\ell(i-j)\), by the admitted cyclotomic depths.
Write \(L_0(Z)=\prod_{i=0}^{H-1}(Z-u^i)\) and
\(L_1(Z)=\prod_{i=H}^{2H-1}(Z-u^i)\).

At each root of \(L_0\), Lagrange interpolation for
\(\ell^{E_W}/L_1(u^i)\) gives a polynomial \(B\) with
integral coefficients by the displayed denominator bound.
Then \(\ell^{E_W}-BL_1\) vanishes at all the roots of
\(L_0\), and monic division gives an integral \(A\) with
\[
    A L_0+B L_1=\ell^{E_W}.
\]
Applying this relation at \(\gamma\) to \(X^H N_2\) puts its
image inside \(X^H\operatorname{im}N_1\). Cancelling \(X^H\)
gives
\[
 \ell^{E_W}N_2\subseteq\operatorname{im}N_1.
\]

For \(H=0\), \(X^H=1\) already gives surjectivity.
\label{cor:full-cyclotomic-lifting-2}

\proofstep{Correct a defect in the kernel through the cyclotomic graded pieces.}

For the Wach lattice \(N\) of the kernel let
\(N_Q=Q^{-h}N\) in the étale module.
The restricted Herr pairs lie in \(N_Q\oplus XN\),
with differential
\[
 d(y,z)=(\gamma-1)y-(\varphi-1)z
\]
taking values in \(XN_Q\). This uses
\(\gamma(Q)/Q\equiv1\bmod X\) and the height
inclusion. The exactness comparison identifies
the kernel of the map on these coordinates by
intersection with the original Wach lattices.
The following additional inputs of that comparison
are used in the correction.

\begin{itemize}
\item \(N\) is \(\ell\)-saturated in \(N_Q\), and
  \(Xy\in N\) for \(y\in N_Q\) implies \(y\in N\).
  These follow by division for the monic cyclotomic
  factors of \(Q\) in the free Wach lattices.
\item \(1-\varphi\) is onto on \(X^{h+1}N\).
  Indeed
\[
 \varphi(X^{h+1}N)\subseteq QX^{h+1}N,
\]

  and \(\varphi(Q)\) and all its iterates lie in
  the maximal ideal. The geometric series converges
  in the complete Wach lattice.
\item With \(U\) the union of source jumps, put
  \(J'=\bigl((U+\{1,\ldots,h\})\setminus\{0\}\bigr)
  \cup\{1,\ldots,h\}\) and
\[
 P(Z)=\prod_{j\in J'}(Z-u^{-j}).
\]

  The polynomial \((Z-1)P(Z)\) annihilates
  \(N_Q/N\) under \(\gamma\).
  The following local calculation explains both the distinct factors and the integral conclusion.
\end{itemize}
To see why repeated jumps do not require repeated polynomial factors, invert \(\ell\) and separate the simple roots of \(Q\). Each is fixed by \(\gamma\), since its cyclotomic coordinate has order dividing \(\ell^s\). In the completed component take a filtration-adapted crystalline basis. In the logarithmic parameter \(t\) the generators are \(t^{-b}v_b\) for the jumps \(b\), with \(\gamma(t)=ut\). Modulo the original lattice, \(Q^{-h}N\) contributes the individual monomials \(t^{-b-j}v_b\), \(1\le j\le h\). In particular
\[
         \gamma(t^{-b-j}v_b)=u^{-(b+j)}t^{-b-j}v_b .
\]
The eigenvectors for equal exponents have the same eigenvalue. The single factor for that eigenvalue kills their whole span. Monic division and freeness make \(Q^{-h}N/N\) torsion-free over the coefficient DVR, so the same annihilating polynomial works integrally. Strictness places the kernel jumps among the source jumps. \label{cor:full-cyclotomic-lifting-3}

Let \(c_G=\prod_{j=1}^h(u^j-1)\).
Successively correcting powers \(X^j\) for
\(1\le j\le h\) shows that for \(z\in XN\) one has
\[
 c_Gz=(\gamma-1)x+t,\qquad x\in XN,\quad t\in X^{h+1}N.
\]
In each degree the correction replaces multiplication
by \(u^j-1\) by \(\gamma-1\) modulo the next degree.
Suppose \(\omega\in XN_Q\) and
\(\ell^k\omega=d(y,z)\) for a restricted pair.
Then
\[
 c_G\ell^k\omega
   =(\gamma-1)\bigl(c_Gy-(\varphi-1)x\bigr)
                      -(\varphi-1)t.
\]
Here the first argument is in \(N_Q\) and the last
term is in \(N\). Applying \(P(\gamma)\) puts the
right side in \(N\) by the quotient comparison.
Saturation gives \(P(\gamma)\omega\in N\), and
since \(XN_Q\) is stable this lies in \(XN\).

Apply the graded correction again to this vector.
The remaining term in \(X^{h+1}N\) is a
\((1-\varphi)\)-image there.
Finally \(P(\gamma)-P(1)\) is divisible by
\(\gamma-1\) with integral polynomial quotient.
It follows that \(c_G P(1)\omega\) is the
differential of a restricted pair.
The valuation loss is
\[
 E_U = \sum_{j\in J'}(s+v_\ell(j))+hs+v_\ell(h!).
\]

\proofstep{Lift a finite extension, then kill the defect of its Herr representative.}
Lifting an integral vector in a Wach module alone need not give the desired finite cohomology class. We apply the preceding correction to an extension representing that class and then pass to its Herr representative.


A finite target cocycle defines a crystalline
extension of the trivial lattice by the target
lattice. A lift of \(\ell^{E_W}\) in the Wach module
of the extension exists by the same cokernel
bound, since the extension adds only the zero jump.
Its Frobenius and cyclotomic differences give the
restricted Herr pair representing the scaled
class by the connecting comparison.
A restricted pair lifted termwise from
the target has a defect in \(XN_Q\) in the kernel,
provided its target pair is a cocycle.
For this defect \(z\), the rational comparison supplies some \(k\ge0\) and a restricted kernel pair \((a,b)\) with
\[
 \ell^k z=d(a,b).
\]

Indeed, for rational crystalline representations \(V_1\to V_2\) in the sequence the crystalline and de Rham functors are exact, the latter strictly filtered. In the Bloch--Kato presentation
\[
 H^1_f(F_s,V_i)=\operatorname{coker}\!\left[
 D_{\rm cris}(V_i)\xrightarrow{(1-\varphi,\,\mathrm{can})}
 D_{\rm cris}(V_i)\oplus D_{\rm dR}(V_i)/\mathrm{Fil}^0\right]
\]
the terms to the right are therefore surjective from \(i=1\) to \(i=2\). Lift the target class rationally. By Lemma~\ref{lem:ramified-restricted-herr}, adjust for the difference of its representative and the given target pair by a coboundary. Equivalently the differential of the termwise lift has zero class in the rational kernel complex. Clearing the finitely many coefficient denominators yields the displayed equality for some \(k\). \cite[Section~3]{bloch-kato-finite}
\label{cor:full-cyclotomic-lifting-4}

The cokernel bound lifts the two coordinates
of the target representative after another
\(\ell^{E_W}\). Subtract the kernel correction
from the scaled lifted pair. This is now a
restricted cocycle in the source and represents
the desired finite lift by naturality.
The lift of the extension, the lift of its coordinates and the correction of the defect require the respective multipliers
\[
 \ell^{E_W}\ell^{E_W}\ell^{E_U}=\ell^{2E_W+E_U}.
\]

Indeed \(J'\subseteq[-H,H]\setminus\{0\}\), so
\[
 \sum_{j\in J'}v_\ell(j)
 \le 2\sum_{j=1}^H v_\ell(j)
 =2\sum_{a\ge1}\lfloor H/\ell^a\rfloor\le2H.
\]
Also \(|J'|\le 2H\), \(h\le H\) and
\(v_\ell((2H)!)\le2H\).
These bounds give \(2E_W+E_U\le14Hs\) for \(s\ge1\)
and in particular the stated allowance \(32(2hs+1)\).
All bounds are independent of the coefficient
degree and of the Wach cokernel's initial
\(\pi\)-power. The restriction--corestriction
argument returns the finite lift to the
original local field and original lattices.
\end{proof}


\subsubsection*{The finite condition after projecting to a character}
\label{subsec:ramified-character}

A finite class lifted through the crystalline lattices must still satisfy the required condition after projection to a character. The extension class below isolates the difference between an integral lift of a generator and a Galois-invariant one.

Let \(F_s=F_0(\mu_{\ell^s})\) be the full cyclotomic field of
Section~\ref{subsec:ramified-ht}, and let \(O\) be the coefficient DVR.
Choose a crystalline lattice extension and a lift \(v\) of its quotient generator:
\[
 0\longrightarrow T\longrightarrow\widetilde T\xrightarrow{p}O\longrightarrow0,\qquad p(v)=1.
\]
Put \(f(g)=gv-v\in T\).
Suppose that \(\psi:T\to E/O(-e)\) is an equivariant
map, where \(E\) is the fraction field of the
coefficient integers \(O\), and \(0\le e<t\le h\).
The kernel jumps lie in a common set
\(U\subseteq[-h,h]\) by strictness of the crystalline
filtration. We bound the loss to make
\(\psi f\) a finite class. Use \(M\) from
Section~\ref{subsec:ramified-ht} with the interval
\([-h,h]\).

\paragraph{The truncated integral periods.}
Write \(A_{\inf}\) for the integral period ring,
\[
 \xi=([\epsilon]-1)/\varphi^{-1}([\epsilon]-1)
\]
for the cyclotomic generator of the kernel of
Fontaine's map, and \(t_{\rm dR}\) for the
cyclotomic period.
Modulo \(\xi^{e+1}\), the cocycle
\(f\otimes1\) is a rational coboundary.
In fact strictness gives a lift of one in filtration
zero of the de Rham invariants of \(\widetilde T\).
Subtract it from \(v\) in the de Rham tensor.
The result lies in filtration zero of the tensor
with \(T[1/\ell]\), and has coboundary \(f\otimes1\).
The quotient \(A_{\inf}/\xi^{e+1}\) after inverting
\(\ell\) is the corresponding de Rham jet, so some
\(\ell\)-power gives an integral cobounding vector.

At level \(M\), the degree \(i\) piece of the jet
is the integral scalar ball in twist \(i\).
For Fontaine's map \(\vartheta\), the ratio of the two generators satisfies
\[
 \vartheta(t_{\rm dR}/\xi)=\zeta_\ell-1.
\]
The difference between the two
generators in degree \(i\) rescales the ball by an
element fixed by \(G_{F_M}\). Use the unprojected
frame \(D\) obtained at absolute depth \(s+n_0\) in
\eqref{eq:full-level-fontaine-map}. The map \(\vartheta\) in the jet
formula is the original depth-zero quotient. Indeed, writing
\(\mu=[\epsilon]-1\), one has
\[
  t_{\rm dR}/\xi
   =\bigl(\log(1+\mu)/\mu\bigr)\varphi^{-1}(\mu)
   \ \xrightarrow{\ \vartheta\ }\ 1\cdot(\zeta_\ell-1).
\]

Both the scalar in this formula and the scalar \(c\) in
\eqref{eq:full-level-fontaine-map} are fixed by \(G_{F_M}\).
The transformed lattice lies between the ball and
its \(\ell^{-1}\) multiple. Multiply a torsion cocycle
once by \(\ell\) to enter the ball. A rational
cobounding vector can also be brought into the ball
by a sufficiently large \(\ell\)-power.
The degree-one assertion of Lemma~\ref{lem:cyclotomic-invariant-loss}
applies even when the auxiliary extension has a zero jump. It bounds
the additional loss by
\[
   1+\max\left(1,\sum_{a\in U,\ a\ne i}(M+v_\ell(a-i))\right).
\]
This gives the summand of \(C(e+1)\) in degree \(i\)
in the original integral tensor lattice.

These losses add through the truncation.
For successive rational jets \(J'\twoheadrightarrow J\), the necessary lifting statement is
\[
 H^0(G_{F_M},J')\longrightarrow H^0(G_{F_M},J)\longrightarrow0 .
\]
An adapted filtration basis in the de Rham comparison reduces surjectivity to scalar filtered periods. In a scalar interval the
invariants come from its degree-zero piece, if
present, by Tate's vanishing in nonzero degrees.
This piece lifts by constants in the field \(F_M\).
After clearing denominators this is an integral
lift up to a power of \(\ell\).

For two successive quotients, kill the projected torsion class and subtract its integral coboundary. The remaining graded cocycle is torsion: some power of \(\ell\) gives a cobounding vector in the whole jet. Its
projection is invariant; lift that invariant
after clearing denominators and subtract the lift.
The remaining cobounding vector belongs to the
graded piece. If the projected class is killed by \(\ell^u\) and the recovered graded torsion by \(\ell^v\), the original class \(c\) satisfies
\[
 \ell^v(\ell^u c)=0.
\]
This is the additivity of the successive exponent bounds. Restriction to \(F_M\) and
corestriction back add
\[
 M-s=v_\ell([F_M:F_s])
\]
once. Thus \(\ell^{C(e+1)}(f\otimes1)\) is a
coboundary over \(A_{\inf}/\xi^{e+1}\).
\label{cor:cyclotomic-character-lifting-1}

\begin{lemma}[Integral dual exponentials of positive twists at two]
\label{lem:dyadic-integral-period}\label{hyp:positive-character-integral-period}
Let \(F_0/\mathbb Q_2\) be finite unramified, \(F_s=F_0(\mu_{2^s})\), \(s\ge2\), and \(e\ge1\). Identify the de Rham line for \(\mathbb Q_2(-e)\) with \(F_s\) in the Tate basis \(t_{\rm dR}^e\). The integral cohomology classes have dual exponentials satisfying
\[
  2^{1-e-v_2((e-1)!)}\mathcal O_{F_s}
           \subseteq\exp^*\!\bigl(H^1(F_s,\mathbb Z_2(-e))\bigr).
\]
\end{lemma}

\begin{proof}
\proofstep{Construct integral classes before taking their moments.}
Choose compatible primitive roots \(\zeta_r\) of order \(2^r\). On \(\mathcal O_{F_0}[[X]]\) write \(Y=1+X\) and let \(\varphi\) act by arithmetic Frobenius \(\sigma\) on coefficients and by \(Y\mapsto Y^2\). Its left inverse \(\psi\) extracts the even powers in the decomposition \(f=\varphi(f_0)+Y\varphi(f_1)\), so \(\psi(f)=f_0\).

At level \(r\ge2\), put \(u=1+2^r\). For a positive odd integer \(a\) and \(b\in\mathcal O_{F_0}\) set
\[
    g=b(Y^a-Y^{au}),\qquad y=\sum_{n\ge0}\varphi^n(g).
\]
Since \(Y^{2^n}-1\in(2,X)^{n+1}\), the series converges in \(\mathcal O_{F_0}[[X]]\). Both exponents in \(g\) are odd, hence \(\psi(g)=0\) and \(\psi(y)=y\). The integral Iwasawa comparison for the trivial lattice attaches to \(y\) a system of integral cohomology classes. Its moment of order \(-e\) at level \(r\) belongs to \(H^1(F_r,\mathbb Z_2(-e))\). At two the torsion-free cyclotomic tower used for this comparison starts at \(r=2\). \cite[Proposition~II.1.2 and Theorem~II.1.3]{cherbonnier-colmez-iwasawa}

\proofstep{Evaluate a finite-level moment.}
Let \(y^{\sigma^{-r}}\) denote the series obtained by acting on its coefficients. In the chosen Tate basis, the dual exponential of the moment is the coefficient of \(t^e\) in
\[
   2^{-r}y^{\sigma^{-r}}(\zeta_r\exp(t/2^r)-1),\qquad t=t_{\rm dR}.
\]
The dual exponential formula of Cherbonnier and Colmez gives this series at all sufficiently high levels. To descend one step from \(j\) to \(j-1\ge2\), trace the series. The two conjugates of \(\zeta_j\) have opposite signs. Their average on \(y\) retains its even part \(\varphi(\psi y)=\varphi(y)\). Hence
\[
 \operatorname{Tr}_{F_j/F_{j-1}}\!
     \left(2^{-j}y^{\sigma^{-j}}(\zeta_j\exp(t/2^j)-1)\right)
 =2^{-(j-1)}y^{\sigma^{-(j-1)}}(\zeta_{j-1}\exp(t/2^{j-1})-1).
\]
Corestriction of the moment commutes with the dual exponential and this unnormalized trace. The coefficient formula therefore holds for each \(r\ge2\). \cite[Theorem~IV.2.1 and its proof]{cherbonnier-colmez-iwasawa}

After division by the odd unit \(a^e\), the computed value is \(C_r S_{r,a}(b)\), where
\[
 C_r=\frac{1-u^e}{2^{r(e+1)}e!},\qquad
 S_{r,a}(b)=\sum_{n\ge0}2^{ne}\sigma^{n-r}(b)\zeta_r^{a2^n}.
\]
Indeed the powers \(a2^n\) and \(au2^n\) give the same root of unity at level \(r\), while the coefficient of \(t^e\) in their exponentials introduces their \(e\)-th powers. The series for this coefficient converges since \(e>0\). For \(r\ge2\), the lifting-the-exponent formula yields
\[
    v_2(C_r)=r+v_2(e)-r(e+1)-v_2(e!)
            =-re-v_2((e-1)!).
\]

\proofstep{Recover each integral direction in the cyclotomic field.}
Replace \(a\) by \(a+2^{r-1}\). In the difference of the two normalized moments all the terms with \(n\ge1\) cancel. Since \(b\) is arbitrary,
\[
 C_r\bigl(S_{r,a}(b)-S_{r,a+2^{r-1}}(b)\bigr)
                   =2 C_r\sigma^{-r}(b)\zeta_r^a .
\]
Thus the image at level \(r\) contains \(2^{1-re-v_2((e-1)!)}\mathcal O_{F_0}\zeta_r^a\) for each odd \(a\).

The ground-field direction is also needed. At \(r=2\), add the normalized moments for \(a=1,3\). Their leading terms cancel, the roots in the next terms equal \(-1\), and all later roots equal one. Consequently
\[
 C_2\bigl(S_{2,1}(b)+S_{2,3}(b)\bigr)
 =2^{e+1}C_2\left(-\sigma^{-1}(b)+
             \sum_{n\ge2}2^{(n-1)e}\sigma^{n-2}(b)\right).
\]
The endomorphism in parentheses is congruent to \(-\sigma^{-1}\) modulo \(2\), hence is bijective on the finite free \(\mathbb Z_2\)-module \(\mathcal O_{F_0}\). The image contains \(2^{1-e-v_2((e-1)!)}\mathcal O_{F_0}\).

Restriction from \(F_r\) to \(F_s\) for \(2\le r\le s\) preserves integrality and includes the dual exponential value into \(F_s\). Each power of \(\zeta_s\) is \(\pm1\) or an odd power of a primitive \(\zeta_r\) with \(2\le r\le s\). These powers span \(\mathcal O_{F_0}[\zeta_s]=\mathcal O_{F_s}\) over \(\mathcal O_{F_0}\). The exponent \(1-re-v_2((e-1)!)\) is at most \(1-e-v_2((e-1)!)\), proving the assertion.
\end{proof}

\begin{proposition}[Returning scalar cocycles to the finite condition]\label{prop:cyclotomic-scalar-finite-condition}
Take the cyclotomic module of Proposition~\ref{prop:ramified-integral-wach-comparison} for the ramified lattice. Let \(f(g)=gv-v\) represent the crystalline lattice extension over \(F_s\),
and let \(\psi:T\to(E/O)(-e)\) be equivariant, with \(0\le e<t\le h\).
With the integer loss \(C(k)\) of \eqref{eq:local-jump-loss}, put
\[
 A_e=(e+2)s+v_\ell(e)+\lfloor e/(\ell-1)\rfloor+1,
 \qquad B_e=C(e+1)+s+v_\ell(e+1)+A_e,
\]
where the \(v_\ell(e)\) term in \(A_e\) is omitted when \(e=0\). Then the pushed-out class satisfies
\[
 \ell^{B_e}[\psi\circ f]\in H^1_f(F_s,(E/O)(-e)).
\]
Here the finite subgroup is the image of the rational finite condition.
For a class defined over the original local field at level \(N\),
corestriction adds at most \(1+\lfloor\log_\ell N\rfloor\) to the exponent.
\end{proposition}

\begin{proof}
\proofstep{Use the Bockstein allowance to lift the divisible cocycle to rational coefficients.}
Lift \(\psi\) on an integral
\(\mathbb Z_\ell\)-basis of \(T\) to the rational
target. The lift need only be equivariant modulo the
target lattice. Expand in an integral
\(\mathbb Z_\ell\)-basis of this target lattice.
Each coordinate of \(\psi f\) modulo \(\mathbb Z_\ell\)
is now a cocycle in
\(\mathbb Q_\ell/\mathbb Z_\ell(-e)\).

By local Tate duality,
\(H^2(F_s,\mathbb Z_\ell(-e))\) is dual to the
invariants in \(\mathbb Q_\ell/\mathbb Z_\ell(e+1)\).
The cyclotomic generator has \(v_\ell(\chi(\gamma)^{e+1}-1)=s+v_\ell(e+1)\). Duality therefore gives
\[
 \ell^{s+v_\ell(e+1)}H^2(F_s,\mathbb Z_\ell(-e))=0.
\]
The integral Bockstein sequence
then gives a rational cocycle representing the scaled
divisible cocycle in each coordinate. Adjust by a
rational coboundary lifting the discrepancy between
the divisible cocycles. The representatives then
agree pointwise modulo
\(\mathbb Z_\ell\).
The equivariance errors of the lifted coefficient
map are integral, so applying it to the jet
coboundary preserves the congruence modulo integral
periods. After scaling further by \(\ell^{C(e+1)}\),
each rational cocycle \(z\) so obtained satisfies
\[
 z(g)\equiv \chi(g)^{-e}g b-b
       \pmod {A_{\inf}+t_{\rm dR}^{e+1}B_{\rm dR}^+}
 \qquad (b\in B_{\rm dR}^+).
\]
The congruence means that the error is in the
image of \(A_{\inf}\) in the jet.

Use the dual exponential in the Tate basis \(t_{\rm dR}^e\). Its value \(\alpha\in F_s\) is
specified by the remaining term
\(\log\chi(g)\alpha t_{\rm dR}^e\) after a rational
twisted coboundary in that jet. To see this, filter the rational jet
\((B_{\rm dR}^+/t_{\rm dR}^{e+1})(-e)\) by powers of \(t_{\rm dR}\).
Its degree-\(i\) quotient is \(\mathbb C_\ell(i-e)\) for \(0\le i\le e\).
For \(i<e\), Tate's theorem gives
\(H^1(F_s,\mathbb C_\ell(i-e))=0\), so the cocycle in that quotient is
a coboundary. In degree \(e\) the twist is trivial, and Tate's calculation gives
\[
 H^1(F_s,\mathbb C_\ell)=F_s\,[\log\chi].
\]

When the dual exponential is zero, the Bloch--Kato
comparison gives a cobounding vector over
\(B_{\mathrm{cris}}\otimes_{\mathbb Q_\ell}K_\ell(-e)\),
where \(K_\ell\) is the coefficient field. This identifies
the corrected class with the crystalline finite condition.
\cite{tate-pdivisible,bloch-kato}\label{cor:cyclotomic-character-lifting-2}

\proofstep{Estimate the dual exponential in the truncated period coordinates.}
The Bockstein argument leaves a rational lift. Control of its projection to the finite subgroup now comes from the integral periods entering the dual exponential.

The displayed integral congruence bounds \(\alpha\).
At degree \(i<e\) use the integral multiplier
\[
 \ell^{1+s+v_\ell(i-e)}
\]
on the modified cobounding vector, then subtract an integral lift of its leading coefficient.
The degree \(i\) scalar belongs to the ball
modulo a twisted invariant, after this multiplication.
This is the rank-one ball argument at depth \(s\)
in twist \(i-e\), using the valuation of
\(\chi(\gamma)^{i-e}-1\).
In the \(t_{\rm dR}^i\) coordinates the integral
leading coefficient is multiplied by
\((\zeta_\ell-1)^i\), from \(t_{\rm dR}/\xi\).
Thus it can be lifted by Fontaine's surjection
and the power \(\xi^i\).
After eliminating all lower degrees the logarithm coefficient is
\[
 \ell^{e(s+1)+v_\ell(e!)}(\zeta_\ell-1)^e\alpha.
\]
Its two factors record the scalar elimination costs and the generator conversion, respectively.

If \(\alpha_1\log\chi(g)+g b_1-b_1\) is integral
and \(\alpha_1\in F_s\), apply the integral
coboundary estimate on the cyclotomic kernel
after multiplying by \(\ell\).
The remaining cobounding vector is fixed by that
kernel. Take the integral normalized trace at the
generator. It kills the difference of this
vector and its translate and gives
\(\ell\alpha_1\log\chi(\gamma)\) integral.
Since \(v_\ell(\log\chi(\gamma))=s\), one obtains
\[
 v_\ell(\alpha)\ge
 -e(s+1)-v_\ell(e!)-1-s-\frac{e}{\ell-1}.
\]

At weight zero the relevant integral image can be computed by local class field theory. Write \(E_0=(\ell-1)\ell^{s-1}\) for the ramification index of \(F_s/F_0\). The different is generated by \(\Phi'_{\ell^s}(\zeta_{\ell^s})\); its absolute valuation, with \(v_\ell(\ell)=1\), is \(s-1/(\ell-1)\). For \(v_\ell(z)\ge1/E_0\) the logarithmic series gives
\[
 v_\ell(\log(1+z))\ge \min_{n\ge1}\{n/E_0-v_\ell(n)\}
                  =1+1/(\ell-1)-s .
\]
For fixed \(v_\ell(n)=j\) the minimum occurs at \(n=\ell^j\); comparing consecutive values gives the equality. For the inverse different \(\mathfrak D^{-1}=\mathfrak D_{F_s/\mathbb Q_\ell}^{-1}\) this gives
\[
          \ell^{-1}\mathcal O_{F_s}\log(1+\mathfrak m_{F_s})
                              \subseteq\mathfrak D^{-1}.
\]

For \(a\in\ell^{-1}\mathcal O_{F_s}\) define on principal units
\[
       h_a(u)=\operatorname{Tr}_{F_s/\mathbb Q_\ell}(a\log u)\in\mathbb Z_\ell .
\]
Extend it by zero on prime-to-\(\ell\) roots of unity and a uniformizer. Local reciprocity gives an integral Galois character. The duality of the Bloch--Kato exponential with Kummer theory identifies its dual exponential with \(a\), up to the common reciprocity sign: pairing on principal units on both sides is the displayed trace. This proves the required inclusion for \(e=0\). \cite[Section~3]{bloch-kato-finite}

For odd \(p=\ell\) and \(e>0\), put \(m=e+1\) and \(L=F_s\). For \(m\ge2\) the rational Bloch--Kato exponential is an isomorphism. In its standard Tate basis put
\[
 \mathcal H_m=\exp_{L,m}^{-1}\!\left(\operatorname{im}\big[H^1(L,\mathbb Z_p(m))
                           \longrightarrow H^1(L,\mathbb Q_p(m))\big]\right)\subset L .
\] Let \(\sigma\) be Frobenius on \(F_0\), and let \(f\in\mathcal O_{F_0}[[X]]\) range over the power series satisfying \(\sum_{\zeta^p=1}f(\zeta(1+X)-1)=0\). Denote their module by \(R\).

The integral Coleman sequence and its explicit evaluation at level \(s\) give
\[
          p^{s+v_p(e)}\mathcal H_m\ \subseteq\ e!\,\Xi_{m,s}(R).
\]
Indeed the quotient of the full inverse-limit classes by the classes from principal units, after passing to the coinvariants for \(F_s\), is \(H^0(F_s,\mathbb Q_p/\mathbb Z_p(e))\). A generator with cyclotomic value \(1+p^s\) acts on that quotient through the relation
\[
       ((1+p^s)^e-1)x=0,\qquad v_p((1+p^s)^e-1)=s+v_p(e).
\] For the principal-unit classes the evaluation diagram sends each integral class to \((m-1)!\,\Xi_{m,s}(f)\) in exponential coordinates. Here \(L=F_s\) is already the entire cyclotomic level, so the trace from that level to \(L\) in the diagram is the identity. \cite[Sections~2.3.1--2.3.2, (4)--(5)]{benois-nguyen-tamagawa}

The explicit evaluation gives more than its determinant. In the formula for \(\Xi_{m,s}\), division by \(p^{m-s}\) makes the terms of index \(0\le k<s-1\) divisible by \(p^m\): their coefficients are \(p^{m(s-1-k)}\). The term for \(k=s-1\) and the constant correction give
\[
 p^{s-m}\Xi_{m,s}(f)\equiv
 f^{\sigma^{-1}}(\zeta_p-1)-(1-p^m\sigma^{-1})^{-1}f^{\sigma^{-1}}(0)
                      \pmod{p^m\mathcal O_L}.
\]
The inverse in the second term is the convergent geometric series and is the identity modulo \(p^m\). Evaluation of an integral series at \(\zeta_p-1\) agrees with its constant coefficient modulo \(\zeta_p-1\). The calculation also covers \(s=1\), when the earlier sum is empty. \cite[Section~2.2.1, definition of \(\Xi_{m,n}\)]{benois-nguyen-tamagawa}

In particular every vector in \(\mathcal H_m\) has valuation at least
\[
 A=v_p(e!)+m-2s-v_p(e)+1/(p-1).
\]
Under local Tate duality the image of \(H^1(L,\mathbb Z_p(1-m))\) in the dual de Rham line is the trace dual of \(\mathcal H_m\). Indeed the integral cup product on the torsion-free quotients is perfect, and the Bloch--Kato exponential and dual exponential are adjoint for this pairing. In the dual Tate bases their de Rham pairing is \((a,b)\mapsto\operatorname{Tr}_{L/\mathbb Q_p}(ab)\). \cite[Section~3]{bloch-kato-finite}

Put \(\beta=s-1-e-v_p((e-1)!)\). The valuation threshold for the inverse different is \(-s+1/(p-1)\); with \(m=e+1\) one has
\[
        \beta+A=-s+1/(p-1).
\]
Thus \(p^\beta\mathcal O_L\mathcal H_m\subseteq\mathfrak D^{-1}\). Every vector of \(p^\beta\mathcal O_L\) has integral trace pairing with \(\mathcal H_m\), proving the required inclusion at every odd prime.

For \(\ell=2\) and \(e>0\), the exponent in Lemma~\ref{lem:dyadic-integral-period} is independent of the cyclotomic level. Since \(s\ge2\), the required ball satisfies
\[
  2^{s-1-e-v_2((e-1)!)}\mathcal O_{F_s}
     \subseteq 2^{1-e-v_2((e-1)!)}\mathcal O_{F_s}
     \subseteq\exp^* H^1(F_s,\mathbb Z_2(-e)).
\]
\label{cor:cyclotomic-character-lifting-3}

Set
\[
 A_e=(e+2)s+v_\ell(e)+\lfloor e/(\ell-1)\rfloor+1,
\]
omitting the \(v_\ell(e)\) term at \(e=0\).
The bound for \(\alpha\) gives
\(\ell^{A_e}\alpha\in\ell^\beta\mathcal O_{F_s}\).
For positive \(e\) use
\(v_\ell(e!)=v_\ell(e)+v_\ell((e-1)!)\) and
\(e/(\ell-1)<\lfloor e/(\ell-1)\rfloor+1\).
For \(e=0\) it follows directly.
Subtract an integral cocycle with this exponential
from \(\ell^{A_e}z\). The difference has zero
exponential and represents the same divisible
class. Reassemble the scalar coordinates in the
coefficient lattice. This puts the character image
in the divisible finite condition after loss
\(C(e+1)+s+v_\ell(e+1)+A_e\).
Corestriction adds at most the original local degree
loss, giving the character allowance used above.
\end{proof}



\subsubsection*{Why the retained localization has an irreducible block}
\label{subsec:retained-galois}
\hypertarget{source-retained-galois}{}

The local estimates on the character summands leave a complementary Galois block. Its residual irreducibility is needed for the integral matrix algebra used in the congruence argument. The two tame type conditions test that irreducibility on the retained support.

\begin{paramdefinition}[Good theta point and the two tame type tests]\label{hyp:retained-type-tests}
For a candidate type at a theta point of Section~\ref{subsec:selected-cut}, let \(g\in\{g_0,g_1\}\) be a type prime and let \(r\) be
the rank of the complementary summand. At each place \(v\mid g\), require that the degree-\(r\) polynomial \(P_v\) have simple nonzero roots. Call the tuple good at this prime when it satisfies all three tests.
\begin{itemize}
\item For each possible constituent dimension, one place excludes it:
\[
 \forall a\in\mathbb Z\cap[1,r/2]\quad\exists v\mid g
 \quad\nexists Q\in\mathbb F_g[X]:\quad Q\mid P_v,\quad \deg Q=a .
\]
\item Away from the bad prime set for \((r,g)\), projection to the
prime-to-\(\ell\) parts is injective on the roots of each \(P_v\).
At every \(\ell\ne g\), those projected roots avoid the projected
labels of the character block and their inverses.
\item At one place there is an irreducible factor of prime degree
\(p\) with \(r/2<p\le r-3\).
\end{itemize}
Retain a pair of types when it is good at both primes. The bad prime set is
\[
 \{g\}\cup\{\ell\text{ prime}:\ g^d<\ell^{4v_\ell(g^d-1)}
                    \text{ for some }1\le d\le r\}.
\]
The projector in Definitions~\ref{def:localized-type-support} and~\ref{def:good-projector-support} retains exactly the jointly good labels within its theta support. Proposition~\ref{prop:two-type-primes} supplies type primes with disjoint exceptional sets. For a root of unity \(z\) of order
\(\ell^a b\), with \(\ell\nmid b\), its prime-to-\(\ell\) part
is \(z^e\), where \(e\equiv0\pmod {\ell^a}\) and
\(e\equiv1\pmod b\).  \label{cor:retained-galois-1}
\end{paramdefinition}

\begin{proposition}[The irreducible rank-complement]\label{prop:retained-rank-complement}
At a good theta point satisfying Definition~\ref{hyp:retained-type-tests},
the complement of the prescribed string is cuspidal of rank \(r\).
Its attached integral representation has absolutely irreducible
reduction at the coefficient prime \(\ell\). Choose a marked type
prime \(g\) with \(\ell\notin B_r(g)\), with \(B_r(g)\) as in Definition~\ref{def:type-exceptional-primes}. Write \(\bar W\) for the reduction over an algebraic closure of its residue field. For every open normal subgroup \(H\) with solvable quotient containing the inertia groups at the selected places over \(g\),
\[
 0\ne U\subseteq\bar W,\quad HU=U\quad\Longrightarrow\quad U=\bar W
\]
for every subspace \(U\).
\end{proposition}

\begin{proof}
\proofstep{Apply the common divisor-degree test to the possible isobaric partitions.}
First the parameter left after removal of the specified character
block is cuspidal of rank \(r\). Indeed, if its parts have ranks
\(a_j\) and multiplicities \(d_j\), the degree calculation is
\[
 r=\deg\!\left(\prod_j P_{v,j}^{d_j}\right)=\sum_j d_j\deg P_{v,j}\le\sum_j d_j a_j=r.
\]
Thus \(\deg P_{v,j}=a_j\) for every part at every place.
The labels have no repetitions, so every \(d_j=1\). If there were at least two parts, choose a smallest one of rank \(a\). Its factor would satisfy
\[
 1\le a\le r/2,\qquad P_{v,j}\mid P_v,\qquad\deg P_{v,j}=a\quad(\text{every }v).
\]
This contradicts the
first test.

\proofstep{At a good type prime, use the distinct prime-to-residue-characteristic inertia labels.}
Choose \(g\) for which \(\ell\) is outside the bad set. The attached
integral representation for this regular algebraic conjugate
self-dual cusp has the following inertia polynomials. In the
depth-zero dictionary a root \(z\) gives the corresponding tame
root of unity on a fixed generator \(T_v\). Transport to
\(\overline{\mathbb Q}_\ell\) followed by reduction kills its
\(\ell\)-part and is injective on the roots of order prime to
\(\ell\). For the prime-to-\(\ell\) order \(b\), these roots remain distinct because
\[
 z^b=1\quad\Longrightarrow\quad b z^{b-1}\ne0
 \quad\text{in residue characteristic }\ell.
\] The unramified
normalizing twist has value one on inertia. Thus the eigenvalues of
\(\bar\rho(T_v)\) are the images of the projected roots just
described. The depth-zero parameter has trivial wild inertia.
For a wild element the polynomial \((X-1)^r\) makes its residual image unipotent, hence of \(\ell\)-power order. Its order is also a power of \(g\), so
\[
 g\ne\ell\quad\Longrightarrow\quad\operatorname{ord}(\bar\rho(\text{wild element}))=1.
\]

Arithmetic Frobenius conjugates \(T_v\) to \(T_v^g\) modulo wild
inertia, since every marked type place splits completely and has
residue field \(\mathbb F_g\). Consequently Frobenius permutes the
residual eigenlines with the same orbit lengths as the roots of
\(P_v\) under \(z\mapsto z^g\). Projection commutes with this
power map and the roots remain distinct. On lines, taking the
inverse permutation gives the same orbit lengths.

\proofstep{Use the two permutation tests and then restrict to the auxiliary subgroup.}
The tame generator has separated the residual inertia lines. Thus an invariant subspace would select a set of lines stable under the permutation induced by Frobenius.

An invariant subspace over \(\overline{\mathbb F}_\ell\)
is a sum of the eigenspaces for \(T_v\), by the polynomial
projectors for its distinct eigenvalues. It is Frobenius stable.
Thus for each selected \(v\) its dimension and codimension lie in
the set
\[
     \left\{\sum_{C\in S}|C|:
          S\text{ a subset of the Frobenius cycles at }v\right\}.
\]
This set consists of the degrees of divisors of \(P_v\)
over \(\mathbb F_g\). The first test excludes a common degree
between \(1\) and \(r/2\), so no proper nonzero subspace
is possible. Scalar normalization preserves invariant subspaces.
At each split unramified Frobenius outside the finite level, denote the full polynomial by \(P_{\rm full}\) and the monic character factors by \(P_1,\dots,P_t\). The comparison with the cusp reads
\[
 P_{\rm full}/(P_1\cdots P_t)=P_{\rm cusp}.
\]
Chebotarev on a common finite quotient and Brauer--Nesbitt give
the same semisimplified reduction. Its irreducibility implies
irreducibility of the actual reduction.

Let \(H\) now be the open normal subgroup containing the selected
inertia groups, with solvable quotient. At the prime-degree
test place a minimal nonzero \(H\)-invariant space is a span of
coordinate lines of \(T_v\). Its translates are also minimal
and distinct ones have zero intersection. Full irreducibility makes these translates a partition of the eigenlines into \(h\) blocks of size \(b\), hence
\[
 \#\{\text{eigenlines}\}=hb=r.
\]
A power of the Frobenius permutation is a single
\(p\)-cycle: the other cycles have lengths smaller than \(p\).
For \(h,b>1\), the group preserving the blocks has order
\((b!)^h h!\), prime to \(p\) since
\[
 h,b\le r/2<p.
\]

This excludes those values of \(h,b\).

If \(b=1\), the action on the lines is transitive and factors
through the solvable quotient. It is primitive by the same
order argument applied to any invariant nontrivial partition.
A minimal nontrivial normal subgroup is elementary abelian
and transitive, and faithfulness makes that abelian action
regular. Hence the action is affine on \(\mathbb F_q^u\)
for a prime \(q\) with \(r=q^u\). Our \(p\)-cycle has a
nonempty affine fixed space, so
\[
 r-p=q^d.
\]
In the test
range \(3\le r-p<p\), giving \(d\ge1\). But then
\[
               p=q^d(q^{u-d}-1)
\]
would force \(q^d=p\), a contradiction. The remaining
partition has \(h=1\), so the original minimal subspace
is the whole space.

This proves the restriction assertion for \(H\). \label{cor:retained-galois-2}
\end{proof}

\subsubsection*{Constructing the two integral projections}

With residual irreducibility established, the next step is integral. The matrix algebra must contain projections that distinguish the complementary block from a character line after specialization.

\begin{proposition}[Integral projections at a retained localization]\label{prop:retained-integral-projections}
At a retained theta localization with absolutely irreducible
residual complement of rank \(r>1\), the local matrix algebra contains orthogonal idempotents \(e_B,e_1\). For the complement \(W\) and every character line \(L\) at a theta point, their ranks are
\[
 (\operatorname{rk}e_B|_W,\operatorname{rk}e_1|_W)=(0,1),\qquad (\operatorname{rk}e_B|_L,\operatorname{rk}e_1|_L)=(1,0).
\]
The marked corner operator may be chosen with simple spectrum at
the retained fibres, keeping its prescribed value at the selected
point and perturbing by any given power of the coefficient uniformizer
in the ideal vanishing there.
\end{proposition}

\begin{proof}
Let \(O\) be the complete coefficient DVR and \(F\) its fraction
field. At a localization containing a good point write \(X\)
for the finite set of systems and \(\mathcal T\subset O^X\)
for the local Hecke algebra. At every \(x\) there is a
continuous semisimple integral model \(R_x\), with characteristic
coefficients in \(\mathcal T\) as functions of \(x\).
Here an integral model means a stable lattice in the semisimple
rational representation. A compact image stabilizes such a lattice
by taking the span of the orbit of an initial lattice. The
coefficient splitting in the selected norm data realizes the
constituents over \(F\). At theta points use the rational
decomposition into the rank-\(r\) complement and the \(t\)
prescribed lines.

Let \(k\) be the residue field and \(\mathfrak m\) the unique maximal
ideal of \(\mathcal T\). Each residual evaluation
\[
 \overline{\operatorname{ev}}_x:\mathcal T\longrightarrow k
\]
is onto, since it contains the reduction of the constant scalars. Its kernel is therefore
\(\mathfrak m\), and the induced isomorphism \(\mathcal T/\mathfrak m\to k\)
is the same for all \(x\) because it fixes the residue scalars. Thus,
for every \(g\) and any two systems \(x,x'\), reduction of the
characteristic coefficients gives
\[
 \det(1-T\overline R_x(g))=\det(1-T\overline R_{x'}(g))
       \quad\text{in }k[T],
 \qquad
 \overline R_x^{\,\mathrm{ss}}\simeq \overline R_{x'}^{\,\mathrm{ss}}
\]
by Brauer--Nesbitt. This includes the systems outside theta.
In particular their distinct simple modules consist of the
absolutely irreducible rank-\(r\) module above and some
one-dimensional characters. Repeated residual characters are
allowed. Here \(r>1\) by the type tests.

\proofstep{The residual algebra and its idempotents.}
The integral residual action is obtained from the simultaneous matrix order
\[
 A=\sum_g\mathcal T\,(R_x(g))_x.
\]

It is a finite torsion-free \(O\)-algebra inside a product of
matrix rings. In its residue action on the distinct simple
modules the image is the product of their full matrix rings.
Indeed each action is onto by absolute irreducibility and the
density theorem, and inequivalent simples have different
maximal two-sided annihilators, which are pairwise comaximal.
This also follows by applying the semisimple algebra theorem
to the joint action. The simple actions for \(A\) here are the
ones for the residual Galois matrices, because the extra Hecke
scalars have common reduction. In the product of the rank-\(r\) matrix factor with the character factors, the desired residual projections are
\[
 (0,1,\ldots,1),\qquad (E_{11},0,\ldots,0).
\]


Let \(k=O/\varpi O\) and let \(J\) be the Jacobson radical
of \(A/\varpi A\). Its semisimple quotient surjects onto
the product acting on the indicated simples. Choose the
character projection and the matrix unit there and give
both zero on any unused simple factors. For an approximate idempotent \(a\), correct its defect by
\[
 d=a^2-a,\qquad b=3a^2-2a^3=a-(2a-1)d.
\]
The polynomial identities
\[
       b-a=-(2a-1)d,\qquad
       b^2-b=d^2(4d-3)
\]
show that the error is squared at each iteration.
This first lifts across the nilpotent ideal \(J\) and then
converges in the \(\varpi\)-adically complete ring \(A\).
After lifting the character projection \(e_B\), lift the
second idempotent \(e_1\) inside the corner with identity
\(1-e_B\). This preserves orthogonality.

For a stable integral fibre, an idempotent \(e\) gives
\(R_x=eR_x\oplus(1-e)R_x\) on its lattice. Reduction
therefore preserves its rank, and applying \(e\) on an
invariant composition series adds the ranks on the
residual factors. There are \(t\) one-dimensional factors \(L_i\)
counted with multiplicity, and one rank-\(r\) factor \(W\).
The chosen residual actions give
\[
 \begin{aligned}
 \operatorname{rk}(e_B R_x)&=\sum_{i=1}^{t}\dim e_B L_i
                                      +\dim e_B W=t,\\
 \operatorname{rk}(e_1 R_x)&=\sum_{i=1}^{t}\dim e_1 L_i
                                      +\dim e_1 W=1 .
 \end{aligned}
\]
This calculation depends only on the residual composition
factors and hence holds for every admitted fibre. At the marked theta point the \(\mathcal T\)-linear
matrix expressions restrict individually to each of the
rational constituents. Applying the same rank check on
each given stable lattice gives zero and rank one on the
rank-\(r\) summand, respectively, and identity and zero
on every scalar line. These are the characteristic-zero
projections asserted above.

\proofstep{An inertia element with the prescribed depth.}
At the split small primes the local host lies inside the compositum of the unramified field and the level-\(s\) cyclotomic field. Hence one can choose
\[
 g_0\in I_\ell,\qquad \epsilon_{\rm cyc}(g_0)=1+\ell^s,
\]
where \(I_\ell\) denotes host inertia and \(\epsilon_{\rm cyc}\) the cyclotomic character. When a type prime \(g_k=\ell\), start with
\(\sigma\) in inertia above \(\ell\) whose cyclotomic value
is \(1+\ell^{s_0}\), with \(s_0\ge1\) the prescribed local
depth.

In the closed procyclic subgroup generated by \(\sigma\), define
\[
 g_0=\operatorname{pr}_{\ell}(\sigma),
\]
where the projection is to the \(\ell\)-primary component in each finite quotient. The original
inertia group is closed, hence contains \(g_0\). At every
place above \(\ell\) whose inertia contains \(g_0\), its
image in tame inertia is trivial since that quotient has
prime-to-\(\ell\) order in every finite quotient. The
cyclotomic image of the closed cyclic group lies in the
pro-\(\ell\) group \(1+\ell^{s_0}\mathbb Z_\ell\).
Projection therefore preserves its value, in particular the
required exact depth
\[
 v_\ell(\epsilon(g_0)-1)=s_0.
\]

In particular the marked \(t\) eigenvalues are distinct,
since consecutive character values differ by the fixed
cyclotomic ratio of infinite order.

\proofstep{Preserving the value at the marked point.}
Lifting an idempotent separates the blocks in the integral order. The determinant calculation also requires a corner element with a prescribed specialization at the marked point.

We construct \(X_0\) in the corner of \(e_B\), specializing
to \(e_BR_\pi(g_0)e_B\) at \(\pi\) and with simple spectrum
on the corner at all \(x\). Over an algebraic closure of
\(F\) the evaluation algebra at each \(x\) is the product
of full matrix algebras on distinct irreducible constituents.
Their multiplicities are one by the weight test of
Section~\ref{subsec:selected-cut}. Its \(e_B\)-corner can
therefore realize a diagonal matrix with \(t\) distinct
entries. Over \(F\), separation by \(\mathcal T\) gives all
individual coordinate functions by interpolation. Hence the
rational span of
\[
  J_\pi e_B A e_B,\qquad
  J_\pi=\{f\in\mathcal T:f(\pi)=0\},
\]
gives all the corner directions away from \(\pi\).
Write \(V_x\) for the rational fibre at \(x\).

On the perturbation lattice consider the polynomial
\[
    F(s)=\prod_{x\ne\pi}\operatorname{disc}
       \det\bigl(T-(e_BR(g_0)e_B+s)|_{e_B V_x}\bigr).
\]
Each factor is nonzero as a polynomial, because the directions
at that point include matrices with distinct eigenvalues.
Thus the product is nonzero, including the empty product. This may be checked after
extending the fraction field by the preceding matrix
description. For any prescribed \(L\ge0\), the lattice scaled
by \(\varpi^L\) is Zariski dense in its rational span. Indeed a
polynomial vanishing on every integral linear combination of
a spanning list is zero there by repeated univariate
interpolation over the infinite ring \(O\). Writing \(D_x\) for the discriminant of the perturbed corner at \(x\), density supplies an integral perturbation \(s\) in the chosen scaled lattice with
\[
 s(\pi)=0,\qquad \prod_{x\ne\pi}D_x(s)\ne0,
\]
with the empty product interpreted as one.
This gives all three marked matrices and the perturbation in
the integral congruence module. The caller takes \(L=0\).

\proofstep{The finite choice of coefficient field.}
To arrange roots in the actual coefficients, first take the
system at a preliminary coefficient level which splits the
constituents and admits the theta decompositions. At each live
localization make the choices just proved and record the roots
of its corner polynomials at all points and all theta marks,
using the tied embedding in \(\mathbb C\). This is a finite set
for that localization. Write \(S_i\) for the finite root set at a live localization. At a coefficient place \(w\) take
\[
 S_w=\bigcup_{j=0}^{m}\ \bigcup_{i\mid(j,w)}S_i,
\]
with the inner union over the live localizations there. The construction of the coefficient rings
allows any prescribed finite set at each such place, along
with the elements required to identify the systems and
their integral Hecke actions at all larger levels.

At the resulting level every live localization lies over
a preliminary one at the same coefficient place. The Hecke
character and theta factorization descend along this
comparison. If \(k\subset k'\) are the two residue fields, cancellation of the common character factors at split Frobenius elements identifies the semisimplified rank-\(r\) complements:
\[
 (\bar W_{\rm old}^{\rm ss}\otimes_k k')^{\rm ss}\simeq \bar W_{\rm new}^{\rm ss}.
\]
Thus a live point upstairs lies over a live one downstairs. The marked roots occur among
the prescribed elements at its place. The integral Hecke
span extends by the identified coefficient map, so the
old matrices and their polynomial factorizations give the
marks there. Restricting to a cut restricts these functions
and matrices, keeps every good point and can only decrease
the finite set of root conditions.  It applies to the coefficients allowed
in the original choice of a norm. No estimate on the degree
of this splitting field enters the normalized lengths.
\label{cor:retained-galois-3}
\end{proof}

\subsubsection*{Inertia tests in the localized Hecke algebra}

The projections separate the two blocks. The congruence presentation also needs inertia relations inside the localized Hecke algebra, where different character blocks can have identical restrictions to inertia.

Use the same localization \(i\), coefficient DVR, character
\(\chi_1\), and the sets \(S,Q\) of its norm choice. All their
places are away from \(\ell\). Write \(\gamma_v\) for the
chosen lift of a topological tame generator. Unramifiedness of the cyclotomic character gives
\[
 \chi_c|_{I_v}=(\chi_1\epsilon^{1-c})|_{I_v}=\chi_1|_{I_v}.
\]

There are two different applications of the type dictionary.

At a type place \(v\in S\), put \(u=\chi_1(\gamma_v)\).
One needs separation of \(\bar u\) from the residual complement
even if \(\ell\) belongs to the bad prime set for this particular
type prime. Here \(\ell\ne g\). The local admissibility condition excludes equality of the prime-to-\(\ell\) part of any root of the complementary label and that of the specified character label. The test also holds for the inverse label. The
change to a conjugate marked place can invert either value in
the comparison.

If \(z\) is a projected complementary label and \(b\) the projected character label, the two exclusions give
\[
 z^\delta\ne b^{\delta'}\qquad(\delta,\delta'\in\{1,-1\}),
\]
by simultaneous inversion as necessary. Changing the tame generator
acts on the finite cyclic image by an automorphism and preserves
the exclusions. The residual Frobenius comparison identifies the
complement in the actual integral lattice with the scalar twist
of the good cuspidal one. This scalar is unramified at the type
places away from \(\ell\). Let \(D^{\rm res}_v\) be the residual characteristic
polynomial of the rank-\(r\) complement at \(\gamma_v\).
Thus
\[
      D_v^{\rm res}(\bar u)\ne0,\qquad
      \bar F_v(X)=(X-\bar u)^tD_v^{\rm res}(X).
\]
The multiplicity of \(\bar u\) is exactly \(t\) on every
system with the same residual characteristic polynomials.
This uses all the character factors and equality of the residual
characteristic polynomials.

At a retained point the type comparison gives the rank-\(t\)
inertia summand on which every \(g\in I_v\) acts by
\(\chi_1(g)\), and an invariant complement of rank \(r\).
Indeed the occurring full label has the character block of
Section~\ref{subsec:selected-cut}. The unramified normalization
preserves it. Its distinct tame labels give the inertia
characters and its monodromy vanishes. The local parameter supplies \(t\) scalar copies and reduction can only increase multiplicity. Thus
\[
 t\le\operatorname{mult}_{u}(F_v)\le\operatorname{mult}_{\bar u}(\overline F_v)=t,
\]
where \(F_v\) is the full inertia characteristic polynomial. Local--global
compatibility passes this decomposition to each Galois fibre
after scalar extension. The common complement is invariant for
the entire inertia group.

Form the full characteristic polynomial \(F_v(X)\in\mathcal T[X]\)
of \(R(\gamma_v)\) and divide by the monic polynomial
\((X-u)^t\). In every fraction-field fibre the remainder
vanishes by the decomposition. Injectivity into the product of fibres kills the remainder already over \(\mathcal T\). Thus
\[
 F_v(X)=(X-u)^t D_v(X),\qquad D_v\in\mathcal T[X]\ \text{monic of degree }r.
\]
Its evaluation
at a fibre is the characteristic polynomial on the complement.
The multiplicity computation shows that \(D_v(u)\) has nonzero
residue in the local Hecke algebra.

On each generic fibre write
\[
 V_x=V_{x,B}\oplus V_{x,C}
\]
for the invariant block and complement. Then
\[
 D_v(R_x(\gamma_v))|_{V_{x,B}}=D_v(u)\,\mathrm{id},
 \qquad D_v(R_x(\gamma_v))|_{V_{x,C}}=0,
\]
the latter by Cayley--Hamilton. On the rank-\(t\)
space every \(R(g)-\chi_1(g)\) vanishes. As the first space
is stable under \(R(g)\), the product
\[
       D_v(R(\gamma_v))(R(g)-\chi_1(g))=0
       \qquad(g\in I_v)
\]
holds integrally on the whole family. This constructs the single
polynomial with unit character value needed for the row cocycles.
The unit test uses the residual multiplicity at the localization.
Squarefreeness in characteristic zero by itself only controls
the rational eigenspaces. \label{cor:retained-galois-4}

\begin{lemma}[Inertia at the auxiliary and order-fifteen places]\label{lem:retained-inertia-types}
At \(Q\) and at the additional distinguished places of \(S\),
the zero-monodromy and character conditions give, on every characteristic-zero fibre,
\[
 R_x|_{I_v}\simeq \mathbf1^{\oplus r}\oplus(\chi_1|_{I_v})^{\oplus t}.
\] The power of the chosen inertia generator representing
\(R_x(g)\), for \(g\in I_v\), can be chosen independently of \(x\).
At the distinguished \(S\)-places \(\zeta_v-1\) is a unit, including
in coefficient characteristics three and five.
\end{lemma}

\begin{proof}
At \(Q\) and at the additional distinguished places of \(S\),
the local parameter has zero monodromy and inertia polynomial
\((X-1)^r(X-\chi_1(g))^t\) for each \(g\in I_v\).
These are the unramified rank-\(r\) and prescribed character
types in the norm tuple. Zero monodromy gives finite inertia
image for each fibre by the Weil--Deligne comparison away from
\(\ell\). The displayed polynomial identifies the representation
of inertia with \(1^r\oplus\chi_1^t\): both representations
factor through a finite quotient and their characteristic
polynomials agree there, so Maschke and Brauer--Nesbitt apply
in characteristic zero. In particular one decomposition per
fibre works for all inertia elements. Merely checking the
spectrum of a single matrix would not prove this assertion.

At \(v\in Q\) the first twist factor is unramified and the ramified auxiliary factor has
\[
 \#\chi_1(I_v)=\ell^a\qquad(a\ge0).
\]
At the extra
distinguished prime of \(S\) the auxiliary character is
unramified, and the tame factor or its inverse has inertia
image of order \(15\), with residue norm \(1\bmod 15\).
The construction uses the second distinguished prime when
the first equals \(\ell\). In both cases wild inertia is
trivial on the character: the wild pro-residue-characteristic
group has no quotient of the indicated order. Hence the image
of inertia is the finite cyclic group generated by
\(\zeta_v=\chi_1(\gamma_v)\).

At \(Q\) ramification forces
\(\zeta_v\ne1\). For \(g\in I_v\) choose an integer \(a\)
such that
\[
 \chi_1(g)=\zeta_v^a.
\]
In the simultaneous
decomposition just proved,
\[
 R_x(g)=R_x(\gamma_v)^a
\]
at every \(x\), using the same
\(a\) for them all. At the additional \(S\)-places
\(\zeta_v\) is primitive of order \(15\).
The identity \(\Phi_{15}(1)=1\) makes \(\zeta_v-1\) a unit:
subtract \(\Phi_{15}(\zeta_v)=0\) and factor the difference.
This argument also covers coefficient characteristics three
and five. The common exponent of the tame generator and the unit at the additional places give the two inertia controls needed in the ensuing congruence presentation. \label{cor:retained-galois-5}
\end{proof}

\subsubsection*{A determinant supported on the theta points}
\label{subsec:congruence-determinant}
\hypertarget{source-congruence-determinant}{}

The projections and inertia relations describe the relevant blocks of the integral Galois matrix algebra. A determinant in the Hecke algebra will convert their cohomology lengths into a congruence estimate supported on the theta points. Fix one coefficient DVR \(O\) with
fraction field \(F\). In this calculation write
\(\operatorname{len}\) for the unnormalized \(O\)-length.
Use a finite parameter set \(X\), the theta subset \(Z\), a
Hecke algebra \(\mathcal T\subset O^X\) and a good point
\(\pi\in Z\). Write \(I_Z\) for the ideal of elements of
\(\mathcal T\) vanishing outside \(Z\). The congruence
length before normalization is
\(\operatorname{len}(O/I_Z(\pi))\).

\begin{established}[Ordered corner and block evaluation]\label{hyp:theta-marked-corner}
The auxiliary parameter set \(X'\) records an ordering of
the eigenvalues of a marked \(t\)-dimensional corner. Its
coefficient algebra \(\mathcal T'\subset O^{X'}\) has a
basis \(b_j\) over the pullback of \(\mathcal T\) along
\(p:X'\to X\). For each \(x\in X\setminus Z\), evaluation gives an injection
\[
 O^{t!}\lhook\joinrel\longrightarrow O^{p^{-1}(x)},\qquad
 (c_j)_j\longmapsto\bigg(\sum_j c_j b_j(x')\bigg)_{x'\in p^{-1}(x)} .
\] At the chosen
point of \(X'\) above \(\pi\) the matrix algebra
\(\mathcal A\) has a specialization into matrices over
\(O\). Its image is block triangular with one block of
size \(r\) and \(t\) blocks of size one, and each block
map \(\rho_k:\mathcal A\to M_{d_k}(O)\) is surjective.
The determinants of square matrices over \(\mathcal A\),
computed on the full matrix spaces at all points of
\(X'\), belong to \(\mathcal T'\).
\end{established}
\begin{proof}
\proofstep{Integral coefficients of the matrix determinants.}
These integral properties can be seen from the marked
corner construction. Let \(\mathcal A_0\) be the
\(\mathcal T\)-span of the Galois matrices \(r(g)\),
which is an algebra since products of Galois matrices
are Galois matrices. The coefficients of the
characteristic polynomials of \(r(g)\) are in
\(\mathcal T\). For any finite number of block rows,
expand a matrix over \(\mathcal A_0\) as a scalar linear
combination of \(e_{ij}\otimes r(g)\), with \(e_{ij}\)
the elementary matrices on the block rows. To see that
its characteristic coefficients are in \(\mathcal T\),
give each summand \(A_i=e_{ij}\otimes r(g)\) an
indeterminate \(X_i\). For a primitive cyclic word \(w\), let
\(A_w\) be its ordered matrix product and \(X_w\) its monomial.

In the completed polynomial ring in characteristic zero,
\[
 \log\det(1-\textstyle\sum_iX_iA_i)
 =-\sum_{a\ge1}\frac{\operatorname{tr}((\sum_iX_iA_i)^a)}a
 =-\sum_{w\ {\rm primitive}}\sum_{b\ge1}
                    \frac{X_w^b\operatorname{tr}(A_w^b)}b .
\]
Indeed \(w^b\) contributes \(|w|\) rotations at length \(b|w|\).
Exponentiating gives \(\prod_w\det(1-X_w A_w)\), with only
finitely many factors in each degree.
A product \(A_w\) is zero, has one off-diagonal
block, or has a single diagonal block carrying a
Galois matrix. The individual factors are therefore
\[
 \det(1-X_w A_w)\in\mathcal T[X_w].
\]
For the zero and off-diagonal cases the determinant equals one. After the identity has been
proved in the fraction fields, its coefficient
expression is integral and uses no division in
\(\mathcal T\). Substitution of the scalar coefficients
proves the assertion, also over any \(\mathcal T\)-algebra
of scalar functions.

\proofstep{The ordered corner roots and their splitting algebra.}
The corner operator in \(\mathcal A_0\) is
\(e_Br(g_0)e_B+s\), where \(e_B\) is the marked
rank-\(t\) idempotent and \(s=e_Bse_B\) is the
perturbation with \(s(\pi)=0\). It is zero on the
complement of the corner. The previous argument
places its corner characteristic polynomial in
\(\mathcal T[X]\), with \(X\) here a polynomial
indeterminate. The marked perturbation gives distinct
roots at every fibre in the tied splitting field.
The successive monic quotients adjoining the ordered roots have degrees \(t,t-1,\ldots,1\). Freeness at each step gives
\[
 \operatorname{rank}_{\mathcal T}(\text{splitting algebra})=\prod_{a=1}^{t}a=t!.
\]

Over each rational fibre this algebra is the
product of \(t!\) copies of the coefficient field,
by the Chinese remainder theorem. The free basis
therefore stays independent when evaluated at all
orderings in each fibre. This also gives an embedding
of the splitting algebra into \(O^{X'}\), since the
roots of integral monic polynomials are integral.
Take \(\mathcal A\) to be the span of the Galois
matrices over this splitting algebra. The determinant
argument applies with these coefficients. All
matrices are integral, so \(\mathcal A\) is finite
over \(O\). At the chosen fibre the marked
decomposition gives the triangular blocks, and
the residual full matrix span of the rank-\(r\)
constituent gives surjectivity on that block by
Nakayama's lemma. The scalar blocks have surjectivity
from the coefficient ring. These algebraic steps
use the marked idempotents and the perturbation
on the given family, constructed in Proposition~\ref{prop:retained-integral-projections}.
\end{proof}

\begin{proposition}[The congruence ideal on the theta locus]
\label{prop:theta-determinant-length}
Take the ordered corner and block maps of Lemma~\ref{hyp:theta-marked-corner}. Let \(M\) be the finite
presentation with the column relations in the proof. If these
relations have the prescribed finite column dimensions away from
\(Z\), then
\[
 \operatorname{len}(O/I_Z(\pi))\le
       \sum_k\operatorname{len}\operatorname{Hom}_{\mathcal A}
                          (M,(F/O)^{d_k}).
\]
The \(k\)-th action is via \(\rho_k\), with the same coefficient
ring and chosen point in all the summands.
\end{proposition}

\begin{proof}
Lemma~\ref{hyp:theta-marked-corner} supplies the splitting algebra, its fibre basis and the full block evaluations.

Write
\[
 B=e_Br(g_0)e_B+s
\]
on the corner, and let its ordered
roots be \(\alpha_0,\ldots,\alpha_{t-1}\). The operator separating
the first eigenline is
\[
 E=\prod_{i>0}(B-\alpha_i e_B),\qquad
 E|_{V_{\alpha_i}}=
 \begin{cases}\Delta\,{\rm id},&i=0,\\0,&i>0,\end{cases}
 \qquad\Delta=\prod_{i>0}(\alpha_0-\alpha_i).
\]
Distinctness of the roots gives \(E^2=\Delta E\) on every fibre.
At \(\pi\) the roots are \(\chi(g_0)\varepsilon(g_0)^{-i}\), where
\(\varepsilon\) is the cyclotomic character. The
marked value is \(\varepsilon(g_0)=1+\ell^s v\)
with \(v\in\mathbb Z_\ell^\times\), at \(s\ge2\) for
\(\ell=2\) and \(s\ge1\) for \(\ell=3\).
The binomial calculation gives
\[
     v_\ell(\Delta_\pi)=(t-1)s+v_\ell((t-1)!).
\]
\label{cor:congruence-determinant-1}

\proofstep{The column relations in a finite presentation.}
Here is the module used in the proof. Write \(U_g\) for the
Galois operator after dividing by the first character.
Use \(E,E_1\in\mathcal A\) with \(E^2=\Delta E\) and
\(E_1^2=E_1\), with \(\Delta\in\mathcal T'\) central.
The two operators have rank one over \(F\) at each
point of \(X'\) above \(X\setminus Z\), and
\(\Delta\) is then nonzero. Put
\(Y=\sum_g\mathcal A (U_g-1)E\subset\mathcal A\).

At each column place \(u\) with decomposition subgroup
\(D_u\), adjoin a generator \(n_u\), \(j_u\) generators
\(s_{u,a}\) and a generator \(s_{u,0}'\) to \(Y\).
Impose the relations
\[
 \begin{gathered}
 E n_u=0,\qquad (E_1-1)s_{u,a}=0,\qquad
 (E-\Delta)s_{u,0}'=0,\\
 (U_g-1)E+(U_g-1)n_u
     =\Delta\left(\sum_{a=1}^{j_u}K_{u,a}(g)s_{u,a}
                  +K_{u,0}(g)s_{u,0}'\right)\qquad(g\in D_u).
 \end{gathered}
\]
Here the first term in the last relation is the
element of \(Y\). The functions \(K_{u,a},K_{u,0}\)
are the matrix cocycles in the column data. Denote the
quotient by \(M\). Since \(\mathcal A\) is left
Noetherian, \(Y\) is finitely generated and \(M\)
has a finite presentation.

\proofstep{A kernel vector on each complementary fibre.}
The candidate determinant has integral coefficients. To lie in the theta congruence ideal it must vanish when evaluated on a complementary fibre, which can be tested by producing a vector in its kernel there.

Consider a rational evaluation space \(V\) at \(x'\)
over \(X\setminus Z\). Fix a nonzero vector \(\eta=Ew\)
in the image of \(E\) and a nonzero \(s\) in the image
of \(E_1\). The two conditions on a vector in the first image are incompatible:
\[
 \operatorname{im}E\ \cap\ \bigcap_g\ker(U_g-1)=0.
\]
At \(u\), the
column condition says that the functions
\(K_{u,a}(\cdot)s\) and \(K_{u,0}(\cdot)\eta\)
lie in a space containing the coboundaries and of
dimension at most \(j_u+\dim V^{D_u}\) modulo
coboundaries. Here invariants and coboundaries use \(U\).
If \(V^{D_u}\ne0\), take \(v_u\ne0\) invariant there
and all scalars zero. Otherwise \(V^{D_u}=0\). Let \(L_u\) be the span of the classes of the \(j_u+1\) functions in the quotient by coboundaries. The column bound gives
\[
 \dim L_u\le j_u<j_u+1.
\]

A nontrivial linear dependence in this quotient gives a nonzero tuple \((v_u,(a_{u,a})_a,a_{u,0})\)
such that
\[
 (U_g-1)v_u
      =\sum_a a_{u,a}K_{u,a}(g)s+a_{u,0}K_{u,0}(g)\eta
           \qquad(g\in D_u).
\]
If at some \(u\) we can take \(Ev_u=0\), map \(Y\)
and all other places to zero and map the generators
at \(u\) to \(\Delta v_u,a_{u,a}s,a_{u,0}\eta\).
This satisfies the relations and gives a nonzero map
\(M\to V\). In the other case rescale each tuple so
that \(Ev_u=\eta\). Map \(y\in Y\) to \(yw\) and
the local generators to
\(\Delta v_u-\eta,a_{u,a}s,a_{u,0}\eta\).
The relations follow from \(E\eta=\Delta\eta\)
and the equation at \(D_u\). The map is nonzero
on \(Y\) because \(\eta\) is moved by some \(U_g\).
Thus
\[
 \operatorname{Hom}_{\mathcal A}(M,V)\ne0
\]
at every evaluation point over \(X\setminus Z\).

\proofstep{Blockwise Smith lengths at the chosen theta point.}
Present \(M\) with \(g\) generators and a finite
matrix of relations \(P\), written so that applying
\(\operatorname{Hom}_{\mathcal A}(-,V)\) identifies it
with the kernel of \(P\) acting on \(V^g\). At the
chosen point above \(\pi\), use the block targets
\[
 H_k=(F/O)^{d_k}
\]
with action via \(\rho_k\).
If their Hom lengths are finite, the expanded relation
matrix \(P_k\) in each block has full column rank.
Let \(\varpi^{a_{k,i}}\) be its elementary divisors, with
\(1\le i\le gd_k\). Since
\[
 (F/O)[\varpi^a]\simeq O/(\varpi^a),
\]
the kernel on the divisible target and its determinantal ideal satisfy
\[
 \begin{aligned}
 \operatorname{len}\operatorname{Hom}_{\mathcal A}(M,H_k)
     &=\sum_{i=1}^{gd_k}a_{k,i}=:s_k,\\
 I_{gd_k}(P_k)&=(\varpi^{s_k}),
 \end{aligned}
\]
where \(I_{gd_k}\) is generated by the maximal minors. Thus one such
minor attains valuation \(s_k\).

For each \(k\), there exists a contraction of the rows
giving a square \(gd_k\)-matrix with determinant
valuation \(s_k\). Surjectivity of \(\rho_k\) lifts
such a contraction to a rectangular matrix \(C_k\)
over \(\mathcal A\). Use the matrix
\(C(\mathbf u)=\sum_j u_j C_j\) over a polynomial
ring with indeterminates \(u_j\).
In block \(k\), the determinant of \(CP\), divided
by a uniformizer to the power \(s_k\), is a polynomial
over \(O\). Indeed Cauchy--Binet expresses the numerator as
\[
 \det(C(\mathbf u)_kP_k)=
   \sum_{|J|=gd_k}\det(C(\mathbf u)_{k,*,J})\det(P_{k,J,*}).
\]
Each term is divisible by \(\varpi^{s_k}\) in \(O[\mathbf u]\).
Its reduction is nonzero since setting \(u_k=1\)
and the other \(u_j=0\) gives a unit. Therefore the
product of the reductions is a nonzero polynomial,
and some coefficient is a unit.

\proofstep{A determinant coefficient in the original congruence ideal.}
The same characteristic coefficient argument applied
over \(\mathcal T'[\mathbf u]\) puts the determinant of
\(C(\mathbf u)P\) in \(\mathcal T'[\mathbf u]\). At every
point above \(X\setminus Z\) the nonzero maps \(M\to V\)
give a kernel vector, so the determinant vanishes
as a polynomial at those points. At the data at
\(\pi\) the determinant is the product of the block
determinants by triangularity. Take a coefficient
\(f'\in\mathcal T'\) corresponding to a unit coefficient
of the reduced product. Its valuation at this point
is \(\sum_k s_k\), and it vanishes off the theta subset.
Write
\[
 f'=\sum_j(f_j\circ p)b_j.
\]

For every
\(x\notin Z\), its evaluations give a vanishing vector:
\[
 0=(f'(\xi))_{\xi\in p^{-1}(x)}
   =\sum_j f_j(x)(b_j(\xi))_{\xi\in p^{-1}(x)} .
\]
Independence on that fibre forces \(f_j(x)=0\) for
all \(j\). Thus \(f_j\in I_Z\), and the values \(b_j(\pi')\)
are integral in the common coefficient ring, so \(f'(\pi')\)
lies in \(I_Z(\pi)\). Its valuation gives
\[
 \operatorname{len}(O/I_Z(\pi))
 \le \sum_k\operatorname{len}\operatorname{Hom}_{\mathcal A}(M,H_k).
\]
If a Hom length is infinite this inequality already
holds in extended lengths. The empty generator case
has determinant one. Division by the ramification
index of \(O/\mathbb Z_\ell\) gives the normalized
inequality.  Choose a residue field
of size greater than \(g\sum_k d_k\) to specialize the
contraction simultaneously. Coefficient extraction gives
the same ideal bound over any residue field. \label{cor:congruence-determinant-2}

\end{proof}

\subsubsection*{Choosing the column cocycles}

The determinant argument requires actual columns in the integral presentation. A finite quotient specifies their evaluations only up to torsion, so we choose lifts and bound the exponent of a relation among them.

\begin{lemma}[Column lifts and their torsion tests]
\label{lem:congruence-column-cocycles}
Under the finite-class lift bound \(e'\) on the designated local
lattices, write \(c_a\) for a lift of the \(a\)-th finite quotient basis class and \(c_*\) for the extra column. Their evaluations are
\[
 c_a\longmapsto\ell^{e'}\beta_a\ (1\le a\le j_u),\qquad c_*\longmapsto\ell^{e'}\Delta_\pi\quad\text{in the unramified trivial block}.
\] These columns
satisfy the finite-class dimension bound away from \(Z\) and their
torsion dependences in the \(W\)-block are killed by \(\ell^{e'}\).
\end{lemma}

\begin{proof}
\proofstep{Integral finite bases at a column place.}
At the selected point let \(T\) be the integral lattice
in the \(W\)-block. Choose a basis \((\beta_a)_a\) of
the free quotient of \(H^1_f(E_u,T)\). This quotient
is the lattice of integral finite classes in
\(H^1_f(E_u,W)\). It spans the rational finite group
since a continuous cocycle on the compact local Galois
group has bounded image and can be scaled into \(T\).
Thus its rank is determined on the rational finite part:
\[
 \operatorname{rank}_O\bigl(H^1_f(E_u,T)/\mathrm{tors}\bigr)=\dim H^1_f(E_u,W)=j_u.
\]


Use the left ideals \(\mathcal A E_1,\mathcal A E\)
with local action by \(U_g\). Rationally these are
subrepresentations of finite direct sums of the full
twisted family, by embedding matrices as columns in
the evaluation spaces. In block \(0\), choose a
generator \(v_1\) of the image of \(E_1\) on \(T\).
Evaluation on \(v_1\) maps \(\mathcal A E_1\) onto
\(T\), because the block image is the full matrix
algebra. Evaluation of \(\mathcal A E\) on the
trivial block has image \(\Delta_\pi O\), on which
the twisted action is trivial. The two block evaluations are thus surjections of stable lattices
\[
 \mathcal A E_1\twoheadrightarrow T,\qquad \mathcal A E\twoheadrightarrow\Delta_\pi O.
\]

The twisted Galois action on the second target is trivial.
The first column lift produces finite cocycles in
\(\mathcal A E_1\) with images representing
\(\ell^{e'}\beta_a\), using lifts of the \(\beta_a\)
to integral finite classes. Use these as \(K_{u,a}\).
The second lift uses the unramified class in
\(H^1_f(E_u,\Delta_\pi O)\) with value \(\Delta_\pi\)
on arithmetic Frobenius. This class is finite since
the trivial representation's rational finite part
is the unramified line. Its lift \(K_{u,0}\) has image
with Frobenius value \(\ell^{e'}\Delta_\pi\). This
value holds on every lift of that Frobenius because
coboundaries of the trivial module vanish.

\proofstep{The local dimension at every evaluation.}
At an evaluation point \(x'\), a vector fixed by \(E_1\)
and an eigenvector \(E\eta=\Delta\eta\) give the equivariant
evaluation maps \(a\mapsto as\) on \(\mathcal A E_1\) and
\[
 a\mapsto a\eta
\]
on \(\mathcal A E\), respectively
(\(E_1s=s\)). Functoriality sends their finite cohomology
classes to finite classes. The negative Hodge jumps of the full
twisted family are constant at the column points. Removing the
cyclotomic characters at \(\pi\) removes no negative jumps.
As \(H^0(E_u,W)=0\), their common total over embeddings is
\(j_u\). The finite-part dimension formula at \(x'\) is therefore
\[
        \dim H^1_f(E_u,V)=j_u+\dim V^{D_u}.
\]
This is the local dimension used for the finite column cocycles.

\proofstep{Torsion relations in the chosen basis.}
The rational dimension leaves the exponent of integral torsion relations undetermined.

Test a torsion relation on the \(W\)-block. Write a linear combination
of the evaluated column cocycles with coefficients
in \(F/O\), and lift those coefficients to \(F\).
If the combination is a coboundary modulo \(T\),
the rational finite class \(\sum_a y_a\ell^{e'}\beta_a\)
comes from \(H^1(E_u,T)\) by exactness for
\(0\to T\to W\to W/T\to0\).
It therefore comes from \(H^1_f(E_u,T)\) by the
definition of the integral finite part. Because
the \(\beta_a\) form its free basis, each
\(y_a\ell^{e'}\) belongs to \(O\). This proves the
killing exponent for such a dependence.   \label{cor:congruence-determinant-3}

\end{proof}

\subsubsection*{The Selmer length comparison for this presentation}
\label{subsec:congruence-selmer}

We can now evaluate the presentation at the chosen theta point. Each block contributes a Selmer length, while the valuation of the specialized determinant records the cost of imposing the integral conditions.

At the chosen point let \(\Delta_\pi\ne0\) be the
specialization of \(\Delta\) and put \(v_\Delta
=v_\ell(\Delta_\pi)\). From now on lengths are
normalized by \(v_\ell(\ell)=1\).
The block targets are \(H_0=(F/O)^r\) with the
action of \(W\), \(H_1=F/O\) with trivial action,
and \(H_k=(F/O)(1-k)\) for \(2\le k\le t\).
Write \(P_k=\operatorname{Hom}_{\mathcal A}(M,H_k)\).
\begin{proposition}[From block lengths to Selmer kernels]
\label{prop:congruence-block-selmer-length}
With the integral presentation and column lifts of
Lemma~\ref{lem:congruence-column-cocycles}, the normalized sum of
block Hom lengths in Proposition~\ref{prop:theta-determinant-length}
is bounded by the global kernel lengths and the local and
presentation costs below. The three global kernels are
the rank-\(r\) finite kernel, the unramified first character and
the strict classes in twists \(1-k\) for \(2\le k\le t\).
If \(c_k\) is the normalized contribution charged to coboundaries in block \(k\), then
\[
 c_1=0,\qquad c_k\le2d_kv_\Delta\quad(k\ne1).
\]
\end{proposition}

\begin{proof}
\proofstep{The map to continuous cocycles.}
A homomorphism \(\varphi\) supplies the cocycle
\[
    c_\varphi(g)=\varphi((U_g-1)E).
\]
The cocycle identity follows by multiplying \(U_g U_h\).
In the integral presentation all the matrix functions
are continuous, and the orders and their modules are
finite over \(O\). A map to \(H_k\) kills a common
power of the uniformizer times the module because the
images of a finite set of generators are torsion.
Thus \(c_\varphi\) is continuous.

\proofstep{Its kernel on the local generators.}
First consider the map from \(P_k\) to cocycles before
taking cohomology. Its kernel consists of maps vanishing
on \(Y\). They are determined by the images of the local
generators, subject to the relations of \(M\).
The specializations on the target blocks satisfy
\[
 E|_{H_0}=0,\qquad E|_{H_1}=\Delta_\pi,\qquad E|_{H_k}=0\quad(k\ge2).
\]
The idempotent \(E_1\) has rank one
in block \(0\) and is zero on the character blocks.
The extra column cocycle \(K_{u,0}(g)\) lies in the
left ideal generated by \(E\).

\begin{itemize}
\item In block \(0\) the column lift input says that
a linear combination \(\sum_a K_{u,a}(g)s_a\) with
\(s_a\in E_1H_0\) which is a coboundary on \(D_u\)
forces \(\ell^{e'}s_a=0\) for every \(a\).
For a map vanishing on \(Y\) the combination with
coefficients \(\Delta_\pi\varphi(s_{u,a})\) is a
coboundary by the relation for \(D_u\). The tuple of images of the \(s_{u,a}\) lies in
\[
   (E_1 H_0)[\ell^{e'}\Delta_\pi]^{\,j_u},
\]
of normalized length at most \(j_u(e'+v_\Delta)\).

For a fixed tuple, differences of choices of \(n_u\) lie in
\(H_0^{D_u}\), of length \(h_W(u)\), and the remaining
image of \(s'_{u,0}\) lies in \(H_0[\Delta_\pi]\), of length
at most \(rv_\Delta\). Kernel plus image for projection onto
the tuple therefore gives
\(j_u(e'+v_\Delta)+h_W(u)+rv_\Delta\).

\item In block \(1\) the matrix \(K_{u,0}\) has some
value on \(D_u\) acting by
\(\ell^{e'}\Delta_\pi\) times a unit, the first-character
column input. The relation for that element gives
\[
 \ell^{e'}\Delta_\pi^2\varphi(s_{u,0}')=0
\]
for a
map vanishing on \(Y\). The image of \(n_u\) is killed
by \(\Delta_\pi\) and that of \(s_{u,a}\) is zero.
The resulting cost is at most \(e'+3v_\Delta\).

\item In block \(k\ge2\) the \(s_{u,a}\) vanish and
\(\Delta_\pi\varphi(s_{u,0}')=0\). The relation for
\(D_u\) forces \(\varphi(n_u)\in H_k^{D_u}\), since
\(K_{u,0}\) acts by zero. With the scalar invariant
bound \(h_0\) the cost is \(h_0+v_\Delta\).
\end{itemize}

\proofstep{The cost of quotienting by coboundaries.}
Coboundaries change the length differently on each block.

Pass next to cohomology. In block \(1\) the coboundaries
are zero. In the other blocks the central scalars \(a_g\) of the marked corner satisfy
\[ E=\sum_g a_g U_g,\qquad \sum_g a_g=\Delta_\pi\quad\text{at the chosen point}. \]
If \(c_\varphi(g)=(U_g-1)v\), the element
\[
    y=\sum_g a_g (U_g-1)E=(E-\textstyle\sum_g a_g)E
\]
has image \(-\Delta_\pi v\) in \(H_k\), since \(E\)
acts by zero on \(H_k\). The algebra identity
\(Ey=\Delta y\) then forces \(\Delta_\pi^2 v=0\).
The space of image cocycles that are coboundaries
has length at most \(2d_k v_\Delta\).

\proofstep{Relative Selmer kernels and local correction terms.}
Finally let \(C_k\) be the image of \(P_k\) in
\(H^1(E,H_k)\). At each local place impose a target
subgroup \(F_{k,u}\) of local cohomology. Suppose
the restrictions of \(C_k\) are in subgroups
\(B_{k,u}\supset F_{k,u}\) and that the classes also
lie in a global subgroup \(Q_k\). The kernel of
\[
   C_k\longrightarrow\bigoplus_u B_{k,u}/F_{k,u}
\]
lies in the subgroup of \(Q_k\) satisfying all the
local conditions. Thus its global length bound and
the sum of the local colength bounds give a length
bound for \(C_k\). In the split comparison, the
column images already satisfy \(F_{k,u}\).

In
block \(0\) the row quotient has the colength
bound proved above, and the global kernel has length
at most \(s_\pi\). In the character blocks the row and
tame quotients have bounds \(([E_w:\mathbb Q_\ell]+2)c_B\)
and \(v_\ell(\zeta_v-1)\). If \(K_1\) is the global kernel in block one, unramified class field theory gives
\[
 K_1\hookrightarrow\operatorname{Hom}(\operatorname{Cl}(E),F/O),\qquad \operatorname{len}_{\rm norm} K_1\le v_\ell(h_E).
\]
In the higher blocks the local
conditions defining this kernel are strict at \(\ell\)
and unramified elsewhere. Its length is bounded by the
strict group \(\operatorname{Sha}^1(E,(F/O)(1-k))\).
The normalization and the row and tame restriction arguments
follow below.

\proofstep{The sum of the normalized block costs.}
Let \(k_c\) be the number of columns. Summing these
bounds over all blocks gives the expression used in
the split-row estimate. The terms from the rank-\(r\)
and character blocks contribute
\[
 \sum_u\{j_u(e'+v_\Delta)+h_W(u)+e'\},
\]
and the \(t-1\) higher blocks contribute at most
\(tk_ch_0\). The remaining denominator terms from
the presentation and the coboundaries are bounded by
\(2k_cv_\Delta+4R(k_c+1)v_\Delta\). Indeed the dimensions \(d_k\) of the blocks sum to \(R\), so their contribution is
\[
 \sum_k d_k k_c v_\Delta=R k_c v_\Delta,
\]
and the coboundaries cost at most \(2Rv_\Delta\).

The row and tame costs of the
character blocks occur at most \(t\) times. Adding the
global kernels accounts for \(s_\pi\), the class
number and the strict cohomology sum. Pairing the row
cost for \(W\) at \(w\) with its conjugate column
\(u\) and using \(j_u\le [E_u:\mathbb Q_\ell]m\)
gives the formula below.
\label{cor:congruence-determinant-4}
\end{proof}

\subsubsection*{Changing the frame at the selected point}
\label{subsec:split-selmer}
\hypertarget{source-split-selmer}{}

To identify those abstract block lengths with the prescribed local conditions, choose integral coordinates adapted to the decomposition at the good point. The surrounding family still includes the points off the theta locus that enter its congruence ideal.

Use the integral family \(R_x\) at the fixed Hecke
localization, including the points outside the theta subset.
At the good theta point \(\pi\) the given rational decomposition
is an isomorphism with \(\rho_\pi\oplus\bigoplus_{k=1}^t\chi_k\).
Here \(\rho_\pi\) is the specified integral rank-\(r\) representation
and \(\chi_k/\chi_1=\epsilon^{1-k}\) for the cyclotomic character
\(\epsilon\). Replace \(R_\pi\) by the direct sum in these
integral bases and keep all other matrices at their own points.
\begin{lemma}[Change of frame in the family algebra]
The rational change of basis identifies the two integral matrix
spans, their marked corners and congruence ideals. At the chosen point \(\pi\), evaluation on the complement gives
\[
 \operatorname{ev}_{\pi,W}(\mathcal A)=M_r(O),
\]
where \(\mathcal A\) is the normalized integral span.
\end{lemma}
\begin{proof}
If \(C_x\) changes the frame at \(x\), conjugation sends each generator
\(R_x(g)\) to \(R'_x(g)\). Thus it carries
\[
 \left(\sum_j a_j(x)R_x(g_j)\right)_x
 \longmapsto\left(\sum_j a_j(x)R'_x(g_j)\right)_x .
\]
The coefficients are the very same Hecke functions. Applying
inverse conjugation proves equality of the integral spans.
Characteristic coefficients are unchanged. Row embeddings and
column quotients transport by rational conjugation, with the
specified lattices of \(\rho_\pi\) and its twisted dual as targets.
The resulting finite family is continuous at each point.
\label{cor:split-selmer-1}
The marked data before the change of frame consist of
orthogonal summand projections at \(\pi\), extended as integral
idempotents in the family span, and a corner matrix \(X_0\) with
distinct roots in every fibre. The fibrewise ranks and the two specializations are
\[
 \operatorname{rank}e_B=t,\qquad \operatorname{rank}e_1=1.
\]

At \(\pi\) the first projects onto the character block and the second onto a line of \(\rho_\pi\). These data can be
transported from an integral marked model on the same coefficient
localization. Write \(R_x^{\rm m}\) for the marked model at \(x\).
The shared Hecke character gives
\[
 \det(1-TR_x(\operatorname{Frob}_v))
   =\det(1-TR_x^{\rm m}(\operatorname{Frob}_v))
\]
at the unramified Frobenius classes. Taking traces, then using
continuity and Chebotarev density, gives equality of the trace
functions on the entire Galois group. Brauer--Nesbitt applied
to the two semisimple rational representations gives an
intertwiner at every \(x\). At \(\pi\), cancelling the specified
character summands identifies the rank-\(r\) complements.
Write \(C\) for the resulting tuple of changes of frame and
\(\mathcal A^{\rm m},\mathcal A'\) for the two integral spans.
The generator identity above gives
\begin{equation}\label{eq:selmer-transport-corner}
 C\mathcal A^{\rm m} C^{-1}=\mathcal A',\qquad
 C(e\mathcal A^{\rm m}e)C^{-1}
      =(CeC^{-1})\mathcal A'(CeC^{-1})
      \quad(e=e_B,e_1).
\end{equation}
By \eqref{eq:selmer-transport-corner} the corner matrix is conjugated in every fibre. Its
characteristic roots, the two idempotent ranks and their
orthogonality are preserved.  \label{cor:split-selmer-2}

The cyclotomic test
\(g_0\) lies in the prescribed inertia and has
\[
 \epsilon(g_0)=1+\ell^s u,
\]
with \(u\) a unit. In the new basis
\(X_0(\pi)\) is \(e_B R_\pi(g_0)e_B\). This identification follows
because both models of the direct sum transport the action of
every sum \(\sum_j a_j(\pi)R_\pi(g_j)\) to the same block action.
Thus the perturbation \(X_0-e_B R(g_0)e_B\) lies in the new span,
is supported on the corner, and vanishes at \(\pi\). The ranks,
corner characteristic polynomials and orthogonality relations
are unchanged by conjugation. The matrix-span and ordered-root
construction consequently applies with the same Hecke localization
and the same marked point.

Let
\[
 B_\pi\subseteq M_r(O)
\]
be the image of the normalized algebra
on the specified rank-\(r\) lattice. Absolute irreducibility of its
reduction gives the full matrix algebra over the residue field by
the density theorem. Writing \(\varpi\) for a uniformizer, this says
\[
 B_\pi+\varpi M_r(O)=M_r(O),\qquad
 M_r(O)/B_\pi=\varpi\bigl(M_r(O)/B_\pi\bigr).
\]
Nakayama's lemma applies to this finite \(O\)-module. Thus \(B_\pi=M_r(O)\): the localized order contains every integral endomorphism of the complementary lattice used in the row calculation.
Section~\ref{subsec:retained-galois} constructs the marked model.
Proposition~\ref{prop:retained-rank-complement} proves residual
absolute irreducibility at this localization.

\end{proof}
\subsubsection*{The lattice used in the rank-\(r\) row quotient}

The new coordinates identify the residual complementary block. Computing its row quotient requires the particular lattice on which its normalized Galois action operates, together with the dual action at the conjugate place.

Write \(U_g=\chi_1(g)^{-1}R(g)\), \(Y=\sum_g\mathcal A(U_g-1)E\),
and \(y(g)=(U_g-1)E\) in the normalized algebra. Let \(T\) be
the specified lattice of the \(U\)-action on the rank-\(r\) block,
so \(W=T[1/\ell]\). Write \(T^*=\operatorname{Hom}_O(T,O)(1)\)
and \(W^*=T^*[1/\ell]\). At a row place the dual finite-class map uses the source lattice and right action
\[
 P=e_1\mathcal A,\qquad g\cdot b=\epsilon(g)bU_g^{-1}.
\]
To see its target, write
\(e_1(\pi)=v\lambda\) with \(\lambda(v)=1\) over \(O\) in the
rank-\(r\) block. An element \(b\in P\) has evaluation \(v\lambda_b\)
on this block. The map \(b\mapsto\lambda_b\) onto \(T^*\) is
surjective because the normalized algebra maps onto \(M_r(O)\).
\begin{lemma}[The local row quotient]
Evaluation \(p:P\to T^*\), \(b\mapsto\lambda_b\), is an integral
surjection for the twisted dual action. With
\(A=W/T,\ L_w=H^1_f(E_w,T^*)\) and
\(J_w=p_*H^1_f(E_w,P)\), put
\(F_w=\operatorname{im}(H^1_f(E_w,W)\to H^1(E_w,A))\).
Local Tate duality induces
\[
 J_w^\perp/F_w\simeq\operatorname{Hom}_O(L_w/J_w,F/O).
\]
\end{lemma}
\begin{proof}
Since the block maps onto \(M_r(O)\), every row \(\lambda_b\) occurs.
If \(\lambda_b\) is a row functional, transposition converts its right action into the column-coordinate action
\[
 \lambda_b^{\mathsf T}\longmapsto\epsilon(g)W(g)^{-\mathsf T}\lambda_b^{\mathsf T}.
\]
This is precisely the action on \(W^*\). For the standard basis \(e_i\), integrality is the coordinate test
\[
 u(O^r)\subseteq O\quad\Longleftrightarrow\quad
 (u(e_1),\ldots,u(e_r))\in O^r
 \qquad(u\in\operatorname{Hom}_F(F^r,F)).
\] Restriction through the
chosen completion identifies the decomposition group with the
local absolute Galois group. The coefficient isomorphism there
carries the fraction field and DVR simultaneously, so it transports
continuous cocycles, coboundaries and the finite conditions with
the lattice. Recording matrix rows at the ordered points embeds the rational source in a finite sum of the evaluated twisted duals. This is the row-colength map.  \label{cor:split-selmer-3}

Use local Tate duality in its \(F/O\)-valued form, obtained from
the usual \(\mathbb Q_\ell/\mathbb Z_\ell\)-valued pairing by the
trace and a generator of the inverse different. An integral class
annihilates \(F_w\) precisely when its rational image pairs to zero
with \(H^1_f(E_w,W)\): any nonzero rational value can be scaled to
a nonintegral one. By Bloch--Kato orthogonality this condition
cuts out \(L_w\). Double annihilation gives
\[
 F_w=L_w^\perp.
\]

The restriction
\[
 J_w^\perp\to\operatorname{Hom}_O(L_w/J_w,F/O)
\]
has this kernel and is onto by injectivity of \(F/O\) over the DVR.
The extension is continuous on the finitely generated integral
cohomology group. Thus the normalized length equals that of
\(L_w/J_w\), bounded by the row index at this same place.
\cite{bloch-kato}
\end{proof}

\begin{lemma}[Annihilation of the integral row classes]
For an \(\mathcal A\)-linear \(\varphi:Y\to A\), put \(c(g)=\varphi(y(g))\). Local Tate duality gives
\[
 \operatorname{res}_w[c]\in J_w^\perp.
\]
\end{lemma}
\begin{proof}
To prove annihilation of \(J_w\), let \(a\) be a continuous
row cocycle with values in \(P\), and let \(c(g)=\varphi(y(g))\)
with values in \(A\). The multiplication cochain
\[ z(g,h)=a(g)U_g y(h) \]
takes values in \(e_1Y\), with the action on this lattice
given by the scalar \(\epsilon\).

The contraction \(z\) is a continuous two-cocycle. This follows by the cocycle identity for \(a\)
with action \(b\mapsto\epsilon(g)bU_g^{-1}\) and the
identity \(y(gh)=y(g)+U_g y(h)\). Its image in
\(e_1\mathcal A E(1)\) is a coboundary. Indeed the
one-cochain \(d(g)=a(g)U_g E\) satisfies
\[
 \epsilon(g)d(h)-d(gh)+d(g)=-a(g)U_g(U_h-1)E=-z(g,h).
\]
The inclusion \(e_1Y\subseteq e_1\mathcal A E\) induces an
injection on continuous \(H^2\) after the Tate twist. Both
lattices are finite free over \(O\), and the local invariant
map
\[
 H^2(E_w,\mathbb Z_\ell(1))\simeq\mathbb Z_\ell,
\]
extended to finite free coefficient modules, identifies the
induced map with the original lattice inclusion. This works
also for a nonsaturated inclusion. Therefore the cochain
\(z\) is already a coboundary in \(e_1Y(1)\).

Apply to that smaller lattice the map \(m\mapsto\lambda(\varphi(m))\).
The pairing is well defined modulo \(O\), and is equivariant
for the scalar Tate action. Since \(\varphi\) is
\(\mathcal A\)-linear, the image of \(z(g,h)\) is
\(\langle p(a(g)),U_g c(h)\rangle\), with
\[
 p:P\to T^*
\]
the row quotient.

This is the cup pairing in the reversed
order, which differs from the desired order by the sign
for two degree-one classes. It follows that \([c]|_{E_w}\)
annihilates the image of every such \([a]\), in particular
\(J_w\). The local class of \(c\) is continuous, as the
universal cocycle is continuous in its finite integral
lattice and a linear map from that lattice into \(A\) has
finite image and is continuous. Consequently the integral index \(\operatorname{len}(L_w/J_w)\) bounds the excess over the finite condition on the complementary row. The character rows require their own local conditions. \label{cor:split-selmer-4}
\end{proof}
\subsubsection*{Character row and column conditions}

The complementary row has now been accounted for. On the remaining rank-one blocks the cyclotomic twists distinguish the finite conditions, and the row and column impose different restrictions even at the same coefficient prime.

\begin{proposition}[Local conditions on the character blocks]
At a column the specialization is finite on block \(0\),
unramified on block \(1\), and zero on blocks \(k\ge2\).
At a row above \(\ell\), for \(1\le k\le t\) let \(L_k\) be the normalized length of the excess over the finite subgroup on block \(k\). Then
\[
 L_k\le([E_w:\mathbb Q_\ell]+2)c_B,\qquad c_B=40(t+1)(n+1)\lambda_N(h).
\]
\end{proposition}
\begin{proof}
At a column, the last defining relation of \(M\) reads in local
cohomology as a combination of the classes \(K_{u,a}\) and
\(K_{u,0}\) after evaluation. In block \(0\) the first classes
are finite and the last term vanishes because \(K_{u,0}\) has a
right factor \(E\). In block \(1\) the first terms vanish because
\(E_1\) acts by zero, and the last class is unramified. In
blocks \(k\ge2\) both sets of terms vanish. The generator
\(n_u\) contributes a coboundary in each case. These conclusions
hold at the normalized point and use the specific column lifts
of the determinant construction.

For the character rows, form the stable lattice \(Y\oplus O\)
with action
\[
  (z,a)\longmapsto (U_g z+a\,y(g),a).
\]
With \(e=(0,1)\), the quotient is the trivial line, and
\(ge-e=(y(g),0)\). Put \(i(z,a)=(z+aE,a)\).
The cocycle is \(y(g)=(U_g-1)E\), so
\[
    i(U_g z+a\,y(g),a)=(U_g(z+aE),a)=g\,i(z,a).
\]
The map \(i\) is injective, with the left \(U_g\)-action
on matrices and the trivial action on the second factor.

Evaluation
on the matrix columns embeds the rational space of this lattice
into a finite sum of the twisted family \(R_x^\circ\) and one
trivial line. The trivial character occurs in \(R_\pi^\circ\). The target of
the embedding is therefore crystalline, and so is its rational
subrepresentation \(V=(Y\oplus O)\otimes_O F\). For each
embedding \(\tau\), strictness of the crystalline filtration gives
\[
 \operatorname{Jumps}_\tau(V)
     \subseteq \bigcup_x \operatorname{Jumps}_\tau(R_x^\circ)
     \subseteq [-h,h].
\]
The union on the right over all embeddings has cardinality at
most \(n\) by the common row bound, with \(t\le h\).

The character lifting bound is used with that lattice, the
surjection onto the trivial line and the lift \(e\) just given.
It applies to an equivariant map of the kernel \(Y\) into the
divisible quotient of a rank-one lattice acted on by
\(\epsilon^{1-k}\), for \(1\le k\le t\). Extend this integral
map \(O\)-linearly to the full rational source using injectivity
of \(F/O\) over \(O\). The comparison only requires equivariance
on the kernel lattice.

Indeed, if \(\phi\) is the equivariant map on \(Y\) and
\(\widetilde\phi\) any such extension, then
\begin{equation}\label{eq:selmer-row-kernel-cocycle}
 c(g)=\widetilde\phi(ge-e)=\phi(y(g)),\qquad
 c(gh)=c(g)+\epsilon^{1-k}(g)c(h).
\end{equation}
Both identities in \eqref{eq:selmer-row-kernel-cocycle} use only the restriction to \(Y\). The completion lies in
the local level-\(\operatorname{lcm}(\ell,N)\) cyclotomic field,
the coefficients contain those roots of unity, and \(4\mid N\)
at \(\ell=2\). With
\[
 c_B=40(t+1)(n+1)\lambda_N(h),
\]
the
comparison places \(\ell^{c_B}[c]\) in the finite subgroup of
the target. For \(k=1\) this subgroup is unramified. For \(k>1\)
it is zero, by the rank-one jump \(k-1>0\) and the absence of
rational invariants in the Bloch--Kato dimension formula.
\cite{bloch-kato}\label{cor:split-selmer-5}

For the length allowance, let \(F_{k,w}\) be this finite
subgroup. It is divisible, as the image of a vector space.
Choose \(b\in F_{k,w}\) with \(\ell^{c_B}b=\ell^{c_B}[c]\).
Then \([c]-b\) represents the same class modulo \(F_{k,w}\) and
belongs to
\[
 H^1(E_w,A_k)[\ell^{c_B}]
   =\operatorname{im}\bigl(H^1(E_w,A_k[\ell^{c_B}])
                          \longrightarrow H^1(E_w,A_k)\bigr),
\]
by the cohomology sequence for multiplication by \(\ell^{c_B}\). For \(M_k=A_k[\ell^{c_B}]\), local Tate duality and the bound on
the coefficient module give \(\operatorname{len}H^i(E_w,M_k)\le c_B\)
for \(i=0,2\). The local Euler characteristic gives
\[
 \operatorname{len}H^1(E_w,M_k)
 = [E_w:\mathbb Q_\ell]c_B+
   \operatorname{len}H^0(E_w,M_k)+\operatorname{len}H^2(E_w,M_k).
\]
This gives an allowance for each character row even when a sharper row estimate is available. It leaves the inertia conditions away from the coefficient prime to be imposed.
\end{proof}
\subsubsection*{Tame tests on the character blocks}

It remains to control ramification away from the split row and column places. The tame inertia relations on each character block bound the quotient by unramified classes at the auxiliary level primes.

For a character block let \(A_k=(F/O)(1-k)\) and let
\(c(g)=\varphi(y(g))\) for an \(\mathcal A\)-linear map
from the column presentation into \(A_k\).

\begin{lemma}[Tame inertia and the remaining global condition]
Away from the split row and column places the character class is
unramified outside \(Q\); at \(v\in Q\) the ramified quotient \(D_v\) satisfies
\[
 \operatorname{len}(D_v)\le v_\ell(\zeta_v-1).
\] After imposing
the row, column and unramified conditions the residual length is at
most \(v_\ell(h_E)\) for \(k=1\); for \(k\ge2\) it is bounded by
the strict length with exponent \(k-1\) over the same field.
\end{lemma}
\begin{proof}
At a place
outside the split row and column places and the auxiliary
sets \(S,Q\), the family inertia acts trivially on every
fibre and \(\chi_1\) is trivial there. The last assertion
can be checked on its summand at \(\pi\). Thus \(y(g)=0\)
on this inertia.

At \(v\in Q\), the chosen inertial element
\(\gamma_v\) maps every fibre into the two eigenspaces for
\(1\) and \(\zeta_v=\chi_1(\gamma_v)\). The defining conditions on the tame generator give \(\zeta_v\ne1\), equality of all
\(\chi_k\) on inertia, and a common power \(n(g)\) of
\(R(\gamma_v)\) and of \(\chi_1(\gamma_v)\) for the actions
of each \(g\) on this inertia. The common power comes from
the finite image of the chosen tame character. In the family
algebra we have
\[
 (R(\gamma_v)-1)y(\gamma_v)
 =\zeta_v^{-1}(R(\gamma_v)-1)(R(\gamma_v)-\zeta_v)E=0.
\]

Evaluating under \(\varphi\) gives
\((\zeta_v-1)c(\gamma_v)=0\). On the target \(U_{\gamma_v}\) acts trivially. On ordinary powers the
cocycle therefore satisfies
\[
 c(\gamma_v^a)=\sum_{i=0}^{a-1}U_{\gamma_v}^ic(\gamma_v)=a\,c(\gamma_v)
 \qquad(a\ge1).
\]
Continuity on the procyclic quotient gives \(c(g)=n(g)c(\gamma_v)\).
Coboundaries vanish on inertia here. Therefore the quotient
by the unramified condition injects, by evaluation at
\(\gamma_v\), into \(A_k[\zeta_v-1]\), of normalized length
\(v_\ell(\zeta_v-1)\). This bound uses only \(\zeta_v-1\ne0\).

At \(v\in S\), there is an integral family operator \(a\)
such that
\[
 a(R(g)-\chi_1(g))E=0
\]
on this inertia and
\(a\) acts invertibly on \(A_k\). At the two type primes,
\(a=P_v(R(\gamma_v))\), where the coefficients of \(P_v\)
are in the localized Hecke algebra. The tame polynomial
condition asserts that on each fibre this polynomial kills
\(R(g)-\chi_1(g)\) and that \(P_v(\chi_1(\gamma_v))\) is
a unit upon evaluation at \(\pi\). Since all the scalar
characters agree on the chosen inertia, the second property
is exactly invertibility on \(A_k\).

At the remaining
\(S\)-places the family is simultaneously diagonal on inertia
with values \(1\) or \(\chi_1\) on the corresponding
summands. The factors kill the trivial and the \(\chi_1\) summands respectively. Hence on this inertia
\[
 (R(\gamma_v)-1)(R(g)-\chi_1(g))=0.
\]
At these places
\(\zeta_v-1\) is a unit, so the former operator is invertible
on the character target. Consequently \(c\) vanishes on
inertia at every \(S\)-place. These restriction bounds apply under their respective tame hypotheses.
\label{cor:split-selmer-6}

Subtract the row and \(Q\)-local allowance by taking the
kernel of restriction to the corresponding quotients. The
strict conditions for \(k\ge2\) at the coefficient prime
are zero on row and column decomposition groups. For \(k=1\)
the local finite condition for the trivial character at those
places is unramified. The declared rows and columns partition
the places over \(\ell\), and the remaining inertia groups
are among \(S,Q\) or the places with trivial family inertia.
In block \(k=1\), global class field theory therefore gives
\[
 \ker(\text{restriction to the local quotients})\hookrightarrow\operatorname{Hom}(\operatorname{Cl}(E),F/O).
\]
Its normalized length is at
most \(v_\ell(h_E)\).

For \(k\ge2\), let \(S_\ell\) consist of
the places over \(\ell\) and the infinite places. Put
\(A=(\mathbb Q_\ell/\mathbb Z_\ell)(1-k)\). The strict group is
\[
 S_{k-1}=\ker\!\left(H^1(G_{E,S_\ell},A)
        \longrightarrow\prod_{w\mid\ell}H^1(E_w,A)\right).
\]
The restriction kernel maps into \(S_{k-1}\otimes O\).
Indeed the local class vanishes at \(\ell\), all other places
are unramified, and finite freeness over \(\mathbb Z_\ell\) lets
tensoring by \(O\) commute with cochains and localization kernels.
Each \(\mathbb Z_\ell\)-length-one factor becomes \(O/\ell O\),
so the normalized length is \(\operatorname{length}_{\mathbb Z_\ell} S_{k-1}\). The strict
cohomology bound below applies with this exponent and on
this same auxiliary field. \label{cor:split-selmer-7}

\end{proof}
For block \(0\), the same vanishing at the places with trivial
family inertia follows directly from \(y(g)=0\). The row and
column conditions in the restriction kernel lie in the finite
condition on \(W/T\) at all places over \(\ell\). These are
the unramified and finite conditions of the relaxed Selmer
group at \(\pi\) used for \(s_\pi\). The length at this fixed point can therefore be compared with the signature Selmer length, provided the auxiliary character is the one used in its theta lattices.

\subsubsection*{Choosing one twist for the Selmer group and the theta lattices}
\label{subsec:retained-twist}
\hypertarget{source-retained-twist}{}

The same localization supplies a Selmer length and theta denominators. A single character must therefore control every signature and theta value occurring there.

Fix a twisted coefficient place of residue characteristic \(\ell\).
Write \(v(\ell)=1\) in every coefficient field. If \(O\) is a
coefficient DVR, normalized length means \(\operatorname{length}_O/e_O\),
where \(e_O=\operatorname{ord}_O(\ell)\). There are finitely many
good signatures \(s\) for the lower cusp at this place, finitely
many abelian theta values indexed by \(\lambda\), and two sides
for each of the \(m\) theta rows. The sign of the twist on each
of these objects is fixed by its side. A choice for each object
separately would leave the congruence comparison on different
Hecke modules. We describe the simultaneous specialization and
its passage to the retained module.

\begin{established}[Common character data and coefficient tests]\label{hyp:simultaneous-retained-twist}
At each nonreference coefficient place of the initial signed row, both theta comparisons use the same ray coordinates. For each occurring good signature the integral conjugate-duality isomorphism and the isotypic and valuation tests hold. Its jumps and type primes satisfy
\[
 \operatorname{HT}_\tau(T_s)\cap[-1,t]=\varnothing
 \quad(\tau\mid\ell),\qquad
 r\ge100,\quad\min(g_0,g_1)\ge2r .
\] The abelian values have nonzero untwisted specialization and an allowed critical integer. The two Hecke-stable integral rows admit commuting projected theta bases with uniformly bounded denominators and the required residual eigensystems.
\end{established}
\begin{proof}
Fix the row and both signs before selecting a ray character. Lemma~\ref{lem:initial-row-signature-tests} applies to the signatures in the images of the two initial type projectors. It proves the assertions on their lattices, jumps and residual valuations. The simultaneous numerical choices are those of Lemma~\ref{lem:window-parameter-choice} and Proposition~\ref{prop:two-type-primes}. The nonzero abelian specialization is the value in \ref{cor:toric-twist-5}, with the critical range and common denominator of Lemma~\ref{lem:toric-common-denominator}. On either signed row Proposition~\ref{prop:signed-theta-data} gives the commuting action on the integral class sums, their universal denominator and the residual character. Proposition~\ref{prop:retwisting-row-basis} extracts a basis from their projected span before specialization. Both signs are functions of the same class coordinates: they apply their respective coefficient maps to the same finite translations. These are precisely the spaces over which the subsequent character is chosen.
\end{proof}

\begin{auxlemma}[A bound for spherical parameters]
\label{lem:initial-spherical-bound}
Let \(q\) be the residue cardinality of a nonarchimedean local field
and \(\beta_1,\ldots,\beta_r\) the Satake parameters of an unramified
unitarizable representation of \(\mathrm{GL}_r\), in the normalized
Satake convention. For \(r\ge1\),
\[
             \max_i|\beta_i|\le4q^{(r-1)/2}\le q^{(r+3)/2}.
\]
\end{auxlemma}
\begin{proof}
Normalize Haar measure so that the maximal compact subgroup has volume one.
The Hecke operator of \(\operatorname{diag}(\varpi I_j,I_{r-j})\)
is a sum over \(\genfrac{[}{]}{0pt}{}{r}{j}_{\!q}\) residue subspaces of dimension \(j\), each giving
one unitary translate of a spherical vector of norm one. The Gaussian
binomial coefficient counts these subspaces by the number of ordered
independent \(j\)-tuples divided by \(|\mathrm{GL}_j(\mathbb F_q)|\). The normalized Satake formula gives
\[
 |e_j(\beta)|\le q^{-j(r-j)/2}\genfrac{[}{]}{0pt}{}{r}{j}_{\!q}
 \le q^{j(r-j)/2}\prod_{h=1}^j(1-q^{-h})^{-1}.
\]
The last product is at most \(2\) for \(j=1\) and less than \(4\)
for \(j\ge2\), since \(\prod_{h\ge3}(1-2^{-h})\ge3/4\).
Thus \(|e_j(\beta)|^{1/j}\le2q^{(r-1)/2}\).

If a root \(z\) of
the Satake polynomial satisfied \(|z|>4q^{(r-1)/2}\), division by
\(z^r\) would give
\[
      1\le\sum_{j=1}^r|e_j(\beta)/z^j|
                   < \sum_{j=1}^r2^{-j}<1 .
\]
Consequently \(|\beta|\le4q^{(r-1)/2}\le q^{(r+3)/2}\).
\end{proof}

\begin{auxlemma}[The signatures on the initial signed rows]
\label{lem:initial-row-signature-tests}
Fix a coefficient index \(i\) at which the marked character differs
from the reference character. Write \(O_i\) for its coefficient DVR and \(\ell\) for its residue characteristic. Use the initial rows and the two type markings of Proposition~\ref{prop:signed-theta-data}. A candidate signature \(s\) records a lift \(j\), a side \(b\) and
split Hecke eigenvalues \(a_{s,k,P}\) of the lower cuspidal row. Write
\(e_b\) for its assigned finite type projector and \(V_{i,j,b}\) for
the rational space of that lower row. Use precisely the signatures with
\[
 \begin{gathered}
 \exists v_s\in V_{i,j,b}\setminus\{0\},\qquad e_bv_s=v_s,\\
 T_{k,P}v_s=a_{s,k,P}v_s\quad(1\le k\le r),\qquad
 P\text{ split outside the excluded set}.
 \end{gathered}
\]

Let \(\rho_s\) denote the integral Galois representation of this
lower cusp, \(\xi_s\) the first character of its prescribed string,
and \(W=\rho_s\otimes\xi_s^{-1}\) on its chosen coefficient lattice.
Write \(c_E\) for complex conjugation. Then
\[
 (W^\vee(1))^{c_E}\simeq W(t),\qquad
 \operatorname{gr}^a D_{\rm dR}(W)_\tau=0
       \quad(-1\le a\le t,\ \tau\mid\ell).
\]
The isomorphism is integral. Both \(W\) and its Tate dual have
zero decomposition invariants over the coefficient field at every
finite place of the host.
\end{auxlemma}
\begin{proof}
The finite-dimensional projected Hecke spaces are sums of their
simultaneous eigenspaces, by Proposition~\ref{prop:rational-row-isotypic-span}.
Taking the nonzero summands thus gives a finite list on both sides and
at all lifts. Each summand has the cohomological cuspidal attachment
of Proposition~\ref{prop:signed-theta-data}. The two selected types
satisfy the divisor-degree and separation tests of
Definition~\ref{hyp:retained-type-tests}. The inertia argument of
Proposition~\ref{prop:retained-rank-complement} gives absolute
residual irreducibility for \(\bar\rho_s\). These assertions concern
the marked type polynomials on the lower row. On the upper theta
locus Proposition~\ref{prop:cut-selector} identifies the matching
points by
\[
                    Z_i\cap S_i=Z_i^+.
\]

Let \(k_i\) be the residue field of the signature. Absolute
irreducibility and the matrix algebra theorem give the first equality
below after extension to \(\overline{k_i}\), hence also over \(k_i\).
The span is already an algebra because \(\bar\rho_s\) is a
group representation. For \(v\mid g_d\) with \(g_d\ne\ell\), the
depth-zero action is trivial on wild inertia \(P_v\). The tame labels
are separated from the first string character by the tests in
Definition~\ref{hyp:retained-type-tests}. Reading them on a marked
tame inertia generator \(u_v\) gives
\[
 \begin{gathered}
 \operatorname{span}_{k_i}\bar\rho_s(G_E)=M_r(k_i),\qquad
 \bar\rho_s(P_v)=1,\\
 \det\bigl(\bar\rho_s(u_v)-\bar\xi_s(u_v)I_r\bigr)\ne0 .
 \end{gathered}
\]
These are the residual tests on the coefficient ring. Multiplication
by the scalar character \(\bar\xi_s^{-1}\) preserves absolute
irreducibility.

Put \(h=(t-1)/2\). By the weight calculation of
Lemma~\ref{lem:marked-signed-jumps}, the labelled jumps in the
complement have the form \(a=h\pm z\), where the interlacing
inequalities give \(|z|\ge h+2\). Consequently
\[
 a\le h-(h+2)=-2\quad\text{or}\quad
 a\ge h+(h+2)=t+1.
\]
At the conjugate embedding the corresponding jump is \(t-1-a\).
The displayed ranges therefore give exactly one negative jump for
each pair of corresponding entries. With multiplicities this yields
\[
 \dim(D_{\rm dR}(W)_\tau/\operatorname{Fil}^{0})+
 \dim(D_{\rm dR}(W)_{\bar\tau}/\operatorname{Fil}^{0})=r.
\]
The untwisted lower cusp is unramified outside \(\ell,g_0,g_1\).
At \(v\mid\ell\) the two attachments \(\rho_s,W\) are de Rham.
They are crystalline if \(\ell\notin\{g_0,g_1\}\), and potentially
crystalline in the marked type case because their local monodromy
vanishes, as in Proposition~\ref{prop:signed-theta-data}.

At a split unramified Frobenius the polarized attachment, with the
same character normalization on both sides, identifies the
characteristic polynomials of \((W^\vee(1))^{c_E}\) and \(W(t)\).
Chebotarev and Brauer--Nesbitt give an intertwiner over the fraction
field. Multiply it by a scalar to make its entries integral with at
least one unit. For its reduction \(\bar U\), residual
irreducibility gives
\[
 \bar U\ne0\quad\Longrightarrow\quad\ker\bar U=0
 \quad\Longrightarrow\quad \det U\in O_i^\times .
\]
This proves
the required isomorphism on the given lattices.

At \(v\mid\ell\), a nonzero fixed vector would give a jump zero.
For the Tate dual the jumps are \(-1-a\). Thus
\[
        H^0(E_v,W\otimes_{O_i}\operatorname{Frac}(O_i))=
        H^0(E_v,W^\vee(1)\otimes_{O_i}\operatorname{Frac}(O_i))=0 .
\] 
At a type prime of residue characteristic different from
\(\ell\), the separated tame labels give a generator with no
eigenvalue one on the normalized complement. Its inverse transpose
has the same property. At the distinguished auxiliary places the
normalized inertia action includes a nonidentity scalar, with its
inverse on the Tate dual. 
At any remaining place \(v\), put \(q_v=\mathrm Nv\). An arithmetic Frobenius
root is \(\alpha=\beta q_v^{(1-t)/2}\), where \(\beta\) is an
unramified Satake parameter of the unitary lower cusp. Lemma~\ref{lem:initial-spherical-bound} gives
\(|\beta|\le q_v^{(r+3)/2}\). Since \(t\ge r+5\), this gives
\[
 |\alpha|\le q_v^{(r+4-t)/2}\le q_v^{-1/2}<1<q_v .
\]

A fixed vector in the representation or its Tate dual would force
\(\alpha=1\) or \(\alpha=\mathrm Nv\), respectively. This covers
type primes above \(\ell\) by the first case.
\cite{blggt-attachment,caraiani-monodromy}
\end{proof}

\begin{proposition}[One character for the Selmer and theta comparisons]\label{prop:one-retained-character}
For the signatures, abelian values and rows of
Lemma~\ref{hyp:simultaneous-retained-twist}, one finite-order anticyclotomic character
with the prescribed coordinate readings gives all the
following bounds on the original local models, with \(\nu_\ell\) as in
\eqref{eq:signature-cyclotomic-minimum}:
\begin{itemize}
\item for the normalized length \(L_s(\chi)\) at signature \(s\) after applying this character \(\chi\), there are errors \(e_s\) with
\[
 L_s(\chi)\le\nu_\ell+e_s<\infty,\qquad e_s\ge0,\qquad\sum_s e_s\le\ell^{-2};
\]
\item the abelian theta values are nonzero and have valuation
  between \(-N_0\) and \(v(\theta_\lambda(1))+\ell^{-2}\),
  with \(N_0\) the common integral denominator bound;
\item the nonnegative row-index excesses over the untwisted
  indices have sum at most \(\ell^{-2}\);
\item the sum of the valuations of the nontrivial marked
  inertia readings, multiplied by \(t\), is at most \(\ell^{-2}\).
\end{itemize}
The twist may avoid a prescribed finite set and can be
transported at conjugate coefficient places with the
prescribed inversion.
\end{proposition}

\begin{proof}
\proofstep{A common parameter space for the three tests.}
Proposition~\ref{prop:ray-coordinate-pairs} supplies the two
groups of pairs. For the fixed finite family there exist integers \(R_*,D_*,M_0\).
Use \(2R_*\) variables, grouped as \(X_i,Y_i\). At every
\(M\ge M_0\), enlarged to be at least one, there are \(2R_*\)
distinct conjugate pairs of split degree-one places, unramified
over \(\mathbb Q\), of norms \(1\bmod 2\ell^M\). Their union
\(Q\) avoids the level and any prescribed finite set.

There are
continuous anticyclotomic coordinate characters \(d_i\) to
\(\mathbb Z/\ell^M\), each surjective on inertia at both places
of its pair, unramified away from the pair and trivial on all
decomposition groups above \(\ell\). The first half of the
selected Frobenius elements are conjugate to a prescribed tame element
modulo an open normal subgroup. If \(v_j\) denotes the selected representative of the
\(j\)-th second-half pair, the first-half characters satisfy
\[
 \bigl(d_i(\operatorname{Frob}_{v_j})\bigr)_{1\le i,j\le R_*}
       =I_{R_*}\quad\text{over }\mathbb Z/\ell^M\mathbb Z .
\] The
open subgroup may be shrunk and the avoidance set enlarged to
satisfy the finitely many abelian and row conditions.

\begin{enumerate}
\item For each good \(s\) there is an integral polynomial
\(F_{s,M}\) in these variables over a fixed coefficient DVR.
One coefficient of total degree at most \(D_*\) has valuation
at most \(\nu_\ell\). For coordinates \(u_i^{\ell^M}=1\), put
\[
 \psi(g)=\prod_i u_i^{d_i(g)},
\]
over both halves of the
coordinates. The signature Selmer length for \(\psi\) equals
\(v(F_{s,M}(u-1))\), with value infinity at zero. To see
the determinant step in this description, take the square
matrix over \(O[(\mathbb Z/\ell^M)^{2R_*}]\) whose kernel on
discrete \(\ell\)-primary modules computes the Selmer functor
with inverse universal twist. Lift its determinant by sending
each variable to the corresponding group element minus one.

Let \(A(u)\) be the specialized matrix, \(O_u\) the
valuation ring at this specialization and \(K_u\) its fraction
field. Write \(a\) for the size of this square presentation. Its Smith entries
\(c_j\) give the kernel on \((K_u/O_u)^a\), one summand at a time:
\[
 \ker(c_j:K_u/O_u\to K_u/O_u)=
 \begin{cases}
    c_j^{-1}O_u/O_u,&c_j\ne0,\\
    K_u/O_u,&c_j=0 .
 \end{cases}
\]
The normalized length of the first summand is \(v(c_j)\).
Hence, with \(v(0)=+\infty\), specialization yields
\begin{equation}\label{eq:selmer-twist-smith-specialization}
 \operatorname{len}_{\rm norm}\ker A(u)
       =\sum_j v(c_j)=v(\det A(u))=v(F_{s,M}(u-1)).
\end{equation}
For the opposite side, evaluate \eqref{eq:selmer-twist-smith-specialization} at \(u^{-1}\).
\item For each abelian value and each row there is an integral
polynomial \(f_M\) satisfying the following condition, over a
fixed DVR for that object. There are integral polynomials
\(B_M,g_M\), an integer \(b\ge0\) independent of \(M\), and a
bound \(n_*\) independent of \(M\) such that
\[
 f_M\equiv B_M g_M\pmod{(Y_i,(1+X_i)^{\ell^M}-1)_i},
 \qquad \operatorname{ord}_O(g_M(0))\le b.
\]
Modulo the uniformizer and the powers \(X_i^{\ell^M},Y_i^{\ell^M}\),
\(B_M\) is a monomial \(\prod_i X_i^{n_i}\), \(n_i\le n_*\),
times a polynomial with nonzero constant term.

For an abelian
label the value of \(f_M(u-1)\) is \(\ell^{N_0}\theta_\lambda(\psi)\),
where
\[
 b/e_O=N_0+v(\theta_\lambda(1))
\]
and the last value
is taken at empty auxiliary level. For a row, at any specialization
with nonzero value, the relative index of any lattice containing
the specified projected twisted generators, in their prescribed
rational span, is at most the untwisted index plus
\(v(f_M(u-1))-b/e_O\). These two formulas are used
with the sign appropriate to the label or row.
\end{enumerate}

We prove that for any \(\delta>0\) the valuations of all these
polynomials at one signed tuple are bounded by their respective
baselines \(\nu_\ell,b/e_O\) plus \(\delta\). The tuple can
also be required to have
\[
 2t\sum_i v(u_i-1)\le\delta,
\]
where \(i\) runs over both
halves and all orders are nontrivial. We give a limiting argument using discrete valuations.
The theta coefficient DVR has algebraically closed residue field,
so compactness of its coefficients is unavailable.

Choose data at a sequence of levels tending to infinity. There
are finitely many choices of the distinguished monomials of the bounded-degree signature and theta coefficient tests in the three families of polynomials. Fix them on a subsequence. In each fixed
coefficient DVR take the sequences of all coefficients of the
polynomials, including \(B_M,g_M\), and the constant sequence
of a uniformizer. The subring of the ring of sequences that
they generate over \(\mathbb Z\) is countable. Enumerate the countably many coefficient sequences as \(a_1,a_2,\ldots\). A diagonal subsequence \(M_j\) then satisfies
\[
 \lim_{j\to\infty}\operatorname{ord}(a_i(M_j))\in\mathbb N\cup\{\infty\}\qquad(\text{each }i).
\]

This is possible because any sequence of nonnegative orders
has a constant subsequence or one tending to infinity.
For such a sequence ring \(R\), put
\[
       I_\infty=\{a\in R:\operatorname{ord}(a_M)\to+\infty\},
       \qquad v_\infty(\bar a)=\lim_M\operatorname{ord}(a_M).
\]
The valuation inequality and nonnegative orders show that
\(I_\infty\) is an ideal. If \(\bar a\ne0\), its order is
finite (hence eventually constant), and adding an element
of \(I_\infty\) does not change it. For nonzero classes
\[
 v_\infty(\bar a\bar b)=v_\infty(\bar a)+v_\infty(\bar b),
 \qquad v_\infty(\bar a+\bar b)\ge
                   \min(v_\infty(\bar a),v_\infty(\bar b)).
\]
In particular \(R/I_\infty\) is a domain. Extend the order
to its fraction field by subtracting orders. The constant
uniformizer sequence has order one, and the order
of \(\ell\) is unchanged. Use the valuation ring of this field
for the limit series. Every fixed expression in finitely many
coefficients has its limiting order by the definition of the
quotient.

In particular each limit series \(F_s\) has a coefficient
of valuation at most \(\nu_\ell\).

For the second condition, set all \(Y_i=0\). In each fixed
total degree the coefficients of \((1+X_i)^{\ell^M}-1\) tend
to zero, including its zero constant term, by the binomial
formula. Multiplication by arbitrary integral ideal multipliers
preserves this convergence in each bounded degree. The limits
thus satisfy
\[
 f(X,0)=B(X,0)g(X,0).
\]
The selected monomial
coefficient of \(B(X,0)\) is a unit, and
\(\operatorname{ord}_O(g(0))\le b\). For an integral nonzero
formal series let its content valuation be the minimum of the
valuations of its coefficients. This minimum is attained by
discreteness. Write \(c(A)\) for this content valuation. After scaling each nonzero series to have content zero, the residue series are nonzero in a domain. Consequently
\[
 c(AB)=c(A)+c(B).
\]

Consequently the content valuation of \(f\) is at most
\(b/e_O\). This proof permits every \(B_M(0)\) to vanish.
For a negative side first substitute \((1+Z_i)^{-1}-1\) for
each variable \(Z_i\) of the limit and its approximants.
This integral formal automorphism preserves content and is
continuous coefficientwise.

Normalize each resulting limit series by its content so that
its reduction is nonzero. For all of these reductions at once,
choose nonnegative integers \(h_i\) so that substitution
\[
 Z_i=T^{\ell^{h_i}}
\]
leaves a nonzero series. Indeed, for
each series take the lexicographically first monomial of its
support. Starting with the last coordinate, choose each earlier
weight \(\ell^{h_i}\) greater than the total contribution of
the later coordinates in every one of the finitely many chosen
monomials. At the first coordinate where a competing monomial
differs, its increase costs more than can be saved at later
coordinates. Thus the chosen term remains uniquely lowest.
In characteristic zero now substitute
\[
          Z_i=(1+T)^{\ell^{h_i}}-1 .
\]
This reduces to the substitution just considered. Let \(d\)
be the degree of the first unit coefficient in a normalized
one-variable series.

For a primitive \(\ell^a\)-th root
\(\xi\), \(w=v(\xi-1)=1/(\ell^{a-1}(\ell-1))\).
Write \(\sum_k a_k T^k\) for the normalized series.
At \(T=\xi-1\) the difference from the valuation \(dw\)
of its degree-\(d\) term satisfies
\[
 v(a_kT^k)-dw\ge
 \begin{cases}
       1/e_O-(d-k)w,&k<d,\\
       (k-d)w,&k>d.
 \end{cases}
\]
Both bounds are positive when \(dw<1/e_O\). The unique
least term gives valuation \(dw\) for the normalized
series. Take \(a>\max h_i\) so large that
each \(dw<\delta\), all preceding strict inequalities hold,
and \(2t\sum_i\ell^{h_i}w\le\delta\).
Set \(u_i=\xi^{\ell^{h_i}}\), for which
\[
 v(u_i-1)=\ell^{h_i}w.
\]
For zero variables there is simply a
finite list of constants and the empty tuple has zero tame cost.

Fix this \(a\). The preceding single-variable computation uses
only orders of coefficients and the value \(w\). It transfers to
the original polynomials as follows. Write \(m\) for the
normalized content valuation of a limit before dividing by its
uniformizer power. Take \(N\) so large that
\[
 (N+1)w>m+dw
\]
for all series, and \(N\ge d\) for all of
them.

The coefficients of the original one-variable substitutions
through degree \(N\) are fixed integral polynomial expressions
in finitely many input coefficients, including for inverse
coordinates. Their orders therefore converge to the orders in
the limit. Write \(b_j\) for the coefficients at one such level and take
\(T=\xi-1\), so \(v(T)=w\).

For all sufficiently large levels on
the subsequence, finite coefficient convergence and integrality
of the tail give
\[
 \begin{aligned}
 v(b_dT^d)&=m+dw,\\
 \min_{\substack{0\le j\le N\\j\ne d}}v(b_jT^j)&>m+dw,\\
 \inf_{j>N}v(b_jT^j)&\ge (N+1)w>m+dw .
 \end{aligned}
\]
An empty minimum has value \(+\infty\). The ultrametric inequality is
therefore an equality on the full series:
\[
 v(\sum_{j\ge0}b_jT^j)=m+dw.
\]
Inverse-coordinate series may be evaluated in the
completion of the coefficient field, where the substitution
agrees with inversion at the root of unity. The test works
for every primitive root of that order because it uses only
the valuation \(w\) and a unique minimal term. Thus the same
level gives the bounds for every power of the character prime
to \(\ell\).

At such a level \(M\ge a\) use its selected primes and its
characters \(d_i\). The ray class character at the chosen level is read on the same coordinates by
\[
 \psi(g)=\prod_i u_i^{d_i(g)}.
\]
Each coordinate is abelian, unramified off its selected pair and tamely ramified there, so the Artin map gives the desired ray character. The universal theta coordinates at the chosen primes agree with the characters \(d_i\) by the restriction property for the chosen coordinate characters. Ramification is exact because each \(u_i\ne1\)
and no other coordinate ramifies at that pair. The image on either inertia group at pair \(i\) is
\(\langle u_i\rangle\), with \(u_i\) or \(u_i^{-1}\) realized
on a generator.

For any integer \(n\),
\(u_i^n-1\in(u_i-1)O'\) in the coefficient valuation ring
\(O'\); for negative \(n\), use
\[
 u_i^{-n}-1=-u_i^{-n}(u_i^n-1).
\]
Thus the two places in
each pair have the same minimum, and the weighted cost is
\[
  t\sum_i 2\min_{\psi(g)\ne1}v(\psi(g)-1)
      =2t\sum_i v(u_i-1)
      =2t\sum_i\ell^{h_i}w \le \delta ,
\]
where the minimum at each \(i\) is restricted to its inertia
group.

\proofstep{Choose one specialization that fits all four allowances.}
All evaluations now use a common auxiliary parameter. Distributing the permitted error among the finitely many signatures and abelian values allows one specialization to satisfy all of the length and index bounds.

Take \(\delta\) smaller than \(\ell^{-2}/(1+2m+\#\{s\})\).
Writing \(L_s\) for each specialized signature length, the common choice gives
\[
 L_s\le\nu_\ell+e_s,\qquad e_s\ge0,\qquad \sum_s e_s\le\ell^{-2}.
\]
In particular every \(L_s\) is finite. The abelian
theta values are nonzero with valuations between \(-N_0\) and
\(v(\theta_\lambda(1))+\ell^{-2}\), the lower bound following
from integrality of \(f_M\). The nonnegative excesses over the
untwisted row indices sum to at most \(\ell^{-2}\).

The tame
cost multiplied by \(t\) is at most \(\ell^{-2}\) as well.
The number of primes, the depth, the coefficient extension and
the coordinate orders may depend on the host and the given
finite family. Only the four displayed allowances are used at
subsequent places. At conjugate coefficient places transport
the character with the prescribed inversion. For finitely many
other choices add all rational primes already used to the
avoidance set, taking all places over them. This gives the
disjoint auxiliary levels of the norm data. \label{cor:retained-twist-1}
\end{proof}

\subsubsection*{The relaxed global length at the chosen point}

Enlarging the coefficient field can realize the chosen character and its values. The Selmer cost in the determinant must still be measured in the original units of normalized length.

\begin{auxlemma}[Normalized torsion length under a DVR extension]
\label{lem:retained-dvr-normalized-length}
Let \(A\to C\) be a local inclusion of discrete valuation rings of mixed
characteristic \((0,\ell)\), let \(v_A,v_C\) give value one to the respective
uniformizers, and let \(D\) be a torsion \(A\)-module. Then
\[
 \frac{\operatorname{len}_C(C\otimes_A D)}{v_C(\ell)}
   =\frac{\operatorname{len}_A D}{v_A(\ell)}
\]
with infinite length allowed.
\end{auxlemma}

\begin{proof}
For a uniformizer \(\pi_A\), put \(e=v_C(\pi_A)\). The extension is flat
because \(C\) is torsion-free over the DVR. Tensoring a composition series
for a finite-length module replaces each simple factor by
\(C/\pi_A C\), of length \(e\). Also \(v_C(\ell)=e v_A(\ell)\).
An infinite-length torsion module has finite-length submodules of arbitrarily
large length. Flatness preserves their inclusions, proving the assertion
in that case.
\end{proof}

\begin{lemma}[Specialization on the actual retained lattice]\label{lem:retained-signature-specialization}
At a good retained point of the marked localization, the
selected lower complement on the prescribed lattice identifies
after coefficient extension with its good signature. The
relaxed Selmer group injects into that of the specialized
signature in the normalized comparison. In particular its
finite length satisfies
\[
 s_\pi\le\nu_\ell+\ell^{-2},\qquad
 t\sum_{v\in Q}v(\zeta_v-1)\le \ell^{-2}.
\]
\end{lemma}

\begin{proof}
\proofstep{Descend the rational intertwiner and use its nonzero integral reduction.}
At a good retained point \(\pi\), take the signature \(s\)
with the same lower cusp, tuple and side. Cancelling the
prescribed character factors in the row parameter
identifies the traces at split unramified Frobenius elements.
Chebotarev and continuity give
\[
 \operatorname{tr}\rho_\pi(g)=\operatorname{tr}\rho_s(g)
\]
for every \(g\), after scalar extension. Absolute irreducibility
of the reduction implies absolute irreducibility in characteristic
zero: a proper invariant subspace would have a proper saturated
intersection with a stable lattice after finite coefficient
extension, and reduction would give a proper nonzero invariant
subspace of the residual lattice. Brauer--Nesbitt gives
a fraction-field intertwiner, which descends because the
intertwiner equations are linear. Multiply it by a scalar so
that its entries are integral with at least one unit.

If \(A\) is the scaled intertwiner, its nonzero reduction has proper stable kernel. Residual irreducibility and equality of ranks give
\[
 \ker\bar A=0\quad\Longrightarrow\quad\det\bar A\ne0\quad\Longrightarrow\quad\det A\in O^\times.
\]
Thus \(A\) is an integral isomorphism. Twisting by
\(\chi_1^{-1}\) identifies \(W=\rho_\pi\chi_1^{-1}\)
with the signature lattice twisted by the inverse auxiliary
character of the chosen side. The signature lattice already
contains the inverse untwisted rank-one factor.

The two character evaluations identify their generated
orders by equality of kernels on the base group algebra.
Take the integral closure in their common fraction field.
Let \(A\to C\) denote either of the two local torsion-free
DVR extensions from this common base and \(\pi_A\) a
uniformizer. Fix lifts \((c_i)_i\) of a basis of \(C/\pi_A C\)
over \(A/\pi_A A\). They give, compatibly for every \(n\ge1\),
\[
 (A/\pi_A^n A)^{(i)}
     \xrightarrow{\ (x_i)\mapsto\sum_i c_i x_i\ }
                         C/\pi_A^n C ,
\]
where the superscript denotes the direct sum on the index set.
For surjectivity lift residues successively through order \(n-1\),
using finitely many terms at each order.

For injectivity, suppose \(\sum c_i x_i\in\pi_A^n C\) and some
\(x_i\notin\pi_A^n A\). For \(m=\min_i v_A(x_i)<n\), torsion-freeness
permits division by \(\pi_A^m\) in \(C\). Reduction gives
\[
 \sum_i\overline c_i\,\overline{x_i/\pi_A^m}=0\quad\text{in }C/\pi_A C.
\]
At least one coefficient is nonzero in \(A/\pi_A A\), contrary to the
choice of the residue basis. The same lifts work on all bounded torsion
submodules. Also \(C[1/\pi_A]=K_C\) since \(\pi_A\) has positive
\(C\)-valuation, so flatness gives
\[
 C\otimes_A(K_A/A)=K_C/C.
\]


If \(D\) is a discrete torsion coefficient module with this
action, continuous cochains commute with the indicated direct
sum. Each cochain has finite image on a finite power of the
compact Galois group, and consequently uses finitely many
copies. Differentials and restrictions act coordinatewise.
The auxiliary character is trivial on each decomposition group
over \(\ell\). The signature finite condition is the scalar span
of the base images, which is the direct sum of the base
conditions under this decomposition. For the finite local
coefficient fields the crystalline condition commutes with
extension by the finite basis decomposition of the rational
period comparison, and taking the image in the torsion quotient
commutes with it as well. Conjugate embeddings over the same
place transport this condition. Writing \(S_A\) for the Selmer module over the common base and \(S_C\) for either extension, compatibility of all local kernels gives
\[
 C\otimes_A S_A\simeq S_C.
\]

Lemma~\ref{lem:retained-dvr-normalized-length}, applied to the common
base Selmer group, identifies both normalized lengths, including the
criterion for finiteness.

Finally identify \(K^r/O^r\) with \((K/O)^r\) coordinatewise.
This carries the global Selmer group defining \(s_\pi\) into
the signature group defining the specialized determinant.
Unramified restrictions and unrestricted conditions at
\(S\cup Q\) away from \(\ell\) are preserved. At
\(v\mid\ell\) a generator of the local finite condition
comes from the reduction of a rational crystalline finite
cocycle. Write \(c\) for the chosen torsion cocycle and \(z\) for a rational finite cocycle. If \(c-\bar z=d a\), lift \(a\) to \(\tilde a\in K^r\); then
\[
 \overline{z+d\tilde a}=\bar z+d a=c .
\] Its coordinates then give the lifts in
the definition of the signature local condition, and taking
the integral span preserves membership. Coordinatewise
identification induces an isomorphism on ambient cohomology,
so the map on these Selmer subgroups is injective. In
particular finiteness of the specialized signature length
gives finiteness at \(\pi\) and
\[
          s_\pi\le \nu_\ell+\ell^{-2},
       \qquad t\sum_{v\in Q} v(\zeta_v-1)\le\ell^{-2}.
\]

\proofstep{Read the norm and signature characters on the same generator of inertia.}
Passing to a larger coefficient field preserves normalized length. We must further check that the character used to define the relaxed condition there is the one selected by the local norm data.

For the second inequality, the generator selected in the
marked model has full inertia image. The character of the
norm data and the twisted signature have the same character
values up to coefficient embedding and inversion, both of
which preserve the valuation of the difference from one.
These are the two global allowances in the local congruence
table. The pair-removal calculation derives the two theta polynomial assertions from their integral level and Euler relations. \label{cor:retained-twist-2}
\end{proof}

\subsubsection*{The square Selmer presentation}

The character choice used a polynomial coming from the Selmer group. Its construction begins with a square presentation over the group ring, so specialization gives a determinant on the prescribed auxiliary coordinates.

\begin{established}[Coefficient algebra and Selmer functor]\label{hyp:signature-presentation}
Put \(\Lambda=O_0[(\mathbb Z/\ell^M)^{2R_*}]\), with \(O_0\) the coefficient DVR of the signature, and let \(k\) denote the residue field. This algebra is finite free, local and complete. For the signature Selmer functor \(S\), with inverse universal character and relaxed conditions at the selected pairs, \(S\) is additive, \(\Lambda\)-linear and left exact. Every class is defined over a finitely generated torsion submodule, and \(S(k)\) is finite. For \(0\to k\to B\to C\to0\) one has
\[
 \#\operatorname{coker}(S(B)\longrightarrow S(C))\le\# S(k).
\]
\end{established}
\begin{proof}
The group elements form an \(O_0\)-basis. Indeed every maximal ideal lies
over the maximal ideal of \(O_0\), and after reduction one has
\[
 k[(\mathbb Z/\ell^M)^{2R_*}]\simeq k[X_1,\ldots,X_{2R_*}]/(X_1^{\ell^M},\ldots,X_{2R_*}^{\ell^M}).
\]
Here \(X_i\) is the \(i\)-th generator minus one, so the augmentation ideal is nilpotent.
Thus the maximal-ideal and base uniformizer topologies agree.
Its residue field \(k\) is finite.

 Additivity and linearity follow from the cocycle maps. Lemma~\ref{lem:signature-left-exact} proves left exactness and finite generation for individual classes. Finiteness over \(k\) follows from global cohomological finiteness for the finite module \(T(k)\). In \eqref{eq:sig-obstruction} a global connecting class and a tuple of local corrections give obstruction groups of orders \(|A|\) and \(|C|\), with product \(|S(k)|\). The fibre count following that identity gives the required bound, before the construction of the square presentation.
\end{proof}

\begin{lemma}[Square presentation of the signature functor]
\label{lem:signature-square-presentation}
For the ring and functor of Lemma~\ref{hyp:signature-presentation}, put
\(a=\dim_k S(k)\). There is one matrix \(D\) over the maximal ideal such that for every discrete torsion coefficient module \(B\) one has, naturally in \(B\),
\[
 D\in M_a(\mathfrak m),\qquad S(B)\simeq\ker(D:B^a\longrightarrow B^a).
\]
\end{lemma}

\begin{proof}
On finite modules use the exact duality
\(D(-)=\operatorname{Hom}_{\mathbb Z_\ell}(-,
\mathbb Q_\ell/\mathbb Z_\ell)\). The modules \(S(B)\) are
finite for finite \(B\) by a composition series and left
exactness. Hence \(J(B)=D S(D B)\) is a right exact linear
functor on finite modules.

For \(A_n=\Lambda/\mathfrak m^n\)
put \(P_n=J(A_n)\). If \(B\) is an \(A_n\)-module of finite
length then \(J(B)\simeq B\otimes_{A_n}P_n\) naturally.
For \(b\in B\) let \(\lambda_b:A_n\to B\) send \(1\) to
\(b\). The map \(b\otimes p\mapsto J(\lambda_b)(p)\)
is balanced and is the identity on \(A_n\). Additivity makes
it an isomorphism on finite free modules.

For a presentation \(A_n^q\xrightarrow{U}A_n^d\to B\to0\),
applying \(J\) gives
\[
       J(B)\cong \operatorname{coker}
                    (P_n^q\xrightarrow{U}P_n^d)
                \cong B\otimes_{A_n}P_n .
\]
The maps are induced by \(b\otimes p\mapsto J(\lambda_b)p\),
so the construction is natural in \(B\). In particular, view \(A_n\)
as an \(A_m\)-module when \(m\ge n\). This gives
\[
      P_n=J(A_n)\simeq A_n\otimes_{A_m}P_m
                       =P_m/\mathfrak m^nP_m .
\]
Write \(a=\dim_k S(k)=\dim_k P_1\), using \(Dk\simeq k\).
Lift a basis of \(P_1\) compatibly in the surjections \(P_m\to P_n\).

By Nakayama this gives surjections \(\Lambda^a\to P_n\) with
descending kernels \(K_n\). For \(m\ge n\),
\(K_n=K_m+\mathfrak m^n\Lambda^a\).
These kernels are compact in the finite base coefficient
topology.

For each \(x\in K_n\) the compact sets
\(K_m\cap(x+\mathfrak m^n\Lambda^a)\) are nested and
nonempty for \(m\ge n\). Consequently for \(K=\bigcap K_n\)
one has
\[
 K_n=K+\mathfrak m^n\Lambda^a.
\]
It follows that
\(P=\Lambda^a/K\) represents the original functor:
on a finite module killed by \(\mathfrak m^n\), dualizing the
tensor description gives
\[
 S(B)=\operatorname{Hom}_\Lambda(P,B)
\]
naturally. Every finite torsion module is killed by such a
power. The finite-generation condition on classes and left
exactness extend this equality to arbitrary discrete torsion
modules, since \(P\) is finitely generated and every finitely
generated torsion submodule is finite.

For the relations set
\(I=\operatorname{Hom}_{\mathbb Z_\ell}(\Lambda,
\mathbb Q_\ell/\mathbb Z_\ell)\). It is injective over
\(\Lambda\) by adjunction and divisibility over \(\mathbb Z_\ell\).
Finiteness of \(\Lambda\) over \(\mathbb Z_\ell\) makes \(I\)
torsion, and dualizing \(\Lambda\to k\) embeds \(Dk\simeq k\)
in it. Apply \(\operatorname{Hom}_\Lambda(P,-)\) to
\(0\to k\to I\to I/k\to0\). The connecting map is surjective
because \(I\) is injective, so the arithmetic cokernel is
\[
 \operatorname{coker}\bigl(S(I)\longrightarrow S(I/k)\bigr)
          =\operatorname{Ext}^1_\Lambda(P,k).
\]
Applying \(\operatorname{Hom}_\Lambda(-,k)\) to
\(0\to K\to\Lambda^a\to P\to0\) identifies this group
with \(\operatorname{Hom}_k(K/\mathfrak mK,k)\) because restriction
from \(\Lambda^a\) is zero when
\(K\subseteq\mathfrak m\Lambda^a\).
Its size bound implies \(\dim_k K/\mathfrak m K\le a\).

The ring is Noetherian, so at most \(a\) generators suffice
for \(K\). Use them as relations and pad by zero relations
to get an \(a\)-square matrix with entries in \(\mathfrak m\).
The kernel of its action on \(B^a\), taking the relations
as rows, is \(S(B)\) naturally in \(B\). The square matrix now represents the Selmer functor. Detecting its determinant in bounded degree is an arithmetic step: the auxiliary Frobenius elements must be realized by primes with prescribed ray coordinates. \label{cor:retained-twist-3}
\end{proof}


\subsubsection*{Ray characters and the two Frobenius prescriptions}
\label{subsec:aux-coordinates}
\hypertarget{source-auxiliary-coordinates}{}

The variables of the group-ring determinant must be realized by ray characters. A new prime must satisfy two Frobenius prescriptions, and both prescriptions have to survive the restrictions imposed at the distinguished twist prime.

Fix the auxiliary CM field \(E\), its conjugation \(c\) (with the chosen
extension of order two), the coefficient prime \(\ell\), and
\(L=\ell^M\), \(M\ge1\). Write \(T\) for the places over the
distinguished twist prime. Its residue characteristic is prime
to \(2L\). We first construct an open normal subgroup \(H\)
of \(G_E\), containing inertia above \(T\), and a finite set
\(S\) such that each split place \(P\notin S\) whose
decomposition group lies in \(H\) admits the coordinate
character of Section~\ref{subsec:retained-twist}.

\begin{proposition}[Prescribed Frobenius elements for the ray coordinates]\label{prop:ray-coordinate-pairs}
For the prescribed finite list of Frobenius conjugacy
conditions, the two groups of split auxiliary pairs can
be chosen outside any enlarged finite avoidance set.
Their anticyclotomic coordinates are surjective on the local
inertia of their own pair, unramified off that pair, and
trivial at decomposition groups above \(\ell\). Writing \(\eta_i\) for the additive first-group coordinates and \(P_j\) for the marked primes of the second group, one has
\[
 \eta_i(\operatorname{Frob}_{P_j})=\delta_{ij}\quad\text{in }\mathbb Z/L.
\]
\end{proposition}

\begin{proof}
\proofstep{Kill the global unit obstruction before prescribing a ray coordinate.}
Choose ideal representatives for the finite class group of \(E\)
prime to \(T\). This is possible by weak approximation at \(T\)
applied to the valuations of each representative. Put their prime
supports and the primes above \(\ell\) in \(S\), disjoint from \(T\).
Every fractional ideal has a generator away from \(S\). Choose
generators \(u_1,\ldots,u_a\) for the \(S\)-units, including
generators for their torsion. Take for \(H\) the subgroup fixing
\(E(\mu_L,u_1^{1/L},\ldots,u_a^{1/L})\). This is a finite
Galois extension. At \(v\in T\) every radicand is a unit
and \(L\) is invertible. The roots of each defining polynomial
lift its simple roots over the maximal unramified extension, so
inertia at \(v\) lies in \(H\).

For \(P\) as above, \(L\mid (\mathrm NP-1)\) and every \(u_i\)
reduces to an \(L\)-th power. Choose a surjection
\[
 \omega:k(P)^\times\to\mathbb Z/L
\]
from the cyclic residue
group. For an ideal \(J\) prime to \(P\), take a generator
\(\alpha\) away from \(S\). Its order at \(P\) is zero, so
\[
 J\mapsto\omega(\bar\alpha)
\]
is defined. Two generators
differ by an \(S\)-unit, and multiplication of generators proves
additivity. A principal ideal generated by \(1\bmod P\) has
value zero. Ray class reciprocity now gives a continuous
character \(\xi_P:G_E\to\mathbb Z/L\), unramified away from
\(P\). Its image on inertia at \(P\) is the image of \(k(P)^\times\):
on principal ideals prime to \(P\) the character evaluates the
residue character, and every residue unit lifts to an element
prime to \(P\).

At every \(v\mid\ell\) take \(1\) for a
generator of \(v\) away from \(S\). Both Frobenius and inertia
therefore act trivially there. Set
\[
       d_P(g)=\xi_P(g)-\xi_P(cgc^{-1}).
\]
Since \(P\ne cP\), only one summand ramifies at each member of
the pair. For the resulting difference character \(d_P\), the two inertia images are
\[
 d_P(I_P)=d_P(I_{cP})=\mathbb Z/L,
\]
including at \(\ell=2\).
This gives the desired anticyclotomic character, trivial on
decomposition groups above \(\ell\). \cite{milne-cft}

The universal coordinates in the theta data can be fixed once
per unordered pair. For each pair take a character with these
four properties (anti-invariance, triviality above \(\ell\),
ramification confined to the pair and surjectivity on its two
inertia groups) if one exists, and the trivial character otherwise.
Use this choice for the coordinate at \(P\) in every
conjugation-stable set containing it. Thus the ray construction
certifies existence at every admitted pair and the actual universal
coordinate there has the required properties. No uniqueness of a
ray character is required.
Restrict the chosen character to the local level of each selected series.

Here is the form of prime selection needed for the Galois
determinant. Let \(U\) be any open normal subgroup and let \(z\)
be the prescribed tame element at \(T\). Let
\[
       W=U\cap H\cap G_{E(\mu_{2L})}.
\]
We can prescribe the Frobenius of each first-half prime \(P_b\)
to be conjugate to \(z\) modulo \(W\), and the Frobenius of
each second-half prime \(P'_a\) to be conjugate modulo \(U\)
to any specified \(h_a\in H\cap G_{E(\mu_{2L})}\). In addition
we can require
\begin{equation}\label{eq:aux-coordinate}
       d_{P_b}(\mathrm{Frob}_{P'_a})=\mathbf1_{a=b}
                         \quad\text{in }\mathbb Z/L .
\end{equation}

To prove this, apply Chebotarev over \(E\) in the finite
quotient by \(W\), avoiding any given finite set, the ramification
loci, their conjugates and the primes above \(2L\). Degree-one
primes suffice: in the prime count over a number field primes of
higher absolute residue degree up to norm \(x\) contribute
\(O_E(\sqrt{x}\log x)\). A degree-one prime unramified over
\(\mathbb Q\) splits over its restriction to the maximal real
subfield because an inert quadratic place would have even residue
degree. The element \(z\) lies in \(H\) and fixes \(\mu_{2L}\)
by the inertia condition on \(T\). The primes selected have
\(\mathrm NP_b\equiv1\bmod 2L\) and are admitted. Their
inertia groups lie in \(W\). After each choice exclude the
whole pair.

\proofstep{Multiply each later Frobenius by inertia before applying Chebotarev.}
The new Frobenius must realize the prescribed values on the earlier ray coordinates as well.

The map
\[
 W\to(\mathbb Z/L)^R
\]
given by the \(R\) first-half
characters is surjective. Inertia at \(P_b\) realizes every
multiple of the \(b\)-th basis vector, and the target is abelian.
Choose \(e_a\in W\) so that the coordinates of \(h_a e_a\)
are the \(a\)-th basis vector. Apply Chebotarev again, now
modulo
\[
 W\cap\bigcap_b\ker d_{P_b},
\]
to \(h_a e_a\).
These primes are admitted by the same norm and splitting tests,
satisfy the prescribed conjugacies modulo \(U\), and satisfy
\eqref{eq:aux-coordinate} for every Frobenius lift. All exclusions
can be enlarged before either choice. This proves the prime and
character part of the construction at every depth. Bounds on
the norms of the auxiliary primes do not enter the length
estimate. The numbers of pairs and the degree tests uniform in
\(M\) come from the residual Selmer argument below.
\end{proof}


\subsubsection*{The Selmer group on a signature lattice}
\label{subsec:signature-selmer}

With ray coordinates available, the remaining inputs to the square presentation concern the Selmer functor. We verify them on the lattice for one good signature, retaining its fixed twist and its local conditions.

\begin{established}[Signature coefficients and local conditions]\label{hyp:signature-module-local}
Let \(T\) denote the rank-\(r\) signature lattice including its
prescribed scalar twist, \(O_0\) its coefficient DVR, and \(V\)
its fraction-field representation. Write \(k\) for the residue
field of the coordinate group ring \(\Lambda=O_0[(\mathbb Z/\ell^M)^{2R_*}]\) and \(\kappa\) for
the inverse universal character. Put
\[
 T(B)=(T\otimes_{O_0}B)\otimes\kappa
\]
for an
\(\mathfrak m\)-power torsion module \(B\). Let
\(\Sigma\) consist of the places above \(\ell\), the specified
type and distinguished places away from \(\ell\), and the
selected pairs \(Q\). Define \(G_\Sigma=\operatorname{Gal}(E_\Sigma/E)\), where \(E_\Sigma\) is the maximal extension unramified outside \(\Sigma\). The global signature classes unramified away from \(\Sigma\) are exactly the classes on \(G_\Sigma\). At the finite non-\(\ell\) places in \(\Sigma\) the conditions are relaxed, so the remaining local tests are at \(v\mid\ell\).
\end{established}
\begin{proof}
The underlying representation is unramified off its type places, the distinguished places and the places above \(\ell\), by Proposition~\ref{prop:signed-theta-data}. The universal character is unramified outside the chosen pairs and the prescribed ray places. Thus inertia outside \(\Sigma\) acts trivially on \(T(B)\). A cocycle representing an unramified class restricts to zero on each such inertia group because every coboundary there is zero. It vanishes on their conjugates and hence by continuity on the closed normal subgroup they generate. It therefore factors through \(G_\Sigma\); inflation gives the inverse, including for discrete torsion \(B\).
\end{proof}

\begin{lemma}[Variation in the signature coefficient module]\label{lem:signature-left-exact}
For \(T(B)\) and \(S(B)\) in Lemma~\ref{hyp:signature-module-local},
the finite local condition at \(v\mid\ell\) lifts under
surjections of \(\mathfrak m\)-power torsion coefficient
modules \(B\) and is reflected under injections. Every short exact coefficient sequence \(0\to B_1\to B_2\to B_3\to0\) induces an exact sequence
\[
 0\longrightarrow S(B_1)\longrightarrow S(B_2)\longrightarrow S(B_3).
\] Each of its classes comes from a
finitely generated coefficient submodule.
\end{lemma}

\begin{proof}
\proofstep{Describe the finite local condition as a span of integral cocycles.}
For \(v\mid\ell\) we have \(\kappa|_{G_v}=1\). Let
\(Z_v\) be the module of \(T\)-valued continuous cocycles on
\(G_v\) whose rational classes belong to \(H^1_f(E_v,V)\).
The local condition \(L_v(B)\) is the span in cohomology of
\[
 g\mapsto z(g)b,
\]
for \(z\in Z_v\) and \(b\in B\),
using a basis of \(T\) to read this formula coordinatewise.
All coboundaries lie in the corresponding span of cocycles:
use the integral coboundaries of the coordinate basis vectors
of \(T\). Thus a class belongs to \(L_v(B)\) precisely when
any representing cocycle is such a finite sum. Maps of
coefficient modules preserve this description. For a coefficient surjection \(q:B_1\twoheadrightarrow B\), take any lift \(b_1\) of \(b\). Termwise lifting follows from
\[
 q(z(g)b_1)=z(g)q(b_1)=z(g)b.
\]


\proofstep{Use saturation to read each local relation before enlarging coefficients.}
The finite local condition was described after passing to the rational representation. Saturation identifies the integral relations that must be preserved on extending the coefficient module.

Injections of coefficient modules reflect this condition.
Here the integral assertion is at the level of functions on
\(G_v\times\{1,\ldots,r\}\). The module \(Z_v\) is saturated
in all integral functions: if \(az\in Z_v\) with \(a\ne0\),
division in the rational vector space preserves the cocycle
law, continuity and the finite local condition. Suppose a
\(B'\)-valued function has image \(\sum_{j=1}^n z_j b_j\)
under an inclusion \(B'\hookrightarrow B\). Discard zero
terms.

Take an entry \(a=z_1(i_0)\) of least valuation and put
\[
 u=z_1/a\in Z_v,
\]
with \(u(i_0)=1\). For the original function \(f\) with values in \(B'\), let
\(\iota:B'\hookrightarrow B\) be the coefficient inclusion.
With \(f(i_0)\in B'\), the subtraction gives
\[
 \iota\!\left(f-u f(i_0)\right)
  =\sum_{j>1}\bigl(z_j-u z_j(i_0)\bigr)b_j .
\]
Each new cocycle belongs to \(Z_v\), and the representation
has one fewer term. Induction puts \(f-u f(i_0)\) in the
required span over \(B'\), which also contains \(u f(i_0)\). This argument works when
\(B'\subset B\) has torsion quotient.

\proofstep{Kill global invariants by reduction to the coefficient socle.}
The absolute residual complement is irreducible of dimension
\(r>1\), by the retained type tests. On a simple coefficient
module the residual scalar twist preserves irreducibility and
\(\kappa\) reduces to one. Hence
\[
 H^0(E,T(k))=0,
\]
also for
the Tate dual. For arbitrary \(B\), a nonzero invariant vector
would generate a cyclic module killed by a power of
\(\mathfrak m\). Its last nonzero maximal-ideal power would
give a nonzero invariant in the socle, where the coefficient
space is a sum of copies of \(k\). This proves the same
vanishing on all torsion modules.

The long exact sequence of
global cohomology now gives injectivity on \(H^1\) for
inclusions and exactness at the middle for a short exact
coefficient sequence. The local reflection just proved gives
left exactness of \(S\). A continuous cocycle on a compact
group with discrete target has finite image, so its coordinates
generate a finitely generated torsion submodule. Local reflection
puts the class into the Selmer group of that submodule. This
proves the finite-generation assertion used in constructing
the representing module.
\end{proof}

\paragraph{Cardinalities of the finite local conditions.}
The good-signature isotypic condition gives \(H^0(E_v,V)=0\)
at \(v\mid\ell\), and \(V\) is de Rham there. Set
\(d_v=\dim H^1_f(E_v,V)\), the dimension over the coefficient
field. The Bloch--Kato dimension formula gives
\[
 d_v=\sum_{\tau\mid v}\dim\big(D_{\rm dR}(V)_\tau/
                                      \operatorname{Fil}^0\big).
\]
Here the coefficients split the embeddings, and in the formula
each summand is the dimension of the indicated quotient.
The signature filtration at conjugate embeddings has complementary
quotient dimensions, with sum \(r\). As \([E:\mathbb Q]=2f\),
let \(h_\tau\) be the quotient dimension at \(\tau\). Summation gives
\[
 \sum_{v\mid\ell}d_v
   =\sum_{\{\tau,\bar\tau\}}(h_\tau+h_{\bar\tau})
   =\frac{[E:\mathbb Q]}2\,r=fr .
\]
\cite{bloch-kato}

The fraction-field space of finite cocycles has dimension
\(r+d_v\), since coboundaries have dimension \(r\). Its
intersection \(Z_v\) with integral functions is a full free
lattice. Compactness bounds each of a finite basis of rational
cocycles and hence clears its denominators. For finite generation,
choose finitely many evaluations whose coordinate functionals
span the dual of the rational space, giving an injection of
the integral functions into a finite free module. Reduction of
\(Z_v\) modulo its uniformizer injects into the functions
over the residue field by saturation. Consequently the span
of these cocycles with coefficients in \(k\) has cardinality
\(|k|^{r+d_v}\). Dividing by the residual coboundaries, of
size \(|k|^r/|H^0(E_v,T(k))|\), gives
\begin{equation}\label{eq:signature-local-size}
  \prod_{v\mid\ell}|L_v(k)|
   =|k|^{fr}\prod_{v\mid\ell}|H^0(E_v,T(k))|.
\end{equation}
This calculation allows the reduction to have local invariants.

\paragraph{Counting the obstruction to lifting a Selmer class.}
Put \(D=T(k)\), \(m=|D|^f\), and evaluate cohomology on
\(G_\Sigma\). Define
\[
 \begin{aligned}
 A&=\ker\Big(H^2(E_\Sigma/E,D)\longrightarrow
                   \prod_{v\mid\ell} H^2(E_v,D)\Big),\\
 C&=\operatorname{coker}\Big(H^1(E_\Sigma/E,D)
       \longrightarrow\prod_{v\mid\ell} H^1(E_v,D)/L_v(k)\Big).
 \end{aligned}
\]
These groups are finite by local and global finiteness for a
finite module, a finite ramification set containing all primes
over \(\ell\), and a totally imaginary field. The last
localization in degree two is surjective by Poitou--Tate:
the cokernel of localization at all places of \(\Sigma\)
is dual to the invariants of the Tate dual, which vanish.
Projection to the places over \(\ell\) preserves surjectivity.
The global Euler characteristic gives \(|H^1|=m|H^2|\)
because there are \(f\) complex places and \(H^0=0\).
The product of the local Euler characteristics at \(\ell\)
gives
\[
 \prod_{v\mid\ell}|H^1(E_v,D)|=
 m^2\prod_{v\mid\ell}|H^0(E_v,D)|\,|H^2(E_v,D)| ,
\]
where \(\sum_v[E_v:\mathbb Q_\ell]=2f\). Taking kernels
and cokernels of the two displayed localization maps and using
\eqref{eq:signature-local-size} yields \cite{milne-duality}
\begin{equation}\label{eq:sig-obstruction}
                        |S(k)|=|A|\,|C|.
\end{equation}

For
\[
 0\to k\to B\to B''\to0,
\]
the connecting map on
\(S(B'')\) has image in \(A\), since every allowed local class
lifts to an allowed local class. On its kernel choose a global
lift to \(H^1(G_\Sigma,T(B))\). Subtract an allowed local
lift at each \(v\mid\ell\): the difference comes from
\(H^1(E_v,D)\). Its tuple of preimages, modulo the local
conditions and global image, records the remaining obstruction.
For the cardinality bound one can allow choices: for each
class in the kernel choose one global lift, local lifts and
preimages.

If two of the resulting classes in \(C\) agree,
subtract a global \(D\)-class and use the local \(L_v(k)\)
to put the difference of the global lifts in \(S(B)\).
There are at most \(|C|\) cosets in the kernel modulo the Selmer image and at most \(|A|\) fibres. With \eqref{eq:sig-obstruction} this gives
\[
 \#\operatorname{coker}\bigl(S(B)\to S(B'')\bigr)\le |A|\,|C|\le |S(k)|.
\]
Together with left exactness and passage
to finitely generated submodules this supplies the hypotheses
of the square-presentation proof for every selected coordinate
group ring. Its size may depend on the signature and the
auxiliary level.

\subsubsection*{The residual coefficient of the signature determinant}
\label{subsec:residual-twist}

The signature lattice now gives a square matrix. To control its determinant at a specialization, we need a coefficient detectable from the mod-\(\ell\) Selmer classes and the Frobenius elements chosen for them.

Keep the notation of Section~\ref{subsec:signature-selmer} and write
\(W=T\otimes k\), \(W'=W^\vee(1)\). Conjugation of a group
element by complex conjugation is denoted by \(h\mapsto h^c\).
The good-signature duality gives an integral isomorphism of
\((T^\vee(1))^c\) with \(T\otimes\epsilon^t\), where \(\epsilon\)
is cyclotomic. The tame marking gives an inertia element \(z\)
acting on \(W\) by \(d\), with \(d^2\ne1\), and on \(W'\) by
\(d^{-1}\). At \(z^c\) the two scalars are reversed.
Both modules are absolutely irreducible. Let \(S,S'\) be the
residual Selmer groups for these modules at empty auxiliary
level, with finite local conditions at \(\ell\) defined by
the integral lattices.

\begin{lemma}[Residual Selmer dimensions]
For the residual primal and dual Selmer spaces of a good signature over its residue field \(k\),
\[
 \dim_k S=\dim_k S'.
\]
\end{lemma}
\begin{proof}
At \(v\mid\ell\), the cup of two integral finite cocycles has
zero rational invariant by Bloch--Kato orthogonality. The local
invariant map embeds integral \(H^2\) in rational \(H^2\) on the
scalar Tate twist, hence the reductions are orthogonal. For the
lattice of \(W\), Section~\ref{subsec:signature-selmer} gives
\[
 \dim L_v(k)-\dim H^0(E_v,W)
   =\dim H^1_f(E_v,V)-\dim H^0(E_v,V).
\]
The right side is the dimension of the de Rham quotient by
\(\operatorname{Fil}^0\). Add this equality to its analogue for
the dual lattice. At each embedding the two quotient dimensions
sum to \(r\) by the dual Tate filtration. Local duality identifies the invariants of the Tate dual with the
dual of \(H^2(E_v,W)\). Thus the local Euler characteristic reads
\[
 \dim L_v(k)+\dim L'_v(k)
 = [E_v:\mathbb Q_\ell]r+\dim H^0(E_v,W)+\dim H^2(E_v,W)
 =\dim H^1(E_v,W).
\]
The orthogonal subspaces therefore have complementary dimensions.

At the type and distinguished places
away from \(\ell\) the inertia test of the signature excludes
invariants on the residual modules: an
inertia element acts with no eigenvalue one. The dual has
the inverse eigenvalues there. The local Euler characteristic
and duality give \(H^1=0\). Away from these places and
\(\ell\), the unramified subgroups are exact annihilators.
Global duality (the finite-module Selmer order formula) now
gives \(|S|/|S'|=1\): the global invariant factors are one,
the complex factors contribute \(|W|^{-f}\), and the
\(\ell\)-adic factors contribute \(|W|^f\) by
\eqref{eq:signature-local-size}. \cite{bloch-kato,milne-duality}

\end{proof}
\paragraph{Forms specified by Frobenius elements.}
There is an open normal subgroup \(J\subset H\), stable under
\(c\), on which the residual representations and the action
on \(\mu_{2L}\) are trivial and which contains all commutators
\([g,z]\). Intersect \(H\) with the cyclotomic and residual
kernels and with the conjugate of the intersection. The
commutators lie in these groups because \(H\) contains
inertia at the distinguished pair, cyclotomic is abelian,
and the marked elements act by scalars. The same construction
works simultaneously for finitely many signatures. If needed
one may replace each residual kernel here by the kernel modulo
\((\ell^{\nu_\ell})\mathfrak m_{O_0}\): the distinguished
inertia acts by scalars integrally on \(T\).

For \(h\in J\) put
\[
       B_h(x,y)=x(h)\cdot y(h)-x(h^c)\cdot y(h^c)
            \qquad (x\in S,\ y\in S').
\]
\begin{lemma}[Corank of the Frobenius pairing]
\label{lem:residual-frobenius-forms}
For every good signature there is a linear combination \(B\) of the second-half Frobenius forms for which
\[
 \dim\ker\bigl(S\longrightarrow(S')^*,\ x\longmapsto B(x,\cdot)\bigr)\le \mathbf1_{\nu_\ell>0}.
\] The same second-half primes can be
used for the finitely many signatures. Before combining their
lists, each requires at most \((\dim S+\dim S')^2\) elements of \(J\).
\end{lemma}
\begin{proof}
The values are independent of the cocycle representatives.
Restriction to \(J\) is injective on cohomology. In fact a
cocycle \(x\) vanishing there factors through a quotient in
which \(z\) is central. The two expressions for \(x(gz)=x(zg)\)
give
\[
 (d-1)x(g)=(g-1)x(z),\qquad
 x(g)=(g-1)\bigl((d-1)^{-1}x(z)\bigr).
\]
Thus \(x\) is a coboundary. On the dual use \(d^{-1}\) in the
same identity; both denominators are nonzero. For independent classes
\(x_1,\ldots,x_a\) with values in \(W\), the vectors
\((x_j(h))_j\) span \(W^a\) over \(k\) as \(h\) varies in
\(J\). Let \(U\) be the span of the evaluation tuples in \(W^a\). It is stable by conjugation. If \(U\) were proper, semisimplicity and absolute irreducibility would give scalars, not all zero, with
\[
 \sum_{j=1}^a c_j x_j(h)=0\qquad(h\in J).
\]
Indeed any map \(W^a\to W\) on this isotypic space is a linear combination of coordinate projections. Restriction injectivity rules it out.

If \(W\not\simeq (W')^c\), then
\(W,W',W^c,(W')^c\) are pairwise inequivalent: the scalars
at \(z\) separate the two groups and conjugation carries the
remaining inequivalence to its partner. For bases of \(S,S'\)
the simultaneous evaluations of these four lists on \(h\)
span the direct sum of their four full isotypic spaces. This
follows from the preceding projection test on each simple
type. To show that the \(B_h\) span all bilinear forms, suppose
\[
 \sum_{i,j}a_{ij}B_h(x_i,y_j)=0
\]
for every \(h\in J\).

Write the four evaluation blocks as \(v_i,w_j,\tilde v_i,\tilde w_j\),
so the relation says that every evaluation tuple is a zero of
\[
 Q(v,w,\tilde v,\tilde w)=\sum_{i,j}a_{ij}
      (v_i\cdot w_j-\tilde v_i\cdot\tilde w_j).
\]
Evaluation on \(J\) is additive, so the bilinear form
\(B_Q(s,t)=Q(s+t)-Q(s)-Q(t)\) vanishes on all pairs of evaluation
tuples and hence on their full linear span. Fix \(i,j\); give
\(s\) just the entry \(v\) in the \(i\)-th first block and
\(t\) just the entry \(w\) in the \(j\)-th second block. Then
\begin{equation}\label{eq:frobenius-polarization-block-test}
 0=B_Q(s,t)=a_{ij}\,v\cdot w .
\end{equation}
Choosing a pair with nonzero pairing in \eqref{eq:frobenius-polarization-block-test} forces \(a_{ij}=0\), including
in characteristic two.

If \(W\simeq(W')^c\), choose an invertible intertwining
matrix \(A\), so that
\(W(g^c)^t A W(g)=\bar\epsilon(g) A\), with superscript
\(t\) on a matrix denoting transpose. Transposing the same
identity at \(g^c\) shows \(A^t=\lambda A\) by Schur's lemma.
Here \(\lambda^2=1\) and \(\lambda^r=1\) by determinants, so
\(A^t=A\) since \(r\) is odd. Transport \(S'\) to a subspace
\(S_2\) of \(H^1(E,W)\) by sending \(y\) to \(A^{-1}y^c\).
Choose a basis of \(S+S_2\) and evaluate it at \(h,h^c\).
The two evaluations span the two full isotypic blocks together
because \(z\) separates \(W\) and \(W^c\).

In this basis the
form has matrix entries
\[
     x_i(h)^t A x_j(h^c)-x_i(h^c)^t A x_j(h).
\]
A relation with coefficients \(a_{ij}\) depends on
\(a_{ij}-a_{ji}\). Polarization and testing on vectors in
the two separate blocks forces \(a_{ij}=a_{ji}\). Conversely
these symmetric relations annihilate the displayed alternating
matrices, including their diagonals. Thus the forms span the restrictions of alternating forms on
\(S+S_2\). Set \(J_0=S\cap S_2\) and take complements
\(S=J_0\oplus U\), \(S_2=J_0\oplus U'\). These complements
have the same dimension. With dual bases for the pairing
of \(U\) with \(U'\), choose the restriction \(S\times S_2\)
with block matrix
\[
 \begin{pmatrix}J_{\dim J_0}&0\\0&I_{\dim U}\end{pmatrix},
 \qquad
 J_d=\bigoplus_{i=1}^{\lfloor d/2\rfloor}
       \begin{pmatrix}0&1\\-1&0\end{pmatrix}
       \quad\text{with a final zero block if \(d\) is odd}.
\]
The opposite off-diagonal pairing on \(U'\times U\) is the
negative transpose, so this matrix extends to an alternating
form on \(S+S_2\). Its left corank is \(d\bmod 2\), including
in characteristic two.

Consequently there is a linear combination of the forms
with left corank at most \(e\), where \(e=0\) if
\(\nu_\ell=0\) and \(e=1\) otherwise. Indeed an equivalence
\(W\simeq(W')^c\) would imply
\[
 \bar\epsilon^{tr}=1
\]
by taking determinants in the signature duality, which
implies \(\nu_\ell>0\) by \eqref{eq:signature-cyclotomic-minimum}.
The number of elements \(h\) needed is at most
\((\dim S+\dim S')^2\), independently of \(J\) and \(M\).
Choose such elements for each good signature and concatenate
the lists, padding with identities if necessary.

For the finite list \(\mathcal Z\) of restrictions of the basis cocycles and their conjugates, use the common kernel
\[
 J_0=\bigcap_{z\in\mathcal Z}\ker(z|_J).
\]
The conjugation formula gives an open normal subgroup of \(J\).
Apply Section~\ref{subsec:aux-coordinates} prescribing these
elements for the second half. On replacing \(h\) by a
Frobenius conjugate to it modulo this subgroup the matrix
\(B_h\) changes only by the nonzero cyclotomic scalar of the
conjugating element. Absorb it into the coefficients. Thus at
the second-half Frobenius elements \(\phi_i\) we have, for each signature,
coefficients \(b_i\) such that \(\sum_i b_i B_{\phi_i}\) has
the stated corank. The same primes work for all signatures.

\end{proof}
\begin{lemma}[A residual coefficient from the Selmer pencil]\label{lem:residual-selmer-coefficient}
Use the two halves of auxiliary places and the Frobenius forms
of Lemma~\ref{lem:residual-frobenius-forms}, with \(R\) places
in each half and \(L>4rR+2\). When \(\nu_\ell=0\), every polynomial determinant lift \(D\) satisfies
\[
 \exists\alpha:\quad |\alpha|\le4rR,\qquad [X^\alpha]D\in O^\times.
\] For positive \(\nu_\ell\) the corresponding
unit coefficient is obtained at this stage in an adjugate
entry, for a nonempty square presentation. An empty presentation
already has determinant one.
\end{lemma}

\begin{proof}
\proofstep{Use the inertia and Frobenius evaluations in the relaxed residual space.}
Let \(R\) denote the common length of each half. At the first
half the untwisted residual module is unramified with Frobenius
scalar \(d\) or \(d^{-1}\); both differ from one. The
cyclotomic character is trivial there modulo \(\ell\).
Thus \(H^0,H^2\), and by the local Euler characteristic also
\(H^1\), vanish. At each second-half prime the Frobenius is
in \(J\), as is inertia. Fix a tame \(\ell\)-generator
\(\gamma_i\) with \(d_{P'_i}(\gamma_i)=1\). This is possible
since the character on inertia is onto. In the following
evaluation pairs, use at the conjugate place the transports
\(\gamma_i^c,\phi_i^c\). Superscript \(c\) on group elements
continues to denote conjugation.

Write \(S_Q\) for the relaxed residual Selmer space and
\[
 D=(k^r\oplus k^r)^R.
\]
Define \(I:S_Q\to D\) by evaluation
at \((\gamma_i,\gamma_i^c)_i\). Its kernel is \(S\), since
on the trivial residual module a homomorphism on inertia
factors through its pro-\(\ell\) tame quotient. The first-half
places contribute nothing. For \(x\in S_Q\), \(y\in S'\),
reciprocity gives units \(a_i\in k^\times\), independent of
\(x,y\), such that
\begin{equation}\label{eq:residual-reciprocity}
 \sum_i a_i\{x(\gamma_i)\cdot y(\phi_i)
             -x(\gamma_i^c)\cdot y(\phi_i^c)\}=0 .
\end{equation}
Indeed the cup invariant at an unramified pair of classes is
zero, and at the primes over \(\ell\) it vanishes by the local
orthogonality above.

At a second-half place the second class
is unramified. The invariant on a tame character of
\(\gamma_i\) paired with an unramified character of \(\phi_i\)
is a nonzero scalar: the two unramified classes cup to zero
and local Tate duality is nondegenerate on the two-dimensional
space for the trivial residual line. Bilinearity gives the
formula for vectors. The transport to the conjugate place
reverses the chosen identification of \(\mu_\ell\) with the
prime field, since \(c\) acts by \(-1\) there. This explains
the common scalar and the minus sign. The sum of local
invariants vanishes by global reciprocity.

Define
\[
 F:S_Q\to D
\]
by taking the pairs of Frobenius
values of each class multiplied by \(b_i/a_i\). The two necessary linear conditions on a lift are
\[
 Ix=0,\qquad Fx=Ix_1\quad(x_1\in S_Q).
\]
The first gives \(x\in S\) and \eqref{eq:residual-reciprocity} applied
to \(x_1\) puts \(x\) in the left radical of
\(\sum b_i B_{\phi_i}\). Consequently
\begin{equation}\label{eq:residual-pencil}
  \dim\{x\in\ker I: Fx\in\operatorname{im}I\}\le e,
  \qquad \dim S_Q\le 2rR+e .
\end{equation}
For the second inequality apply rank-nullity to
\(\ker I\to D/\operatorname{im}I\) induced by \(F\).

We give the specialization producing a coefficient of bounded
degree; setting all group variables to the same variable would
lose the coefficients \(b_i/a_i\). Put \(K=k((s))\).
In the polynomial presentation matrix reduce modulo the
uniformizer and set \(X_j=\alpha_j Z\), \(Y_i=\beta_i Z\),
where
\[
  \beta_i=s\, b_i/a_i,\qquad
  \alpha_j=1-\sum_i\beta_i\,\overline{d_{P'_i}(\phi_j)} .
\]
Write \(P(Z)\) for this square matrix over \(K[Z]\).
Modulo \(Z^2\) the substitution factors through the group ring,
since \(L\) is a positive power of \(\ell\). On the second-half
decomposition group the coefficient of \(Z\) in the inverse
universal character is \(-1\) at \(\phi_i\) and \(-\beta_i\)
at \(\gamma_i\), by \eqref{eq:aux-coordinate}. At the
conjugate place both signs reverse.

A first-order lift of \(v\in\ker P(0)\) is a vector \(v+Zv_1\) satisfying
\[
 P(Z)(v+Zv_1)\equiv0\pmod{Z^2}.
\]
Naturality of the square presentation gives
a Selmer cocycle in the dual numbers lifting its residual
class \(x\in K\otimes_k S_Q\). This identification commutes
with taking its constant term. Indeed additive functoriality
gives the assertion for finite sums of \(k\), and each class
over a discrete extension field is supported on a finite
coefficient subspace.

At each second-half place write the
cocycle as \(x+Zx_1\), and the character as \(1+Z\lambda\).
The tame relation
\[
 \phi\gamma\phi^{-1}=\gamma^q,
\]
where
\(q=\mathrm NP'_i\), gives
\[
             \lambda(\phi)x(\gamma)=\lambda(\gamma)x(\phi).
\]
For this calculation the crossed term of the first-order
cocycle at \(g,h\) is \(\lambda(g)x(h)\). In its value at
\(\gamma^q\) the coefficient \(\binom q2\) vanishes in \(k\)
because \(q\equiv1\bmod 2L\); the coefficient \(q\) is one.
The prime-to-\(\ell\) inertia kernel has zero image in the
successive additive groups used in this calculation.
Thus \(Ix=sFx\), including at the conjugate places.

The subspace of \(K\otimes S_Q\) satisfying this equation
has dimension at most \(e\). Intersect it with the lattice
\(S_Q[[s]]\). This saturated sublattice has reduction of
dimension equal to its rank and may be given a basis with
independent constant terms by elimination over \(k[[s]]\).
Coefficient comparison in \(Ix=sFx\), with \(x=x_0+sx_1+\cdots\), gives
\[
 Ix_0=0,\qquad Ix_1=Fx_0.
\]
Now apply
\eqref{eq:residual-pencil}. This argument can be used over
any finite residue field, including \(\mathbb F_2\).

\proofstep{Use the first-order kernel to bound the number of invariant factors of order at least two.}
The evaluations have bounded the kernel of the first-order matrix. Reading the same matrix through its invariant factors converts that kernel bound into the restriction on the determinant's low-degree coefficients.

Take Smith form over \(K[[Z]]\), with orders
\(0\le a_1\le\cdots\le a_d\le\infty\), where a zero
diagonal entry has order \(\infty\). Reduce the kernel of the \(i\)-th diagonal entry on \(K[[Z]]/(Z^2)\)
modulo \(Z\). Its image in \(K\) is
\[
 \begin{cases}
 0,&a_i=0\text{ or }a_i=1,\\
 K,&2\le a_i\le\infty.
 \end{cases}
\]
For \(a_i=1\) its kernel is precisely the ideal \((Z)\), whose
residue in \(K\) is zero.
The dimensions just computed give
\[
 \#\{i:a_i>0\}=h=\dim S_Q\le2rR,
 \qquad \#\{i:a_i\ge2\}\le e.
\]
If \(e=0\), all Smith exponents are zero or one. Therefore
\[
 \operatorname{ord}_Z\det P=\sum_i a_i=\#\{i:a_i=1\}=h .
\]
If \(e=1\) and \(d>0\), the ideal generated by the submaximal
minors has order
\[
 \sum_{i=1}^{d-1}a_i\le\max(0,h-1)\le2rR.
\]

At least one adjugate entry has that order. For an empty matrix
the determinant is already one. Both conclusions suffice with
the degree bound \(4rR\). For \(L>4rR+2\) the
choice of polynomial lift cannot change any reduced coefficient
in these degrees: the coordinate relations reduce to
\(X_j^L=Y_i^L=0\). A nonzero coefficient of the line
specialization requires a nonzero coefficient of the same total
degree before substitution. Hence every lift of the determinant
has a unit coefficient of bounded degree when \(\nu_\ell=0\).
For \(\nu_\ell>0\) the conclusion at this stage concerns one
adjugate entry. Its conversion to a determinant with valuation
at most \(\nu_\ell\) requires one further arithmetic step.
\end{proof}


\subsubsection*{The positive cyclotomic valuation}
\label{subsec:positive-twist}

Suppose the cyclotomic minimum \(\nu=\nu_\ell\) of
\eqref{eq:signature-cyclotomic-minimum} is positive. Write \(O=O_0\), choose a
uniformizer \(\pi\), and put
\[ e=\operatorname{ord}_O(\ell^\nu),\qquad I=\ell^\nu\mathfrak m_O=(\pi^{e+1}). \]
The residual calculation detects an entry of the adjugate. One elementary divisor can still carry the entire remaining loss, so the determinant itself must now be detected modulo \(I\).

Let \(b\) be the number of good signatures at this coefficient
place, and \(R_0\) the sum of the bounds for the number of Frobenius elements in their
residual tests, padding each list with identities. Use \(R=R_0+b\) first pairs, and initially
\(n=R_0\) second pairs. Prescribe the first \(R_0\) Frobenius elements of
the second list as in Section~\ref{subsec:residual-twist}.
Future second pairs will use the remaining first coordinates
one at a time. Write \(p=2r\) and set
\[
 w_0=4rR_0,\qquad w_{j+1}=2(e+1)w_j+2+p .
\]
Only the finite signature data enter these constants. Throughout
take \(M\) so large that \(\ell^{M-\nu-1}>w_b+2\) and
\(M\ge\nu+2\).

\begin{lemma}[Independence of low-degree coefficients]
Let \(F,G\) lift the same group-ring element. After reduction to \(O/I\), or to the residue field, their coefficients agree in the range
\[
 [X^\alpha]F=[X^\alpha]G\qquad(|\alpha|\le w_b+1).
\]
\end{lemma}
\begin{proof}
For \(0<i<\ell^M\) one has
\[
 i\binom{\ell^M}{i}=\ell^M\binom{\ell^M-1}{i-1},\qquad
 v_\ell\binom{\ell^M}{i}=M-v_\ell(i),
\]
since \(\binom{\ell^M-1}{i-1}\) is a unit by expansion modulo \(\ell\).
For \(1\le i\le w_b+1<\ell^{M-\nu-1}\) the coefficient belongs
to \(I\). Each coordinate relation \((1+U)^{\ell^M}-1\), and hence
each polynomial multiple, has zero coefficients in these degrees.
\end{proof}

The induction keeps a unit coefficient in a cofactor through
degree \(w_j\) for every signature with a nonempty presentation,
after setting to zero the unused first variables. For each of
the first \(j\) signatures it also keeps a determinant coefficient
outside \(I\) through degree \(w_j\), with all variables allowed.
To start, set the extra first variables to zero. At these places
the residual Frobenius is scalar different from one on either
side, and its norm is one modulo \(\ell\). Writing \(D\) for that residual representation, local duality and the local Euler characteristic give
\[
 H^j(E_v,D)=0\qquad(j=0,1,2).
\]

Successive extensions give the same vanishing on finite
coefficient modules after this specialization. Deleting the
extra pairs therefore gives the old Selmer functor. Its representing
module after scalar change is the old module by naturality.
The first Fitting ideal (the ideal generated by the submaximal
minors of a square presentation) commutes with this base change.
The unit cofactor from Section~\ref{subsec:residual-twist}
therefore supplies the initial assertion for any presentation.
An empty presentation has determinant one; it can also be
stabilized by a unit diagonal entry.

\paragraph{Preserving a coefficient at a larger level.}
Consider a new second pair satisfying the coordinate conditions
and lying in the subgroup \(J\) for every signature.
\begin{lemma}[Fitting ideals at an added pair]
After setting the new second variable to zero, let \(P_+\) and
\(P_0\) be the new and old modules reduced modulo \(I\).
For the Frobenius \(\sigma\) at the new prime put \(u=\eta(\sigma)-1\).
Then
\[
 u^p\operatorname{Fitt}_i(P_0)\subseteq\operatorname{Fitt}_i(P_+)\qquad(i=0,1).
\] A nonzero determinant coefficient modulo \(I\), or
a unit cofactor coefficient modulo \(\pi\) with unused variables
zero, can be retained with degree increased by at most \(p\)
in the asserted low-degree range.
\end{lemma}
\begin{proof}
First set its own variable to zero and reduce modulo \(I\).
Write \(P_+\) for the new representing module after this
specialization and \(P_0\) for the old module. There is a
surjection
\[
 P_+\to P_0
\]
with kernel generated by \(p\)
inertia coordinates. In fact in the new Selmer functor the
unramified classes at the new pair are exactly the old ones.
The old action is unramified there; the only possible
inertia homomorphisms on finite primary modules factor through
the pro-\(\ell\) tame quotient. Evaluating a representative at
one generator on each side gives the claimed rows in \(P_+\),
independent of coboundaries. The kernel description now follows
by testing on all finite modules. Here the specialized rings
are finite.

Put \(u=\eta(\sigma)-1\) for the old universal character on
a Frobenius at the new prime, choosing the direct character
\(\eta\). For the inertia generators \(t_1,\ldots,t_p\) of the kernel this gives the scalar relations
\[
 (\eta(\sigma)-1)t_i=u t_i=0\qquad(1\le i\le p).
\]

Indeed the untwisted Frobenius matrix is one modulo \(I\) on
\(J\) and the tame relation multiplies an inertia value by
the residue norm, which is one in \(O/I\). Inversion on the
conjugate side or in the twist changes this annihilator by
a unit. If \(A_0\) is a presentation matrix for \(P_0\), lift its
relations to the enlarged list of generators for \(P_+\).
Together with the inertia relations these give a block of the
new presentation of the form
\[
        \begin{pmatrix} \widetilde A_0 & *\\ 0 & u I_p\end{pmatrix}
\]
when rows write relations. Adjoining all \(p\) inertia rows
and columns to a minor for \(\operatorname{Fitt}_i(P_0)\)
multiplies it by \(u^p\).

For \(i=0,1\) this proves
\[
        u^p\operatorname{Fitt}_i(P_0)
                \ \subseteq\ \operatorname{Fitt}_i(P_+).
\]
It suffices that these are relations; further relations enlarge
the ideal. One may stabilize either presentation for \(i=1\).
The polynomial for \(u\) has constant term zero and its
linear coefficient at the newly used first variable is one
by the Frobenius coordinate prescription. Even after deleting
the still unused first variables its linear form is monic in
that variable. Write \(U_1\) for its linear part. Over \(O/I\), a monic
polynomial in the newly used first variable is a non-zero-divisor.
If \(F_d\ne0\) is the lowest nonzero homogeneous component of an
element \(F\) in the old ideal, then
\[
 (u^pF)_{d+p}=U_1^p F_d\ne0.
\]

All smaller degrees are zero. The degree test for the coordinate
relations above then shows that any old detection through
degree \(w\) gives some new generator detected through
\(w+p\). Apply this modulo \(I\) to the determinant, and
modulo \(\pi\) to the adjugate with unused first variables
zero. In the latter test the old cofactor survives when the
newly used first variable is restored, since it survived even
with that variable zero.

\end{proof}
\paragraph{Producing a dual class.}
Process one remaining signature, with current bound \(w\).
Put \(D=(e+1)w\). If its old determinant is already detected
through degree \(2D+1\), take a new pair with trivial prescribed
Frobenius modulo the common subgroup and use the preceding
comparison. Otherwise choose an old square presentation with
matrix \(\Phi\) of positive size \(a\) and a polynomial lift
of its determinant with all coefficients through degree
\(2D+1\) in \(I\).

Choose a cofactor, written as an adjugate
entry \(H\), with a unit coefficient through degree \(w\)
after the unused first variables are set to zero. Choose a nonzero homogeneous part \(\bar H_j\) of the residual cofactor in the used variables, with \(j\le w\). The slopes \(\bar a_i\) must satisfy
\[
 \bar a_i\ne0\quad (i\text{ used}),\qquad \bar H_j((\bar a_i)_i)\ne0.
\]
Such a vector exists on the torus over
the algebraic closure and has finitely many algebraic coordinates.
Lift the slopes to units \(a_i\) in a finite unramified DVR
extension \(O'/O\), putting \(a_i=0\) at unused variables.
Work in
\[
          A=(O'/I O')[\tau]/(\tau^{2D+2}),\qquad
          U_i=a_i \tau,\qquad c_0=\ell^\nu \tau^D .
\]
\begin{lemma}[A dual Selmer class in the truncated ring]
For the resulting specialization \(\Psi\) with zero determinant
and unit cofactor coefficient in degree at most \(w\), let
\(H=(\operatorname{adj}\Psi)_{jl}\) denote the selected cofactor.
Column \(l\) gives a primal Selmer class. On the dual side,
\[
 \exists\,y\in S^\perp(A):\qquad c_0y\ne0.
\]
\end{lemma}
\begin{proof}
The coordinate relations hold by the binomial estimate. The
specialized matrix \(\Psi\) has determinant zero. Suppose the
first unit coefficient of the specialized \(H\) occurs at
degree \(d\le w\). Then \(H=\pi B+\tau^d V\) with \(V\) a
unit. Modulo \(H\), the relation \(\tau^d=-\pi B/V\) gives
\(\tau^{(e+1)d}=0\). Thus \(H\mid\tau^D\), since \(D=(e+1)w\).
Because \(\Psi\,\operatorname{adj}\Psi=(\det\Psi)I_a=0\),
column \(l\) represents \(x\in S(A)\).
Writing \(\tau^D=H H'\), one has
\[
 (c_0\tau H)H'=\ell^\nu\tau^{2D+1}\ne0
 \quad\text{in }(O'/IO')[\tau]/(\tau^{2D+2}).
\]

We use the natural identification of kernel vectors with
classes for every coefficient module. Representatives can be
chosen from one cocycle with coefficients in the representing
module, by applying the functor isomorphism to its identity.
A chosen set of lifts of its entries to the free presentation
then gives matrices whose contraction with a kernel vector
represents the corresponding cocycle.

Let \(Y=\operatorname{Hom}(A,\mathbb Q_\ell/\mathbb Z_\ell)\)
be the additive dual with the balanced evaluation pairing.
Write \(S^\perp\) for the Selmer condition on the Tate dual
with inverse character, orthogonal to the signature condition
on \(Y\). The finite-module Selmer order formula gives
\begin{equation}\label{eq:pos-order}
 |\ker\Psi(Y)|=|S^\perp(A)|
             \prod_{v\in Q}|H^0(E_v,T^\vee(1)\otimes\kappa^{-1} A)|.
\end{equation}
Here \(\kappa\) is the specialized inverse universal
character, and the product includes both places of each pair.

To check the local factors for arbitrary finite coefficients,
use an integral basis of the saturated finite-cocycle module
at \(v\mid\ell\). Its residual functions are independent:
a dependency would divide a primitive integral combination
by \(\pi\) in the same lattice. Finitely many evaluations
give an invertible minor. The span on a finite module \(Y\)
thus has order \(|Y|^{r+d_v}\), where \(d_v\) is the de Rham
quotient dimension in Section~\ref{subsec:signature-selmer}.

Let \(Z_v(Y)\) be this span of cocycles, \(B^1_v(Y)\)
its coboundaries, and \(L_v(Y)=Z_v(Y)/B^1_v(Y)\).
The coboundary map has source of order \(|Y|^r\) and
kernel \(H^0_v\), so
\[
 |L_v(Y)|=\frac{|Z_v(Y)|}{|B^1_v(Y)|}
          =\frac{|Y|^{r+d_v}}{|Y|^r/|H^0_v|}
          =|Y|^{d_v}|H^0_v|.
\]
The product of the factors \(|Y|^{d_v}\) cancels the
complex-place factor.
At the full local conditions away from \(\ell\) the local
Euler characteristic leaves the dual invariants. They vanish
at the exceptional places by the inertia test. Outside the
ramification set unramified conditions are paired. Both global
invariant groups vanish by the scalar marking. This proves
\eqref{eq:pos-order}, and the same proof applies to the
perfectly paired modules \(Y/c_0\tau Y\) and \(A[c_0\tau]\).
\cite{milne-duality}

There exists \(y\in S^\perp(A)\) with \(c_0[y]\ne0\).
Indeed at a first pair the dual invariants vanish by the
residual Frobenius test and finite dévissage. At each current
second pair the inertia of its own coordinate has value
\(1+a_i \tau\) with \(a_i\) a unit. For any invariant \(v\) at such a place,
\[
 (1+a_i\tau)v=v\quad\Longrightarrow\quad \tau v=0 .
\] Suppose all the dual classes were killed
by \(c_0\). For a representative \(y\), write
\[
 c_0y=dm.
\]

The image of \(m\) in the coefficient quotient by \(c_0\)
is invariant globally, hence zero. Subtracting a coboundary
from \(y\) gives values in \(A[c_0]\). At a strict place a
local expression \(y=dm_v\) then has \(c_0m_v\) invariant,
so \(c_0\tau m_v=0\). At \(\ell\) every finite local class
on \(Y/c_0\tau Y\) lifts from \(Y\) by the local description;
naturality of the pairing gives the required orthogonality.
The class thus lifts to the dual Selmer group on \(A[c_0\tau]\).
The invariant products in \eqref{eq:pos-order} coincide,
so
\[
 |\ker\Psi(Y)|\le |\ker\Psi(Y/c_0\tau Y)|.
\]

For an endomorphism \(U\) of a finite module \(M\),
\[
 |\ker U|=|M|/|U(M)|=|\operatorname{coker}U|.
\] Accordingly, the kernel inequality implies
that the quotient surjection
\[
 \frac{Y^a}{\Psi Y^a}
 \longrightarrow\frac{Y^a}{\Psi Y^a+c_0\tau Y^a}
\]
is bijective. Hence \(c_0\tau Y^a\subseteq\Psi Y^a\). Applying
row \(j\) of \(\operatorname{adj}\Psi\), which annihilates this
image, to \(c_0\tau Y\) in coordinate \(l\) shows
\((c_0\tau H)Y=0\). The evaluation pairing of \(A\) with \(Y\)
is perfect, implying \(c_0\tau H=0\), a contradiction.

\end{proof}
\paragraph{A Frobenius element with the prescribed evaluations.}
Intersect \(J\) with the kernels of the used coordinates to
obtain \(J'\). This subgroup is normal and conjugation-stable.
The two twisted actions are trivial on it, the commutators
\([g,z]\) lie in it, and \(\kappa(z)=1\) by the ramification
condition. Both restrictions of the cocycles to \(J'\) are
therefore additive, independent of representatives and
equivariant under conjugation.

\begin{lemma}[A second-half Frobenius with nonzero pairing]
For the primal and dual classes just chosen, there is a new prime
pair, avoiding the previous exceptional set, with the prescribed first-coordinate values, satisfying at the chosen prime \(P\), with Frobenius \(\sigma\),
\[
 \mathrm NP\equiv1\pmod{2\ell^M},\qquad B_\sigma\ne0.
\] Its untwisted Frobenius and that
of its conjugate act trivially modulo \(I\).
\end{lemma}
\begin{proof}
Some \(h\in J'\) has
\[
 B_h=x(h)\cdot y(h)-x(h^c)\cdot y(h^c)\ne0\quad\text{in } A .
\]
For suppose all these differences vanished. Polarizing at
\(h,k\in J'\) gives equality of the two cross sums. Replacing
\(h\) by \(zhz^{-1}\) multiplies its \(x\)-value by \(d\)
and its \(y\)-value by \(d^{-1}\), with the inverse multipliers
on the conjugate pair. Since \(d-d^{-1}\) is a unit, subtraction
separates the two terms and yields
\[
            x(h)\cdot y(k)=x(k^c)\cdot y(h^c).
\]

Let \(R\) and \(R'\) be the twisted primal and dual actions.
Conjugate duality and anti-invariance of \(\kappa\) identify
\((R')^c\) with \(\chi R\), \(\chi=\epsilon^{\text{tw}}\) with \(\text{tw}\) the
signature's fixed twist exponent.
Replace \(h\) by \(ghg^{-1}\) in the last identity.
Equivariance on \(J'\) gives, for every \(g\),
\[
 \operatorname{tr}(R(g)U)=\operatorname{tr}(\chi(g)R(g)V),
 \qquad U=x(h)y(k)^{\mathsf T},
\]
where \(V\) is the conjugate of \(y(h^c)x(k^c)^{\mathsf T}\)
by the fixed intertwiner. The span of \(R(g)\) is the full
matrix algebra over \(A\), by residual absolute irreducibility
and Nakayama. The same holds for \(\chi R\).

Consider the subalgebra spanned by the pairs \((R(g),\chi(g)R(g))\) in
\(M_r(A)^2\). It surjects onto both factors. Its intersections
with the separate factors are two-sided ideals, of the form
\(M_r(I_1)\) and \(M_r(I_2)\) by multiplication with matrix
units. If \(\mathcal A\) is this subalgebra, a pair \((B,C)\in\mathcal A\)
associates the class of \(B\) modulo \(M_r(I_1)\) to the class
of \(C\) modulo \(M_r(I_2)\). Surjectivity and the definitions
of \(I_1,I_2\) give a well-defined inverse in the other direction.
It is an isomorphism
\[
      M_r(A/I_1)\xrightarrow{\sim} M_r(A/I_2)
\]
of \(A\)-algebras, so their scalar annihilators coincide and
\(I_1=I_2=I'\).

If the quotient is nonzero its matrix algebra
automorphism is inner. Indeed the images of the diagonal matrix
units cut the standard module over the local ring \(A/I'\)
into isomorphic projective summands, necessarily free of rank
one. Transport a generator of the first by the images of the
off-diagonal matrix units; the resulting basis conjugates
all matrix units at once. Taking determinants in the quotient
gives \(\chi(g)^r-1\in I'\) for every \(g\), also true for the
zero quotient. By definition of \(\nu\), these scalars generate
\((\ell^\nu)\) after coefficient extension. The trace identity
on \((\ell^\nu E_{ab},0)\) gives
\[
 \ell^\nu x(h)_i y(k)_j=0
\]
for all coordinates.

Let \(J_x\) be the ideal generated by the entries of all
\(x(h)\), \(h\in J'\). Modulo \(J_x\) the cocycle factors
through \(G_E/J'\), on which \(z\) is central and acts by
a scalar differing from one by a unit. Comparing the cocycle law at \(zg\) and \(gz\) gives
\[
        (d-1)x(g)=(g-1)x(z)
\]
over \(A/J_x\). As \(d-1\) is a unit, this expresses the
cocycle as a coboundary with constant \((d-1)^{-1}x(z)\). Naturality and injectivity
of the kernel description put every entry of column \(l\)
of the adjugate in \(J_x\).

In particular \(\tau^D\in J_x\)
since \(H\mid\tau^D\). As \(\ell^\nu J_x y(k)=0\) for \(k\in J'\),
this gives
\[
 c_0y|_{J'}=0.
\]
In the quotient
\(G_E/J'\), the central element \(z\) acts on the dual by \(d^{-1}\)
with \(d^{-1}-1\) invertible. For \(f=c_0y\) the two evaluations of \(f(zg)=f(gz)\) give
\[
 (d^{-1}-1)f(g)=(g-1)f(z).
\]
Dividing by the unit exhibits \(f\) as a coboundary, contradicting the preceding lemma.

Choose an open normal subgroup in \(J'\) on which both
cocycles and their conjugates vanish. This is possible since
their restrictions are continuous homomorphisms to finite
modules; intersect the kernels and take the normal core.
This subgroup can contain all inertia at the unused first
pairs. There the used characters and signature actions are
unramified and the residual Frobenius test gives vanishing
of local \(H^1\) for both modules by finite dévissage.
Local coboundaries vanish on inertia, also at every conjugate
place and for every conjugate of its inertia group.
By Chebotarev choose a new second pair with Frobenius
\(\sigma\) conjugate to \(h\) modulo this subgroup, prescribing the first-coordinate vector
\[
 (0,\ldots,0,1,0,\ldots,0)
\]
with the single \(1\) in the next unused first coordinate (additive notation). Indeed \(h\) is in
the kernels of the used coordinates and the unused first
inertia groups can correct all the remaining coordinates.

All old used characters then vanish on the new decomposition
group. The Frobenius can be adjusted by its own inertia to
give value zero also in the new coordinate. Equivariance and
the dual pairing multiply \(B_h\) under conjugation by a
cyclotomic unit, equal on the two conjugate terms. Thus
\(B_\sigma\ne0\). The new prime avoids all exceptional
places, has norm one modulo \(2\ell^M\), and its untwisted
action and that at its conjugate are trivial modulo \(I\).
At each stage the open subgroup and avoidance set may include
the common earlier Frobenius requirements of the other signatures.

\end{proof}
\begin{proposition}[A determinant coefficient after all positive-valuation steps]\label{prop:positive-signature-determinant}
Starting from Lemma~\ref{lem:residual-selmer-coefficient},
under the hypotheses in this construction, after choosing the
reserve pairs, there is a common degree bound \(D\) such that every signature determinant \(F_s\) has a coefficient satisfying
\[
 \exists\alpha:\quad |\alpha|\le D,\qquad v_\ell([X^\alpha]F_s)\le\nu_\ell.
\] The number of pairs, degree and depth threshold depend
only on the original finite data. For each sufficiently large
\(M\), one choice of pairs serves all signatures. If
\(\nu_\ell=0\), no reserve pair is needed.
\end{proposition}

\begin{proof}
\proofstep{Make the new coordinate substitutions before taking the Fitting ideal.}
Use the finite ring
\[
 B=A[X,Y]/(X^2,Y^p),
\]
over \(A\) by
constant inclusion and with projection \(B\to A\).
Keep the old used variables \(a_i \tau\), set the reserve first
variable corresponding to this pair to \(Y\), and set its
own second variable to \(X\). These assignments satisfy all
coordinate relations by the binomial bound. At the new
prime take a tame generator \(\gamma\) mapped to one in the
new coordinate. The inverse universal character has values
\((1+X)^{-1}\) and \((1+Y)^{-1}\) on \(\gamma,\sigma\);
the conjugate pair has inverse values. On projection to \(A\)
all these values become one.

Let \(P'\) be the new representing module after this base
change. For its universal cocycle write \(u_i,t_i\in P'\)
for the Frobenius and inertia entries at the two places,
\(1\le i\le p\). The tame relation gives
\[
          X u_i=Y\eta_i t_i,\qquad
             \eta_i\in B^\times,\quad \eta_i\bmod(X,Y)=1 .
\]
To verify it write \(q\) for the norm. On the pro-\(\ell\)
tame quotient
\[
 \sigma\gamma=\gamma^q\sigma.
\]

For an inertia scalar \(1+v\), \(v^2=0\), the cocycle on
\(\gamma^q\) is \((q+\binom q2 v)\) times its value on
\(\gamma\), hence just its value on \(\gamma\), since
\(q\equiv1\bmod 2\ell^M\). Wild and prime-to-\(\ell\) tame
inertia in the kernel of the action have zero cocycle values.
Expanding the relation now gives the displayed identity,
with
\[
 \eta_i=(1+Y)^{-1}
\]
or one.

In \(P'/(X,Y)P'\), quotienting by the inertia entries gives
the old module with presentation \(\Psi\), by the unramified
condition. Choose lifts of its \(a\) generators. The projected
Frobenius rows modulo inertia are the old rows \(U_i\)
evaluated at \(\sigma,\sigma^c\); an ambiguity of the
representative is a coboundary and vanishes at these elements
of trivial action. Define \(k_i\) to be the entries of
\(y(\sigma)\) on the first side and \(-y(\sigma^c)\) on the
second. Global reciprocity gives a further relation
\[
                \sum k_i t_i=0 \pmod{(X,Y)P'} .
\]
Indeed test against the additive dual of \(A\). A new relaxed
class on that dual with the old character pairs with the old
dual class \(y\). The local cup invariants vanish at the old
places by their conditions, and outside the level by
unramifiedness.

At each of the two new places the action is
trivial and \(y\) is unramified. Only the new inertia value pairs with \(y\)'s Frobenius value. To determine the two signs, identify the scalar Tate coefficients by a primitive \(\ell^{\nu+1}\)-st root. By local reciprocity the unramified character taking Frobenius to one pairs to a generator with the Kummer class of a uniformizer, whose tame pro-\(\ell\) inertia value is a unit. Unramified parts pair to zero.
Conjugation takes the chosen primitive root to its inverse. Thus, in the additive identification of the scalar Tate coefficients,
\[
 \zeta^a\longmapsto\zeta^{-a},\qquad a\longmapsto-a.
\]
This gives the opposite sign at the conjugate place. Applying the
sum of invariants and varying the homomorphism to the additive
dual proves the relation, since that dual is a cogenerator
for finite \(A\)-modules. \cite{milne-duality,milne-cft}

The lifted old generators and \(\eta_i t_i\) generate \(P'\)
by nilpotence of \((X,Y)\). Use them in this order. In the
reduction to \(A\), write the Frobenius rows as \((U_i,*)\),
and lift to rows over \(B\). The top \(a\) relations may be
chosen to reduce to \((\Psi,C)\) for some \(C\), except that
row \(l\) is replaced by \((0,k)\) by reciprocity. Each
reduced relation lifts: its sum lies in \((X,Y)P'\), so
subtract an expression in the same generators with
coefficients in \((X,Y)\).

The additional \(p\) exact tame rows are
\[
 Y(0,E_i)-X(\widetilde U_i,*)\qquad(1\le i\le p),
\]
where \(E_i\) is the \(i\)-th coordinate row. The determinant of this square
array lies in \(\operatorname{Fitt}_0(P')\). Expand in the
last \(p\) rows. The terms with no \(X\) or with at least two copies vanish, since \(Y^p=X^2=0\). For the surviving terms
\[
 X\,(XY^{p-1})=0,\qquad Y\,(XY^{p-1})=0,
\]
so all their other coefficients
may be reduced to \(A\). Expanding now on row \(l\) and the
\(p-1\) unchanged coordinate rows gives, up to one common sign,
\[
       XY^{p-1}\sum_i k_i(U_i\operatorname{adj}\Psi)_l
                      =XY^{p-1}B_\sigma\ne0.
\]
The last inequality uses the monomial basis over \(A\).

\proofstep{Use the bounded nilpotence degree to locate a coefficient.}
Consequently the specialized new determinant, which generates
this Fitting ideal, is nonzero. A monomial of total degree
greater than \(2D+1+p\) vanishes on this substitution:
the old variables are multiples of \(\tau\), and the two others
have exponents at most one and \(p-1\). Thus every lift of
the determinant has a coefficient outside \(I\) in the
asserted degree range. This reasoning applies to any new
square presentation for the same functor.

\proofstep{Use the common degree bound while preserving earlier signature detections.}
The final degree bound must retain the coefficients already detected at earlier signatures.

Repeat this construction for the \(b\) signatures in turn.
The Fitting comparison at each new pair preserves all earlier
detections and the available cofactors within the next bound
\(w_{j+1}\). For each signature determinant \(F_s\), the final common degree bound reads
\[
 \min_{|\alpha|\le w_b} v_\ell([X^\alpha]F_s)\le\nu_\ell.
\]
Here the multi-index ranges over all the variables. For each sufficiently large
\(M\), the same choice of pairs serves all signatures, and the
pairs may depend on \(M\). For \(\nu_\ell=0\)
the residual argument finishes without the reserve pairs.
The number of pairs, the degree and the depth threshold may
depend on the original finite data at the coefficient place;
they are fixed before increasing \(M\). The subsequent
specialization of Section~\ref{subsec:retained-twist} retains
the valuation allowance \(\nu_\ell+\delta\).
\end{proof}


\subsubsection*{The abelian theta coefficients}
\label{subsec:toric-twist}
\hypertarget{source-toric-twist}{}

The Selmer determinant controls the signature contribution to the simultaneous twist. The abelian theta values require a separate argument: changing the toric order changes both the coefficients and their period normalization. Let \(E/F\) be the CM extension of the
selected host, with \(F\) totally real, and let \(\varphi\)
be its chosen CM type. Fix a toric label, its integer
\(s\ge1\), the nonnegative harmonic exponents \(\kappa\),
its base conductor \(\mathfrak f\), and its unitary Hecke
character \(\chi\), trivial on the ideles of \(F\).
Its infinity exponents are \(\kappa_i+s\) at the places
of \(\varphi\) in the unitary normalization. At the base
conductor its local conductor is prescribed and nontrivial.
Enlarge the order conductor to
\[
 \mathfrak c=\mathfrak f\prod_{v\in Q}v
\]
in the real
field, at split places prime to the base conductor and to
the coefficient prime \(\ell\).
Let \(\operatorname{Pic}(\mathfrak c)\) denote the
ideal class group of the resulting order.

For a finite idele \(y\) representing a class, let
\(S_y(z)\) be the continuation of the lattice sum on
\(E\cap y\widehat{\mathcal O}_{\mathfrak c}\), with the
harmonic numerator and shifted exponent of the fixed label,
over nonzero vectors modulo \(\mathcal O_F^\times\).
The toric term is
\[
       T_y=\frac{S_y(0)}
          {\chi_f(y)\,|y|_f^{s}},
\]
where the finite modulus \(|y|_f\) is the one in the
unitary Hecke normalization. In the convergent region the
weights are homogeneous under
multiplication of the lattice by an element of \(E^\times\).
The character at infinity and the norm exponent give
the inverse of this change in the finite prefactor.
The completed order units preserve the lattice and
the finite character is trivial on them. The resulting
normalized term therefore depends on the class. These
identities pass to the continued value.

\paragraph{The denominator after changing the order.}
Choose the fixed CM period lattice with basis
\(1,(1+\sqrt{-d})/2\) at the imaginary quadratic subfield,
and choose a twelfth root of its discriminant. Incorporate the
fixed power of this root and the normalization constant
of the harmonic numerator into a nonzero scalar
\(\Omega_*\). In this normalization
\[
 \beta(T)=(T/\Omega_*)^{12}.
\]

The local class-value estimate for lattices with
\(y_w=1\) at every \(w\mid\ell\) has the following form.
Put \(W=\operatorname{wtM}(2s,\kappa)\), the degree of
the period used in the twelfth power. Let \(a_0\) be the
maximum conductor exponent of \(\mathfrak f\) over \(\ell\).
Let \(c_E\) be the largest denominator exponent at
\(\ell\) of the partial zeta values of \(F\) at \(1-2s\),
divided by \(2^{[F:\mathbb Q]}\) before taking the
denominator. Put \(f_0=|I|=[F:\mathbb Q]\) for the size
of the chosen CM type, \(K_\kappa=\sum_{i\in I}\kappa_i\),
and \(W=2sf_0+2K_\kappa\).
Write \(c_\ell=a_0\) if \(a_0>0\) and
\(c_\ell=v_\ell(D_E)/[E:\mathbb Q]\) otherwise. The
ramified allowance is
\[
 n_\ell =
 \begin{cases}
 0,&c_\ell=0,\\
 2sf_0c_\ell+K_\kappa(2c_\ell+1+
       \mathbf 1_{\ell=2}+\tfrac12\mathbf 1_{\ell=d}),&c_\ell>0 .
 \end{cases}
\]
\begin{proposition}[Uniform denominator of the CM class values]\label{prop:toric-class-denominator}
At every embedding into \(\overline{\Q}_\ell\), for a representative
with \(y_w=1\) at all \(w\mid\ell\), the normalized twelfth
power is algebraic and satisfies
\[
    v_\ell\bigl(\beta(S_y(0))\bigr)\ge-12(c_E+n_\ell+Wa_0).
\]
Auxiliary split places disjoint from \(\ell\) may vary.
\end{proposition}
The proof follows in the local jet calculation below. First we pass
from one lattice to the coefficients of the theta element.

\begin{lemma}[Clearing the abelian theta coefficients]\label{lem:toric-common-denominator}
With the CM type, base conductor, critical weight and coefficient
prime fixed, there is an integer \(c\ge0\) such that every coefficient \(a\) of a normalized abelian theta element, at any allowed auxiliary level and after any allowed finite-order ring class twist, satisfies
\[
 v_\ell(a)\ge-c.
\]
\end{lemma}

\begin{proof}
Every ideal class admits a representative that is one
above \(\ell\). Indeed, use weak approximation in \(E\)
to put the components above \(\ell\) in the local
completed order-unit subgroups and then divide by those
units. On such a representative \(\chi_f(y)\) is a unit
at \(\ell\) in the algebraic normalization. This follows
from the formula for its infinity type on principal
ideles, applied after taking a power to clear the ideal
class, and from the roots of unity on the compact unit
group.

The element of \(E^\times\) used in this formula
has valuations zero above \(\ell\), so every conjugate
factor is a unit. The character and norm factors of such a representative have
\[
 v_\ell(\chi_f(y))=0,\qquad v_\ell(|y|_f)=0,
\]
where \(|y|_f\) denotes the finite norm factor in this normalization.
The same bound thus holds for \(\beta(T_y)\) on all
classes. If a finite order ring class character is added,
its values again have valuation zero. If the separately cleared summands \(a_y\) have nonnegative valuation, the only estimate required for their sum is
\[
 v_\ell\!\left(\sum_y a_y\right)\ge\min_y v_\ell(a_y)\ge0.
\]
Thus the number of classes of the varying order does not change the clearing power.

For the abelian theta module the coefficient at an Artin
label is obtained from the terms \(T_y\) in its fiber,
multiplied by the same nonzero normalization scalar
specified by the toric construction. The fiber map
uses the inverse ideal class. Choose the complex
embedding in the algebraic closure of the row fraction
field as in the toric comparison. Taking twelfth roots gives the individual integral bound
\[
 v_\ell(T_y/\Omega_*)\ge -N\qquad(N\in\mathbb Z,\ N\ge c_E+n_\ell+Wa_0).
\]

One more fixed power clears the product of
\(\Omega_*\) and the normalization scalar in that
algebraic closure. To see this for any algebraic element,
take a monic polynomial over the fraction field and
choose the power large enough that the rescaled monic
polynomial has integral coefficients. Each coefficient
of the theta module is already in the fraction field
of \(O\). Integral closure of the DVR puts its cleared
value in \(O\). For each abelian label with an allowed
\(s\), use one such \(s\) for its fixed coefficient
functions and then take the maximum over the finitely
many labels of this row. This gives the exponent \(N_0\)
in the abelian part of the simultaneous specialization.
\label{cor:toric-twist-1}
\end{proof}

\paragraph{The local jet bound in period coordinates.}
Let \(\omega^{12}\) be the fixed lattice discriminant
and let \(g\) be the defining holomorphic lattice sum
before inserting the powers of the conjugate variables.
The field data here use \(E/\mathbb Q\) abelian,
\(E/F\) unramified at all finite places, and a CM type
containing the chosen image \(i\sqrt d\) of \(\sqrt{-d}\)
at each selected embedding. In particular
\([F:\mathbb Q]\ge2\). If \(F=\mathbb Q\), the
discriminant formula would give \(|D_E|=1\) for an
imaginary quadratic field. The condition on the real
degree applies at weight two as well.

\begin{proof}[Proof of Proposition~\ref{prop:toric-class-denominator}]
\proofstep{Integral idempotents and tangent directions.}
We give the integral comparisons on the lattice
\(E\cap y\widehat{\mathcal O}_{\mathfrak c}\) with
\(y=1\) above \(\ell\). Work over the valuation ring
\(R\) of a finite algebraic coefficient field in
the prescribed embeddings. We enlarge it finitely
to contain the CM embeddings and the good-reduction
model constructed below. The order action extends
to that abelian scheme by the extension property
for homomorphisms. Use \(v_\ell(\ell)=1\) in every
coefficient extension.
Use a generator of the fractional polarization
module locally at \(\ell\). If \(H\) denotes the integral de Rham lattice, perfection of the resulting alternating form means
\[
 H\xrightarrow{\;v\mapsto B(v,-)\;}H^\vee\ \text{is an isomorphism}.
\]

Indeed that local generator gives an isomorphism to
the dual at the coefficient prime, and the integral
de Rham duality comparison gives its perfect form.
The comparison with the ordinary polarization
pairing uses \(2B\), with \(B\) itself the perfect
integral pairing. The generic CM eigenspaces have
rank one, with Hodge subspace given by one
embedding from each conjugate pair. The perfect
pairing is between the opposite eigenspaces.

Let \(D\) be a scalar clearing every primitive CM
idempotent, with
\[
 v_\ell(D)=c_\ell.
\]
For a basis
\((x_j)_j\) of the maximal order and its trace-dual basis
\((x_j^\vee)_j\), the idempotent at an embedding \(\iota_i\)
is \(\sum_j\iota_i(x_j^\vee)x_j\) after coefficient
extension. Multiplication by an element of the different
clears the trace-dual coordinates. Galois uniformity gives
\[
 c_\ell=v_\ell(D_E)/[E:\mathbb Q]
\]
away from the conductor.
At a conductor prime the maximal order is unramified at the
odd prime \(\ell\), so those idempotents are integral there.
Multiplication by \(\ell^{a_0}\) puts their coefficients in
the conductor order. The extra order places are prime to
\(\ell\), so the same local order works in both comparisons.

Let \(e_i,f_i\) be the primitive generators of the two
opposite integral lines in the de Rham lattice and put
\(P_i=B(e_i,f_i)\). They are saturated line lattices.
Perfection of \(B\) supplies an integral \(x\) with
\(B(e_i,x)=1\). Write \(D E_{\bar i}x=s_i f_i\), where
\(E_{\bar i}\) is the idempotent on the opposite line and
\(s_i\in R\). Then
\[
      s_iP_i=B(e_i,D E_{\bar i}x)=D ,
      \qquad 0\le v_\ell(P_i)\le c_\ell .
\]
The equality uses the oppositely paired eigenspaces.
Multiplying the rank-one map \(e_i\mapsto f_i\) by \(D\)
gives an integral operator \(M_i\) killing other eigenlines,
as \(D E_i\) maps into the primitive lattice on the
line of \(e_i\), where \(E_i\) denotes that idempotent.
Each \(M_i\) commutes with the real order and is skew for \(B\).

If \(M=\sum_i t_i M_i\), then \(M^2=0\) and
\[
 B((1+M)u,(1+M)v)
  =B(u,v)+B(Mu,v)+B(u,Mv)+B(Mu,Mv)=B(u,v).
\]
The linear terms cancel by skewness and the last is
\(-B(u,M^2v)=0\). The same expansion applies to the other
polarization forms, since they are obtained by real
multiplication and the \(M_i\) commute with that multiplication.

\proofstep{Comparison with the reference CM periods.}
The primitive integral tangent eigenforms have
complex period scalars \(a_i\). For their comparison
with \(\omega\), let \(A_{\rm ref}\) be the elliptic curve
of the fixed quadratic lattice. Its CM \(j\)-invariant
is integral, so it has good reduction after finite
extension at \(\ell\). Choose an integral generator of
its tangent line. Its analytic coordinate \(t\) in the
given lattice has \(t^{12}\omega^{12}\) equal to the
discriminant of a good Weierstrass model. This is the
weight-twelve change of Weierstrass invariants, using
the full Weierstrass equation at two and three.
Thus \((t\omega)^{-1}\) is algebraic and integral at
the selected place.

Let
\[
 L_y=E\cap y\widehat{\mathcal O}_{\mathfrak c}.
\]

Since the local lattice at each place above \(\ell\)
is the conductor order, \(\ell^{a_0}\mathcal O_E\)
is contained in it locally. At the remaining places
one integer \(m_0\) prime to \(\ell\) clears the
necessary lattice index. Thus
\[
 n_0\mathcal O_E\subset L_y
\]
with
\(n_0=\ell^{a_0}m_0\). Choose \(\theta\in\mathcal O_E\)
separating the complex embeddings and take the
matrix with columns \(n_0(\varphi_i(\theta^j))_i\),
\(0\le j<f_0\).

For \(V=(n_0\varphi_i(\theta^j))_{i,j}\), the embeddings
give distinct values \(\varphi_i(\theta)\), so
\[
 \det V=n_0^{f_0}\prod_{i<i'}
               (\varphi_{i'}(\theta)-\varphi_i(\theta))\ne0 .
\]
The reference lattice consists of simultaneous embedded
quadratic integers. Thus column \(j\) maps such an
integer \(a\) to the lattice element \(n_0a\theta^j\)
of \(n_0\mathcal O_E\subseteq L_y\). The matrix defines an
isogeny from \(A_{\rm ref}^{f_0}\) to the CM variety.
After finite extension its finite generic kernel has
a finite flat schematic closure in the torsion group
of the good source. Quotient by this closure gives
the good target model and an integral isogeny.
The integral tangent vector on the factor in column
\(j=0\) has image with analytic coordinates
\(w_i^0=n_0t\) for all \(i\). Consequently
\[
 v_\ell((w_i^0\omega)^{-1})\ge-a_0.
\]

Evaluating the integral eigenform on this vector
gives \(a_iw_i^0\in R\). This implies the first
of the two bounds below. For the determinant,
the same cleared CM idempotents on the tangent
dual put \(D\) times the full lattice in the span
of the primitive eigenforms. Their wedge is
\(u_0\) times a determinant generator, where
\(0\le v_\ell(u_0)\le f_0c_\ell\) by Smith normal
form. The complex determinant scalar is
\[
 C=(\prod_i a_i)/u_0.
\]


For the lowering coordinates let \(w_i\) be the
complex tangent coordinate of \(f_i\) and choose
the integral unit \(r_i\) with \(P_i=r_i a_i w_i\).
This equality follows by using \(B(e_i,-)\) as
the primitive covector on the tangent quotient.
The Riemann form in the polarized uniformization
has real scale \(h_i\) on this coordinate.
If \(\psi\) is the complex polarization pairing
and \(\xi\) is the imaginary generator of the
polarization module locally at \(\ell\), the
comparison of alternating forms gives
\[
  2i\varphi_i(\xi)\psi(e_i,f_i)=h_i P_i .
\]
The local ideal of polarizations is dual to the
lattice under the trace pairing.

At the chosen
places, \(y=1\) and the extra level is a unit.
The local generator consequently satisfies
\(v_\ell(\varphi_i(\xi))\le0\).
Equivalently one can choose a local anti-invariant
inverse-different generator \(\xi'\) of nonpositive
valuation and use
\(m\xi'=a\xi\) with \(m\) prime to \(\ell\)
and \(a\in\mathcal O_F\) to compare the two generators.
The scale of the lowering direction and its auxiliary
period are
\[
  A_i=\pi h_i\frac{D}{\psi(e_i,f_i)}w_i^2,\qquad
  b_i=\frac{a_i A_i}{D}.
\]
In particular \(b_i/(a_i A_i)=D^{-1}\). Writing
\(\delta=\sqrt{-d}\) with \(\varphi_i(\delta)=i\sqrt d\),
the other period quotient is
\[
   (a_i b_i\sqrt d/(2\pi))^{-1}
       =\frac{r_i^2}{P_i\,\varphi_i(\xi\delta)} .
\]
All of these right-hand sides are algebraic.
Since \(v_\ell(\varphi_i(\delta))=\frac12\mathbf1_{\ell=d}\),
we obtain
\[
 v_\ell(a_i/\omega)\ge-a_0,\qquad
 v_\ell\bigl((a_i b_i\sqrt d/(2\pi))^{-1}\bigr)
     \ge-c_\ell-\tfrac12\mathbf1_{\ell=d}.
\]
The determinant and inverse-idempotent bounds just
proved supply the other two valuations.

\proofstep{The divided-power deformation.}
Use the divided-power thickening of the chosen abelian scheme. In horizontal coordinates the prescribed variation of its Hodge submodule is
\[
 \operatorname{Fil}^1_{\boldsymbol t}=(1+\sum_i t_iM_i)\operatorname{Fil}^1_{0},
\]
with the original filtration extended to the jet ring.
The real action and polarizations lift with the
filtration because the indicated operators commute
with that action and preserve the forms.
Here use the divided-power jet lifting theorem for the integral level.
Its base is a characteristic-zero local domain and its
thickening is the truncated divided-power jet ring,
with split augmentation and nilpotent divided-power
kernel. The input endomorphisms are admissible for the
real action, filtration and alternating forms, as are the
lowering operators above. The theorem gives a polarized
abelian scheme with real multiplication over that jet
ring, together with the pullback to the augmentation.
Under the horizontal de Rham identification, its
Hodge submodule is the span of the prescribed jet lifts
of the original filtration.

If \(s_{\det}\) is the determinant lift and \(s_0\) the given generator at the origin, its normalization in the horizontal identification is
\[
 s_{\det}\bmod(t_1,\ldots,t_d)=s_0,
\]
where \(d\) is the number of deformation parameters.
Complex base change identifies the integral jet and its lifted filtration with the analytic deformation used to calculate the periods.
On the reductions modulo powers of \(\ell\) the
underlying deformation comparison has \(\ell\) nilpotent,
as in Grothendieck--Messing theory.
\cite{messing}\label{cor:toric-twist-2}

Transport the periods horizontally in that family.
In the \(i\)-th real factor take complex combinations
of the parallel cycles with holomorphic coordinate one
and with zero holomorphic coordinate and antiholomorphic
coordinate one, respectively. The latter cycle
represents the Hodge direction at the origin.
After the lowering it lies in the annihilator of
the lifted holomorphic forms.

Let \(P_i^{\rm hol}\) be the nonzero period on the first cycle and \(\varepsilon\) the period pairing sign. Evaluation on the second cycle gives
\[
          P_i^{\rm hol}(\zeta_i-\varepsilon A_i t_i)=0 .
\]
Its constant term is nonzero, so \(P_i^{\rm hol}\) is a unit
in the complex jet ring. Hence \(\zeta_i=\varepsilon A_i t_i\) there.
The determinant scale on the horizontal frame is
\(C\). Indeed the parallel cycles with holomorphic
coordinates the standard basis are on the opposite
CM eigenlines, where the lowering operator is zero.
Evaluating the lifted determinant on their wedge
gives exactly the value at the origin.

\proofstep{Extracting the coefficient of the prescribed jet.}
Put \(g_\kappa=\partial^\kappa g(0)\) and
\[
 J_\kappa=g_\kappa\prod_i((b_i/a_i)^{\kappa_i}/a_i^{2s}).
\]

Let \(m=\ell^{c_E}\) and let \(R\) be the local ring
of algebraic coefficients for the integral CM model,
contained in the chosen valuation ring. It contains
the rational \(\ell\)-integers. The weight-\(2s\) Hilbert
Eisenstein form on each polarized component has
integral nonconstant Fourier coefficients in the
algebraic normalization. Its constant coefficient
is the corresponding partial zeta value divided
by \(2^{f_0}\). The partial zeta rationality theorem
applies since \(F\) is totally real and \(s\ge1\).
The integral model of the Hilbert moduli problem
and the \(q\)-expansion principle on all components
then gives the form multiplied by \(m\) over \(R\).
For weight two the Fourier statement uses the value
at zero of the continued family in real degree at
least two. \cite{deligne-pappas,serre-tate}\label{cor:toric-twist-3}

Use the divided-power algebra in the deformation parameters truncated at total degree
\[
 |\kappa|=\sum_i\kappa_i.
\]

The admitted de Rham operators of the integral CM
model determine its deformation over this algebra.
The lifted determinant generator reads the section
as a scalar in the divided-power basis, with
coefficient \(x\in R\) at the multi-index \(\kappa\).
After base change to \(\mathbb C\), the lowering
direction in the \(i\)-th coordinate has scale
\(\varepsilon A_i\), where \(\varepsilon\in\{1,-1\}\).
The determinant scale is \(C\). Comparing the divided
power coefficient after base change gives
\[
       C^{2s}x=m g_\kappa\prod_i(\varepsilon A_i)^{\kappa_i}.
\]

In particular there is no factorial in this coefficient
identity. The integral model comparison gives
\(v_\ell(C/\prod_i a_i)\ge-f_0c_\ell\) and
\(v_\ell(b_i/(a_i A_i))\ge-c_\ell\). Thus
\[
       J_\kappa=x m^{-1}
         \left(\frac{C}{\prod_i a_i}\right)^{2s}
         \prod_i\left(\frac{b_i}{\varepsilon a_i A_i}\right)^{\kappa_i}
\]
has valuation at least \(-c_E-2sf_0c_\ell-K_\kappa c_\ell\).
At zero depth this is \(-c_E\).
At positive depth the sum of the factorial valuations satisfies
\[
                  \sum_i v_\ell(\kappa_i!)\le K_\kappa/(\ell-1).
\]
Subtract this from the bound for \(v_\ell(J_\kappa)\) to obtain
\[
 v_\ell\!\left(J_\kappa/\prod_i\kappa_i!\right)
    \ge -c_E-2sf_0c_\ell
             -K_\kappa(c_\ell+1+\mathbf1_{\ell=2}) .
\]
\proofstep{Returning to the harmonic lattice sum.}
The normalized coordinate derivative for the harmonic
lattice sum is \((-1)^{K_\kappa}c_{2s}
\prod_i(2s)_{\kappa_i} S_y(0)\), where \(c_{2s}\) is
the defining lattice constant and \((2s)_a\) is the
rising factorial. This identity follows by differentiating
the continued defining sum. In the convergence region each
derivative increases the relevant denominator exponent by
one, supplying its successive factor, and continuity in
the exponent gives the formula at the evaluation point.
The definition of \(\Omega_*\) absorbs the same constant
and rising factorials. Multiplying \(J_\kappa\) by
\[
   \prod_i (a_i b_i\sqrt d/(2\pi))^{-\kappa_i}
                  (a_i/\omega)^{2s+2\kappa_i}
\]
leaves, up to the sign, \(S_y(0)/\Omega_*\).

In the positive-depth case use the divided-power jet
and also multiply by \(\prod_i\kappa_i!\). These
factorials are integral at every prime. The two period
factors cost \(K_\kappa(c_\ell+\tfrac12\mathbf1_{\ell=d})\)
and \(Wa_0\), respectively. At zero depth \(\ell\ne d\)
since the auxiliary field contains the imaginary quadratic
subfield of discriminant \(-d\). The same multiplication
then costs only \(Wa_0\). Adding these contributions
and taking the twelfth power proves the stated local
lattice estimate from the normalized jet inequalities.
The coordinate and jet comparisons use the fixed CM
type and base conductor. The auxiliary order and the
class label may vary with these data and the coefficient
prime fixed. \label{cor:toric-twist-4}
\end{proof}

\paragraph{Distribution at an added split place.}
Write \(q\) for the size of the residue field of a finite place \(v\) of the totally real field
that splits in the CM extension, and \(P,P^c\) for the places above it.
The local completed order at depth one consists of
pairs with equal residues. For a character unramified
at \(v\), let \(u\) be its value on a uniformizer at
\(P\), so its value at \(P^c\) is \(u^{-1}\). In the
absolutely convergent half-plane write \(X=q^{-z}\),
where \(z\) is the complex argument of the Hecke series.

\begin{lemma}[Distribution at a split finite place]\label{lem:toric-split-distribution}
With inverse Artin reciprocity on the ideal classes, removal of an
auxiliary pair multiplies a value unramified at the pair by
\[
       (1-uX)(1-u^{-1}X)+(q-1)X^2,\qquad X=q^{-z},
\]
where \(u\) is the corresponding twisted uniformizer value on
one factor. This relation holds on the remaining group-ring
coordinates with the same period normalization.
\end{lemma}

\begin{proof}
The two unit coordinates with equal residues give one
class modulo the order units. Pairs with both
coordinates divisible by a uniformizer give \(q-1\)
unit cosets for each pair of positive valuations.
There are no pairs with exactly one positive valuation.
The local sum is consequently
\[
     1+(q-1)\frac{X^2}{(1-uX)(1-u^{-1}X)} .
\]
Relative to the original two Euler factors this is the
polynomial
\[
          (1-uX)(1-u^{-1}X)+(q-1)X^2 .
\]
If the tested character is ramified at this place,
only the simultaneous unit class contributes. The other
terms sum to zero because the character on the ratio of
residue units is nontrivial. The usual local factor of
the Hecke series there is one as well. The exponent-one
order is essential to the unit coset count.

\proofstep{The unit group of the order.}
To justify applying this computation to the lattice sum,
we identify the global units of the order precisely.
The base conductor is nonzero and proper and has an odd
prime divisor of residual characteristic \(p\) unramified
in \(E\) over \(\mathbb Q\). Let \(\mathfrak P\) be a prime of
\(E\) above that divisor. If \(u\in\mathcal O_{\mathfrak c}^\times\),
then \(u,\bar u\) are algebraic units and
\[
 u-\bar u\in\mathfrak c\mathcal O_E\subseteq\mathfrak P,
\]
since the image of the order modulo the conductor is real.
All complex conjugates of \(u/\bar u\) have absolute value one.
An algebraic integer with that property is a root of unity,
as the monic polynomials of all its powers have bounded integral
coefficients. At \(\mathfrak P\) the root lies in the principal unit
group. Here the normalized local valuation satisfies \(v(p)=1\).

If \(v(t)\ge1\) and \(t\ne0\), then \(v(pt)=1+v(t)\). The
intermediate terms \(\binom pj t^j\), \(2\le j<p\), have valuation
at least \(1+2v(t)\), and the final term has valuation \(pv(t)>1+v(t)\)
because \(p\) is odd. Thus
\[
 v((1+t)^p-1)=v(pt)=1+v(t).
\] For an exponent \(a\) prime to \(p\) the linear term
gives \(v((1+t)^a-1)=v(t)\). For any positive exponent
\(m=p^j a\) with \((a,p)=1\), iteration therefore gives
\[
                  v((1+t)^m-1)=v(t)+j<\infty .
\]
This is incompatible with \(m\)-torsion.
Thus the principal unit group has no nonidentity torsion, and
\(u=\bar u\). Both \(u\) and its inverse belong to \(\mathcal O_F\),
so
\[
 \mathcal O_{\mathfrak c}^\times=\mathcal O_F^\times.
\]


With this equality,
in the convergence half-plane group the nonzero lattice
vectors by the corresponding integral finite ideles
modulo the completed order units. Dividing each
idele into a class representative and an element of
\(E^\times\) identifies the fiber with the lattice vectors
modulo the global order units. Homogeneity identifies
its summand with the Hecke weight. The local
order conditions and the multiplicativity of the weight
give the Euler product. At a place of the base conductor
the character has the prescribed exact conductor.
In its nonunit layers the sum over the last available
residue unit quotient vanishes by that conductor, leaving
the primitive local factor. At all other unchanged
places the full order gives the usual geometric series.
The equality of the continued sums follows from the
identity theorem, and at the desired value one
substitutes
\[
 X=q^{-s}.
\]


\proofstep{The sign in the inverse-class map.}
Pair removal will compare coefficients at a split prime through their ray classes. The transfer must therefore use the same convention for inversion as the distribution relation.

For removal of one pair of auxiliary places, evaluate
the abelian group-ring data on any character of the
remaining Artin labels. Such a character is unramified
at the removed place. Summing the removed coordinate
amounts to forgetting it in the inverse-class map.
A uniformizer at either place above \(v\) acts on the
remaining labels through its inverse Artin class.
Define \(Y\) to be the inverse of the image of \(\chi(P)\)
in the coefficient field.

With this reciprocity convention the unramified twisting factor at \(P\) is
the value of the remaining character on its Frobenius.
Thus evaluation at a remaining character \(\psi\) gives
\[
 u_\psi=Y^{-1}\psi(\operatorname{Frob}_P),\qquad
 (1-u_\psi X)(1-u_\psi^{-1}X)+(q-1)X^2 .
\]
The second local root is \(u_\psi^{-1}\).
The polynomial gives the same
multiplicative relation at every tested character, so
finite Fourier inversion proves the group-ring
distribution. With all auxiliary places removed, the
sum of the base coefficients is a nonzero scalar times
the untwisted Hecke \(L(s,\chi)\). The character is unitary
and is ramified at the base conductor, so the classical
nonvanishing theorem for unitary Hecke \(L\)-functions
on \(\RePart z\ge1\), away from the pole of a norm
character, applies. Both the coefficient bound and this
distribution refer to the same inverse-class map and
the same period. \label{cor:toric-twist-5}
\end{proof}

\subsubsection*{The theta polynomials from pair removal}
\hypertarget{source-twist-euler}{}

The toric distribution relation describes the effect of adding a split prime. To use it in the simultaneous choice of Section~\ref{subsec:retained-twist}, encode pair removal by specialization in the group ring before evaluating any character.
The first and second halves use the multiplicative group-ring coordinates
\[
 U_i=1+X_i,\qquad V_i=1+Y_i,\qquad U_i^L=V_i^L=1,\qquad L=\ell^M.
\]

The group ring of the coordinates is the quotient by
\(U_i^L-1,V_i^L-1\). In an integral coefficient polynomial one
can use the exponent \(L-1\) for the inverse of a group
coordinate. Equality at all \(L\)-torsion tuples after setting
\(V_i=1\) gives equality modulo \((Y_i,U_i^L-1)_i\) over the
original integers. Indeed division in each variable by the
monic polynomial leaves a remainder of degree less than \(L\)
in each variable. In a characteristic-zero splitting field
successive univariate interpolation kills that remainder.
This proof involves no division by a group order in the
coefficient ring.

\begin{lemma}[Removal of pairs for abelian theta values]\label{lem:abelian-pair-removal}
For an abelian label at an admissible critical integer \(s\ge0\),
use the uniform denominator \(\ell^{N_0}\) and the two halves of the
chosen auxiliary pairs. The cleared full and first-half theta values
have integral polynomial representatives \(f_M,g_M\) such that
\[
 f_M\equiv B_Mg_M\pmod{(Y_i,U_i^L-1)_i},\qquad
 v(g_M(0))=N_0+v(\theta_\lambda(1)).
\]
Each second-half factor of \(B_M\) reduces to a polynomial in
\(X_i\) of order zero or two at zero.
\end{lemma}

\begin{proof}
First consider an abelian label at an admissible critical
integer \(s\ge0\). The input denominator \(\ell^{N_0}\)
makes each coefficient of its ray-class theta element integral,
uniformly in the auxiliary level. Push this element to the
coordinate group ring by the chosen characters and choose a
polynomial lift \(f_M\). More explicitly, lift each ray class
to an Artin element, use its character coordinates as exponents
in \([0,L-1]\), and sum the resulting monomials in \(U_i,V_i\)
with the cleared coefficients. The characters factor through
the ray group. Evaluating gives exactly
\(\ell^{N_0}\theta_\lambda(\psi)\). Use \(g_M\) for the
same construction at the union of the first half of the pairs,
using just the first-half characters.

\proofstep{The second-half factors.}
The pair-removal identity of the theta input at a pair of
norm \(q\) and at a character unramified on that pair gives
the factor
\[
 E_q(A)=(1-Aq^{-s})(1-A^{-1}q^{-s})+(q-1)q^{-2s},
 \qquad A=\lambda(\mathrm{Frob})\psi(\mathrm{Frob}).
\]
All these coefficients are integral at \(\ell\) because
\(q\equiv1\bmod 2L\). Remove the second half of the pairs
with all \(V_i=1\). At pair \(i\) the Frobenius character
value is exactly \(U_i\) on evaluating, by the identity
matrix condition on the first-half characters. Thus \(B_M\)
is the product over the second half of the integral polynomial
representatives of \(E_q(\lambda(\mathrm{Frob})U_i)\).
Successive removal and the interpolation argument above give
\[
 f_M\equiv B_Mg_M\bmod(Y_i,U_i^L-1)_i.
\]


\proofstep{The first-half unit test.}
To compute \(g_M(0)\), remove the first half at the trivial
character. Choose the open normal subgroup inside the kernel
of the residue character \(\bar\lambda\). On the prescribed
tame element \(z\), \(\bar\lambda(z)\ne1\). Indeed at the
distinguished place of norm \(q_0\) the character label there
agrees with \(\delta^{\pm1}\) after taking a \(q_0\)-power
iterate (power \(q_0^n\) for some \(n\)). The finite
quotient values of \(z\) have order prime to \(q_0\ell\).
Writing a bar for reduction at \(\ell\), the choice of \(z\) gives
\[
 (\overline{\delta(z)}-1)(\overline{\delta(z)}+1)=\overline{\delta(z)}^2-1\ne0.
\]
These are the tests on \(z\): take the
prime-to-\(q_0\ell\) part of a generator at the distinguished
order-15 place. Its value retains the prime-to-\(\ell\)
part of a primitive fifteenth root, of order divisible by
three or five and prime to \(q_0\). It follows that
\(\bar\lambda(z)^{q_0^n}\ne1\).

By the first Frobenius test, the factors at the first-half
places reduce at the trivial character to
\((1-\bar\lambda(z))(1-\bar\lambda(z)^{-1})\), a unit.
Consequently
\[
 v(g_M(0))=N_0+v(\theta_\lambda(1)).
\]
Here the value on the right is at empty auxiliary level.
For a second-half factor write \(x=\bar\lambda(\mathrm{Frob})\).
Its reduction is \((1-xU_i)(1-x^{-1}U_i^{L-1})\).
For \(x\ne1\) it has unit constant term. For \(x=1\)
it equals \(X_i^2\sum_{j=0}^{L-2}(1+X_i)^j\), whose last
factor has constant \(L-1\ne0\) in the residue field.
The denominator bound and the pair relation prove the abelian coefficient condition with \(n_*=2\). The distinction between the
two halves of the primes is essential to the constant
coefficient estimate. \label{cor:twist-euler-1}
\end{proof}

\begin{proposition}[Determinant and signed index for the row theta values]\label{prop:row-pair-removal}
Let \(B\) be the prescribed row lattice of rank \(n\), with its
\(n\) independent projected vectors at empty level and the uniform
clearing scalar \(x\). Write \(C_Q\) for their cleared coordinate
matrix at level \(Q\), \(Q_1\) for the first half of the auxiliary
pairs, and \(b=\operatorname{ord}_O(\det C_0)\). Then
\(f_M=\det C_Q\) and \(g_M=\det C_{Q_1}\) satisfy the same
pair-removal congruence as in Lemma~\ref{lem:abelian-pair-removal},
with the second-half determinant factors displayed below, of residual
order at most \(rn\) in each variable. For any lattice containing the specialized columns, let \(\Delta\) denote its excess signed index over the original column index. Then
\[
 \operatorname{ord}_O g_M(0)=b,\qquad f_M(u-1)\ne0\ \Longrightarrow\ \Delta\le v(f_M(u-1))-b/e_O.
\]
\end{proposition}

\begin{proof}
\proofstep{Clearing the columns on both sides.}
For a row, let \(B\) be the prescribed free lattice of rank
\(n\), \(A\) the commuting central operator of the row
calculation, and \(\theta_1,\ldots,\theta_n\) the selected
projected vectors at empty level. They are independent and
span the rational space. Choose a basis of \(B\).
There is a single denominator
\[
 x=\ell^c\ne0
\]
for the
universal lifts of this finite list, valid at every \(M,Q\)
disjoint from their levels. Projection takes the lattice
containing these cleared lifts into \(B\).

Take their basis
coordinates, with the universal characters still present,
and push the group coordinates to the polynomial presentation
above. This gives an integral square matrix \(C_Q\) such
that at every torsion character its columns are the coordinates
of \(x\) times the projected specializations. At empty level
it gives the constant matrix \(C_0\) of the \(x\theta_j\).
In particular \(\det C_0\ne0\). With \(Q_1\) the first-half set, the two determinant polynomials and the base valuation are
\[
 f_M=\det C_Q,\qquad g_M=\det C_{Q_1},\qquad b=\operatorname{ord}_O(\det C_0).
\]

This choice cancels the denominator in the signed index.
Indeed, if \(\det C_Q(u-1)\ne0\), those specialized columns
before clearing span a lattice of normalized relative index
\(v(\det C_Q(u-1))-n v(x)\) in \(B\). Every lattice in
the same rational space that contains them has relative
index at most this, by elementary divisors. Subtract the
corresponding index for the original \(\theta_j\) to get
exactly \(v(f_M(u-1))-b/e_O\). The formula also holds for
fractional lattices extending beyond \(B\).

\proofstep{The primitive local function.}
For a place \(P\) of norm \(q\), with integral Hecke matrices
\(T_{a,P}\) in the basis of \(B\), put
\[
 H_P(Y)=I+\sum_{a=1}^r(-1)^a q^{a(a-1)/2}Y^aT_{a,P}.
\]
Work in the split local model at the chosen member \(P\).
To distinguish the matrix variable from the rank of the row
lattice, denote its number of columns by \(n_{\rm loc}\).
It has \(r\) rows. The split base frames
identify the vacuum with the characteristic function of
integral matrices. Sum over the character coordinate at a
deleted pair. If the reduced rows span an \(r\)-plane, its
complements are the graphs of maps from a fixed complement
to that plane. For a matrix \(z\) the normalized count is
\[
 q^{-r(n_{\rm loc}-r)}
 \#\{\text{complements of dimension } n_{\rm loc}-r\}
 =
 \begin{cases}1,&\operatorname{rk}\bar z=r,\\0,&\operatorname{rk}\bar z<r.\end{cases}
\]
For smaller rank the dimension condition prevents complementarity.
This is the primitive local function on the integral vacuum.

Let \(z\) be integral and \(c=\dim\ker(\bar z^{\mathsf T})\).
The right cosets of the index-\(a\) minuscule Hecke double
coset parameterize the lattices \(D\) between
\(\varpi O_P^r\) and \(O_P^r\) of residual codimension \(a\):
send a representative \(\gamma\) to \(\gamma O_P^r\), and
use elementary divisors for surjectivity. The condition
\(\gamma^{-1}z\) integral means that each column is in
\(D\). Equivalently the annihilator of \(D/\varpi O_P^r\)
is an \(a\)-plane in \(\ker(\bar z^{\mathsf T})\).
Thus there are \(\genfrac{[}{]}{0pt}{}{c}{a}_q\) translates.
The subspace recursion obtained by intersection with a hyperplane
gives the polynomial identity
\[
 \sum_{a=0}^c (-1)^a q^{a(a-1)/2}
    \genfrac{[}{]}{0pt}{}{c}{a}_q Z^a
   =\prod_{j=0}^{c-1}(1-q^jZ).
\]
At \(Z=1\) the product is one for \(c=0\) and zero for \(c>0\).
Hence the alternating Hecke sum is precisely the primitive
indicator. Off the integral vacuum every term vanishes since
\(D\subset O_P^r\).

\proofstep{The normalized Hecke translates.}
In the split Weil action the right translate of the vacuum
by an index-\(a\) representative has the indicator just
counted with modulus and splitting-character multiplier
\[
 q^{a n_{\rm loc}/2}\mu(\varpi)^a.
\]
Here \(n_{\rm loc}\) is even. This is
the change-of-variables factor \(|\det\gamma|^{-n_{\rm loc}/2}\)
on the matrix argument and the unramified splitting
character on its determinant. Dividing the sum of translates
by this same multiplier is the normalized \(T_{a,P}\) on the
row.

The other local components carry the shorter coordinate
datum: the coordinates agree by restriction, and replacing
the deleted test by the constant one on its integral support
restores the vacuum. The identity of local functions passes
to theta series by linearity and right translation. If \(b\) is the remaining Artin label, the determinant of a representative in degree \(a\) changes it by
\[
 b\longmapsto b+a\delta,
\]
where \(\delta\) is the Artin-coordinate vector of the uniformizer at \(P\). This follows
by taking determinants of the representatives, all of which
are \(\varpi^a\) times a local unit and all remaining
characters are unramified at \(P\). With the inverse Artin map of Lemma~\ref{lem:toric-split-distribution} \(\delta\) is minus the vector of
their arithmetic Frobenius values.

For a character descending to the shorter set, write \(Y\) for its Frobenius value. Translating the index back by \(-a\delta\) supplies the factor
\[
 \chi(-a\delta)=\chi(-\delta)^a=Y^a.
\]
Injectivity of
the row realization transfers this equality to the prescribed
Hecke matrices. This gives \(H_P(Y)\) at the signed character
value on the representative selected by the row.

Projection commutes with the
Hecke operators. Each of the \(n\) columns therefore changes
by the very same matrix. Applying the relation successively
and taking determinants gives
\[
 B_M=\prod_{i=1}^{R_*} \det H_{P'_i}(U_i^{c_i}),
       \qquad c_i\in\{1,L-1\},
\]
with \(P'_i\) in the second half. The identity-matrix
Frobenius test gives \(U_i\) on the specified representative,
conjugating the place inverts the value by anticyclotomicity,
and the sign of the row can invert it once more. This proves
the formula at torsion tuples with trivial second half and
hence the required ideal congruence by interpolation.

Over the residue field
\[
 h(Y)=\det H_{P'_i}(Y)
\]
satisfies
\[
 h(0)=1,\qquad\deg h\le rn.
\] The order of
\(h((1+X)^{c_i})\) at \(X=0\) is the multiplicity of the
root \(Y=1\) of \(h\), including multiplicity zero. Indeed
\((1+X)^{c_i}-1\) has nonzero linear coefficient there.
Each factor is therefore \(X_i^{n_i}\) times a polynomial
with unit constant, \(n_i\le rn\), uniformly in \(M\).

\proofstep{Units at the first half of the pairs.}
The factors from the second set of auxiliary primes have already been identified. The determinant coefficient also depends on what survives when the coordinates from the first set are specialized.

It remains to check the first-half determinant units at the
trivial character. For each signature \(s\), write \(H_{P,s}\) for its unramified Hecke polynomial with the prescribed character \(\eta_s\). Evaluation at one gives
\[
 H_{P,s}(1)=\det(1-\eta_s\rho_s(\mathrm{Frob}^{-1})).
\]

At the distinguished place the twist \(\eta_s\rho_s\)
acts by the scalar \(\delta\) or \(\delta^{-1}\) on inertia.
Intersect the kernels of all the reduced twisted matrices and
the conjugate of this intersection. This is an open normal
subgroup because there are finitely many continuous residual
actions. At either place of a first-half pair the inverse residual Frobenius is scalar. Its two possible values are
\[
 \overline{\delta(z)},\qquad \overline{\delta(z)}^{-1},
\]
since conjugation transports inertia and inverts \(\delta\). Thus the displayed determinant is a
unit for every signature.

Suppose \(\bar H_P(1)\) is singular over the algebraically
closed residue field. Its kernel is nonzero and stable under
all commuting Hecke operators. Restrict successively to a
nonzero eigenspace of a nonscalar operator. This gives a common
eigenvector \(w\) in the kernel. Its residual signature \(s\)
therefore satisfies
\[
       0=\bar H_P(1)w=\overline{H_{P,s}(1)}\,w .
\]
The first-half test just proved makes \(H_{P,s}(1)\) a unit,
a contradiction. Thus \(\det H_P(1)\) is a unit on the row.
Deleting the first-half pairs at the trivial character gives
\[
                   g_M(0)=\det C_0\prod_P\det H_P(1).
\]
This proves the coefficient and index bounds for every allowed
lattice. The following comparison identifies the integral lattice,
residual signature and the common geometric denominator.
\label{cor:twist-euler-2}
\end{proof}

\subsubsection*{Uniform integrality on the ray covers}
\label{subsec:theta-cover-integrality}
\hypertarget{source-theta-cover-integrality}{}

The polynomial argument requires an integral bound on the row coefficients as the squarefree ray level varies. We obtain it first on the cover, using its prescribed frames, and then descend to the level on which the determinant is evaluated.

\begin{paramdefinition}[Admitted levels and coefficient place]\label{hyp:ray-cover-integrality}
We prove the coefficient bound on the ray covers used in the
row determinant. Fix the vacuum pair, the base weight and
coefficient map, the pinned orthonormal frame, and the finite
source translate \(s\). The base level and the varying ray level are
\[
 M_0=g_0g_1,\qquad M=M_0d,
\]
where \(d\) is odd,
squarefree, prime to \(M_0\ell\) and to the conductor and
discriminant exclusions of the host. We use the uniformized
cover with its compatible moduli and adelic frames. Its
coefficient field contains the ray field of level \(M\), the
harmonic coefficients under consideration, and the normalized
values of the splitting characters. A cover with these levels and frames is called admissible for the row at the chosen coefficient place.
\end{paramdefinition}

At the fixed place over \(\ell\), the exponent of \(\ell\) in the level is unchanged. The new marking is a basis of prime-to-\(\ell\) torsion. This torsion is finite étale on the good model, and the compatible bases form a finite étale cover there. Pullback to that cover therefore preserves the good-model condition used in the coefficient bound.

\begin{auxlemma}[Integral finite Fourier expansion]\label{lem:cover-vacuum-fourier}
For a level \(M=M_0d\) of Definition~\ref{hyp:ray-cover-integrality},
let \(\Phi\) be an integral level function on the selected half
lattice, periodic modulo \(M\). If \(k\) generators suffice for
the reduction of the tensor lattice, the vacuum \(v\) and its Heisenberg translates \(W(y)v\) satisfy
\[
 M^k f_\Phi\in\sum_y O\,W(y)v,
\]
with \(O\) the coefficient valuation ring.
The clearing cost at \(\ell\) is at most \(k v_\ell(M_0)\).
\end{auxlemma}
\begin{proof}
Let \(\Lambda\) be the adelic tensor lattice of the vacuum
and write \(\Lambda_e\) and \(\Lambda_{\bar e}\) for its
two projected half lattices. Here \(e\) is the sum of the
idempotents at the places of the level selected by the finite
test vector. A level function \(\Phi\) is supported on
\(\Lambda_{\bar e}\) and is periodic modulo
\(M\Lambda_{\bar e}\). Its values belong to the chosen
valuation ring. Fix a list of \(k\) generators for the
reduction of \(\Lambda\) modulo any positive integer, using
generators for the fixed source and target lattices.
Project these to \(\Lambda_e\) and divide them by \(M\), giving
\(y_i\). The character of \(t\) paired with
\[
 y(\alpha)=\sum_i\alpha_i y_i,
\]
for
\(\alpha\in(\mathbb Z/M\mathbb Z)^k\), is a product of
\(M\)-th roots of unity.

By vacuum self-duality the two half lattices pair perfectly.
For \(t\in\Lambda_{\bar e}\), put \(z_i=\psi(\langle t,y_i\rangle)\).
Then \(z_i^M=1\) and
\[
 \sum_{\alpha}\psi(\langle t,y(\alpha)\rangle)
 =\prod_{i=1}^{k}\sum_{a=0}^{M-1}z_i^a
 =M^k\mathbf 1_{t\in M\Lambda_{\bar e}} .
\]
Indeed, a factor vanishes unless \(z_i=1\); all the \(z_i\) equal one
precisely when \(t\) pairs trivially with the generators. This argument
allows repetitions among the generators. Let \(W(y)\) be
Heisenberg translation by \(y\). Fourier inversion
writes the finite theta vector \(f_\Phi\) as
\[
 M^k f_\Phi
   =\sum_{\alpha}
       \left(\sum_{t\in\Lambda_{\bar e}/M\Lambda_{\bar e}}
             \Phi(t)\,\psi(-\langle t,y(\alpha)\rangle)\right)
             W(y(\alpha))f_0 ,
\]
The coefficients are integral, including at a prime dividing
\(M\). Since \(d\) is prime to \(\ell\), the clearing
exponent in this step is \(k v_\ell(M_0)\).
\end{proof}

\paragraph{Moving to the geometric half lattice.}
The moduli level selects a half lattice at each split
place. Switch the halves that differ from this choice by
the central finite Weil operators. A switch exchanges a
place with its conjugate in the idempotent and carries
each translated vacuum to a translated vacuum in that half,
times the splitting-character scalar. At primes of \(d\)
these switches have identity component at \(\ell\). The
scalar there is integral at \(\ell\) because the local norm factors
are powers of rational residue cardinalities prime to \(\ell\),
and the character uses the norm-normalized values in the
cover field. At primes of \(M_0\) there are finitely many
subsets of places to switch. If \(a_J\) denotes the scalar for such a subset \(J\), a common nonnegative clearing exponent satisfies
\[
 c_{\rm flip}\ge\max_J\{-v_\ell(a_J),0\},
\]
and the associated source
translations form a fixed finite set \(Z\).

For each translate \(sz\), \(z\in Z\), use the split
class representatives away from \(2\ell M\). Concretely
the finite adelic point, after the auxiliary switches, is
the product of a rational isometry \(\gamma\), a finite
point \(h\) supported at split primes away from \(2\ell M\),
and a compact element fixing both the standard lattice and
the level translates of the vacuum. Such a factorization
comes from split-place approximation in the open
stabilizer of the lattice and level vectors. The archimedean
factor is the unique correction to the finite adelic
identity. Write the finite factorization as \(sz=\gamma h k\), with \(k\) the compact factor. At \(\ell\) it gives
\[
 \gamma=(sz)k^{-1}h^{-1},\qquad \gamma^{-1}=hk(sz)^{-1},
\]
with \(h,k\) and their inverses integral there. Thus the
entries of \(\gamma\) and its inverse have a bound on the
\(\ell\)-part of their denominators from \(sz\), even though
the rational isometry varies.

When the Weil kernel is
evaluated at zero, rational automorphy removes \(\gamma\)
and the compact element acts trivially. At infinity this is
a linear substitution in each block of the harmonic
polynomial. It preserves the defining harmonic relations
and weights. In each degree \(h\) its coefficients are
sums of products of \(h\) substitution entries. The finitely
many degrees and the finite set \(Z\) therefore give a
common additional clearing exponent \(c_{\rm sub}\).
The rational source space here is the standard one and
the local entries are tested at all embeddings used in
the cover. These conditions are needed in the split
factorization and the bound on both matrices.

\begin{lemma}[The clearing exponent for the short sections]
\label{lem:theta-short-clearing}
Use the cover and fixed data in Definition~\ref{hyp:ray-cover-integrality}
and the split representative away from \(2\ell M\). At a good frame, the short section for an integral
translate mean and harmonic data of total derivative degree \(h\)
is integral after multiplication by \(\ell^{c_{\rm jet}}\), where
\[
 c_{\rm jet}=v_\ell\!\left(
     M_0^{[E:\mathbb Q]nr+n k_W}(2D_E M_0)^h\right).
\]
Here \(n,r\) are the source and target ranks and \(k_W\) is the
\(\mathbb Z\)-rank of the target marking lattice. The section and
its pinned reading use the same harmonic polynomial and
coefficients of the translate mean.
\end{lemma}

\begin{proof}
Let \(\mathcal T\) be the tuple set of torsion translates of the
source letters at level \(M\). A half-level translate has the
same reading as its representative in \(\mathcal T\): their
difference belongs to the integral lattice of that half, and
their Heisenberg pairing vanishes. The reading of a single
translate at \(t_0\) is the normalized mean with coefficients
\[
 u(t)=|\mathcal T|\,\mathbf 1_{t=t_0}.
\]
In particular its value
in the mean does not acquire a factor of \(|\mathcal T|^{-1}\)
at \(\ell\). Lemma~\ref{lem:cover-vacuum-fourier}, after its indicated
clearing, gives integral linear combinations of these coefficient
arrays.

Let \(R\) be the cover valuation ring and \(P\) the polynomial
in the matrix letters of the moduli frame. The finite module over
\(R\) generated by the coefficients of the array \((u(t)P)_t\)
is torsion-free in the coefficient extension. Choose a basis
\((b_j)_j\) and write the array as \(\sum_j b_j P_j\) with
\(P_j\) arrays over \(R\). Write \(D_\lambda\) for any of the linear operators cutting out
the prescribed harmonic weight over the fraction field of \(R\).
Coordinate extraction in the basis \((b_j)\) gives
\[
 D_\lambda\big((u(t)P)_t\big)
        =\sum_j b_jD_\lambda(P_j)=0
        \quad\Longrightarrow\quad D_\lambda(P_j)=0\quad\text{for all \(j\)}.
\]
Each \(P_j\) has the prescribed weight and derivative degree. By linearity of the jet construction \(J\),
\[
 J\big((u(t)P)_t\big)=\sum_j b_j J(P_j).
\]

The short-section construction on the split representative lattice
supplies the algebraic section and its evaluation by that Weil
kernel average, with the fixed Lie, component, and splitting-character
factors. The hypotheses are the full level marking, the standard
source lattice at \(2\ell M\) with self-duality at \(\ell\), the
integral torsion indicators and constant section, and the prescribed
row degrees. Cuspidality uses the dominant weight and strictly
positive last exponent. In the scalar on this class, the determinant
idele has an ideal representative away from \(2\ell M\).
The normalized unramified character times the square root of its
norm is a local unit at the test place. The remaining inverse
norms and lattice indices, including the index from the split
lattice sandwich, are prime to \(\ell\).

Let
\[
 N=Mu_0
\]
clear the lattice displacement, where \((u_0,2\ell)=1\).
The level index divides \(N^{[E:\mathbb Q]nr}\) and the translate
normalization divides \(M^{n k_W}\) by the standard-lattice and
marking conditions. The latter denominator is harmless for the
indicator array computed above, and may still be included in a
common bound. The good-germ calculation of
Lemma~\ref{lem:regular-theta-jets} costs \((2D_E)^h\).
The additional cost of pullback by \([N]\) is \(N^h\) since
\[
 d[N]=N\,\mathrm{id}
\]
and every derivative monomial has total
degree \(h\). Thus the possible denominator divides
\[
      N^{[E:\mathbb Q]nr} M^{n k_W}(2D_E N)^h.
\]
Its \(\ell\)-valuation is \(c_{\rm jet}\), since \(M/M_0\) and
\(N/M\) are prime to \(\ell\). Its prime-to-\(\ell\) part is a
unit at the frame. This proves the bound there. Compatibility
with degenerating integral frames is proved separately below.
\end{proof}

\paragraph{Good-reduction germs and the connection.}

\begin{lemma}[Denominators of derivatives at a good frame]\label{lem:regular-theta-jets}
For an integral regular germ in the specified theta
trivialization at a good frame, let \(h\) be the order of a jet
in integral tangent directions. If \(a_h\) is its value in the modified polarization gauge and \(R\) is the valuation ring of the frame, then
\[
 (2D_E)^h a_h\in R.
\] The bound descends from a finite extension carrying
the required torsion points.
\end{lemma}

\begin{proof}
Work over a complete valuation ring \(R\) for a good
frame. Use the compatible abelian scheme and the rigidified
symmetric theta bundle. A finite extension carries all
\(N^2\)-torsion used in the short sum and the chart with the
pulled-back frame. Multiplication by \(M\) kills the
level sections. Thus if \([N]t\) is a level section and
\(M\mid N\), then \([N^2]t=0\).
The degree-one section is the integral generator specified
by the frame. The cubical isomorphism
\[
 [N]^*\mathcal L\simeq\mathcal L^{N^2}
\]
preserves the
rigidification and extends over \(R\). It is the normalized
pullback in the short-section formula. After this pullback
the two moves are symmetric theta group translates at
\(N^2\)-torsion. The oddness of \(N\) allows the symmetric
lifts to be chosen of finite order. Each torsion section
extends to the abelian scheme by properness.

The extended torsion section lies in the kernel of the polarization of \(\mathcal L^{N^2}\)
over \(R\), because this equality holds on the generic
fiber and the dual abelian scheme is separated. The
theta group over this section is a multiplicative-group
torsor over the local base, so it has an integral lift.
Write
\[
 g_{\rm fin}=c g_R,
\]
where \(g_{\rm fin}\) is the
finite-order generic lift and \(g_R\) is the integral lift.
Take \(m\) divisible by the order of \(g_{\rm fin}\) and the
order of the underlying torsion section. Then \(g_R^m\) is
a scalar in \(R^\times\), and
\[
  1=g_{\rm fin}^m=c^m g_R^m,\qquad m v(c)=-v(g_R^m)=0 .
\]
Thus \(c\in R^\times\) and both translates extend in the
prescribed normalization. This argument permits
torsion of order divisible by the residue characteristic.
Completing in the integral rigidified trivialization
gives power series over \(R\) for these germs.

\proofstep{Changing the biextension trivialization.}
Write \(B(X,Y)\) for the formal biextension in an integral
chart and \(g(X,Y)\) for a logarithmic gauge over the
fraction field. It vanishes on the axes. The identities
for its two addition defects, transpose defect and CM
balance defect are the logarithms of the corresponding
integral units with constant term one. In particular
their constant first derivatives and their mixed
quadratic coefficients are integral. Write \(\beta\)
for the bilinear coefficient matrix of \(g\), and
\(A(a)\) for the action of \(a\in\mathcal O_E\) on the tangent
space. If the tangent rank is \(m\), integrality of the rigidified mixed coefficients gives
\[
 A(a),\quad \beta^t-\beta,\quad A(a)^t\beta-\beta A(\bar a)\ \in M_m(R).
\]

Choose integral pairs \(x_i,y_i\) giving the separability
element times \(D_E\) in \(E\otimes_{\mathbb Q}E\).
One can take a basis and its trace-dual basis with
the latter multiplied by \(D_E\). They satisfy
\[
 \sum_i x_iy_i=D_E,\qquad \sum_i ax_i\otimes y_i=\sum_i x_i\otimes ay_i\quad(a\in E).
\]
Put
\[
  \gamma=\sum_i A(x_i)^t\beta A(\bar y_i),\qquad
  \beta_0=(\gamma+\gamma^t)/(2D_E).
\]
The latter is symmetric and \(A(a)^t\beta_0=\beta_0 A(\bar a)\).
Put \(T=\gamma-D_E\beta\). The balance defect makes \(T\) integral.
Symmetrizing now gives
\[
 2D_E(\beta-\beta_0)=D_E(\beta-\beta^t)-T-T^t .
\]
Each term on the right is integral in the prescribed frame.

Subtract the symmetric balanced bilinear form \(\beta_0\)
in the logarithmic coordinates of the formal group.
The resulting gauge still satisfies the required identities.
In taking the harmonic value this change contributes
contractions of opposite CM types. These are zero
because the polynomial is in the joint harmonic kernel.
The same contraction rule gives the invariance of the
harmonic value under change of chart with the fixed
tangent and theta frames.

\proofstep{Integral derivatives in the modified gauge.}
Let \(D_u\) be the invariant derivation on the formal group
in an integral tangent vector \(u\), and let \(q(X)\)
be the logarithmic correction obtained by evaluating the
gauge at \((X,[w_0]X)\), with \(w_0\) the fixed integral
polarizing element. Let \(d_0=2D_E\). We have
\(q(0)=D_uq(0)=0\) and
\[
 d_0 D_uD_v q\in R[[X]].
\]

For the last statement, differentiate the two addition
identities at the two identity arguments. The
nonconstant terms are derivatives of the logarithms of
the integral units in the addition defects. Such first
derivatives are integral by \(d\log B=B^{-1}dB\).
The remaining terms are the bilinear coefficient
evaluated on \(u,[w_0]v\) and \(v,[w_0]u\), cleared by \(d_0\).
The algebra generated by \(R[[X]]\) and all
\(d_0D_uq\) for the finitely many tangent directions
has integral constant coefficients. It is stable under
each \(d_0D_u\).

For a regular germ \(f\) and integer
bundle power \(j\), applying \(d_0D_u\) to
\(a\exp(jq)\) acts on the coefficient by
\[
  a\longmapsto d_0D_ua+j(d_0D_uq)a .
\]
Starting from \(a=f\) proves integrality of every
degree-\(h\) derivative after multiplication by \(d_0^h\).
The exponential contributes constant term one, so
no bound on its higher coefficients is used. Apply this
letter by letter and use the integral polynomial
coefficients. A finite extension multiplies normalized
discrete valuations by a positive ramification index,
so the resulting inequality descends to the original
good frame.
\end{proof}

\paragraph{Integral tests at degenerating frames.}\label{par:theta-degenerate-arc}
We use the good-at-\(\ell\) level condition of the
admitted cover and the toroidal moduli model over the
selected coefficient place. The degeneration description
uses the split Raynaud extension and polarized period
lattice after finite extension. The relative toroidal
construction over the complete base turns these data
back into a semiabelian family. \cite{faltings-chai}
For a definite signature
all the test frames have good reduction. In the other
signature consider any integral test frame. By the
valuative criterion on the compactification it extends
to a trait.

If the trait lies in the abelian locus, the
preceding estimate applies. Otherwise, after a finite
extension, split the degeneration data and write its
period pairing as
\(\varpi^{Q(\lambda,\mu)}u_{\lambda,\mu}\) with
unit factors and the integral degeneration form \(Q\).
The period pairing on the arc and its specialization are
\[
 X^{Q(\lambda,\mu)}u_{\lambda,\mu}\ \;\longmapsto\;\ \varpi^{Q(\lambda,\mu)}u_{\lambda,\mu}.
\]
The linear
period and polarization data use the same rescaling.
The integral frame lifts on this arc with constant
tangent coefficients and the same degree-one
theta generator. This uses the derivative
identification for the semiabelian uniformization
and the integral generator at the base point.
Inverting \(\varpi X\) puts the family in the
generic abelian locus. Compatibility of algebraic
evaluation with the arc and its specialization
expresses the value of a scaled section as
\(H(X)/(\varpi X)^b\), where \(H\in R[[X]]\)
and \(b\ge0\).

\begin{lemma}[Integrality after specialization of the toroidal arc]\label{lem:theta-arc-integrality}
On the framed toroidal arc of Paragraph~\ref{par:theta-degenerate-arc},
specializing at \(X=\varpi\), let a scaled section be
\(H(X)/(\varpi X)^b\), with \(H\in R[[X]]\) and \(b\ge0\).
Suppose the section is bounded on the normalized complex frame
and integral at the algebraic good-point tests with the same
clearing power. On the arc and the original trait, respectively, one has
\[
 H\in \varpi^b X^b R[[X]],\qquad H(\varpi)/\varpi^{2b}\in R.
\]
\end{lemma}

\begin{proof}
\proofstep{Test the order at \(X=0\) on the complex frame.} For the complex test,
embed the coefficient fraction field abstractly into
\(\mathbb C\) over the specified coefficient embedding.
Use the holomorphic uniformization of the arc with
\(X=\exp(2\pi i z)\) near the puncture. To compare
the two evaluations, take an affine open chart of the
generic moduli scheme containing the Laurent-series point
and use its algebraic frame. The chart-coordinate
comparison for the holomorphic uniformization of the
split degeneration gives convergent Laurent series for
the arc's chart coordinates. The framed comparison uses
the constant tangent periods and the nonzero degree-one
theta generator of this same arc.

The frame comparison gives convergent Laurent series for the entries of the change-of-frame
matrices and for the ratio of theta generators, with
values equal to the corresponding analytic quantities.
The determinants of the frame matrices and the theta
ratio are invertible formally. Each such denominator germ can be written
\[
 d(X)=X^a u(X),\qquad u(0)\ne0,\qquad d(X)^{-1}=X^{-a}u(X)^{-1}.
\]
After shrinking the punctured disk the last germ converges. In the chart, evaluation
of a regular section is evaluation in the chart frame
followed by its algebraic weight factor. The latter uses
just the entries of the frame change, inverse
determinants and inverse theta ratio. Sums and products
of the convergent readings therefore give the analytic
value with Laurent coefficients equal to the transported
formal coefficients.  \label{cor:theta-cover-integrality-1}

To bound its value, normalize the complex
frame by a Borel change to unit Gram frame.
The cusp section is bounded there in its invariant
norm. Along the arc the diagonal tangent ratios
and the generator ratio in the weight character are
bounded as follows. In the symplectic normalization of
the arc the imaginary period matrix is
\[
 A+\operatorname{Im}(z)Q.
\]

Positivity for all sufficiently large \(\operatorname{Im}(z)\)
implies \(Q\ge0\). For a fixed \(z_0\) in that region the
matrix dominates its positive value at \(z_0\). More explicitly, writing
\[
 Y(z)=A+\operatorname{Im}(z)Q,
\]
the order inequalities for \(\operatorname{Im}z\ge\operatorname{Im}z_0\)
are
\[
 Y(z)\succeq Y(z_0)>0,
 \qquad Y(z)^{-1}\preceq Y(z_0)^{-1},
 \qquad \det Y(z)\ge\det Y(z_0).
\]
The inverse bound controls the tangent ratios and the determinant
bound controls the generator norm in the same period frame. The squared Hermitian length of each constant
tangent vector is computed using the inverse imaginary
period matrix. The squared diagonal ratios from
Gram orthogonalization are bounded by these lengths
times the fixed block scalars.

For the theta generator \(e\), the Gaussian integral over a symplectic period parallelepiped is
\[
 \|e\|^2=\frac{|\kappa|^2}{\sqrt{\det(2\operatorname{Im}\tau)}} ,
\]
where \(\kappa\) is the arc's constant generator scalar. The unit frame uses a generator of squared norm
one. Thus the generator ratio is bounded as well.
The dominant weight uses these ratios only with
nonnegative exponents. Thus the complex Laurent
function is bounded as \(\operatorname{Im}z\to\infty\).
Its negative coefficients vanish by Cauchy's formula.
Injectivity of the coefficient embedding gives
\(X^b\mid H\) in \(R[[X]]\).

\proofstep{The nonarchimedean valuation.}
Holomorphic testing on complex frames bounds the order of vanishing. To use this information in division, we still need the valuation of the coefficient in the nonarchimedean frame.

Write
\[
 H=X^b H_1.
\]
For the second test use the
valuation at \((\varpi)\) in \(R[[X]]\), extended to the
completion of its fraction field. The coefficient formula for this valuation is
\[
 v(\varpi)=1,\qquad v(X)=0,\qquad
 v\!\left(\sum_i a_iX^i\right)=\min_i v(a_i)\quad(a_i\in R).
\]
The arc is now in the abelian locus. The frame still
gives an integral point there, since its tangent and
theta generators lift on the arc. The good-point
bound extends to this valuation by the smooth
integral flag charts. Concretely, on an open
abelian chart the section is a rational function
in integral chart coordinates. If it had a pole
at an integral test with arbitrary residue extension,
one can retain its finite pole order on a point with
algebraic coefficients by smooth lifting of the
finite congruences in that chart. The condition to
stay in the open locus is one of these congruences.
Such a pole contradicts the algebraic good-point tests.

Since \(X\) is a unit in this valuation, the valuation of the
section is \(\min_i v([X^i]H_1)-b\). The good-point tests give
\[
 \min_i v([X^i]H_1)\ge b.
\]
Thus \(H=\varpi^bX^b H_2\)
with \(H_2\in R[[X]]\). The value on the original framed trait
is \(H_2(\varpi)\in R\). Both tests used the section after its
fixed multiplication by a power of \(\ell\).
\label{cor:theta-cover-integrality-2}
\end{proof}

Lemma~\ref{lem:theta-arc-integrality} thus extends the good-point
bound to the degenerating frames of the same cover.

\begin{proposition}[One phase and one denominator on the admitted covers]
\label{prop:theta-uniform-cover-clearing}
Fix the data at the start of this subsection and a finitely generated
integral lattice of input polynomials. There is a phase \(\eta\) from the
base frames and a clearing power \(\ell^c\) such that the comparison
with the algebraic reading of each admitted coefficient is integral
at every prescribed place above \(\ell\), after scaling by
\(\eta\ell^c\). One can take
\[
    c=k v_\ell(M_0)+c_{\rm flip}+c_{\rm sub}+c_{\rm jet}+c_P,
\]
with a common maximum over the degrees of the finitely many
polynomial generators when needed. The varying cover, level function,
twist, places and class label do not enter this choice.
\end{proposition}

\begin{proof}
The source frame in the norm comparison changes by a unitary map
preserving the vacuum lattice. On each component compare the
vacuum identities on a nonzero theta value. If \(c_1,c_2\)
are the two component constants, and \(A\) is the product of
that value, the Gaussian and the nonzero character factor, the
comparison of phases is the comparison of \(c_1 A\) and \(c_2 A\).
The common nonzero factor cancels in the phase ratio:
\[
 \frac{c_1 A}{c_2 A}=\frac{c_1}{c_2}.
\]
This reads the base vacuum phase. The dominant pins
on the cover read these same constants after pullback. Multiplying
its inverse by the inverse phase of the fixed lowest weight vector
in the archimedean harmonic map gives \(\eta\).

The matrix coefficients of that harmonic map on a fixed set of
lattice generators form a finite set of algebraic numbers in each
admitted coefficient field. For each coefficient \(a\), take a
positive integer \(d_a\) for which \(d_a a\) is integral. Then
\[
 c_P=\max_a v_\ell(d_a)
\]
works at all places above \(\ell\),
since \(d_a/\ell^{v_\ell(d_a)}\) is a unit there. The short-section
construction over the field gives the rational section. The local
construction of Lemma~\ref{lem:theta-short-clearing}, extended
by Lemma~\ref{lem:theta-arc-integrality}, gives integral comparison
sections. They have the same pinned adelic reading after the indicated
scaling, so injectivity identifies each with the rational section
at that place.

The determinant twist is a finite-order character on the
norm-one determinant quotient, trivial on the rational points
and the cover congruence subgroup. The indicated ray field gives
its values on the components. Each such value is a root of unity,
so multiplication by it preserves the integral section lattice.
It costs no clearing power. Summing the exponents from Fourier
inversion, the half switch, the change of variables, the short
section and the fixed harmonic map gives the displayed value of
\(c\). \label{cor:theta-cover-integrality-3}
\end{proof}

\subsubsection*{Descent to the original row}
\label{subsec:theta-cover-descent}

The preceding coefficient bounds were obtained in frames on the covers. To measure a coordinate in the original cusp lattice we must also account for the different of the coefficient extension in which those frames were evaluated.

\begin{proposition}[Uniform descent from ray covers]\label{prop:ray-cover-descent}
Suppose a section on the prescribed ray cover descends as a
holomorphic form and is integral at all good frames above
\(\mathfrak q\mid\mathfrak p\mid\ell\) of its coefficient
extension \(L/K_0\). Each coordinate \(c\) in an integral cusp basis at \(\mathfrak p\) satisfies
\[
 c\in\mathfrak D_{L_{\mathfrak q}/(K_0)_{\mathfrak p}}^{-1}
\]
under the chosen local embedding.
Unramified descent preserves integrality. For the universal row
covers contained in the fraction field of the fixed coefficient
DVR \(O\), multiplication by \(\ell^{e_O^2}\) suffices at the
chosen place, where \(e_O=v_O(\ell)\).
\end{proposition}

\begin{proof}
\proofstep{Descent of the holomorphic span.}
Let a form on the cover be invariant under the base level group.
The character relating its adelic and holomorphic realizations is invariant
under that group by the level condition on the splitting
character. Hence the weighted holomorphic form factors
through the base quotient. The bounded cuspidal holomorphic
space of the base is spanned by the complex realizations of its algebraic
cuspidal sections in the prescribed dominant weight with
positive last exponent.

In the application its degree is
at least \(8fr+4\), where \(r\) is the target rank and
\(f\) is the base degree parameter.
Pull the sections back to the cover. Linear independence persists under
extension to \(\mathbb C\). Among the cover pullbacks,
choose a basis over the cover field \(L\). If \(s_L\) lay outside the \(L\)-span of independent pullbacks \(F_1,\ldots,F_t\), then
\[
 \dim_L\langle F_1,\ldots,F_t,s_L\rangle=t+1.
\]
Their complexifications would remain independent, contradicting complex descent. Thus the rational
descent expresses the section in the \(L\)-span of the
original base sections, over the base field \(K_0\).

For the integral statement used on the ray coefficients,
let \(\mathfrak q\mid\ell\) be the selected place of \(L\)
and \(\mathfrak p\) its restriction to \(K_0\).
The base cusp sections integral at \(\mathfrak p\) form a
full lattice over \({\cal O}_{K_0,\mathfrak p}\). The
lattice is defined by the valuation bounds at all good
algebraic frames. It is finite over the localized base ring
by the unit ball bound on the finite-dimensional space.
For any section \(v\) in the base-field space, a common scalar clears its finitely many coordinates, so
\[
 \exists\,a\in\mathcal O_{K_0,\mathfrak p}\setminus\{0\}:\quad av\ \text{lies in the integral section lattice}.
\]
Thus one can use a
free lattice basis \(F_1,\ldots,F_t\). Write the descended
cover section as \(\sum_i y_i F_i\) with \(y_i\in L\),
using the pullbacks for this equality.

\proofstep{Testing the dual lattice by local traces.}
Write
\[
 s_L=\sum_i y_iF_i.
\]
For a base frame \(x\),
pullback commutes with evaluating each \(F_i\). At a place
\(\mathfrak q\) where \(s_L\) is integral, this is true
in every local embedding of \(L_\mathfrak q\) into a common
valued field. At the other places over \(\mathfrak p\), choose
\(b\in\mathcal O_L\) with \(v_{\mathfrak q'}(b)\) large
enough to clear the finitely many \(y_i\). Then every local
evaluation of \(b\,s_L\) on the lift of \(x\) is integral.
The trace section is
\[
    \operatorname{Tr}_{L/K_0}(b\,s_L)
       =\sum_i\operatorname{Tr}_{L/K_0}(by_i)F_i .
\]
Its evaluation at \(x\) is the sum over the embeddings
of \(b\,s_L\), so the base lattice test applies to it.
Since the \(F_i\) form an integral basis, each coefficient
of this trace section is integral at \(\mathfrak p\).

By weak approximation take \(b\)'s approaching any fixed
\(b_\mathfrak q\in\mathcal O_{L_\mathfrak q}\) at \(\mathfrak q\)
and approaching zero at the other local factors. Local
trace and the closedness of the integral base ring give
\[
   \operatorname{Tr}_{L_\mathfrak q/K_{0,\mathfrak p}}
             (b_\mathfrak q y_i)\in
                   \mathcal O_{K_{0,\mathfrak p}}^{\wedge}
      \quad(b_\mathfrak q\in\mathcal O_{L_\mathfrak q}).
\]
This is precisely
\[
 y_i\in\mathfrak D_{\mathfrak q/\mathfrak p}^{-1},
\]
with \(\mathfrak D\) the local different. The test used
all frames at the chosen place, as needed for the full trace.
For an unramified extension the different is the unit ideal.

For a finite prescribed set of unramified places, take one such
local decomposition \(s_L=\sum_i y_{v,i}F_{v,i}\) per place \(v\).
Rescale each \(F_{v,i}\) by a base scalar that is a unit
at the place where it was selected and clears its poles at
the other base places, and rescale \(y_{v,i}\) inversely.
The Chinese remainder theorem gives \(\rho_v\in L\) with
\[
 \sum_v\rho_v=1,
\]
\(\rho_v\) a unit at its place and
divisible by sufficient powers at all the others. Indeed one
may choose all but one by their local congruences and set
the last to one minus their sum.

At a prescribed place \(w\), the choice of powers is
\[
 v_w(\rho_v)\ \ge\
      \max_i\max\{0,-v_w(y_{v,i})\}\quad (v\ne w).
\]
The \(F_{v,i}\) are integral there by the initial rescaling.
For \(v=w\), both \(F_{w,i}\) and \(y_{w,i}\) are integral by
local descent and the triviality of the unramified different.
It follows term by term that the section
\[
 \sum_{v,i}\rho_v y_{v,i}F_{v,i}=s_L
\]
is integral at \(w\).

\proofstep{A uniform bound on the different.}
The descent formula expresses the new denominator through a local different. It remains to bound that different uniformly as the allowed ray levels change.

In the application to the universal row coefficients the
cover can have ramification at \(\ell\). The actual cover
field embeds in the fixed fraction field of the row DVR
\(O\). Take \(e_O=v_{\!O}(\ell)\) with \(v_{\!O}(O^\times)=0\)
and a uniformizer of value one. The restriction of \(v_{\!O}\)
to \(L\) is a positive integral multiple of the discrete
valuation at \(\mathfrak q\). Writing \(a\ge1\) for the integer multiplying the valuation at \(\mathfrak q\), evaluation on \(\ell\) gives
\[
 e_O=a\,v_{\mathfrak q}(\ell)\ge v_{\mathfrak q}(\ell).
\]
A bound of
\(e_O^2\) suffices for the relative different exponent.

For instance, over the maximal unramified subextension
of \(L_{\mathfrak q}/\mathbb Q_\ell\), a uniformizer has
an Eisenstein polynomial of degree \(a\), the absolute
ramification index. The different has the valuation of
its derivative. The nonzero terms of the derivative
have distinct valuation residues modulo \(a\), so its
valuation is at most that of its leading term,
\[
 a-1+a v_\ell(a)<a^2,
\]
including \(a=1\).
Transitivity of the different then bounds the relative
exponent by the same number. Multiplying by
 \(\ell^{e_O^2}\) makes the inverse-different coordinates
integral. In fact, writing \(e_{\rm abs}=v_{\mathfrak q}(\ell)\)
and \(d_{\rm rel}=v_{\mathfrak q}(\mathfrak D_{\mathfrak q/\mathfrak p})\),
\[
 v_{\mathfrak q}(\ell^{e_O^2}y_i)
       \ge e_O^2 e_{\rm abs}-d_{\rm rel}\ge0
\]
by \(d_{\rm rel}<e_{\rm abs}^2\) and \(e_{\rm abs}\le e_O\). These two numerical inequalities give the uniform estimate
\[
 d_{\rm rel}<e_{\rm abs}^2\le e_O^2.
\]
Its role is important even
when the additional level is prime to \(\ell\), since
the coefficient field also contains finite-order values.

\proofstep{The fields used before specialization.}
To see which covers this argument serves, let \(Q\) and
\(m_{\rm ray}\) be the auxiliary places and the depth
for the universal coefficient. The construction takes
the split rational primes below the selected conjugate
half of \(Q\), with the designated prime of the row twist.
The avoidance conditions make their product a squarefree
auxiliary level, prime to \(\ell\), the base level and the
host exclusions.

Write \(K_{\rm ray}\) for the ray class field at the joint level, \(A\) for the finite set of character values of the finite type and distinguished twists, and \(K_{\rm row}\) for the row coefficient field under the fixed embedding. The ray-field condition and coefficient choice give
\[
 K_{\rm ray}\subseteq K_{\rm ray}(A)\subseteq K_{\rm row}.
\]
The uniformized, dominantly
pinned cover is formed over the resulting number
field. Thus both the field conditions for integrality
on the cover and the ramification bound for descent
hold for these choices. In particular this construction
uses the cover field before specializing a character of
the variable depth.

The finite universal data on this cover are indicators
times finite-character values. Let \(A_Q\) be the class-label
group, \(b(g)=\operatorname{detchar}_Q(g)\) the determinant
label and \(S_{a'}\) the section for label \(a'\). Then the
assembled section is
\[
              \Xi_a(g)=\sum_{b\in A_Q}
                     \mathbf1_{b(g)=b}\,S_{a-b}(g).
\]
Each determinant fibre is a union of components by the
ray-field comparison. The indicators preserve the section
lattice, and the finite-character values are integral. Thus
the sum uses the same integrality bound without a denominator
\(|A_Q|\). The prescribed finite type and the trivial
upper-left reductions make \(\Xi_a\) invariant under the
original level group, so the trace descent applies.

The resulting base sections are still bounded at their
chosen base place. Under the faithful row realization,
their coordinate vectors lie in the localization of the
supremum lattice at that place. For a coefficient vector \(x\), localization gives \(x=y/s\) with \(y\) a bounded global vector and \(s\) a base scalar of valuation zero. Under the chosen embedding,
\[
 v_\ell(s)=0\quad\Longrightarrow\quad s^{-1}\in O^\times.
\]
Thus passing from \(y\) to \(x\) costs no further denominator. Taking the additional clearing power to be
\(e_O^2\), the descent yields a bounded row for every
coefficient, with the complex phase and the fixed
harmonic comparison scalar supplied by the pinned
coefficient map. The clearance depends on the datum
and on the fixed \(O\), independently of the level,
depth and class label. The passage from this complex
reading to the actual rational row lattice is given next.
\label{cor:theta-cover-integrality-4}
\end{proof}

\subsubsection*{The span of the original theta rows}
\label{subsec:theta-row-rank}

Integrality after descent bounds the coefficients of the theta rows. For their determinant to detect the original row space, these rows must also span its good type at the base tuple with the prescribed unipotent character.

\paragraph{Stable range and the lattice operator.}
Write \(n\) for the source rank, \(r\) for the target
rank and \(D_E=[E:\mathbb Q]r\) for the dimension of a
target lattice over \(\mathbb Q\). The source dimension
condition is \(n\ge 4D_E+4\), with \(n\) even as in the
row construction. In the finite Weil model
for the prescribed level, let \(\Phi(L,A)\) denote the
normalized test vector for a self-dual target lattice
\(L\) and a source lattice class \(A\). The test vectors
are summed over the finite source classes with their
stabilizer weights. Their pairings define the matrix
\[
      C(L,L')=\sum_A w(A)
                \langle\Phi(L,A),\Phi(L',A)\rangle .
\]
With the same weights on both entries, its diagonal
is a positive scalar. If \(m=[L':L\cap L']\), the
ratio of the off-diagonal bound to the diagonal is
\(m^{-n/2}\). We give the common-label calculation
that makes this bound independent of the size of the level.

\begin{lemma}[Invertibility of the overlap matrix in the stable range]\label{lem:theta-overlap-invertible}
At the specified finite level let \(n\ge4D_E+4\) be even.
For distinct lattices let \(a(L,L')\) denote their relative overlap. Then
\[
 |a(L,L')|\le [L':L\cap L']^{-n/2}.
\] The resulting matrix \(C\) has
a bounded inverse on bounded automorphic holomorphic frame functions
of the prescribed weight.
\end{lemma}

\begin{proof}
Let \(\mathfrak m\) be the level ideal, coprime to both
\(2\) and \(\bar{\mathfrak m}\), and write \(L_*\) for the reference
target lattice. A target lattice is commensurable with
\(L_*\) by multiplication in both directions by some
\(s\in\mathcal O_E\) prime to \(\mathfrak m\bar{\mathfrak m}\).
The shifted lattice is
\[
 L^{\sharp}=(\mathfrak m L:\bar{\mathfrak m}).
\]

For a source lattice \(A\), let
\(B_L=A\otimes L^{\sharp}\) be its self-dual tensor lattice
in the Weil space, where the tensor convention uses
conjugate source contractions. The normalized shifted
vacuum \(v_L\) is fixed by the operators
\[
 U(\delta)=\varepsilon(\delta)W(\delta)
\]
for
\(\delta\in B_L\). Here \(W\) is the Heisenberg action,
and \(\varepsilon\) is the common lattice semicharacter.
The level sum for \(L\) uses translates \(W(y)v_L\),
with \(y\) labeled by source rows modulo \(\mathfrak m L_*\)
after localization at \(\mathfrak m\).

The allowed labels
are the level subgroup in the first \(r\) rows and zero
in the remaining rows. The reduced matrices and their
unitary character give the same weights for all \(L\).
The two sums admit common representatives: the intermediate lattices agree at the level because the intervening scalars are prime to \(\mathfrak m\). The Chinese remainder theorem supplies representatives in their intersection. Denominator-lattice shifts fix the translate. Both the lift \(y\) and the difference \(d\) contain the factor \(2\) in the level quotient. Hence the vacuum semicharacter and their integral pairing phase equal one. Consequently
\[
 \varepsilon(d)=1,\qquad W(y+d)v_L=W(y)v_L.
\]


\proofstep{Orthogonality from the Heisenberg characters.}
Suppose \(\langle W(y)v_L,W(y')v_{L'}\rangle\ne0\). For every
\(\delta\in B_L\cap B_{L'}\), the two vectors are eigenvectors
of the same unitary \(U(\delta)\), with eigenvalue ratio
\(\exp(2\pi i\, E_A(\delta,y-y'))\) up to the convention for
the sign of the alternation \(E_A\). Put \(\Delta=y-y'\) and \(B_0=B_L\cap B_{L'}\).
Unitarity gives
\[
 \bigl(1-e^{2\pi i E_A(\delta,\Delta)}\bigr)
         \langle W(y)v_L,W(y')v_{L'}\rangle=0
       \qquad(\delta\in B_0).
\]
For a nonzero overlap this gives
\(\Delta\in B_0^\#\), with \(\#\) denoting the alternating
dual.

Choose \(s,s'\) prime to \(\mathfrak m\bar{\mathfrak m}\)
from the two comparisons with \(L_*\). The inclusion
\(ss' B_L\subseteq B_0\) implies
\[
 B_0^\#\subseteq(ss' B_L)^\#
             =(\overline{ss'})^{-1}B_L ,
\]
using \(B_L^\#=B_L\) and the conjugate adjoint rule. The source lattices in the
level sum contain \(N\mathcal O_E^n\) for an integer \(N\)
prime to a level integer in \(\mathfrak m\).
Contracting with \(Ne_k\) consequently gives
\(N\overline{ss'}(y_k-y'_k)\in L^{\sharp}\) in each row.
Take \(a\in\mathfrak m\), \(b\in\bar{\mathfrak m}\) with
\(a+b=1\). The definition of \(L^\sharp\) gives
\[
  bN\overline{ss'}(y_k-y'_k)\in\mathfrak m L,\qquad
  sbN\overline{ss'}(y_k-y'_k)\in\mathfrak m L_* .
\]
The scalar in the second inclusion is a unit modulo
\(\mathfrak m\), hence \(y_k,y_k'\) have the same level
label. This holds for every row.

For equal labels the same \(W(y)\) occurs in each argument.
Since it is unitary and the level-character weights have
modulus one, the normalized level contribution is
\(\langle v_L,v_{L'}\rangle\) times the common number of
labels. For its absolute value put
\(D'=B_{L'}/(B_L\cap B_{L'})\). The projection on the
\(B_{L'}\)-vacuum line, evaluated on \(v_L\), is
\[
      \Pi_{L'}v_L=\frac1{|D'|}\sum_{\delta\in D'}U(\delta)v_L .
\]
The terms are well defined because the denominator fixes
\(v_L\).

Distinct terms have different \(B_L\)-characters
by self-duality of \(B_L\), hence are orthogonal and each
has norm one. Locally the source lattice is free of rank
\(n\), so after the common shift the relative tensor index
is \(|D'|=[L':L\cap L']^n=m^n\). Therefore
\[
 |\langle v_L,v_{L'}\rangle|^2
   =\|\Pi_{L'}v_L\|^2=|D'|^{-1}=m^{-n}.
\]
The relative bound is the same for every source class. Summation
with the identical nonnegative stabilizer weights gives it for
the matrix \(C\).

\proofstep{Counting neighboring self-dual lattices.}
The number of self-dual \(L'\)
with index \(m\) is at most \(m^{2D_E}\). Set
\(J=L\cap L'\). Self-duality gives equal covolumes
and hence both indices over \(J\) equal to \(m\).
The intersection is an index-\(m\) sublattice of
\(L\). The number of index-\(m\) sublattices of
a free abelian group of rank \(D_E\) is at most
\(m^{D_E}\), by triangular Hermite normal form.
Indeed the recurrence in rank \(d\) is
\[
 a_d(m)=\sum_{k\mid m}k^{d-1}a_{d-1}(m/k)\le\sum_{k\mid m}k^{d-1}(m/k)^{d-1}=m^{d-1}\tau(m).
\]
There are at most \(m\) positive divisors, giving the asserted bound.

Now \(L'\) is an index-\(m\) sublattice of the
dual of \(J\), by the same covolume calculation.
The same count bounds the choices for \(L'\) for
each \(J\). Index one occurs only for \(L'=L\).
The normalized absolute row sum off the diagonal
is therefore at most
\[
       \sum_{m\ge2} m^{2D_E-n/2}
          \le \sum_{m\ge2} m^{-2}<0.65 .
\]
The estimate applies to every target lattice in the
system and every rank meeting the size condition.

\proofstep{The inverse and its multiplier.}
Let \(D\) be the positive diagonal of \(C\).
The unweighted bound on the off-diagonal row sum
makes \(D^{-1}C=1+T\) invertible on the bounded
automorphic functions by the series
\(\sum_{j\ge0}(-T)^j\). Each diagonal entry is
the fixed test-vector norm with the finite class
weights, so its inverse is bounded. The coefficient
matrix is compatible with the unit-modulus
automorphy multiplier on the target lattices.
It acts on bounded holomorphic frame functions
with the prescribed weight. The series converges
uniformly on compact subsets, since the weighted
supremum bound gives an ordinary supremum bound
on compact subsets of the frame domain. Thus its
inverse preserves holomorphicity and the multiplier.
Given the original bounded holomorphic family
\(F\), use \(F'=C^{-1}F\) in the unfolding.
\end{proof}

\paragraph{Reproduction on the thin slab.}

\begin{lemma}[Reproduction of the holomorphic frame functions]\label{lem:slab-reproducing-kernel}
For the Gaussian attached to the given stable-range weight, there is
a sufficiently thin slab and a constant \(\rho>0\) such that its
kernel reproduces each bounded holomorphic frame function of the
specified weight there. At an orthonormal frame \(x_0\) the scalar is
\[
 \rho=\frac{\int_{\rm slab}|K(x_0,y)|^2\,dy}{K(x_0,x_0)} .
\]
\end{lemma}

\begin{proof}
Use the Gaussian kernel with the row highest
polynomials in the prescribed weight. The stable
range contains these polynomials in the harmonic
tensor model. Take a finite family orthonormal
for the Fock pairing on this component. At a real
place use coordinates in which the Hermitian Gram
matrix of an orthonormal frame is
\(J=\operatorname{diag}(I_p,-I_q)\), \(p+q=r\).
The frame domain has invertible frames and
positive first \(p\) columns for this form. A thin
slab consists of frames with Gram matrix in a small
open neighborhood of \(J\), at every real place.
Indefinite Gram--Schmidt writes a frame as \(u b\)
in orthonormal coordinates, with \(u\) unitary for
\(J\) and \(b\) upper triangular.

For the first \(j\) columns of the frame \(y\), let
\(G_j=y_{[j]}^*J y_{[j]}\) and put \(\det G_0=1\). Since
\[
 G_j=b_{[j]}^*J_{[j]}b_{[j]},
\]
where the brackets on the square
matrices denote their leading blocks,
\[
 |b_{jj}|^2=\left|\frac{\det G_j}{\det G_{j-1}}\right|.
\]
On a sufficiently thin slab each \(\det G_j\) remains close to
\(\det J_{[j]}=\pm1\). The diagonal entries and their inverses
are therefore bounded. The Gram--Schmidt recursion then bounds
the remaining entries of \(b\) using these pivots and the
bounded entries of the Gram matrix.
The left unitary action preserves the slab and
Lebesgue measure because each unitary matrix has
determinant of absolute value one.

Let \(K(x,y)\) be the polarized Gaussian after
applying the highest polynomials in both arguments.
It is holomorphic in \(x\), and its right upper
triangular covariance is the weight character
\(\chi_\lambda\). For a unitary change \(u\), the two invariances are
\[
 K(y,x)=\overline{K(x,y)},\qquad K(ux,uy)=K(x,y).
\]
In the Gaussian
these changes preserve the Gram factors. Swapping
the two frames conjugates the polarized factors.
In the first argument the upper triangular change
contributes the determinant power of the Gaussian
and the diagonal powers of the two blocks of the
highest polynomial. Terms from same-argument
quadratic factors are killed by the harmonic
contractions.

At an orthonormal frame \(x_0\) the
diagonal Gaussian, in the harmonic letters \(z_l,w_l\),
is a positive determinant factor times
\(\exp(\sum_l t_l z_l w_l)\), with \(t_l>0\) from the
positive imaginary part of the polarized parameter.
For a nonzero polynomial \(P(z)=\sum_\alpha P_\alpha z^\alpha\)
in the chosen harmonic basis, coefficient extraction gives
\begin{equation}\label{eq:slab-gaussian-diagonal}
 \left.P(\partial_z)\overline P(\partial_w)
       \exp\!\left(\sum_l t_l z_l w_l\right)\right|_{z=w=0}
      =\sum_\alpha |P_\alpha|^2\alpha!\prod_l t_l^{\alpha_l}>0 .
\end{equation}
Only matching multi-indices survive the two differentiations.
Summing \eqref{eq:slab-gaussian-diagonal} over the basis and restoring the positive determinant
factor gives \(K(x_0,x_0)>0\).

\proofstep{The two torus characters.}
Consider the difference \(F\) of two functions with the same
value at \(x_0\), the prescribed right upper triangular weight
and the matching left diagonal unitary character. In coordinates
at \(x_0\), Gaussian elimination reduces a frame in the dense
cell to a lower triangular matrix \(L\). Define
\[
 D(z)=\operatorname{diag}(z,z^2,\ldots,z^r),
 \qquad L(z)=D(z)LD(z)^{-1}.
\]
The \((i,j)\) entry on or below the diagonal is \(z^{i-j}L_{ij}\),
so \(L(z)\) extends holomorphically to an invertible diagonal
matrix at zero.

For \(0<|z|\le1\), the first \(p\) diagonal entries
of \(D(z)\) have squared modulus at least \(|z|^{2p}\);
the others have squared modulus at most \(|z|^{2p}\).
Hence \(D(z)^*JD(z)\succeq |z|^{2p}J\), and the positive
Gram block stays positive after left multiplication. The
right diagonal changes that block by a congruence, so the
arc stays in the domain. For \(|z|=1\) the left and right
characters cancel and
\[
 F(L(z))=F(L).
\]
The maximum modulus
principle applied to \(F(L(z))-F(L)\) gives the equality at
zero. The value there vanishes by the prescribed weight and
\(F(x_0)=0\). The dense cell and continuity give \(F=0\).

For a fixed \(x_0\) this uniqueness expresses \(K(y,x_0)\)
as \(K(x_0,x_0)\) times the product of leading minors
\(\Delta_i\) of the coordinate matrices of \(y\) with
powers \(\lambda_i-\lambda_{i+1}\), with the weight
extended by zero at \(r+1\). This product has precisely
the required two characters. Possible negative powers
occur only at \(p,r\), where the minors do not vanish.
At a place with \(p,q>0\) the weights decrease within
each block and
\[
 \lambda_{p+1}-\lambda_p\ge n.
\]

On an orthonormal coordinate matrix \(U\), all the
leading minors are bounded by \(|\Delta_p(U)|\ge1\)
and \(|\det U|=1\). Write \(P=\mathrm{diag}(I_p,0)\) and \(Q=I-P\).
For \(U^*JU=J\) set \(N=PU+Q\) and \(N'=P+QU\).
Then
\[
 N^*N=U^*PU+Q=U^*(P-Q)U+U^*QU+Q=P+U^*QU=N'^*N'.
\]
The first
\(p\) columns of \(U\), written \(\binom{A}{C}\), satisfy
\(A^*A=I_p+C^*C\), so \(|\det N|=|\det A|\ge1\).
Writing \(W=N'N^{-1}\), we obtain
\[
 W^*W=(N^{-1})^*N'^*N'N^{-1}=(N^{-1})^*N^*NN^{-1}=I .
\] For \(N_i\) retaining the first
\(i\) rows of \(U\) and taking the other rows from the identity,
each row of \(N_iN^{-1}\) is a row of \(I\) or \(N'N^{-1}\).
All have norm one, so Hadamard gives
\[
      |\Delta_i(U)|=|\det N_i|\le|\det N|=|\Delta_p(U)|.
\]

The sum of the minor exponents up to \(r-1\) telescopes
to \(\lambda_1-\lambda_r\le-n\).
It follows that \(|K(y,x_0)|\) is bounded by a
constant times the product of \(|\Delta_p(y)|^{-n}\)
on the slab. The triangular factors there change
this bound only by a bounded factor. At a definite
place the slab is bounded and separated from singular
frames. A left unitary change and right triangular
covariance give the same bounds with any fixed first
frame \(x\).

\proofstep{Integrability in the noncompact directions.}
For integrability at an indefinite place write the
first \(p\) columns of \(y\) as \(\binom{A}{C}\).
The matrix
\[
 G=A^*A-C^*C
\]
stays in a bounded
neighborhood of \(I_p\) in the positive cone.
At fixed first columns use a left unitary isometry
to move them to \(\binom{G^{1/2}}{0}\). On the last
\(q\) columns this change has Jacobian one. The
cross Gram block bounds their top block and the
last Gram block then bounds their lower block.
Their total volume is bounded independently of
the first columns. For a square complex matrix \(A\),
polar coordinates have measure proportional to
\(d(A^*A)\) times compact unitary measure. This follows
from the triangular decomposition with positive
diagonal entries \(r_{ii}\), where both the matrix
and Gram Jacobians have the factor
\(\prod_i r_{ii}^{2(p-i)+1}\).

Since
\[
 A^*A=C^*C+G,
\]
integration in \(A\) amounts
to integration in a bounded set of \(G\) and is bounded
by a constant times \(\det(I_p+C^*C)^{-n/2}\).
To integrate \(C\), remove one row \(v\) and put
\(H=I_p+C_{\rm rem}^*C_{\rm rem}\). The determinant lemma
and the substitution \(v=wH^{1/2}\) give
\[
 \int_{\mathbb C^p}\det(H+v^*v)^{-a}\,dv
    =\det(H)^{1-a}
          \int_{\mathbb C^p}(1+\|w\|^2)^{-a}\,dw .
\]
The last integral is finite for \(a>p\). Starting with
\(a=n/2\), the last of the \(q\) removals has exponent
\(n/2-q+1>p\), since \(n/2\ge p+q\).
This proves absolute integrability of \(K(x,\cdot)\) on the
product slab. The pointwise bounds make the kernel bounded,
and the weighted frame-function norm bounds the ordinary norm
on the slab.


\proofstep{Evaluation by the torus average.}
Average a holomorphic frame function \(F\) against
\(\chi_\lambda^{-1}\) on the left diagonal unitary torus
at \(x_0\). Compactness gives a holomorphic average.
The average retains the right weight and has value
\(F(x_0)\) at \(x_0\). It also leaves the integral against
\(K(x_0,y)\) unchanged, by simultaneous invariance
and covariance in the first frame. Both this average
and \(K(y,x_0)\) have the two characters from the
uniqueness argument. Therefore
\[
 \int_{\rm slab}K(x_0,y)F(y)\,dy
 =\frac{\int_{\rm slab}|K(x_0,y)|^2\,dy}{K(x_0,x_0)}
   F(x_0).
\]
The numerator is finite and positive by boundedness,
integrability and continuity at \(x_0\). Move \(x_0\)
to any orthonormal frame by a left unitary change of
variables, and then to any frame by right weight
covariance. This proves reproduction with the same
positive constant \(\rho\).
\end{proof}

\paragraph{Unfolding the reproducing integral onto the original rows.}

\begin{proposition}[Spanning by the original theta rows]\label{prop:span-original-theta}
For the level and stable range \(n\ge4D_E+4\) of
Lemma~\ref{lem:theta-overlap-invertible}, unfolding
the doubled Weil kernel spans the bounded holomorphic target of the
prescribed weight by the original theta functions. At the base tuple write \(V_{\rm good}\) for the rational good row and \(\Theta_{\rm good}\) for the projected definite-source lift. After the prescribed scalar identifications,
\[
 \operatorname{im}\Theta_{\rm good}=V_{\rm good}.
\]
\end{proposition}

\begin{proof}
In the doubled Weil calculation the sum of the product
kernel over rational changes of target frame is
the finite sum of theta products, with the first
factor read at \((L,x)\) and the second at the other
target lattice and frame. On the unfolded side
the term for \(L'\) is
\[
    C(L,L')\int_{\rm slab} K(x,y)F'(L',y)\,dy .
\]
To obtain the doubled identity, apply the definite
Hermitian Siegel--Weil identity to the admissible
doubled test function at zero in the auxiliary
variables. The difference from the Eisenstein sum
solves the doubled heat equation. The invariant
sum of highest-polynomial differential operators
therefore also kills that difference.

In the
Eisenstein sum a maximal isotropic subspace of
the doubled space is either a unitary graph or
contains a nonzero isotropic vector in one summand.
Writing \(I(U)\) for the term attached to a maximal isotropic space \(U\), the no-coinvariants condition and the support of the finite Weil pairing give, respectively,
\[
 I(U)=0\quad\text{if }U\text{ is not a graph},
 \qquad I(\operatorname{graph}g)=0\quad\text{if }g\text{ is outside the level group}.
\]
On a graph the finite pairing is the overlap of the two test functions. For the
remaining graphs the Weil action on the translated
frame and the conjugate multiplier give the
integrand above. The double sum and
integration may be interchanged.

For fixed \(x\),
the Gaussian at \(x\) is integrable in \(y\) on a
sufficiently thin slab, the functions \(F'\) are
uniformly bounded on that slab, and
\[
 \sum_{L'}|C(L,L')|<\infty
\]
by the overlap
estimate. The group changes preserve the measure
and the slab, and the unit-modulus multiplier
cancels against the conjugate in the theta product.
Integrate over the fundamental domain for these
group changes. The finite sum of theta products
then yields a finite sum of theta functions in
\((L,x)\), whose scalar coefficients are the integrals
of the second factors against \(F'\).
The unfolded expression is
\[
    \rho\sum_{L'} C(L,L')F'(L',x)=\rho F(L,x).
\]
Dividing the coefficients by \(\rho\) gives the
theta spanning assertion. Only the strict bound
on \(T\) is used for an estimate on the rank.
The overlapping lattice system and the doubled
kernel both use the prescribed finite level in this
application. \label{cor:theta-row-rank-1}

\proofstep{Identification of the lowest-vector realization.}
In the adelic realization, evaluate the target
discrete-series map on its nonzero lowest weight
vector. If \(v_0\) is that lowest vector, cyclicity gives the test for an intertwiner \(A\):
\[
 Av_0=0\quad\Longrightarrow\quad A(U(\mathfrak g)v_0)=0\quad\Longrightarrow\quad A=0.
\]
Here \(\mathfrak g\) denotes the complexified Lie algebra on the Harish--Chandra module.
The bottom evaluation has the given compact
torus weight and is killed by the lowering
operators. After extracting the common
determinant character it gives the holomorphic
frame function of the preceding argument. Cusp
rapid decay gives its weighted boundedness. The
nonzero finite Weil vector of the base pair is
a finite translate of the level test vector for
that argument. Conjugating the frames and
undoing this translation preserve the allowed
harmonic source span. The lowest-vector identity spans the unprojected range of \(\Theta\): unfold each term as a definite-source lift summed over source classes. For the good-type projector \(e\),
\[
 \operatorname{im}(e\Theta)=e(\operatorname{im}\Theta).
\]

\proofstep{Projection onto the good row.}
Spanning before taking finite types gives access to all row coordinates. The determinant at the selected place uses the good isotypic component, so the projection has to be applied to the same spanning family.

Now take the base tuple on the actual row. The
trivial auxiliary character and empty auxiliary
level identify its projected lift with the class
sum in the row construction. For its finite unipotent group \(U\) and character \(\chi\),
the projection is
\[
 e_\chi=|U|^{-1}\sum_{u\in U}\chi(u)^{-1}u.
\]

Changing variables \(v=u_0u\) in the sum yields
\[
 u_0e_\chi=|U|^{-1}\sum_{v\in U}\chi(u_0^{-1}v)^{-1}v
            =\chi(u_0)e_\chi .
\]
Conversely \(uf=\chi(u)f\) implies \(e_\chi f=f\). For the equivariant theta map \(T\),
\[
                  T(e_{\chi,\mathrm{src}}v)
                         = e_\chi T(v),
\]
where the source operator uses the dual transformation
rule of the row reading. This imposes that rule on each
preimage of a good row vector. The equivariance of the row reading
makes the average fix the given target vector.
On the source the average imposes the dual
unipotent rule. Thus a good row with this type
is the reading of a class sum of functions on
the source with the base finite transformation
law.

On a finite list of source classes, the
harmonic coefficient lattice spans its allowed
value space over \(\mathbb C\). For the class
stabilizer, the finite trace with its level
multipliers preserves that lattice. For a stabilizer \(H\) on the complex value space, its trace and projector satisfy
\[
 \operatorname{Tr}_H=\sum_{h\in H}h,\qquad |H|^{-1}\operatorname{Tr}_H=\operatorname{pr}_{H\text{-invariants}}.
\]
The operators include the specified level multipliers. Using the unnormalized trace on
integral vectors therefore spans all such class
functions. Their class sums span the given good
row.

A summand at one source class may be moved
to a finite split-supported representative
away from the excluded places. Rational
automorphy and the finite transformation
rule transport its polynomial in the harmonic
lattice. The target character factor is
absorbed by that rule. Each resulting term
is the complex reading of an admitted original
theta datum. This uses the split-support and
level conditions in the definition of the
allowed original data, and the base-series
rational comparison on the same row.

Finally, any vector of the rational good row
whose complex reading is in the span of those
data lies in their rational span. Indeed if \(W\) is their rational span, extension of
scalars over a field is exact, so the kernel of
\[
         V_0\longrightarrow (V_0/W)\otimes_K\mathbb C
\]
is \(W\). The comparison identifies the complex reading
with this scalar extension. Applying the type projection, which
fixes the good row, gives the full span of the
projected original theta elements in \(V_0\).
The finite set of source classes and the
finite harmonic lattices give a finitely generated
row module.

Equivalently, in the rational
comparison the original theta elements lie
in the fixed finite level lattice. Over its DVR,
a generating set for a finite module contains
a subset projecting to a basis modulo the
maximal ideal. Since this module is torsion-free,
the chosen generators form a lattice basis.
This proves the original full-rank comparison
on the fixed tied data and in the stated stable
range. Its use for the division rank later
also requires the identification with the
selected point's period and row dimension.
\label{cor:theta-row-rank-2}
\end{proof}

\paragraph{Containment for the twisted rows.}
Fix the tested character at one admitted auxiliary
level, and work over its coefficient DVR \(O_k\).
The finite type multiplier on the source level
group is a product of its base character, the
splitting character at the two type primes, and
the local components of the fixed and varying
twists. The base character agrees with the
given row character there. All its values are
roots of unity in the coefficient field.
At a place over a type prime \(g_i\) the field
assumption includes the roots of order dividing
\(g_i(g_i-1)\), since the host places there have
residue field \(\mathbb F_{g_i}\).
At the distinguished and auxiliary places use
the prescribed character values in \(O_k\).
If \(\mu\) denotes the product character and \(u\) a level
element, each of its factors has finite order in the admitted
coefficient field. Thus the normalization is
\[
       \mu(u)=\mu_{\rm base}(u)\mu_{\rm spl}(u)
                          \mu_{\rm fixed}(u)\mu_{\rm aux}(u)
                       \in O_k^\times .
\]

For a twisted generator the source translate
differs from its original representative by a
change supported at the chosen split places.
Those places avoid the coefficient prime and
the fixed exclusions. In the harmonic coefficient
lattice the two readings agree over \(O_k\) after
this change. At places over the coefficient prime
the self-dual lattices agree, and all rational
denominators of the other changes are units there.
Write \(h\) for the resulting finite source point
and \(P\) for its integral harmonic polynomial.
The finite trace of \(P\) on the double coset
of \(h\) is the sum of its rational source
translates with the finite level multipliers.

At a given source point the summands are indexed
by a coset of a finite stabilizer, so this is
a finite sum in its harmonic coefficient lattice.
The reindexing uses the bijection of finite indexing sets
\[
 H\longrightarrow h_0H,\qquad h\longmapsto h_0h,
\]
where \(H\) is the stabilizer and \(h_0H\) is the coset indexing a translated sum. Including the level multiplier gives the dual transformation rule. It belongs to the integral module
of class functions for this level.

Unfold its class sum using the finite source
representatives and the same stabilizer weights
as in the original theta formula. On the kernel,
rational left translation fixes the value at
zero and the right level action is the inverse
of the multiplier in the trace. The two factors
cancel.

For one contributing class let \(n_h\) be its stabilizer
order and \(\Theta_h(P)\) the value of the corresponding
flat series term. Each of the \(n_h\) translated terms
has that same value after the two finite multipliers cancel.
The resulting term in the class sum is
\[
       n_h^{-1}\sum_{i=1}^{n_h}\Theta_h(P)=\Theta_h(P).
\]
The reciprocal arises only in the characteristic-zero
class-sum evaluation, after forming the integral trace. The result is the desired single flat
twisted series, multiplied by its norm-one
determinant character on the source point.
That scalar is a root of unity in \(O_k\).
Including it in the integral trace therefore gives
the correct twisted generator. The trace takes
no inverse stabilizer order on the integral
module.

By the comparisons on the fixed row, the class
sum is the reading of an element of the finite
level lattice. The twisted generator has the same
reading. Injectivity over the extension field
identifies these vectors, so the generator belongs
to that lattice before applying the good-type
projection. This proves the twisted containment
used in the signed-index comparison. It applies
to both sides and every admitted character of
the specified depth. \label{cor:theta-row-rank-3}

\subsubsection*{The row lattice measured by the twist}
\label{subsec:twist-row-lattices}

The spanning argument is rational. For a bound on division one must also compare the lattice of bounded forms in that span with the lattice actually used to evaluate the twist.

Fix one row and side at the chosen coefficient place. Let
\(F_0\) be the rational subspace on which the prescribed
unipotent subgroup acts by the prescribed character, and
let \(B_{\rm sup}\) be the integral lattice of bounded
forms. Write \(e\) for the rational central type idempotent,
\(g\) for the order of the finite row group, and put
\[
 M_{\rm all}=B_{\rm sup}\cap F_0,\qquad V_0=eF_0,\qquad
 B=M_{\rm all}\cap V_0,\qquad A=ge .
\]
\begin{lemma}[The original integral type lattice]
\label{lem:projected-row-type}
Over the fraction field \(K\) of the coefficient DVR the intersections satisfy
\[
 M_{\rm all}\otimes K=F_0,\qquad B\otimes K=V_0,\qquad A(M_{\rm all})\subseteq B\quad(A=ge).
\] Any residual
simultaneous Hecke eigencharacter of \(B\) is the reduction of
one of the rational signature characters on the same row
and side.
\end{lemma}

\begin{proof}
\proofstep{Integral intersection and type projection.}
These intersections use the given embeddings in the rational
space. In the row model \(B_{\rm sup}\) is a full finitely
generated lattice stable under the finite group and the
admitted Hecke operators, which commute with the group
action. Clearing denominators in this lattice shows that
\(M_{\rm all}\) spans \(F_0\) and \(B\) spans \(V_0\);
both intersections are finitely generated over the DVR.
The Hecke operators preserve each intersection by commutation.
The integral central representative of \(ge\) preserves \(B_{\rm sup}\) and sends \(F_0\) onto \(V_0\) since \(g\ne0\). Writing \(K\) for the fraction field of the coefficient DVR, we obtain
\[
 ge(M_{\rm all})\subseteq B_{\rm sup}\cap V_0=B,\qquad
 \operatorname{span}_{K}ge(M_{\rm all})=V_0 .
\]
These give the lattice and the operator used in the row
determinant. Dividing by \(g\) on the integral lattice at
this stage could alter a length.

\proofstep{One actual residual signature.}
The residual character test used after pair removal follows
from the rational signature decomposition on this same
\(V_0\). Let \(x\in B\) lift a nonzero residual simultaneous
eigenvector of the Hecke operators. Decompose in the rational simultaneous eigenspaces as
\[
 x=\sum_a y_a,
\]
where the signature \(a\) on this row and side has eigenvalues \(e_a(T)\), integral at the coefficient place. Suppose each one differs from the
residual character on some admitted operator \(T_a\). Then
\[
 \prod_a(T_a-e_a(T_a))
\]
kills \(x\), by commutativity
on the sum of these eigenspaces. Each factor acts on
\(\bar x\) by a nonzero scalar, so its product has nonzero
value on \(\bar x\). This is a contradiction since each
factor preserves \(B\). Thus some signature has every
required residual eigenvalue. A decomposition of the
integral lattice into its characteristic-zero eigenspaces
is unnecessary for this test.

\end{proof}

\begin{lemma}[Clearing the universal row coefficients]
\label{lem:universal-row-denominator}
For the finite original basis there exists \(c\ge0\) such that every individual universal coefficient function \(\Xi\) at every allowed level satisfies
\[
 \ell^c\Xi\in B_{\rm sup}.
\]
This uses the boundedness and common-phase comparisons. The same exponent
clears the sum of coefficient functions whose class labels identify
in the row.
\end{lemma}

\begin{proof}
\proofstep{A denominator for all individual class coefficients.}
For the common denominator, write the universal lift of a
chosen datum at \(Q,M\) as a finite group-ring sum
\(\sum_\delta \Xi_\delta[\delta]\). In the row realization,
the coefficient functions for the individual class labels
lie in the rational row and have the prescribed unipotent
type. If the bounded-coefficient estimate gives
\[
 \ell^{c(a)}\Xi\in B_{\rm sup}
\]
for each such coefficient
of datum \(a\), with \(c(a)\) independent of \(Q,M\),
then the same exponent clears all the \(\Xi_\delta\)
in \(M_{\rm all}\). Indeed the map from class labels to
coordinates can identify labels, and its coefficient at
\(\delta\) is the sum of their coefficients. Thus the exponent used in \(C_Q\) is uniform in the level and can be chosen by
\[
 c=\max_{a\text{ in the original basis}}c(a).
\]
The bound needed here is on the individual
universal coefficients as the modulus varies.

\proofstep{The common phase and its inverse in the coefficient field.}
Propositions~\ref{prop:theta-uniform-cover-clearing}
and~\ref{prop:ray-cover-descent} give one phase and one
clearing power on each generator of the harmonic coefficient
lattice. A finite linear combination of these generators
admits the maximum of their clearing powers, since the phase
is common and the coefficient scalars are integral.
The harmonic polynomial realization is injective in the
range where the sum of the two signature dimensions is at
most the target rank. Consequently the expression for the
coefficient function applies to its given polynomial.

Let
\(\mathrm{rd}\) be the complex realization of the rational
row. With its fixed comparison scalar incorporated, the
integral combination gives a single \(\lambda\in\mathbb C^\times\)
such that, for every coefficient \(x\) under consideration,
there is \(x'\in B_{\rm sup}\) with
\[
 \mathrm{rd}(x')=\lambda\mathrm{rd}(x).
\]
Faithfulness
here includes linear independence over \(\mathbb C\) of
the images of a rational basis. If some \(x\ne0\), comparing
a nonzero coordinate gives \(\lambda\) in the embedded
fraction field. Clear its inverse once by a power of
\(\ell\); injectivity gives that power times every \(x\)
in \(B_{\rm sup}\). If every \(x\) vanishes, exponent zero
suffices. This gives the coefficient estimate in
\(B_{\rm sup}\) used above.

\end{proof}

Let \(L_0=e\Theta_0\) be the projected original theta module
of the specified row indices, and \(L_\psi=e\Theta_\psi\)
the projected integral module generated by their allowed
twisted series on the chosen side, over the DVR \(O_k\)
of the character. Write \(B_k\) for the scalar extension
of \(B\). For full lattices use the signed index
\([U:V]_\ell\), the valuation of the determinant of a basis
of \(V\) in a basis of \(U\), normalized by \(v(\ell)=1\).
Negative values are allowed.

\begin{proposition}[Signed index of the actual twisted row]
\label{prop:twisted-row-index}
With the faithful coefficient comparisons and the selected
character and row defined above, \(L_\psi\) is full in \(V_{0,k}\)
and
\[
          [(L_0)_k:L_\psi]_\ell\le\epsilon_{j,b}.
\]
Its inclusion in the projected level lattice persists under
finite coefficient extension, and the index bound persists after
extension in the normalization \(v(\ell)=1\).
\end{proposition}

\begin{proof}
\proofstep{Specialization of the universal series.}
The chosen basis and full-rank
condition give \(L_0\) full in \(V_0\). The twisted type law puts \(L_\psi\subset V_{0,k}\). It is finitely generated, being a submodule of the projected finitely generated level lattice over a DVR. At a character of order dividing \(\ell^k\)
factor evaluation through its chosen ray coordinates. Avoidance makes every basis datum in \(C_Q\) prime to \(Q\) in level. The specialized row series belongs to the integral span of that datum's twisted series by the following calculation. At a chosen place of residue
field \(k_P\), size \(q\), put \(r_0\le n_0\) for the two
matrix dimensions and let \(y\) have \(r_0\) rows. Sum over
codimension-\(r_0\) planes \(U\subset k_P^{n_0}\).

For each plane the normalized local contribution is
\[
 q^{-r_0(n_0-r_0)}\begin{cases}[d_U(y)],&\operatorname{row}(y)\oplus U=k_P^{n_0},\\ 0,&\text{otherwise}.\end{cases}
\]
Here \([d_U(y)]\) is the group-ring element attached to the quotient determinant. Choose a basis adapted
to each plane. Then
\[
 d_U(y)=c_U\det((yg_U)_{[r_0]})
\]
for a fixed residue unit \(c_U\), an invertible change of
basis \(g_U\), and the first square block of \(yg_U\).
The flat local function at the coordinate character \(\chi\) is
\[
 B\longmapsto\begin{cases}\chi(\det B),& B\text{ invertible},\\
 0,& B\text{ singular}.
 \end{cases}
\]
The preceding formula expresses the evaluation as a sum of its translates. The normalization times the residue character of \(c_U\) is integral since \(\ell\nmid q\). Lift the changes of basis integrally at the chosen places. Multiplication and linearity of theta summation now give an integral combination of the allowed translates. The character of the determinant
agrees with the specialization by Artin reciprocity and the
coordinate restriction law.

In the rational row comparison
both forms have these same complex evaluations; injectivity
after scalar extension gives membership on the prescribed
row. After applying \(A\),
its value belongs to \(g L_\psi\). Hence every full lattice
\(T\) containing \(g L_\psi\) satisfies the same estimate
from the determinant,
\[
           [B_k:T]_\ell\le [B:gL_0]_\ell+\epsilon_{j,b}.
\]
These comparisons concern the whole specialized theta
module; equality between a single universal column and a
chosen basis vector of that module is not needed.

\proofstep{Full rank and scalar cancellation in the index.}
The lattice index only measures a determinant if the specialized rows have full rank. It remains to verify this rank directly before cancelling the scalar common to the two covolumes.

In particular \(L_\psi\) spans \(V_{0,k}\). To see this
without assuming its rank in the index inequality, take
\[
 T_n=L_\psi+\ell^n B_k.
\]
If two consecutive terms of this
descending chain coincide, Nakayama on \(T_n/L_\psi\)
gives \(T_n=L_\psi\), which has full rank. Otherwise every
strict inclusion increases \([B_k:T_n]_\ell\) by at least
\(1/e_{O_k}\); every \(T_n\) contains \(g L_\psi\), so
this contradicts the common bound. We may now take
\(T=gL_\psi\). The valuation of a determinant is unchanged
under normalized finite DVR extension, and indexing two
full lattices after multiplying both by the same nonzero
scalar leaves their relative index unchanged. Additivity
therefore gives
\[
           [(L_0)_k:L_\psi]_\ell\le\epsilon_{j,b}.
\]

For any further finite extension \(O'/O_k\) and full
reference lattice \(M'\) in that space the bound is
\(
 [M':(L_\psi)_{O'}]_\ell
 \le [M':(L_0)_{O'}]_\ell+\epsilon_{j,b}.
\)
The inclusion into the projected level lattice persists under
this flat extension. These are the lattice, inclusion and
index assertions used at every permitted power of the
single selected character. Their geometric inputs are the
bounded-coefficient estimate, type preservation, the faithful row
comparison and integral containment, and the original
full-rank basis on these very levels. \label{cor:twist-row-lattices-1}
\end{proof}


\subsubsection*{Conversion of the split row bound}

The auxiliary character and both row lattices are now fixed. We return to the small split primes and insert their local exponents into the normalized congruence length.

\begin{proposition}[Weighted split loss at the small primes]\label{prop:small-split-loss}
On these split row and column declarations, let
\(H=b+S_c+\eta_c+t+2\), use the full cyclotomic depth \(s\) in Lemma~\ref{lem:small-split-finite-colength}, and
put \(\Lambda=s+\log_\ell(2H)+1\).
Let \(c_{\rm loc}\) denote the normalized local cost after subtracting the scalar and finite Selmer baseline. Then
\[
 c_{\rm loc}\le 6000 f m\,H s+30000 f^2t^2 R\Lambda.
\]
For the ranks and weights on the selected family this contributes
\(O(fmbz^2)\) after summation over \(\mathcal S_{\rm sp}\)
with weight \(\log\ell\). Constants and thresholds are uniform in
the selected point and good eigencharacter.
\end{proposition}

\begin{proof}
\proofstep{Put the lifting exponents in the weight and depth ranges of the selected point.}
Put \(H=b+S_c+\eta_c+t+2\), where \(S_c,\eta_c\) are the
two shift and weight sizes in the selected parameter choice.
Here \(b\le H\le6b\). Write
\[
                    \Lambda=s+\log_\ell(2H)+1 .
\]
The local lifting and invariant estimates before numerical
conversion use
\[
 \lambda_N(h)=s+\lfloor\log_\ell(2h)\rfloor+ \lfloor\log_\ell N\rfloor+1,
\]
with \(h\) the absolute
jump cutoff in each estimate. The selected level satisfies
\(N\le8H\), and both jump cutoffs used here satisfy \(h\le4H\).
Therefore \(\lambda_N(h)\le3\Lambda\), since
\[
 \log_\ell(8H)\le\log_\ell(2H)+2\le\Lambda.
\]

The local lift exponent is \(32(2hs+1)+\lambda_N(h)\).
For the invariant and block estimates, the bounds are
\(20(4nR+1)\lambda_N(h)\) and
\(40(t+1)(4nR+1)\lambda_N(h)\), respectively.
Using \(f=2n\), \(t\ge1\), and
\[
 \Lambda\le4Hs
\]
gives the coefficient \(1000\) below.

\proofstep{Subtract the scalar and rank-complement terms on the same cut.}
The split row inequality bounds the normalized localized
congruence term for a good Hecke character by a sum of a
baseline and the local row and column costs described below.
In the bound used for the numerator, the baseline is
\(s_\pi+V_\ell+v_\ell(h_E)+\ell^{-2}\).
Here \(s_\pi\) is the relaxed global Selmer length,
\(V_\ell\) is the sum of strict localization lengths through
twist \(t-1\), and \(h_E\) is the class number of \(E\).
The scalar and finite Selmer terms have been separated at this point. The remaining tame torsion term is controlled by the local killing condition, which bounds
\[
 t\sum_v v_\ell(\zeta_v-1)
\]
in this formula by \(\ell^{-2}\),
where the \(\zeta_v\ne1\) are the indicated tame test values.

Let \(D_c,D_r\) denote the sums of the local degrees of the
column and row places, and let \(k_c,k_r\) be their numbers.
All these places are over \(\ell\). The degree formula
gives \(k_c\le D_c\le2f\) and \(k_r\le D_r\le2f\).

Write \(e'\) for the lift exponent, \(c_B\) for the block
bound, \(v_\Delta\) for the factorial denominator length
and \(h_0\) for the scalar invariant bound above. The two
column terms involving local weights have bounds
\(r c_0(u)\) and \(h_W(u)\) at a column place \(u\).
The cyclotomic exponent and local invariant estimates used
in the split bound have the following size requirements
with \(c=1000\):
\[
 \begin{gathered}
  e'\le cHs,\qquad c_B\le ctfR\Lambda,\qquad
  c_0(u)\le cfR\Lambda,\qquad h_W(u)\le crfR\Lambda,\qquad
  h_0\le c\Lambda,\\
  v_\Delta\le ts+v_\ell(t!)\le t(s+1)\le t\Lambda .
 \end{gathered}
\]
The inequality for \(h_0\) follows from the scalar
calculation above. The second inequality for \(v_\Delta\)
uses \(v_\ell(t!)=\sum_{a\ge1}\lfloor t/\ell^a\rfloor\le t\).

\proofstep{Sum the column and row terms with their respective local degrees.}
The split row formula then has total cost above its
baseline at most
\begin{align*}
 D_c m(2e'+v_\Delta)
 &+k_c(e'+2v_\Delta+2crfR\Lambda)\
 &+t(D_r+2k_r)c_B+4R(k_c+1)v_\Delta+tk_ch_0 .
\end{align*}
The terms containing \(e'\) contribute at most \(6 c f m\,H s\).
In order, the six remaining terms contribute at most
\(2,4,4,6,12,2\) times \(cf^2t^2R\Lambda\).
For instance \(4R(k_c+1)v_\Delta\le12fRt\Lambda\),
and \(f,t,c\ge1\). Hence the excess over the baseline is
bounded by
\begin{equation}\label{eq:split-row-cost}
             6000 f m\,H s+30000 f^2t^2R\Lambda .
\end{equation}
This calculation uses the cohomological split row
inequality on the same row-and-column declaration.

\proofstep{Take the logarithmically weighted sum only after the local conversion.}
For a uniform size comparison put \(z=x^\epsilon\), as in
Section~\ref{sec:comparison}, and take a fixed \(C_0\ge1\)
with \(f,t\le x^{C_0}\), \(R\le2x^{C_0}\),
and \(Q_S\le n^2\) with \(n\le x^{C_0}\).
For \(\ell\ne q,d\), \(v_\ell(N)=v_\ell(Q_S)\), so
\(s\le3C_0\log x\) for \(\log x\ge1\).
Also \(\log_\ell(2H)\le2z+6\), using \(H\le6b\)
and \(\log\ell>1/2\).
For \(z\ge1\), this gives \(\Lambda\le(3C_0+9)xz\).

For any fixed nonnegative coefficients \(A,B\) we obtain
\[
 A f m\,H s+B f^2t^2R\Lambda
 \le18AC_0 fmb\log x+
             2B(3C_0+9)x^{5C_0+1}z .
\]
Thus \eqref{eq:split-row-cost} is bounded by
\(C fmb\log x+Cx^Cz^2\) with a single \(C\) depending on
\(C_0\). The second term is at most \(C fmb\) past one
threshold since \(b=\exp(x^\epsilon)\). These constants
and the threshold are independent of the point, the
coefficient field and the particular good eigencharacter.
The weights of all the primes just treated satisfy
\[
    \sum_{\ell\in\mathcal S_{\rm sp}}\log\ell
    =\log\!\prod_{\ell\mid 3Q_S}\ell\le\log(3Q_S).
\]
The resulting weighted cost is \(O(fmbz^2)\) beyond the
threshold, using \(\log x\le z\).

 The normalized Galois algebra, tame tests and
global restriction argument apply at the retained points.
The character and finite-class estimates use the same local
lattices as the invariant estimate.
The simultaneous twist selection is given in
Section~\ref{subsec:retained-twist}.
The remaining prime ranges enter the same division numerator.
\label{cor:small-prime-lengths-4}
\end{proof}


\subsection{Integral local models for the selected lattices}
\label{subsec:cut-wach-models}

At unramified primes we need Wach modules realizing the signature lattices and their labeled jumps.

Keep the lift, the twisted coefficient place and its assigned
tuple fixed. At an odd rational prime \(p\) unramified in the auxiliary field \(E\), let \(F\) be a completion of \(E\) above \(p\). Use a finite extension of
the coefficient field containing all embeddings of \(F\) and
the first \(p\)-th roots of unity. Write \(O\) for its integers,
\(\varpi\) for its uniformizer and \(k_O\) for its residue field.
Write \(a\) for an ordinary coefficient embedding of
\(O_F\otimes_{\mathbb Z_p}O=\prod_a O\), and set
\(\rho(a)=a\circ\operatorname{Frob}_F\), with
\(\operatorname{Frob}_F\) the arithmetic Frobenius on \(O_F\).
The terminal Hodge label of the ordinary target \(a\) is
\(i=\rho(a)\). Expressing the same Frobenius arrow in the two
sets of labels gives
\begin{equation}\label{eq:wach-terminal-label-arrow}
 \begin{array}{rclcl}
 \text{ordinary embeddings}&\rho(a)&\xrightarrow{\ \varphi\ }&a,\\
 \text{terminal labels}&i^-=\rho(i)&\xrightarrow{\ \varphi\ }&i .
 \end{array}
\end{equation}

By \eqref{eq:wach-terminal-label-arrow} the ordinary factor carrying label \(i\) is
\(\rho^{-1}(i)\). This is the convention of
Paragraph~\ref{par:wach-level-models}. The de Rham evaluation at
the first root uses \(\varphi^{-1}\), so we transport the
crystalline coefficients of the target by inverse crystalline
Frobenius to read the filtration at \(i\).
The factors over distinct completions of \(E\) are
independent. The reflected arrays at conjugate embeddings use
the same predecessor permutation, since conjugation commutes
with the unramified Frobenius on the embeddings of \(E\).

\begin{auxlemma}[Integral Wach factors]\label{lem:covariant-wach-factor-free}
Let \(F/\mathbb Q_p\) be unramified, with \(p\) odd, and take the
coefficient ring \(O\) splitting its embeddings and containing the
first \(p\)-th roots of unity. For a stable free \(O\)-lattice \(T\)
in a crystalline representation, put
\[
 N(T)=D(T)\cap N(T[1/p]),\qquad N_a=e_aN(T)
\]
in the rational \'etale module, where \(e_a\) is the ordinary
idempotent of \(O_F\otimes_{\mathbb Z_p}O\). Each \(N_a\) is
free over \(O[[\pi]]\) of rank \(\operatorname{rk}_O T\).
\end{auxlemma}
\begin{proof}
Here \(D\) is the cyclotomic \'etale equivalence for integral
and finite \(\mathbb Z_p\)-representations. Use its
exactness, tensor compatibility and compatibility with scalar
actions \cite[Theorem A.3.4.3]{fontaine-etale}. Berger's Wach
comparison applies because \(F\) is unramified. It gives an
integral free lattice before extending coefficients, a basis of
\(D(T)\) after inverting the cyclotomic variable and completing,
trivial cyclotomic action at the origin and a Frobenius matrix
and inverse with at most powers of the first cyclotomic factor
as denominators. The ring action of \(O\) is carried along by
functoriality. For any coefficient idempotent \(e\) of \(O_F\otimes O\), the corresponding summand satisfies
\[
 eN\ \text{is finite over}\ O[[\pi]].
\]
Here \(N\) denotes the full integral Wach module under consideration. Before the coefficient
action is taken into account the integral Wach lattice is free
over \(O_F[[\pi]]\), so its reduction modulo \(\pi\) is
\(\mathbb Z_p\)-torsion-free. The direct summand \(N_a/\pi N_a\)
has the same property.

Since a power of \(\varpi\) is associated
to \(p\) in \(O\), it is also torsion-free over \(O\), hence free.
Choose lifts \(b_1,\ldots,b_s\) of an \(O\)-basis of this reduction. Nakayama gives a surjection
\[
 O[[\pi]]^s\to N_a
\]
with these images. Let \(K\) be the kernel of the indicated surjection. Reducing
a relation modulo \(\pi\) puts each coordinate in
\(\pi O[[\pi]]\). As \(N_a\) is \(\pi\)-torsion-free, division
by \(\pi\) gives another relation, so \(K\subseteq\pi K\).
Iterating gives
\[
       K\subseteq\bigcap_{k\ge1}\pi^k O[[\pi]]^s=0 .
\]
Thus the surjection is an isomorphism. The \'etale comparison identifies
\(s\) with the dimension of the coefficient representation.
\cite[Proposition III.4.2 and Corollary III.4.5]{berger-wach}
\end{proof}

For a finite extension \(O'/O\), choose an \(O\)-basis of
\(O'\). The integral and rational comparison modules of
\(T\otimes_O O'\), and their intersection, are the finite
direct sums determined by that basis. This yields
\[
       N(T\otimes_O O')=N(T)\otimes_O O'
\]
with its natural scalar action. The same finite direct sums on
continuous cochains commute with taking cohomology, and with the
images defining the integral finite condition. If \(e'\) is
the relative ramification index, tensoring a composition factor
\(O/\varpi O\) gives \(O'/\varpi O'\) of length \(e'\).
Dividing the respective lengths by \(\operatorname{length}_O(O/pO)\)
and \(\operatorname{length}_{O'}(O'/pO')\) therefore gives the
same normalized length. Faithful flatness
also detects the required inclusions over \(O\). We use
these comparisons only at the coefficient place already fixed.

\paragraph{The norm coordinate and the labeled heights.}
On each factor of Lemma~\ref{lem:covariant-wach-factor-free},
use the ordinary cyclotomic variable \(\pi\), with
\(\gamma(1+\pi)=(1+\pi)^{\chi(\gamma)}\) and
\(Q=\big((1+\pi)^p-1\big)/\pi\). With \(\Delta\) the tame
group of order \(p-1\), put
\[
     Z=\prod_{\delta\in\Delta}\delta(\pi).
\]
Since \(\delta(\pi)\equiv\chi(\delta)\pi\pmod{\pi^2}\),
we have \(Z=v\pi^{p-1}\) for a unit \(v\in O[[\pi]]^\times\).
For an invariant series \(f\) with first term \(a_n\pi^n\ne0\),
comparison on the \(n\)-th graded piece gives
\[
 a_n\chi(\delta)^n=a_n\quad(\delta\in\Delta),
 \qquad n=(p-1)m\quad(m\ge0).
\]
Subtract \(a_n v(0)^{-m}Z^m\) to raise the order. Iteration
in the complete ring, together with the invariance of \(Z\),
gives
\[
 O[[\pi]]^\Delta=O[[Z]].
\]
Lifting a basis of \(N(T)/\pi N(T)\) and averaging
it over \(\Delta\) gives an invariant basis on each factor.

In
this basis the tame-invariant lattice \(N_i\) is free over
\(O[[Z]]\) and gives the full ordinary Wach lattice after
extension to \(O[[\pi]]\). At \(\pi=\zeta^a-1\), with
\(\zeta\) primitive of order \(p\) and \(1\le a<p\), we have
\(Z=p\), by the product of the nonzero cyclotomic differences.
Weierstrass preparation thus makes \(Z-p\) associated to \(Q\)
over \(O[[\pi]]\). The constant term and the initial unit
coefficient give the norm parameter used in the lift comparison.

Choose a cyclotomic generator \(\gamma\) with
\[
 u=\chi(\gamma)=1+pc,\qquad c\in\mathbb Z_p^\times,\qquad v_p(u-1)=1.
\]
The quotient
\(\gamma(Z)/Z-1\) is a series in \(Z\) which vanishes at
\(Z=p\), since \(\gamma\) fixes each primitive root. Hence
\(\gamma(Z)=Z(1+(Z-p)A(Z))\) for an integral series \(A\).
The constant term of \(\gamma(Z)/Z\) is \(u^{p-1}\), from
the leading ordinary term of \(Z\). Consequently
\[
 1-pA(0)=u^{p-1}\equiv 1+(p-1)pc\pmod{p^2},
 \qquad A(0)\equiv c\pmod p .
\]
In the residue ring Frobenius substitutes \(Z^p\). In the
ordinary lifting calculation put \(X=-Z\), \(q=X+p\). With bars
denoting reduction in the residue ring the quadratic terms are
\begin{equation}\label{eq:wach-norm-cyclotomic-sign}
 \overline{\gamma(Z)}\equiv Z+\bar c Z^2\pmod{Z^3},\qquad
 \overline{\gamma(X)}\equiv X-\bar c X^2\pmod{X^3}.
\end{equation}
Taking \(u=1+p\) in \eqref{eq:wach-norm-cyclotomic-sign} gives
quadratic coefficient \(-1\) in \(X\). This change of variable also changes the sign
of the divided matrix \((\gamma-1)/Z\) on the constants.

In the terminal labels write \(P_i\) for the Frobenius matrix
from \(i^-\) to \(i\) and \(1+Z D_i\) for the cyclotomic matrix
on \(i\). If the Hodge jumps at \(i\) are \(d_{i,1},\ldots,d_{i,r}\),
the elementary divisors of \(P_i\) at a primitive root have
orders \(d_{i,a}\). For an unramified quadratic host and a
line of jumps \((0,1)\), the orders in ordinary target order
are \((1,0)\). The terminal indexing gives \((0,1)\).
\begin{lemma}[First-level heights in the terminal index]
\label{lem:first-level-terminal-heights}
At each terminal label \(i\) the elementary-divisor orders of the
first-level matrix \(P_i\) are the signed jumps \(d_{i,a}\). For
\(-M_i\le d_{i,a}\le H_i\) its two integral height bounds are
obtained by clearing \(P_i\) by \((Z-p)^{M_i}\) and \(P_i^{-1}\)
by \((Z-p)^{H_i}\). For nonnegative jumps each \(a\)-minor is
divisible by the power given by the sum of the \(a\) least
jumps. In particular the characteristic polynomial, including multiplicities, is
\[
 \det(T-\bar D_i(0))=\prod_a(T+\bar c\,d_{i,a}).
\]
\end{lemma}

\begin{proof}
For the general comparison, first take nonnegative jumps and
compare the Wach lattice to the constant crystalline lattice
over the ring of convergent series. Let \(C_i\) express the
Wach basis in the constant basis. At the root, in the
filtration-adapted basis, the Wach lattice has generators
\(s^{-d_{i,a}}\) times the constant generators. Thus \(C_i\)
has elementary-divisor orders \(-d_{i,a}\). The matrix \(U_i\)
of the Frobenius image of the source Wach basis in the target
constant basis is invertible over the root power series ring.
Indeed the root maps to \(\pi=0\), where the comparison has
orders zero, and crystalline Frobenius is invertible.
Write
\(C_i=A_i\operatorname{diag}(s^{-d_{i,a}})_a B_i\) with both
outer factors invertible in the root power series ring.

Since
\(C_iP_i=U_i\) and \(U_i\) is invertible there,
\[
       P_i=B_i^{-1}\operatorname{diag}(s^{d_{i,a}})_a
                               A_i^{-1}U_i .
\]
So its elementary-divisor orders are exactly \(d_{i,a}\).
The inverse Frobenius transport of the constant basis puts these
orders and the cyclotomic matrix at the same terminal label.
A common twist adds its rank-one order to all \(d_{i,a}\).
Cancelling the rank-one factor gives the same formula for signed jumps.
In the same root chart, with local variable
\(s=\log((1+\pi)/\zeta^a)\), the action is \(s\mapsto us\).
The constant crystalline basis is fixed, and the matrix on the
Wach lattice at the root has eigenvalues \(u^{-d_{i,a}}\)
with multiplicity. These are the first-level cases of
Berger's Propositions II.3.1, III.2.1 and III.2.2.
\cite{berger-wach,dousmanis-wach}\label{cor:cut-wach-models-1}

In particular, if \(-M_i\le d_{i,a}\le H_i\), then
\[
 (Z-p)^{M_i}P_i,\qquad (Z-p)^{H_i}P_i^{-1}
                \quad\text{have entries in }O[[Z]].
\]
For descent of these entrywise conditions, start with the
coarse Wach common denominator, a power of \(Q\). In ordinary
coordinates multiply by it and remove the excess factors by
monic division at every primitive root. A monic linear factor
with a root of positive valuation in \(O\) gives integral
division with remainder the value at that root. Its prime ideal
has quotient \(O\), so distinct factors can be removed with their
specified multiplicities. The factors have product \(Q\).
The resulting bounded entries are invariant, so they descend to
\(O[[Z]]\).

For nonnegative jumps, apply the same monic division to any \(a\)-minor \(D_a\) of \(P_i\). Ordering the jumps increasingly gives
\[
 (Z-p)^{\sum_{b=1}^{a}d_{i,(b)}}\mid D_a\quad\text{in }O[[Z]].
\]
At \(Z=p\) subtract the identity
from the cyclotomic matrix and divide by \(p\) before reduction.
Since \(Z=p\) reduces to \(0\), this gives
\[
     \det(T-\bar D_i(0))=\prod_a(T+\bar c d_{i,a})
\]
in the residue field. This uses that the coefficients
of the characteristic polynomial commute with reduction and
that \((u^{-d}-1)/p\equiv-cd\pmod p\). On a scalar line
the Frobenius ideal is \((Z-p)^d\), by the two height bounds
with the sole jump \(d\).
\end{proof}

\begin{lemma}[Height bounds on the induced residue lattices]
\label{lem:localized-induced-wach-heights}
Put \(R_0=k_O[[Z]]\) and \(K_0=k_O((Z))\). Let
\(L_i=N_i/\varpi N_i\) with the terminal labels of Lemma~\ref{lem:first-level-terminal-heights}. A Galois-stable flag in
\(T/\varpi T\) gives a flag in the corresponding \(K_0\)-factors,
of the same ranks. At each factor, intersecting with \(L_i\) and
projecting to a flag quotient gives a free \(R_0\)-lattice. Let \(B_i\) denote the Frobenius span from \(i^-\) to \(i\).
If \(Z^{H_i}L_i\subseteq B_i\subseteq Z^{-M_i}L_i\), write \(L_i'\) for any induced flag-subquotient lattice and \(B_i'\) for its Frobenius span. Then
\[
 Z^{H_i}L_i'\subseteq B_i'\subseteq Z^{-M_i}L_i'.
\] Each marked line is identified by its specified
crystalline rank-one comparison in the residual \'etale realization.
\end{lemma}

\begin{proof}
Exactness of the \'etale equivalence at multiplication by \(\varpi\)
identifies the reduced module with the module of \(T/\varpi T\).
The invariant field of the tame action on \(k_O((\pi))\) is \(K_0\).
Indeed \(Z\) has order \(p-1\), and division by this series gives
extension degree \(p-1\). The \(p-1\) substitutions are distinct by
their linear coefficients. Field descent preserves dimension, and
taking the tame invariants is exact by averaging. Also the map
from the Wach lattice to the \'etale module is \(\varpi\)-saturated:
for \(b\in D(T)\) one has
\[
 \varpi b\in D(T)\cap N(T[1/p])\quad\Longrightarrow\quad b\in D(T)\cap N(T[1/p]),
\]
since the rational Wach module is closed under \(\varpi^{-1}\). Reduction therefore gives the
specified inclusion of \(L_i\) into the residual \'etale realization.

For the maps of a flag, take inverse images of its subspaces in
\(T\). Each inverse image contains \(\varpi T\) and has the same
crystalline rationalization as \(T\). Functoriality and exactness
of the \'etale equivalence identify the images of these maps after
reduction. The given rank-one comparison determines the individual
marked line in this flag. In particular this determines its
position before forming any saturated power-series lattices.
\label{cor:cut-wach-models-2}

For a stable rational subspace \(W_i\) of a flag, set
\(S_i=L_i\cap W_i\). Its quotient \(L_i/S_i\) is torsion-free,
so both are free over \(R_0\). Denote by \(B_i\) the Frobenius
span of \(L_{i^-}\) in the terminal factor. The height conditions
reduce to
\[
    Z^{H_i} L_i\subseteq B_i\subseteq Z^{-M_i}L_i.
\]
Flat Frobenius base change and a basis adapted to \(S_{i^-}\)
identify the Frobenius span of \(S_{i^-}\) with \(B_i\cap W_i\).
Intersection of the displayed inclusions with \(W_i\) gives
\[
    Z^{H_i} S_i\subseteq B_i\cap W_i\subseteq Z^{-M_i}S_i.
\]
Taking images in a stable rational quotient preserves the two
inclusions. This gives the assertion for a subquotient by first
intersecting and then projecting.
\end{proof}

The ordinary row and column comparisons use actual surjections
from the matrix span at the retained point. Their rational sources
are subrepresentations of finite sums of the normalized family and
its dual twists. Strictness of crystalline comparison puts the
Hodge jumps of source, target and kernel in the same labeled
boxes. If a positive twist is used, the same shift applies to
source and target. For a finite class let \(T'\) be its extension of the trivial lattice by the target \(T_2\). Strictness gives the multiset identity
\[
 \mathrm{Jumps}_\tau(T')=\mathrm{Jumps}_\tau(T_2)\sqcup\{0\}.
\]
Thus the same boxes apply whenever they contain zero. These are the data for the integral Herr pairs.

At the complementary primes use the heights of
Lemma~\ref{lem:first-level-terminal-heights}. The local comparison
uses the full
crystalline lattices at the same marked point and the induced
models of Lemma~\ref{lem:localized-induced-wach-heights}. Its
reduced matrix on the full lattice is the divided cyclotomic
matrix computed above. The common containing lattice permits
the various integral presentations of a class. For the least
exponent comparison choose the full lattice of the characteristic-zero
eigensystem lifting the surviving weight. Its marked lines are
those of the actual residual congruence. The Wach bases may differ between comparisons. Keeping the terminal labels fixed identifies the residual extension being tested, and normalization returns its length bound to the original coefficient place.

\subsection{The type primes and the ramified auxiliary primes}
\label{subsec:special-prime-lengths}
\hypertarget{source-special-prime-lengths}{}

The unramified models leave two different sources of exceptional local behavior. A coefficient prime can equal one of the levels imposing the tame type, or it can ramify in one of the auxiliary quadratic fields. We retain the original row in both calculations.

Keep the cut, its good eigenspaces, and its attached
integral Galois representations fixed. Write \(R=r+t\)
for the full rank and normalize the length at a coefficient
prime \(\ell\) by \(v_\ell(\ell)=1\). Both calculations retain the global terms of the numerator: the auxiliary signature length \(s_\pi\), the strict cohomology sum \(V_\ell\), and the class number \(h_E\).

\subsubsection*{The type-prime hypotheses}

At a type prime the local lifting problem retains both tame characters and labeled Hodge jumps. Complete splitting allows them to be specified at one embedding at a time.

Fix a twisted localization at \(\ell=g_i\), with \(i\in\{0,1\}\), and suppose
it has a good point. All completions of \(E\) at \(\ell\)
have degree one. Each completion determines the unique
labeled embedding of \(E\) inducing that place. The CM
type gives a sign at each such embedding, reversed at
the conjugate place. In the family at this localization
the jump set has three parts: the cyclotomic jumps
\(0,\ldots,t-1\), the jumps within \(\eta+1\) of
\(c=(t-1)/2\), and the jumps in the signed outer range.
The negative sign reflects about the central point \(c=(t-1)/2\):
\[
 a\longmapsto 2c-a=t-1-a.
\]

This preserves the central interval and cyclotomic
set and reverses the signed range about \(c\).
The bound on absolute jumps is \(2b\), so this gives
a common set of at most \(R\) distinct integer jumps
for all members of the family at the place. The jumps
avoid \(-1,t\) by the window conditions.

Take \(\eta=\lfloor\eta_c\rfloor\),
\[
 s_T=\lfloor3b\rfloor+2
\]
and \(S_T=0\) in the signed outer range. Then
\(\eta+t+1<s_T\) since
\[
 \eta_c+S_c+t+2\le b.
\]

In each sign the distance of an outer jump in its
specified direction from \(c\) is at most \(s_T-c\).
For the size requirement of the paired-character
argument, put \(y=kb\), where \(k\ge5\) is the ambient rank parameter. For every fixed \(C_T\le b\)
we have \(s_T\le5b\le y\), \(R+1\le3y\) and \(y\ge5\).
Consequently
\[
 C_T\bigl((s_T+S_T)(R+1)^2+(R+1)^3\bigr)
 \le36y^4<y^8\le g_i.
\]
Both type primes lie in the same interval \([Y,2Y]\).
In particular
\[
 g_{1-i}\le2g_i\le g_i^2.
\]

These bounds and the sign calculation supply the original-field and size conditions on the cut family.
\label{cor:special-prime-lengths-1}

In the semisimplification at a retained point the
rank-\(r\) part and the \(t\) cyclotomic characters
carry the specified common inertia and weight labels.
The inertia comparison is transported along the actual
norm choice, as in Section~\ref{subsec:selected-cut}.
The tame generators and wild triviality at the auxiliary
and signature places come from the local type there.
The nontrivial inertia value at each signature place
is part of the marked data. The kill comparison gives
finiteness of the signature length and the strict
cohomology bound gives finite \(V_\ell\) at all the
required positive exponents. Thus the paired-character estimate may be applied on this cut whenever a good point occurs.
If there are no good points, the projected and intersected
lattices are both zero and the contribution is zero.
\label{cor:special-prime-lengths-2}

\subsubsection*{Lifting on a common tame type}

The zero-monodromy and jump conditions at a type place control each paired character. For the numerator we need a uniform bound on the span of the corresponding pushouts, with one exponent for all characters of the common tame type.

\begin{lemma}[A normalized bound on the tame pushout span]\label{lem:common-tame-pushout}
At a type place of degree one let the source and target of the
surjection of stable lattices be rational subquotients of sums from
one potentially crystalline family with zero monodromy. Use a
Frobenius-stable tame character set of size at most \(N=R+1\)
and at most \(N+1\) jumps in \([-u,H]\), each set containing zero.
Assume \(1\le u<H\) and \(C_T(HN^2+N^3)<\ell\) in the integral
tame lifting comparison. If the finite lifting allowance is \(\ell^e\), let \(D(T_0)\) denote the image of the pushout span modulo the finite condition for the target lattice \(T_0\). Its normalized length satisfies
\[
 \operatorname{len}D(T_0)\le2e\operatorname{rk}T_0.
\] The bound includes the torsion in
integral finite cohomology.
\end{lemma}

\begin{proof}
\proofstep{Put the source and target on one tame extension and bound the periods of character differences.}
At a completion of degree one, zero monodromy and
the tame inertia polynomial give a common finite
extension on which the family is crystalline.
One can form it using the kernel of the finite
inertia characters and an unramified extension for
the Frobenius action on them. The Artin character
of the cyclotomic block is included before taking
its inverse twist. Thus the common character set
for the twisted family contains the trivial
character, has at most \(N=R+1\) elements, and
is Frobenius stable. The dual uses the inverse
set and the same cardinality bound. The
coefficient comparison uses a tame exponent
dividing \(\ell^{(\ell-1)!}-1\).

Write the tame characters additively, so Frobenius
acts by multiplication by \(\ell\). If two labels have periods \(a,b\), Frobenius to the power \(\operatorname{lcm}(a,b)\) fixes both. Thus
\[
 \operatorname{period}(\chi-\chi')\mid\operatorname{lcm}(a,b).
\]
In distinct orbits with
lengths \(a,b\), we have \(a+b\le N\) and hence
\[
 \operatorname{lcm}(a,b)\le ab\le\tfrac14(a+b)^2\le N^2/4.
\]
In one orbit, the difference has period dividing
the common period \(a\). A nonzero orbit omits the
fixed point zero, so \(a\le N-1\le N^2/4\).
The zero difference has period one. The maximum \(A\)
over differences thus satisfies \(1\le A\le
\lfloor N^2/4\rfloor\), including differences with
zero.

The integral tame period comparison is used for a
surjection of stable lattices whose rational source
and target are subquotients of sums of the common
potentially crystalline family. It uses the
Frobenius-stable tame character set of size at most
\(N\), including zero, and a jump set with at most
\(N+1\) entries in \([-u,H]\), including zero,
with \(1\le u<H\).
The comparison uses both the bound on \(\ell\)
\[
                  C_T(HN^2+N^3)<\ell
\]
and the degree-one base. It gives the filtered
integral lift after multiplying by
\(\ell^{C_T(NA+u(N+A))}\). Here is the exponent
calculation. The Breuil--Kisin comparison with tame
descent gives an exact triple for the kernel,
source and target lattices.

Write \(A\) for the
maximum period above. The integral Frobenius
coordinates represent the finite target class
after loss \(e_W=2(N-1)A\) by an element in the
kernel of monodromy. A further factor \(\ell^u\)
puts the element, up to a Frobenius coboundary,
in the sum of the two integral coefficient
submodules used for lifting. The exact-triple
comparison lifts it with zero monodromy and loss
\[
             k_{\mathrm{lift}}
               =u(A+1)+e_W+u+(N+1)u+e_W .
\]
The integral cocycle comparison on the source
needs one more loss \(u(A+1)\). Its image on the
target is the prescribed multiple of the class
modulo an integral coboundary, by functoriality
of the Frobenius-coordinate calculation.
The sum of the four losses is
\[
                   6(N-1)A+u(2A+N+5)
                   \le 7\{NA+u(N+A)\}.
\]

In particular \(C_T\ge40\) suffices. An integral
linear retraction of the enlarged coefficients
preserves the cochain and lattice identities and
returns the lift to the original coefficients.
In our use \(u=\eta+t+1\), \(H=s_T+S_T\),
with the Tate dual on the negative side.
\label{cor:special-prime-lengths-3}

For the local cocycle length, take a finite
collection of cocycles from the cut. Write the
underlying cocycles as images of \(gE-E\), with \(E\)
a vector in a finite sum of the twisted family and
\(gE-E\) in a stable finitely generated lattice. Their
representations can be combined in one such sum
and one stable rational subrepresentation \(X\) of it.
The stable source lattices \(L\) in \(X\)
give finite cocycles \(y\) with values in \(L\),
since \(gE-E\) is a rational coboundary.

The residual maps in the presentation give equivariant
maps
\[
 \phi:L\to U/T_0,
\]
where \(T_0\) is the
target lattice in its rational representation \(U\).
The corresponding classes are represented by \(\phi y\).
The rational subrepresentation comparison puts \(X,U\),
or their Tate duals for the dual lift, in the same sums
of the common potential family. The period estimate
can therefore be applied to a surjection of lattices
between their direct sums.

Write \(K\) for the local base and \(F\) for the coefficient
field. To see the integral comparison, let \(\psi:X\to U\)
be a coefficient-field linear lift of \(\phi\) on \(L\).
Choose it on a basis of the free lattice \(L\). Its
equivariance defect takes values in \(T_0\) on \(L\).
The lattice
\[
 L_\psi=\{(a,b)\in L\oplus U:b-\psi a\in T_0\}
\]
in \(X\oplus U\) is stable for the product action and
surjects onto \(L\). If the lift with loss \(\ell^e\)
is available on this side, lift \(\ell^e[y]\) to an
integral finite cocycle of \(L_\psi\). Its second
coordinate reduces to \(\ell^e\phi y\) modulo a
coboundary in \(U/T_0\) and belongs to the finite image.

For the dual lift, use \(X'=X^\vee(1)\),
\(U'=U^\vee(1)\), and the dual lattices \(L'\) and \(T_0'\)
with respect to the Tate pairings. The transpose
\(\psi^{\rm t}:U'\to X'\) has equivariance defect in \(L'\)
on \(T_0'\), as follows by evaluating the defect of \(\psi\)
on \(L\). Use the lattice
\[
 \{(a',b')\in T_0'\oplus X':b'-\psi^{\rm t}a'\in L'\}
\]
and its projection onto \(T_0'\).

For any
\(d\in H^1_f(K,T_0')\), lift \(\ell^e d\) to a finite
cocycle \((a',b')\) in this lattice. Take the cup product with the dual cocycle first.
For \(g,h\in G_K\) its value in \((F/O)(1)\) is computed by
the following congruences in \(F(1)\):
\begin{equation}\label{eq:local-dual-lift-cup}
 \begin{aligned}
 \langle a'(g),g\psi(y(h))\rangle
 &\equiv\langle a'(g),\psi(gy(h))\rangle\\
 &=\langle\psi^{\rm t}a'(g),gy(h)\rangle
 \equiv\langle b'(g),gy(h)\rangle \pmod{O(1)} .
 \end{aligned}
\end{equation}
The first difference in \eqref{eq:local-dual-lift-cup} pairs \(a'(g)\in T_0'\) with the
equivariance defect in \(T_0\). The second pairs an element of
\(L'\) with \(gy(h)\in L\). Thus the cup cochain for
\(a'\) and \(\phi y\) is the reduction of the rational cup
cochain for \(b'\) and \(y\). The latter rational
cup product vanishes in cohomology by Bloch--Kato
orthogonality.

Integral local Tate duality identifies the
finite image in \(H^1(K,U/T_0)\) with the annihilator of
\(H^1_f(K,T_0')\). It follows that \(\ell^e[\phi y]\)
is in the finite image also in this comparison.

\proofstep{Apply finite-condition duality and count the free and torsion generators.}
The common tame type has supplied an exponent for the local lifts. To convert that exponent into a length for the span, one must count both the free generators of the finite condition and its torsion.

For the length, write
\[
 M=H^1_f(K,T_0')
\]
for the
inverse image of the rational finite subspace. It is
finite over \(O\) by local cohomology finiteness.
The quotient of the pushout span by its intersection
with the finite image injects by the cup pairing into
\(\operatorname{Hom}_O(M/\ell^e M,F/O)\).
The torsion of \(M\) has at most \(\operatorname{rk}T_0\)
generators. Indeed the coefficient exact sequence identifies this torsion module as
\[
 H^0(K,U'/T_0')/\operatorname{im}H^0(K,U').
\]
The latter rational subspace meets \(T_0'\)
in a saturated sublattice, so after quotienting by
its image the ambient divisible module is still a direct
sum of at most \(\operatorname{rk}T_0\) copies of \(F/O\).

The rank of \(M\) is
\(\dim_F D_{\mathrm{dR}}(U')/\mathrm{Fil}^0+
\dim_F H^0(K,U')\) by the finite dimension formula.
The invariant vectors inject into \(\mathrm{Fil}^0\),
so the rank is at most \([K:\mathbb Q_\ell]\operatorname{rk}T_0\).
Put \(r_0=\operatorname{rk} T_0\) and
\(e_O=\operatorname{ord}_O(\ell)\). A free generator of \(M\)
contributes \(ee_O\) to the quotient length, and a torsion
summand \(O/(\varpi^a)\) contributes \(\min(a,ee_O)\).
The two generator bounds therefore give
\[
 \operatorname{length}_O(M/\ell^e M)
 \le ee_O\bigl([K:\mathbb Q_\ell]r_0+r_0\bigr)
 =2ee_O r_0\quad(K=\mathbb Q_\ell).
\]
Its dual has the same length. Dividing by \(e_O\)
gives \(2e\,\operatorname{rk}T_0\) in the normalized
comparison. The bound holds for every finite span with
the same exponent and rank, hence for the full span.
\cite[Proposition 3.8 and Corollary 3.8.4]{bloch-kato}
\label{cor:special-prime-lengths-4}
\end{proof}

\subsubsection*{Passing the paired-character bound to the numerator}

The estimate on a span is made at one local character. To put it in the global numerator, sum it first with the actual good multiplicities and then with the normalized weights of the places above the type prime.

\begin{proposition}[Sum of type-prime lengths on the selected row]\label{prop:type-prime-row-length}
At \(\ell=g_i\), write \(M_\ell=C_*x^{C_*}\eta_c+
2\ell^{-2}+V_\ell+v_\ell(h_E)\), with the fixed \(C_*\) below.
The good multiplicities on the cut have
\[
 \sum_{w\mid\ell,\ w\text{ twisted}}
       e_w\omega_w\eta_{j,w}\le d_j(\log\ell)M_\ell.
\]
The term from \(C_*x^{C_*}\eta_c\) is \(d_j\,o(b)\)
on the two type primes together with one threshold.
\end{proposition}

\begin{proof}
\proofstep{Use the label projector before applying the common tame pushout bound.}
For each good \(\pi\), form a point-dependent central
type operator on the cut space. It projects to the
series labels of \(\pi\) at both type primes, multiplied
by the row group order \(h\). The scalar action at \(\pi\)
uses the common isobaric partition for its Hecke
character. The support \(Y_\pi\) of this operator
contains \(\pi\). The agreement of the two series labels on the support gives
\[
 \pi\in Y_\pi,\qquad Y_\pi\cap Z_i\subseteq Z_i^+.
\]
This comparison
uses the finite series projectors, their realization on
the point spaces, and the span of those spaces in
each generic eigenspace. The operator acts integrally
and kills the eigenspaces outside \(Y_\pi\).

On \(Y_\pi\) the squarefree labels exclude the two
prescribed character values. The local--global
comparison supplies zero monodromy and the same
tame inertia type throughout this family. The inertia
polynomial of the full representation includes the
rank-\(t\) scalar character block. The corresponding
tame characters are inverted at the conjugate place.
Apply the type-prime period lifting comparison to the
twisted family on the positive side and to its Tate
dual on the negative side. The dual sends a jump \(a\) to \(-a-1\). Thus in both applications the jump parameters are
\[
 (\text{lower parameter},\text{upper parameter})=(\eta+t+1,s_T+S_T).
\]
The base field is \(\mathbb Q_\ell\) at every such use.
For \(N=R+1\), put
\[
 A_0=\lfloor N^2/4\rfloor.
\]

The integral period lift has loss
\[
 e_*=C_T\{N A_0+(\eta+t+1)(N+A_0)\}.
\]
The local pushout comparison sends this to normalized
length at most \(2u e_*\) modulo the finite
cohomology condition in rank \(u\). For the rank-\(r\)
constituent use \(u=r\), and for each of the \(t\)
character constituents use \(u=1\).
The same exponent applies throughout the locus \(Y_\pi\).
\label{cor:special-prime-lengths-5}

The global marked congruence comparison on \(Y_\pi\)
now gives an element of the Hecke order supported
on \(Z_i\cup(X_i\setminus Y_\pi)\). Its value at
\(\pi\) has valuation bounded by \(4fR e_*\)
plus the signature, strict cohomology, class-number,
mark and tame allowances. The local term is
\(2f(2re_*+2te_*)\) since all completions have
degree one. The marked determinant is evaluated at the very same normalized specialization. Over all places above \(\ell\) the local rank and character terms sum with degree factor
\[
 \sum_{w\mid\ell}[E_w:\mathbb Q_\ell]=[E:\mathbb Q]=2f.
\]

In this application the row
quotients are the local pushout images. With the
finite condition imposed, the three global kernels
are the signature, unramified class-number and
strict Tate-twist kernels. For the character block,
the coefficient character after twisting is
\(\chi_\ell^{\,1-c}\) at index \(c\).
The divisible finite image is unramified at \(c=1\).
It vanishes at \(c>1\) by the finite cohomology
formula for negative Tate twists. These conditions
give the class-number and strict kernels after
the tame restrictions away from \(\ell\).
The marked and tame costs
are \(4R v_\delta\) and
\[
 t\sum_{v\in\mathcal Q_i}v_\ell(\zeta_v-1).
\]

Multiplying the Hecke element by the integral
type operator gives support on \(Z_i^+\).
Apply the lattice filtration comparison with this
product for the graded piece at \(\pi\).
The other graded pieces may use different type
operators.

Since \(A_0\le R^2\) and \(N\le2R\), we have
\[
 e_*\le8C_TR^2(R+\eta+t).
\]

The valuation of the row group order at \(g=g_i\)
is also explicit. Its factor from \(g\) is
\(f\binom r2\) by the order formula for
\(\mathrm{GL}_r(\mathbb F_g)\). The factor from
\(g'=g_{1-i}\) is at most \(fr(r+1)\), because
\(g'^j-1<g^{2j}\) and hence
\[
 v_g(g'^j-1)\le2j.
\]
Each of the factors has
exponent \(f\), one place for each conjugate pair.
Their total valuation is therefore at most
\[
 f\binom r2+f\sum_{j=1}^r2j=\tfrac12 f(3r^2+r)\le2fr^2 .
\] Thus the whole
localization length \(\eta_i\) is bounded by the
multiplicity-weighted sum, over good \(\pi\), of
\[
 4fR\,C_T'R^2(R+\eta+t)+2fr^2
   +s_\pi+V_\ell+v_\ell(h_E)+4R\,v_\delta
   +t\sum_{v\in\mathcal Q_i}v_\ell(\zeta_v-1).
 \tag{T}\label{eq:local-prime-t}
\]
The constants in this inequality can be
taken to be \(C_T=40\) and \(C_T'=8C_T\).
Here \(v_\delta\) measures the cyclotomic mark of the
localized order. Its depth is one at either type
prime, so the marked determinant gives
\[
 v_\delta\le(t-1)+v_\ell((t-1)!)\le2t.
\]

The two killing allowances give
\(s_\pi\le\ell^{-2}\) and
\[
 t\sum_{v\in\mathcal Q_i}v_\ell(\zeta_v-1)\le\ell^{-2}.
\]

The cyclotomic invariant term vanishes because the
host is unramified at the odd prime \(\ell\) and \(tr\) is odd.
For a lift of \(-1\) the scalar difference is
\[
 (-1)^{tr}-1=-2\in\mathbb Z_\ell^\times .
\] For the signature
allowance, take the character of the chosen signature of
\(\pi\) in Section~\ref{subsec:signature-selmer} and its
specialization at the auxiliary character of this coefficient
place. The specialization is chosen simultaneously for the
good signatures. The lattice comparison identifies the
rank-\(r\) piece of the same marked representation with the
specialized signature lattice. The Selmer injection and the
normalized coefficient-extension comparison of the retained
twist give finiteness and transfer the length allowance.
The full-image generator of auxiliary inertia gives the
tame allowance on this same localization. The strict
cohomology lengths for \(1\le e<t\) are the separate
finite lengths in the definition of \(V_\ell\).

The uniform bounds used in the estimate are
\[
 R\le2t,\qquad t\le\eta_c,\qquad n,t\le x^{C_1}.
\]
Hence
the first, second and sixth terms of \eqref{eq:local-prime-t} have sum at most
\[
 (256C_T'+20)nt^3\eta_c
 \le (256C_T'+20)x^{4C_1}\eta_c .
\]
Put \(C_* = 2000(C_T'+1)+5C_1\). For \(x\ge1\) and
\(C_1\ge1\) the last bound is at most
\(C_*x^{C_*}\eta_c\). Thus the uniform bound per
copy of a good character at \(\ell\) is
\[
 M_\ell=C_*x^{C_*}\eta_c+
            2\ell^{-2}+V_\ell+v_\ell(h_E).
\]
\proofstep{Weight the lengths by the actual good multiplicities and the local degree identity.}
The local inequality applies to a single good character. In the numerator each occurs with its actual eigenspace multiplicity; the bound for their sum uses the dimension of the same selected row.

Use the local comparison including the empty case
\(\eta_i=0\) when \(Z_i^+=\varnothing\).
The ties defining the twisted row identify \(\eta_{j,w}\)
with \(\sum_{i\ {\rm over}\ (j,w)}\eta_i\). Apply
\eqref{eq:retained-multiplicity} to this sum:
\[
 \eta_{j,w}\le d_j M_\ell,\qquad
 \sum_{\substack{w\mid\ell\\w\ {\rm twisted}}}
        e_w\omega_w\eta_{j,w}\le d_j(\log\ell)M_\ell .
\]
The second inequality uses the degree sum over all places
above \(\ell\). These bounds use the actual multiplicities
and require no uniform upper bound for the number of
localizations.

In particular the contribution of the first, second and
sixth terms to the numerator, summed over both rational
type primes, is at most
\[
 d_j C_*x^{C_*}\eta_c(\log g_0+\log g_1)=d_j\,o(b),
\]
since \(\log g_i=O(\log b)\), \(\log\eta_c\le x^{\epsilon/4}\)
and \(b=\exp(x^\epsilon)\). The \(o(b)\) bound per unit of rank
is uniform on the chosen windows for fixed \(\epsilon\). The division
comparison uses \(d_j=D_0\) before dividing by \(D_0\).
The group-order factor in \((T)\) is already included at
these two primes. The separate cost of returning to the
uncut length and the common budgets for strict cohomology
and class number enter the full numerator comparison.
\label{cor:special-prime-lengths-6}
\end{proof}

\subsubsection*{The two ramified auxiliary primes}\label{subsec:ramified-auxiliary-row}

Fix \(\ell\in\{q,d\}\) and let \(\ell'\) denote the other prime.
The three factors left after omitting \(\sqrt{-\ell}\) from
\eqref{eq:host-character-compositum} have a common cyclotomic modulus
prime to \(\ell\). In fact, once the lower endpoint in
Proposition~\ref{prop:quadratic-prime-pair} is large, both auxiliary
primes exceed the fixed conductor \(f_\psi\). For each \(w\mid\ell\) put
\[
 M_\ell=\operatorname{lcm}(2^{e+2},f_\psi,\ell'),\qquad
 B_\ell=K_\psi C_e(\sqrt{-\ell'}),\qquad F_w=(B_\ell)_{w|_{B_\ell}}.
\]
The completion embeds in \(\mathbb Q_\ell(\mu_{M_\ell})\). Roots of
unity of order prime to the residue characteristic generate an
unramified extension of a complete discretely valued field
\cite[Corollary~10.16]{sutherland-local-fields}. In this application
the base is \(\mathbb Q_\ell\) and \(\ell\nmid M_\ell\). Hence
\begin{equation}\label{eq:host-auxiliary-completion}
 F_w/\mathbb Q_\ell\text{ is unramified},\qquad
 E_w=F_w(\sqrt{-\ell}),\qquad e(E_w/\mathbb Q_\ell)=2 .
\end{equation}

Indeed \(X^2+\ell\) is Eisenstein over \(F_w\), since
\(v_{F_w}(\ell)=1\). Since
\(\ell\equiv3\bmod4\), the unique quadratic subfield of
\(\mathbb Q_\ell(\mu_\ell)\) is \(\mathbb Q_\ell(\sqrt{-\ell})\)
by the quadratic Gauss sum. Writing \(\chi_{\rm cyc}\) for the cyclotomic character, we obtain
\[
 [F_w(\mu_\ell):E_w]=\frac{\ell-1}{2},
 \qquad
 \chi_{\rm cyc}(I_{E_w})\bmod\ell\ \supseteq\
           (\mathbb F_\ell^\times)^2 .
\]
The size choice gives \((\ell-1)/2>tr\). For every nonzero tested exponent \(a\),
\[
 0<|a|\le tr<(\ell-1)/2.
\]
The character with exponent \(a\) cannot be trivial on the square subgroup, which has order \((\ell-1)/2\). Thus \(\nu_\ell=0\).

The source describes the CM type by the sign of
\(\sqrt{-d}\). At \(q\), the congruence on the auxiliary discriminant gives
\[
 -d\equiv1\pmod q,\qquad -d\in\mathbb Q_q^{\times2}.
\]
Thus its decomposition group preserves the sign of \(\sqrt{-d}\). At \(d\) the nonidentity inertia element reverses it.
In either case the local jump list is taken at each
embedding, allowing its sign to depend on the
embedding of \(E_w\). These primes are distinct from
the two type primes, which split completely in the host.
The integral global congruence comparison away from
type primes applies. The declaration used in this ramified comparison is
\[
 \mathcal R_\ell=\{w:w\mid\ell\},\qquad \mathcal C_\ell=\varnothing.
\]

This choice requires a bound on each row at \(d\)
with mixed embedding signs. A count of rows obtained
by taking one sign for the whole completion would
lose the role of the inertia involution at \(d\).
\label{cor:special-prime-lengths-7}

At each embedding the original jumps are
\(0,\ldots,t-1\), \(m\) central jumps and \(k-1\) jumps
in a signed interval of width \(S_c\), separated from
the center by more than \(\eta_c+t+1\).
Put \(J_\tau\) for the original outer integer jump set at
embedding \(\tau\). The Tate operation on the outer jumps is
\[
 J_\tau\longmapsto -1-J_\tau=\{-1-j:j\in J_\tau\},
\]
with
\[
 \begin{aligned}
 \max(-J_\tau-1)-\min(-J_\tau-1)
    &=\max J_\tau-\min J_\tau\le\lfloor S_c\rfloor,\\
 \min_{a\in J_\tau}|-a-1|
    &\ge\min_{a\in J_\tau}|a|-1>\lfloor\eta_c\rfloor+t .
 \end{aligned}
\]
Both the original jump list and its Tate-dual image have
absolute bound \(\lfloor2b\rfloor\). Put \(I=[-\lfloor\eta_c\rfloor-t,\lfloor\eta_c\rfloor+t]\). For each embedding \(\tau\), removal of the cyclotomic summand leaves exactly the simple central jumps of the rank-\(r\) dual, so
\[
 \#\{a\in\mathrm{Jumps}_\tau(T_2):a\in I\}=m.
\]

For the row source, a subrepresentation of a sum of
the full duals is crystalline and has its jumps in
their union. This uses strictness of the crystalline
filtration and the common jump set at the embedding.
These are the family and target hypotheses of the
integral row comparison.

The local result used here is the crystalline
row-denominator theorem with ramification two.
For a surjection of stable lattices as above with
rank-\(r\) target, let the source jumps at each
embedding lie in the central interval of radius
\(\eta_0\) together with a signed interval of width
\(S\), lying at distance greater than \(\eta_0\)
from zero. Let all jumps have absolute value at
most \(h_0\). The target has at most \(m\) jumps,
counted with multiplicity, in the central interval.
Then the normalized finite-cohomology colength
for this surjection is at most
\[
 160\bigl((D_w+1)r(2h_0/\ell+2S+
             2\log_\ell h_0+1+\eta_0)
                   +3D_w m(h_0+1)\bigr).
\]
We derive this estimate from the integral period comparisons in the
next paragraphs. The coefficient valuation throughout has
\(v(\ell)=1\). \label{cor:special-prime-lengths-8}

\paragraph{The integral maps at one cyclotomic level.}
\label{par:row-metric}
\hypertarget{source-row-metric}{}
\begin{paramdefinition}[Local field, signed intervals and integral map]\label{hyp:first-level-row}
Write \(p=\ell\), let \(K/\mathbb Q_p\) be the local field, and let
\(K_0\) be its maximal unramified subfield. Its residue degree,
ramification index and absolute degree are
\[
 n=[K_0:\mathbb Q_p],\qquad e=[K:K_0],\qquad D=[K:\mathbb Q_p]=ne.
\] In the application \(e=2\).
The calculation also applies at \(e=1\).
A first-level row has odd residue characteristic,
\[
 K\subset K_0(\mu_p),
\]
and a coefficient
DVR \(O\) containing the images of the rings of integers
under all embeddings of \(K_0(\mu_p)\). All the following lengths are
divided by \({\rm length}_O(O/pO)\). Let \(F:T_1\twoheadrightarrow T_2\)
be the given integral crystalline surjection, with \(r={\rm rank}_O T_2\).
Write \(h,S,\eta\) for the integral height, width and central radius.
If the sign at \(\tau\) is \(d_\tau\in\{+1,-1\}\), the corresponding outer jumps \(a\) satisfy
\[
 b_\tau\le d_\tau a\le b_\tau+S.
\]

The signed jump data include a central interval \([-\eta,\eta]\) and a bound \(m\) on the number of central target jumps at each embedding. Call these data a first-level row when \(\eta<b_\tau\) at every embedding.
\end{paramdefinition}

For the two ramified auxiliary primes \(p=q,d\), \eqref{eq:host-auxiliary-completion} gives \(K=K_0(\sqrt{-p})\) with \(p\equiv3\pmod4\). The quadratic Gauss sum gives
\[
 K=K_0(\sqrt{-p})\subseteq K_0(\mu_p),\qquad [K:K_0]=2 .
\]
At an unramified place \(K=K_0\). A finite coefficient extension contains all the indicated embedding images. All actual jumps lie in \([-h,h]\); clipping an empty outer interval or the central interval outside this range therefore preserves support and separation. The normalized endpoints are
\[
 \eta' = \min(\eta,h),\qquad b'_\tau=\min(b_\tau,h+1),
 \qquad \eta'<b'_\tau\le h+1 .
\]
Use the original symbols for these normalized values. Then \(H=h+S+1\) satisfies \(b_\tau+S\le H\). Strict exactness of the crystalline filtration puts the target jumps in the same support.

We describe the integral comparison used here. At full level one
there are \(n\) components indexed cyclically by the embeddings of
\(K_0\). For the tame group \(G\), one has
\[
 |G|=s=(p-1)/e,\qquad O[[\pi]]^G=O[[X]],\qquad
 X\in\pi^s O[[\pi]]^\times,\qquad
 \frac{1}{s}\sum_{g\in G}g\in O[G].
\]
The last element is the integral projection to its invariants.

Let \(E=E_c(X)\)
be the degree-\(e\) polynomial associated to the first cyclotomic
factor on component \(c\), up to a unit. Its roots are the images of
\(X\) at the first cyclotomic points. They belong to the maximal
ideal of \(O\). We use a generator \(\gamma\) of the principal
cyclotomic direction, \(u=\chi_{\rm cyc}(\gamma)\), \(v(u-1)=1\).
On the invariant ring
\[
 \gamma X=u^s X\bmod X^2
\]
and
\(\varphi X\) is a unit times \(X E^s\), using the target component
of Frobenius.

\(N_c(T)\) is the tame-invariant Wach lattice and \(P_c(T)\)
is the invariant Frobenius span from the preceding component. Both
are free over \(A\), with
\[
             E^H P_c\subset N_c\subset E^{-H}P_c .
\]
The action of \(\gamma\) on \(N_c/XN_c\) is trivial. After inverting
\(p\) this quotient gives the corresponding crystalline constants.
Frobenius identifies the constants from \(c-1\) with those used
to evaluate \(P_c\).

\begin{lemma}[Filtration test on bounded numerators]
On each labelled first-level component with logarithmic parameter
\(t\), write \(a\) for the numerator of the element \(E^{-H}a\in E^{-H}P_c\). The test is
\[
 E^{-H}a\in N_c\quad\Longleftrightarrow\quad [t^i]a\in {\rm Fil}^{H-i}_\tau\quad\text{for every first-level root and }i\ge0.
\]
\end{lemma}
\begin{proof}
At a primitive root use a logarithmic parameter
\(t\) on which \(\gamma t=ut\). On bounded-pole Frobenius coordinates,
membership in \(N_c\) is tested by requiring the coefficient of
\(t^i\) to lie in \({\rm Fil}^{-i}_\tau\) at every point of the
component. In numerator coordinates the same condition for a vector \(z=E^{-H}a\) becomes
\[
 [t^i]a\in{\rm Fil}^{H-i}_\tau\qquad(\text{each degree }i).
\]
Indeed \(E\) has a simple zero, and multiplication by a power-series unit preserves the coefficient tests for a decreasing filtration. The evaluation identifies the rational constant
space with the labeled de Rham space using the crystalline comparison;
lifts of the same constant give the same constant value, and the
kernel on the preceding integral lattice is \(XN_{c-1}\).
One uses each embedding over \(c\), including both at ramification
two.

In an adapted de Rham basis of jumps \(a_j\), the completed
filtered lattice is generated by \(t^{-a_j}d_j\), which gives
the coefficient test. Its integral use on already bounded
numerators follows by division by the distinguished factors of
\(E\): if a remainder vanishes at the full set of roots with the
specified multiplicities over the coefficient field, the integral
monic remainder is zero. The height inclusions supply the initial denominator.
\label{cor:row-metric-1}
\end{proof}

We use the integral Wach comparison,
étale exactness, and the finite Herr class criterion for these
lattices. The finite criterion sends a tame-invariant restricted
cocycle pair \((w,z)\), with \(w\in E^{-H}P_c\) and \(z\) sufficiently
deep in \(N\) as below, to a finite class, naturally for \(F\).
The extension comparison sends the connecting pair of a lift of the
unit in a crystalline extension to its integral extension class
up to the fixed unit of the trivial frame.
\label{cor:row-metric-2}

Put
\[
 a=\lceil 2H/s\rceil,\quad M=\lceil4H/s\rceil,\quad
 A_i=i+v_p(i!),\quad \rho=A_{2a-1},\quad \sigma=A_{M-1}.
\]
\begin{proposition}[Integral comparison at the first cyclotomic level]
\label{prop:one-level-integral-row-maps}
For the local field, signed heights and maps in Definition~\ref{hyp:first-level-row},
the map on tame-invariant Wach lattices satisfies
\[
 p^\rho N_c(T_2)\subseteq F(N_c(T_1)).
\]
A first coordinate invariant modulo \(XN_c\) can be multiplied
by \(p^\sigma\) and corrected to become invariant modulo
\(X^MN_c\). It then has a
unique second coordinate in \(X^MN\). The two filtered coordinate
maps \(\alpha,\beta\) have composite \(p^{2\sigma}\) modulo the
specified error modules and preserve the tests modulo \(X^j\)
for \(1\le j\le M\). The same exponents work on the source,
target and unit extension.
\end{proposition}

\begin{proof}
\proofstep{From the tame average to a constant cokernel annihilator.}
First \(p^\rho N_c(T_2)\subset F N_c(T_1)\). For the initial parameter
bound use the ordinary Wach cokernel computation of
Section~\ref{subsec:ramified-lift} on the full cyclotomic field at level
one. Under the étale and Wach identification after restriction to
\(K_0(\mu_p)\), the underlying full-level map is the original map
with the tame action forgotten. The two height bounds of size \(H\)
give \(\pi^{2H}\) in its cokernel annihilator by extracting successive
Frobenius coordinates of a polynomial in \(1+\pi\).

If \(y\) belongs
to the tame-invariant target, \(X^a y\) therefore has a preimage in
the full-level component, since \(sa\ge 2H\). If \(Fz=X^a y\) and \(T_{\rm tame}\) is the tame group, then
\[
 z_0=|T_{\rm tame}|^{-1}\sum_{g\in T_{\rm tame}}gz,\qquad
 Fz_0=|T_{\rm tame}|^{-1}\sum_{g\in T_{\rm tame}}g(X^ay)=X^ay.
\]
The group order is a unit, so this gives an integral invariant preimage. This proves
\(X^a N_c(T_2)\subset F N_c(T_1)\).

On \(X^iN_c/X^{i+1}N_c\) the operator \(\gamma\) is the scalar
\(u^{si}\). This is true on the corresponding filtration of the
image as well, by equivariance. On \(X^ay\) for \(y\in N_c(T_2)\),
apply the two polynomials with roots \(u^{si}\) for \(0\le i<a\)
and \(a\le i<2a\). Both results lie in \(X^aFN_c(T_1)\), by the
image filtration for the first and the target filtration for the
second.

Put \(r_i=u^{si}\) and
\[
 L_i(Z)=\prod_{j\ne i}(Z-r_j)/(r_i-r_j)
\]
on the \(2a\) nodes.
The sum of \(L_i\) for \(i<a\) vanishes on the \(a\) upper
nodes, and that for \(i\ge a\) vanishes on the lower nodes.
Also \(\sum_i L_i=1\). The bound below gives \(p^\rho L_i\in O[Z]\),
so monic division by the products of the opposite roots
gives an identity
\[
         p^\rho=A(Z)\prod_{i<a}(Z-r_i)
                    +B(Z)\prod_{i\ge a}(Z-r_i),
                    \qquad A,B\in O[Z].
\]
Applying it at \(\gamma\) proves the claimed inclusion.

For distinct nodes \(v_p(r_i-r_j)=1+v_p(i-j)\), so the
denominator of the \(i\)-th interpolation polynomial costs
\[
 \sum_{j\ne i}v_p(r_i-r_j)
 =2a-1+v_p\bigl(i!(2a-1-i)!\bigr)
 \le 2a-1+v_p\bigl((2a-1)!\bigr)\le\rho .
\]
The first inequality uses the integrality of the binomial coefficient.
Cancel \(X^a\) to obtain the assertion. The same computation applies
to the surjection from an extension by \(T_2\) to the unit object,
whose extra jump is zero.

\proofstep{Raising the precision and completing the Herr pair.}
A second use of these scalars lifts invariance. If
\(w\in E^{-H}P_c\) satisfies \((\gamma-1)w\in XN_c\), one can
correct \(p^\sigma w\) by an element of \(XN_c\) so that its
difference under \(\gamma\) belongs to \(X^MN_c\). If \(R_i=(\gamma-1)w_i\in X^iN_c\), set
\(w_{i+1}=(u^{si}-1)w_i-R_i\). The scalar action on the \(i\)-th
graded piece gives the next precision:
\[
 (\gamma-1)w_{i+1}=(u^{si}-\gamma)R_i\in X^{i+1}N_c .
\]
The product over \(1\le i<M\) is a unit times \(p^\sigma\).
Also an element of \(XN_c\) invariant to order \(M\) already lies in
\(X^MN_c\), by its first nonzero graded coefficient and
\(u^{si}\ne1\).

On \(X^MN\) the operator \(1-\varphi\) is invertible
by its geometric series. Indeed \(sM-H\ge1\) leaves a factor \(E\)
after applying Frobenius and clearing its height denominator, so
iteration goes to zero in the complete \((p,\pi)\)-adic topology.
Consequently a first Herr coordinate invariant to order \(M\)
has a unique second coordinate in \(X^MN\) solving
\[
 (\varphi-1)z=(\gamma-1)w.
\]
Changing its first coordinate by
\(X^MN\) changes the pair by a coboundary, by the same inverse.

\proofstep{Transporting the horizontal bases by Frobenius.}
For later use we describe the coordinate maps together with their
precision. Let \(L=N_{c-1}/XN_{c-1}\). In an integral basis of
\(N_{c-1}\) lifting a basis of \(L\), construct a matrix \(U(X)\)
with \(U(0)=1\) so the corrected basis is \(\gamma\)-invariant
modulo \(X^M\) over the fraction field.

If the residual first occurs
in degree \(i\), the coefficient equation divides that residual by
\(1-u^{si}\). Thus the degree \(i\) coefficients of \(U\) have
denominator at most \(p^{A_i}\) for \(i<M\). The inverse has the
same bound, since
\[
 \sum A_{i_j}\le A_{\sum i_j}
\]
by factorial
divisibility. Frobenius transport of \(p^\sigma U\) and
\(p^\sigma U^{-1}\) furnishes two integral maps between the model
\(E^{-H} A\otimes_O L\) and \(E^{-H}P_c\), modulo
\(X^M E^{3H}\) times the respective undivided modules. Their composites are \(p^{2\sigma}\) in these quotients because
\[
 \frac{\varphi(X)^M}{E^H}\in X^M E^{sM-H}A
                         \subseteq X^M E^{3H}A
 \qquad(sM\ge4H).
\]
They commute in the quotients with \(\gamma\) and lattice maps.
For the latter property the map in horizontal bases agrees through
degree \(M-1\) with its constant matrix, by equivariance and
\(u^{si}\ne1\). Denote the forward and backward maps by \(\alpha,\beta\).

\proofstep{The filtration tests at all first-level roots.}
In the first-level expansion the corrected basis specializes to its
constant through orders below \(sM\) before division by \(E^H\).
Indeed the pulled-back invariant equation forces all intermediate
positive coefficients to vanish, since the degree \(i\) cyclotomic
eigenvalue is \(u^i\ne1\). The precision and jump-height bounds give
\[
 i\ge sM\ \Longrightarrow\ i\ge4H\ \Longrightarrow\ i\ge2H.
\]
All first-level filtration tests for numerator coefficients in those degrees are automatic. Hence both maps carry the filtered model lattice,
denoted \(N^{\rm mod}\), into \(N_c\) and back, respectively, modulo
the indicated error modules. They preserve the tests for membership
in \(X^jN\) for \(1\le j\le M\) in the forward or backward direction
from a member, by linearity in the invariant coefficient ring and
the inclusion of the error in \(X^MN\). In particular an invariant
modulo \(XN\) has an invariant partner. These constructions apply
with the same \(H,M,\rho,\sigma\) to the source, target and unit
extension.
\label{cor:row-metric-3}
\end{proof}

\paragraph{The number of exceptional coordinates.}
\label{par:row-restricted}
At component \(c\) of the target, use the lattice
\(L=N_{c-1}/XN_{c-1}\) of Section~\ref{par:row-metric}.
Identify \(L\otimes_O{\rm Frac}(O)\)
with the crystalline constant space and pull every labeled filtration
back to it. At a plus embedding, \(C_\tau={\rm Fil}^{b_\tau}_\tau\)
has codimension at most \(m\). At a minus embedding,
\(B_\tau={\rm Fil}^{1-b_\tau}_\tau\) has dimension at most \(m\).
These inequalities follow by counting the central jumps, including
their multiplicities, on the two sides of the respective cuts.
The filtered spaces for distinct roots belonging to the same embedding
give the same test on the crystalline constants.

\begin{proposition}[A bounded set of exceptional coordinates]\label{prop:row-exceptional-coordinates}
Assume \(\operatorname{codim} C_\tau\le m\) at the plus
embeddings and \(\dim B_\tau\le m\) at the minus embeddings,
with at most \(e\) embeddings in total. There is a direct summand \(B\) with complement \(L_{\rm main}\) such that
\[
 L=L_{\rm main}\oplus B,\quad L_{\rm main}\subseteq\bigcap_\tau C_\tau,\quad B_\tau\subseteq B\otimes K\quad(\tau\text{ minus}),\quad \operatorname{rk}_O B\le em,
\]
where \(K\) is the fraction field and the intersection is over the plus embeddings.
The image \(I_B\) of the invariant residual coordinates defined
below is free of rank at most \((e+1)\operatorname{rk}_O B\).
\end{proposition}

\begin{proof}
\proofstep{Splitting off the exceptional directions.}
There is an integral splitting \(L=L_{\rm main}\oplus L_{\rm err}\)
such that \(L_{\rm main}\) lies in every \(C_\tau\) and
\(L_{\rm err}\otimes{\rm Frac}(O)\) contains every \(B_\tau\), with
\({\rm rank}\,L_{\rm err}\le em\). Let \(K=\operatorname{Frac}(O)\), \(C=\bigcap_{\tau,+}C_\tau\),
and write \(L=(C\cap L)\oplus D\) using saturation. Define
\[
  L_{\rm err}=L\cap\Big(D_K+\sum_{\tau,-}B_\tau\Big),
  \qquad
  \operatorname{rk}L_{\rm err}\le \operatorname{codim} C
                          +\sum_{\tau,-}\dim B_\tau\le em .
\]
Choose \(L_{\rm main}\) complementary to \(L_{\rm err}\cap C\)
in \(C\cap L\). Then \(L=L_{\rm main}\oplus L_{\rm err}\)
since \(D\subseteq L_{\rm err}\). All splittings used here exist over the DVR because
a saturated submodule of a finite free module has free quotient.
In particular the projection to \(L_{\rm main}\) has integral entries
and kills all the \(B_\tau\).

\proofstep{The polynomial that separates the main coordinates.}
Consider a model coordinate \(z\in E^{-H}A\otimes L\) with
\[
 (\gamma-1)z\in XN^{\rm mod},
\]
for example \(\beta(w)\) with \((\gamma-1)w\in XN_c\). At a primitive root the coefficient of
\(t^i\) in \(z\), for \(i\ne0\), belongs to \({\rm Fil}^{-i}_\tau\).
Indeed the coefficient of its difference is multiplied by \(u^i-1\),
nonzero in the fraction field, and multiplication by \(X\), a unit
at this root, preserves the filtration inequalities. On projecting
to \(L_{\rm main}\) it follows that the coefficients vanish in degrees
\(i<-b_\tau-S\) at plus embeddings, and in nonzero degrees
\(i<b_\tau\) at minus embeddings. In the latter case the coefficient
is in \(B_\tau\), which the projection kills.

Let \(\Omega\) be the union over these embeddings of the integer
intervals \([-b_\tau-S,-b_\tau)\) at plus and
\([b_\tau,b_\tau+S)\) at minus. The exceptional degrees determine the polynomial
\[
 P(Z)=\prod_{i\in\Omega}(Z-u^i).
\]
The operator \(P(\gamma)\) preserves \(E^{-H}A\) since \(\gamma E/E\)
is a unit. In the Laurent parameter at that root, write the projection
to the main summand as \({\rm pr}_{\rm main}z=\sum_i z_i t^i\).
Since this projection commutes with \(\gamma\),
\[
 [t^i]{\rm pr}_{\rm main}P(\gamma)z=P(u^i)z_i,\qquad
 P(u^i)=0\quad(i\in\Omega).
\]
Below the indicated intervals the projected coefficients already vanish,
with the possible exception of the constant at a minus embedding.
The polynomial removes the degrees inside the intervals. Write the main part of \(P(\gamma)z\) in the integral
basis obtained above as \(\sum f_j\otimes l_j\).

At plus points \(f_j\) has order at least \(-b_\tau\); at minus
points
\[
 f_j=x_{\tau,j}+O(t^{b_\tau+S})
\]
for a constant
\(x_{\tau,j}\) in the coefficient field (with the value at the
chosen root understood). Also \((\gamma-1)f_j\in X E^{-H}A\).
The latter inclusion follows by commuting projection and the polynomial
with \(\gamma\) and using the stability of \(XE^{-H}A\).
Each simple tensor is invariant modulo \(XN^{\rm mod}\).
For its difference divided by \(X\), the required filtration at
plus points follows from \(l_j\in{\rm Fil}^{b_\tau}\) and the lower
order just proved. At minus points the constant cancels and all
of \(L\) is in \({\rm Fil}^{-b_\tau-S}\).
Thus the residual part is also invariant modulo \(XN^{\rm mod}\).

Finally
\[
 P(\gamma)z\equiv P(1)z\bmod XN^{\rm mod},
\]
by stability
and \((Z-1)\mid(P(Z)-P(1))\).
The scalar cost to replace \(P(1)\) by a power of \(p\) is at most
\[
             \kappa=e\{S+v_p(S!)+\lfloor\log_p H\rfloor\}.
\]
Each interval has \(S\) nonzero integers of magnitude at most \(H\).
For one interval \(I\), the principal-unit formula therefore gives
\[
 \sum_{i\in I}v(1-u^i)
 = |I|+\sum_{a=1}^{\lfloor\log_p H\rfloor}
                  \#\{i\in I:p^a\mid i\}
 \le S+v_p(S!)+\lfloor\log_p H\rfloor .
\]
There are at most \(e\) intervals; overlaps can only reduce
the sum over their union. An empty union has cost zero.

\proofstep{Counting the residual directions.}
We need the number of integral residual classes, including the
coefficient from the ring variable. For \(B=L_{\rm err}\), consider
the image \(I_B\) in \((E^{-H}A\otimes L)/XN^{\rm mod}\) of those
\(z\in E^{-H}A\otimes B\) with \((\gamma-1)z\in XN^{\rm mod}\).
The two assertions to prove are torsion-freeness and the bound
\[
 \operatorname{rank}_O I_B\le(e+1)\operatorname{rank}_O B.
\]
Over the coefficient DVR they imply the asserted freeness once finiteness is known.
The ambient quotient is finite over \(O\): \(N^{\rm mod}\) contains
\(E^HA\otimes L\) by the jump bounds, and \(XE^{2H}\) is a unit times
a distinguished polynomial.

The ambient quotient is also torsion-free. If \(a z\in
XN^{\rm mod}\) for \(0\ne a\in O\), the constant coefficient of the
integral numerator of \(a z\) at \(X=0\) vanishes. Thus \(z/X\) still
has an integral numerator. At every primitive root the coefficient test for \(az/X\) lies in a
filtration space over \(K=\operatorname{Frac}(O)\). For each degree \(i\),
\[
  a[t^i](z/X)\in {\rm Fil}^{-i}_\tau
      \quad\Longrightarrow\quad [t^i](z/X)\in {\rm Fil}^{-i}_\tau .
\]
Together with integrality at zero this gives \(z/X\in N^{\rm mod}\).
The ambient quotient is consequently finite torsion-free over the DVR.
Its submodule \(I_B\) is free.

For the rank, use the finite filtration starting with \(E^{-H}A\otimes B\).
Multiply successively by \(X-r_j\), \(2H\) times for each of the
\(e\) roots \(r_j\) of \(E\), and finally by \(X\). The last term
is contained in \(XN^{\rm mod}\). Every grade has rank at most
\({\rm rank}\,B\), by division with remainder at its root, including
the last root zero. The filtration is \(\gamma\)-stable. At \(r_j\)
the local parameter is a unit times the cyclotomic logarithm. To see
this, write
\[
 X=\pi^s v(\pi)
\]
with \(v\) a unit.

At
\(\pi=\zeta-1\) its derivative is
\(\pi^{s-1}(sv(\pi)+\pi v'(\pi))\ne0\), since \(s\) is a \(p\)-adic
unit. The \(p-1\) simple roots in the ordinary parameter give \(e\)
simple roots in the invariant ring: the tame orbits have size \(s\)
and \(E(X(\pi))\) is a unit times the cyclotomic polynomial at
\(1+\pi\). The identical multiplicity calculation also shows that
different orbits give distinct roots. Alternatively \(E\) has
degree \(e\) and its derivative is nonzero at each image, and division
over each fibre of \(X-r_j\) in the ordinary disc has degree \(s\).

On the grade with preceding multiplicity \(a\) at \(r_j\), \(\gamma\)
acts by \(u^{a-H}\); factors coming from other roots are units there.
Since \(u\) has infinite order, the invariant-grade test at this root is
\[
 u^{a-H}=1\quad\Longleftrightarrow\quad a=H.
\]
The
grade at zero contributes at most one more copy of \(B\).
Over the fraction field filter the image in the quotient by
\(XN^{\rm mod}\). Its graded pieces are quotients of the original
grades, and a fixed vector has a fixed leading term. The action in the graded factors at \(r_j\) gives
\[
 \operatorname{rk}_O I_B
 \le \sum_{j=1}^{e}\sum_{\substack{0\le a<2H\\u^{a-H}=1}}
           \operatorname{rk}_O B+\operatorname{rk}_O B
 \le (e+1)\operatorname{rk}_O B\le e(e+1)m .
\]
Thus an \(O\)-basis of \(I_B\) furnishes a single choice of
representatives at this component for all target classes. Taking a different basis for each input
would give no bound on the number of exceptional generators.
\label{cor:row-coordinates-1}
\end{proof}

\paragraph{Lifting an outer coordinate through a quadratic component.}
\label{par:row-mixed}
\hypertarget{source-row-mixed}{}

\begin{lemma}[Lifting the main tensors]\label{lem:mixed-row-tensors}
For the main tensors on a one-sign component, or on a mixed
component with \(e=2\), put
\(y_c=\lceil\eta+1/2+\log_pH\rceil\).
For each such tensor, let \(y\in N_c(T_2)\) represent its image
under the forward comparison \(\alpha\) of Section~\ref{par:row-metric},
modulo \(XN_c(T_2)\). There exists \(v\in N_c(T_1)\) with
\[
 (\gamma-1)v\in XN_c(T_1),\qquad Fv\equiv p^{\rho+\max(2\sigma,y_c)}y\pmod{XN_c(T_2)}.
\]
\end{lemma}

\begin{proof}
Let \(f\otimes x\) be one of the simple tensors above. We require an
invariant in the source with the same image as a power of \(p\) times
its image under \(\alpha\) modulo \(XN_c(T_2)\). Two operations suffice.
First, if an invariant model tensor has an invariant source partner
under \(F\), apply \(\alpha\) and use naturality. Changing the partner by \(u\in X^M E^{3H}P_c\) preserves its cut class because
\[
 X^M E^{3H}P_c\subseteq XN_c.
\]

Second, if a model
tensor belongs to \(N^{\rm mod}_2\), its image under \(\alpha\) belongs to
\(N_c(T_2)\), hence has an integral lift after multiplication by
\(p^\rho\). Any such lift is invariant modulo \(XN_c(T_1)\).
Here a model tensor with a scalar coefficient and a preceding constant
can be lifted by multiplying the scalar by a lift of that constant,
provided the resulting tensor meets the invariance test. Actual
preceding constants have lifts after \(p^\rho\) by the same Wach
cokernel bound on the preceding component.

\proofstep{Components with one sign.}
At a component with only plus tests, the product of each
coefficient of \(f\) with \(x\) lies in the required
filtration, since \(f\) starts in degree \(-b_\tau\) or
later and \(x\in\operatorname{Fil}^{b_\tau}\). Thus
\[
 f\otimes x\in N^{\rm mod}_2.
\]

At a component with only minus tests, choose a preceding constant
\(x_1\) with
\[
 F x_1=p^\rho x,
\]
and take \(f\otimes x_1\).
The source conditions are
\[
  E^H f\in A,\qquad
  (\gamma-1)f\in XE^{-H}A,\qquad
  \operatorname{ord}_{\tau}\big((\gamma-1)f/X\big)
      \ge b_\tau+S .
\]
The last inequality suffices because the entire source constant
space is in \(\operatorname{Fil}^{-b_\tau-S}\) at each minus test.

\proofstep{The two roots of the quadratic component.}
For opposite signs assume \(e=2\). Write \(r_U,r_D\) for the
two parameter values at the plus and minus tests, and put
\[
 y_U=X-r_U,\ y_D=X-r_D.
\]
Division by \(X-r\) at any such
value is integral since \(r\) belongs to the maximal ideal and
\(O\) is complete. Over the fraction field the changes to \(y_U\) and \(y_D\) have invertible linear term at their respective roots. Each preserves the order of vanishing and transports the filtration test. In particular equal \(X\)-values impose identical scalar Taylor conditions.

If a plus test has the minus
root as well, the high vanishing of \((\gamma-1)f/X\) at the
minus root also suffices at that test, since its positive degree
is greater than \(\eta\). This gives the preceding-lift argument
when every plus test shares a minus root. Otherwise there are just
two distinct roots, the minus tests at \(r_D\) and the remaining
plus tests at \(r_U\), with the same labeled filtration at all
points over the corresponding embedding.

Write
\[
 b=b_U,\qquad B=b_D+S,\qquad
 E=a_2 y_U y_D,\qquad v(a_2)=0,\qquad d_*=r_D-r_U.
\]
Both roots have valuation \(s/(p-1)=1/2\). Also
\(v(d_*)=1/2\): at \(r_U\) differentiate the cyclotomic identity
\(Q(\pi)={\rm unit}\cdot E(X(\pi))\), where
\(Q=((1+\pi)^p-1)/\pi\). The derivative of \(Q\) at a primitive
point has valuation \(1-1/(p-1)\), and \(X'(\pi)\) has valuation
\((s-1)/(p-1)\). Hence
\[
 v(E'(r_U))=v(Q'(\pi))-v(X'(\pi))=1-\frac1{p-1}-\frac{s-1}{p-1}=\frac12.
\]
This accounts for the half power in the following allowance.
Let \(y_c=\lceil\eta+1/2+\log_p H\rceil\).

\proofstep{Nearby breaks.}
Suppose first \(b\le B+2\eta\). The numerator \(a=E^H f\) is
divisible by \(y_U^{H-b}\) in \(A\). At the down root the series
of \(f\) equals a constant \(x_0\) modulo degree \(B\).
Consequently the coefficients of \(a/y_U^{H-b}\) there, through
degree \(H+B-1\), agree with those of
\(x_0 a_2^H y_D^H(y_D+d_*)^b\). They are integral in \(y_D\)
by division in \(A\). In particular
\[
           x_0a_2^H \binom{b}{j} d_*^{\,b-j}\in O
                   \quad (0\le j<B,\ j\le b).
\]

If \(b<B\) take \(j=b\). Otherwise take \(j=B-1\); then
\(b-j\le2\eta+1\) and
\[
 v_p\binom{b}{j}\le\lfloor\log_p b\rfloor
\]
by the factorial formula, a difference of at most one in each
power-\(p\) count. In either case \(p^{y_c}x_0\in O\).
Subtract this constant from \(p^{y_c}f\). The remainder times
\(x\) belongs to \(N^{\rm mod}_2\). At the minus points its
order is at least \(B\); at the plus points its order is at
least \(-b\) and the target constant is in \({\rm Fil}^b\).
The constant term itself lifts by an arbitrary preceding lift,
since a constant is fixed by \(\gamma\). Applying the two operations
gives an invariant lift in the source Wach lattice at exponent \(\rho+y_c\).

\proofstep{Separated breaks.}
Suppose now \(B+2\eta<b\). A division by \(d_*^{b-B+1}\) could
cost an amount proportional to the outer height. Instead form the
auxiliary numerator
\[
 a_0=y_U^{H-b}y_D^{H+B}.
\]
The tensor
\(a_0 x/E^H\) belongs to \(N^{\rm mod}_2\) by its orders at both
roots. Apply \(\alpha\), lift its image in the target Wach lattice after \(p^\rho\)
in \(N_c(T_1)\), and apply \(\beta\).

Write the result with denominator
\(E^H\) as \(\sum_j a_j v_j/E^H\), \(a_j\in A\), where \(v_j\)
are integral preceding constants of the source. The sum in every
numerator coefficient of degree \(i\) at a root lies in
\({\rm Fil}^{H-i}_\tau\). On applying \(F\) its expansion in degrees
through \(H\) agrees with
\[
                  p^{\rho+2\sigma}a_0 x .
\]
Indeed the forward and backward composites give this scalar and the
coordinate modulus, after multiplying by \(E^H\), contains the
factor \(E^{4H}\) in every error term. Integral bases on the preceding
component give a representation in this common constant space for
both filtrations.

Expand each \(a_j\) at \(r_D\) through degrees \(<H-\eta\), with
coefficients in \(O\), and remove this polynomial in \(y_D\).
In degree \(i<H-\eta\) the corresponding sum against \(v_j\)
belongs to
\[
 {\rm Fil}^{H-i}_D=0,
\]
since the largest source
jump there is at most \(\eta\). These polynomial coefficients
therefore give zero sums against both \(v_j\) and \(Fv_j\) at the
plus root also, by the common crystalline constant identification.
Divide the remaining functions by \(y_D^{H-\eta}\), obtaining
integral \(a'_j\). At \(r_U\) this divisor is a unit over the
fraction field and the source coefficient condition
\({\rm Fil}^{H-i}_U\) persists. The new target expansion agrees
through degree \(H\) with
\[
       p^{\rho+2\sigma}y_U^{H-b}
                  (y_U+r_U-r_D)^{B+\eta}x .
\]

Take its coefficient of degree \(j_*=H-b+B+\eta\) in \(y_U\).
The assumption gives \(j_*<H-\eta\); changing from \(t\) to \(y_U\)
preserves the filtration test and precision, by triangular substitution.
This is the leading coefficient of the displayed polynomial. Hence
the same integral linear combination of the \(v_j\) maps to
\(p^{\rho+2\sigma}x\) and belongs to
\({\rm Fil}^{H-j_*}_U\subset{\rm Fil}^{\eta+1}_U
={\rm Fil}^b_U\). The equality uses the empty interval of jumps
strictly between \(\eta\) and \(b\). This combination supplies a
preceding lift compatible with the plus test. Keeping \(f\) gives
an invariant model lift at both signs, hence an invariant Wach-lattice
lift after \(\alpha\). Altogether a single exponent
\[
                       y_1=\rho+\max(2\sigma,y_c)
\]
serves every main tensor on every component, including the one-sign
cases.  \label{cor:row-mixed-1}
\end{proof}

\paragraph{From the Herr pairs to normalized colength.}
\label{par:row-colength}
On each component use the module \(I_B\) of exceptional coordinates
from Proposition~\ref{prop:row-exceptional-coordinates}. Choose lifts
of an \(O\)-basis of \(I_B\).
Apply \(\alpha\) to each lift, choosing a representative invariant
modulo \(XN_c\). By the order-raising argument in
Section~\ref{par:row-metric}, its \(p^\sigma\)-multiple can be
corrected modulo \(XN_c\) to an invariant to order \(M\).
Complete it by the unique deep second Herr coordinate. This gives
at most \(e(e+1)m\) fixed finite classes on this component.
Write \(H^1_f(K,T)\) for the inverse image of the rational
finite condition in integral cohomology, including torsion.
Write \(J\) for the module generated by all these classes on all
components, in the quotient
\[
              C_F=H^1_f(K,T_2)/F_*H^1_f(K,T_1).
\]
There are at most \(ne(e+1)m=D(e+1)m\) generators of \(J\).

\begin{proposition}[Normalized colength of the finite local condition]\label{prop:local-row-colength}
Let \(C_F\) be the cokernel on integral finite cohomology
of the map of Definition~\ref{hyp:first-level-row}. With
\(J\) the span of the exceptional classes constructed from
Proposition~\ref{prop:row-exceptional-coordinates}, the
normalized length satisfies
\[
 \frac{\operatorname{length}_O C_F}{\operatorname{length}_O(O/pO)}
       \le 2(D+1)rY+D(e+1)m\,18h .
\]
For \(p\ge3\), \(h\ge1\) and \(e\in\{1,2\}\), the common exponent
\(Y\) below is at most \(80(eh/p+eS+e\log_p h+1+\eta)\).
The additional \(18h\) allowance occurs on at most \(D(e+1)m\)
exceptional directions.
\end{proposition}

\begin{proof}
\proofstep{Lifting modulo the fixed exceptional classes.}
Let \((w,z)\) be a target Herr pair with \(w\in E^{-H}P_c\)
invariant to order \(M\) and \(z\in X^MN\). In the model
\(\beta(w)\), apply \(P(\gamma)\) and use the splitting into
main and residual coordinates. Expand the residual modulo
\(XN^{\rm mod}\) in the chosen basis of \(I_B\).
Return from the model to the Wach lattice using \(\alpha\). The integral multiplier and the cost of the two comparison maps yield
\[
 \frac{p^\kappa}{P(1)}\,P(1)p^{2\sigma}w=p^{2\sigma+\kappa}w.
\]

At exponent \(y_1\) every main tensor on its right lifts from
the source modulo \(XN_c\), by Section~\ref{par:row-mixed}.
The residual terms are integral combinations of the fixed
representatives. Multiply by \(p^\sigma\), raise all these
invariants to precision \(M\) and complete the source pair.
The difference of the first coordinates is invariant to order \(M\)
and lies in \(XN_c\). It lies in \(X^MN_c\) by uniqueness of
invariant lifting. Naturality and the deep coboundary criterion now give
\[
              p^{\,3\sigma+\kappa+y_1}[(w,z)]\in J
                       \quad\hbox{in } C_F .
\]
In particular the precision raising costs one \(\sigma\) for the
entire relation, irrespective of its number of terms.

\proofstep{Reduction of a finite class to the tested form.}
It remains to put a general finite target class into this form.
Its extension of the unit by \(T_2\) is crystalline over the fraction
field. Lift \(p^\rho\) times the unit generator in each component
to \(x_c\) in the extension Wach lattice. Its cyclotomic difference
lies in \(XN_c(T_2)\) since the quotient generator is fixed.
Here one needs the kernel intersection: the integral étale lattice
sequence is exact and the subrepresentation has the induced Wach
lattice. The same intersection holds in the Frobenius span, as the
coefficient ring is free over its Frobenius image. In the bounded
pole module it follows after multiplying by \(E^H\).
The graded correction produces \(x'_c\) with
\[
 x'_c-p^\sigma x_c\in XN_c(T_2),\qquad (\gamma-1)x'_c\in X^M N_c(T_2).
\]

The correction uses the graded scalar calculation on that kernel.
Its projection is still \(p^{\rho+\sigma}\) times the unit.
Consequently
\[
                        w_c=\varphi x'_{c-1}-x'_c
\]
belongs to the kernel in \(E^{-H}P_c\) and is invariant to order
\(M\), since \(\gamma\) commutes with \(\varphi\) and \(\varphi\)
preserves \(X^MN\). Choose the unique \(z_c\in X^MN\) with
\((\varphi-1)z_c=(\gamma-1)w_c\). On summing in \(c\), uniqueness
gives \(\sum z_c=(\gamma-1)\sum x'_c\), since both sides are deep
and have the same Frobenius difference. Thus
\(\sum(w_c,z_c)\) represents \(p^{\rho+\sigma}\) times the original
class by the connecting map, allowing its fixed unit. We have proved
\[
       p^Y C_F\subset J,\qquad
       Y=2\rho+4\sigma+\kappa+\max(2\sigma,y_c).
       \tag{R0}\label{eq:tag-r0}
\]

\proofstep{Bounding the exceptional submodule.}
A uniform exponent on \(J\) follows from the coarser full-level lift.
Over \(K_0(\mu_p)\) the two height bounds are \(h,h\), and the jumps
of sources and kernels remain in \([-h,h]\). Restriction and
corestriction cost the degree \(s\), a unit in \(O\).
In the full-level computation of Section~\ref{subsec:ramified-lift}
the two Wach losses are at most \(4h+v_p((4h)!)\) each.

The
kernel correction has loss at most
\[
 h+v_p(h!)+\sum_{j\in\mathcal J}\bigl(1+v_p(j)\bigr),
 \qquad
 \mathcal J\subset[-h,2h]\cap\mathbb Z,\quad 0\notin\mathcal J .
\]
Indeed the shift of the jump set uses degrees \(0,\ldots,h\) and
the positive degrees for the trivial component are at most \(h\);
zero eigenvalue differences are discarded in the rational coboundary
test. Even using the full set of nonzero signed integers in this
interval gives valuation sum at most \(v_p(h!)+v_p((2h)!)\).
Since \(p\ge3\), \(v_p(j!)\le j/2\). The total exponent is at most
\(18h\). Thus \(p^{18h}\) kills \(J\).

For clarity about the other rank factor, the rational Bloch--Kato
dimension formula bounds \(\dim H^1_f(K,V_2)\) by
\[
 \dim H^0(K,V_2)+\dim(D_{\rm dR}(V_2)/{\rm Fil}^0) \le (D+1)r.
\]
For a uniformizer \(\varpi\) of \(O\), the
cohomology sequence of multiplication by \(\varpi\) exhibits
\(H^1(K,T_2)[\varpi]\) as a quotient of \(H^0(K,T_2/\varpi T_2)\).
Its residue dimension is at most \(r\). Local integral cohomology
is finite type by local Tate finiteness. Its torsion and the
torsion of its finite-condition submodule therefore need at most
\(r\) generators by the structure theorem over a DVR.
Hence \(2(D+1)r\) generators suffice for the finite-condition
module and its quotients. Taking lengths first in \(C_F/J\) and
then in \(J\) gives
\[
 \frac{\operatorname{length}_O C_F}{\operatorname{length}_O(O/pO)}
     \le 2(D+1)rY+D(e+1)m\,18h .
\]
In particular the loss proportional to \(h\) per generator only
occurs on the \(D(e+1)m\) exceptional directions.

\proofstep{The common scalar exponent.}
For the numerical comparison put \(v=H/s\). Then
\(v\le 3eh/(2p)+eS/2+1\), as \(p\ge3\) and \(e\le p-1\).
The factorial bound gives \(\rho\le6v+3/2\) and \(\sigma\le6v\).
Therefore
\[
 Y\le48v+3+e(3S/2+\log_p H)+\eta+3/2+\log_p H .
\]

For \(h\ge1\), \(h+S+1\le h p^{S+1}\), so this is at most
\(80(eh/p+eS+e\log_p h+1+\eta)\) for \(e=1,2\).
The host always has \(h\ge1\). Substitution gives the ramified
estimate above and the same formula with \(e=1\) on an unramified
one-sign component. All choices precede the finite target class
apart from its own lift and expansion. The constants are independent
of the number of coefficient places, their ramification, the number
of Frobenius components, the source rank, and the outer jumps.
\label{cor:row-colength-1}
\end{proof}


\begin{proposition}[Weighted cost at the two ramified auxiliary primes]\label{prop:auxiliary-prime-row-length}
At \(\ell=q\) or \(d\), use the local field
\(E_w=F_w(\sqrt{-\ell})\) of \eqref{eq:host-auxiliary-completion}
for each \(w\mid\ell\), and
use the signed jump windows at each individual embedding
from Section~\ref{subsec:ramified-auxiliary-row}.
The local bound \eqref{eq:tag-r4}, after allowing the separate scalar
cohomology and class number budgets, is \(O(fmb)\).
Writing \(c_\ell\) for this local cost, the combined bound on the selected interval is
\[
 q,d\le\ell_*\quad\text{eventually},\qquad \sum_{\ell\in\{q,d\}}(\log\ell)c_\ell\ll fmb\log b.
\]
\end{proposition}

\begin{proof}
In this comparison put \(D_w=[E_w:\mathbb Q_\ell]\),
\[
 X_\ell=4b/\ell+2S_c+2\log_\ell(2b)+1+\eta_c+t,\qquad
 Y_b=2b+1.
\]
On the original row lattices the integral denominator
bound for a crystalline source and quotient with
the preceding jump conditions gives
\[
 L_w\le 160\bigl((D_w+1)r X_\ell+3D_wmY_b\bigr)+1 .
 \tag{R1}\label{eq:tag-r1}
\]
Here the local field is admitted at full level one,
with its coefficient field containing the embeddings
of \(F_w(\mu_\ell)\), and \(e(E_w/F_w)=2\).
The raw bound uses \(h_0=\lfloor2b\rfloor\),
\(S=\lfloor S_c\rfloor\) and
\(\lfloor\eta_c\rfloor+t\) in place of \(2b,S_c,\eta_c+t\).
Indeed the first two terms in \(X_\ell\) bound
\(2h_0/\ell+2S\), and its logarithmic term bounds
\(2\log_\ell h_0\). Rounding the resulting row bound up
adds at most one.

In the character rows again put \(h_0=\lfloor2b\rfloor\)
and write
\[
 \lambda=\lfloor\log_\ell(2h_0)\rfloor.
\]

The distinct jump set \(U\) of a source for a character
row lies in the union of the labeled jump sets. It
has at most \(D_wR\le2fR\) elements and lies in
\([-h_0,h_0-1]\). In the character lifting argument
at depth \(M=3+n_0\), take \(n_0\) least with
\[
 2h_0\le(\ell-1)\ell^{n_0},
\]
so \(M\le4+\lambda\).
The base change to \(F_w(\mu_\ell)\) costs no
\(\ell\)-power, as its relative degree is
\((\ell-1)/2\). The relative degree from
the first full level to depth \(M\) has valuation
\(M-1\).

For a twist exponent \(0\le e<t\), each nonzero
difference \(a-i\), with \(a\in U\) and
\(0\le i\le e\), has \(v_\ell(a-i)\le\lambda\).
The distinct-jump sum in each degree is bounded by
\(2fR(4+2\lambda)\). The exponent of the
truncation, scalar Bockstein and integral dual
exponential of Section~\ref{subsec:ramified-character}
is consequently at most
\[
 c_B=6+\lambda+7t+2tfR(4+2\lambda).
\]
Here \(6+\lambda\) allows the depth term and the
constant scalar terms, and \(7t\) allows the
per-degree and scalar terms linear in \(t\).
Since \(\lambda\le4\log b\) for \(\ell\ge3\), this gives
\[
                  c_B\le60tfR\log b .
 \tag{R2}\label{eq:tag-r2}
\]
The application is to the cocycle of the chosen
integral lift of the trivial quotient of a crystalline
extension. The character comparison moves its image
into finite cohomology. The rank-one local length
estimate then costs \((D_w+2)c_B\).
\label{cor:special-prime-lengths-9}

With no columns the global determinant bound of the
previous section has no column terms, and its row
length is exactly \(L_w\). The character row cost is
at most \(t\sum_{w\mid\ell}(D_w+2)c_B\).
The cyclotomic mark has depth one and \(\ell>t\).
Thus its normalized valuation is \(t-1\), giving a
determinant allowance at most \(4Rt\). The tame
character allowance is at most \(\ell^{-2}\). Together
with the signature length and the strict and unramified
global kernels this gives
\[
 v_{cc}(i,\pi)\le
 \sum_{w\mid\ell} L_w+t\sum_{w\mid\ell}(D_w+2)c_B
   +4Rt+\ell^{-2}+s_\pi+V_\ell+v_\ell(h_E).
 \tag{R3}\label{eq:tag-r3}
\]
Every \(D_w\ge2\) and \(\sum_wD_w=2f\).

In particular the two row sums are bounded by
\[
 160(3fr X_\ell+6fm Y_b)+f,\qquad
 240 f^2t^2R\log b
\]
respectively. The additional \(4Rt\) and \(f\)
are bounded by \(20 f^2t^2R\log b\).
For \(\ell\ge3\) and \(b\ge6\), each coefficient
in \(3X_\ell\) is at most \(13\) times its
corresponding coefficient in
\(b/\ell+S_c+\eta_c+t+\log b\), absorbing the constant
into the logarithm. Also \(6Y_b\le13b\) once \(b\ge6\).
Hence at this
threshold \eqref{eq:tag-r3} yields
\[
 \begin{split}
 v_{cc}(i,\pi)\le{}&
 2080 f\bigl(r(b/\ell+S_c+\eta_c+t+\log b)+mb\bigr)
     +260 f^2t^2R\log b\\
 &+2\ell^{-2}+V_\ell+v_\ell(h_E).
 \end{split}
 \tag{R4}\label{eq:tag-r4}
\]
For the second \(\ell^{-2}\) we used the signature
bound. This supplies the entry at the ramified quadratic prime on the same point
as its \(v_{cc}\) in the cut lattice.

\proofstep{Use the prime range to absorb the non-leading two terms in \eqref{eq:tag-r4}.}
For the selected auxiliary prime pair,
\[
 \log q,\log d=o(\log b)
\]
and both primes exceed
every fixed power of \(x\), uniformly in the scale
interval of that pair. In particular the condition
\(\ell>x^{2C_t+5}\) at the pointwise ramified caller
holds for the fixed rank exponent \(C_t\).
The ranks grow polynomially,
and \(\eta_c,S_c,t\) are at most
\(x^{O(1)}\exp(x^{\epsilon/4})\).
Thus the two terms of \eqref{eq:tag-r4} other than \(2080fmb\)
and the common budgets are \(o(fmb)\), uniformly
over both primes and all coefficient localizations.

In particular the table cost is at most \(2Cfmb\)
at a fixed large-scale threshold.
Also
\[
 q,d\le\ell_*,
\]
since the denominator of
\(\ell_*\) is \(x^{O(1)}\exp(O(x^{\epsilon/4}))\)
and its numerator is \(mb\). Summing with
\(\log q,\log d\le\log b\) gives \(O(fmb\log b)\),
which fits the stated \(O(fmb\log^2 b)\) budget.
\label{cor:special-prime-lengths-10}
\end{proof}


\subsubsection*{Ordinary split primes and the short band}
\label{subsec:ordinary-split-length}
\hypertarget{source-unramified-split-length}{}

After removing the type and ramified primes, the auxiliary field is unramified at the remaining primes in this range. At split primes the signed jumps lead to a short exceptional band; its placement must be compatible with the previously selected top positions.

\begin{established}[Unramified split range and signed jump window]\label{hyp:unramified-split-row}
Let \(\ell\) be outside the prime divisors of \(Q_S\), the primes
\(3,q,d\), and the two type primes. Consider first the range
\(\ell\le\ell_*\) together with the exceptional primes above
\(\ell_*\). At a prime split in the distinguished quadratic field, each column \(u\) is paired with its conjugate row, the column degrees sum to \(f\), and every completion is unramified. Put \(z=\log b\) and \(a=\lfloor b+S_c+\eta_c+t+2\rfloor\). For each column \(u\), both signed crystalline maps have jumps in \([-a,a]\), and their respective targets obey
\[
 \max\left\{\operatorname{rk}_O H^1_f(E_u,T_{2,u}),
          \operatorname{rk}_O H^1_f(E_{\bar u},T_{2,\bar u})\right\}
               \le [E_u:\mathbb Q_\ell]m .
\]
\end{established}
\begin{proof}
The excluded list contains every prime dividing the conductor of the host. Let \(U\) index the places on the selected side and \(D_u=[E_u:\mathbb Q_\ell]\). Conjugation gives the other side, so
\[
 2\sum_{u\in U}D_u=[E:\mathbb Q]=2f,\qquad |U|\le f .
\]

Lemmas~\ref{lem:marked-row-column-lattices} and~\ref{lem:marked-signed-jumps} identify the maps and the labelled jumps on this cut. At a good point use the direct normalized family on a column and
the Tate dual family on its conjugate row. Their jumps are in
\([-a,a]\) at every embedding by the jump boxes for the common
twist. Taking rational subquotients of finite direct sums preserves
crystallinity and this bound. For \(W\) on the column and
\(W^\vee(1)\) on the row, the free ranks of integral finite
cohomology are at most \(D_um\). Indeed the respective target
jumps avoid zero.

Among them \(k-1\) are positive on the column
and arise from negative jumps of \(W\) under \(j\mapsto-j-1\)
on the row. Write \(h_\tau^-\) for the number of negative target jumps at \(\tau\), including multiplicities. The split finite-rank formula in Section~\ref{subsec:split-local-length} gives, on either side,
\[
 \operatorname{rk}_O H^1_f(E_u,T_{2,u})=\sum_{\tau\mid u}h_\tau^-
                        \le D_um .
\] The column and row surjections in the congruence determinant
are on lattices in these subquotients. \label{cor:unramified-split-length-1}
\end{proof}

\begin{proposition}[Uniform finite lifting over the unramified completion]\label{prop:unramified-split-finite-lift}
For a crystalline lattice surjection in the signed window
\([-a,a]\) of Lemma~\ref{hyp:unramified-split-row}, the cokernel on integral finite cohomology
is killed by \(\ell^{e_A}\), where
\[
 e_A=\left\lceil23\bigl(2+2a/\ell+\log_\ell(2a+1)\bigr)\right\rceil .
\]
The allowance is independent of the degree of the coefficient
extension and the number of embeddings of the unramified field.
\end{proposition}

\begin{proof}
We show why the exponent on these surjections can be taken to be
\[
 e_A=\left\lceil23\bigl(2+2a/\ell+\log_\ell(2a+1)\bigr)\right\rceil .
\]
The estimate applies to a surjection of crystalline lattices with
the indicated window over any of these unramified fields. Use an
admitted coefficient extension containing the local embeddings and
the first cyclotomic level as in Section~\ref{subsec:cut-wach-models}.
The invariance of normalized length there returns the result on the
row or column surjection to its original coefficient ring. Use the
integral Wach--Herr comparison over the unramified field, with
Wach parameter \(\pi\) and height width \(h_1=2a\).
The Frobenius coefficient extraction of
Section~\ref{subsec:ramified-lift}, here with \(X=\pi\), puts
\(\pi^{h_1}N_2\) in the image of \(N_1\). On the full Wach
lattice the cyclotomic action is trivial modulo \(\pi\).
Choose a cyclotomic operator with character \(v\) generating
\(\mathbb Z_\ell^\times\), so its reduction has order
\(\ell-1\) and
\[
 v_\ell(v^{\ell-1}-1)=1.
\]

It acts on degree \(j\) by \(v^j\). Use the two interpolation
intervals \([0,h_1-1]\) and \([h_1,2h_1-1]\).
For a root \(v^i\) the only other roots giving nonunit
differences have indices congruent to \(i\) modulo \(\ell-1\).
On the two sides of \(i\) their numbers are
\[
 d_1=\lfloor i/(\ell-1)\rfloor,\qquad
 d_2=\lfloor(2h_1-1-i)/(\ell-1)\rfloor.
\]
The total valuation of the interpolation denominators satisfies
\[
 d_1+d_2+v_\ell(d_1!d_2!)\le \frac{2h_1\ell}{(\ell-1)^2}.
\]
The integral polynomial identity from the two intervals therefore puts \(\ell^{E_W}N_2\) in the image, for
\(E_W=\lfloor2h_1\ell/(\ell-1)^2\rfloor\).
For \(a=0\) the parameter-power inclusion already gives the lift.
Averaging over the subgroup \(\Delta\) of order \(\ell-1\)
preserves this bound on the invariant lattices. This elementary
interpolation bound suffices at the split caller.

Write \(N=N_0^\Delta\) for the invariant Wach lattice of the
kernel. The invariant coordinate \(P\) is a unit times
\(\pi^{\ell-1}\). Put \(Q'=P+\ell\), and let \(u\) be the
character of the principal cyclotomic generator \(\gamma\), with
\(v_\ell(u-1)=1\). The comparison gives the restricted coordinates
\((Q')^{-a}N\oplus PN\) for pairs and differential
\[
 d(y,z')=(\gamma-1)y-(\varphi-1)z'.
\]

It gives the following two facts on the kernel. First,
\(N\) is \(\ell\)-saturated in \((Q')^{-a}N\), and
\(\partial=(\gamma-1)/P\) preserves \(N\).
Second, if \(U\) is the set of source jumps, put
\[
 J=\big((U+\{1,\ldots,a\})\setminus\{0\}\big)
                \cup\{1,\ldots,a\},\qquad
 R_0(Z)=\prod_{j\in J}(Z-c_j),\quad c_j=(1-u^{-j})/\ell .
\]
Then \(\partial R_0(\partial)y\in N\) for \(y\in(Q')^{-a}N\).
The latter is the cyclotomic semisimplicity input on the
bounded-pole quotient. The comparison at the cyclotomic roots
reads the jumps shifted by the pole degrees \(1,\ldots,a\).
In the divided operator the eigenvalue for shift \(j\) is
\(c_j\), since the invariant coordinate at the root is
\(-\ell\). The extra factor \(Z\) accounts for shift zero.
Strictness puts the kernel jumps among the source jumps.
\label{cor:unramified-split-length-2}

Here is the correction of a defect in lifting a pair. Put
\(d_0=\lfloor a/(\ell-1)\rfloor\) and
\(G=\prod_{i=1}^{d_0}(1-u^{(\ell-1)i})\). On
\(P^iN/P^{i+1}N\), \(\gamma\) acts by \(u^{(\ell-1)i}\).
Successive subtraction through the graded factors gives
\[
 GPN\subseteq(\gamma-1)PN+P^{d_0+1}N.
\]

The height inclusion and the factor \((Q')^{\ell-1}\) in
\(\varphi(P)/P\) make Frobenius contracting on
\(P^{d_0+1}N\), since \((\ell-1)(d_0+1)>a\).
In this complete submodule the geometric series solves the inhomogeneous equation:
\[
 (1-\varphi)\sum_{n\ge0}\varphi^n x=x\qquad(x\in P^{d_0+1}N).
\] The integral restricted Herr
correction applied with these data kills every rationally
trivial defect by \(G R_0(0)\). To see the role of division
by \(P\), write the defect as \(P\beta\).

In the equation
for a multiple which is a boundary, use the graded correction
on the second coordinate. Dividing by \(P\) leaves
\(\partial\) on the first coordinate and an integral term
on the other side. Applying \(R_0(\partial)\) makes this
first-coordinate difference integral, by the bounded-pole
identity. More explicitly this gives \(\ell^j G R_0(\partial)\beta\in N\)
when the original boundary equation holds after multiplication by \(\ell^j\). Saturation in the
bounded-pole lattice and \(G\mid\ell^{v_\ell(G)}\) up to a unit
give
\[
 R_0(\partial)\beta\in N.
\]

Using \(R_0(Z)=R_0(0)+Z S_0(Z)\) and
\(P\partial=\gamma-1\), the corrected defect satisfies
\[
  R_0(0)P\beta
       =P R_0(\partial)\beta-(\gamma-1)S_0(\partial)\beta .
\]
Here \(P R_0(\partial)\beta\in PN\). The second summand is
the indicated cyclotomic difference. The positive-degree
correction by \(G\) and the deep correction by \(1-\varphi\)
apply to the first summand in \(PN\).

\label{cor:unramified-split-length-3}

An integral finite class is represented by a crystalline extension
of the trivial lattice. Lifting its quotient generator on the
invariant Wach module costs \(E_W\) and gives a restricted pair.
Lifting the two coordinates through the given surjection costs
a second \(E_W\). The defect lies in the kernel coordinates
by exactness. It is rationally trivial by surjectivity on rational
finite classes for crystalline quotients and the rational Herr
comparison. The integral correction gives a finite source class.
The resulting exponent is at most
\[
  2E_W+\sum_{j\in J}v_\ell(j)
       +(d_0+1)+v_\ell((d_0+1)!).
\]
Indeed \(v_\ell(c_j)=v_\ell(j)\) and
\(v_\ell(G)=d_0+v_\ell(d_0!)\).
The distinct indices lie in \([-a,2a]\setminus\{0\}\), so their
sum of valuations is at most \(3a/(\ell-1)\).

For \(\ell\ge3\), the three displayed contributions are at most
\(9h_1/\ell\), \(9h_1/(4\ell)\), and
\(3(a/(\ell-1)+1)/2\), respectively. Their sum is at most
\(13h_1/\ell+2\le e_A\). In particular the exponent has no
factor depending on the number of embeddings or the coefficient
ramification. \label{cor:unramified-split-length-4}
\end{proof}

\paragraph{The determinant on the split family.}
\begin{lemma}[Split determinant bound at the same lift]
\label{lem:unramified-split-determinant-total}
Use the unramified split range and the maps of
Lemma~\ref{hyp:unramified-split-row}, with \(z\ge1\). If \(L\) denotes the total lifting term and \(D\) the sum of the remaining local determinant terms, then
\[
 L\le (2fm+f)e_A,\qquad D\le7200fnt^2Rz^2.
\]
The separate global and signature allowances are unchanged.
For the inert declaration only \(B_0\) from this calculation
is used.
\end{lemma}

\begin{proof}
Use Proposition~\ref{prop:unramified-split-finite-lift} on the
row and column surjections specified by the same lift.
Normalize every length by \(\operatorname{length}_O(O/\ell O)\).
Write \(R=r+t\), \(h=\lfloor2b\rfloor\),
\(N_0=4nR\), \(K_0'=\lfloor\log_\ell(2h)\rfloor\) and
\(K_1'=\lfloor\log_\ell(2(h+1))\rfloor\).
There are at most \(D_wR\le N_0\) distinct jumps in the
character calculation. Passing to \(E_w(\mu_\ell)\) has
prime-to-\(\ell\) degree. Thus the character calculation giving
\eqref{eq:tag-r2} applies, with depth at most \(4+K_0'\), and gives exponent
\[
 c_B=6+K_0'+7t+tN_0(4+2K_0').
\]
For the column and dual row torsion comparison the window lies in
\([-h-1,h+1]\) with zero absent.

In the invariant formula in
Section~\ref{subsec:split-local-length}, take one truncation
degree and full cyclotomic depth one. The depth term is at most
\(3+K_1'\), the sum over jumps is at most \(N_0(4+2K_1')\),
and the two integral comparison losses in this degree add two.
The bound for the invariant-killing exponent is therefore
\[
 c_0=5+K_1'+N_0(4+2K_1').
\]
If the jump sum is empty, the extra unit allowance in the
maximum is also covered by this bound since \(N_0\ge1\).
Consequently \(h_W(u)\le rc_0\).
The free and torsion parts of the dual row give colength at
most \(D_um e_A+rc_0\). The torsion bound follows from the
invariant subgroup of the divisible target modulo its lattice.
These are estimates on the given lattices. \label{cor:unramified-split-length-5}

Write \(k_c\) for the number of columns. At depth one the
marked scalar satisfies \(v_\Delta\le 2t\) by the cyclotomic
difference valuations. On the nontrivial scalar twists
\((K_\ell/O)(-j)\) with \(1\le j<t\), the unramified base has
an inertia element of character \(1+\ell\), so its invariants are killed by \((1+\ell)^{-j}-1\). Their
normalized length is bounded by
\[
 v_\ell((1+\ell)^{-j}-1)=v_\ell((1+\ell)^j-1)
                         =1+v_\ell(j)\le t .
\] Use \(h_0\le t\) for
their maximum, or zero for an empty set.

Summing the three
blocks of the determinant gives a common local term
\[
 B_0=t\sum_{w\ {\rm row}}(D_w+2)c_B
       +4R(k_c+1)v_\Delta+tk_ch_0 .
\]
The column and conjugate row at \(u\) add
\((2D_um+1)e_A\) and the error
\(D_umv_\Delta+2v_\Delta+h_W(u)+rc_0\).
In fact the column contributes \(j_u(e_A+v_\Delta)+e_A\)
and the row contributes \(D_um e_A\) before adding their
torsion costs. The other column and coboundary denominators
are those of Section~\ref{subsec:split-local-length}.

Using \(z\ge1\), \(K_0'\le2+2z\) and \(K_1'\le3+2z\), the two scalar losses are
\[
 c_B\le120t n\,R z,\qquad c_0\le120 n\,R z.
\]
Put
\(U_0=120fnt^2Rz\). The sum of row degrees plus twice their
number is at most \(6f\), and \(k_c\le2f\). Thus the three
terms of \(B_0\) are at most \(6U_0,24U_0,2U_0\).
The additional split errors are at most \(4U_0,8U_0,4U_0\)
for the degree term, the scalar term, and the two invariant
terms, using \(m\le r\le t\).

Adding the six errors bounds their sum by
\[
        (6+24+2+4+8+4)U_0
           \le 7200 fnt^2Rz^2 ,
\]
since \(z\ge1\). The sum of the lifting terms is at most
\((2fm+f)e_A\). Together these give the split and boundary
entries in Section~\ref{subsec:prime-cost-sums}, since
\(e_A\ll 1+H/\ell+\log_\ell H\). The global kernels, the
signature bound and the tame character costs retain their
separate allowances. In the inert case \(k_c=0\) and the
same calculation bounds \(B_0\) alone. The rank-\(r\) row
colengths there depend on the inert jump word.
\label{cor:unramified-split-length-6}
\end{proof}

\paragraph{Transport of the band at an inert place.}
\begin{lemma}[Passage from the row band to the place bound]
\label{lem:inert-place-band}
In the last two inert ranges let \(b_{\rm row}\) be the marked row band indicator and \(b_\ell\) the place indicator used in the weighted sum at its rational prime. Then
\[
 b_{\rm row}\le b_\ell\qquad(b_{\rm row},b_\ell\in\{0,1\}).
\]
\end{lemma}

\begin{proof}
On a row in either of the last two inert ranges, suppose the
local band indicator is one. In the canonical data this gives
a good point, a row word at a place over \(\ell\ge11\), and
a position with smaller jumps of local weight \(a_0\) with
\[
 |\ell a_0-b_0|\le K\{\ell\eta+\ell S+b_0/\ell+a_0+
                                          \ell\log_\ell(2h_b)\}.
\]
Here \(h_b\) is the height of the word, and \(b_0\) its large
weight. The weight box on the same word supplies a place weight
\(s\) between \(\omega\) and the upper place-weight size with
\[
 |a_0-s|\le T,\qquad b-T\le b_0\le b-T+1,\qquad T=(t+1)/2 .
\]
The prime beneath this row equals the table prime, as the
contraction of its maximal ideal contains a unique rational
prime.

We have \(\eta\le\eta_3=\eta_c+t+2\),
\(S\le S_c\), \(h_b\le b_0+\eta+S\le 2b\) on the
selected range, and \(\log_\ell(4b)\le z\) since \(z\ge1\)
and \(4e<11\). Put \(A_s=\ell(\eta_3+S_c+z)+b/\ell+s\).

The braces in the local inequality are at most
\(A_s+\ell\eta_3\), since \(a_0\le s+T\) and
\(b_0/\ell\le b/\ell+1/\ell\).
The two changes on the left cost at most
\(\ell T+T\le\ell\eta_3\). Hence
\[
                 |\ell s-b|\le(2K+1)A_s .
\]
The table constant dominates \(6K+1\), and therefore the
indicator in the weighted sum dominates every local band
indicator. Hence the global harmonic sum bounds the frequency of the band on each of these rows. Its contribution per occurrence still requires an inert-row colength estimate.
\label{cor:unramified-split-length-7}
\end{proof}


\subsection{Ordinary inert primes on the retained cut}
\label{subsec:inert-rows}

At the ordinary inert primes the split column condition disappears. The row calculation must instead follow the signed jumps around a complete Frobenius orbit in the unramified local field.

Use the ordinary prime range of Section~\ref{subsec:ordinary-split-length}.
An inert prime in this range has \(\ell\ge11\) once the family scale exceeds its lower threshold. Every place \(\lambda\mid\ell\) is a row. Its completion
\(F=E_\lambda\) is unramified over \(\Q_\ell\), with degree \(d_\lambda\).
The cyclic index at \(\lambda\) runs over
\[
 \operatorname{Hom}_{\mathbb Q_\ell}(E_\lambda,\overline{\mathbb Q}_\ell),\qquad \#\operatorname{Hom}_{\mathbb Q_\ell}(E_\lambda,\overline{\mathbb Q}_\ell)=d_\lambda.
\]
The ordering comes from Frobenius on the actual host, including its quadratic factors.

We describe the Wach levels and compare the arcs at these
inert primes.

Use the marked row surjection \(T_1\twoheadrightarrow T_2\) of
Section~\ref{subsec:split-local-length}. The target is the rank-\(r\)
Tate dual at the selected point and \(T_1\) is a lattice in a
subrepresentation of a finite sum of the Tate duals of the normalized
retained representations. The rational comparison is crystalline.
The local jump windows hold on all of the summands on this cut.
Strictness for the rational crystalline filtration passes them to
\(T_1\) and to \(T_2\). For the marked rank-\(r\) target the central multiplicity bound is pointwise in the letter:
\[
 \#\{\text{central target jumps at }i\}\le m\qquad(\text{every }i),
\]
where multiplicities are counted.
Section~\ref{subsec:cut-wach-models} gives the signed integral
Wach inclusions on these same two lattices and their local finite
extensions, at the labeled places of the retained tuple.

The common word is obtained by restricting the host weight chart
to the Frobenius cycle of \(\lambda\). Its signs are the two CM
signs for the Tate dual. At an inert prime Frobenius exchanges
these signs, which accounts for the alternating words and the
even local degree. Rounding the upper interval and its width
gives local bounds \(b_0,w,S,\eta\) as below, within the common bounds
\[
 w\le w_P,\qquad S\le S_c,\qquad \eta\le\eta_3,\qquad b_0+S+\eta\le H.
\]
The shifted place weights obey the same
central gap at every retained point. These comparisons concern
the surjection and jump data used by the congruence determinant,
including the case of a ramified extension of the coefficient DVR.
\label{cor:inert-row-lengths-1}

\paragraph{The inert contributions to the weighted prime sum.}
\label{sec:cut-length-rows}
Use \(H,G_P,w_P,\eta_3,z\) from
Section~\ref{subsec:prime-cost-sums}. Write \(A_{\rm sec}\) for
the secondary threshold denoted \(A\) there and \(o_\ell\) for
the order in the host's real cyclotomic quotient. Case B has
\(\ell<A_{\rm sec}\) and \(o_\ell>4L_w\). Set
\[
 j_\ell=\max(1,\lceil\log_{\ell-1}A_{\rm sec}\rceil),\qquad
 U_B=(d_\lambda+1)r
     \left(\frac{H}{(\ell-1)^{j_\ell+1}}+
          (G_P+1)(j_\ell+1)+z\right).
\]
The remaining ordinary inert primes are case C if
\(\ell\ge A_{\rm sec}\), and case D if \(\ell<A_{\rm sec}\) and
\(o_\ell\le4L_w\). The local comparison distinguishes two
cases. Larger jumps occur at every position precisely when
Frobenius acts as complex conjugation and the starting
position has larger jumps.

To see both alternatives, first consider \(o_\ell\le4L_w\).
Two visits to the large set with labels \(s_1,s_2\) satisfy
\[
 D\mid o_\ell t_0(s_2-s_1),\qquad
 n_0\mid o_\ell(s_2-s_1),\qquad
 |o_\ell(s_2-s_1)|<n_0 .
\]

The first relation follows from \eqref{eq:translated-dyadic-visits},
whose coordinate formula holds at every order. Thus \(s_1=s_2\).
Each label specifies a single real place, represented by the two
conjugate embeddings of \(E\). Consecutive larger positions would
then give the same real place, so Frobenius fixes an embedding of
\(E^+\). Since \(E/\mathbb Q\) is abelian, it belongs to
\(\operatorname{Gal}(E/E^+)=\{1,c\}\). Inertness in
\(\mathbb Q(\sqrt{-d})\) rules out 1. For Frobenius \(c\),
a larger starting place remains larger throughout its cycle.
In the long-order case the three-visit argument of
Lemma~\ref{lem:long-orbit-spacing} applies for any ordinary inert
prime, since it uses only \(o_\ell>4L_w\) and the designated
labels. As \(L_w\ge3\), at most two consecutive positions are larger.

Use the place indicator \(B_\lambda\)
that a good point over this place exhibits the short band of
Section~\ref{subsec:ordinary-split-length}. Define
\[
\begin{aligned}
 \mathcal P_\ell&=1+H/\ell+\log_\ell H,&
 F_{\ell,\lambda}&=(w_P+\eta_3)/\ell+S_c+z+\eta_3B_\lambda,\\
 U_H&=(d_\lambda+1)rH/\ell^2,&
 U_F&=(d_\lambda+1)r F_{\ell,\lambda},&
 U_M&=d_\lambda m\mathcal P_\ell .
\end{aligned}
\]
The following estimates yield one rational bound at each place,
valid simultaneously for its good points and coefficient lattices.
The case B entry is \(O(U_B)\), and the common case C and D
entry is \(O(U_H+U_F+U_M)\). Their constants are absolute once
the host constants are fixed. In particular they are independent of
the degree of the coefficient field and the jump heights.
\label{cor:inert-row-lengths-2}

Every \(d_\lambda\ge1\), so the sum over all row places satisfies
\[
 \sum_{\lambda\mid\ell}(d_\lambda+1)\le2\sum_{\lambda\mid\ell} d_\lambda=4f.
\]
The global indicator \(I_\ell\) in
Section~\ref{subsec:prime-cost-sums} bounds every \(B_\lambda\)
by the band transport there. Thus the \(U_B\), \(U_H\) and \(U_F\)
terms give the three inert costs in that section, after
enlarging the fixed common constant. The \(U_M\) terms contribute
at most \(2fm\mathcal P_\ell\) before that constant and go into
the same sum as the split lifting term. The factor in the
global determinant is absorbed by the common constant.
At these inert primes the column set is empty, so the remainder
is the boundary term with \(k_c=0\) in the split calculation.

For case B, it remains to obtain the required row colength from the length of the Frobenius orbit; this is Proposition~\ref{prop:inert-long-colength}.

\subsubsection{The long-orbit row bound}\label{sec:inert-long-orbit}
\hypertarget{source-inert-long-orbit}{}

Among the inert orbits, a long Frobenius orbit allows the marked jumps to be cut into shorter arcs. The colength can be bounded using those arcs only if the correction at each cut is controlled integrally.

\begin{established}[Case B and the cyclic signed jump windows]\label{hyp:long-orbit-data}
Let \(\ell<A_{\rm sec}\) be an ordinary inert prime whose Frobenius
has order \(o_\ell>4L_w\) in the real cyclotomic quotient
(case B of Section~\ref{sec:cut-length-rows}). Write \(F=E_\lambda\) for the
unramified completion of the auxiliary field and
\[
 d=[F:\Q_\ell]=n_\lambda.
\]

Work over the integers \(O\) of a local coefficient field splitting \(F\).
Write \(T_2\) for the marked dual Tate lattice, of rank \(r\), and
\(T_1\) for a lattice in a subrepresentation of a finite sum of the
twisted duals from the same cut. The maps to be estimated, including the marked matrix row, are precisely the indicated integral Galois surjections
\[
 f:T_1\twoheadrightarrow T_2,\qquad \operatorname{rank}_O T_2=r .
\]

The source and target are crystalline with central jumps \(|a|\le\eta_0\) and one outer signed interval of width \(S_0\) at each Frobenius embedding. The small letters have height at most \(G_P\) and all letters have height at most \(H\), in the notation of Section~\ref{sec:cut-length-rows}.
\end{established}
\begin{proof}
The ordinary inert branch excludes all primes ramifying in the field or in the rows. Conjugating its marked row preserves the coefficient lattice and takes its dual Tate jumps to the reflected signed intervals. Apply Lemma~\ref{lem:marked-row-column-lattices} to obtain the displayed surjection. Lemma~\ref{lem:marked-signed-jumps} gives the central and outer bands before passage to its quotient. Both rational representations are crystalline with the
local jump windows from the cut. Enumerate the embeddings of \(F\)
cyclically by Frobenius. The residue characteristic is \(\ell\ge11\)
in this inert branch.

Use \(b_0,w_0,S_0,\eta_0\) for the local size parameters, with
\(S_0,\eta_0\ge1\), and \(s_i\) for the size at position \(i\).
The source jumps at \(i\) lie in \([-\eta_0,\eta_0]\) together
with either \([s_i,s_i+S_0]\) or its negative. The larger range is
\[
 b_0\le s_i\le b_0+\eta_0.
\]

At every other position,
\(s_i\le w_0\). Put \(g=w_0+S_0+\eta_0\) and set
\[
 h=\begin{cases}b_0+\eta_0+S_0,&\text{if some position has larger jumps},\\
                  g,&\text{otherwise}.\end{cases}
\]
Thus the positions with smaller jumps have jumps in \([-g,g]\) and all positions
have jumps in \([-h,h]\). The target satisfies the same bounds by
strict exactness of the rational crystalline filtration.
The two upper endpoints in Section~\ref{sec:cut-length-rows} bound \(b_0+\eta_0+S_0\le H\)
and \(g\le G_P\), with \(H,G_P\) as in Section~\ref{sec:cut-length-rows}.
The length of the cycle here is the unramified local degree, including
the quadratic components of Frobenius.
\end{proof}

\begin{auxlemma}[Spacing of the larger jumps]\label{lem:long-orbit-spacing}
On the unramified cycle of Lemma~\ref{hyp:long-orbit-data}, let
\(J\) be the positions with larger jumps and \(L_w\) the length
parameter for \(T_2\) in Lemma~\ref{lem:second-set-progressions}.
Write \(\theta_\ell\in\mathbb Z/D\mathbb Z\) for the action of
arithmetic Frobenius on \(C_e\). If
\(o_\ell=\operatorname{ord}(\theta_\ell)>4L_w\), then
\[
 \#(J\cap I)\le2\qquad\text{for every interval of }L_w
                  \text{ consecutive positions }I .
\]
If \(1\le j_\ell\) and \(2j_\ell+2\le L_w\), the cycle admits
cuts with \(j_\ell\) smaller positions following each cut, at most
two larger positions per arc, and maximum arc length \(3j_\ell+1\).
\end{auxlemma}
\begin{proof}
In \(\operatorname{Gal}(C_e/\mathbb Q)\), the step \(\theta_\ell\)
is the class of \(\ell\bmod 2^{e+2}\) up to sign, expressed in the
generator 5. Thus
\[
   o_\ell=\min\{a\ge1:\ell^a\equiv\pm1\pmod{2^{e+2}}\},
    \qquad o_\ell\mid [E_\lambda:\mathbb Q_\ell] .
\]
The divisibility holds because the order of the full Frobenius
equals the unramified local degree and restriction preserves powers.

By Lemma~\ref{hyp:translate-weights} the larger places are the
translated labels \(e_1(s+u-t)\) and their complex conjugates
at the two embeddings of a real place when the coordinates are
centred at the selected top \(t\). Restriction to the real
field \(C_e\) kills complex conjugation. Let \(w\) be the
coordinate of the translate, equal to \(\rho(u-t)\) with this
origin or \(\rho(u)\) before recentering. With \(s_i\) the integer representative of the larger label at
position \(i\), \eqref{eq:dyadic-coordinate-orders} gives
\begin{equation}\label{eq:translated-dyadic-visits}
 x_0+i\theta_\ell=w+t_0 s_i\pmod D,\qquad
                  0\le 4L_w s_i<n_0\quad(i\in J).
\end{equation}

Here \(x_0\) is the starting coordinate. Reversing the orientation
replaces \(\theta_\ell\) by its negative and preserves its order.
If \(a<b<c\) in one interval of length \(L_w\) are larger
positions, their coordinates are distinct by the order bound.
Put \(u'=b-a\), \(v'=c-a\). Eliminating the step gives
\[
 t_0\big(v'(s_b-s_a)-u'(s_c-s_a)\big)\equiv0\pmod D,
 \qquad
 v'(s_b-s_a)-u'(s_c-s_a)\equiv0\pmod{n_0}.
\]

The second implication uses \(D=t_0n_0\). The four terms have
absolute value less than \(n_0/4\) by \eqref{eq:translated-dyadic-visits},
so the last congruence is an integer equality between three
distinct selected labels. The sphere construction in
Lemma~\ref{lem:second-set-progressions} excludes this equality.
Taking integer lifts of the positions includes intervals crossing
the end of the cycle.

Let \(j=j_\ell\). The numerical host choice allows some exponent
\(j'\) with \(A_{\rm sec}\le(\ell-1)^{j'}\) and \(2j'+2\le L_w\).
Here \(\ell<A_{\rm sec}\), so the least exponent \(j\) is at least two
and satisfies
\[
 2j+2\le L_w.
\]
Thus the local cycle has \(d>2j\).
Join successive positions with larger jumps when their forward gap is at most
\(j\). A chain of three would occupy at most \(2j+1\) consecutive
positions and is excluded. The clusters thus have one or two
positions, with at least \(j\) intervening positions with smaller jumps between
clusters. Cut just after the last position with larger jumps of each cluster,
using that position as the cut index \(c\). Each of \(c+1,\dots,c+j\)
then has smaller jumps.

If the next cluster starts at forward distance \(a\),
insert cuts at the multiples of \(j\) up to \(a-j-1\), measured
from the current cut. Include the original cut at distance zero. If
\(k=\lfloor(a-j-1)/j\rfloor\), the remaining gap is
\[
 a-kj=j+1+\bigl((a-j-1)\bmod j\bigr)\in[j+1,2j].
\] Adding the span of the next
cluster gives an arc length at most \(3j\). Each arc is taken
half-open, with its terminal cut included.

Each such arc contains at most two
positions with larger jumps. When there are no positions with larger jumps, partition the
cycle into arcs of length \(j\) with the last remainder added to one
arc. Each arc has length at least \(j\), since inserted cuts are
spaced by \(j\) and each gap to a cluster of positions with larger jumps starts with its
\(j\)-letter buffer of positions with smaller jumps. The maximum length \(A\) satisfies
\(j\le A\le3j+1\).
\end{proof}

\paragraph{The exponent in the Wach lifting estimate.}\label{par:inert-cut-exponents}
The correction at a cut in Paragraph~\ref{par:wach-cut-correction} and the integral level tests of Proposition~\ref{prop:wach-arc-membership} apply to the
cyclotomic Wach modules of a surjection of crystalline lattices over
an unramified field. Its arc hypotheses are a cyclic partition with
\(j\) positions with smaller jumps after each cut, at most two positions with larger jumps per
arc, and specified size inequalities when their signs differ.
Lemma~\ref{lem:long-orbit-spacing} supplies the partition.

We now compute the exponents for the symmetric jump boxes of
Lemma~\ref{hyp:long-orbit-data}.
Put \(q_0=\ell-1\), \(D_0=q_0^{j+1}\), and
\[
 \begin{split}
 M&=\left\lfloor\frac{h+2g}{q_0}\right\rfloor+1,\\
 \sigma&=\left\lceil 1+\frac{\ell}{q_0^2}
                      \left(3g+\frac{M}{\ell^{j-1}}\right)\right\rceil,\\
 H_i&=\begin{cases}2h,&i\text{ with larger jumps},\\2g,&i\text{ with smaller jumps},\end{cases}\\
 c_H&=\left\lceil 1+\log_\ell(1+\max_i H_i)\right\rceil,\\
 S(c)&=\sum_{a\ge0}\frac{H_{c+a+1}+c_H+3}{q_0^a},\\
 R&=\max_c\left(c_H+\left\lfloor\frac{\ell}{q_0^2}
           \min\{\lfloor S(c)\rfloor,\max_i H_i\}\right\rfloor\right).
 \end{split}
\]
The initial exponent is \(a_0=2\sigma+R+2gA\). Define \(\kappa\)
as the maximum of
\[
       A+\sum_{u\in J\setminus\{0\}}(A+1+v_\ell(u))
\]
over unions \(J\) of three intervals of at most \(3g\) consecutive
integers in \([-2h,2h]\).

In Theorem E the inputs for the source and
target can both be taken to be \(H_i\). If the positions with larger jumps of
opposite sign have magnitudes \(x\) at backward offset \(s_1\) and
\(y\) at offset \(s_0<s_1\), the required comparisons are
\(y<\ell x\) when the earlier position in the cycle is negative,
and \(2y\le\ell x\) and \(2a_0<x\) when it is positive. With these
inputs the theorem lifts every integral finite class of \(T_2\) after
multiplication by \(\ell^e\), where
\begin{equation}\label{eq:long-orbit-exponent}
                 e=2R+4\sigma+2a_0+2\kappa.
\end{equation}

For the proof with cyclotomic evaluation coordinates, extend the
coefficient field to contain the primitive \(\ell^d\)-th roots of
unity. This extends the coefficients while keeping \(F\) and its
Frobenius cycle fixed. The embeddings over the larger coefficient
field are the same after composition, since the original field
already splits \(F\). The finite cohomology group of a lattice
is the preimage of the rational finite group.

A basis of the
extended coefficient ring over \(O\) identifies its continuous
cochains with finitely many copies of the original complex.
Rationalization and the crystalline period comparison also commute
with this finite direct sum. Thus integral finite cohomology, its
map and the cokernel extend by tensor product. For the cokernel \(C\) and extended coefficient DVR \(O'\), faithful flatness gives
\[
 \ell^e(C\otimes_O O')=0\quad\Longrightarrow\quad\ell^e C=0.
\]


The geometric estimate for \(R\) exhibits the dependence on the
buffer of positions with smaller jumps. The first \(j\) terms of the sum at each cut have
\(H_{c+a+1}=2g\), hence
\[
 S(c)\le \frac{q_0}{q_0-1}
          \left(2g+c_H+3+\frac{2h}{q_0^j}\right).
\]
Since \(\ell\ge11\) and \(h\ge1\), we have
\(c_H\le3+\log_\ell h\) and
\[
  R\le8g+3h/D_0+2\log_\ell h,
  \qquad \sigma\le3g+2h/D_0.
\]
Indeed \(\ell/(q_0(q_0-1))\le11/90\) and
\(2\ell/(q_0-1)\le22/9\) give the first inequality.

For the second, \(\ell/q_0^2\le11/100\), and the coefficient
of \(h+2g\) is at most \((11/10)/D_0\) for \(j\ge1\).
Choose the common constant \(C_2\ge3200\) in the size condition
\(b_0\ge C_2(j+1)(g+1)\). This condition holds uniformly for
the rounded local word. In fact \(b_0\ge b-(t+1)/2\ge b/2\)
at large scale, \(g+1\le G_P+1\le11b'\), and
\(j+1\le3+\log A_{\rm sec}=O(\log x)\). Hence
\[
 C_2(j+1)(g+1)/b_0=O(C_2\log x\, b'/b),
\]
which tends
to zero since
\[
 C_2\log x\,b'/b=C_2x\log x\exp(-x^{\varepsilon/3})\longrightarrow0.
\]

If there are positions with larger jumps,
\(g\le b_0\), \(h\le2b_0\), and both magnitudes are between
\(b_0\) and \(b_0+\eta_0\). Therefore \(2y\le4b_0<\ell x\).
The bounds on \(\sigma,R\) in the expression \(a_0=2\sigma+R+2gA\) give
\[
 2a_0\le40g(j+1)+14h/D_0+4\log_\ell h<b_0.
\]
For example, the three terms are bounded by \(0.04b_0\),
\(0.28b_0\), and \(4\sqrt{2b_0}\) using \(b_0\ge2000\) and
\(\log_\ell h\le\sqrt h\). The two mixed-sign comparisons follow.
When all positions have smaller jumps these comparisons are vacuous.

To bound \(\kappa\), each interval with at most \(3g\) integers
contains at most \(3g/\ell^v+1\) nonzero multiples of \(\ell^v\),
and such multiples require \(v\le\log_\ell(2h)\). Summing over
the three intervals gives
\(\kappa\le25gA+3\log_\ell h\) since \(A,g\ge1\).
Using \eqref{eq:long-orbit-exponent},
\[
 e\le110gA+28h/D_0+14\log_\ell h
   \le C_2\left(h/D_0+(g+1)(j+1)+\log_\ell h\right).
\]
In particular the high jump bound carries the factor
\((\ell-1)^{-j-1}\) even when the cycle has many clusters. The
exponents use a maximum over the arcs. The factor for the host
degree enters only in the dimension of finite cohomology.

\paragraph{Correcting the lattice at a Frobenius cut.}\label{par:wach-cut-correction}
In the ordinary cyclotomic variable \(\pi\), let \(\Delta\)
be the tame group of order \(\ell-1\) and take
\(X=-\prod_{\delta\in\Delta}\delta(\pi)\).
\begin{lemma}[The tame-invariant parameter]
The tame-invariant coefficient ring is \(O[[X]]\). Put \(q=X+\ell\)
and \(p_s=\phi^s(q)\). Up to a unit, \(p_s\) is a distinguished
polynomial of degree \(\ell^s\) with constant term associated to
\(\ell\), and for each \(j\ge1\),
\[
 (\phi^j X)=(X)\prod_{0\le s<j}(p_s)^{\ell-1}.
\]
\end{lemma}
\begin{proof}
The defining product is a unit times \(\pi^{\ell-1}\). On the invariant ring
write \(\mathcal R=O[[X]]\). Indeed the action on degree \(k\) in the \(\pi\)-adic
filtration is by the \(k\)-th power of the tame character.
An invariant series can only have its first nonzero term in
degree a multiple of \(\ell-1\). Successively subtracting
a constant multiple of the corresponding power of \(X\)
expresses it as a series in \(X\).
At a primitive \(\ell\)-th root of unity minus one,
\(X\) evaluates to \(-\ell\), by the product formula for
all primitive \(\ell\)-th roots. Thus \(q=X+\ell\) is
associated to the usual cyclotomic factor in \(O[[\pi]]\)
by Weierstrass preparation.

As a result
\[
 \phi(X)=uXq^{\ell-1},\qquad
 \gamma(X)=\omega_X X\pmod{X^2},\qquad
 v_\ell(\omega_X-1)=1 ,
\]
where \(u\) is a tame-invariant unit and
\(\omega_X=\chi(\gamma)^{\ell-1}\).
Here \(\phi\) acts on coefficients in
\(\mathcal R\) and \(\Phi\) is the corresponding semilinear operator
between letters. Put \(p_s=\phi^s(q)\). Iteration gives
\(\phi^j(X)\) associated to \(X\prod_{s<j}p_s^{\ell-1}\), and
\(p_s\) lies in \((\ell,X^{\ell^s})\). The latter follows by
reducing the Frobenius formula modulo \(\ell\). By Weierstrass
preparation, \(p_s\) is a unit times a monic polynomial of degree
\(\ell^s\) whose other coefficients are divisible by \(\ell\)
and whose constant coefficient is associated to \(\ell\).
Different \(s\) have different cyclotomic zero sets and their
distinguished polynomials are relatively prime over
\(\operatorname{Frac}(O)\).

\end{proof}
For the free Wach lattices \(N_i\), use the height bounds
\(q^h\Phi N_i\subset N_{i+1}\) everywhere and
\(q^g\Phi N_i\subset N_{i+1}\) in the first \(j\) steps after a
cut. The cyclotomic action is trivial on \(N_i/XN_i\). Write
\(P_c=\mathcal R\Phi^j N_c,\ D=X^M\prod_{s<j}p_s^{3g}\).

\begin{lemma}[Correction at a cut]
Use the chosen positive integers \(M,j\) and nonnegative integer
\(\sigma\) with \((\ell-1)M>h,\ (\ell-1)M\ge3g\) and
\((\ell-1)^2\sigma/\ell-3g\ge M/\ell^{j-1}\). In the odd-prime
range of the cut, each \(z\in XN_c\) has a correction \(\tau\in P_c\)
satisfying
\[
 \ell^\sigma\Phi^j z-(\gamma-1)\tau\in DP_c\subset X^M N_{c+j}.
\]
Its class modulo \(DP_c\) is uniquely determined by the two
congruences for \(v,\tau\) below, is additive in \(z\), and is
compatible with maps of the Wach lattices.
\end{lemma}
\begin{proof}
The height bounds and the parameter formula give
\[
 \Phi(X^M N_i)\subset q^{(\ell-1)M-h}X^MN_{i+1},\qquad
 \Phi^j(X^MN_c)\subset DP_c\subset X^MN_{c+j}.
\]
The first inclusion yields contraction of \(\Phi\) on the cyclic
product of these deep lattices. For the second use
\[
 (\prod_{s<j}p_s^g)P_c\subset N_{c+j},
\]
by iteration.
Frobenius sends a basis to an independent family after inverting
the cyclotomic factors by the two-sided height comparison. It thus
gives a basis of \(P_c\). The image of \(D\) modulo a uniformizer
is a nonzero power series. Componentwise reduction in this basis
then shows that \(P_c/DP_c\) has no \(\ell\)-torsion.
Both lattices are \(\gamma\)-stable by commutation with \(\Phi\)
and stability of the ideals \((X),(p_s)\).
Given \(z\in XN_c\), first solve the equation for
\(\gamma-1\) modulo \(X^M\).

On degree \(r\) it acts by
\(\omega_X^r-1\). With
\(A_r=\sum_{k=1}^r(1+v_\ell(k))\) and \(K=A_{M-1}\), there are
\(d_r\in N_c\) such that
\[
 v=\sum_{1\le r<M}\ell^{K-A_r}X^r d_r,\qquad
 (\gamma-1)v-\ell^K z\in X^M N_c.
\]
To construct \(v\), at degree \(r\) multiply the previous
approximation by \(\ell^{1+v_\ell(r)}\) and subtract an
\(X^r\)-term given by the inverse of the unit part of
\(\omega_X^r-1\) times the degree-\(r\) residual coefficient.
This eliminates the residual at that degree. Legendre's formula
gives \(A_r\le\ell r/(\ell-1)\).

For \(1\le r<M\) one can divide
\(\ell^\sigma(\phi^j X)^r\) by \(\ell^{A_r}\) modulo \(D\).
For \(A_r\le\sigma\) this is literal division. Otherwise let
\(d_j=\ell^{j-1}\) and
\(m=(\ell-1)r-3g-(A_r-\sigma)\). From the bound on \(A_r\)
and the definition of \(\sigma\),
\[
 m\ge\left(\frac{(\ell-1)^2}{\ell}-1\right)A_r-3g+\sigma
   \ge\frac{(\ell-1)^2}{\ell}\sigma-3g\ge M/d_j.
\]

The binomial expansion of \(p_{j-1}^{A_r-\sigma+m}\), using
\(p_{j-1}\in(\ell,X^{d_j})\), lies in
\((\ell^{A_r-\sigma},X^{d_j m})\).
After multiplying by the other factors in \((\phi^j X)^r\),
the second summand is divisible by \(D\). Choose \(b_r\in
\mathcal R\) with
\(\ell^{A_r} b_r\equiv\ell^\sigma(\phi^jX)^r\pmod{D}\).
Then
\[
 \tau=\sum_{1\le r<M} b_r\Phi^j d_r\in P_c,\qquad
 \ell^K\tau-\ell^\sigma\Phi^j v\in DP_c .
\]
Commuting \(\gamma\) past \(\Phi^j\) and using the equation for
\(v\) shows that
\[
 \ell^K(\ell^\sigma\Phi^j z-(\gamma-1)\tau)\in DP_c.
\]

Saturation removes \(\ell^K\), and
\[
 DP_c\subset X^M N_{c+j}
\]
gives the required correction. If two corrections solve the
two displayed congruences for the same \(z\), take their
difference and bring their powers of \(\ell\) to a common power.
The associated difference of the \(v\)'s is in \(XN_c\) and
has \(\gamma-1\) image in \(X^MN_c\).

If \(v'\) denotes that difference, reduction in each successive degree gives
\[
 (\omega_X^r-1)[v']_r=0\ \Longrightarrow\ [v']_r=0
 \qquad(1\le r<M).
\]
The scalar is nonzero on a free graded piece, so induction gives \(v'\in X^MN_c\). Saturation gives a difference of corrections in
\(DP_c\). This also proves compatibility with maps and addition
of the congruences, by applying the equations before comparing
corrections. For a change of the cut vector by \(u\in XN_c\) the
correction to \((\gamma-1)u\) can be taken to be
\(\ell^\sigma\Phi^j u\), since one may use \(v=u\), \(K=0\)
in those congruences. This is the comparison modulo \(XN_c\)
needed in the Herr assembly.

\end{proof}
\input{wach-levels}

\paragraph{A level comparison for the mixed-sign arc.}
\label{par:wach-mixed-constant}
Use the arc length \(a\le A\), the powers \(H_s^{\mathrm{arc}}\),
and the level data of Paragraph~\ref{par:wach-level-models}.
Distinguish the transported ranks \(r_i\), the coordinate map
\(L:O^{r_1}\to O^{r_2}\) on the starting bases modulo \(X\),
and its coefficientwise extension \(L_{\mathcal R}\).
The polynomial comparison uses, at each \(s\), evaluation
maps from \(W_i\) and from \(\mathcal R^{r_i}\) to a common level
space. Denote their values by \(w_s\) and \(y_s\), respectively.
The latter refers to a vector \(y\) of coordinate series, not to
its coefficients at \(X=0\). The following are the lattice
properties used in this comparison.

\begin{itemize}
\item There are two local lattices \(B_{i,s},J_{i,s}\) with
\(p_s^g B_{i,s}\subset J_{i,s}\) and
\[
 p_s^g J_{i,s}\subset B_{i,s}.
\]
A submodule \(F_{i,s}\) lies
in both. Every vector on either side of the evaluation has a partner
on the other side, whose evaluation agrees with \(\ell^\sigma\)
times the original modulo \(F_{i,s}\), simultaneously at all
levels. The two evaluation maps commute with the corresponding
maps induced by \(f\) and \(L_{\mathcal R}\). The level map
preserves the \(F_{i,s}\).
\item The preimage of \(N_{i,c}\) under multiplication by
\(\Theta\) lies in \(W_i\). On \(W_i\) this preimage is tested
by
\[
 p_s^g w_s\in J_{i,s}
\]
for all \(s\).
The map \(f\) lifts on this preimage with cost \(\ell^R\).
If \(y_s\in F_{2,s}\), then \(p_s^{H_s^{\mathrm{arc}}}\) divides
\(\varpi^k y\) for some \(k\), where \(\varpi\) is a
uniformizer of \(O\).
\item There are saturated constant submodules \(U_{i,s}\) of
\(O^{r_i}\) and subgroups \(D_{i,s}\) of crystalline constants
in the level spaces. Membership \(y_s\in B_{i,s}\) is equivalent,
over all levels, to the conditions
\[
 y\in
 \begin{cases}
 p_s^{H_s^{\mathrm{arc}}}\mathcal R^{r_i} &\text{smaller jumps},\\
 p_s^{H_s^{\mathrm{arc}}-b_s}\mathcal R U_{i,s}+p_s^{H_s^{\mathrm{arc}}}\mathcal R^{r_i}
                          &\text{up},\\
 p_s^{H_s^{\mathrm{arc}}}\mathcal R U_{i,s}+p_s^{H_s^{\mathrm{arc}}+b_s}\mathcal R^{r_i}
                          &\text{down},
 \end{cases}
 \tag{\(*\)}\label{eq:wach-level-condition}
\]
after allowing multiplication of \(y\) by a power of \(\varpi\)
at each level. Membership \(y_s\in D_{i,s}+B_{i,s}\) is tested
by the same conditions on \(\varpi^k y-\Lambda d\) at each level,
for some \(k\ge0\), \(d\in O^{r_i}\). Write \(\mathcal N_i\)
and \(\mathcal I_i\) for these two modules of coordinate series.
\end{itemize}

\begin{lemma}[The up subspace after a down letter]\label{lem:up-after-down}
Suppose the two positions with larger jumps are a down letter
at \(s_0\) and an up letter at \(s_1>s_0\).
Assume \(L(U_{1,s_1})\subset U_{2,s_1}\),
\(2b_{s_0}\le\ell b_{s_1}\) and \(2a_0<b_{s_1}\), where
\(a_0=2\sigma+R+2gA\) is the initial exponent of
Section~\ref{sec:inert-long-orbit} and \(a\le A\) is the arc length. Then
\[
 \ell^{2a_0}U_{2,s_1}\subset L(U_{1,s_1}).
\]
\end{lemma}

\begin{proof}
Put \(Q=p_{s_1}^{b_{s_1}}\), \(P=p_{s_0}^{b_{s_0}}\) and
\(\Lambda_1=\Lambda/Q\). If \(v\in U_{2,s_1}\), the series
\(y_2=P\Lambda_1 v\) belongs to \(\mathcal N_2\), directly
by \eqref{eq:wach-level-condition} at each of the two positions with larger jumps and at every position with smaller jumps.

Choose its partner \(w_2\in W_2\). The error at each level lies
in \(F_{2,s}\) and \((y_2)_s\in B_{2,s}\), so
\[
 p_s^g(w_2)_s\in J_{2,s}.
\]
The preimage test gives a lift
\(w_1\) on the \(\Theta\)-divisible lattice, with
\(f(w_1)=\ell^R w_2\). Choose the coordinate partner \(y_1\)
of \(w_1\). Each \(p_s^{2g}(y_1)_s\) belongs to \(B_{1,s}\),
so \(\Theta^2 y_1\in\mathcal N_1\).
The partner errors and compatibility with \(f\) give
\[
 (L_{\mathcal R}y_1-\ell^{2\sigma+R}y_2)_s\in F_{2,s}.
\]

Every \(p_s^{H_s^{\mathrm{arc}}}\) therefore divides the difference after
multiplication by a power of \(\varpi\). These divisibilities
hold integrally: each quotient by a distinguished polynomial
power is free over \(O\) by Weierstrass division. Their product
\(\Lambda\) divides by relative primality after inverting
\(\varpi\), followed by the same saturation for the product.
This uses no integral Chinese remainder decomposition.

\proofstep{Testing cyclic quotients of the image.}
Take an integral linear form \(\lambda\) on \(O^{r_2}\)
such that \(\lambda(v)=1\) and
\(\lambda L(U_{1,s_1})\subset tO\).
For any \(y\in\mathcal N_1\),
\[
                 \lambda L_{\mathcal R}(y)
                   \in \Lambda_1(t,Q).
\]
Indeed split \(O^{r_1}\) along the saturated \(U_{1,s_1}\).
On its complement the up condition implies divisibility by
\(p_{s_1}^{H_{s_1}}\), and on the summand \(U_{1,s_1}\) it
implies divisibility by \(p_{s_1}^{H_{s_1}-b_{s_1}}\).
All coordinates are divisible by \(p_s^{H_s^{\mathrm{arc}}}\) at the other
levels. Multiplying the powers of the distinct cyclotomic factors, the complementary
summand gives a multiple of \(\Lambda\) and the up summand gives
a multiple of \(t\Lambda_1\). Applying this to
\(\Theta^2 y_1\) and using the congruence for
\(L_{\mathcal R}y_1\) yields, after cancelling \(\Lambda_1\),
\begin{equation}\label{eq:wach-gauss-ideal}
                 \ell^{2\sigma+R}\Theta^2P\in(t,Q).
\end{equation}

\proofstep{The Gauss valuation and Smith factors.}
Normalize \(v_\varpi(\varpi)=1\), and put \(e_O=v_\varpi(\ell)\).
For \(x>0\) the Gauss valuation of \(\sum a_m X^m\) is
\(\min_m(v_\varpi(a_m)+mx)\). If \(i,j\) are the least minimizing indices for \(f,g\), then in
the coefficient of \(X^{i+j}\) only \(f_i g_j\) attains the minimum.
Indeed a different pair with the same index sum uses an index below
\(i\) or below \(j\). Therefore
\[
 v_\varpi([X^{i+j}]fg)+(i+j)x=v_x(f)+v_x(g)=v_x(fg).
\] Preparation gives
\(v_x(p_s)=\min(e_O,\ell^s x)\).

If \(t\ne0\), set
\(k=v_\varpi(t)\); for \(t=0\) take \(k=+\infty\).
The assertion for \(k=0\) is immediate. Otherwise set
\(b_1=b_{s_1}\), \(b_0=b_{s_0}\),
\(m=\min(k,e_O b_1)\), and \(x=m/(b_1\ell^{s_1})\). Then
\(v_x(Q)=m\), and \(v_x(P)\le b_0 m/(\ell b_1)\).
The ideal relation and \(v_x(p_s)\le e_O\) give
\[
 m\le e_O(2\sigma+R+2ga)+\frac{b_0}{\ell b_1}m
   \le e_O a_0+m/2 .
\]
Thus \(m\le2e_Oa_0<e_Ob_1\), which excludes \(t=0\) and
proves \(t\mid\ell^{2a_0}\).

Since \(U_{2,s_1}\) is a direct
summand, Smith coordinate forms extend to the ambient
\(O^{r_2}\). Apply the divisibility assertion to its Smith
basis vectors for the inclusion of \(L(U_{1,s_1})\).
A missing image direction would give \(t=0\); the invariant
factors in the image directions all divide \(\ell^{2a_0}\).
Consequently
\[
                   \ell^{2a_0}U_{2,s_1}\subset L(U_{1,s_1}).
\]
The use of \(s_0<s_1\) gives the factor \(1/\ell\) in the
valuation comparison. The argument measures divisibility by a
scalar in a ramified coefficient ring and so is uniform under
the coefficient extension.
\end{proof}

\paragraph{From constant coordinates to the arc lattice.}
\begin{lemma}[Lifting constant coordinates across an arc]\label{lem:wach-constant-lift}
With the integral lattices and scalar \(a_0\) of this comparison,
each \(y\in\mathcal I_2\) has a lift modulo \(\mathcal N_2\):
\begin{equation}\label{eq:wach-const-coord-lift}
 \ell^{2a_0}y=L_{\mathcal R}y_1+\ell^{2a_0}\nu,\qquad
 y\in\mathcal I_2,\quad y_1\in\mathcal I_1,\quad
 \nu\in\mathcal N_2 .
\end{equation}
\end{lemma}
\begin{proof}
Here is the constant comparison for all the words. For either
representation, omit its index and let \(V=O^r\). If there is
no down letter, the level conditions are unchanged by adding a
local constant times \(\Lambda\), so \(\mathcal I=\mathcal N\).
This includes a word with two up letters. Otherwise there is at
most one up letter. Write \(Q=p_s^{b_s}\) at the up letter and
\(Q=1\) if there is none. Divisibility at every level gives
\(y=(\Lambda/Q)Y\), with \(Y\in\mathcal R^r\).
At an up letter the condition for \(Y\) is
\[
             Y\in\mathcal R U_s+Q\mathcal R^r .
\]

At a down letter with \(P=p_s^{b_s}\) it is the condition that
the image of \(Y\) in \((\mathcal R/(P))\otimes_O(V/U_s)\)
lie in the rational constant line factor \(\overline Q\) times
\((V/U_s)\otimes_O\operatorname{Frac}(O)\). To see this,
divide by the full local power \(p_s^{H_s^{\mathrm{arc}}}\), then cancel the
other level factors modulo \(P\) after inverting \(\varpi\).
The difference between \(\mathcal I\) and \(\mathcal N\) is
that \(\mathcal N\) requires the image at each down letter to
be zero. The quotients in this test are free over \(O\) by
Weierstrass division and saturation of \(U_s\).

With a single down letter and no up letter, primitivity means the coefficient map splits the constant inclusion:
\[
 O\xrightarrow{\ a\mapsto a\ }\mathcal R/(P)\xrightarrow{\lambda}O,\qquad \lambda(1)=1.
\]
In fact, if the projected class is \(1\otimes v\), choose an integral
linear form \(\lambda\) on \(\mathcal R/(P)\) with \(\lambda(1)=1\).
Then
\[
 (\lambda\otimes1)(1\otimes v)=v\in V/U_s .
\]
This recovers the constant integrally in the saturated quotient. If the up letter
has smaller offset than the down one, its power \(Q\) reduces
modulo \(\varpi\) to a unit times a monomial of smaller degree
than the distinguished polynomial for \(P\), by the inequality
\(b_{\rm up}<\ell b_{\rm down}\). The class of \(Q\) modulo
\(P\) is again primitive, which proves integrality of the
constant by the same linear-form test. In both cases lift the
constant to \(d\in V\). Then
\[
 y-\Lambda d\in\mathcal N.
\]


\proofstep{Two down letters.}
For two down letters let \(P,Q\) be their powers in increasing
offset order, and \(U_0,U_1\) the corresponding submodules.
There are integral constants \(d_0,d_1\in V\) by the separate
primitive test for \(1\) in each quotient. Their difference
belongs to the ideal \((P,Q)\) in
\(\mathcal R\otimes_O V/(U_0+U_1)\). Let \(\rho\) be least with
\[
 \varpi^\rho\in(P,Q),
\]
which exists by coprimality after
inverting \(\varpi\). The constant components in a cyclic
decomposition of \(V/(U_0+U_1)\) belong to
\(\varpi^\rho V/(U_0+U_1)\). On a free summand this follows
from the definition of \(\rho\).

On a summand of order
\(\varpi^k\), if an integral representative \(a\) in
\((P,Q,\varpi^k)\) had valuation below \(\min(k,\rho)\), then
\(a-\varpi^k f\in(P,Q)\) for some \(f\in\mathcal R\).
The factor \(1-(\varpi^k/a)f\) is a unit, a contradiction
to the choice of \(\rho\).
Write
\[
 d_0-d_1=u_0+u_1+\varpi^\rho d'
\]
with
\(u_i\in U_i\). Choose \(f_\rho\in(Q)\) with
\(f_\rho-\varpi^\rho\in(P)\). Put \(Y'=Y-d_1-u_1-f_\rho d'\).
The two reductions are
\[
 \begin{aligned}
  Y'&\equiv d_0-d_1-u_1-\varpi^\rho d'=u_0
        &&\pmod {P\mathcal R^r+\mathcal R U_0},\\
  Y'&\equiv d_1-d_1-u_1=-u_1
        &&\pmod {Q\mathcal R^r+\mathcal R U_1}.
 \end{aligned}
\]
Both images vanish. Multiplying back by \(\Lambda\) gives the
two factors \(\Lambda\) and \(\Lambda f_\rho\).

\proofstep{A down letter before an up letter.}
For a down letter at the smaller offset and an up letter at
the larger one, write \(U=U_{s_1}\) and \(V_0=U_{s_0}\),
with \(P,Q,\Lambda_1\) as in the mixed-sign estimate. Write
\(Q=\varpi^k\alpha+P\beta\), with \(k\) largest possible.
There is such a finite \(k\) because the image of \(Q\) in
the finite free \(\mathcal R/(P)\) is nonzero. The class
of \(\alpha\) is primitive. The down condition for \(Y\)
therefore gives
\[
 \overline Y=\overline\alpha d_0
\]
in
\((\mathcal R/(P))\otimes_O(V/V_0)\) for some \(d_0\in V/V_0\).
The up condition, projected further to \(V/(U+V_0)\),
shows that \(\overline\alpha d_0\) belongs to
\(\varpi^k\overline\alpha((\mathcal R/(P))\otimes_O V/(U+V_0))\).

A linear form taking value one on \(\overline\alpha\)
gives \(d_0=v+\varpi^k d\) modulo \(V_0\), with
\(v\in U,d\in V\). This test works on the tensor product
with \(V/(U+V_0)\) even if that module has torsion. Now
\[
 \overline{Y-\alpha v-Qd}
    = \overline\alpha(d_0-v-\varpi^k d)=0
     \quad\text{in }(\mathcal R/(P))\otimes_O(V/V_0).
\]
The same difference still lies in the required up lattice:
\(\alpha v\in\mathcal R U\) and \(Qd\in Q\mathcal R^r\).
Upon multiplying back by the coordinate factors, the resulting
normal form uses \(\Lambda\) times an ordinary constant and
\(\alpha\Lambda_1\) times a vector of \(U\).
In each of the cases the displayed representatives themselves
satisfy the \(\mathcal I\) conditions, by reducing at the same
levels.

On those normal forms, the ordinary constants lift using the
cost \(\ell^R\) on the starting constant map and \(R\le2a_0\).
In the mixed case the additional constant in \(U_{2,s_1}\)
lifts by the Gauss valuation and Smith argument. This gives
\eqref{eq:wach-const-coord-lift}.
\end{proof}

\begin{lemma}[Integral lift between the two polynomial actions]
Let \(x'\) be a target arc vector evaluating in \(D_{2,s}+B_{2,s}\)
at each level of the cut. There is one source vector \(z\) for the arc map \(f\) such that
\[
 f(z)=\ell^{2\sigma+2a_0+R}x',\qquad z_s\in D_{1,s}+p_s^{-g}J_{1,s}\quad\text{at every level }s,
\]
where \(z_s\) is its evaluation.
\end{lemma}
\begin{proof}
Choose the coordinate partner \(y\) of \(x'\). Its error belongs
to \(F_{2,s}\subset B_{2,s}\), so \(y\in\mathcal I_2\).
Apply \eqref{eq:wach-const-coord-lift} to \(y\) and choose the
source arc partner \(x_1\) of \(y_1\). Put
\(\Delta=\ell^{2\sigma+2a_0}x'-f(x_1)\).
In its evaluation the target partner error has coefficient
\(-\ell^{\sigma+2a_0}\), the mapped source partner error has
coefficient \(-1\), and the remaining term is
\(\ell^{\sigma+2a_0}\nu_s\). All belong to \(B_{2,s}\);
the preimage test therefore puts \(\Delta\) in the
\(\Theta\)-divisible lattice. Lift \(\ell^R\Delta\) there to
\(\Delta_1\). Then
\[
 f(\ell^Rx_1+\Delta_1)
 =\ell^R(f(x_1)+\Delta)=\ell^{2\sigma+2a_0+R}x'.
\]
The source partner congruence and
\(B_{1,s}\subset p_s^{-g}J_{1,s}\) give the stated evaluation
at every level, with the same exponent.
\end{proof}
\begin{lemma}[The lift on invariant classes]
The invariant first-coordinate test used in the cut comparison is
\[
 (\gamma-1)z\in XN.
\]
The source and target polynomials in \(\gamma\) and the intervening lattice lift preserve this test in the sense of that comparison.
\end{lemma}
\begin{proof}
To apply this lift, use the level properties of
Paragraphs~\ref{par:wach-level-models} and~\ref{par:wach-polynomial-actions}. The level actions commute with the
evaluation of arc vectors. If the difference under \(\gamma-1\)
is in \(J_{i,s}\) at a level, the vector there lies in
\(D_{i,s}+J_{i,s}\). Writing \(P_2,P_1\) for the target and source polynomials, the actions at a level are
\[
 \begin{aligned}P_2(\gamma)(D_{2,s}+p_s^{-g}B_{2,s})&\subseteq D_{2,s}+B_{2,s},\\ P_1(\gamma)(D_{1,s}+p_s^{-g}J_{1,s})&\subseteq D_{1,s}+J_{1,s}.\end{aligned}
\]
The action fixes \(D_{i,s}\) and preserves
\(J_{i,s}\). Finally \(XN_{i,c}\) consists of the vectors in
\(XW_i\) whose evaluations belong to all the \(J_{i,s}\).
These are statements on the two level models at the same cut.

The action on \(W_i/XW_i\) is trivial. To see this, use the
trivial action on the starting Wach lattice modulo \(X\) and
the commutation with \(\Phi\), together with
\(\phi^a(X)\in(X)\). Division by \(\Lambda\) preserves the trivial action at the origin, since
\[
 \gamma(\Lambda)/\Lambda\equiv1\pmod X.
\]
For an arc vector \(w\) with
\((\gamma-1)w\in XN_{2,c}\), the first level property gives
evaluations in \(D_{2,s}+J_{2,s}\).

Since
\(J_{2,s}\subset p_s^{-g}B_{2,s}\), the target polynomial
produces the input \(x'\) of the intervening lattice lift.
The output, after applying the source polynomial, has
evaluations in \(D_{1,s}+J_{1,s}\). Its \(\gamma-1\) image
therefore has evaluations in all \(J_{1,s}\) and belongs to
\(XW_1\). The last level property puts the image in
\(XN_{1,c}\). An operator polynomial on a vector fixed
modulo \(XN_{2,c}\) agrees with its value at one modulo that
submodule. Writing \(P_2,P_1\) for the target and source
polynomials and \(c_{\rm lift}=2\sigma+2a_0+R\) for the intervening exponent,
choose \(y\) with
\[
 f(y)=\ell^{c_{\rm lift}} P_2(\gamma)w.
\]
Then
\[
 f(P_1(\gamma)y)
 =\ell^{c_{\rm lift}} P_1(\gamma)P_2(\gamma)w
 \equiv\ell^{c_{\rm lift}} P_1(1)P_2(1)w\pmod{XN_{2,c}} .
\]

\end{proof}
\paragraph{The colength at the cuts of the long orbit.}\label{par:wach-cut-colength}

\begin{lemma}[The scalar annihilator at a cut]\label{lem:wach-cut-annihilator}
Let \(f\) be the integral crystalline surjection on the local
arc of Paragraph~\ref{par:wach-level-models}. For \(R\) defined
in Paragraph~\ref{par:inert-cut-exponents}, write \(f_N\) for the Wach map at the cut and \(f_N^{\rm tame}\) for its map on tame invariants. Then
\[
 \ell^R\operatorname{coker} f_N=0,\qquad \ell^R\operatorname{coker} f_N^{\rm tame}=0.
\]
In the extension it gives a lift of \(\ell^R e_0\), and gives
the corresponding constant-map bound after reduction modulo \(X\).
\end{lemma}

\begin{proof}
For an integral crystalline surjection, let \(s_i\) be the
least exponent of the ordinary cyclotomic parameter \(\pi\)
that kills the Wach cokernel at letter \(i\), before tame
invariants. The global height box in this case is symmetric
and has width
\[
 H_{\max}=\max_i H_i=2h.
\]
A common Tate
twist moves both representations into the positive range
of Berger's convention. This multiplies both Wach lattices
by the same power of \(\pi\) in the twisted coordinates and
preserves the cokernel. Berger's integral comparison
\cite[Theorem III.3.1]{berger-wach} places the base change of
each positive Wach lattice between the corresponding base
change of \(T_i\) and its sublattice multiplied by
\(\pi^{H_{\max}}\). Apply the comparison over
\[
 \widetilde{\mathbf A}^{+}=W(\widetilde{\mathbf E}^{+}),
\]
which contains the integral period ring of that comparison.

Since the integral Galois lattice map is onto,
\(\pi^{H_{\max}}N_2\) lies in the image after this base change.
This base change is faithfully flat over
\(\mathcal O_F[[\pi]]\). Modulo \(\ell\) the tilt valuation
ring is torsion-free over the DVR \(k_F[[\overline\pi]]\),
and \(\overline\pi\) is a nonunit in it. The ring of Witt
vectors is \(\ell\)-torsion-free and \(\ell\)-adically complete,
so the adic flatness criterion gives flatness
\cite[Lemma 15.28.5]{stacks-flat}. The residue of the maximal
ideal is proper after extension, giving faithfulness.
Extending by \(O\) and taking the embedding factors preserve
faithful flatness. In the ordinary cokernel
\[
 C_i=N_{2,i}/f(N_{1,i}),
\]
the inclusion
after base change reads
\[
 (\pi^{H_{\max}}C_i)\otimes_{\mathcal O_F[[\pi]]}
      \widetilde{\mathbf A}^+=0.
\]
Faithful flatness implies \(\pi^{H_{\max}}C_i=0\), hence \(s_i\le H_{\max}\). The same argument applies to every crystalline
extension of the trivial lattice by \(T_2\), since its
jumps are the union of the target jumps and zero.

\proofstep{The scalar loss from the cyclotomic action.}
Use a topological generator \(\gamma_*\) of the full
cyclotomic group whose cyclotomic value \(u_*\) generates
\(\Z_\ell^\times\). This action is trivial on Wach lattices
modulo \(\pi\). Suppose \(\ell^k\pi^r N_{2,i}\subset f(N_{1,i})\)
for \(r>0\). For any \(v\in N_{2,i}\), choose \(x\) lifting
\(\ell^k\pi^r v\). The vector \((\gamma_*-1)x/\pi\) is
Wach-integral. Its image, modulo \(\ell^k\pi^rN_{2,i}\), is
\[
                   \ell^k\pi^{r-1}(u_*^r-1)v .
\]
This follows by reducing \(\gamma_*(\pi)/\pi\) and the action
on \(v\) modulo \(\pi\).

The error is already in the image.
One therefore lowers \(r\) by one at the cost
\(v_\ell(u_*^r-1)\). The cost is zero unless
\[
 \ell-1\mid r,
\]
and it is \(1+v_\ell(r/(\ell-1))\)
otherwise. Thus the initial parameter exponent \(s\) gives
a scalar exponent at most \(\ell s/(\ell-1)^2\) by
Legendre's formula.

\proofstep{The recurrence in the letter exponents.}
To pass through Frobenius, put \(d_0=\ell-1\) and let
\(s=s_i\). From \(s\) down to \(r<s\), the exact cost is
\(n-m+v_\ell(n!/m!)\), with \(n=\lfloor s/d_0\rfloor\)
and \(m=\lfloor r/d_0\rfloor\). For \(n=0\) this is zero.
Otherwise it is at most
\[
         \frac{\ell(s-r)}{d_0^2}+\frac{\ell}{d_0}
                         +\lfloor\log_\ell n\rfloor .
\]
Indeed each term in the factorial valuation differs from
\((n-m)/\ell^v\) by less than one, and there are at most
\(\lfloor\log_\ell n\rfloor\) nonzero terms.

Take
\(a=\lceil s/d_0\rceil+c_H+2\). Since
\(q_{\rm cyc}\in(\ell,\pi^{d_0})\), the power
\(\phi(\pi)^a=\pi^a q_{\rm cyc}^a\) is in the ideal
generated by \(\ell^{a-j'}\pi^{a+d_0j'}\) for
\(0\le j'\le a\).
Those terms with \(r=a+d_0j'\ge s\) kill the cokernel.
For \(r<s\),
\[
 a-j'-\frac{\ell(s-r)}{d_0^2}
 =a-\frac{\ell(s-a)}{d_0^2}+\frac{j'}{d_0}
 \ge\frac{a}{d_0^2}+\frac{\ell(c_H+2)}{d_0}.
\]
This exceeds \(\ell/d_0+\lfloor\log_\ell n\rfloor\) since
\(n\le H_{\max}\). Hence each term kills the cokernel by
the cost bound, and so does \(\phi(\pi)^a\).

Write \(H_i=M_i+h_i\) for the letter's width, with the two ordinary one-step height inclusions
\[
 q_{\rm cyc}^{M_i}\Phi N_{2,i-1}\subset N_{2,i},\qquad q_{\rm cyc}^{h_i}N_{1,i}\subset\operatorname{span}(\Phi N_{1,i-1}).
\]

Lift \(\phi(\pi)^a q_{\rm cyc}^{M_i}\Phi v\) for
\(v\in N_{2,i-1}\) and multiply its lift by
\(q_{\rm cyc}^{h_i}\pi^{H_i}\).
It now lies in the source Frobenius span and maps to
\(\Phi(\pi^{a+H_i}v)\).
Frobenius linearization is injective on both lattices,
as it is invertible on the ambient \((\Phi,\gamma)\)-modules
after localization. The scalar ring is free over its
Frobenius image with a basis starting at \(1\), by
Weierstrass division by \(\phi(\pi)-Y\).
Compare the coefficient of \(1\) in the Frobenius-span
expression for the lift.

This constant coefficient gives a lift of
\(\pi^{a+H_i}v\) in the preceding letter. Consequently
\[
                  s_{i-1}\le H_i+c_H+3+\frac{s_i}{\ell-1}.
\]
Iterating the recurrence leaves a remainder bounded by a
geometric fraction of \(H_{\max}\), which tends to zero.
This gives \(s_i\le S(i)\), in addition to
\(s_i\le H_{\max}\). The scalar annihilator above therefore
has exponent at most \(R\) at each cut.

\proofstep{Passage to tame invariants.}
Averaging a lift of an invariant vector by the tame group costs
a unit. Thus the ordinary-module annihilator yields the desired
annihilator on the tame-invariant Wach lattices. This gives
the lift of \(\ell^R e_0\) in the extension Wach lattice and
the ordinary constant-map bound at the preceding cut after
reducing the surjection modulo \(X\). Dividing by \(\Theta\)
commutes with the map and is possible in the ambient module
because the factors of \(\Theta\) are units there. It gives
the required lift on the \(\Theta\)-divisible lattices at
the terminal cut.
\end{proof}


\begin{proposition}[Finite-class colength in the long-orbit range]\label{prop:inert-long-colength}
Use the marked crystalline surjection and signed jump windows
of Lemma~\ref{hyp:long-orbit-data}, with the cyclic spacing
and arc exponents of Paragraph~\ref{par:inert-cut-exponents}.
Then \eqref{eq:long-orbit-exponent} gives a normalized colength on
integral finite classes at most \(3C_2U_B\), with \(U_B\) as in
Section~\ref{sec:cut-length-rows}.
There exists one rational place entry \(c_w\) for all good characters at \(w\), satisfying
\[
 3C_2U_B\le c_w\le(3C_2+1)U_B.
\]
\end{proposition}

\begin{proof}
The lifting argument behind \eqref{eq:long-orbit-exponent} is
in the two operators \(\Phi-1\) and \(\gamma-1\) of the Herr
complex. Here \(\Phi\) advances one letter and \(\gamma\) acts
within a letter. Work in the tame-invariant Wach coordinates,
with parameter \(X\) and Wach lattices \(N_i\) at the letters.
The first filtration estimate places \(\Phi(X^M N)\) in
\(X^M(\ell,X)N\) on the product \(N\) of the letters. On \(X^MN\) contraction therefore gives the convergent inverse
\[
 (\Phi-1)^{-1}=-\sum_{a\ge0}\Phi^a.
\]
The coefficients are complete and
\(\Phi\) preserves the ideal \((\ell,X)\).

\proofstep{Represent the scaled class by a restricted Herr pair and lift its arcs.}
A finite class in the target has a Herr representative attached to
a crystalline extension of the trivial lattice by \(T_2\). If \(E_0\) lifts the trivial generator in period coordinates, the target pair representing the connecting class is
\[
 ((\Phi-1)E_0,(\gamma-1)E_0).
\]
The
cut colength bound provides, at each cut \(c\), a vector \(x_c\)
in the extension Wach lattice over \(\ell^R\) times the trivial
generator. The difference \((\gamma-1)x_c\) belongs to \(XN_c\)
of the target.

In fact it belongs to \(X\) times the extension
Wach lattice by triviality modulo \(X\), and to the target
period module because its image in the trivial quotient vanishes.
Since \(X\) is regular on the trivial quotient and the
extension Wach lattice intersects the target period module in the
target Wach lattice, this gives the assertion. After advancing
\(j\) letters and multiplying by \(\ell^\sigma\), the correction
of Paragraph~\ref{par:wach-cut-correction} gives
\[
       Q_c=\ell^\sigma\Phi^j x_c-\tau_c,
       \qquad (\gamma-1)Q_c\in X^M N_{c+j}.
\]
Use the target Wach lattice as its image in the extension.

For consecutive cuts \(c',c\), write \(a\ge j\) for their
forward distance. At the letters from \(c'+j\) up to the
letter before \(c+j\), set \(E_{c'+j+t}=\Phi^tQ_{c'}\).
This gives a global vector \(E\) with trivial quotient
\(\ell^{R+\sigma}\) times the generator. The first component
\(\ell^\sigma(\Phi-1)E\) vanishes at all the letters within
each propagated arc. At a shifted cut it is
\[
 \begin{split}
   y_c&=\ell^\sigma(\Phi^a Q_{c'}-Q_c)
       =\ell^\sigma(\Phi^j w_c+\tau_c),\\
   w_c&=\Phi^{a-j}Q_{c'}-\ell^\sigma x_c.
 \end{split}
\]
The two terms defining \(w_c\) have the same trivial
quotient, so \(w_c\) is in the target. Since \(\gamma\)
commutes with \(\Phi\), the deep-lattice defects of the
\(Q_c\)'s propagate along each arc and
\((\gamma-1)y_c\in X^M N\).

Set \(z_c=(\Phi-1)^{-1}(\gamma-1)y_c\) there.
Uniqueness of the inverse on \(X^MN\) gives
\[
 \sum_c(y_c,z_c)
    =\bigl(\ell^\sigma(\Phi-1)E,\ell^\sigma(\gamma-1)E\bigr).
\]
The difference from \(\ell^{R+2\sigma}\) times the original
pair is the Herr coboundary of
\(\ell^\sigma E-\ell^{R+2\sigma}E_0\). This is a target
vector because its two terms have equal trivial quotient.

Let \(c'\) be the preceding cut and \(a\) the length of its arc
to \(c\). The propagated \(Q_{c'}\) at \(c\) belongs to the span
of \(\Phi^a\) of the extension Wach lattice at \(c'\), since its
correction at \(c'+j\) is in the span of \(\Phi^jN_{c'}\).
The signed bound into the Frobenius span in
Paragraph~\ref{par:wach-level-models} gives the same membership
after multiplying \(x_c\) by
\[
 \Lambda=\prod_{s<a}p_s^{H_s^{\mathrm{arc}}}.
\]

The filtration of this extension has the same bounds as the target
with the jump at zero included.

For \(\Lambda w_c\), membership in the extension Frobenius
span and in the target period module already holds. To
intersect these two modules, put \(\mathcal R=O[[X]]\)
and use a basis \((b_t)_t\) of \(\mathcal R\) over
\(\phi^a(\mathcal R)\). Such a basis comes from Weierstrass
division by \(\phi^a(X)-Y\) over \(O[[Y]]\), since
\(\phi^a(X)\) first has a unit coefficient in degree \(\ell^a\).
Group the extension expression as \(\sum_t b_t\Phi^a u_t\).
If \(d_t\) is the trivial quotient of \(u_t\), then
\[
       \sum_t b_t\Phi^a u_t\longmapsto
                    \left(\sum_t b_t\phi^a(d_t)\right)e_0=0 .
\]
Independence of the basis and injectivity of \(\phi^a\)
give \(d_t=0\) for every \(t\). Therefore every \(u_t\)
lies in the target Wach lattice at \(c'\). This proves
\(w_c\in W_2\). The two defects to be subtracted belong respectively to \(X^MN_c\) and \(XN_c\). Therefore
\[
 (\gamma-1)w_c\in X^M N_c+XN_c\subseteq XN_c.
\]


For a change of cut vector by \(u\in XN_c\), the correction
of Paragraph~\ref{par:wach-cut-correction} changes by
\(\ell^\sigma\Phi^j u\) modulo the transported error.
It cancels the change of the uncorrected advanced vector.
Hence the corrected first component changes by a vector
\[
 v\in X^M N.
\]
The cocycle equation determines the change
of the second component as \((\Phi-1)^{-1}(\gamma-1)v\).
Together these give the coboundary of \((\Phi-1)^{-1}v\).
Thus comparison of cut vectors modulo \(XN_c\) suffices
for the finite classes of the corrected arc pairs.

Choose a product \(D_{\rm prev}\) of cyclotomic factors
clearing \(\Lambda w_c\) into \(N_c\), by iterating the
one-step height inclusions on the preceding arc. Choose
\(D_{\rm adv}\) clearing \(\Phi^j N_c\) into \(N_{c+j}\)
and \(D_\tau\) clearing \(P_c\) into \(N_{c+j}\).
Then the product
\[
            D=\phi^j(D_{\rm prev}\Lambda)\,D_{\rm adv}D_\tau
\]
clears \(y_c=\ell^\sigma(\Phi^j w_c+\tau_c)\) into \(N_{c+j}\)
by semilinearity and \(\tau_c\in P_c\).
The second component \(z_c\) and the defect
\((\gamma-1)y_c=(\Phi-1)z_c\) are in \(X^MN\).
Thus the two components satisfy the actual finite
cyclotomic-denominator and deep-lattice tests.

Pass to the ordinary cyclotomic parameter \(\pi\). Put
\(V=T\otimes_O K\), with \(K=\operatorname{Frac}(O)\), and set
\[
 q_{\rm cyc}=\phi(\pi)/\pi,\qquad t=\log(1+\pi),\qquad
 \lambda=t/\pi .
\]
The factors \(p_s\) occurring in \(D\) are associates of
\(\phi^s(q_{\rm cyc})\). Also \((X^M)=(\pi^{(\ell-1)M})\).
Thus, for some finitely supported family of integers \(b_s\ge0\),
the Herr pair \((y,z)\) satisfies
\begin{equation}\label{eq:long-arc-herr-lattices}
 \Bigl(\prod_s\phi^s(q_{\rm cyc})^{b_s}\Bigr)y\in N(T),\qquad
 z\in\pi N(T),\qquad (\gamma-1)y=(\Phi-1)z .
\end{equation}
The integral Herr isomorphism gives a class in \(H^1(F,T)\).
Its rational image is represented by the same pair in the Robba-ring
Herr complex \cite[Theorems 2.3 and 2.8]{liu-herr}.
In the corresponding extension \(0\to V\to V'\to K\to0\), choose
a tame-invariant lift \(e_0\) of the trivial basis. The extension
actions are \(\Phi e_0=e_0+y\) and \(\gamma e_0=e_0+z\).

Let \(\mathcal B^+\) denote the analytic functions on \(|\pi|<1\).
Use crystalline coordinates in the kernel. For a positive crystalline
representation the Wach comparison and its elementary divisors put the
change of basis and its inverse over \(\mathcal B^+[1/\lambda]\)
\cite[Propositions II.2.1 and III.2.1]{berger-wach}.
A Tate twist gives the same conclusion for \(V\): the ratio of
the two twist generators is
\[
 \pi^{-k}/t^{-k}=\lambda^k.
\]

Since \(\lambda(0)=1\), the change of basis is regular at zero
in both directions. Writing brackets for crystalline coordinates gives
\[
 \lambda=\prod_{s\ge0}\frac{\phi^s(q_{\rm cyc})}{\ell},
 \qquad \lambda^B[y]\in(\mathcal B^+)^{\dim V}\quad(B\gg 0),
 \qquad [z](0)=0 .
\]
By \eqref{eq:long-arc-herr-lattices}, a sufficiently high power of the product absorbs all finitely
many cyclotomic denominator factors of \(y\) and the poles of the
change of basis. The factor \(\pi\) in \(z\) gives the last equality.

In the crystalline basis the equation \((\gamma-1)v=z\)
has a solution analytic near zero. Use \(t\) as local
coordinate and put \(\omega=\chi(\gamma)\). If
\(z(t)=\sum_{k\ge1} z_k t^k\), the solution with zero
constant is
\[
          v(t)=\sum_{k\ge1}\frac{z_k}{\omega^k-1}t^k .
\]
The valuations \(1+v_\ell(k)\) of the denominators grow
logarithmically in \(k\). Thus any strictly smaller positive disk of convergence of
\(z\) is a disk of convergence of \(v\). For \(u=y-\Phi v+v\)
the Herr equation and \(\gamma\Phi=\Phi\gamma\) give
\[
 (\gamma-1)u=(\Phi-1)z-(\Phi-1)(\gamma-1)v=0 .
\]

In the crystalline basis the coefficient \(u_k\) in degree \(k\)
is therefore killed by \(\chi(\gamma)^k-1\). For \(k>0\) this
factor is nonzero, so the coordinates of \(u\) are constant. For every \(m\ge1\)
the identity of germs is
\[
 v=\Phi^m v-\sum_{i=0}^{m-1}\Phi^i y+\sum_{i=0}^{m-1}\Phi^i u.
\]
On every disk of fixed radius less than one,
\(\phi^m(\pi)=(1+\pi)^{\ell^m}-1\) tends uniformly to zero.
Hence \(\Phi^m v\) is analytic there for sufficiently large
\(m\). Also \(\lambda^B\Phi^i y\) is analytic on the
whole open disk for every \(i\), since
\[
 \lambda/\phi^i\lambda=\prod_{k<i}\phi^k(q_{\rm cyc}/\ell).
\]

The displayed identity therefore extends \(\lambda^B v\)
to the whole disk. The continuations agree on overlaps by
their Taylor series at zero. It follows that \(v\) belongs
to the Robba module with \(t\) inverted. The identity
\[
 (\gamma-1)v=z
\]
holds there by its power series identity
at zero after clearing powers of \(\lambda\), so \(e_0-v\)
is \(\gamma\)-invariant. Its average over the tame
group is invariant under the full cyclotomic group and still
maps to the trivial basis. Berger's comparison
\cite[Theorem 0.2]{berger-differential} identifies these
invariants with \(D_{\rm cris}(V')\). The rank increases
by one from \(D_{\rm cris}(V)\), so \(V'\) is crystalline.
Thus the integral Herr class is finite. This criterion applies
to both lattices and is compatible with the same row map.

The quantitative estimate at an arc supplies a source vector
of this kind mapping to \(m_c w_c\) modulo \(XN_c\), where
\[
 m_c=u_c\ell^{2\sigma+2a_0+R+k_c},\qquad u_c\in O^\times,
                 \qquad 0\le k_c\le2\kappa.
\]
It applies the polynomials in \(\gamma\) to the target and then
to its lift to impose the required local exponent windows.
Modulo \(XN_c\) their actions reduce to their values at \(1\),
since \(w_c\) is \(\gamma\)-invariant modulo this submodule.

Let \(\omega\) be the cyclotomic value of \(\gamma\), with
\(v_\ell(\omega-1)=1\), and let \(E=a-1\le A\).
The polynomials have the form
\[
 P_{E,J}(Z)=\left(\sum_{v=0}^{\ell^E-1}Z^v\right)
            \prod_{s\in J\setminus\{0\}}
                 (Z^{\ell^E}-\omega^{s\ell^E}),
\]
At \(1\) the valuation is
\(E+\sum_{s\in J\setminus\{0\}}(1+E+v_\ell(s))\)
by the valuation of powers of a principal unit. The window
bound gives at most \(\kappa\) for each polynomial. Between the two
polynomials the comparison of Paragraph~\ref{par:wach-mixed-constant}
costs \(2\sigma+2a_0+R\) on the level inputs.
Paragraph~\ref{par:wach-cut-colength} gives the cut colength
on the two-sided Wach height bounds. The simultaneous level
realization and the two polynomial actions are those of
Paragraphs~\ref{par:wach-level-models}--\ref{par:wach-polynomial-actions}.

Multiply the finite source class for arc \(c\) by
\(u_c^{-1}\ell^{2\kappa-k_c}\), and sum over the cuts. Its image
is \(\ell^{2\sigma+2a_0+R+2\kappa}\) times the sum of the target
arc classes. Together with the initial decomposition this gives
\[
 e=2R+4\sigma+2a_0+2\kappa.
\]
The sum over the cuts costs no
additional power of \(\ell\).

\proofstep{Include the torsion generators and then take the normalized colength.}
The free part of the arc calculation bounds the failure of lifting along the orbit. Torsion in the generators must be included before the resulting length can enter the inert-prime table.


Let \(V_2=T_2\otimes_O\operatorname{Frac}(O)\) and
\(r_0=\dim V_2^{G_F}\). By the Bloch--Kato dimension formula the
free rank of the integral finite group is
\[
 \dim(D_{\mathrm{dR}}(V_2)/\mathrm{Fil}^0)+r_0\le dr+r_0.
\]

For a uniformizer \(\varpi\) of \(O\), the long exact sequence
for multiplication by \(\varpi\) gives
\[
 0\longrightarrow T_2^{G_F}/\varpi T_2^{G_F}
   \longrightarrow (T_2/\varpi T_2)^{G_F}
   \longrightarrow H^1(F,T_2)[\varpi]\longrightarrow0.
\]
The first term has dimension \(r_0\) and the middle term has
dimension at most \(r\). Thus at most \(r-r_0\) torsion generators
are needed. The rank of the free quotient and the number of torsion generators together cost at most
\[
 (dr+r_0)+(r-r_0)=(d+1)r.
\]
Since its quotient by the image of
the finite group of \(T_1\) is killed by \(\ell^e\), its length
divided by \(\operatorname{length}_O(O/\ell O)\) is at most
\((d+1)re\). This is the normalization used for the actual row
cokernel. It handles all coefficient ramification degrees.

Finally, for \(h\ge1\) use
\[
 \log_\ell h\le h/D_0+j+1,
\]
by splitting the logarithm at
\(D_0\) and using \(\log_\ell t\le t\) for \(t>0\).
When the word has a position with larger jumps, the bracket in the bound for
\(e\) is at most \(2H/D_0+2(G_P+1)(j+1)\). Otherwise \(h=g\le G_P\)
and it is at most \(3(G_P+1)(j+1)\).

In both cases the resulting
normalized colength is at most \(3C_2 U_B\), with \(U_B\) as in
Section~\ref{sec:cut-length-rows}. Choose a rational table entry \(u_B\) satisfying
\[
 3C_2U_B<u_B<(3C_2+1)U_B.
\]
The lower bound makes \(u_B\) valid for the cokernel at every good character. Such an entry exists since \(U_B>0\).
This gives case B of the ordinary congruence comparison on the
tied local data, with one constant and cutoff for every place,
coefficient field and marked character in the window.

Use the single word chosen for the large-orbit estimate.

Sum the arc bounds and convert the result to the rational entry in its marked row.
\label{cor:inert-long-orbit-1}

The filtration statement for the Wach module
of a crystalline representation over an unramified field is
Berger's filtered comparison \cite[Corollary III.4.5]{berger-wach}.
The exponent in \eqref{eq:long-orbit-exponent} belongs to the
submission's quantitative lifting argument.
\end{proof}


\subsubsection{Alternating sequences with larger jumps and the finite-class length}
\label{sec:inert-thick-word}
\hypertarget{source-inert-thick-word}{}

For an alternating sequence whose larger jumps dominate, the finite-class estimate needs information about the image of each successive lattice map. We follow a lifted class through the sequence and control the restrictions on its coefficients at the even positions.

\begin{established}[Marked alternating word]\label{hyp:alternating-thick-word}
Let \(F/\Q_\ell\) be unramified of even degree \(n\), with
\(\ell\ge11\). Let \(f:T_1\to T_2\) be the marked crystalline surjection of
free stable lattices over the coefficient ring \(O\).
Take its field \(K_O\) to contain the embeddings of \(F\) and
the primitive \(\ell\)-th roots. The descent
of the finite-class bound to the original coefficients uses the same
faithfully flat base change as in Section~\ref{sec:inert-long-orbit}.
Write \(r=\rk_O T_2\). Index the embeddings cyclically and absorb the
parity shift in the choice of origin. If \(J_i\) denotes the source jumps at letter \(i\), their boxes are
\[
 J_i\subseteq[-\eta,\eta]\cup\begin{cases}[-b_i-S,-b_i],&i\text{ even},\\ {}[b_i,b_i+S],&i\text{ odd}.\end{cases}
\]
At most \(m\) target jumps, counted
with multiplicity, lie in the central interval at each embedding.
The target jumps are among the source jumps by strictness of the
crystalline filtration.

Put
\[
\begin{gathered}
 h=\max_i b_i+S,\quad H_{\rm loc}=h+\eta,\quad
 L_0=1+\log_\ell(2H_{\rm loc}),\\
 A_0=\left\lceil\frac{(2\ell-1)H_{\rm loc}}{(\ell-1)^2}
                +3+L_0\right\rceil .
\end{gathered}
\]
In the retained alternating case one has
\(b_i\ge h/2\) and
\(b_i\ge 2(A_0^{\rm real})+\eta+1\), where \(A_0^{\rm real}\) is
the expression inside the ceiling. Define
\[
 c_D=1+\frac{H_{\rm loc}+1}{\ell-1},\qquad
 C_D=rS+2mc_D .
\]
For an \(O\)-module of finite length, we use its length divided by
\(k_O=\operatorname{length}_O(O/\ell O)\).
\end{established}
\begin{proof}
At an inert prime the two embeddings of the distinguished quadratic field are exchanged by Frobenius. On an orbit carrying larger jumps at every position, Lemma~\ref{lem:marked-signed-jumps} therefore gives the displayed signs after a cyclic reindexing. The central multiplicity is the number of central columns in the marked rank-\(r\) weights. The target inherits the source support by strict exactness. On the retained word
\(b_0\le b_i\le b_0+\eta\) with \(b_0=b+O(t)\), whereas
\(S+\eta+\log H=o(b_0)\) uniformly in the host. Thus
\(h\le b_0+o(b_0)\), and the coefficient
\(2(2\ell-1)/(\ell-1)^2\le 42/100\) for \(\ell\ge11\)
gives the second inequality beyond a common threshold.
In particular \(b_i>\eta\).
\end{proof}

\paragraph{Passing from a chain to a length.}
Set \(M_2=H^1_f(F,T_2)\) and \(I=f_*H^1_f(F,T_1)\).
Choose a free \(O\)-module \(A\) of rank at most \((n+1)r\)
with a surjection \(A\to M_2\), using the generator argument in
Section~\ref{sec:inert-long-orbit}. Representatives for a basis
give a linear map to the finite cocycles themselves.
\begin{lemma}[Length from the successive restrictions]
Let \(\ell^{e'}M_2\subset I\). Suppose the final submodule of a
chain in \(A\) maps into \(I\) after multiplication by \(\ell^e\).
Let \(N_g,N_d\) be the numbers of the generator-bounded and length-bounded successive quotients, respectively. If each of the former requires at most \(m\) generators and each of the latter has normalized length at most \(C_D\), the numerical hypotheses
\[
 N_g\le4n,\qquad N_d\le3n/2
\]
give the bound \eqref{eq:thick-chain-mass} on \(\operatorname{len}(M_2/I)\).
\end{lemma}
\begin{proof}
Suppose \(A=A^{(0)}\supset\cdots\supset A^{(k)}\) and
\(\ell^e A^{(k)}\) maps into \(I\).
Suppose also that \(\ell^{e'}M_2\subset I\).
If a successive quotient of the chain is generated by at most
\(m\) elements, its image in the filtration of \(M_2/I\) costs at
most \(me'\) in normalized length. If a successive quotient has
normalized length at most \(C_D\), the corresponding image costs
at most \(C_D\). Indeed these successive images are quotients of
the successive quotients of the chain.

At the last term the image
is killed by \(\ell^e\) and has at most \(\rk_O A\) generators,
since \(A^{(k)}\) is free of rank at most \(\rk_O A\).
Thus \(4n\) steps of the first kind and \(3n/2\) of the second
give normalized length at most
\begin{equation}\label{eq:thick-chain-mass}
           (n+1)r e+4nm e'+\tfrac32 n C_D .
\end{equation}
For later use, a different sufficient condition is
\(\ell^e M_2\subset I+\langle Z\rangle_O\), where \(Z\) is a finite
set of classes in \(M_2\). The normalized length is then at most
\((n+1)r e+|Z|e'\). Take the quotient by the image of
\(\langle Z\rangle_O\) first, and use \(\ell^{e'}M_2\subset I\)
on that image. These arguments apply to the full integral finite
group, including its torsion.

\end{proof}
\paragraph{Wach coordinates and the initial finite classes.}
\label{par:thick-initial-classes}
Use the tame-invariant parameter \(X\) and \(q=X+\ell\) of
Paragraph~\ref{par:wach-cut-correction}, and put \(p_j=\phi^j(q)\).
Write \(N_i(T)\) for the tame-invariant lattice, \(N(T)\) for their
product, and \(P_i^{(j)}(T)=\langle\Phi^jN_i(T)\rangle\) for its
integral Frobenius span at \(i+j\). Let \(a=A_0\),
\(B_i=b_i+S\) and \(H=H_{\rm loc}\). The two signed height bounds
at a destination \(i\) are
\[
 q^{M_i}P_{i-1}^{(1)}(T)\subset N_i(T),\qquad
 q^{H_i}N_i(T)\subset P_{i-1}^{(1)}(T),\qquad
 (M_i,H_i)=
 \begin{cases}(B_i,\eta)&i\text{ even},\\
 (\eta,B_i)&i\text{ odd}.\end{cases}
\]
Use the cyclotomic lattice comparison over the unramified field
and the filtered comparisons at the primitive roots
\cite[Proposition III.2.1 and Corollary III.4.5]{berger-wach}.
\begin{lemma}[Annihilators at the initial letters]
For the given surjection, and for the crystalline extension of
the trivial lattice by the target with its quotient map, let \(C_i\) be the cokernel on the \(i\)-th tame-invariant letter. Then
\[
 q^a C_i=\ell^R C_i=0,\qquad R=\left\lfloor\frac{\ell H}{(\ell-1)^2}+3+\log_\ell(H+1)\right\rfloor.
\]
\end{lemma}
\begin{proof}
The following image bound uses the maximum \emph{letter width}
\(H=\max_i(M_i+H_i)\). Do the calculation before tame invariants
in the ordinary parameter \(\pi\).
The common character with labelled filtration jump \(+M_i\) shifts both lattices by
\[
 [-M_i,H_i]+M_i=[0,M_i+H_i]\subseteq[0,H].
\]

The jumps of the decreasing filtration are now in \([0,H]\),
as required in Berger's positive comparison. Such a character is a product
of powers of the labeled Lubin--Tate characters of the
unramified field, with values in the coefficient field.
\cite[Proposition 2.3]{berger-lubin-tate}
Tensoring by this rank-one lattice preserves the annihilator
of the Wach cokernel. Rational Wach comparison is tensor
compatible. Integrally, multiplication by a generator of the
rank-one Wach lattice identifies the integral \'etale modules,
since the generator is a basis after \'etale scalar extension.
Taking their intersections with the rational Wach modules gives
the tensor identification on integral Wach lattices.
\cite[Proposition III.4.2]{berger-wach}

Apply Berger's positive comparison and the faithfully flat
descent in Paragraph~\ref{par:wach-cut-colength} after the common
twist. For the untwisted Wach cokernel \(C\), the common rank-one twist preserves the annihilator and hence gives
\[
 \pi^H C=0.
\]

The argument of that paragraph which lowers the \(\pi\)-degree
by the divided cyclotomic action gives a scalar exponent at most
\(\lfloor\ell H/(\ell-1)^2\rfloor\). The two input lattices also
occur as the extension of the trivial lattice by the target
and its trivial quotient. The jumps of the extension are those
of the target together with zero, so the same boxes apply.

After tame averaging the annihilators on each letter include
\(X^{\lfloor H/(\ell-1)\rfloor+1}\) and \(\ell^R\), with the following integer exponent
\[
 R=\left\lfloor \frac{\ell H}{(\ell-1)^2}
                 +3+\log_\ell(H+1)\right\rfloor .
\]
Set \(d=\lfloor H/(\ell-1)\rfloor\), so the two annihilators are
\(X^{d+1}\) and \(\ell^R\). In the expansion of \((X+\ell)^a\),
each term with \(X\)-degree \(k\le d\) has \(\ell\)-exponent
\(a-k\ge R\); every other term is divisible by \(X^{d+1}\).
Thus \(q^a\) annihilates the same cokernel.

The bounds are uniform in the number of letters,
the rank and the ramification of the coefficient ring.
\label{cor:inert-thick-initial-1}

\end{proof}

\begin{lemma}[Integral Herr representatives for the basis extensions]
For every basis cocycle \(y\) the two annihilators give vectors \(\mu_i(y),\nu_i(y)\) satisfying \eqref{eq:thick-extension-data} on each letter. On the free span, the choice is extended by
\[
 (\mu_i,\nu_i)\!\left(\sum_j b_j y_j\right)=\sum_j b_j\big(\mu_i(y_j),\nu_i(y_j)\big).
\]
\end{lemma}
\begin{proof}
For each basis cocycle in the free family take its extension
\(0\to T_2\to Y\to O\to0\). Its rational extension is crystalline
by the definition of the finite group. Use exactness of the
integral \((\varphi,\Gamma)\)-module functor and the Wach
intersection with the kernel. Tame averaging, of order
\(\ell-1\), preserves integral lifts.

Choose a tame-invariant
lift \(E\) of the trivial generator in the \'etale module of
\(Y\), and take the Herr pair \(y=(\Phi-1)E\),
\(z=(\gamma-1)E\) in that of \(T_2\). Denote their letter
components by \(y_i,z_i\), and use the convention
\((\gamma-1)y=(\Phi-1)z\).
The two annihilators provide Wach lifts \(E_i,E_i'\) of
\(\ell^R\) and \(q^a\) times the trivial generator in letter \(i\).
Set
\[
 \mu_i=E_i-\ell^RE|_i,\qquad \nu_i=E_i'-q^a E|_i,\qquad
 A_i=\ell^R y_i+\Phi\mu_{i-1}-\mu_i .
\]
Both differences defining \(\mu_i,\nu_i\) belong to the target
\'etale module. The Wach intersection and the first signed height
bound give
\begin{equation}\label{eq:thick-extension-data}
\begin{gathered}
 \ell^R z_i+(\gamma-1)\mu_i\in XN_i(T_2),\qquad
 q^{M_i} A_i\in N_i(T_2),\\
 \ell^R q^{H_i} A_i\in P_{i-1}^{(1)}(T_2),\qquad
 \ell^R\nu_i-q^a\mu_i\in N_i(T_2).
\end{gathered}
\end{equation}

For the third inclusion, \(A_i=\Phi E_{i-1}-E_i\) in the
extension, so \(q^{H_i}A_i\) belongs to its Frobenius span.
Expand in a free basis of \(O[[X]]\) over \(\phi(O[[X]])\),
as in the proof of the arc estimate. Each quotient coefficient
of the resulting Wach vectors vanishes by independence.
Thus the Frobenius span meets the target in the target span.
This even gives the inclusion before multiplication by \(\ell^R\).
The data in \eqref{eq:thick-extension-data} now extend linearly in the free family,
even though the individual extensions \(Y\) were chosen only
on its basis.

\end{proof}
\paragraph{First restriction at an even letter.}
\begin{lemma}[The first restriction on the even letters]
At each even letter one can impose a single kernel condition \(z\in\ker\pi_i\) for which
\[
 \operatorname{im}\pi_i\text{ is generated by at most }m\text{ elements},\qquad z\in\ker\pi_i\ \Longrightarrow\ u_i(z)\in N_i(T_2),
\]
where \(u_i(z)\) is the corrected vector constructed in the proof. If \(m=0\) the assertion holds on
the full free family.
\end{lemma}
\begin{proof}
For \(x\in q^{-a}N_i(T_2)\), let
\(\partial x=X^{-1}(\gamma-1)x\) in the \'etale module.
It preserves \(q^{-a}N_i(T_2)\) and \(N_i(T_2)\).
Indeed the action on the invariant Wach lattice modulo \(X\)
is trivial, and \(\gamma(q)/q\) is a unit congruent to one
modulo \(X\). At a primitive root, put
\[
             c_j=\frac{1-\chi(\gamma)^{-j}}{\ell}.
\]
On the associated Laurent coefficient of degree \(-j\) the
divided operator acts by \(c_j\) since \(X=-\ell\) there.
For a Hodge jump \(u\), the quotient of the allowed pole lattice
by the unshifted lattice has degrees \(-u-k\), \(1\le k\le a\).

To see why one factor per distinct shift suffices, choose a basis
of the crystalline space adapted to the filtration. In each
direction the Laurent action is triangular, with distinct
diagonal entries indexed by the degrees in that direction.
The action of \(1/X(t)\) is scalar and the cyclotomic action
in the crystalline basis is substitution
\[
 t\mapsto\chi(\gamma)t.
\]

Thus the product vanishes separately on each direction even
when a shift repeats in another direction. This is a computation
over the coefficient field. The Hodge criterion at the cyclotomic
roots and saturation for \(\ell\) in \(q^{-a}N_i(T_2)/N_i(T_2)\)
give the same inclusion integrally.

Let \(J_i\) be the nonzero shifts that occur. They belong to
\([1-B_i,a-b_i]\cup[1-\eta,a+\eta]\).
By the valuation of powers of a principal unit,
\[
 v_\ell(c_j)=v_\ell(j)
\]
for \(j\ne0\). In the first
interval there are at most \((a+S)/\ell^k+1\) multiples of
\(\ell^k\), and no such nonzero multiples for
\(k>\log_\ell(2H)\). For the central interval split at zero
and apply the factorial bound on both sides.

Taking the maximum
over even letters, choose \(s_1\) and polynomials \(R_i(U)\)
proportional over \(O\) to \(\prod_{j\in J_i}(U-c_j)\) with
\[
 R_i(0)=\ell^{s_1},\qquad
 s_1\le(2a+2\eta+S)/(\ell-1)+L_0,\qquad
 \partial R_i(\partial)q^{-a}N_i(T_2)\subset N_i(T_2).
\]
Including \(\partial\) kills degree zero if that shift occurs.
The scaling of each polynomial is integral by the choice of the
maximum. Write \(R_i(U)=\ell^{s_1}+U R_i^*(U)\), and set
\[
 m_i=R_i(\partial)(q^{-a}\nu_i)+R_i^*(\partial)(z_i/X),
 \qquad u_i=\ell^R m_i-\ell^{s_1}\mu_i .
\]
Put \(v_i=\ell^R q^{-a}\nu_i-\mu_i\) and
\(d_i=\ell^R z_i/X+\partial\mu_i\).

By \eqref{eq:thick-extension-data},
\(q^a v_i,d_i\in N_i(T_2)\). The polynomial identity gives
\(u_i=R_i(\partial)v_i+R_i^*(\partial)d_i\). Thus
\(q^a u_i\in N_i(T_2)\) and \((\gamma-1)u_i\in XN_i(T_2)\).
Also
\[
 \ell^R\bigl(\ell^{s_1}z_i+(\gamma-1)m_i\bigr)
       =X\ell^{s_1}d_i+(\gamma-1)u_i .
\]
The bracket belongs to \(XN_i(T_2)\) by cancellation
of \(\ell^R\) in the integral \'etale module, where the Wach
lattice and its \(X\)-multiple are \(\ell\)-saturated. In coordinates this follows from the
description of the \'etale module as the completed localization
of the lattice, with injective reduction of the lattice modulo
the coefficient uniformizer.

At a primitive root, the only Laurent coefficient of \(u_i\)
that can obstruct membership in the Wach lattice is degree zero.
Writing \([t^j]\) for coefficient extraction at the root gives
\[
 [t^j](\gamma-1)u_i=(\chi(\gamma)^j-1)[t^j]u_i\qquad(j\ne0).
\]
The scalar on the right is nonzero. The degree-zero obstruction is therefore
the class in \(\mathrm{Fil}^{-a}D_i/\mathrm{Fil}^{0}D_i\).
All jumps in \([-a,0)\) are central since \(a<b_i\).
The quotient has dimension at most \(m\). Each obstruction
map is linear in the free family and its image is an
\(O\)-module generated by at most \(m\) elements.
Successively take these kernels for even \(i\). This leaves
a family on which every \(u_i\) belongs to \(N_i(T_2)\), using
one rank restriction per even letter. When \(m=0\) the entire
family survives.
\end{proof}

\paragraph{Cyclotomic correction at fixed precision.}
\label{par:thick-integral-correction}
Put
\[
 M=1+\lfloor h/(\ell-1)\rfloor,\qquad
 \sigma_1=\left\lceil 1+(M+\eta)/(\ell-1)\right\rceil .
\]
For a vector \(v\in XN_i(T)\) there is an integral correction
\(\tau(v)\in P_i^{(1)}(T)\) well-defined modulo
\(X^Mq^\eta P_i^{(1)}(T)\), with
\[
  \ell^{\sigma_1}\Phi v-(\gamma-1)\tau(v)
             \in X^Mq^\eta P_i^{(1)}(T).
\]
We use it at an even \(i\), where the odd destination gives
\(q^\eta P_i^{(1)}(T)\subset N_{i+1}(T)\).
There is also a two-step correction with precision \(Xq^{2H}\)
on \(P_i^{(2)}(T)\) and exponent
\(\sigma_*=\lceil1+2H/(\ell-1)^2\rceil\).

Here are the divisions that give these two claims. In either
case solve \((\gamma-1)d=v\) modulo a large power \(X^K\),
over the coefficient field. The degree-by-degree solution
of Paragraph~\ref{par:wach-cut-correction} puts the term in
degree \(y\) in \(\ell^{-e_y}X^yN_i(T)\), where
\[
 e_y=\lceil\ell y/(\ell-1)\rceil.
\]
Take \(d\) with zero
constant coefficient. Write \(J\) for the exponent of \(q\)
in the precision ideal and \(a_y=(\ell-1)y\). In one step
\(\phi(X)^y\) is a unit times \(X^yq^{a_y}\).

If \(a_y<J=\eta\), then \(e_y\le\sigma_1\), since
\(J\le (\ell-1)M\). If \(a_y\ge J\) and \(y<M\),
expand \(q^{a_y-J}=(X+\ell)^{a_y-J}\).
The terms which survive modulo \(X^Mq^J\) are divisible by
\(\ell^{\max(0,\ell y-J-M+1)}\). For \(e_y>\sigma_1\)
we have \(\ell y>J+M+\ell-1\); hence
\[
 \ell y-J-M+1\ge e_y-\sigma_1.
\]

For \(a_y\ge J\) and \(y\ge M\) the entire term vanishes.

For two steps use \(M=1\), \(J=2H\), and
\(\phi^2(X)^y\) associated to
\(X^yq^{a_y}p_1^{a_y}\).
If \(a_y\ge J\) the terms with \(y>0\) vanish. Otherwise
\(p_1\in(\ell,q^\ell)\), by reduction modulo \(\ell\).
After factoring out \(X^yq^{a_y}\), the terms of the binomial
expansion that can survive have \(q^\ell\) used at most
\(\lfloor(J-a_y-1)/\ell\rfloor\) times. If
\(e_y>\sigma_*\), then
\(a_y>J/\ell+(\ell-1)^2/\ell\), and
\[
 a_y-\frac{J-a_y-1}{\ell}
      \ge \frac{\ell a_y-J}{(\ell-1)^2}
      \ge e_y-\sigma_* .
\]
The last inequality uses the defining roundings. Thus the
required division is integral. Here one may take \(K>\max(M,J)\), since
\(\phi^j(X)^K\) then belongs to the precision ideal in both
cases.

The integral coefficient quotient by \(X^Mq^J\) is free
over \(O\) by monic Weierstrass division. For either free Frobenius span \(P\) and its precision submodule \(P_{\rm prec}\), saturation gives
\[
 (P_{\rm prec}[1/\ell])\cap P=P_{\rm prec}.
\]

The resulting representative of \(\ell^{\sigma_1}\Phi d\)
or \(\ell^{\sigma_*}\Phi^2d\) satisfies the claimed
cyclotomic congruence, by commutation and saturation.
Uniqueness of \(d\) modulo \(X^K\) over the coefficient field
follows from the nonzero diagonal \(\chi(\gamma)^{(\ell-1)y}-1\)
in positive degrees. Hence the corrections are linear modulo
the precision and natural in maps of crystalline lattices.
They may be made linear on a free family by choosing
representatives on its basis. In one step the correction to a
vector of \(X^MN_i(T)\) can be taken zero. The calculation of
\(d\) in that case starts at degree \(M\), and
\[
 (\ell-1)M\ge\eta.
\]
For \(v=(\gamma-1)w\) with
\(w\in XN_i(T)\) one can use \(d=w\).

\paragraph{Second restriction and the alternating sum.}
Write \(N_i=N_i(T_2)\) for this construction. For even \(c\)
set
\[
 \xi_c=\ell^{s_1}z_c+(\gamma-1)m_c\in XN_c.
\]

Choose \(t_c=\tau(\xi_c)\) linearly on the restricted family.
At the odd letter \(c+1\) put
\[
 L_c=\ell^{\sigma_1}\Phi m_c+
             \ell^{s_1+\sigma_1}y_{c+1}-t_c,\qquad
 \Xi_c=\ell^{s_1+\sigma_1}z_{c+1}+(\gamma-1)L_c .
\]
The Herr relation gives
\(\Xi_c=\ell^{\sigma_1}\Phi\xi_c-(\gamma-1)t_c\).

Consequently \(\Xi_c\in X^MN_{c+1}\). It lies more precisely
in \(X^Mq^\eta P_c^{(1)}(T_2)\).
To make the next difference integral, write
\[
 U_{c+1}=\ell^R L_c-\ell^{s_1+\sigma_1}\mu_{c+1}
        =\ell^{\sigma_1}\Phi u_c+\ell^{s_1+\sigma_1}A_{c+1}
           -\ell^R t_c .
\]
The signed height and \eqref{eq:thick-extension-data} show that \(q^\eta U_{c+1}\in N_{c+1}\).
Applying \(\gamma-1\) gives \(\ell^R\Xi_c\) minus
\(\ell^{s_1+\sigma_1}(\ell^R z_{c+1}+(\gamma-1)\mu_{c+1})\).
This is in \(N_{c+1}\). Its only possible Laurent obstruction
is again degree zero, this time in
\(\mathrm{Fil}^{-\eta}D_{c+1}/\mathrm{Fil}^0D_{c+1}\).
Only central jumps contribute, so taking one kernel of rank
at most \(m\) per odd destination makes each \(U_{c+1}\) integral.
There are at most \(n\) restrictions including the first ones.

On the resulting family define at each even \(c\)
\[
 w_c=\Phi L_{c-2}+\ell^{s_1+\sigma_1}y_c-\ell^{\sigma_1}m_c,
 \qquad Y_c=\Phi w_c+t_c .
\]
All indices are cyclic, including when the word has length two.
There are three useful tests on \(w_c\). First,
\[
 \ell^R w_c=\Phi U_{c-1}-\ell^{\sigma_1}u_c
                        +\ell^{s_1+\sigma_1}A_c .
\]
The height bounds place the right side, after multiplication
by \(q^{B_c}\), in \(N_c\). Multiplication instead by \(q^\eta\)
places it in \(P_{c-1}^{(1)}(T_2)\) after a possible
additional factor \(\ell^R\), by \eqref{eq:thick-extension-data}.
Both lattices are \(\ell\)-saturated inside the integral
\'etale module.

For the Frobenius span, reduce in a Wach basis
and extend scalars by the finite free Frobenius coefficient map.
The original basis stays independent over the residue Laurent
field because the linearized Frobenius on the \'etale module
is invertible. Cancel these powers of \(\ell\).
Second, \((\gamma-1)w_c=\Phi\Xi_{c-2}-\ell^{\sigma_1}\xi_c\)
belongs to \(XN_c\), since
\(\Phi(X^MN)\subset q X^MN\). Third,
\[
                (\gamma-1)Y_c=\Phi^2\Xi_{c-2}-\Xi_c
                                 \in X^MN_{c+1}.
\]
These calculations give exactly the pole, first pullback
and cyclotomic tests on \(w_c\) used by the even-letter estimates.

Define \(Z_c=(\Phi-1)^{-1}(\gamma-1)Y_c\) in \(X^MN\),
where the inverse exists by contraction on the product of all
letters. Then the sum of \((Y_c,Z_c)\) over even \(c\)
represents \(\ell^{s_1+\sigma_1}(y,z)\). Indeed let
\[
  B=\sum_{c\text{ even}}\bigl(L_c+\Phi L_c
                                  +\ell^{s_1+\sigma_1}y_{c+2}\bigr).
\]
Summing first components gives
\(\sum_cY_c=\ell^{s_1+\sigma_1}y+(\Phi-1)B\)
by invariance of the even labels under the shift by two.
For the second components
\(\ell^{s_1+\sigma_1}z+(\gamma-1)B
 =\sum_c(\Xi_c+\Phi\Xi_c)\in X^MN\).
The Herr relation and uniqueness of the inverse prove the same
coboundary identity for the whole pair.
\label{cor:inert-thick-herr-1}

\paragraph{Integral comparison with the lifted vectors.}
For an admissible \(w\) at an even \(c\), use the first
correction to form
\[
             V(w)=\ell^{\sigma_1}\Phi w
                          -\tau((\gamma-1)w).
\]
Its cyclotomic defect lies in \(X^MN\). The second component is uniquely determined at this precision by
\[
 z=(\Phi-1)^{-1}(\gamma-1)V(w),\qquad z\in X^MN.
\]

Its first component has denominators in finitely many
cyclotomic factors, by the height bounds, and the second
belongs to \(X^MN\). The crystalline criterion proved in
Section~\ref{sec:inert-long-orbit} therefore makes its
integral Herr class finite. A change of representative of
the correction changes both components by a coboundary.
Indeed if the first component of a pair is in \(X^MN\)
along with its second component, the first is the
\((\Phi-1)\)-image of a vector in \(X^MN\). Commutation
and injectivity on \(X^MN\) give the other component of
the coboundary. This also shows that the resulting class
depends linearly on \(w\) and is natural in the marked map.
Writing \([V(w)]\) for the resulting class, the last property of the correction gives
\[
 w\in XN_c\quad\Longrightarrow\quad[V(w)]=0.
\]
Thus integral congruence modulo \(XN_c\) is
sufficient for comparing classes.

For \(w=w_c\), the defect at letter \(c\) is
\(\Phi\Xi_{c-2}-\ell^{\sigma_1}\xi_c\).
Its first summand lies in \(X^MN_c\), on which the
correction is zero. Its second summand has correction
\(-\ell^{\sigma_1}t_c\) modulo the precision.
Thus
\[
 V(w_c)\equiv\ell^{\sigma_1}Y_c\pmod{X^MN}.
\]

The resulting finite classes, denoted by \(\kappa_c\), sum
to \(\ell^{s_1+2\sigma_1}\) times the original class.
The displayed cocycle uses \(\Phi-1\) and \(\gamma-1\)
in this order and has the displayed plus sign. If the
opposite convention for the extension class is used, replace
the sum of the source lifts by its negative.

In particular, if
\[
 \ell^e w_c-f(u_c')\in XN_c(T_2)
\]
for an admissible vector \(u_c'\) in the source, the
finite class attached to \(u_c'\) maps to
\(\ell^e\kappa_c\). All these comparisons take place on
the original rank block and the same map; they do not
require exactness of the Wach functor on the right or a
choice of integral splitting of the surjection.
\label{cor:inert-thick-herr-2}

\paragraph{Filtered comparison at a positive letter.}
\label{par:thick-filtered-comparison}
Write \(i\) for an odd letter, \(b=b_i\) and \(A=B_{i-1}\).
Suppose \(y\in N_{i-2}(T_2)\) has
\(z=\Phi^2\operatorname{sp}_{i-2}(y)\in\mathrm{Fil}^bD_i\).
\begin{proposition}[Integral preimage of the positive filtered class]
In the two-step height and jump range at \(i-2,i-1,i\), put
\[
 E'=\left\lfloor\frac{3a+S+2\eta-1}{\ell-1}
                         +\log_\ell H\right\rfloor .
\]
There exists a vector of the source Wach lattice at \(i\) whose
specialization maps to \(\ell^{E'+\sigma_*+A-b}z\).
The exponent \(E'\le E_G\) is independent of \(y\).
\end{proposition}
\begin{proof}

Use the two-step correction
of Paragraph~\ref{par:thick-integral-correction} to form
\[
 \lambda=\ell^{\sigma_*}\Phi^2 y-
                      \tau((\gamma-1)y)\in P_{i-2}^{(2)}(T_2).
\]
It has cyclotomic defect in \(Xq^{2H}P_{i-2}^{(2)}(T_2)\).
At a primitive root the span is regular in \(t\) and the
constant of \(\lambda\) is \(\ell^{\sigma_*}z\).
The correction has zero constant there because each transported
positive \(X\)-power vanishes at the root and its division modulo
\(Xq^{2H}\) preserves this. By the cyclotomic defect all
positive coefficients below \(2H\) vanish. The correction
likewise has zero specialization at \(X=0\).
These conclusions hold for the same integral representative
at every primitive root.

The signed height at \(i-1\) puts \(p_1^A\lambda\) in
\(P_{i-1}^{(1)}(T_2)\). Set \(x=q^{-b}p_1^A\lambda\).
Then
\[
 q^{b+\eta}x\in N_i(T_2)
\]
by the height at \(i\).
At a primitive root, the coefficients of \(x\) of degree
\(k<-b\) vanish. In degrees \(-b\le k<2H-b\) they lie
in the span of \(z\), hence in \(\mathrm{Fil}^{-k}D_i\).
The filtration is full for \(k\ge 2H-b\ge\eta\), since
all jumps at the odd letter are at least \(-\eta\).
The Hodge lattice test therefore gives \(x\in N_i(T_2)\).
The same signed height at \(i\) applied to the defect
gives
\[
           p_1^A(\gamma-1)\lambda\in Xq^{a+b}N_i(T_2),
\]
because \(2H\ge a+b+\eta\).
In particular
\(\operatorname{sp}_i(x)=\ell^{\sigma_*+A-b}z\).

It remains to lift a controlled multiple of \(x\) in the
integral Wach lattice. Let \(V_q=\gamma(q)/q\),
\(V_1=\gamma(p_1)/p_1=\phi(V_q)\).
Then \(V_q\) is one modulo \(X\) and \(\chi(\gamma)\) modulo \(q\).
The last equality follows by expanding at a primitive root
in the parameter \(t\) of degree one. Also
\(V_1\equiv1\pmod{Xq}\) by the formula for \(\phi(X)\).
The preceding defect equation implies
\[
 \gamma x=V_q^{-b}V_1^A x
                    \pmod{Xq^aN_i(T_2)}.
\]
Thus \(\partial x=d x\bmod q^aN_i(T_2)\) with
\(d=(V_q^{-b}V_1^A-1)/X\) integral and \(d\equiv c_b\bmod q\).

For \(g\in O[[X]]\), \(\gamma(g)\equiv g\bmod q\), and
the divided operator preserves \(q^aN_i(T_2)\).
Applying it to \(q^k g x\) shows that
\[
    (\partial-c_{b-k})q^k g x
           \in q^{k+1}O[[X]]x+q^aN_i(T_2)\qquad(0\le k<a).
\]
Indeed the multiplier on the surviving degree is
\(\chi(\gamma)^k c_b+(\chi(\gamma)^k-1)/(-\ell)=c_{b-k}\).
It follows that \(P_-(\partial)x\in q^aN_i(T_2)\), where
\(P_-(U)=\prod_{j=b-a+1}^{b}(U-c_j)\).
The image input of Paragraph~\ref{par:thick-initial-classes}
gives an integral source preimage of this vector.

For a second filter, first lift \(q^a x\) integrally and divide
the lift by \(q^a\) in the source \'etale module.
The resulting vector \(v\in q^{-a}N_i(T_1)\) maps to \(x\).
The shifted-jump calculation of the first restriction, now
on the odd source, shows that \(P_+(\partial)v\in N_i(T_1)\),
where
\[
 P_+(U)=\prod_{j\in[b+1,b+S+a]\,\cup\,[1-\eta,\eta+a]}(U-c_j).
\]
The two sets of indices for \(P_-\) and \(P_+\) are disjoint
because \(b\ge 2a+\eta\). This integer inequality follows
from \(b\ge 2A_0^{\rm real}+\eta+1\) and
\(a<A_0^{\rm real}+1\), with equality of \(a\) and the
real bound allowed when the latter is integral.

For a root index \(j\) of \(P_-\), the valuation of
\(P_+(c_j)P_-'(c_j)\) is the sum of \(v_\ell(j-j')\)
over the other indices. Indeed, if \(I_-\) and \(I_+\) denote the two index sets, factoring
the polynomials at the root gives
\[
 v_\ell(P_+(c_j)P_-'(c_j))
 =\sum_{j'\in(I_-\cup I_+)\setminus\{j\}} v_\ell(c_j-c_{j'})
 =\sum_{j'\in(I_-\cup I_+)\setminus\{j\}} v_\ell(j-j') .
\]
The indices in or above the first interval give factorials
with total weight \(2a+S-1\). In the central interval the
differences form an interval of length \(a+2\eta\) with
largest absolute value at most \(H\).
Thus the valuation is at most
\[
           \frac{3a+S+2\eta-1}{\ell-1}+\log_\ell H .
\]
Let \(E'\) be the floor of the displayed bound.

For the index set \(J_-\) of \(P_-\) take
\[
 B(U)=\sum_{j\in J_-}
 \frac{\ell^{E'}}{P_+(c_j)P_-'(c_j)}\,\frac{P_-(U)}{U-c_j}.
\]
The valuation bound gives \(B\in O[U]\). Since \(P_-\) is monic
and the numerator vanishes at every \(c_j\), there is \(C\in O[U]\)
with \(BP_++CP_-=\ell^{E'}\). Choose integral source preimages
\(v_+,v_-\) of \(P_+(\partial)x,P_-(\partial)x\), respectively,
as above. Stability under \(\partial\) then gives an integral vector
\(B(\partial)v_++C(\partial)v_-\). By commutation with \(f\), its
image is \(\ell^{E'}x\) and its specialization maps to
\(\ell^{E'+\sigma_*+A-b}z\).

The exponent \(E'\le E_G\) is independent of \(y\)
and of any choice of an integral splitting. This is the
filtered estimate used at each odd letter in the two
even-letter comparisons.
\label{cor:inert-thick-filtered-1}
\end{proof}


\paragraph{The zeroth coefficient at an even letter.}\label{par:thick-constant-data}
Use \(N_i\) for the target Wach lattice at \(i\), in the tame-invariant
coordinate ring. Put \(B_i=b_i+S\) at each letter.
Use the associate \(q=X+\ell\) of the first cyclotomic factor
and \(t\) for the logarithmic parameter at a primitive \(\ell\)-th root.
The parameter \(X\) of the
tame-invariant Wach ring is a unit in the completed de Rham ring at
that root. Write \(D_i\) for the \(i\)-th crystalline space and
\(\operatorname{sp}_i:N_i\to D_i\) for the reduction at \(X=0\).

Let \(\Phi\) also denote the induced Frobenius \(D_i\to D_{i+1}\).
The one-step Wach height at \(c-2\) gives
\[
\begin{aligned}
 P_{c-2}&=XN_{c-2}+q^{B_{c-2}}\langle\Phi N_{c-3}\rangle
                 \ \subset N_{c-2},\\
 \operatorname{sp}_{c-2}(P_{c-2})
       &=\ell^{B_{c-2}}\Phi\,\operatorname{sp}_{c-3}(N_{c-3}),
\end{aligned}
\]
Here the bracket denotes the span over the tame-invariant ring.
The equality follows by reduction at \(X=0\), since \(q(0)=\ell\)
and all scalar constants in \(O\) occur as reductions.
Take
\[
  C=\mathrm{Fil}^{b_{c-1}}D_{c-1},\qquad
  L_\eta=\mathrm{Fil}^{-\eta}D_c,\qquad L_0^{\rm fil}=\mathrm{Fil}^0 D_c .
\]
Both \(D_{c-1}/C\) and \(L_\eta/L_0^{\rm fil}\) have dimension at most
\(m\). The first quotient keeps only jumps in the central interval at
the odd letter. The second keeps only the negative central jumps at
the even letter.

We give the coefficient comparison with its quantitative inputs.
On a submodule \(A'\) of the free family of cocycles, suppose that each
\(w(a)\) is a tame-invariant vector at \(c\), linear in \(a\), with
\((\gamma-1)w(a)\in XN_c\) and
\[
 q^{\eta+B_c}w(a)\in N_c.
\]

Its normalized Laurent coefficients \(d_j(a)\) are obtained by
expanding at one primitive \(\ell\)-th root after dividing by the
expansion of \(q^{\eta+B_c}\), which is \(t^{\eta+B_c}\) times a
unit. They are linear in \(a\). The action of the principal generator
fixes the root and on the \(j\)-th logarithmic coefficient acts by
\(\chi(\gamma)^j\). For \(j\ne 0\), the coefficient of
\((\gamma-1)w(a)\) is in \(\mathrm{Fil}^{-j}D_c\) and
\(\chi(\gamma)^j-1\ne0\). Hence
\[
 d_j(a)\in\mathrm{Fil}^{-j}D_c
\]
for \(j\ne0\).

\paragraph{Transport to the first cyclotomic level.}
Use either Wach module and choose an integral vector \(z_1\)
at \(c-1\). For a positive precision \(D\), put
\(J=\lceil D/(\ell-1)\rceil\) and
\[
 \sigma_D=\lceil 1+(D+1)/(\ell-1)\rceil.
\]

On a Wach basis at \(c-1\), take a polynomial vector \(v\)
of degree \(<J\) with constant coefficient equal to that of
\(z_1\) and with \((\gamma-1)v\) in \(X^J\) times the Wach
module after inverting \(\ell\).
The degree-by-degree construction of
Paragraph~\ref{par:wach-cut-correction} bounds the denominator
in degree \(a\) by \(\ell^{A_a}\), \(A_a=a+v_\ell(a!)\).

Consider \(\ell^{\sigma_D}\Phi v\)
in the Frobenius span at \(c\) modulo \(Xq^D\).
In degree \(a\ge1\) its scalar before the coefficient denominators
is a unit times \(\ell^{\sigma_D}X^a q^{(\ell-1)a}\).
In the expansion \(X^a=X(q-\ell)^{a-1}\), the terms retained
modulo \(Xq^D\) have \(\ell\)-content at least
\(\max(0,\ell a-D)\). Since \(A_a\le\ell a/(\ell-1)\),
\[
 \sigma_D+\max(0,\ell a-D)\ \ge\ A_a .
\]
For \(\ell a\le D\) this follows from
\(\sigma_D\ge D/(\ell-1)\). For \(\ell a\ge D\) it
follows because \(\ell a/(\ell-1)-\ell a+D\) decreases
as \(a\) increases. Thus the transported vector has an
integral representative \(y\) in the Frobenius span modulo
\(Xq^D\) after inverting \(\ell\).

Since \(\phi(X)^J\) is divisible by \(Xq^D\), the representative
has \((\gamma-1)y\) in \(Xq^D\) times this span after inverting
\(\ell\). For the free Frobenius span \(P\), the quotient by the distinguished factor \(Xq^D\) is \(O\)-torsion-free. Consequently
\[
 (Xq^D P[1/\ell])\cap P=Xq^D P.
\]
This returns the rational divisibility to the integral span.
The construction is in the tame-invariant span. The coordinates
of \(y\) at \(X=0\) on the Frobenius basis are
\(\ell^{\sigma_D}\) times the coordinates of \(z_1\) at \(X=0\)
on the preceding basis.

At each primitive \(\ell\)-th root the
Frobenius span evaluates in the crystalline space with
coefficients in \(K_O[[t]]\). Its constant term is obtained
by \(\Phi\) from specialization at the preceding letter.
The \(\gamma-1\) congruence gives zero coefficients for
\(0<j<D\), because
\(\chi(\gamma)^j-1\ne0\) at each such degree.
Thus through degree \(D-1\) the expansion consists of the one
constant \(\ell^{\sigma_D}\Phi\operatorname{sp}_{c-1}(z_1)\).

The filtered crystalline comparison at the intervening odd letter
provides a number \(E'\le E_G\), where
\[
\sigma_2=1+\frac{2H_{\rm loc}}{(\ell-1)^2},\qquad
E_G=\frac{3A_0+S+2\eta}{\ell-1}+L_0 .
\]
For \(z\in \Phi^2\operatorname{sp}_{c-3}(N_{c-3})\cap C\) there
is \(z_1\in N_{1,c-1}\) whose specialization maps to \(\ell^k z\),
where \(k=\lceil E'+\sigma_2\rceil+B_{c-2}-b_{c-1}\).
Here \(N_{1,i}\) denotes the source Wach lattice.
Apply the transport in the source with \(D=B_c\).
The signed height gives \(q^{B_c}y\in N_{1,c}\) and
\((\gamma-1)y\in XN_{1,c}\). Its image has constant coefficient
\(\ell^\kappa\Phi z\) and all nonzero coefficients below \(B_c\)
vanish, with
\[
 \kappa\le (B_c+1)/(\ell-1)+3+E_G+\sigma_2+
                   B_{c-2}-b_{c-1},\qquad
 \kappa=\sigma_{B_c}+k.
\]

\paragraph{The lattice at an even letter.}
\begin{lemma}[Covolume after retaining the central directions]
\label{lem:thick-even-covolume}
At an even letter \(i\) of this word, start with
\(P_i=XN_i+q^{B_i}\langle\Phi N_{i-1}\rangle\).
Adjoining at most \(m\) vectors to \(P_i\) leaves normalized
colength at most \(C_D=rS+2mc_D\) in \(N_i\).
In particular
\[
 m=0\quad\Longrightarrow\quad\ell^S N_i\subseteq P_i.
\]
\end{lemma}

\begin{proof}
At an even \(i\), choose a Wach basis \(e_j\) at \(i-1\) whose
Frobenius specializations give an \(O\)-basis of the specialized
Frobenius lattice adapted to the Hodge filtration at \(i\).
Such a basis exists because the intersections of a lattice
with the subspaces of a flag are saturated. Each successive
quotient is free over \(O\). Lift the resulting basis from
the specialization to the Wach module, where it is a basis
by Nakayama's lemma. Write \(-\alpha_j\) for the jump of the
corresponding crystalline vector. Divide the indices according to the signed jumps of the adapted basis:
\[
 J_0=\{j:|\alpha_j|\le\eta\},\qquad |J_0|\le m,\qquad b_i\le\alpha_j\le B_i\quad(j\notin J_0).
\]


Let \(D=B_i+\eta\). For each central direction \(j\), apply the
first-level transport with precision \(D\) to \(e_j\), obtaining
\(y_j\) in the Frobenius span. The vector \(g_j=q^{\alpha_j}y_j\)
belongs to \(N_i\). To see this, at each root its constant
direction starts in degree \(\alpha_j\) since \(q=t\) times
a unit, and its crystalline direction is in
\(\mathrm{Fil}^{-\alpha_j}\). Its error starts in degree
\[
 D+\alpha_j\ge B_i,
\]
where the filtration is full.
The \(q\)-power denominator is removed by the same integral
Weierstrass test as above.

For a Wach basis \((e_j)\) at \(i\), factoring out the common powers from the columns gives
\[
 \det(q^{B_i}\Phi e_1,\ldots,q^{B_i}\Phi e_r)
       =u q^{rB_i-\sum_j\alpha_j},\qquad u\in O[[X]]^\times .
\] The exponent follows at the
cyclotomic level because the Hodge lattice in the adapted
directions has orders \(\alpha_j\).
The two-sided signed height makes this matrix invertible
after inverting \(q\) over \(O[[X]]\), and \(q=X+\ell\)
generates a prime ideal. These facts give the asserted unit
and power over the integral ring. Leave the noncentral
columns as they are and replace each central one by \(g_j\).
Relative to the old columns the replacement is given by the
transport matrix with the central columns multiplied by
\(q^{\alpha_j-B_i}\). At \(X=0\), before the individual
multipliers the transport matrix has determinant
\(\ell^{m'\sigma_D}\),
where \(m'\le m\) is the number of central columns.

If \(C_{\rm old},C_{\rm new}\) denote the specialized
column matrices, then
\[
 \det C_{\rm new}=\det C_{\rm old}\cdot
       \ell^{m'\sigma_D+\sum_{j\ {\rm central}}(\alpha_j-B_i)}.
\]
Its normalized valuation is therefore
\[
\sum_{j\ {\rm noncentral}}(B_i-\alpha_j)+m'\sigma_D
 \le rS+2mc_D.
\]
Indeed \(D\le H_{\rm loc}\) gives
\(\sigma_D\le\lceil c_D\rceil\le2\lfloor c_D\rfloor\).
The noncentral columns belong to \(P_i\), and adding the
at most \(m\) central classes to \(P_i\) thus leaves a quotient
of normalized length at most \(C_D\).
When \(m=0\) the completed Hodge lattice is contained in
\(q^{b_i}\) times the completed Frobenius span. Weierstrass
saturation gives this inclusion integrally. Multiplication
by \(q^S\) then lands in \(P_i\), and
\(q^S\equiv\ell^S\bmod X\) gives the sharper
\(\ell^S N_i\subset P_i\).
\end{proof}

\paragraph{Recovering the target coefficient integrally.}
\begin{lemma}[An integral representative of the constant coefficient]
Use Lemma~\ref{hyp:alternating-thick-word} at an even letter \(c\).
Let \(w(a)\) be the tame-invariant vector with
\((\gamma-1)w(a)\in XN_c\) and \(q^{\eta+B_c}w(a)\in N_c\),
with the coefficient \(d_0(a)\) and filtered space \(L_\eta\)
of Paragraph~\ref{par:thick-constant-data}. There is
\(x\in N_{c-2}\) such that
\[
 \ell^{e_c}d_0(a)-\Phi^2\operatorname{sp}_{c-2}(x)\in L_\eta,\qquad
 e_c\le\sigma_2+B_{c-1}-b_c/\ell+L_0+3.
\]
\end{lemma}
\begin{proof}
Write \(\overline d\) for the class of \(d_0(a)\) modulo \(L_\eta\)
and \(y_0=B_{c-1}\). The coefficients of \(w(a)\)
in nonzero degrees \(j<b_c\) lie in \(L_\eta\). For \(j\le\eta\)
this follows from \(\mathrm{Fil}^{-j}\subset L_\eta\).
For \(\eta<j<b_c\) the two filtrations agree, by the gap
between the negative and central jumps at the even letter.
Take \(D=2b_c+S+\eta\), \(P=\phi(q)\), and
\(\sigma_H=\lceil 1+D/(\ell-1)^2\rceil\).
We first obtain the integral coefficient relation
\[
 \ell^{\sigma_H}q^\eta P^{y_0}\overline d
       \equiv \sum_j a_j(q)\mu_j \pmod{q^{b_c+\eta}} .
\]
The coefficients \(a_j\) are in \(O[[q]]\) and each \(\mu_j\) is
the image of \(\Phi^2\operatorname{sp}_{c-2}(x_j)\) for some
\(x_j\in N_{c-2}\).

The pole input \(q^\eta w(a)\in\langle\Phi N_{c-1}\rangle\)
and the opposite height at \(c-1\) give
\(q^\eta P^{y_0}w(a)\in\langle\Phi^2 N_{c-2}\rangle\).
Transport an entire truncated horizontal basis and its inverse
from \(c-2\), now by two Frobenius steps. We give the scalar
integrality for both matrices. Truncate at degree
\(J=\lceil D/(\ell-1)\rceil\) so the first omitted degree is at
least the degree needed for the \(D\)-jet test. The inverse has the
same denominator bound \(A_a=a+v_\ell(a!)\) in degree \(a\),
since
\[
 A_{i+j}-A_i-A_j=v_\ell\binom{i+j}{i}\ge0.
\]
Iterating this inequality controls every product in the inverse matrix series.

The degree-\(a\) scalar \((\phi^2 X)^a\) is a unit times
\[
        X^a q^{(\ell-1)a}P^{(\ell-1)a}.
\]
In the \(q\)-parameter \(P=\ell\epsilon+q^\ell v\) with
\(\epsilon,v\in O[[q]]^\times\), by
\(\phi(X)=Xq^{\ell-1}\) times a unit and \(X=q-\ell\).
In the binomial expansion of \(P^{(\ell-1)a}\), an index \(j\)
contributes a factor \(q^{\ell j}\ell^{(\ell-1)a-j}\).
For \(a\ge1\), the term with index \(j\) is divisible by \(Xq^D\) as soon as \((\ell-1)a+\ell j\ge D\). Thus a surviving term has
\[
 j<\frac{D-(\ell-1)a}{\ell},\qquad
 (\ell-1)a-j>(\ell-1)a(1+1/\ell)-D/\ell .
\]
The constant term in the transported matrix has no denominator. The retained terms
have \(\ell\)-content at least
\[
\max\bigl(0,(\ell-1)a(1+1/\ell)-D/\ell\bigr).
\]

Put \(a_*=D/((\ell-1)(\ell+1))\). For \(a\le a_*\) the cost bound
\(\ell a/(\ell-1)\) increases. For \(a\ge a_*\) one subtracts the
linear content term, so the slope becomes
\[
 \frac{\ell}{\ell-1}-(\ell-1)(1+1/\ell)<0\qquad(\ell\ge3).
\]
Extending the cost bound from integral degrees to real arguments, its maximum is attained at the junction of the two linear bounds:
\[
 \max_{a\in\mathbb R_{\ge0}}\left\{\frac{\ell a}{\ell-1}
       -\max\left(0,(\ell-1)a(1+1/\ell)-D/\ell\right)\right\}
 =\frac{\ell a_*}{\ell-1}
 =\frac{\ell D}{(\ell+1)(\ell-1)^2}
 \le \frac{D}{(\ell-1)^2}.
\]

Thus \(\ell^{\sigma_H}\) clears both transported matrices
modulo \(Xq^D\). Denote their integral representatives by \(Y,Y'\),
in the order of the horizontal basis and its inverse. The first
omitted degree contributes nothing because \(Xq^D\mid\phi^2(X)^J\).
For the integral coordinate vector \(\mathbf v\) of
\(q^\eta P^{y_0}w(a)\) in the two-step Frobenius span, put
\begin{equation}\label{eq:thick-two-step-inverse}
 \mathbf a=Y'\mathbf v,\qquad
 Y\mathbf a=YY'\mathbf v\equiv\ell^{2\sigma_H}\mathbf v
                                              \pmod{Xq^D}.
\end{equation}

Evaluate at a primitive \(\ell\)-th root. Two Frobenius steps take
the initial cyclotomic argument to the identity. The columns of
\(Y\) in the Frobenius basis are \(\gamma\)-invariant to precision
\(Xq^D\), as in the one-step construction. Their jets through
order \(D-1\) therefore equal the constants
\(\ell^{\sigma_H}\Phi^2\operatorname{sp}_{c-2}(x_j)\).
In the quotient by \(L_\eta\) these become
\(\ell^{\sigma_H}\mu_j\).

Write a bar for projection to this quotient.
The jump gap proved above gives
\[
 \overline{w(a)}\equiv\overline d\pmod{q^{b_c}},\qquad
 q^\eta P^{y_0}\overline{w(a)}
    \equiv q^\eta P^{y_0}\overline d\pmod{q^{b_c+\eta}} .
\]
These congruences use \((q)=(t)\) after evaluation.
Divide the evaluation of \eqref{eq:thick-two-step-inverse} by \(\ell^{\sigma_H}\) over
the coefficient field. Since \(D\ge b_c+\eta\), it gives the claimed
relation with the integral coordinates \(a_j(q)\) of \(\mathbf a\).
The translation \(q=X+\ell\) in the integral series is valid by
completeness of \(O\).

Let \(L\) be the \(O\)-span of the \(\mu_j\). It is finite free.
The coefficient of \(q^\eta\) shows
\[
 \overline d\in K_O L
\]
because \(P(0)\ne0\). Pass to an \(O\)-basis of \(L\), write
one coordinate of \(\overline d\) as \(\theta\), and divide the
resulting congruence by \(q^\eta\). This gives
\(\ell^{\sigma_H}\theta P^{y_0}\) integral modulo \(q^{b_c}\).
Put \(j_*=\lfloor(b_c-1)/\ell\rfloor\le y_0\), the inequality
following from \(b_{c-1}\ge h/2\). Among
\(j=0,\ldots,j_*\), choose the least index \(j_1\) minimizing
\(v_\ell\binom{y_0}{j}+y_0-j\).

Larger indices contribute only at degrees \(>\ell j_1\),
and smaller ones have strictly larger valuation by the minimal
choice of \(j_1\). The coefficient from \(j_1\) is a unit
multiple of \(\binom{y_0}{j_1}\ell^{y_0-j_1}\), since the
constant terms of the two units are units. Thus
\[
 \begin{aligned}
 v_\ell([q^{\ell j_1}]P^{y_0})
   &= y_0-j_1+v_\ell\binom{y_0}{j_1}\\
   &\le y_0-j_*+v_\ell\binom{y_0}{j_*}
    \le y_0-b_c/\ell+1+\log_\ell y_0 .
 \end{aligned}
\]
The last inequality uses the factorial floor sum and
\(j_*=\lfloor(b_c-1)/\ell\rfloor\).
It follows that every coordinate of \(\ell^{e_c}\overline d\)
is integral for
\[
e_c=\sigma_H+\lceil y_0-b_c/\ell+2+\log_\ell y_0\rceil .
\]
Lifting the resulting \(O\)-linear combination of the \(\mu_j\)
gives \(x\in N_{c-2}\) with
\[
 \ell^{e_c}d_0(a)-\Phi^2\operatorname{sp}_{c-2}(x)\in L_\eta .
\]
Since \(\sigma_H\le\sigma_2+1\) and \(y_0\le H_{\rm loc}\),
this exponent satisfies
\[
 e_c\le \sigma_2+B_{c-1}-b_c/\ell+L_0+3 .
\]

\end{proof}
\paragraph{Restrictions on the coefficients at the even letters.}
\begin{lemma}[Reduction to a Wach residual]
At an even letter let \(N_g,N_d\) count the restrictions by generator number and by length, respectively. To reduce to an integral Wach residual it suffices to take
\[
 N_g=3,\qquad N_d=1,
\]
with each of the three generator quotients generated by at most \(m\) elements. For \(m=0\) the reduction holds on the
whole family at exponent \(E_{\rm lev}+S\).
\end{lemma}
\begin{proof}
The height cost cancels before taking a bound over the letters.
Indeed \(e_c+\kappa-B_{c-2}\le B_{\rm lev}\), where
\[
 B_{\rm lev}=\frac{h}{\ell(\ell-1)}+\frac{S+1}{\ell-1}
              +S+E_G+2\sigma_2+6+L_0 .
\]
The calculation uses
\((B_c+1)/(\ell-1)-b_c/\ell
 =b_c/(\ell(\ell-1))+(S+1)/(\ell-1)\).
Put \(E_{\rm lev}=\lfloor B_{\rm lev}\rfloor\).
As \(e_c+\kappa-B_{c-2}\) is integral, the bound gives
\[
  E_{\rm lev}-e_c-\kappa+B_{c-2}\ \ge0,\qquad
  \ell^{E_{\rm lev}}d_0(a)
   \equiv \ell^{\kappa-B_{c-2}}\Phi^2\operatorname{sp}_{c-2}(x)
                      \pmod{L_\eta}
\]
for \(x\in N_{c-2}\). The correction is obtained by multiplying
the earlier choice of \(x\) by the nonnegative power
\(\ell^{E_{\rm lev}-e_c-\kappa+B_{c-2}}\).

Let \(K_\eta\) be the kernel of
\(\Phi^2\operatorname{sp}_{c-2}\bmod L_\eta\) on \(N_{c-2}\).
Multiplication by \(\ell^{\kappa-B_{c-2}}\) is invertible
on the coefficient-field quotient. Hence two choices of \(x\)
differ by \(K_\eta\), and addition of congruences gives an
\(O\)-linear map to \(N_{c-2}/K_\eta\). Composing with the
further quotient gives
\[
 A'\ \longrightarrow\ N_{c-2}/K_\eta\
       \longrightarrow\ N_{c-2}/(P_{c-2}+K_\eta).
\]
Write \(Q=N_{c-2}/(P_{c-2}+K_\eta)\) and \(I\subset Q\) for
the image of \(A'\). The kernel consists of the elements for which
one can choose \(x\in P_{c-2}\).

By Lemma~\ref{lem:thick-even-covolume}, there is a submodule
\(U\subset Q\) generated by at most \(m\) distinguished classes
with normalized length \(\ell_{\rm norm}(Q/U)\le C_D\).
The two restrictions on \(A'\) have successive quotients read from
\begin{equation}\label{eq:thick-restriction-quotients}
 0\longrightarrow I\cap U\longrightarrow I
       \longrightarrow (I+U)/U\longrightarrow0,\qquad
 \ell_{\rm norm}((I+U)/U)\le C_D .
\end{equation}
For the first term of \eqref{eq:thick-restriction-quotients}, the preimage of \(I\cap U\) in a free cover \(O^m\to U\) is free
of rank at most \(m\). Its image generates \(I\cap U\); killing that
image costs one generator restriction.

On the resulting kernel the new coefficients satisfy
\[
 \ell^{E_{\rm lev}}d_0(a)\equiv\ell^\kappa\Phi z\pmod{L_\eta},
 \qquad z\in Z:=\Phi^2\operatorname{sp}_{c-3}(N_{c-3}).
\]
Set \(K'=\ker(Z\xrightarrow{\,\Phi\,}D_c/L_\eta)\), so \(z\) is determined modulo
\(K'\). Requiring a representative in \(C\) costs at most \(m\)
generators: the quotient \(Z/((Z\cap C)+K')\) is an image of
\[
 Z/(Z\cap C)\hookrightarrow D_{c-1}/C,\qquad
 \dim_K(D_{c-1}/C)\le m ,
\]
where \(K\) denotes the coefficient field. The source of the
injection is finitely generated over \(O\), since \(Z\) comes from
a Wach lattice modulo \(X\). It is therefore free of rank at most
\(m\), and submodules of its quotient also need at most \(m\)
generators. For \(m=0\) the inclusion
\[
 \ell^S N_{c-2}\subset P_{c-2}
\]
allows \(\ell^S x\) as representative and
\(C=D_{c-1}\).

Choose the source vectors \(y\) from the first-level construction on
an \(O\)-basis of the restricted family, then extend linearly.
The constant of \(\ell^{E_{\rm lev}}w(a)-f(y(a))\) lies in
\(L_\eta\). Requiring it to lie in \(L_0^{\rm fil}\) costs one more
generator restriction of size at most \(m\). When \(m=0\) these
two filtration steps agree. For \(j\ne0\) the difference of
the coefficients lies in \(\mathrm{Fil}^{-j}D_c\), since the
source image has vanishing coefficients for \(j<B_c\) and
\[
 \mathrm{Fil}^{-j}D_c=D_c
\]
for \(j\ge B_c\).
Tame transport carries this property from the chosen root to
all primitive \(\ell\)-th roots. The rational cyclotomic
evaluation test now puts the vector difference in \(N_c\).

For the integral test, start with the known denominator
\(q^{\eta+B_c}\) and divide its coordinates in a Wach basis by the
distinguished polynomial associated to that power. At each root the
required coefficient vanishing follows from membership in the
completed Hodge lattice. The remainder is therefore zero by
Weierstrass division. This proves the reduction to a Wach vector
residual with three generator restrictions and one length restriction.
For \(m=0\) it works on the whole family with exponent
\(E_{\rm lev}+S\).

\end{proof}
\paragraph{The residual modulo \(X\).}
\begin{lemma}[Lifting the remaining class modulo \(X\)]
Three generator restrictions and two length restrictions lift the residual modulo \(XN_c\) after multiplication by \(\ell^{e_6}\), where
\[
 e_6=S+\lceil E_G+\sigma_2\rceil.
\] If \(m=0\), exponent
\(e_6+2S\) suffices on the entire family.
\end{lemma}
\begin{proof}
Use the same filtered crystalline comparison and
Lemma~\ref{lem:thick-even-covolume} at both even letters \(c,c-2\). Let \(v\) be the
integral Wach residual on the family.
Restrict first to \(v\in P_c\), costing one length and one
generator restriction, as above. There is a vector \(v_1\in N_{c-1}\)
with \(\operatorname{sp}_c(v)=\ell^{B_c}\Phi\operatorname{sp}_{c-1}(v_1)\).
At the odd letter the opposite Wach height gives
\(q^{B_{c-1}}N_{c-1}\subset\langle\Phi N_{c-2}\rangle\).
Specializing at \(X=0\) gives \(v_2\in N_{c-2}\) with
\[
 \ell^{B_{c-1}}\operatorname{sp}_{c-1}(v_1) =\Phi\operatorname{sp}_{c-2}(v_2).
\]

The specializations of \(v_1,v_2\) are linear in the family,
because Frobenius is invertible on the crystalline spaces and the
scaling factors are nonzero. The kernel of specialization is
\(XN_{c-2}\), so the map \(v_2\bmod P_{c-2}\) is well defined.
Restricting it to zero costs one length and one generator
restriction. Then for some \(v_3\in N_{c-3}\) one has
\(\operatorname{sp}_{c-2}(v_2)
=\ell^{B_{c-2}}\Phi\operatorname{sp}_{c-3}(v_3)\).
Restrict to the kernel of the image of
\(z=\Phi^2\operatorname{sp}_{c-3}(v_3)\) in \(D_{c-1}/C\).
This is one generator restriction of size at most \(m\).

Now \(e_6=S+\lceil E_G+\sigma_2\rceil\) satisfies
\[
 e_6+B_{c-2}=u_0+k+B_{c-1},\qquad
 u_0=\lceil E_G+\sigma_2\rceil-\lceil E'+\sigma_2\rceil\ge0 .
\]
Choose \(z_1\in N_{1,c-1}\) with
\(f(\operatorname{sp}_{c-1}z_1)=\ell^k z\).
Combining the three specialized equalities gives
\[
 \operatorname{sp}_c(\ell^{e_6}v)
   = f\bigl(\operatorname{sp}_c(\ell^{u_0}q^{B_c}\Phi z_1)\bigr).
\]
The vector \(\ell^{u_0}q^{B_c}\Phi z_1\) belongs to \(N_{1,c}\)
by the source height, so it is tame invariant and has
\(\gamma-1\) difference in \(XN_{1,c}\).
The residual is therefore lifted modulo \(XN_c\) at the cost of
three generator and two length restrictions. When \(m=0\), one
uses \(\ell^S N_i\subset P_i\) twice and \(C=D_{c-1}\).
This gives a lift on the entire family with exponent \(e_6+2S\).

\end{proof}
\begin{proposition}[Finite-class loss on an alternating word of larger jumps]\label{prop:thick-inert-colength}
For the marked crystalline surjection and signed jump word of
Lemma~\ref{hyp:alternating-thick-word}, the chain uses at most \(4n\) generator restrictions of
size \(m\) and \(3n/2\) length steps. Let \(A_{\rm end}\) be the last submodule in the restricted chain and
\(\pi:A\to M_2=H^1_f(F,T_2)\) the chosen surjection. There is an exponent \(e\) such that
\[
 \ell^e\pi(A_{\rm end})\subseteq f_*H^1_f(F,T_1),\qquad
 e\le \frac{32H_{\rm loc}}{\ell^2}+4S+\frac{10\eta}{\ell-1}+25L_0 .
\]
For \(m=0\) the construction gives a lift on the full finite
group with the additional allowance \(3S\). With the separately
available coarse lift, \eqref{eq:thick-chain-mass} converts these
losses to the normalized place colength.
\end{proposition}

\begin{proof}
Paragraphs~\ref{par:thick-initial-classes} and
\ref{par:thick-integral-correction} decompose the target finite
classes after \(s=s_1+2\sigma_1\), with at most \(n\) generator
restrictions of size \(m\). They realize the vectors \(w_c\)
on the target and compare their finite classes under the
same marked map modulo \(XN_c\).

\proofstep{Restrict the same free family at all even letters, then sum the cut classes.}
Now apply the letter
comparisons successively on the submodules obtained so far.
They are stable under restriction
of the family. The first lift \(y\) and the lift of the residual
combine by adding \(\ell^{e_6}y\) to the latter source vector. There
are \(n/2\) even letters. Counting at every other letter gives the respective maxima
\[
 n+6(n/2)=4n\quad\text{generator steps},\qquad 3(n/2)=3n/2\quad\text{length steps}.
\]
Each cut class has a finite lift after the same power.
Summing their classes lifts the scaled original class, since the
decomposition holds modulo coboundaries. This costs
\(e=s+E_{\rm lev}+e_6\), independently of the number of letters.
For \(m=0\) the extra powers on the two comparisons give \(e+3S\)
on the entire finite group.

\proofstep{Combine the decomposition and residual losses, using \(\ell\ge11\).}
The successive restrictions have been imposed on one family of lifts. Their losses can now be added without counting a different family at each even position.

Here is the numerical bound from this construction. Put
\(u=1/(\ell-1)\) and
\(\sigma=1+(\lfloor hu\rfloor+1+\eta)u\).
The resulting bounds for the decomposition and its remainder are
\[
\begin{aligned}
 s&\le(2A_0+2\eta+S)u+L_0+2\sigma+2,\\
 e_6&\le S+E_G+\sigma_2+L_0 .
\end{aligned}
\]
Adding these bounds and the bound for \(E_{\rm lev}\), using
\(\sigma\le1+hu^2+u+\eta u\), gives
\[
 e\le 8A_0u+8\eta u+(4S+3)u+9H_{\rm loc}u^2+2S+5L_0+13 .
\]
Since \(u\le1/10\) and
\(A_0u\le (2+u)H_{\rm loc}u^2+(4+L_0)u\), we obtain
\[
 e\le25.8H_{\rm loc}u^2+8\eta u+2.4S+5.8L_0+16.5
   \le \frac{32H_{\rm loc}}{\ell^2}+4S+
              \frac{10\eta}{\ell-1}+25L_0 ,
\]
using \(L_0\ge1\) and \(\ell/(\ell-1)\le 11/10\).
With the common coarse finite-class lift \(e'\),
\eqref{eq:thick-chain-mass} gives the bound for sequences with larger jumps.

The target coefficient bound uses the estimate on the rank of the kernel. The filtered crystalline estimate gives the realization at the first level. Apply the two lattice bounds for even letters in their respective cases.
 Combine the squared Frobenius bound on the residual lattice for an even letter with the coefficient lift.
 The resulting lattice bounds give the length at the chosen prime.
The local argument applies the integral Wach and Herr comparisons
to the marked crystalline models and uses the corresponding
coarse class lift for the generator restrictions.
\end{proof}


The common entry in cases C and D is
Proposition~\ref{prop:ordinary-inert-place-entry}.

\subsubsection{Smaller jumps and the two estimates at ordinary inert primes}
\label{sec:inert-row-table}
\hypertarget{source-inert-row-table}{}

When smaller jumps occur, the uniform restriction used for the larger jumps does not cover the whole sequence. Decomposing it into short arcs gives the missing estimate; it can then be combined with the larger-jump case at the same inert place.

Here \(\ell\ge11\) and \(n=n_\lambda\) is the local unramified degree. In the
case with some smaller jumps, let \(s_i\) be the sizes in the local word.
Use \(b_0,w,S,\eta\) for its threshold for the larger jumps, bound for the smaller jumps,
spread and central width. Each
\[
 s_i>\eta.
\]
A size in the larger range is
in \([b_0,b_0+\eta]\), and a size in the smaller range is at most \(w\).
Use the signed jump intervals of Section~\ref{subsec:inert-rows}.

At every letter at most \(m\)
target jumps, counted with multiplicity, lie in \([-\eta,\eta]\).
On the marked row the signed windows come from the source lattice
in the sum of twisted duals. The bound on the central multiplicity
comes from the twisted dual target at each labeled embedding,
applied with the same enumeration of the embeddings as the word.
Both statements therefore apply to a size in the smaller range.
\label{cor:inert-row-table-1}
There is a position with smaller jumps, the signs alternate cyclically,
and there are at most two consecutive positions with larger jumps.
The numerical hypotheses are
\[
 20w,20S,20\eta\le b_0.
\]

Put \(h=b_0+\eta+S\) if there is a position with larger jumps and \(h=w+S\)
otherwise, and put
\[
 \gamma_h=\log_\ell(2h),\qquad
 S_g=h/\ell^2+(w+S+\eta)/\ell+1+\gamma_h .
\]
Let \(B\) be the indicator that some position with smaller jumps \(i\)
has a neighbor with larger jumps and satisfies
\[
 \abs{\ell s_i-b_0}\le 320
              (\ell\eta+\ell S+b_0/\ell+s_i+\ell\gamma_h).
\]
The common place indicator in \(F_{\ell,\lambda}\) is the
existence of such a word on the marked data, over all good
characters and compatible local coefficient data. Therefore every
local \(B\) is at most this place indicator.

Cut before each position with smaller jumps, numbering the cut by
its preceding letter. Begin the enumeration at one such position, so
the construction includes the arc that crosses the end of the cycle.
Let \(\mathcal C\) be the set of cuts. The arc ending at \(c\)
consists of one smaller letter followed by \(k_c\) larger letters.
Consequently
\[
 a_c=1+k_c,\qquad 0\le k_c\le2,\qquad
 \sum_{c\in\mathcal C}a_c=n .
\]

Suppose first that a larger letter occurs and \(40S_g>b_0/4\).
The two jump bounds of the coarse lift are at most \(h\); its
maximum width is at most \(h+\eta\). The coarse exponent calculation
in the proof of Proposition~\ref{prop:ordinary-inert-place-entry}
therefore bounds the exponent \(e'\) of a full finite lift by
\(7.1 h/\ell+7.1+2\gamma_h\). Since \(h\le1.1b_0\),
\(b_0<160S_g\), and \(\ell\ge11\),
\[
 7.1 h/\ell <113.6 S_g,\qquad
 7.1+2\gamma_h\le7.1S_g,\qquad e'<160 S_g .
\]

In the remaining range, use Lemmas~\ref{lem:two-letter-wach-lift},
\ref{lem:three-letter-negative-lift} and~\ref{lem:three-letter-positive-lift}
on arcs of length two and three. To compare the arc lift with the cut
decomposition we need the exponents \(R+2\sigma\) and \(R+3\sigma\).
Put \(g=w+S+\eta\), \(q_0=\ell-1\), and
\(M=\lfloor2h/q_0\rfloor+1\). There is one smaller letter
immediately after the cut. The corresponding correction satisfies
\[
 \begin{aligned}
 \sigma
 &=\left\lceil1+\frac{\ell(2(w+S)+M)}{q_0^2}\right\rceil\\
 &\le 2+\frac{\ell}{q_0^2}
       +\frac{2\ell g}{q_0^2}+\frac{2\ell h}{q_0^3}
 \le 3\left(1+\frac{g}{\ell}+\frac{h}{\ell^2}\right).
 \end{aligned}
\]
Indeed \(\ell/q_0\le11/10\) bounds the two variable coefficients
after normalization by \(2(11/10)^2\) and \(2(11/10)^3\).
For the recurrence of Section~\ref{sec:inert-long-orbit}, take
symmetric widths at most \(2h\) throughout and at most \(2g\)
at the first letter after the cut. With
\(c_H=\lceil1+\log_\ell(1+\max_i H_i)\rceil\) one has
\[
 R\le c_H + \frac{\ell}{q_0(q_0-1)}
                    (2g+c_H+3+2h/q_0).
\]

Here \(c_H\le 2+\tfrac15+\gamma_h\) by
\(1+2h\le3h\) and \((3/2)^5<11\).
The coefficients of \(g/\ell\) and \(h/\ell^2\) in the bound
are at most \(2\cdot11^2/(10\cdot9)<3\) and
\(2\cdot11^3/(10^2\cdot9)<3\), respectively. Since
\(\ell/(q_0(q_0-1))\le11/90\), the remaining terms
are at most \(3(1+\gamma_h)\). Combining these estimates gives
\begin{equation}\label{eq:inert-short-cut-budget}
 R\le3S_g,\qquad \sigma\le3S_g,\qquad
 R+2\sigma\le9S_g,\quad R+3\sigma\le12S_g .
\end{equation}
With the cut costs of \eqref{eq:inert-short-cut-budget}, the arc lift at \(c\) itself has exponent at most
\[
 \sigma+\kappa_c+\lfloor 40S_g+\eta B\rfloor ,
 \qquad
 \kappa_c\le 13S+3(1+\gamma_h).
\]
It lands in the finite image modulo a span of at most
\(m r_c(a_c+1)\) target finite classes, with \(r_c\le a_c\).
We derive the coordinate lift and the fixed exceptional span below.

\paragraph{Coordinates on a short arc.}\label{par:short-wach-frame}
\begin{proposition}[Common coordinates for short-arc comparisons]\label{prop:short-arc-coordinates}
For a short arc, use the same integral vertical and
horizontal frames and the corresponding cyclotomic levels.
For an arc ending at \(c\) of length \(a=a_c\), the offset \(0\le j<a\)
has size \(u_j=s_{c-j}\) and cyclotomic factor \(P_j=\phi^j(X+\ell)\). The horizontal and vertical comparisons use the same positive pole and negative zero factors
\[
 D=\prod_{j\ {\rm positive}}P_j^{u_j},\qquad V=\prod_{j\ {\rm negative}}P_j^{u_j+S}.
\]
The two ordinary lifting rules use these coordinates. In the single-letter edge cases the
positive rule applies to every input, and the negative rule
uses the empty pole set and the initial constant lift.
\end{proposition}

\begin{proof}
\proofstep{Fix the same cut coordinate and ordered cyclotomic levels for both rules.}
Write \(a=a_c,\ c'=c-a,\ G=w+S,\ u_j=s_{c-j}\), and
\(P_j=\phi^j(q)\), so \(P_0=q=X+\ell\).
Let \(\mathcal R=O[[X]]=O[[q]]\). All ideals of fractions below
are in the component at \(c\).
Call an offset positive or negative according to its large Hodge
jumps. This sign is used even at the first letter of small size.
The full denominator for transport over \(a\) steps and the two
smaller products determined by the jumps are
\[
 \Lambda=P_{a-1}^{G}\prod_{j<a-1}P_j^h,\qquad
 D=\prod_{j\ {\rm positive}}P_j^{u_j},\qquad
 V=\prod_{j\ {\rm negative}}P_j^{u_j+S}.
\]
Choose a basis of the tame-invariant Wach lattice at \(c'\).
The unique formal \(\gamma\)-fixed basis with the same constants is
its horizontal basis.

Write \(U(X)=I+\sum_{r>0}U_rX^r\) for the change
to the horizontal basis, and let \(F_r(U_1,\ldots,U_{r-1})\)
be the degree-\(r\) residual computed with \(U_r=0\).
The recursion is
\[
       (\omega_X^r-1)U_r=-F_r(U_1,\ldots,U_{r-1}).
\]
Its new denominator has order \(1+v_\ell(r)\) at step \(r\).
Thus through that step the denominators divide
\(\ell^{r+v_\ell(r!)}\). This also holds for the inverse change
of basis, by the factorial inequality for products of coefficients.
Use the basis truncated before degree \(M\). These are coordinates
on the given lattice map: a map of Wach lattices acts on horizontal
columns by its matrix of constants at \(c'\), since both sides solve
the same equation with the same constant. Frobenius transports these
constants to the crystalline space at \(c\).

Here are details on the precision of these coordinates. Put
\[
 P(T)=\langle\Phi^aN_{c'}(T)\rangle_{\mathcal R}.
\]

Represent \(v\in\Lambda^{-1}P(T)\) modulo \(X\Lambda P(T)\)
in the transported horizontal basis, allowing powers of \(\ell\)
to clear coefficients in this comparison. Equivalently multiply
by \(\Lambda\) and compare modulo \(X\Lambda^2 P(T)\).
The tails beyond the truncation have a factor
\(\phi^a(X)^M\), divisible by \(X\Lambda^2\) since
\((\ell-1)M>2h\). The finite coefficients, and the inverse matrix,
can be transported integrally at cost \(\sigma\). Indeed on
the first step one reduces \(\ell^\sigma\phi(X)^r/\ell^{r+v_\ell(r!)}\)
modulo \(X^Mq^{2G}\).

The term \(X^r q^{(\ell-1)r}\) pays the
denominator after this precision, by the binomial division of
Paragraph~\ref{par:wach-cut-correction} with one forward step and
precision \(2G\). More explicitly, for \(k=r+v_\ell(r!)-\sigma>0\)
one needs
\[
 (\ell-1)r\ge 2G+k+M-r-1.
\]
Then modulo \(\ell^k\)
the binomial expansion of \((X+\ell)^{(\ell-1)r-2G}\) is divisible
by \(X^{M-r}\). If \(k>0\), the coefficient bound gives
\(r>(\ell-1)\sigma/\ell\). Together with the definition of
\(\sigma\) this yields the preceding power inequality.

In
the binomial expansion the terms of \(\ell\)-degree \(\ge k\)
vanish modulo \(\ell^k\); every other term has \(X\)-degree
at least \((\ell-1)r-2G-k+1\ge M-r\). Thus the termwise
division is legitimate at the first precision. The next
\(a-1\) Frobenius maps give
\[
             \phi^{a-1}\bigl((X^Mq^{2G})\bigr)
                                     \subseteq (X\Lambda^2)
\]
by the same level bounds defining \(\Lambda\). Thus the transported matrices are represented by
integral lifts of \(\ell^\sigma\phi^a(U)\) and
\(\ell^\sigma\phi^a(U^{-1})\) modulo \(X\Lambda^2\)
in the Frobenius basis. If \(A,B\) denote these two integral lifts in the displayed order, their common precision gives
\[
 AB\equiv BA\equiv\ell^{2\sigma}I\pmod{X\Lambda^2}.
\]


For clarity, the coefficient test in these coordinates is as
follows. If a constant direction \(d\) at the end of the arc
belongs to \(\mathcal F_j^{-\alpha_j}\) for each transported
filtration, then \(d\) with coefficient in
\((\prod P_j^{\alpha_j})\mathcal R\) passes the Wach test at
each level. Sums of such terms after scalar extension give the
rational test; the bounded integral numerators are required as
well.

To see the rational assertion, separate the levels by the
Chinese remainder theorem for their distinguished polynomials
and take a basis of the filtration at each level. For a basis vector \(d\) of jump \(u\), the corresponding generator and its Laurent order are
\[
 t^{-u}d,\qquad \operatorname{ord}_t(t^{-u})=-u,
\]
as in
Paragraph~\ref{par:wach-level-models}. Multiplying the other
level factors to sufficiently high orders realizes each of these
finite jets without imposing further conditions there.
Divisibility of bounded numerators by a product of these
polynomials descends to \(\mathcal R\) by monic division.

Let \(u\) denote a coordinate vector satisfying all the
local tests and let \(n\) be a vector in \(N_c\).
The forward comparison and its inverse give, in their
respective quotients by the transported precision,
\[
   f_{\rm hor}(u)=[\ell^\sigma u]_{\rm Wach},
   \qquad g_{\rm hor}(n)=[\ell^\sigma n]_{\rm coordinates},
   \qquad g_{\rm hor}f_{\rm hor}=\ell^{2\sigma}.
\]
Thus an admissible coordinate vector produces a Wach class,
and the restricted inverse sends a Wach vector back into the
tested coordinate module. Indeed the transported forward
and inverse matrices above give integral lifts with the bounded
denominator, and changes by \(X\Lambda P(T)\) lie in \(XN_c\)
by the height bounds.

For an admissible vector one tests
\((\gamma-1)/X\) in the same manner. This gives an integral
vector with \((\gamma-1)v\in XN_c\), since its horizontal
columns are fixed to the indicated precision. A product of the
\(P_j\) clears its denominators in \(N_c\). These statements
apply to subarcs beginning at \(c'\) with their own product
\(\Lambda\). \label{cor:inert-short-frame-1}

We will lift a coordinate \(f\) on a target horizontal column
\(d\) which belongs to \(\mathcal F_j^{u_j}\) at every positive
level. Its scalar hypotheses are
\begin{equation}\label{eq:short-scalar}
 Df\in\mathcal R,\qquad
 D(\ell^{k_j}f-x_j)\in P_j^{u_j+S}\mathcal R
 \quad(j\ {\rm negative})
\end{equation}
for some \(k_j\ge0,\ x_j\in O\). A lift at cost \(e\) represents
\(\ell^e f d\) by the image of an admissible source vector up to
\(XN_c(T_2)\) and the coordinate precision. There are two ways
to construct it.
\begin{enumerate}
\item If \(Df\in V\mathcal R\), then \(fd\) passes the target
Wach test. The poles at the positive levels are allowed by the
filtration of \(d\); the zeros at the negative levels accommodate
the lowest possible jump. The coordinate comparison and scalar image bound at \(c\) use successive multipliers
\[
 \ell^\sigma\ell^R=\ell^{R+\sigma}.
\]
Tame averaging preserves the lifted image, so no extra exponent enters this case.
\item Suppose the poles are confined to a subset of the positive
levels and an integral source direction maps to \(\ell^t d\) and
lies in \(\mathcal F_j^{u_j}\) there. At each omitted positive
level also require that \(f\) is constant modulo \(P_j^\eta\)
after clearing scalars and the retained denominator. Then the
coefficient \((\gamma f-f)/X\) times this source direction passes
the level tests. For each level factor, write \(\gamma P_j/P_j=1+Xh_j\) with \(h_j\in\mathcal R\). Then
\[
 (\gamma P_j-P_j)/X=h_jP_j.
\]

For every \(n\in\mathbb Z\) the ratio \((1+Xh_j)^n\) belongs to \(1+X\mathcal R\). Thus
\[
 (\gamma P_j^n-P_j^n)/X\in P_j^n\mathcal R ,
\]
including for negative powers. The constant disappears
at each omitted positive or negative level. Equation~\eqref{eq:short-scalar}
therefore supplies the required zeros at negative levels, and
\(\eta\) supplies them at omitted positive levels. The retained
positive poles pass by the source filtration. The admissible-vector
construction gives the lift at \(t+\sigma\). For an integral constant, transport the constant term of the initial target basis lift at \(c'\). Its coefficient differences vanish, so its cost is at most \(R+\sigma\).
\end{enumerate}

In the second rule the congruences use the denominator for the
retained positive levels. Costs of two summands may be padded
to their maximum, by linearity of the source map and of the
coordinate comparison. For a single positive letter the first
rule applies to every \(f\). For a single negative letter use
the second with empty pole set and the initial constant lift.
\end{proof}


\begin{auxlemma}[Separating a scalar by Weierstrass division]\label{lem:short-arc-content}
Let \(O\) be a coefficient DVR with fraction field \(K\),
\(v_p(p)=1\), and let \(J\subset\mathcal R=O[[q]]\) be generated by
a distinguished polynomial. Take \(D\in\mathcal R\) and
\(f\in\operatorname{Frac}(\mathcal R)\). Assuming that the residue
\(\bar D\in\mathcal R/J\) is nonzero, denote by \(c_J(D)\) the minimum of the
valuations of its Weierstrass remainder coefficients. Let \(k,d\)
be nonnegative integers with \(c_J(D)\le d\). If
\(Df\in\mathcal R\), \(x\in K\), and \(p^k Df-Dx\in J\), then there is an integral constant satisfying
\[
 \exists x'\in O:\qquad D(p^d f-x')\in J.
\]
\end{auxlemma}
\begin{proof}
Write \(H=\mathcal R/J\), a free \(O\)-module with the monomials
below \(\deg J\) as basis, where \(\deg J\) is the degree of the
distinguished generator. In \(H\otimes_O K\) the given equality
implies
\[
 x\bar D=p^k\overline{Df}\in p^k H,\qquad
 v_p(x)+c_J(D)\ge k .
\]
Thus \(x'=p^{d-k}x\) is integral, including when \(k>d\).
Both \(p^d Df\) and \(x'D\) lie in \(\mathcal R\), and their
images agree after tensoring \(H\) with \(K\). Freeness makes
\(H\to H\otimes_O K\) injective, which proves the assertion.
\end{proof}

\begin{lemma}[Two-letter lifting on a short arc]\label{lem:two-letter-wach-lift}
In the models of Proposition~\ref{prop:short-arc-coordinates},
write \(e_2\) for the exponent of a two-letter correction. One can achieve
\[
 e_2\le15L+1+\eta B,
\]
where \(L\) is the local size from the short-arc comparison and \(B\) its
exceptional indicator.
\end{lemma}

\begin{proof}
\label{par:short-two-levels}
Use \(u=s_{c-1}\) for the small size and \(v=s_c\) for the large
one. Put \(p=\ell,\ L=S_g\) within this calculation. Denote by
\(A_i\) the ceiling of the least ordinary Wach image exponent in
\(\pi\) at \(i\), divided by \(p-1\) before taking the ceiling.
Then \(X^{A_i}\) kills the tame-invariant cokernel. In signed
boxes of widths \(s_i+S+\eta\), the recurrence of
Paragraph~\ref{par:wach-cut-colength} gives
\begin{equation}\label{eq:short-A}
 A_i\le(s_{i+1}+S+\eta)/p+\tfrac52 L.
\end{equation}
Here one can replace \(5/2\) by \(2\) if \(i+2\) is small.

\proofstep{Compute the denominators of the two neighboring level evaluations.}
To check the constants, write \(d=p-1,\ H=h+\eta\),
\(W_j=s_j+S+\eta\) and \(c_H=\lceil1+\log_p(1+H)\rceil\).
The ordinary image exponents satisfy
\(s_i^{\rm im}\le W_{i+1}+c_H+3+s_{i+1}^{\rm im}/d\),
and are at most \(H\) by the signed twist. Thus
\[
 A_i-W_{i+1}/p\le 1+(c_H+3)/d+H/(pd)+H/d^2.
\]
The general bound uses \(H=h+\eta\), \(H+1\le2h\), and
\(p/d\le1.1\).

If \(i+2\) is small, iterating the recurrence
once replaces its last summand by
\[
 \frac{H}{d^2}\quad\rightsquigarrow\quad
 \frac{w+S+\eta+c_H+3}{d^2}+\frac{H}{d^3}.
\]
The respective coefficients of \(h/p^2\) are bounded by
\[
 \frac{p}{d}+\frac{p^2}{d^2}\le2.31,\qquad
 \frac{p}{d}+\frac{p^2}{d^3}\le1.24 .
\]
The \(\eta/p^2\) terms are absorbed in \((w+S+\eta)/p\).
Together with \(c_H\le2+\log_p(2h)\), these give the factors
\(5/2\) and \(2\) in the two versions of \eqref{eq:short-A}.
\label{cor:inert-short-two-1}

Apply Lemma~\ref{lem:short-arc-content} with
\(P=P_1=p\varepsilon+Uq^p\), where \(\varepsilon,U\in\mathcal R^\times\).
We estimate the contents for the two possible denominators.
For \(D=q^v,J=(P^y)\), put \(t=\lfloor v/p\rfloor+1\).
Multiply by \(q^{pt-v}\) and expand
\[
 (q^p)^t=U^{-t}(P-p\varepsilon)^t.
\]
Modulo \(P^y\) the
content of this power is the minimum over \(j<y,\ j\le t\)
of \(t-j+v_p\binom tj\). After division by the lowest
uniformizer power, the least \(j\) attaining it has strictly
smaller order in the residue parameter than the other surviving
terms and than \(P^y\). This proves the minimum formula even for
nonconstant units.

Since multiplication by \(q^{pt-v}\) can only
increase content in the first quotient, any admissible index gives
an upper bound for \(c_{(P^y)}(q^v)\).
For \(D=P^u,J=(q^C)\), use the binomial expansion of
\((p\varepsilon+Uq^p)^u\) modulo \(q^C\). The two bounds are
\begin{equation}\label{eq:two-letter-wach-contents}
 \begin{aligned}
 c_{(P^y)}(q^v)&\le t-i_1+v_p\binom{t}{i_1},
       & i_1&=\min(t,y-1),\\
 c_{(q^C)}(P^u)&\le u-i_2+v_p\binom{u}{i_2},
       & i_2&=\min(u,\lfloor(C-1)/p\rfloor).
 \end{aligned}
\end{equation}
For \(pu<C\) the second line has \(i_2=u\) and the content is zero.
For the binomial terms in \eqref{eq:two-letter-wach-contents}, the carry formula gives \(v_p\binom ti\le\log_p t\).
In particular the two bounds are at most
\[
 \max(0,v/p-y+2+\log_p v)+1,\qquad
 \lceil\max(0,u-C/p+2+\log_p u)\rceil ,
\]
respectively. In the first bound the case \(i=t\) contributes
zero. All contents are normalized by \(v_p(p)=1\).

Suppose first that the large terminal letter is positive. If
\(p(u+S+\eta)+\eta<v\), construct a filtered source constant.
On the target take the coordinate \(P^{u+S}/q^v\) in direction
\(d\), realize it at cost \(\sigma\) and lift at the cut at
cost \(R\). Express the lift in source horizontal coordinates
at the cost of a further \(\sigma\). Multiply these coordinates
by \(q^{v+S}P^\eta\), obtaining integral series. In fact every
nonzero filtered summand has exponents at least \(-(v+S)\)
and \(-\eta\) at the two levels. Divisibility after rational
expansion descends by saturation. In the coordinate \(q\) take
the coefficient of degree
\[
 j=S+p(u+S+\eta).
\]

A filtered summand with level-zero exponent \(\alpha\) can
contribute only if
\[
 v+S+\alpha\le j.
\]
This forces
\(-\alpha\ge\eta+1\); the gap puts its direction in
\(\mathcal F_0^v\). The constant matrix of the map sends this
coefficient to the same coefficient of
\(p^{R+2\sigma}q^S P^{u+S+\eta}d\). Indeed the coordinate error
after clearing these denominators is divisible by \(q^h\) and
\(j<h\). The coefficient of degree \(p(u+S+\eta)\) of
\(P^{u+S+\eta}\) is a unit, by reduction modulo the maximal
ideal. Dividing by this unit supplies the filtered preimage of \(p^{R+2\sigma}d\). The additional coordinate comparison in the second rule raises the exponent to
\[
 (R+2\sigma)+\sigma=R+3\sigma.
\]


In the complementary range take \(D=q^v\), \(y=u+S\). The
first content bound separates the constant at the negative level
from the remainder. Each is lifted at cost \(R+\sigma\) by the
two rules. Writing the extra cost as \(c_{\rm cont}\), the two
ranges give
\[
 c_{\rm cont}\le
 \begin{cases}
    3+\gamma_h,& v<py,\\
    \eta+\eta/p+3+\gamma_h,& py\le v\le p(y+\eta)+\eta .
 \end{cases}
\]
In the second range, combining the two bounds on \(v\) with
\(b_0\le v\le b_0+\eta\) gives
\[
 b_0/p-S-\eta-\eta/p\le u\le(b_0+\eta)/p-S,
\]
the prescribed band.

Now suppose the first small letter is positive. Put
\(A=A_{c-1}\). When \(u\ge A+\eta+1\), take the partial
arc of length one at this letter and the coordinate \(q^{-u}d\)
there, with \(d\) read before the last Frobenius transport.
Realize it at cost \(\sigma\), lift \(X^A\) times this integral
vector, and read the source coordinates at cost \(\sigma\).
Multiply by \(q^G\) and take the coefficient in degree \(G-u+A\).

The degree \(G-u+A\) is below \(G-\eta\), so the degree and gap test places the coefficient in \(\mathcal F^u\). The target coefficient in the chosen degree is one, as
\[
 [q^{A+G-u}](q-p)^A q^{G-u}=1.
\]
The precision includes \(q^{2G}\) after denominators are cleared, so the image of the coefficient is exactly \(p^{2\sigma}d\). Tame averaging costs a unit.
Transport forward and apply the second rule at pole set
\(\{1\}\). The cost is \(3\sigma\).

When \(u\le A+\eta\), apply the second content bound with
\(D=P^u,\ C=v+S\), and lift the separated remainder and
constant at cost \(R+\sigma\). There is no extra cost for
\(pu<C\). Otherwise \eqref{eq:short-A} with its improved
constant gives
\[
 u-C/p\le\eta+\eta/p+2L.
\]

If the band is absent, \(pu\ge C\ge b_0\) and
\(pu-b_0>2p\eta\). But
\[
 u-b_0/p\le \eta+(2\eta+S)/p+2L,
\]
so \(\eta<4L\) using
\((2\eta+S)/p\le L\) (since \(w>\eta\)).
Thus the two-letter exponent in all cases is at most
\(15L+1+\eta B\). \label{cor:inert-short-two-2}
\end{proof}


\begin{lemma}[Three-letter lifting with negative outer signs]\label{lem:three-letter-negative-lift}
With \(L=S_g\) and \(B\) the short-band indicator of Section~\ref{sec:inert-row-table},
for the chronological signs \(-,+,-\), one can choose a three-letter correction in the short-arc frames with exponent \(e_3\) satisfying
\[
 e_3\le22L+1+\eta B
\]
in the given DVR normalization.
\end{lemma}

\begin{proof}
\label{par:short-three-levels}
Put
\[
 u=s_{c-2},\ v=s_{c-1},\ z=s_c,
\]
so \(u\le w\)
and \(v,z\in[b_0,b_0+\eta]\). Here \(40L\le b_0/4\).
At a partial arc use level numbers relative to its own endpoint.

\proofstep{Start with the middle positive letter and retain the two negative neighbors.}
First consider signs \(-,+,-\) in chronological order.
If \(p(u+S+\eta)\le v\), work at the middle letter and
write \(A=A_{c-1}\), \(y=u+S+\eta\).
As in the two-letter coefficient extraction use the coordinate
\(P_1^{u+S}/q^v\) in direction \(d\). At this intermediate
point the scalar image bound at a cut is unavailable. Instead
lift \(X^A\) times the realized vector. After the two horizontal
comparisons and multiplication by \(q^{v+S}P_1^\eta\), the
scalar on the target, in the variable \(q\), is
\[
 F=p^{2\sigma}(q-p)^A q^S P_1^y .
\]
Below degree \(v+S-\eta\) every source coefficient lies in
\(\mathcal F^v\) by the gap. To recover an integral direction
it suffices that some such coefficient of \(F\) divides a small
power \(p^t\) in \(O\).

Take
\[
 t=1+\left\lfloor\tfrac{11}{10}
 (2\sigma+(\eta+A)/p+2+\gamma_h)\right\rfloor .
\]
Then \(2\sigma<t\le10L\), \(A+t+\eta\le v\).
For the first bound use \(\sigma\le3L\) and
\(A\le (h+\eta)/(p-1)+1\), giving
\[
 (\eta+A)/p\le 1.1(h+\eta)/p^2+(\eta+1)/p.
\]

Substitution in the floor bound gives \(t\le10L\).
For the second, \(A\le (b_0+2\eta+S)/(p-1)+1\),
\(10L\le b_0/16\) and \(20\eta,20S\le b_0\).
If all the indicated coefficients were divisible by \(p^t\),
put \(k=t-2\sigma\). Modulo \(p^k\) one has
\(q^{A+k-1}\in((q-p)^A)\) by the binomial expansion.
Cancel \(p^{2\sigma}\), remove the initial power of \(q\),
and use this divisibility. The result would be
\[
 P_1^y\in(p^k,q^{v-\eta-(A+t-1)}).
\]
The second content calculation, before rounding, gives
\[
 k\le\max\{0,\ y-(v-\eta-A-t+1)/p+2+\log_p y\}.
\]
Indeed using the index \(\min(y,\lfloor(C-1)/p\rfloor)\)
costs at most \(y-C/p+1+1/p+\log_p y\) when positive.

Since \(py\le v\) and \(p/(p-1)\le11/10\), the definition
of \(t\) contradicts the last display. Choose a coefficient
of valuation less than \(t\), and multiply its source vector by
the integral quotient of \(p^t\) by that coefficient.
Moving one step forward gives the positive-level preimage.
The sole pole of \(f\) is at this level, so the second lifting
rule gives cost
\[
 t+\sigma\le13L.
\]


Suppose next that \(v<p(u+S+\eta)\). The down-level constants
in \eqref{eq:short-scalar} can be different. Eliminate the one
at level two first. The denominator there is \(P_1^v\), the
modulus is \(P_2^{u+S}\). Put \(y=u+S\). Since
\(P_2=p\varepsilon_1+U_1P_1^p\) with unit coefficients, the
first content calculation (using \(P_1,P_2\) in place of
\(q,P_1\)) gives a splitting
\[
 t_2f=x_2+f_2,\quad D f_2\in P_2^{y}\mathcal R,\quad
 D(p^{k_0}f_2-x_0)\in q^{z+S}\mathcal R,\qquad D=P_1^v .
\]

If \(v<py\), the content of \(P_1^v\) modulo \(P_2^y\)
is zero by reduction and comparison of orders. Otherwise use
the ceiling multiple \(p\lceil v/p\rceil\) in its place.
For \(py\le v<p(y+\eta)\), the expansion in \(p,P_1^p\) of \(P_2^y\) costs
at most \(\eta+2+\gamma_h\), and the small letter belongs to
the band. Thus \(v_p(t_2)\le(\eta+2+\gamma_h)B\).

Write \(C=z+S\). For \(x_0\ne0\), the second congruence gives
\begin{equation}\label{eq:short-last-constant}
 v_p(x_0)-k_0\ge y-v+C/p-C/p^2-2-\gamma_h .
\end{equation}
Indeed \(P_2^y\) has content at least \(y-C/p^2\) modulo
\(q^C\). Expand first in \(p,P_1^p\), then in \(p,q^p\):
a term with less scalar valuation has \(q\)-degree at least
\(C\). The content of \(P_1^v\) modulo \(q^C\) is at most
\(v-C/p+2+\gamma_h\) by the second calculation, since
\(C\le pv\). Comparing contents of
\[
 x_0P_1^v\equiv p^{k_0}P_2^y U\bmod q^C
\]
proves the claim.

Here is how to use this scalar inequality with the middle
filtered lattice. Transport the truncated horizontal column to
the middle point, paying \(\sigma\) at the first small letter,
and denote its integral representative by \(\lambda\). It
represents \(p^\sigma d\) there.

Set \(s=P_1^y/q^v\) at this
point. The vector \(s\lambda\) lies in the target Wach lattice.
The level test uses \(d\in\mathcal F^v\) at the middle and
the zero of order \(y\) at the initial negative letter.
In addition
\[
 \gamma(s)(\gamma-1)\lambda\in Xq^{a_*}N_{c-1}(T_2)
\]
where \(a_*=A_{c-1}+r_{c-1}\) and \(r_i\) is the least
scalar Wach image exponent at \(i\). Indeed the one-step
horizontal correction has defect in \(X^M q^{2G}\) in the
first Frobenius span. The next transport gives factors
\(X^M q^{(p-1)M}P_1^{2G}\); after multiplication by
\(\gamma(s)\) these pay \(Xq^{a_*}\) and the signed height losses
\(q^\eta P_1^{u+S}\), since \((p-1)M\ge v+a_*+\eta\).

The powers required in the ordinary parameter and the scalar
give the \(q^{a_*}\) image by the binomial theorem.
The two-filter argument of Paragraph~\ref{par:thick-filtered-comparison}
therefore lifts \(p^E s\lambda\) integrally, where
\[
 E=\left\lfloor\frac{3a_*+S+2\eta-2}{p-1}
               +\log_p(v+\eta-1)\right\rfloor_{\!+}\le8L .
\]
For this application the lower filter has indices \([v-a_*+1,v]\)
and the other has indices in \([v+1,v+S+a_*]\cup[1-\eta,\eta+a_*]\).
The factors adjacent to \(j\) in the first two intervals
give factorials of combined length \(2a_*+S-1\).

In the remaining
interval of length \(2\eta+a_*\), isolate a term of largest
valuation; the others are bounded by factorials of combined
length \(2\eta+a_*-1\). The isolated valuation is at most
\(\log_p(v+\eta-1)\). If \(a_*=0\) the map is already onto
at this letter. The global ordinary width bound \(H\le h+\eta\)
and degree lowering give
\[
 a_*\le H(2p-1)/(p-1)^2+\lceil1+\log_p(1+H)\rceil+1.
\]

These imply \(2a_*+\eta+1\le v\) under \(40L\le b_0/4\)
and give the displayed estimate for \(E\), by \(p\ge11\) and
\(H\le b_0+2\eta+S\). In particular
\[
 v+a_*+\eta\le2v\le2h<(p-1)M,
\]
as used in the defect estimate. The separation gives the
disjointness required by interpolation. \label{cor:inert-short-three-1}

Let \(n_1\) be this lift. Take its constant in a horizontal
basis at \(c-1\) and transport it to \(c\) with precision
\(Xq^C\). This gives \(y_1\) in the last Frobenius span,
with \((\gamma-1)y_1\in XN_c(T_1)\) and \(q^C y_1\in N_c(T_1)\).
Use the truncation order
\[
 R_0=\lceil C/(p-1)\rceil
\]
and
take \(\sigma_1=R_0\). In positive degree \(r<R_0\) the
denominator costs \(r+v_p(r!)\). The transported monomial is
a unit times \(X^r q^{(p-1)r}\). Keep one \(X\) and expand
the others as \((q-p)^{r-1}\).

Modulo \(Xq^C\) the surviving
terms contain at least \(\max(0,pr-C)\) powers of \(p\):
a term of extra degree \(j\) must satisfy \(j+(p-1)r<C\).
The loss is therefore at most
\[
 pr/(p-1)-\max(0,pr-C)\le C/(p-1);
\]
the maximum occurs at the crossing \(pr=C\). The transported
tail is zero to this precision. The last-level jet of \(y_1\)
is horizontal in the last Frobenius span: its difference is
zero modulo \(Xq^C\) there, by the truncation and the integral
correction. Over the coefficient field use the local logarithmic
parameter at the primitive level. On a horizontal column the
action on degree \(k\) has eigenvalue \(\chi(\gamma)^k\), so
an invariant jet has only its constant coefficient. This gives
the jet of \(F(y_1)\) from the transported constant of \(n_1\)
with multiplier \(p^{\sigma_1}\).

To compare with \(f_2\Phi\lambda\), use the common horizontal
column at the last level. At this level the transported multiplier
\(\phi(s)=P_2^y/P_1^v\) has constant term
\[
 (P_2^y/P_1^v)(0)=p^{y-v},
\]
since \(P_j(0)=p\) for \(j\ge1\) in the \(q\)-coordinate.
Thus the two scalar constants to be matched, after multiplying
by \(p^{k_0+v}\), are \(x_0p^v\) and
\(p^{k_0+\sigma_1+E+y}\). In the second expression the
common horizontal normalization cancels; its \(p^\sigma\) is
present in both comparisons.

Choose \(t_3,B'\in O\) with
\[
 t_3 x_0p^v=B' p^{k_0+\sigma_1+E+y},\qquad v_p(t_3)\le14L .
\]
For \(x_0=0\) use \((1,0)\). Otherwise take \(t_3=1\) if
\(v_p(x_0)+v\ge k_0+\sigma_1+E+y+\sigma\). In the other
case choose equality with \(B'=p^\sigma\). Equation
\eqref{eq:short-last-constant} bounds \(v_p(t_3)\) by
\[
 \sigma+E+\sigma_1-C/p+C/p^2+2+\gamma_h\le14L .
\]
In fact \(\sigma_1\le C/(p-1)+1\), \(C\le h\), and
\(1/(p(p-1))+1/p^2\le 2.1/p^2\).

For the marked map \(F\), the required Frobenius-span congruence is
\[
 t_3 f_2\Phi\lambda-B' F(y_1)\in q^C\langle\Phi N_{c-1}(T_2)\rangle.
\]

For the first term the numerator at level two supplies the
height loss from the first negative letter; equivalently
\(f_2\Phi\lambda\) is an integral multiple of
\(\Phi(s\lambda)\). For the divisibility assertion compare
the two horizontal jets after inverting \(p\) and the higher
level factors. Modulo \(q^C\) the first has constant multiplier
\(x_0/p^{k_0}\) on \(\Phi\lambda\), by the last congruence;
the horizontal jet of the second has constant multiplier
\(p^{\sigma_1+E+y-v}\). Higher level factors are units at
\(q=0\) over the fraction field, and the bounded integral
Frobenius span is free.

For the free bounded integral Frobenius span \(P\), monic division yields
\[
 P\cap q^C P[1/p]=q^C P.
\]
Thus the congruence of rational horizontal jets is integral. The negative terminal
height puts this difference in \(N_c(T_2)\). Lift \(p^R\)
times the difference there and add \(p^R B' y_1\).
The result is admissible and maps to \(p^Rt_3 f_2\Phi\lambda\);
its denominators in \(N_c(T_1)\) are cleared by \(q^C\).
Since \(\Phi\lambda\) represents \(p^\sigma d\), this
lifts \(t_3f_2\) at cost \(R+\sigma\). The summand \(t_3x_2\)
admits the constant lift at the same cost. The scalar clearing cost is
\[
 \lceil v_p(t_3t_2)\rceil=\lceil v_p(t_3)+v_p(t_2)\rceil.
\]
Adding the lifting costs gives the bound \(22L+1+\eta B\).
\end{proof}


\begin{lemma}[Three-letter lifting with positive outer signs]\label{lem:three-letter-positive-lift}
For the same \(L=S_g\) and short-band indicator \(B\) as in
Lemma~\ref{lem:three-letter-negative-lift}, and for chronological signs \(+,-,+\), the adjacent-level correction can be chosen with exponent \(e_3\) satisfying
\[
 e_3\le34L+1+\eta B.
\] With the other short-arc branches, the
single exponent \(E=\lfloor40S_g+\eta B\rfloor\) works on
every eligible column of the chosen horizontal frames,
in powers of \(\ell\) in the original DVR.
\end{lemma}

\begin{proof}
\proofstep{Use the middle congruence with the two positive neighbor denominators.}
Now consider \(+,-,+\). The denominator of \(f\) is
\(D=q^zP_2^u\) and its middle congruence is modulo
\(P_1^{v+S}\). Write \(C=v+S\), \(a_0=\lceil z/p\rceil\).
We first show that the middle rational constant \(x/p^k\)
has denominator at most \(p^s\), with
\begin{equation}\label{eq:three-denominator}
 s\le\max\{1,u-C/p+16L/5\}+1 .
\end{equation}
Put \(Q=u-(C-a_0+1)/p+2+\log_p u\), and use
\(s=1+\lfloor\max(0,Qp/(p-1))\rfloor\).

If \(D\in(p^s,P_1^C)\), multiply by \(q^{pa_0-z}\).
Expanding \(P_1=p\varepsilon+Uq^p\) gives
\(P_1^{a_0+s-1}\in(q^{pa_0},p^s)\).
Hence \(P_1^{a_0+s-1}P_2^u\in(p^s,P_1^C)\).
Cancelling this power of \(P_1\) modulo \(p^s\), and descending
by \(\phi\), gives
\[
 P_1^u\in(p^s,q^{C-a_0-s+1}).
\]

Cancellation holds by reduction to a distinguished polynomial,
and descent holds since the ring is free over its Frobenius
image (Weierstrass division by \(\phi(q)-Y\)).
The content inequality from the preceding paragraph would force
\[
 s\le\max(0,u-(C-a_0-s+1)/p+2+\log_p u),
\]
contrary to the choice. Here \(C-a_0-s+1>0\):
\(a_0\le z/p+1\), \(20u\le b_0\) and the bound on \(s\)
in \eqref{eq:three-denominator} gives \(a_0+s<C+1\) using
\(160L\le b_0\).

To verify that bound directly, use
\((a_0-1)/p\le h/p^2\), \(u\le w\) and \(\log_p u\le\gamma_h\).
Expanding the factor \(p/(p-1)=1+1/(p-1)\) charges at most
\((w+h/p^2+2+\gamma_h)/(p-1)\) in addition to
\(u-C/p+h/p^2+2+\gamma_h\), which fits the displayed allowance.
This proves the bound on the content of \(D\); the middle
congruence then gives
\[
 p^sx/p^k\in O.
\]


Put \(A=A_{c-2}\). If \(u<A+\eta+2\sigma+2\), separate the
constant at cost \(s\) and lift both summands at \(R+\sigma\).
Indeed the new remainder has numerator divisible by \(P_1^C\).
We check that the size-dependent loss in
\eqref{eq:three-denominator} fits the band allowance.
By \eqref{eq:short-A},
\[
 u-C/p<\eta+\eta/p+\tfrac52 L+2\sigma+2.
\]
Inside the band this gives \(s\le\eta+14L+1\). Outside it,
if \(pu\le b_0\) the excess \(u-C/p\) is nonpositive.
If \(pu>b_0\), absence of the band gives
\(u-b_0/p>2\eta\), while the last upper bound and
\((C-b_0)/p\le(\eta+S)/p\) imply \(\eta<12L\).
These bounds give
\[
 s\le 28L+1+\eta B
\]
on every arc.

If \(u\ge A+\eta+2\sigma+2\), construct a filtered constant
at the initial positive letter. The one-letter calculation with
the coordinate \(q^{-u}d\) gives scalar comparison
\(p^{2\sigma}(q-p)^A q^{G-u}\) modulo \(q^{G-\eta}\).
The degree \(G-u+A\) is in the comparison range. Its leading coefficient is
\[
 [q^{G-u+A}]\,p^{2\sigma}(q-p)^A q^{G-u}=p^{2\sigma}.
\] Equivalently, some coefficient
in that range divides \(p^{2\sigma+1}\): expanding
\(q^{A+2\sigma}\) in \((q-p)+p\) shows that divisibility of all
the low coefficients by \(p^{2\sigma+1}\) is impossible.
All these source coefficients lie in \(\mathcal F^u\). Thus
after transport a source direction at positive level two maps
to \(p^{2\sigma+1}d\). Only its filtration at level two has
been prescribed. We choose the part of \(f\) to which it applies.

Let
\[
 j=\max(0,\lfloor u-C/p-\eta/p^2\rfloor).
\]
There is a
decomposition
\[
 P_2^u=p^j U_0+P_1^C q^\eta U_1,\qquad U_0,U_1\in\mathcal R .
\]
In the expansion of \((p\varepsilon_1+U'_1P_1^p)^u\), the
term with \(P_1^{pr}\) is paid by its scalar coefficient if
\(r\le C/p+\eta/p^2\). Otherwise factor off \(P_1^C\) and
expand the remaining \(pr-C\) powers of \(P_1\).
Any resulting term with \(q\)-degree below \(\eta\) has
scalar valuation at least
\(u-r+(pr-C)-\eta/p\ge u-C/p-\eta/p^2\).
This proves the decomposition. Put \(e_4=(s-j)_+\) and
\(x'=p^s x/p^k\).

Take
\[
 f_a= p^{e_4+j-s}x' U_0/P_2^u.
\]
Then \(p^sf_a-p^{e_4}x'\) belongs to
\(P_1^Cq^\eta P_2^{-u}\mathcal R\).
The second rule lifts \(f_a\) at cost \(3\sigma+1\), using
the source filtration at level two and the constant congruence
at the omitted positive level zero. The remainder
\(p^{e_4}f-f_a\) has numerator divisible by \(P_1^C\) over
\(D\), by the middle congruence and saturation. It costs
\(R+\sigma\). Since \(j\ge u-C/p-\eta/p^2-1\),
\eqref{eq:three-denominator} gives
\[
 e_4\le 16L/5+\eta/p^2+2
\]
when the second term in the maximum is used; otherwise
\(e_4\le2\). In particular \(R+3\sigma+1+e_4\le18L+1\).
The two initial-positive cases therefore both fit
\(34L+1+\eta B\). \label{cor:inert-short-positive-1}

All the bounds for the individual scalar lifts are at most
\(40S_g+\eta B\), since \(S_g\ge1\).
They hold on every eligible column of every chosen horizontal
frame; consequently the integer
\(E=\lfloor40S_g+\eta B\rfloor\) can be used simultaneously.
The exponents are in powers of \(\ell\) in the original DVR,
including when its uniformizer has fractional normalized valuation.
\end{proof}


\begin{lemma}[A fixed exceptional span for a short arc]\label{lem:short-arc-exceptional-span}
For the short arc of length \(a\) in the cut comparison,
one may choose an exceptional set \(Z_c\) of target classes, independent of the input, with
\[
 \#Z_c\le(a+1)ma.
\] Outside its span the scaled lift has
exponent \(\sigma+E+\kappa_c\). The signed count is
\(r_c=a\) because each letter has an off-central window.
\end{lemma}

\begin{proof}
\label{par:short-exceptions}
On the marked target over the coefficient field let \(U\) be the
intersection of the transported \(\mathcal F_j^{u_j}\) over
positive levels. Let \(V_-\) be the sum of
\(\mathcal F_j^{-\eta}\) over negative levels. By the bound on
central multiplicities,
\[
 \operatorname{codim}U+\dim V_-\le ma .
\]
In the initial integral basis, the constants form the lattice
\(D_O\) with the transport of scalars to the terminal crystalline
space understood. Split the saturated sublattice \(D_O\cap U\)
off with complement \(K\), and project \(V_-\) along \(K\)
to \(U\). The intersection of the projected space with the
lattice of \(U\) is again saturated. Use adapted bases of these
two splittings.

Let \(e_j\) denote the adapted basis and let \(J\) index the complement \(K\) and this last intersection. Then
\[
 |J|\le ma,\qquad V_-\subseteq\operatorname{span}_{K_O}\{e_j:j\in J\}.
\]
Every other column belongs to \(U\).
The two changes of basis are integral and invertible, so they
give a Wach frame at \(c'\). Take the unique horizontal columns
with these initial values.

Apply the cut construction of Section~\ref{sec:inert-long-orbit}
to a crystalline extension representing a finite target class.
Its cut vector is the difference in that extension of a vector
in the incoming Frobenius span and a terminal Wach vector.
At each cyclotomic level the incoming span is regular; a pole
of order \(k>0\) therefore has coefficient in the extension's
\(\mathcal F^k\), and its image lies in the kernel of projection
to the trivial representation. Strictness identifies this
intersection with \(\mathcal F^k\) of the target.

For a polar term \(a_{-k}t^{-k}\) with \(k>0\), strictness therefore gives the two respective tests
\[
 a_{-k}\in V_-\quad(\text{negative letter}),\qquad k\le u_j+S\quad(\text{positive letter}).
\]
The statement for an entire polar part is obtained
coefficient by coefficient in the Laurent expansion, followed
by Chinese remaindering on the finite jets. At negative levels
all negative degrees thus disappear on indices outside \(J\).
Likewise \((\gamma-1)w\in XN_c\) imposes the Wach filtration
on the divided differences. Their projection outside \(J\)
vanishes at a negative level in degrees below \(u_j\):
for lower degrees the required filtration is contained in
\(\mathcal F_j^{-\eta}\), or equals it by the gap.

Write \(\mathcal J_c\) for the union of signed degree intervals
used in \(\kappa_c\).
Apply to these horizontal coordinates the integral operator
\[
 K(\gamma)=\left(\sum_{j=0}^{p^a-1}\gamma^j\right)
     \prod_{k\in\mathcal J_c\setminus\{0\}}
            (\gamma^{p^a}-\chi(\gamma)^{kp^a}) ,
\]
On a Laurent coefficient of
degree \(k\) at level \(t<a\), the principal action raised to
\(p^a\) is the scalar \(\chi(\gamma)^{kp^a}\). The sum
projects degree zero to a common value on the primitive roots
after tame invariants. Outside \(J\) this removes the excess
positive-level poles and the negative-level degrees from \(u_j\)
through \(u_j+S-1\).

At a negative level the divided difference condition on a nonconstant coefficient \(a_k\) below \(u_j\) reads
\[
 (\chi(\gamma)^{kp^a}-1)a_k=0,\qquad k\ne0.
\]
The scalar is nonzero, hence \(a_k=0\).
Thus the transformed scalar coefficients outside \(J\) satisfy
\eqref{eq:short-scalar}. The remaining constant at each negative
level is in the coefficient field, so a scalar power clears it.
The operators preserve the bounded integral denominator
\(\Lambda\); the improved pole orders descend to \(D\) by
division by the distinguished factors. Outside \(J\), the only poles allowed by \(D\) are positive with directions in \(U\); the negative zeros have the required orders. Write \(f=f_J+f_{J^c}\) for the support decomposition and \(\Delta\) for divided difference. The Wach test is linear, so
\[
 \Delta f_J=\Delta f-\Delta f_{J^c}
\]
passes as well.

For the latter part let \(M_J\) consist of the tuples with
denominator \(\Lambda\) supported on \(J\); let \(I_J\) be
the submodule whose divided differences pass the Wach test.
Let \(P_J\) consist of \(X\) times the tuples passing the Wach test
with bounded denominator, in \(M_J\). For numerators supported on \(J\),
a change by \(X\Lambda^2h\) induces the following change of tuple:
\[
 \frac{f+X\Lambda^2 h}{\Lambda}-\frac{f}{\Lambda}
       =X\Lambda h\in P_J .
\]
Weierstrass division reduces each numerator modulo \(X\Lambda^2\);
hence \(M_J/P_J\) is finite over \(O\).

The quotient \(M_J/P_J\) is also torsion-free.
Suppose \(0\ne c\in O\) and \(c(f/\Lambda)\in P_J\). Monic division by
\(X\) gives \(cf(0)=0\), hence \(f(0)=0\). The tuple
\(f/(X\Lambda)\) now has bounded denominator. Its rational filtration
tests hold since \(c f/(X\Lambda)\) satisfies them and the filtration
spaces are vector spaces over the coefficient field. Thus
\[
 f/\Lambda\in P_J.
\]


Let \(E_J\) be the image of \(I_J\) in \(M_J/P_J\).
After inverting \(p\), Chinese remainders for \(X\) and the
distinct cyclotomic factors of \(\Lambda\) split the numerator
jets at \(0\) and at the \(a\) levels. At \(0\) the quotient by
\(X\) leaves at most one entry per coordinate in \(J\).
At a cyclotomic level \(X\) is invertible, so the quotient is
by the filtered Laurent lattice.

In coefficient degree \(k\),
the principal power fixing the roots acts on that quotient by
\(\chi(\gamma)^{k p^a}\). A nonzero fixed class therefore needs
\(k=0\), since \(\chi(\gamma)\) has infinite order. In degree
zero the tame and principal root actions are transitive. Invariants
have at most one copy of each supported coordinate, and passing
to the quotient cannot increase this invariant dimension by
the characteristic-zero averaging projection. Thus
\[
    \dim_{\operatorname{Frac}(O)}(E_J[1/p])
          \le |J|+\sum_{t=0}^{a-1}|J|=(a+1)|J| .
\]
Since \(M_J/P_J\) is finite torsion-free over \(O\), this is
also a bound on the number of generators of \(E_J\).

Take an \(O\)-basis of the image and choose its representatives
in \(I_J\) once for the cut. Realize \(p^\sigma\) times each
by the horizontal comparison and take its cut Herr class. The
resulting set \(Z_c\) consists of finite classes, since a
product of cyclotomic factors carries the representative into
the Wach lattice. Using one representative per basis vector of \(E_J\) gives the uniform bound
\[
 |Z_c|\le\operatorname{rank}_O E_J\le (a+1)ma.
\]

\proofstep{Reduce every class against the same chosen representatives.}
For any input the exceptional coordinates are an integral
combination of these representatives modulo \(P_J\).
Terms of \(P_J\) represent elements of \(XN_c\) after the
same comparison and give zero under the cut class map.
The other coordinates lift at the common cost \(E\). To see the
exponent accounting, the inverse horizontal comparison reads
coordinates of \(p^\sigma w\). Realizing the split coordinates
requires another \(p^\sigma\), while on the outside part the
vector lifting relation already represents \(p^E\) times the
coordinates.

Since \(E\ge\sigma\), pad the realized split parts
by \(p^{E-\sigma}\). Their sum now matches
\[
 p^{\sigma+E} K(\gamma)w
\]
modulo \(XN_c\). Naturality and the vanishing on
\(XN_c\) follow from Frobenius contraction in the Herr complex,
as in the cut construction. Write \(K(T)-K(1)=(T-1)Q(T)\). Stability and the cut-vector condition give
\[
 K(\gamma)w-K(1)w=Q(\gamma)(\gamma-1)w\in XN_c.
\]
Thus the two vectors give the same cut class. In the scalar reading
\[
        p^{\sigma+E}K(1)w
             = u_c p^{\sigma+E+\kappa_c}w,\qquad u_c\in O^\times
\]
since \(v_p(K(1))=\kappa_c\). This gives the exponent
\(\sigma+E+\kappa_c\) with the same set \(Z_c\) for every
input class. For the present signed arcs every letter has a
positive or negative off-central window, so the signed level count
\(r_c=a\). \label{cor:inert-short-exceptions-1}
\end{proof}



\paragraph{Adding the short-arc colengths.}
\begin{lemma}[The common colength from the short cuts]
\label{lem:ordinary-inert-short-cuts}
In either ordinary inert range with a position having smaller jumps, use its
marked arcs. The normalized finite-class colength is bounded by
\eqref{eq:nonthick-colength}; in the range requiring exceptional representatives, there is one set \(Z\) for all the comparisons such that
\[
 \#Z\le4nm.
\]
\end{lemma}

\begin{proof}
Apply Lemma~\ref{lem:short-arc-exceptional-span} with \(p=\ell\) and \(a=a_c\).
On the finite jets at levels \(t<a_c\), the action of \(\gamma^{\ell^{a_c}}\)
on a Laurent coefficient \(c_k\) of degree \(k\) is multiplication by
\(\chi(\gamma)^{k\ell^{a_c}}\). For each specified nonzero degree
the corresponding factor of \(K(\gamma)\) therefore satisfies
\[
 \bigl(\gamma^{\ell^{a_c}}-\chi(\gamma)^{k\ell^{a_c}}\bigr)c_k=0 .
\]

To estimate \(\kappa_c\), count these degrees with their signs. For a positive
jump window these degrees lie in \([-s_i-S,-s_i)\) and for a
negative window in \([s_i,s_i+S)\). Thus there are at most
\(a_c S\) degrees, each nonzero with absolute value at most \(h\).
In any one of these intervals of \(S\) integers the count of
multiples of \(\ell^j\) is at most \(S/\ell^j+1\).
For such an interval \(I\) excluding zero the resulting valuation bound is
\[
 \sum_{k\in I}v_\ell(k)=\sum_{j=1}^{\lfloor\log_\ell h\rfloor}\#\{k\in I:\ell^j\mid k\}\le S/(\ell-1)+\log_\ell h.
\]

The principal trace contributes at most \(a_c\), and the
nonzero degrees \(k\) contribute at most \(a_c+1+v_\ell(k)\) each,
by the principal-unit valuation used in
Section~\ref{sec:inert-long-orbit}. Hence
\[
 \kappa_c\le a_c+(a_c+1)a_c S
                +a_c(S/(\ell-1)+\log_\ell h)
            \le13 S+3(1+\gamma_h).
\]
Given the cut decomposition and the arc lift, pad the powers to
their common maximum and sum the lifted cut classes.
The set \(Z\) of exceptional classes
is contained in the union of the sets from the cuts, so
\[
 |Z|\le\sum_c m r_c(a_c+1)\le4m\sum_c a_c=4nm .
\]
The cost of the scaling on the original class is at most
\(55S_g+13S+\eta B\), since \(1+\gamma_h\le S_g\).
Thus the length argument of Section~\ref{sec:inert-thick-word}
applies to the same integral finite group.

In the larger range
one uses the full lift and takes \(Z\) empty. In both ranges the
normalized colength is at most
\begin{equation}\label{eq:nonthick-colength}
 320(n+1)r\bigl(h/\ell^2+(w+\eta)/\ell+
                      S+1+\gamma_h+\eta B\bigr)+320nm e' .
\end{equation}
Indeed \(S/\ell\le S\), all the terms in the parentheses
are nonnegative, and \(55+13\le320\). Lemma~\ref{lem:short-arc-exceptional-span}
supplies the fixed representatives for this estimate. The coarse exponent
\(e'\) is for the actual marked surjection and is used only on the
image of the span of the exceptional finite classes.
\end{proof}

\begin{proposition}[One normalized entry for all ordinary inert rows]\label{prop:ordinary-inert-place-entry}
In the ordinary inert range, with the marked word data and
\(U_H,U_F,U_M\) defined in Section~\ref{sec:cut-length-rows},
there is one rational place entry \(c_w\) for every good character at the place. Its range is
\[
 0\le c_w-(368U_H+2082U_F+22365U_M)\le U_F.
\]
Here the indicator entering \(U_F\) is the common place indicator,
and every normalized length uses the same coefficient valuation.
\end{proposition}

\begin{proof}
The common unramified coarse finite-class lifting theorem uses the
negative and positive bounds \(M_i,H_i^+\) at each embedding.
Put \(M_*=\max_i M_i\), \(H'=\max_i(M_i+H_i^+)\)
and \(a_*=\lfloor M_*/(\ell-1)\rfloor+1\).
If \(J\) is the set of nonzero degrees formed by adding
\(1,\ldots,M_*\) to jumps and also including \(1,\ldots,M_*\),
the exponent for this theorem is
\[
  2\left\lfloor\frac{H'\ell}{(\ell-1)^2}+3+
                        \log_\ell(H'+1)\right\rfloor
       + a_*+v_\ell(a_*!) +\sum_{j\in J}v_\ell(j).
\]
The theorem supplies a finite lift for every target class after
multiplication by this power of \(\ell\).

In the general boxes \(J\subset[-H',2H']\setminus\{0\}\).
Legendre's formula bounds the valuation sum by \(3H'/(\ell-1)\).
Also
\(a_*+v_\ell(a_*!)\le(\ell/(\ell-1))(H'/(\ell-1)+1)\).
Using \(\ell/(\ell-1)\le11/10\) gives
\[
 e'\le23\bigl(2+H'/\ell+\log_\ell(H'+1)\bigr).
\]
For the large-\(S_g\) boxes above, substitute \(J\subset[-h,2h]\),
\(M_*\le h\), and \(H'\le h+\eta\le21h/20\) into the same
expression. Adding the coefficients of \(h/\ell\) from the three terms gives
\[
 2(21/20)(11/10)^2+(11/10)^2+33/10<7.1.
\]

The constants other than the logarithm total at most \(7.1\),
and \(H'+1\le2h\) since \(H'\) is integral and \(h\ge1\).
This proves the exponent estimate used in that range.

Use the global \(H,w_P,S_c,\eta_3,z\) and the quantities
\(U_H,U_F,U_M,\mathcal P_\ell\) of
Section~\ref{sec:cut-length-rows}. The local transport gives
\(S\le S_c\), \(\eta\le\eta_3\), \(w\le w_P\), and \(h\le H\)
in both branches. In the case in which every position has larger jumps this last inequality
means \(\max_i b_i+S\le H\). The polynomial size bounds give
\(1\le H\le2e^z\), \(z\ge1\) and \(\eta_3\le H\).
When some jumps are in the smaller range, \(H'\le2H\). In the
case in which all jumps are in the larger range
\[
 H'\le H_{\rm loc}\le H+\eta_3\le2H.
\]

As \(H'+1\le3H\le\ell H\) for the logarithm's argument, it follows
that \(e'\le69\mathcal P_\ell\).

For \eqref{eq:nonthick-colength}, \(\gamma_h\le3+z\) follows from
\(2h\le4e^z\) and \(\log\ell\ge1\).
Consequently \(1+\gamma_h\le5z\). The row length is at most
\(320U_H+1600U_F+22080U_M\).

For Proposition~\ref{prop:thick-inert-colength},
\(L_0\le8+z\le9z\), since
\(2H_{\rm loc}\le8e^z\). Further,
\[
 32H_{\rm loc}/\ell^2\le32H/\ell^2+3\eta_3/\ell,\qquad
 10\eta/(\ell-1)\le11\eta_3/\ell .
\]
The exponent \(e\) is therefore at most
\(32H/\ell^2+225F_{\ell,\lambda}\).
Also \(H_{\rm loc}+1\le3H\), which gives
\(c_D\le4\mathcal P_\ell\).
Substituting in \eqref{eq:thick-chain-mass} bounds the colength
by \(32U_H+227U_F+288U_M\).
Both branches thus have common majorant
\[
                       368U_H+2082U_F+22365U_M .
\]
The place entry is a rational number between this value and its
sum with \(U_F\). Here \(U_F>0\) by \(z\ge1\) and \(r>0\).
All comparisons use the same place indicator and coefficient
normalization. Hence the entry bounds the row for every
good character at once.
\end{proof}



\paragraph{Primes above the cutoff outside the exceptional windows.}
\label{subsec:outside-windows}
\hypertarget{source-large-prime-length}{}
\begin{established}[Signed jumps outside the exceptional windows]\label{hyp:offwindow-local}
Fix the same family and good point, and let \(\ell>\ell_*\) be outside
the exceptional set and different from the two type primes.
At each Frobenius letter the shifted rank-\(r\) jumps and scalar string lie in \([1,h_i-1]\). For the end interval of size \(\eta'_i\) and carry \(\sigma_i\) computed below one has
\[
  \sigma_i\le\lfloor 2b/\ell_*\rfloor,\qquad \eta'_i+4\sigma_i+6<\ell .
\]
Every rank-\(r\) jump is at distance greater than \(\sigma_i+2\) from each scalar jump. For a jump outside the end interval its distance from the chosen end is incongruent modulo \(\ell\) to every integer in \([-2\sigma_i-2,\eta'_i+2\sigma_i+3]\).
\end{established}
\begin{proof}
Once the family scale is sufficiently large, every such prime satisfies
\[
 \ell>t+1,\qquad \ell\ge5,\qquad \ell\nmid Q_Sqd.
\]
The estimates here use the jump conditions and the norm choice on
the selected cut. We describe the conversion of the exceptional-set
test to a local cohomology condition. The integral classification below uses the residue models of the
extensions and their specified jumps.

At a place over \(\ell\), each embedding of the maximal unramified
subfield has a unique extension to the local field, since the field
itself is unramified. Order these embeddings by Frobenius, cyclically.
At a letter let \(F_l,F_c\) be the finite sets of integral central and
outer jumps of the rank-\(r\) representation, before shifting the
first scalar character. Put \(p=(t-1)/2\),
\[
 B=b/\ell_*.
\]

The weight boxes give, with a center \(W_0\) from the exceptional list,
\[
 \begin{gathered}
 2B+t/2+3\le |j-p|\le\eta_c \quad(j\in F_l),\\
 W_0\le \pm(j-p)\le W_0+S_c+\eta_c+1 \quad(j\in F_c).
 \end{gathered}
\]
The sign is fixed at this letter. Here
\(W_0\ge 2B+t/2+3\), \(t\le\eta_c\), and
\(2\eta_c+8B+7<\ell\) on a common threshold. Indeed the
sizes \(B,\eta_c,t\) are at most powers of the
\(U\) in Section~\ref{subsec:exceptional-range}, and
that section has
\[
 U^j=o(\omega),o(\ell_*)
\]
for every fixed
\(j\). The central lower separation is a jump-box condition.
The reflected set at the opposite embedding uses \(j\mapsto t-1-j\).

Let \(l_0=\min(F_l\cup\{0\})\),
\(l_1=\max(F_l\cup\{t-1\})\), \(\eta'=l_1-l_0+1\).
Shift all jumps by \(w=1-l_0\) in the positive case
and by \(w=1-\min(F_l\cup F_c\cup\{0\})\) in the negative case.

In the positive case take
\[
 h=\max(F_l\cup F_c\cup\{t-1\})+w+1
\]
and in the negative case take \(h=l_1+w+1\). Then the
shifted rank-\(r\) jumps \(U_i=(F_l\cup F_c)+w\) and the scalar
block \([w,w+t-1]\) lie in \([1,h-1]\). The scalar block and
shifted central jumps lie in \([1,\eta']\) in the positive case
and \([h-\eta',h-1]\) in the negative case. Refer to this
interval as the end interval. The union records the rank-\(r\)
multiplicities since each retained rank-\(r\) jump is simple.
The shift is integral even when some of the input windows have
nonintegral endpoints.

For letter \(i\), abbreviate these quantities by \(w_i,h_i,\eta'_i\)
and form the carry bound
\[
 \sigma_i=\left\lfloor\sum_{j\ge1}\frac{h_{i+j}}{\ell^j}\right\rfloor .
\]
The two formulas for \(h_i\), the boxes and \(W_0\le b\) on the
selected family (the centers other than \(b\) are \(o(b)\))
give \(h_i\le b+S_c+2\eta_c+3\). Summing the
geometric series gives \(\sigma_i\le 2B\), hence
\(\sigma_i\le s_+=\lfloor 2B\rfloor\). For this bound use
\[
 (b+S_c+2\eta_c+3)/(\ell_*-1)\le 2b/\ell_*,
\]
valid
eventually by the preceding sizes. Also
\(\eta'_i+4\sigma_i+6<\ell\). The distance from a rank-\(r\)
jump to each scalar jump is greater than \(\sigma_i+2\), by
subtracting the half-width \((t-1)/2\) from its distance to \(p\).

There is a separate residue test for a shifted jump outside the
end interval. Such a jump comes from \(j\in F_c\). Its distance
from the chosen end, measured from 0 or \(h_i\), is
\(d=j+1-l_0\) on the positive side and \(d=l_1+1-j\) on the
negative side. If \(-2\sigma_i-2\le v\le\eta'_i+2\sigma_i+3\),
then in both cases
\[
 W_0-\eta_c-2s_+-3\le d-v\le
                    W_0+S_c+2\eta_c+2s_++4 .
\]

For the lower bound insert \(\eta'=l_1-l_0+1\) before estimating.
This cancels \(l_0\) on the positive side and \(l_1\) on the
negative side. The inequalities \(p-\eta_c\le l_0\) and
\(l_1\le p+\eta_c\) then suffice. The upper bound uses the
opposite endpoint for \(v\). A congruence \(d\equiv v\pmod\ell\)
would make \(\ell\) exceptional by its defining window.
We have proved every letter's end-interval size, separation, and
residue-avoidance conditions at the given \(\ell\).
\label{cor:large-prime-length-1}
\end{proof}

\paragraph{Putting the residue extensions on one model.}
In the residue calculation write \(k'\) for the residue field and
use \(X\) for the characteristic-\(\ell\) cyclotomic parameter
throughout this calculation. Let \(D_i\) be the corresponding
étale \((\varphi,\Gamma)\)-module at letter \(i\), with Frobenius
\(\Phi_i\) coming from letter \(i-1\) and a principal cyclotomic
generator \(\gamma_i\) acting on letter \(i\).
Write \(\Phi_i N_{i-1}\) for the \(k'[[X]]\)-span of the
Frobenius images of a lattice. The convention for the residue
models after untwisting by \(\chi_1\) is
\[
 X^{w_i^*}N_i\subseteq \Phi_i N_{i-1}
       \subseteq X^{-w_i}N_i,\qquad
 (\gamma_i-1)N_i\subseteq XN_i,\qquad w_i^*=h_i-w_i.
\]
The model of the untwisted trivial quotient is the standard
rank-one lattice in its invariant basis. Its uniqueness among
the admissible models follows from the scalar comparison below.

\paragraph{Reduction of the marked Wach lattices.}
\label{par:residue-wach-realization}
Use the field, coefficient extension and reduced cyclotomic
realization of Section~\ref{subsec:cut-wach-models} with
\(p=\ell\), \(u=1+p\), and \(X\) the reduction of the integral
parameter \(X_0=-Z\). The residue field is \(k'\).
All crystalline and residual factors use the terminal Hodge
labels of that section.
The residual parameter actions, with Frobenius from the preceding letter, are
\[
 \varphi(X)=X^p,\qquad \gamma(X)\equiv X-X^2\pmod{X^3}.
\]


The first scalar character \(\chi_1\) has jump \(w_i\);
the scalar characters are \(\chi_c=\chi_1\epsilon^{1-c}\)
for \(1\le c\le t\), with \(\epsilon\) the cyclotomic character.
A full representation is crystalline with the labeled jumps
of the marked rank block together with those of the scalar
characters. Its residual characteristic polynomials agree with
those of that direct sum on the Galois group.

For a stable lattice
in such a full sum, take the Wach lattice after twisting by
\(\chi_1^{-1}\). With \(U_i\) the multiset of rank-block jumps before the twist, the full twisted multiset is
\[
 B_i=(U_i-w_i)\sqcup\{0,\ldots,t-1\}.
\]
Both parts retain their multiplicities. The separate rank and line lattices
are treated by the same twist. Their labeled height bounds,
after reduction by the coefficient uniformizer, give
\[
      X^{w_i^*}N_i\ \subset\ \Phi_i N_{i-1}
                    \ \subset\ X^{-w_i}N_i,
              \qquad (\gamma_i-1)N_i\subset XN_i .
\]
Here \(N_i\) is the reduction of the tame-invariant Wach lattice
inside the residual 'etale module on the corresponding letter.
A scalar line of sole jump \(e\) has Frobenius ideal \(X^e\).
For the trivial quotient line the invariant basis is \(1\).
\cite{berger-wach}\label{cor:residue-wach-realization-1}

\paragraph{The polynomial on a high-end letter.}
\begin{lemma}[Simple zero factor on a high-end letter]
\label{lem:residue-high-end-polynomial}
For the upper-end letter and \(s=w_i^*+1\), write \(D\) for the divided action on \(N_i/X^sN_i\). For \(P\) as in the proof the minimal polynomial satisfies
\[
 m_D(T)\mid TP(T),\qquad P(0)\ne0,\qquad \ker(D^j)=\ker(D)\quad(j\ge1).
\] Modulo \(X\)
the characteristic polynomial still records every jump residue
with its multiplicity.
\end{lemma}

\begin{proof}
In the integral tame parameter \(X_0\), with reduction \(X\),
put \(q=X_0+p\) and \(\nu=(\gamma-1)/X_0\). This is
\(O'\)-linear and preserves the invariant Wach lattice, since
\(X_0\) generates the invariant augmentation ideal and the
action is trivial on constants. For \(s\ge1\) define
\[
 I_s=\{a-b:\ 0\le a<s,\ b\in B_i\},\qquad
 d(m)=\frac{1-(1+p)^m}{p}\in\mathbb Z_p .
\]
On the characteristic-zero quotient by \(q^s\), the operator
is killed by
\[
 \prod_{m\in I_s}(T-d(m)),
\]
using each distinct
integer once. At the primitive level choose crystalline directions
adapted to the filtration. In a direction of jump \(b\) the
quotient has Laurent degrees \(-b,\ldots,s-1-b\) in the
local logarithmic parameter. The scalar substitution multiplies
this parameter by \(1+p\), and \(X_0=-p\) at the root.

In a direction of jump \(b\) the triangular operator \(T\) has distinct diagonal entries \(d(a-b)\), hence
\[
 \prod_{a=0}^{s-1}\big(T-d(a-b)\big)=0
 \quad\text{on this direction}.
\]
The product over the union of the eigenvalues thus kills every direction, including repeated jumps.
This is the finite-level filtration comparison of
Paragraph~\ref{par:wach-level-models}. The quotient of the
free invariant Wach lattice by \(q^s\) is free over \(O'\)
by monic division. The polynomial identity therefore holds
integrally. For \(s=1\) the characteristic polynomial has each
\(d(-b)\) with its multiplicity in \(B_i\), by the same
calculation and injectivity on coefficients.
\label{cor:residue-wach-realization-2}

Reduction gives \(q=X\), \(\bar d(m)=-m\), and operator
\[
 \nabla=(\gamma_i-1)/X.
\]
For negative \(m\) the congruence
for \(d(m)\) follows by inverting \((1+p)^{-m}\) modulo \(p^2\)
before division. The characteristic polynomial modulo \(X\)
therefore counts the residues of \(B_i\) with multiplicity.
At the upper end take \(s=w_i^*+1\). If \(0\le a<s\) and
\(a\equiv b\bmod p\), then \(a=b\). For a scalar or central
jump, writing \(b=u-w_i\), both \(u\) and \(w_i+a\) are in
the interval of width \(\eta'_i\) next to \(h_i\), including
the endpoint \(h_i\) for the latter.

For an outer jump the
congruence would give \(h_i-u\equiv w_i^*-a\bmod p\) in the
forbidden end window. It follows that
\[
   P(T)=\prod_{m\in I_s\setminus\{0\}}(T+m)\in k'[T],
   \qquad P(0)\ne0,\qquad
   \nabla P(\nabla)N_i\subset X^sN_i .
\]
If zero is absent from \(I_s\), \(P(\nabla)\) already has
image in \(X^sN_i\); applying \(\nabla\) gives the
displayed identity since the divided action preserves that
lattice. In the finite quotient, \(P(\nabla)\) is invertible
on the generalized zero space, since \(P(0)\ne0\). Thus
\(\nabla P(\nabla)=0\) there forces \(\nabla=0\). Using
the characteristic polynomial with all multiplicities would
lose this conclusion. \label{cor:residue-wach-realization-3}
\end{proof}

Finally an admissible residual representation is a quotient
of the reduction of its presenting full lattice. Its presenting
map is onto and kills \(\lambda T\), and untwisting preserves
that map. More generally take the intersection with a vector
subspace for a subrepresentation and the image for a quotient.
The intersection lattice is saturated. In a basis adapted to it,
Frobenius base change is exact, since the series ring is free
over its Frobenius image on \(1,X,\ldots,X^{p-1}\). Intersecting
the height inclusions or taking images gives both inclusions
for the induced lattices. The divided cyclotomic differences
intertwine the maps, so their polynomials restrict and descend:
for a stable vector subspace \(W\) one has
\[
 (X^sN_i)\cap W=X^s(N_i\cap W).
\]
The induced lattice sequence stays
exact modulo \(X\). The integral Wach intersection gives the
direct sum for a sum of lattices.

For the standard reference
use the split sum of the chosen rank lattice and the character
lines. Its reductions give the required models on those
summands, with their original labels. The prescribed residual
characteristic polynomials fix the composition factors of
each full sum, and untwisting multiplies them all by the
same character. This supplies the models used below with
their two heights, polynomials and target Sen counts.
\label{cor:residue-wach-realization-4}


\paragraph{A common lattice for the extension classes.}
\begin{lemma}[A largest compatible residue model]
\label{lem:largest-compatible-residue-model}
For two models \(N^0,N\) of the same residual representation,
each component is bounded by \(N_i\subseteq X^{-\sigma_i}N_i^0\). For the fixed residual representation the largest model is
\[
 N^{\max}=\sum_{N\ {\rm admissible}}N;
\]
the sum is itself admissible.
\end{lemma}

\begin{proof}
If \(N^0,N\) are two models of the same representation, let \(a_i\)
be the least nonnegative \(a\) with \(X^a N_i\subseteq N_i^0\).
The height inclusions give
\[
 \Phi_i N_{i-1}\subseteq X^{-a_i-w_i-w_i^*}\Phi_i N^0_{i-1}.
\]
In a Frobenius basis, the images of the coefficients have orders
multiplied by \(\ell\). Hence this inclusion implies
\(\ell a_{i-1}\le a_i+h_i\).
After \(m\) iterations this reads
\[
        a_i\le \ell^{-m}a_{i+m}
                 +\sum_{j=1}^m h_{i+j}/\ell^j .
\]
The first term tends to zero because the \(a_i\) form a
finite cyclic list. This gives \(a_i\le\sum_{j\ge1}h_{i+j}/\ell^j\).
Thus every model lies in \(X^{-\sigma_i}N_i^0\) in each letter.
The sum of two models is a model.

The sum model arises by direct sum
followed by the addition map of the common target. Exactness
and stability of the admissible subquotients give the induced
image lattice. The finite product of the containing lattices
is Noetherian over \(k'[[X]]\), so the sum of all models is
already a finite sum and is a largest model \(N^{\max}\).
This argument applies with the standard model as \(N^0\).
\label{cor:residue-model-length-1}
\end{proof}

In a residue extension, use the preimage lattice of the target
under the inclusion into the extension's model. Choose at each
letter a vector \(e_i\) mapping to \(1\) in the standard quotient.
The Herr pair on the target model is
\[
 b_i=\Phi_i e_{i-1}-e_i,\qquad g_i=(\gamma_i-1)e_i.
\]
Write \(P_i=\langle\Phi_i N_{i-1}\rangle\) for the target
Frobenius span. The cocycle equation and the two height tests
give
\[
 \begin{aligned}
 (\gamma_i-1)b_i&=\Phi_i g_{i-1}-g_i, &g_i&\in XN_i,\\
 X^{w_i}b_i&\in N_i &&\text{at a lower end},\\
 X^{w_i^*}b_i&\in P_i &&\text{at an upper end}.
 \end{aligned}
\]
Indeed \(\gamma_i\) acts trivially modulo \(X\) on the extension
lattice and \(N_i\) is the kernel there of the trivial quotient.
At the upper end apply the height inclusion to \(X^{w_i^*}e_i\).
Its image in the trivial line is \(X^{w_i^*}\), the coefficient
of \(\Phi_i e_{i-1}\) in the Frobenius span. Subtracting this
term leaves a vector in \(P_i\), by exactness on the induced
spans.

For the extra test at the upper end set \(s=w_i^*+1\).
Use the polynomial of Lemma~\ref{lem:residue-high-end-polynomial}
on the extension model, and put
\[
    e_i'=P(0)^{-1}P(\nabla)e_i,\qquad x_i=e_i-e_i'.
\]
The quotient map intertwines \(\nabla\) with zero, so \(e_i'\)
has the same quotient coordinate as \(e_i\), and \(x_i\in N_i\).
The identity \(\nabla P(\nabla)e_i\in X^s N_i^{\rm ext}\)
therefore yields
\[
     g_i-(\gamma_i-1)x_i=(\gamma_i-1)e_i'
                          \in X^{w_i^*+2}N_i .
\]
Here \(N_i^{\rm ext}\) denotes the extension lattice, and the
last containment uses its intersection with the target. Passing
to \(N^{\max}\) preserves the tested inclusions.
The exact functor on residue modules preserves addition of
extension classes and split extensions. A linear combination of
their Herr pairs that is a boundary on the common model thus
gives a coboundary for the original Galois classes.
\label{cor:residue-model-length-2}

\begin{auxlemma}[Sen multiplicities at the two ends]\label{lem:end-sen-multiplicities}
Use the residue positions and exclusions of Lemma~\ref{hyp:offwindow-local}.
Let \(s_i(y)\) be the generalized multiplicity of the residue Sen
operator at \(y\) on \(N_i^{\max}/XN_i^{\max}\) for the model of
Lemma~\ref{lem:largest-compatible-residue-model}.
Put \(d_i=\#\{u\in U_i:u<w_i\}\), with the labelled jumps counted
with multiplicity. At the indicated end letters one has
\[
 \begin{array}{ll}
 \displaystyle\sum_{y=-w_i}^{-1}s_i(y)=d_i &\text{at a lower end},\\[4pt]
 \displaystyle\sum_{y=0}^{w_i^*}s_i(y)=r-d_i &\text{at an upper end}.
 \end{array}
\]
\end{auxlemma}
\begin{proof}
Write \(N=N^{\max}\), and let \(N^0\) be the standard target
model. At letter \(i\), \(\nabla\) induces endomorphisms on
\(N_i/XN_i\) and \(N_i^0/XN_i^0\), with generalized
eigenvalue multiplicities \(s_i(y)\), \(s_i^0(y)\) at
\(y\in k'\). These are the Sen multiplicities in the
chosen normalization. For the standard target,
\(s_i^0(y)\) counts \(u\in U_i\) with
\[
 u-w_i\equiv y.
\]

The residue model comparison gives the same support condition
for the Sen weights of any target model, without a prescribed
multiplicity. Here is the support argument, which also works
for a subquotient of the marked rank block. The model induced on
that subquotient by \(N^0\) has Sen eigenvalues among \(u-w_i\)
for \(u\in U_i\). Indeed the subobject lattice is saturated and
the image lattice in the quotient is free. Exactness modulo \(X\)
and stability under \(\nabla\) give the weight inclusion.
For two models \(L,L'\) of that subquotient the carry bound
in both directions places \(L'/XL'\) as a subquotient of
\(X^{-\sigma_i}L/X^{\sigma_i+1}L\).

In particular each eigenvalue on \(L'/XL'\) lies in the set
\[
 \{b+a:b\text{ an eigenvalue on }L/XL,\ a\in\mathbb Z,\ |a|\le\sigma_i\}
\]
in the residue field.
On the other hand, any admissible model is induced from a
full crystalline sum, so its Sen weights are in the union of
the shifted rank-block jumps and the scalar residues
\(0,\ldots,t-1\). If an eigenvalue belongs to the scalar
block, the comparison with the induced model of \(N^0\)
would put a rank-block residue within \(\sigma_i\) of
that block. This is exactly the forbidden congruence from
the position tests. Hence all the model's Sen eigenvalues
come from the rank-block residues.
\label{cor:residue-model-length-3}

Let \(q_i(y)\) count generalized eigenvalues on \(N_i/N_i^0\).
On \(X^a N_i^0/X^{a+1}N_i^0\) the eigenvalues shift by \(-a\),
since the scalar action is normalized by
\[
 \gamma(X)=X-X^2\bmod X^3.
\]

In particular \(XN_i/XN_i^0\) has multiplicity \(q_i(y+1)\)
at \(y\). To compute the multiplicities in \(H=N_i/XN_i^0\), compare its
two filtrations. Their short exact sequences are
\[
 \begin{aligned}
 0&\longrightarrow XN_i/XN_i^0\longrightarrow H
        \longrightarrow N_i/XN_i\longrightarrow0,\\
 0&\longrightarrow N_i^0/XN_i^0\longrightarrow H
        \longrightarrow N_i/N_i^0\longrightarrow0 .
 \end{aligned}
\]
Adding generalized multiplicities along the two rows gives
\[
 s_i(y)+q_i(y+1)=s_i^0(y)+q_i(y).
\]
The filtration of \(X^{-\sigma_i}N_i^0/N_i^0\) shows that
if \(q_i(y)\ne0\), then \(y\equiv u-w_i+j\) for
some \(u\in U_i\) and \(1\le j\le\sigma_i\).
This step counts generalized eigenspaces, which have additive
dimensions in exact sequences even for nonsemisimple operators.

Let \(d_i\) count \(u\in U_i\) with \(u<w_i\).
The position tests give \(s_i(y)=0\) for
\[
 0\le y\le\sigma_i+1,
\]
and \(q_i(0)=0\).
Indeed a residue in these ranges would require \(u\)
congruent to \(w_i+y\) or \(w_i-j\) with distance
at most \(\sigma_i+2\) from the scalar block after checking
the end interval. For an outer jump the corresponding
congruence is precisely forbidden by the end-window test.

For a lower end letter sum the displayed identity for
\(-w_i\le y<0\). At the lower endpoint \(q_i(-w_i)=0\),
since any residue would require \(u\equiv -j\) with
\(1\le j\le\sigma_i\). For a central jump \(1\le u\le\eta'_i\), the end-width bound of
Lemma~\ref{hyp:offwindow-local} gives
\[
  0<u+j\le\eta'_i+\sigma_i<\ell ,
\]
which excludes \(u+j\equiv0\pmod\ell\). For an outer jump it contradicts
the residue exclusion. In the interval \([-w_i,0)\), the
standard multiplicities count exactly the central jumps below
\(w_i\). An outer jump cannot contribute, since its residue
would lie in the excluded lower end window. The endpoint vanishings \(q_i(0)=q_i(-w_i)=0\)
eliminate the correction in the telescoping identity:
\[
 \sum_{y=-w_i}^{-1}s_i(y)=
   \sum_{y=-w_i}^{-1}s_i^0(y)+q_i(-w_i)-q_i(0)=d_i.
\]

For an upper end letter sum from \(0\) through
\(w_i^*+\sigma_i+1\). At the upper endpoint after telescoping,
\(q_i(w_i^*+\sigma_i+2)=0\): a contributing \(u\) would
be congruent to \(h_i+\sigma_i+2-j\), \(1\le j\le\sigma_i\).
For an outer jump this is precisely the residue exclusion.

At this end a central jump satisfies \(h_i-\eta'_i\le u\le h_i-1\), so
\[
 0<h_i+\sigma_i+2-j-u\le\eta'_i+\sigma_i+1<\ell .
\]
This excludes the congruence also for a central jump. The same test gives
\(s_i(y)=0\) for \(w_i^*<y\le w_i^*+\sigma_i+1\).
Throughout the summation range the standard multiplicities
count exactly the jumps \(u\ge w_i\), all of which are in
the upper end interval. Put \(B=w_i^*+\sigma_i+1\). The vanishing of
\(s_i(y)\) for \(w_i^*<y\le B\) gives
\[
 \sum_{y=0}^{w_i^*}s_i(y)=\sum_{y=0}^{B}s_i(y)
     =\sum_{y=0}^{B}s_i^0(y)+q_i(0)-q_i(B+1)=r-d_i .
\]
All the interval counts include multiplicities and use bounds
strictly shorter than \(\ell\) at the relevant end. The number
\(d_f\) in the local comparison is \(\sum_i d_i\).
\label{cor:residue-model-length-4}
\end{proof}

\paragraph{From the polar parts to the residue rank.}
The filtration of \(N_i/X^{\sigma_i+2}N_i\) has successive
eigenvalues shifted by \(0,-1,\ldots,-(\sigma_i+1)\).
The Sen vanishing in the proof of Lemma~\ref{lem:end-sen-multiplicities}
makes \(\nabla\) invertible on this quotient.
For each pair choose \(y_i\in N_i\) so that
\[
 g'_i=g_i-(\gamma_i-1)y_i\in X^{\sigma_i+3}N_i,\qquad
 b'_i=b_i-\Phi_i y_{i-1}+y_i .
\]
This is a change by a Herr boundary.
Set \(c_i=b_i-\Phi_i y_{i-1}+x_i\), using \(x_i=y_i\)
at a lower end letter and the deeper cyclotomic correction
at an upper end letter.

The corrected vector \(c_i\) has the same image as \(b'_i\)
modulo \(N_i\).
The relation \(\ell(\sigma_{i-1}+1)\ge h_i+\sigma_i+1\)
follows by shifting the carry series. Consequently
\(\Phi_i g'_{i-1}\in X^{w_i^*+2}N_i\).
The height inclusions and the cocycle equation give
\[
\begin{array}{ll}
c_i\in X^{-w_i}N_i,&(\gamma_i-1)c_i\in XN_i
       \quad\text{at a lower end letter},\\
c_i\in X^{-w_i^*}\Phi_i N_{i-1},&
(\gamma_i-1)c_i\in X^{w_i^*+2}N_i
       \quad\text{at an upper end letter}.
\end{array}
\]
If a linear combination of the \(c_i\) is integral at
every letter, the combination of the \(b'_i\) is integral.
Its cyclotomic difference belongs to \(X^{\sigma_i+3}N_i\),
using the deeper \(g'_i\) and the same carry inequality.
Invertibility modulo \(X^{\sigma_i+2}\) places \(b'_i\) in
\(X^{\sigma_i+2}N_i\).

On the product of these deep lattices
\(\Phi\) strictly increases the order at each step by the
height and carry inequalities. Solve \((\Phi-1)v=b'\) there
by \(-\sum_{a\ge0}\Phi^a b'\). Commutation with \(\gamma\)
and injectivity of \(\Phi-1\) on the deep product give
\[
 (\gamma-1)v=g'.
\]
Thus the original combination is a
Herr boundary. In particular the dimension of a space of
independent pairs is at most the sum of the bounds for the
images of their \(c_i\)'s modulo \(N_i\).

At a lower end letter the possible images lie in the kernel
of \(\nabla\) on \(X^{-w_i}N_i/N_i\).
On its graded piece of degree \(a\) the generalized zero
multiplicity is \(s_i(a)\). The sum for \(-w_i\le a<0\)
is \(d_i\), so the kernel has dimension at most \(d_i\).

At an upper end letter put \(M=\Phi_i N_{i-1}\),
\[
 V=X^{-w_i^*}M
\]
and \(J=X^{w_i^*+1}N_i\).
Then \(J\subseteq N_i\subseteq V\), and \(V/J\) is finite
dimensional over \(k'\). The required image is that of
\[
 \ker(\nabla:V/J\longrightarrow D_i/J)
\]
in \(D_i/N_i\).
The codimension estimate for \(\nabla\) on \(V/J\)
uses the cyclotomic Frobenius left inverse \(\psi_i\).
In the range \(1\le w_i^*\le\ell-7\), its scalar companion
maps \(X^{1-w_i^*}k'[[X]]\) into \(k'[[X]]\), and
\(\psi_i\Phi_i=1\). On the part of \(X^{1-w_i^*}M\) killed
by \(\psi_i\), the operator \(\gamma_i-1\) can be inverted
in \(X^{-w_i^*}M\).

\begin{lemma}[Frobenius coefficient and divided-action inverse]
\label{lem:residue-frobenius-left-inverse}
For an upper-end letter, set \(d=w_i^*\) with \(1\le d\le\ell-7\)
and \(M=\langle\Phi_i N_{i-1}\rangle\) in the residue model of
Lemma~\ref{lem:largest-compatible-residue-model}. Write \(\Psi_i\) for the Frobenius left inverse and \(\Psi\) for its scalar counterpart. Then
\[
 \Psi_i(X^{1-d}M)\subseteq N_{i-1},\qquad\Psi(X^{1-d}k'[[X]])\subseteq k'[[X]].
\]
On the kernel in \(X^{1-d}M\) the \(\gamma_i-1\) equation
has a solution in \(X^{-d}M\) in the kernel.
\end{lemma}

\begin{proof}
Extend scalars temporarily to
\(k'((\pi))\), with \(X\) the invariant parameter of
Paragraph~\ref{par:residue-wach-realization}.
Write \(p=\ell,\ d=w_i^*\) and \(A=1+\pi\).
Use the extended lattices \(\widetilde N=k'[[\pi]]N_{i-1}\)
and \(\widetilde M=k'[[\pi]]M\). A basis of \(\widetilde N\)
from \(N_{i-1}\) gives a Frobenius basis of \(\widetilde M\).

In this basis every vector in the ambient space has a unique
expansion
\[
             \sum_{j=0}^{p-1} A^j\Phi_i(y_j).
\]
Indeed \(1,\pi,\ldots,\pi^{p-1}\) form the usual basis over
the image under \(\pi\mapsto\pi^p\). The triangular change
to \(1,A,\ldots,A^{p-1}\) is integral with unit diagonal.
Define \(\psi_i\) to return \(y_0\). This definition is
independent of the Frobenius basis, by uniqueness of the
direct sum expansion. A tame substitution of exponent \(z\)
sends the \(j\)-summand to the summand with index \(zj\bmod p\),
since
\[
 A^{zj}=A^{j'}\phi(A^{(zj-j')/p}).
\]
Thus the
projection descends to the tame invariants, commutes with
the principal action and satisfies \(\psi_i\Phi_i=1\).
The same expansion in rank one gives the scalar companion.

If the input lies in \(X^{1-d}M\), its ordinary order in
\(\widetilde M\) is at least \((p-1)(1-d)\ge p(1-d)\).
The integral change of basis gives
\(y_j\in\pi^{1-d}\widetilde N\) for every \(j\).
The zeroth component is tame-invariant. Its nonzero
coefficients have orders divisible by \(p-1\) in the
invariant basis. Since \(1-d>-(p-1)\), it belongs to
\(N_{i-1}\). This proves in particular the scalar bound.

For an input with \(y_0=0\), solve separately for \(1\le j<p\).
On the \(j\)-summand the difference is
\[
   (\gamma_i-1)\bigl(A^j\Phi_i(z)\bigr)
      = A^j\Phi_i\bigl((A^j\gamma_{i-1}-1)z\bigr),
\]
since \(\gamma(A)=A^{1+p}\). The operator on the right is
multiplication by \(A^j-1\), of ordinary order one, plus
\(A^j(\gamma_{i-1}-1)\). On each
\(\pi^h\widetilde N\), \(h\in\mathbb Z\), the latter term
increases order by at least \(p-1\): use
\(\gamma(\pi)/\pi\equiv1\bmod\pi^{p-1}\) and
\((\gamma_{i-1}-1)N_{i-1}\subset XN_{i-1}\).
After dividing by \(A^j-1\), the error raises order by at
least \(p-2>0\). Put \(B_j=(A^j-1)^{-1}A^j(\gamma_{i-1}-1)\).

The operator \(B_j\) increases order by at least \(p-2\), so the solution is
\[
  z_j=\sum_{k\ge0}(-B_j)^k(A^j-1)^{-1}y_j
                       \ \in\ \pi^{-d}\widetilde N .
\]
Multiplication by \(A^j-1\) after applying \(1+B_j\)
gives \((A^j\gamma_{i-1}-1)z_j=y_j\).
The resulting sum has ordinary order at least \(-pd\) in
\(\widetilde M\). Tame averaging preserves that order, the
equation and the vanishing under \(\psi_i\). For the averaged
vector, the lower order bound rounds up to \(-(p-1)d\), since
\(d<p-1\) and the coefficients in the invariant basis have
orders in \((p-1)\mathbb Z\). This gives the preimage in
\(X^{-d}M\) with the required one-power loss.
\label{cor:residue-frobenius-inverse-1}
\end{proof}



By Lemma~\ref{lem:residue-frobenius-left-inverse}, for \(u\in V\) apply \(\psi_i\) to \(Xu\), obtaining
\(y\in N_{i-1}\). Subtract \(\Phi_i y\) from \(Xu\)
and use the inverse on the \(\psi_i=0\) part.
Divide by \(X\) to express \(u\) as
\(X^{-1}\Phi_i y+\nabla v\), with \(v\in V\).
In a basis of \(N_{i-1}\), write \(y=y_0+Xy_1\)
with \(y_0\) in the span of the constant basis vectors.
The preceding letter's Sen vanishing solves
\[
 \nabla z\equiv y_1\bmod X^{\sigma_{i-1}+2}N_{i-1}.
\]

Then \(X^{-1}\Phi_i(Xy_1)\) differs from
\(\nabla\Phi_i z\) by an element of \(J\), by commutation
and the carry inequality. It follows that \(V\) is covered
by \(\nabla V+J\) and the span of the \(r\) vectors
\(X^{-1}\Phi_i\) of the preceding basis.
Rank-nullity gives dimension at most \(r\) to
\(\ker(\nabla:V/J\to D_i/J)\), since \(V/J\) lies in the
sum of that \(r\)-dimensional contribution and the image
of this map.

The zero generalized eigenspace on \(N_i/J\) has dimension
\[
 \sum_{a=0}^{w_i^*}s_i(a)=r-d_i
\]
by its graded pieces.
The polynomial with \(P(0)\ne0\) gives
\(\nabla P(\nabla)=0\) on \(N_i/J\). Hence on the generalized
zero eigenspace \(\nabla=0\), since \(P(\nabla)\) is invertible
there. These \(r-d_i\) directions of the kernel have zero
image modulo \(N_i\). The image has dimension at most \(d_i\).
This proves the rank bound \(\sum_i d_i=d_f\) on the
common residue model.
\label{cor:residue-model-length-5}

\paragraph{The residue invariants and the scalar lines.}
\begin{proposition}[The scalar lines on the residue model]
\label{prop:offwindow-scalar-residue-model}
In the residue model of Lemma~\ref{lem:largest-compatible-residue-model},
with heights and exclusions of Lemma~\ref{hyp:offwindow-local}, none of \(C_e=\bar\epsilon^{-e}\) for
\(-1\le e<t\) occurs as a subquotient of the marked rank block.
The marked rank block therefore has zero local cohomology in
degrees zero and two. For \(0\le e<t\), the induced lattice on a scalar
subquotient \(C_e\) is its standard model. If \(A_e\) is the space of admissible residue extension classes with trivial quotient and target \(C_e\), then
\[
 \dim A_0\le1,\qquad A_e=0\quad(1\le e<t).
\]
\end{proposition}

\begin{proof}
\proofstep{The carry estimate and the excluded scalar residues.}
Let \(\bar\epsilon\) be the residue cyclotomic character and
write \(C_e\) for the rank-one representation
\(\bar\epsilon^{-e}\). The standard model \(L_i^{(e)}\) has
Sen eigenvalue \(e\) and its Frobenius span is \(X^eL_i^{(e)}\).
For \(-1\le e<t\) these lattices satisfy the two height
inclusions. Throughout this range the endpoints satisfy
\[
 w_i\ge1,\qquad w_i+e\le h_i-1.
\]
The standard models can therefore
serve as comparison lattices whether or not the line
is a subquotient under study.

Suppose \(C_e\) were a subquotient of the marked rank block
for such an \(e\). The induced lattice on that subquotient
has a Sen eigenvalue congruent to \(e+a\) with
\(|a|\le\sigma_i\), by the two-direction carry comparison
with the standard line. Its eigenvalue comes from the
Sen weights of \(N_i^0\) by exactness.
The residue exclusion near the scalar block forbids
\[
 u\equiv w_i+e+a
\]
for \(u\in U_i\), since this point
is within \(\sigma_i+2\) of the scalar interval.
For \(e=0\) this rules out invariants. For \(e=-1\)
it rules out a quotient of cyclotomic character,
and local Tate duality gives vanishing in degree two.
\label{cor:residue-model-length-6}

\proofstep{Uniqueness of the induced scalar lattice.}
For \(0\le e<t\) the standard line is the unique
model of \(C_e\) induced from the full crystalline sums.
To see this, refine the subquotient presentation of
the line to a composition series of the full residue
representation. Use the induced lattices at each step.
At the level of semisimplifications, one full copy has the same factors as
\[
 \bar W^{\rm ss}\ \oplus\ \bigoplus_{e=0}^{t-1} C_e,\qquad C_e=\bar\epsilon^{-e},
\]
where \(\bar W\) denotes the marked residual rank block after twisting by \(\bar\chi_1^{-1}\). Direct sums repeat this list with multiplicity.
The preceding exclusion keeps the two lists disjoint, and
the scalar characters are distinct since \(t<\ell-1\).
The Sen weights of a factor from the rank block are among
its shifted jump residues by the support argument above.

A lattice on a scalar factor \(C_a\) is \(X^{\kappa_i}L_i^{(a)}\)
for an integer \(\kappa_i\). The two-direction bound gives
\[
 |\kappa_i|\le\sigma_i.
\]
Its height conditions read
\[
 -w_i+\kappa_i\le \ell\kappa_{i-1}+a
                      \le w_i^*+\kappa_i .
\]
The end-width bound of Lemma~\ref{hyp:offwindow-local}
gives \(\eta'_i+\sigma_i<\ell\). Use the left height inequality
at a lower end, and the right one at an upper end, to obtain
\begin{equation}\label{eq:scalar-end-signs}
 \begin{array}{ll}
 \ell\kappa_{i-1}\ge -(w_i+a)+\kappa_i
          \ge-\eta'_i-\sigma_i>-\ell
     &\text{at a lower end},\\
 \ell\kappa_{i-1}\le (w_i^*-a)+\kappa_i
          \le\eta'_i+\sigma_i<\ell
     &\text{at an upper end}.
 \end{array}
\end{equation}

At a fixed letter \(i\), index the scalar composition factors by
\(a_\nu\), with repetitions. Every shift \(\kappa_{i,\nu}\)
is therefore nonnegative if the next letter is a lower end and
nonpositive if it is an upper end. All shifted weights
\(a_\nu-\kappa_{i,\nu}\) belong to
\([-\sigma_i,t-1+\sigma_i]\). This interval has length less than
\(\ell\), contains the nominal scalar weights and excludes the
rank-block weights. Additivity of generalized Sen multiplicities
therefore gives equality of integer multisets in this interval.
In particular
\[
 \sum_\nu (a_\nu-\kappa_{i,\nu})=\sum_\nu a_\nu
 \quad\Longrightarrow\quad\sum_\nu \kappa_{i,\nu}=0 .
\]
The sign from \eqref{eq:scalar-end-signs} forces each shift to vanish. This also applies to the line in
the original subquotient presentation.
\label{cor:residue-model-length-7}

\proofstep{Removing pole terms in an admissible Herr pair.}
Now take an admissible extension with target \(C_e\),
\(0\le e<t\), and trivial quotient. Its Herr pair \(b_i,g_i\)
is on these standard lattices. We have \(g_i\in XL_i^{(e)}\).
At a lower end letter \(b_i\in X^{-w_i}L_i^{(e)}\), and
at an upper end letter
\[
 b_i\in X^{e-w_i^*}L_i^{(e)}.
\]

In either case the permitted pole degrees \(j>0\)
satisfy \(e+j<\ell\), by the short end width:
\(w_i+e\le\eta'_i\) on the lower side and
\(w_i^*\le\eta'_i\) on the upper side.
The cocycle equation implies \(\nabla b_i\in L_i^{(e)}\),
since Frobenius takes \(XL_{i-1}^{(e)}\) to
\(X^{\ell+e}L_i^{(e)}\). On the negative degree \(-j\)
the leading eigenvalue of \(\nabla\) is \(e+j\ne0\).
Inductively eliminating the negative coefficients shows
\[
 b_i\in L_i^{(e)}.
\]
For \(e>0\), \(\Phi\) is contracting
on \(\prod_iL_i^{(e)}\). Inverting \(\Phi-1\) there kills
the pair by the same cochain argument as above.

\proofstep{The constant class on the cycle.}
For \(e=0\), reduce \(b_i\) modulo \(X\) to a constant
\(c_i\) in the invariant basis. Invert \(\Phi-1\)
on \(\prod_i XL_i^{(0)}\) for \(b_i-c_i\).
The cocycle equation again identifies the other
coordinate as the cyclotomic difference of the resulting
vector, since \(\gamma_i\) fixes the \(c_i\).
Let \(q\) be the number of letters and \(k\) the residue
field. The last coordinate space is a subquotient of
\[
 k^q/\{(d_{i-1}-d_i)_i : d\in k^q\}.
\]
On the cycle every difference vector has zero sum. Conversely
if \(\sum_i c_i=0\), start with \(d_0=0\) and solve
\(d_i=d_{i-1}-c_i\) successively. The last equation closes
because the sum is zero. Thus summation identifies the displayed
quotient with \(k\). This calculation uses no division by \(q\)
and gives a dimension bound of one in every residue characteristic.
\label{cor:residue-model-length-8}
\end{proof}


\begin{proposition}[Vanishing of the local rank-complement loss above the cutoff]\label{prop:offwindow-local-vanishing}
Let \(\ell>\ell_*\) lie outside the exceptional set and the
two type primes, at a sufficiently large point satisfying the signed windows and residue
tests of Lemma~\ref{hyp:offwindow-local}. At each \(w\mid\ell\) let \(D_w\) be the image of the marked rank-\(r\) classes modulo their finite local condition. Then
\[
 \operatorname{len}D_w=0\qquad(w\mid\ell).
\]
\end{proposition}

\begin{proof}
The required local assertion will follow by reducing the admissible extension classes to the residue field, then lifting through their successive coefficient quotients. Work over a finite unramified extension \(F\) of
\(\mathbb Q_\ell\) with \(\ell\ge5\), in a splitting coefficient
DVR \(O'\) with uniformizer \(\lambda\), whose fraction field contains
the images of the local field and the first cyclotomic level. Take a crystalline
representation with lattice \(L\) and jumps \(U_i\), and
integral crystalline characters
\[
 \chi_c=\chi_1\epsilon^{1-c}
\]
for \(1\le c\le t\) with jumps \(w_i+c-1\), where
\(\epsilon\) is the cyclotomic character. Assume the position
conditions just proved. Call a representation full if it is crystalline, its residual characteristic polynomials are the reductions for the rank-\(r\) representation plus the scalar characters, and its labeled jump multisets are
\[
 U_i\sqcup\{w_i+c-1:1\le c\le t\}.
\]

\proofstep{Use the residue models in the signed jump windows of this completion.}
For a finite coefficient module \(X\), consider a cocycle with
target \(L\otimes X\) twisted by \(\chi_1^{-1}\), or with the
analogous character target. Form the extension with quotient
line of character \(\chi_1\). Suppose all sufficiently deep
reductions of the quotient line, at depths killing \(X\), give
quotients of nested stable lattices in finite direct sums of full
representations. Then the rank-\(r\) cocycle comes from integral
finite cohomology with coefficients in \(X\), up to coboundary.
Writing \(\Lambda\) for the untwisted rank-\(r\) target lattice, the asserted containment is
\[
 [c]\in\operatorname{im}\bigl(H^1_f(F,\Lambda)\otimes X\longrightarrow H^1(F,\Lambda\otimes X)\bigr).
\]

For characters \(c\ge2\)
the cocycle is a coboundary, and for \(c=1\) its restriction to
inertia is a coboundary.  The residue bound for the character targets is Proposition~\ref{prop:offwindow-scalar-residue-model}. For the rank-\(r\) target the end multiplicities of Lemma~\ref{lem:end-sen-multiplicities} give the dimension bound below. The following coefficient argument converts these bounds into assertions about the integral classes.
\label{cor:large-prime-length-2}

Let \(\Lambda\) denote the lattice after removing the quotient
character \(\chi_1\), and put \(k'=O'/\lambda O'\).
Write
\[
 A(X)\subset H^1(F,\Lambda\otimes X)
\]
for the classes
satisfying the finite-quotient presentation just stated. At a
common depth \(b\) killing the coefficient modules, suppose
\[
   0\longrightarrow\Lambda\otimes X
    \longrightarrow E_X\longrightarrow O'/\lambda^b
    \longrightarrow0
\]
is the extension after untwisting. For \(f:X\to X'\) the pushout along \(1\otimes f\) is
\[
 E_{X'}=\big(E_X\oplus(\Lambda\otimes X')\big)/
  \{(a,-(1\otimes f)a):a\in\Lambda\otimes X\}.
\] Both the
new target and the old quotient characters have presentations
from the split full module. Quotients of nested stable lattices
are closed under sums and subquotients, by taking preimages
between the two lattices. This proves preservation under
pushout.

If \(f\) is injective, the old extension is the
subextension with kernel \(\Lambda\otimes X\) of the pushout,
since \(\Lambda\) is free. Thus injection reflects admissibility.
The same operations work at any larger \(b\). Pullback along the
diagonal on \(O'/\lambda^b\) followed by addition on the two
kernels gives closure under sums. These operations make \(A(X)\)
an \(O'\)-submodule.

Write \(H_1=H^1(F,\Lambda)\). For the rank-\(r\) target the
residue calculation gives
\[
 H^0(F,\Lambda/\lambda)= H^2(F,\Lambda/\lambda)=0
\]
and \(\dim_{k'} A(k')\le d_f\),
where \(d_f\) counts the negative jumps of \(\Lambda\).
The cohomology sequences for the successive reductions
give the same two vanishings at finite depth, and passage to
the complete lattice gives the integral vanishings.
The multiplication sequence for \(\lambda^a\) therefore yields
\[
    H_1[\lambda^a]=0,\qquad
    H_1/\lambda^a H_1\xrightarrow{\sim}
                H^1(F,\Lambda/\lambda^a\Lambda)
            \qquad(a\ge1).
\]
The cyclic decomposition of finite \(O'\)-modules gives
\[
 H^1(F,\Lambda\otimes X)=H_1\otimes X.
\]

Let \(H_{1,f}\subset H_1\) be the inverse image of the rational
finite part. It is saturated, so \(H_1=H_{1,f}\oplus C\)
over the DVR. The crystalline dimension formula gives rank \(d_f\)
for \(H_{1,f}\). Retwist the crystalline extension of a finite
class and add the other scalar lines. Strictness of the
crystalline filtration adds the stated jumps, and the upper
triangular action gives the prescribed residual polynomial.
Its successive reductions are presentations for the class.
Thus
\[
 H_{1,f}\otimes X\subseteq A(X),
\]
with equality at \(k'\)
by the residue dimension bound.

To pass to arbitrary depth, first use the projection onto
each cyclic summand of \(X\). For \(X=O'/\lambda^n\), subtract
the finite component and take a complementary element
\(c\in C/\lambda^n C\). If \(c\ne0\), write \(c=\lambda^j c_0\)
where \(j<n\) and \(c_0\bmod\lambda C\ne0\). Then
\[
    \lambda^{n-1-j}c=\lambda^{n-1}c_0
\]
comes from the nonzero \(c_0\bmod\lambda C\) under
\(k'\hookrightarrow X\), \(1\mapsto\lambda^{n-1}\).
Reflection along this injection makes \(c_0\bmod\lambda C\)
admissible, contrary to equality at \(k'\). This proves
\[
 A(X)=H_{1,f}\otimes X
\]
for the rank-\(r\) target.

For the character ratios \(\epsilon^{1-c}\), \(2\le c\le t\),
the same reduction sequence applies. Use zero as the distinguished
summand in the propagation argument.
The residual invariant spaces of both this ratio and its
Tate dual are zero: the exponents \(1-c\) and \(c\) are
nonzero modulo \(\ell-1\) by \(t+1<\ell\), and residual
cyclotomic inertia has order \(\ell-1\) on the unramified
base. Local Tate duality gives the required vanishing in
degree two. The classification gives \(A(k')=0\), so
the cyclic-module argument gives \(A(X)=0\).

For the trivial ratio, reciprocity identifies the full
cohomology with
\[
 H^1(F,X)=\operatorname{Hom}_{\rm cont}
        \bigl(\mathbb Z_\ell\oplus U_F^1,X\bigr).
\]
The first factor is valuation and the pro-\(\ell\) principal
unit group \(U_F^1\) is free over \(\mathbb Z_\ell\), since
\(\ell\) is odd and \(F\) is unramified. To see that the whole
valuation summand is admissible, write \(x\in X\) for a desired
Frobenius value. Choose \(b\ge1\) with \(\lambda^b X=0\).
Take two split full representations, changing the first scalar
line of the second by an unramified character of Frobenius
value \(1+\lambda^b\). This preserves its jumps and residual
polynomial. With first scalar basis \(e_1,e_2\), set
\(v=(e_1+e_2)/\lambda^b\).

After removing the common character,
Frobenius on the lattice \(O'v\oplus O'e_2\) has
\[
           v\longmapsto v+e_2,\qquad
           e_2\longmapsto(1+\lambda^b)e_2 .
\]
Modulo \(\lambda^b\), Frobenius fixes \(e_2\) and acts on \(v\) by a
translation. Thus for each \(n\ge0\),
\[
  (\operatorname{Frob}^n-1)v=n e_2 .
\]
Sending \(e_2\) to \(x\) is well defined since \(\lambda^b x=0\).
The induced unramified homomorphism \(h\) is determined on a dense subgroup by
\[
 h(\operatorname{Frob}^{n})=nx\qquad(n\in\mathbb Z).
\]
Continuity determines \(h\) on the unramified quotient. Add the other summands to both
lattices of the presentation.
The residue upper bound is one, so equality at \(k'\)
with the valuation summand follows. Splitting off this
summand in the displayed homomorphism module and applying
the same cyclic-module argument kills the complementary
principal-unit part. All three calculations use the residue
tests at the same cut.
\label{cor:large-prime-length-3}

We verify the presentation hypothesis on the cut cocycles. Choose a
crystalline character \(\theta\) with labeled jumps \(w_i\).
Over an unramified local base with coefficients containing each
embedding such a character exists by the labeled crystalline character
comparison (equivalently use the algebraic characters of the units
with the matching reciprocity convention).
Twist the normalized family by \(\theta\). After twisting back, the scalar strings read
\[
 \chi_c=\theta\epsilon^{1-c}\quad(1\le c\le t),\qquad\chi_1=\theta.
\]
Exactness and twisting of the crystalline
filtration give the prescribed rank-\(r\) and full jumps.
At each family point the common residual polynomial follows from
localization of the Galois characteristic polynomials. Thus all
the family representations used in a presentation are full.
\label{cor:large-prime-length-4}

In the original normalized family let \(E\) be the presenting vector in a finite sum of fibres. For a finitely generated stable submodule \(Y\) and an equivariant map \(\phi\) on it, the cut cocycle has the presentation
\[
 gE-E\in Y,\qquad c(g)=\phi(gE-E).
\]

This is the local definition of the cut, obtained by applying
the functional to the corresponding differences in the global
matrix lattice. Extension of coefficients preserves it by extending
\(\phi\) on the free span of \(Y\). Its values form a finite set
in the divisible target, since the cocycle is locally constant
on a compact group. They are killed by some \(\lambda^a\).
In the enlarged coefficients use the cyclic module
\[
 X=\lambda^{-a}O'/O'\simeq O'/\lambda^a O'
\]
and choose its generator \(x_0\). Reading coordinates
identifies \(L\otimes X\) with the corresponding torsion in the
divisible target. The construction applies also to a stable
presentation mapping to that torsion. If \(D\) is the divisible target, replace \(Y\) by
\[
 Y_a=\ker\bigl(Y\xrightarrow{\ \phi\ }D\longrightarrow D/D[\lambda^a]\bigr).
\]
Since \(\lambda^a\phi(gE-E)=0\), each \(gE-E\) lies in \(Y_a\).
Equivariance makes \(Y_a\) stable. It is finitely generated because
it is a submodule of \(Y\) over a Noetherian DVR.

For clarity one can build the requisite finite quotient as follows.
Let \(i_L\) be the rational embedding of the target into the
good-point representation and \(e\) a nonzero vector in its
trivial line, before twisting. For a scalar target use its
eigenline for \(i_L\). Embed \(Y\oplus L\oplus O'\) in the
finite family sum with two extra copies of the good fiber by
\[
        (y,v,s)\longmapsto (y+sE,\ i_L(v),\ se).
\]
These last two coordinates make the map injective. After twisting
the whole sum by \(\theta\), the first coordinate supplies the
term \(s\theta(g)(gE-E)\), and the last coordinate transforms by \(\theta(g)\).

Consequently
\[
 (y,v,s)\mapsto(\phi(y)+v\otimes x_0,s)
\]
maps the
stable image equivariantly onto the extension with the cocycle
in question. Reduce its last coordinate modulo \(\lambda^b\) for
any \(b\ge a\). This gives a finite quotient of the stable image.
To realize it as a quotient of two full lattices, choose a stable
full lattice in the ambient representation. If \(M\) is the stable image, \(K\subset M\) its kernel and \(L\) the full ambient lattice, choose \(n\) such that \(\lambda^n L\cap M\subseteq K\). Adding this same deep lattice to both modules preserves the target quotient:
\[
 (M+\lambda^n L)/(K+\lambda^n L)\simeq M/K.
\]
Such a scale exists
by elementary divisors for the image in its rational span and
because the finite quotient is killed by \(\lambda^b\).
This proves the integral presentation required in the classification.
\label{cor:large-prime-length-5}

\proofstep{Retract the enlarged coefficient field before returning to the original divisible lattice.}
The residue calculation was made after enlarging the coefficient field. Its contribution to the numerator is a length for the original divisible lattice, so the same classes must be detected after descent.

Read back the resulting local conclusions by untwisting
\(\theta\). A finite class in the enlarged coefficient lattice
gives a finite class in the original divisible target up to
coboundary. Indeed the original integers embed as a saturated
rank-one direct summand of the enlarged integers. If \(K\subset K'\) are the two local coefficient fields with integers \(O,O'\), then
\[
 K\cap O'=O,
\]
because integrality is detected by the restriction of the absolute valuation. A linear retraction preserving
the integers and fixing one thus exists, and extends to the
fraction fields. Apply it to the rational finite cocycle and to
the cobounding vector. It commutes with the matrices over the
original field and preserves the crystalline period condition by
linearity. The same retraction descends the character coboundary
and the condition on inertia. The span of rank-\(r\) restrictions
therefore has length zero modulo its finite condition at every
place above \(\ell\). \label{cor:large-prime-length-6}
\end{proof}

\paragraph{Global terms after local vanishing.}
The cyclotomic character on inertia at an unramified completion
surjects on \(\mathbb F_\ell^\times\). Choose a generator in
this residue image. On the rank-\(t\) corner at the good point
its character roots are pairwise distinct modulo \(\ell\),
because their ratios have exponents of absolute value between
one and \(t-1<\ell-1\). Use the orthogonal idempotents from
the residual full matrix span of the rank-\(r\) block and the
character blocks. In the corner the coefficients of the
characteristic polynomial belong to the complete local Hecke
algebra by the matrix-span argument. Its simple residual factors
lift by Hensel's lemma. All the root-difference denominators at
the good point are units. Using these operators in the
determinant comparison therefore makes the coboundary costs
proportional to \(v_\Delta\) vanish. The absence of the first
character off the theta locus permits the same supported
determinant bound there. \label{cor:large-prime-length-7}

\begin{lemma}[A global valuation bound after local vanishing]
\label{lem:offwindow-global-budget}
On the same good family and outside the exceptional windows
and the two type primes above the cutoff, the supported determinant
satisfies
\[
       v_{cc}(i,\pi)\le s_\pi+V_\ell+v_\ell(h_E)
                           +t\sum_{v\in Q}v_\ell(\zeta_v-1).
\]
The numerator table uses \(\mathcal L_\ell=0\) there, retaining
the four summands displayed above as separate allowances.
\end{lemma}

\begin{proof}
At the prescribed tame places with unit root difference the
quadratic inertia relation forces the character cocycle to
vanish on inertia. At the other prescribed places the polynomial
inertia condition gives the same vanishing. On the auxiliary
set \(Q\) the nonzero difference \(\zeta_v-1\) gives cost
\(v_\ell(\zeta_v-1)\) for each scalar block. After removing
these restrictions the first character gives an everywhere
unramified homomorphism and costs at most \(v_\ell(h_E)\) by
global reciprocity. The other scalar characters are locally
coboundaries above \(\ell\) and unramified away from it, giving
the strict cohomology allowance \(V_\ell\). The rank-\(r\)
part satisfies the prescribed Selmer conditions and costs
\(s_\pi\). Thus on the good family
\[
 v_{cc}(i,\pi)\le s_\pi+V_\ell+v_\ell(h_E)
                   +t\sum_{v\in Q}v_\ell(\zeta_v-1).
\]
The tame sum with its factor \(t\) is at most \(\ell^{-2}\)
on the selected twist. The signature estimate bounds \(s_\pi\)
as before.

In the table for the numerator this gives
\(\mathcal L_\ell=0\) outside the exceptional windows above the
cutoff, with all four displayed allowances entering the comparison.
The three thresholds used for the jump boxes, admissible
parameters and conductor bound are independent of the place,
the point and the index of the chosen lift. Their maximum
suffices throughout this range. \label{cor:large-prime-length-8}
\end{proof}


\subsection{The cohomology contribution to the division estimate}
\label{subsec:division-length}
\hypertarget{source-division-length}{}

The local cases bound lengths for individual characters. In the division lattice these lengths occur with eigenspace multiplicities and normalized place weights. We now sum them at one point of the auxiliary family at scale \(x\). Write \(z=x^\epsilon\),
\(b=e^z\), \(f=2n\) and \(B_0=fmbz\), with \(N_c\ge0\) the rank
budget for this point. Put \(\mathcal E=z(B_0+N_c)\).
\begin{definition}[Weighted local congruence numerator]\label{def:weighted-congruence-numerator}
Fix a lift \(1\le j\le m\). Let \(K\) be the coefficient field
and let \(\mu_0\) be the base tuple. The sum below ranges over
places \(w\) where the local tuple differs from \(\mu_0\).
Write \({\rm N}w\) for the norm of \(w\), \(e_w\) for its
absolute ramification index, and \(\eta_j(w)\ge0\) for the sum
of the normalized lengths at the local Hecke algebras.
Define the numerator \(H_j\) by
\[
 H_j=\sum_{w\ {\rm twisted}}
       \frac{\log {\rm N}w}{[K:\mathbb Q]} e_w\eta_j(w).
\]
For the coefficient DVR \(O_w\) at \(w\mid p\) and a finite-length module \(L\), normalization means
\[
 \operatorname{len}_w(L)=\frac{\operatorname{length}_{O_w}(L)}
 {\operatorname{length}_{O_w}(O_w/pO_w)}
 =\frac{\operatorname{length}_{O_w}(L)}{e_w}.
\]
\end{definition}
The normalized local congruence length and its behavior under
removal of eigenspaces admit the following direct description.

Let \(O\) be such a DVR and \(M\) the integral module at one index
over \(j,w\), viewed through its image in its generic fiber \(V\).
The finite Hecke algebra acts through an algebra of \(O\)-valued
functions on a finite set of characters and separates those characters.
For a set \(Z\) of good characters write \(V_Z=\bigoplus_{x\in Z}V_x\),
where \(V_x\) is the generic eigenspace, and write \(p_Z\) for the
corresponding rational projector.

\begin{lemma}[Return from the cut to the original Hecke lattice]
Let \(S\supseteq Z\) be the characters retained by the integral operator of Proposition~\ref{prop:cut-selector}. At a fixed
\(j,w\) let \(\eta_{\rm cut},\eta_{\rm original}\) be the normalized
congruence lengths for \(p_SM\) and \(M\), respectively. Then
\[
 0\le\eta_{\rm original}-\eta_{\rm cut}
           \le(\dim V_Z)v_p(h),
\]
where \(h\) is the finite row group order in
Proposition~\ref{prop:cut-selector}. If the total good multiplicity
at each twisted place is at most \(d_j\), the weighted increase in
the numerator is at most \(d_j\log h\).
\end{lemma}
\begin{proof}
The generic Hecke algebra contains the projector onto \(V_x\):
for each \(y\ne x\) choose \(a_y\) separating \(x,y\), and take the
product of \((a_y-a_y(y))/(a_y(x)-a_y(y))\). Hence
\(V=\bigoplus_x V_x\). For \(M\) the upper and lower lattices are
\[
 U=p_ZM=V_Z\cap(M+V_{Z^c}),\qquad B=M\cap V_Z.
\]
Both have full rank in \(V_Z\); at \(Z=\varnothing\) both are zero.

Now retain a set \(S\supset Z\) of eigenspaces and project the
integral module onto them. The integral lattice of the quotient
in its generic fiber is \(p_S M\). The upper lattice for \(Z\) there
is still \(U\) and its lower lattice becomes
\(B'=p_S M\cap V_Z\supset B\). The integral type operator
\(a\) of Proposition~\ref{prop:cut-selector} acts on \(V_Z\) by
the order \(h\) of the finite row group and vanishes on
\(V_{S^c}\).

For any \(v\in B'\), write \(v=p_S u\),
\(u\in M\). Then \(a u=h v\) lies in \(B\). Thus \(B'/B\) is
annihilated by \(h\) and generated by at most \(\dim V_Z\) elements,
since \(B'\) is a lattice over a DVR. Consequently the sequence
\[
 0\longrightarrow B'/B\longrightarrow U/B
        \longrightarrow U/B'\longrightarrow0
\]
computes the difference of the normalized lengths. Only the scalar
actions on \(V_Z\) and \(V_{S^c}\) enter: for \(p_Su=v\in V_Z\)
the other retained components of \(u\) vanish.

Suppose the total good multiplicity at each twisted \(w\) is at
most \(d_j\), the small-level rank at \(j\). For \(w\mid p\),
\(\log {\rm N}w=f_w\log p\) and
\(\sum_{w\mid p}e_w f_w=[K:\mathbb Q]\).
Consequently the total loss on returning to the uncut numerator is
at most \(d_j\sum_{p\mid h} v_p(h)\log p=d_j\log h\).
The admissible parameter bound \(\log h\le C_1 x^{C_1} z\)
is absorbed by \(C_1 fmbz^2\) beyond a single threshold, even
though the group may vary at every point. Thus an estimate on the
cut numerator with bracket
\[
 d_j\left\{C\big(fmbz^2+(\log x)^2N_c\big)+V\right\}
 \tag{*}\label{eq:local-division-star}
\]
gives the same form on the original numerator, with \(C\) increased
by \(C_1\).

\end{proof}
\paragraph{The residual local condition on the selected cut.}
Use Proposition~\ref{prop:weighted-local-return} on the good cut
of Proposition~\ref{prop:cut-selector} at each
point index above the family threshold. At a retained localization
let \(k\) be the residue field of \(O\), \(\bar\rho\) the
rank-\(r\) residual constituent, and \(\bar\epsilon\) the
residual cyclotomic character. Write \(S_i,Q_i\) here for the two
sets of places in the norm choice, with \(Q_i\) the auxiliary
twist level. Both lie away from \(\ell\).
\begin{lemma}[Local vanishing on the retained row]\label{lem:cut-residual-vanishing}
For each retained localization satisfying the cut conditions, every
\(v\in S_i\), and every integer \(e\), one has
\[
 H^a\big(E_v,\bar\rho\,\bar\chi_1^{-1}\bar\epsilon^{\,e}\big)=0
                       \qquad (a=0,1,2).
\]
\end{lemma}

\begin{proof}
\proofstep{Separate the chosen tame eigenvalue.}
At either type prime that occurs in \(S_i\), the projected
roots of the residual rank-\(r\) inertia polynomial avoid the
projected character label. Thus at the chosen tame generator
\(\gamma_v\), \(u=\bar\chi_1(\gamma_v)\) is not an eigenvalue
of \(\bar\rho(\gamma_v)\).
At the additional prime of \(S_i\), use the choice away from
\(\ell\) of the two distinguished order-\(15\) primes.

On the
full representation the inertia polynomial reduces at \(\gamma_v\)
to \((X-1)^r(X-u)^t\), by the simultaneous inertia decomposition
of Section~\ref{subsec:retained-galois}. The residual factorization
also gives the rank-\(r\) polynomial times \((X-u)^t\), since all
the cyclotomic twists agree on inertia. Cancelling that monic
factor identifies the rank-\(r\) polynomial with \((X-1)^r\).
The exclusion of the eigenvalue one also holds at \(\ell=3,5\), since
\[
 \Phi_{15}(1)=1\not\equiv0\pmod\ell.
\]
Thus \(u\ne1\). These arguments use the inertia representation at
the actual localization, so passage to its retained eigensystems
preserves the test.

\proofstep{Use local duality and the Euler characteristic.}
Cyclotomic inertia is trivial at \(v\). Hence \(\gamma_v-1\)
acting on \(M_e=\bar\rho\,\bar\chi_1^{-1}\bar\epsilon^{\,e}\)
is invertible, independently of \(e\). On its Tate dual the
inertia matrix is the inverse transpose, which has the same
fixed-space vanishing. For the dual over \(\mathbb F_\ell\)
one may identify the linear dual over \(k\) by the nondegenerate
finite-field trace pairing. Local Tate duality gives \(H^2=0\).
The Euler characteristic of a finite \(\ell\)-primary module
over this nonarchimedean local field of residue characteristic
different from \(\ell\) is one. Finiteness and \(H^0=H^2=0\)
therefore give \(H^1=0\). This proves the rank-\(r\) vanishing
at all the \(S_i\)-places used in the row presentations.
\cite{milne-duality}\label{cor:division-length-1}
\end{proof}

At \(Q_i\) the twist can have \(\ell\)-power inertia. The tame allowance at these places is the separate quantity
\[
 t\sum_{v\in Q_i}v(\zeta_v-1),
\]
as in Section~\ref{subsec:retained-twist}.
\cite{milne-duality}\label{cor:division-length-2}

The cut can be made at all indices of one family. Below the
threshold retain the whole space. Above it choose the good
subset in each localization as in Section~\ref{subsec:selected-cut}.
The quotient to the retained generic eigenspaces preserves the
point readings, the local twists and the block and weight labels.
It therefore carries each chosen auxiliary character to its
restriction on the surviving tests. The same character still
satisfies the simultaneous inequalities on these tests. The
level comparison and the group order in the return to the
original numerator are taken on that original point. In particular
the threshold for the restriction is independent of \(j,w\)
and of the translate. The constants in the local table on this
cut are chosen for the whole resulting family, before evaluating
it at the selected point. \label{cor:division-length-3}

\paragraph{Strict classes in the marked presentation.}
\begin{lemma}[Scalar strict contribution]\label{lem:scalar-strict-cost}
On the selected cut, let \(c_{\rm sc}\) be the contribution to the bracket in \eqref{eq:local-division-star} per unit of good multiplicity, from blocks \(2\le c\le t\) summed with the normalized coefficient-place weights. Then
\[
 c_{\rm sc}\le V\le\sqrt b.
\]
\end{lemma}

\begin{proof}
\proofstep{Identify the twists on the original field.}
In each local block the other conditions on the cocycle
are imposed before applying the strict budget. If \(\chi\) is the common auxiliary character on the two lines, the target-to-quotient ratio in block \(c\) is
\[
 (\chi\chi_\ell^{\,1-c})\chi^{-1}=\chi_\ell^{\,1-c}.
\]

The normalization is the same when the family is cut into the
retained eigenspaces. After the estimated local restriction maps
at the auxiliary primes and the row and column conditions have
been imposed, the scalar kernel for \(c\ge2\) is unramified
away from \(\ell\) and locally zero at every place dividing it.
Take \(a=c-1\), which runs through \(1,\ldots,t-1\).
The \(c=1\) kernel is bounded by the class number, and the
rank-\(r\) constituent uses its specialized signature Selmer
length. The same scalar ratios and kernels occur at both
type primes and in the complementary-prime calculation.

\proofstep{Normalize coefficient lengths.}
For the actual coefficient DVR \(O\), with fraction field \(F\),
there is a Galois-equivariant coefficient isomorphism
\[
 (F/O)(-a)
 \simeq O\otimes_{\mathbb Z_\ell}
             (\mathbb Q_\ell/\mathbb Z_\ell)(-a).
\]
The cochain complexes for these discrete coefficients commute
with this tensor product by choosing a finite
\(\mathbb Z_\ell\)-basis of \(O\). Kernels of localization
commute as well. A cyclic factor of length \(h\) over
\(\mathbb Z_\ell\) becomes \(O/\ell^h O\), of length
\(e_Oh\) over \(O\). Thus the normalized strict length of
the scalar kernel costs at most the corresponding summand
of \(V_\ell\) on this field. This also proves independence
of the choice of splitting coefficient extension.

\proofstep{Sum the strict bounds at the selected rank.}
The character ratios identify individual strict twists. Summation must stop below the selected rank on the original field, so comparison with the full interval estimate uses precisely its strict cutoff.

The bound for the weighted sum over primes from Section~\ref{sec:family-window} holds
for all finite subsets and its terms are nonnegative.
Taking their supremum gives
\[
 V=\sum_\ell(\log\ell)V_\ell\le\sqrt b.
\]

Each rational prime is charged after the sum of the good
multiplicities over localizations has been bounded by \(d_j\)
and after the coefficient place weights have been summed.
This puts \(V\) in the bracket of \eqref{eq:local-division-star} with coefficient one.
\label{cor:division-length-4}
\end{proof}

The small-level rank \(d_j\) equals the base row
dimension \(D_0\) by Proposition~\ref{prop:base-theta-rank}. The base-level theta
map surjects onto the base row space and its adjoint has
rank \(D_0\). The integral projection denominators are
included in the length comparison above.

\subsubsection*{Summing the local bounds}
\hypertarget{source-numerator-sums}{}

A bound for a congruence ideal must first be turned into a bound for the corresponding integral module. This conversion keeps the rank of the retained eigenspace in the sum over good characters.

\begin{established}[Hecke congruence ideal and normalized length]\label{hyp:eigencharacter-lattices}
First consider the passage from a Hecke congruence ideal
to the congruence length of its module. Work over one
coefficient DVR \(O\), normalize \(v(\ell)=1\), and write
\(Y\) for the theta eigensystems and \(Z\subseteq Y\) for
the good ones. Let \(T\) be the Hecke algebra inside the
algebra of functions on the generic eigensystems. Write \(I_Y=\{f\in T:f|_{Y^c}=0\}\). For each \(\pi\in Z\) its evaluation and normalized colength satisfy
\[
 0\ne I_Y(\pi)=(f_\pi(\pi))\subseteq O,\qquad
 v_{cc}(\pi)=v(f_\pi(\pi)),\qquad f_\pi\in I_Y .
\]
\end{established}
\begin{proof}
For each generic character \(y\ne\pi\), choose \(t_y\in T\) with \(t_y(\pi)\ne t_y(y)\). The product \(\prod_{y\ne\pi}(t_y-t_y(y))\) vanishes away from \(\pi\) and has nonzero evaluation there. The scalars \(t_y(y)\) belong to the splitting DVR \(O\). Hence the image of the supported ideal is a nonzero ideal of \(O\). Choose a generator and lift it to a supported element of \(T\).
\end{proof}

\begin{lemma}[From eigencharacter valuations to module colength]\label{lem:eigencharacter-colength}
With the ideal and length of Lemma~\ref{hyp:eigencharacter-lattices},
let \(M\) be the Hecke-stable lattice and \(V_Z\) its good generic
subspace. Put \(U=p_ZM\), \(B=M\cap V_Z\)
and let \(a\) be the integral type operator of
Proposition~\ref{prop:cut-selector}. Then the normalized
colength of \(B\) in \(U\) is at most
\[
 \sum_{\pi\in Z}\dim(V_\pi)\bigl(v_{cc}(\pi)+v_\ell(h)\bigr).
\]
\end{lemma}

\begin{proof}
The type operator \(a\) acts by \(h\) on the good
eigenspaces and by zero on the remaining eigenspaces
of \(Y\). The product \(af_\pi\) kills every eigenspace
outside \(Z\). It also preserves the integral module,
so it maps \(U=p_ZM\) into \(B=M\cap V_Z\), by applying
it to a lift in \(M\). The products for the different
\(\pi\) may all act on multiple good eigenspaces.
Order those eigenspaces and filter \(U\) and \(B\) by
intersection with their successive partial sums.

For
the step adding \(V_\pi\), the graded lattice of \(U\)
is the projection of the intersection at this step
onto \(V_\pi\), and likewise for \(B\). Multiplication
by \(h f_\pi(\pi)\) takes the first graded lattice into
the second, since \(af_\pi\) preserves the partial sum.
For the step adding \(V_\pi\), the DVR index formula therefore bounds the normalized graded length by
\[
 \dim V_\pi\bigl(v_{cc}(\pi)+v_\ell(h)\bigr).
\]
With \(U_i,B_i\) the intersections with the first \(i\) eigenspaces,
the successive index is read in the projection to the new eigenspace.
Equivalently
\[
 \operatorname{len}(U/B)=
 \sum_i\operatorname{len}\bigl((U_i/U_{i-1})/(B_i/B_{i-1})\bigr)
 \le\sum_{\pi}\dim V_\pi\bigl(v_{cc}(\pi)+v_\ell(h)\bigr).
\]
This proves the local comparison, with the Hecke
algebra and the theta set of the cut when it is
applied to the cut.
\end{proof}

\begin{proposition}[Weighted local return on the cut]\label{prop:weighted-local-return}
Fix a lift on the selected point family and use its good
multiplicities from Proposition~\ref{prop:cut-multiplicity}.
The numerator of Definition~\ref{def:weighted-congruence-numerator} satisfies \eqref{eq:local-division-star} of Section~\ref{subsec:division-length} on the cut, with the allowed additional term \((\log x)^2N_c\) and a threshold uniform over the selected points. In the notation of the local table its contribution at an ordinary localization \(i\), away from the two type primes, satisfies
\[
 \eta_i\le \sum_{\pi\in Z_i^+}m_{i,\pi}
               \bigl(v_{cc}(i,\pi)+v_\ell(h)\bigr).
\]
Here \(v_{cc}(i,\pi)\) is the local congruence bound, \(h\) the order of the finite row group and \(m_{i,\pi}\) the good multiplicity. The separate loss \(d_j\log h\) from returning to the original Hecke module is absorbed at this scale.
\end{proposition}

\begin{proof}
\proofstep{Weight each localization by its normalized local degree and good multiplicities.}
Apply Lemma~\ref{lem:eigencharacter-colength} on the cut at
a single selected point. Fix a lift \(j\) and write \(d_j\) for the rank after
the small-level rank identification. Let \(\mathcal T\)
be the finite set of rational primes below the twisted
coefficient places in its support. The type and auxiliary prime cases use
Propositions~\ref{prop:type-prime-row-length}
and~\ref{prop:auxiliary-prime-row-length}. The local length
comparison at a prime \(\ell\) other than
\(g_0,g_1\), is the bound
\[
 \eta_i\le \sum_{\pi\in Z_i^+}
      m_{i,\pi}\bigl(v_{cc}(i,\pi)+v_\ell(h)\bigr).
\]
Here \(\eta_i\) is the normalized congruence length of the
cut localization, \(m_{i,\pi}\) is its multiplicity, and
\(v_{cc}\) is the local congruence bound attached to the
Galois family at \(\pi\). The factor \(h\) belongs to this
comparison on the cut itself. Returning to the original
module incurs the separate loss proved above.

Lemma~\ref{lem:marked-row-column-lattices} gives the required integral presentation at each good character of this cut. Proposition~\ref{prop:marked-signature} supplies its complement estimate. The resulting local table has
\[
 v_{cc}(i,\pi)\le
 \nu_\ell+2\ell^{-2}+\mathcal L_\ell+V_\ell+v_\ell(h_E),
\]
where \(\nu_\ell\) and \(\mathcal L_\ell\) are independent
of \(i,\pi\) for fixed \(\ell\) at this point. The table
incorporates the uniform bound on \(s_\pi\) separately
from the split row estimate. At each twisted \(w\mid\ell\), \eqref{eq:retained-multiplicity}
gives
\[
 \eta_{j,w}=\sum_{i\text{ over }(j,w)}\eta_i
 \le d_j\bigl(\nu_\ell+2\ell^{-2}+\mathcal L_\ell+V_\ell
                 +v_\ell(h_E)+v_\ell(h)\bigr).
\]
All terms in this upper bound are nonnegative and independent of
\(w\). Multiply by \(e_w\log\mathrm Nw/[K:\mathbb Q]\) and sum
over the chosen places. The factor from their weights is at most
\(\log\ell\), by the residue-degree formula on all places above
\(\ell\). At the two type primes the
paired-character estimate bounds the weighted lengths
directly by
\(d_j\log\ell(\mathcal L_\ell+2\ell^{-2}+V_\ell+v_\ell(h_E))\).
Up to nonnegative terms the same summation applies there.

The class number term contributes at most \(\log h_E\),
since
\[
 \sum_{\ell\in\mathcal T}v_\ell(h_E)\log\ell=\log\!\left(\prod_{\ell\in\mathcal T}\ell^{v_\ell(h_E)}\right)\le\log h_E.
\]
The factor \(h\) costs at most \(\log h\)
by the same argument. Also
\(\sum_{\ell\in\mathcal T}\ell^{-2}\log\ell\) is absolutely
bounded by the series over the integers. The strict
cohomology terms contribute at most the global budget
\[
 V=\sum_\ell V_\ell\log\ell
\]
once, since \(\ell\) is
summed only after the multiplicity and degree estimates.
This bound is independent of the number of localizations
and the coefficient-field degree.

The scalar cyclotomic term has its own bound for the weighted sum over primes.
Put \(a=tr\), which is odd since \(t\) is odd, \(k,m\)
are even, and \(r=k+m-1\). The cyclotomic minimum \(\nu_\ell\) in \eqref{eq:signature-cyclotomic-minimum}
is taken over the absolute Galois group of this same field \(E\). If \(\ell\mid Q_S\), the image
contains \(1+\ell^{v_\ell(Q_S)}\) since \(E\) lies in
the cyclotomic field of conductor \(Q_Sqd\). The
binomial argument in Section~\ref{subsec:split-local-length}
gives
\[
 \nu_\ell\le v_\ell(Q_S)+v_\ell(a).
\]

The depth at \(2\) is at least two. If
\(\ell\nmid Q_Sqd\), the cyclotomic image is full by
the Chinese remainder theorem. Then \(\ell\) is odd
and \(-1\) gives \(\nu_\ell=0\).
At \(q,d\), the host Artin action of the factor at
that prime has order dividing two. Thus the cyclotomic
image modulo that prime contains the subgroup of
squares, of order \((\ell-1)/2\). This order exceeds
\(a\) by the auxiliary prime choice, so a generator raised
to \(a\) differs from \(1\). This again gives
\(\nu_\ell=0\). Therefore
\[
 \sum_{\ell\in\mathcal T}\nu_\ell\log\ell
 \le \log Q_S + \log(tr)=O(\log x).
\]
Here \(Q_S\le4 f_\psi n\le n^2\) eventually by
\eqref{eq:host-character-compositum} and \eqref{eq:dyadic-coordinate-orders}.
The target is fixed before the scale and the ranks grow polynomially.

The multiplicity estimate and the local table join at
one common threshold in \(x\). The arguments of
Section~\ref{subsec:special-prime-lengths}
give the type and ramified entries. Section~\ref{subsec:ordinary-split-length}
proves the ordinary split entry and the implication from a local
inert band to the indicator in the sum. The cyclic row estimates
in Section~\ref{subsec:inert-rows} supply the long-orbit and
the two remaining inert entries. The comparison above the
cutoff in Section~\ref{subsec:outside-windows} uses the
integral local classification on the same jump windows. The table
assigns \(\mathcal L_\ell\) in the following ranges.

\begin{itemize}
\item At each type prime \(g_i\) the table entry and sizes satisfy
\[
 \mathcal L_{g_i}\le Cx^C\eta_c,\qquad \log\eta_c\le x^{\epsilon/4},\qquad \log g_i=O(z).
\]

\item At each dyadic or odd conductor prime and at three,
Proposition~\ref{prop:small-split-loss} bounds the cost by \(C fmb\log x+Cx^Cz^2\). The weighted
sum is \(O(fmbz^2)\) by the size comparison of
Section~\ref{subsec:split-local-length}.
\item At \(q,d\), the cost is at most \(2C fmb\).
The weighted sum is \(O(fmbz)\) by the conductor of the auxiliary field
bound on the chosen prime pair.
\item The remaining table uses a positive cutoff
\(\ell_*\). The summation below gives \(O(fmbz^2)\)
for the boundary, split, inert and short-band costs
weighted by \(\log\ell\) at primes at most \(\ell_*\).
\item Above \(\ell_*\) the remainder is supported on
the exceptional set. With
\[
 U=x^{C+2}\exp(x^{\epsilon/4}),
\]
its cardinality is at most \(388U^8\). The weighted cost
per exceptional prime is at most \(3000(C+1)^2U^{10}\),
giving at most \(1164000(C+1)^2U^{18}\) in total.
\end{itemize}

\proofstep{Sum the same table before absorbing the exceptional-prime cost.}
The table is now weighted by the local degrees and actual multiplicities. Summing those entries first keeps the exceptional primes and the scalar cohomology on the same retained family.

For fixed \(C\) and \(\epsilon>0\),
\[
 \log(U^{18}/b)=18(C+2)\log x+18x^{\epsilon/4}-x^\epsilon
 \longrightarrow -\infty .
\]
In particular the exceptional sum and the two type-prime
costs are \(o(b)\) uniformly. The same comparison absorbs
any fixed polynomial in \(x\). The bound on \(\log h_E\)
can be obtained separately. All character conductors of the abelian
host divide \(N\), so the conductor--discriminant formula
gives \(\log|\Delta_E|\le [E:\mathbb Q]\log N\).
Minkowski's bound supplies an ideal in every class of
norm at most \(2^{[E:\mathbb Q]}|\Delta_E|^{1/2}\).
\cite[Theorems 3.27 and 4.3 respectively]{milne-cft,milne-ant}
With \(n_E=[E:\mathbb Q]\) and
\(D=2^{n_E}|\Delta_E|^{1/2}\), count these representatives by
\[
 h_E\le\sum_{\mathrm N\mathfrak a\le D}1
  \le D^2\sum_{\mathfrak a}(\mathrm N\mathfrak a)^{-2}
  =D^2\zeta_E(2)\le D^2\zeta(2)^{n_E}.
\]
The last inequality is Euler factor by Euler factor, since there
are at most \(n_E\) prime ideals above each rational prime and
their residue degrees are at least one.

Taking logarithms and
using \(\log|\Delta_E|\le n_E\log N\) gives
\(\log h_E=O(n_E(1+\log N))\), which is polynomial in \(x\)
under the selected size bounds. With the local table,
these estimates give the cut numerator bound (*) with the
allowed additional
nonnegative term \((\log x)^2 N_c\). The loss \(d_j\log h\)
from the cut to the original level can be absorbed
at this same scale.

The local inputs are assembled on the cut family
before applying the weighted sum. The estimate uses
the complete local table and the rank and period
identification at the chosen points. \label{cor:numerator-sums-1}
\end{proof}

\subsubsection*{Summation below the prime cutoff}
\label{subsec:prime-cost-sums}
\hypertarget{source-prime-cost-sums}{}

The remaining sum below the cutoff contains contributions from both split and inert places as well as boundary terms. Their individual bounds have to be summed with the prime weights before the total can be compared to the allowed numerator. Write \(z=x^\epsilon\), \(b=\exp z\), and
\(H=b+S_c+\eta_c+t+2\). Let \(k\) be the ambient rank parameter,
and put \(N(m)=k+m(t+1)\).
Write \(\delta\) for the density of the selected indices with the larger weights.
\begin{established}[Size inequalities for the local table]\label{hyp:small-prime-table-ranges}
The cutoff and the secondary threshold are
\[
 \ell_*=\frac{mb}{n t^2(N(m)+t)z^2+k n^2(S_c+t+z)},
 \qquad A=\frac{kz}{m\delta}.
\]
Write \(X'=\exp(x^{\epsilon/3})\), \(b'=xb/X'\), and
\(\omega=b/(nX')\). Let \(\eta_3=\eta_c+t+2\), let \(w_P\) be
the upper place-weight size plus \((t+1)/2\), and put
\(G_P=w_P+S_c+\eta_3\). For the selected ranks and weights, beyond a threshold depending on the fixed parameters,
\[
  \ell_*\le b/z,\quad z\le A,\quad r/A\le 2m/z,\quad(3+\log A)^2\le z,\quad H\le2b,\quad z\le b'.
\]
\[
 \begin{gathered}
 \eta_3\le4\eta_c,\qquad w_P+\eta_3\le9b',\qquad G_P\le10b',\\
 r b'A\log x\le mb,\qquad r\eta_3\le4m\omega,\qquad
 C(\eta_3+S_c+z)\le\omega/8 .
 \end{gathered}
\]
\end{established}
\begin{proof}
Lemma~\ref{lem:window-parameter-choice} fixes the ranks before the scale and the weights of the retained positions before choosing the place. In this
calculation \(x^{-\epsilon}\le\delta\le1\), \(r\le2k\),
\(R\le N(m)+t\), \(m\le k\), and the ranks are at most a fixed
power of \(x\). In particular \(\ell_*\le b/z\), and
\[
 z\le A,\qquad r/A\le 2m/z,\qquad
 (3+\log A)^2\le z
\]
beyond one threshold. The last inequality follows from
\(A\le k z^2\) and \(z/(\log x)^2\to\infty\).
The size inputs also give \(H\le2b\).

At the smaller places the two possible centers satisfy
\[
 \omega\le 5b'/4,\qquad b_v+\lfloor b'/x\rfloor\le5b'/4+b'/x .
\]
Here the bound for \(b_v\) is from Section~\ref{subsec:exceptional-range}. The widths \(S_c,\eta_c,t+2\) are \(o(b')\), and \(t+2\le3\eta_c\). Including these widths and the half shift gives \(w_P+\eta_3\le9b'\) and \(G_P\le10b'\) uniformly.

For the fourth inequality the selected reserve bounds
\(r(b'+S_c)A\log x\) by \(\delta mb\). The fifth follows from
the reserve \(N(m)\eta_c\le\delta m\omega\) and \(r\le N(m)\).
For the last inequality, the sizes of \(\eta_c,S_c,z\) are
\(O(\omega/\sqrt x)\), with constants fixed before taking \(x\)
large. For fixed \(\epsilon>0\) this follows also by comparing
their bounds, \(x^{O(1)}\exp(x^{\epsilon/4})\), with \(\omega\).
We have \(b'\ge z\) at large \(x\).
\end{proof}

\begin{proposition}[Summation up to the small-prime cutoff]\label{prop:small-prime-cost-total}
Under the cutoff and size inequalities of Lemma~\ref{hyp:small-prime-table-ranges},
let \(c_\ell\) be the sum of the local boundary, split and inert charges in the table at \(\ell\), before multiplication by \(\log\ell\). Then
\[
 \sum_{\ell\in P}(\log\ell)c_\ell\ll_C fmbz^2\qquad\bigl(P\subseteq\{\ell:5\le\ell\le\ell_*\}\text{ finite}\bigr).
\] The threshold
is uniform in the retained point and the subset of primes.
\end{proposition}

\begin{proof}
For a set of integers \(2\le a\le L\), summing over all such
integers gives
\[
 \sum \log a\le L\log L,\qquad
 \sum \frac{\log a}{a}\le (1+\log L)^2.
\]
The second bound uses the harmonic sum and \(\log a\le\log L\).
Replacing \(a\) by \(a-1\) in the denominator costs at most a
factor two. These bounds apply to any subset of the indicated primes.
The boundary cost \(Cfn t^2 R z^2\) consequently contributes at
most \(C fmbz\), by the first term of the cutoff denominator.
\proofstep{Split places.}
The split cost is \((2fm+f)e_P(\ell)\), where
\(e_P(\ell)=C(1+H/\ell+\log_\ell H)\). The other uses of \(e_P\)
have coefficient at most \(2Cfm\). After multiplication by
\(\log\ell\), the three terms in \(e_P\) sum to
\[
 O_C\!\left(\ell_*\log\ell_*+
             H(1+\log\ell_*)^2+\ell_*\log H\right)
       =O_C(bz^2).
\]

\proofstep{Inert places below the secondary threshold.}
For inert \(\ell<A\) of order exceeding \(4L_w\) in the
host's real cyclotomic quotient, the cost is \(4Cfr\) times
\[
 J_\ell=\frac{H}{(\ell-1)^{j_\ell+1}}
         +(G_P+1)(j_\ell+1)+z,\qquad
 j_\ell=\max(1,\lceil\log_{\ell-1}A\rceil).
\]
For \(\ell\ge5\), \((\ell-1)^{j_\ell}\ge A\) and
\(j_\ell+1\le3+\log A\). The logarithm-weighted sum of the
first term, including \(r\), is
\(O(rH(1+\log A)^2/A)=O(mb)\).
For the other terms it is
\(O(rb'A(3+\log A)^2)=O(mbz)\). This uses the reserve for
\(rb'A\) and \(G_P+1=O(b')\).

The other inert primes have a contribution \(4CfrH/\ell^2\).
At \(\ell\ge A\) use
\[
 \sum_{a\ge A}(\log a)/a^2=O((1+\log A)/A)
\]
by integral comparison, and \(r/A\le2m/z\). At \(\ell<A\)
the order is at most \(4L_w\); here the ordinary tail bound at
\(A\) gives no control. Put \(R_0=n/z\). The parameters of the auxiliary construction
give \(3\le R_0\), \(k/m\le R_0^2\), \(n\le o2^e\) and
\(2oL_w\le z\), where \(e\) is the host's dyadic exponent
and \(o\) its fixed cyclic order parameter. Thus
\[
 2^e/L_w\ge2R_0.
\]

The quotient
\((\mathbb Z/2^{e+2}\mathbb Z)^\times/\{\pm1\}\) has order \(2^e\),
so every order \(h\) on it is a power of two. By
Lemma~\ref{lem:short-real-dyadic-order} an integer of order \(h\) is
\(\pm1\pmod{2^{e+2-v_2(h)}}\).

Consequently \(h\le4L_w\) forces the original integer to
be \(\pm1\bmod M\), where
\[
 M=2^{e+2-\lfloor\log_2(4L_w)\rfloor}
       \ \ge\ 2^e/L_w\ \ge\ 2R_0 .
\]
The same congruence holds if \(a=1\). There is no term at
\(1\) in the sum of interest. All the other positive integers
in these two residue classes are \(jM\pm1\) with \(j\ge1\).
They are between \(jM/2\) and \(2jM\), so
\[
 \sum_{\substack{a\ge2\\a\equiv \pm1\bmod M}}
            \frac{\log a}{a^2}
   \le \frac8{M^2}\sum_{j\ge1}\frac{\log(2jM)}{j^2}
   =O\!\left(\frac{1+\log M}{M^2}\right).
\]
The function \((1+\log u)/u^2\) is decreasing for \(u\ge1\).
Multiplication by \(r\le2mR_0^2\) therefore gives
\(O(m(1+\log R_0))\). This bound accounts for the small
inert primes using only the dyadic quotient, with no estimate
for primes in arithmetic progressions.

The remaining inert terms cost \(4Cfr\) times
\[
  \frac{w_P+\eta_3}{\ell}+S_c+z+\eta_3 I_\ell .
\]
Here \(I_\ell\) is the indicator that some \(s\) between
\(\omega\) and the upper place-weight size before the
\((t+1)/2\) shift satisfies
\[
 |\ell s-b|\le
   C\bigl(\ell(\eta_3+S_c+z)+b/\ell+s\bigr).
\]
For \(\ell\ge2C\) the right side is at most
\(\ell s/8+b/2+\ell s/2\), by the selected reserves.
It follows that \(\omega\le s\le4b/\ell\) in this band.
Hence \(r\eta_3 I_\ell\le16mb/\ell\).

At the finitely
many \(\ell<2C\) we use \(r\eta_3\le4m\omega\le4mb\).
Also \(r(w_P+\eta_3)\le9rb'\le9mb\), and the second
term of the cutoff denominator gives
\[
 \ell_*r(S_c+z)\le2mb.
\]
The first and band terms
are summed using \(\sum (\log\ell)/\ell\), and the
constant term using \(\sum\log\ell\). Their total
contribution is \(O_C(fmbz^2)\).

The estimates above apply together to any finite subset of
the primes \(5\le\ell\le\ell_*\). They give the desired
\(O_C(fmbz^2)\) bound for the numerical table used in the
weighted export. The threshold is common to all the retained points with the specified size inputs.
The sum is therefore available for the actual row once each localization satisfies the corresponding local Galois estimate. \label{cor:prime-cost-sums-1}
\end{proof}

\subsubsection*{The exceptional range}
\label{subsec:exceptional-range}

Above \(\ell_*\), only the exceptional jump windows require this additional sum. We count the primes they can contain and bound the total weighted cost uniformly in the window data. Write
\[
 \mathcal D=mb/\ell_*,\qquad
 U=x^{C+2}\exp(x^{\epsilon/4}),\qquad
 s_+=\lfloor2b/\ell_*\rfloor .
\]
The selected parameters give \(U\ge1\),
\(n,k,t,m,z,\eta_c\le U\), \(f,r\le2U\),
\(R,N(m)+t\le4U^2\) and \(S_c\le4U\).
Fixed polynomial rank bounds suffice after taking
\(C\) large enough. The two summands of \(\mathcal D\)
are at most \(4U^7\) and \(6U^4\) respectively, hence
\[
 \mathcal D\le10U^7.
\]
Since \(U^{14}=o(b)\), this
gives \(\ell_*^2>8b\) beyond a common threshold.
Also \(\ell_*>A\), since \(A\le k z^2\) grows
polynomially in \(x\).

The centers \(W\) of the exceptional windows are
\(b,\omega\), and \(b_v+\lfloor b'/x\rfloor\) for the
top places \(v\). Their number is at most \(2n+2\),
using the bound \(2n\) for the number of top places.
The weight construction gives
\(\omega\le b_v\le 5b'/4\), so each center lies in
\([\omega,2b]\) for large \(x\). The window about \(W\)
contains the integers in
\[
[W-\eta_c-2s_+-3,\,
  W+S_c+2\eta_c+2s_++4].
\]
All these integers are positive and less than \(3b\).
Indeed \(s_+\le2\mathcal D/m\) and
\(U^7=o(\omega)\) since \(n\) is polynomial,
\(\omega=b/(nX')\) and \(\log X'=x^{\epsilon/3}\).
This controls both endpoint displacements.

The exceptional set consists of primes \(\ell>\ell_*\)
dividing one of these integers. Two distinct such
divisors would have product greater than
\(\ell_*^2>8b\), so each integer contributes at most
one prime. An interval of length \(D\) contains at
most \(D+1\) integers. Thus the set has size at most
\[
 (2n+2)(S_c+3\eta_c+4s_++8)
 \le(4U)(97U^7)=388U^8.
\]
Each prime in the set is less than \(3b\), so
\(\log\ell\le3U\).
\label{cor:prime-cost-sums-2}

For these primes the boundary cost is
\(Cfn t^2 R z^2\le8C U^8\). The envelopes for the
pointwise row costs are the split cost and the second
inert cost, since \(\ell>\ell_*>A\). For \(\ell>7\),
\[
\begin{split}
 H/\ell^2 &\le H/\ell \le 2\mathcal D\le20U^7,\\
 e_P(\ell)&=C(1+H/\ell+\log_\ell H)\le23C U^7.
\end{split}
\]
For the flat part of the inert cost, the numerical
expression is
\((w_P+\eta_3)/\ell+S_c+z+\eta_3 I_\ell\),
where \(I_\ell\in\{0,1\}\) is the band indicator.
Here \(b'<\ell_*\) at the common threshold, by
\(\mathcal D x/X'\to0\), and \(w_P+\eta_3\le9b'\).
Thus the flat part is at most
\(9+4U+U+4U\le18U^7\).

The split envelope is at most \(138C U^9\).
The inert envelope is at most
\(608C U^9+92C^2U^9\), using
\[
 4Cfr(20+18)U^7+2Cfm(23C U^7)
       \le (608C+92C^2)U^9.
\]
Adding the boundary cost and multiplying by
\(\log\ell\) gives the bound
\(3000(C+1)^2 U^{10}\) per prime.
This proves the exceptional contribution to the
weighted sum from the pointwise row comparisons.
The complementary large-prime range uses the
congruence estimate with zero additional length
outside these windows.
\label{cor:prime-cost-sums-3}



\begin{proposition}[The numerator on the original lift]\label{prop:original-weighted-numerator}
Use the completed family with local contributions from
Lemma~\ref{lem:cut-residual-vanishing},
Lemma~\ref{lem:scalar-strict-cost} and the summation in
Section~\ref{subsec:division-length}. Include the correction
for the return from the cut in \(C\). Then beyond a common
threshold, at every point of the family and for every lift \(j\), the
weighted numerator on the original module satisfies
\[
           H_j\le D_0(|C|+1)z(B_0+N_c).
\]
\end{proposition}

\begin{proof}
Apply \eqref{eq:local-division-star} on the original module. The rank equality and the
inequalities \((\log x)^2\le z\), \(z\ge1\), and \(V\le\sqrt b\le B_0\)
give the stated bound. Choose the maximum of the thresholds for
Lemma~\ref{lem:scalar-strict-cost}, the cut comparison,
Proposition~\ref{prop:base-theta-rank} and the logarithm estimate.
Each threshold applies to the family before choosing its point or
lift.
\label{cor:division-length-5}
\end{proof}

The bound controls the cohomological term in the dual-lattice index
below. The first division error uses the transported row indices and
Weyl inclusion on the actual integral lattices of the same member in
Section~\ref{sec:word-division}. Both estimates and the period identities
are used in \eqref{eq:division-slopes} with the fixed exponent \(e_D=2\).

\paragraph{The contribution of the numerator to the dual index.}
\begin{established}[Nested lattices at the twisted place]\label{hyp:signed-dual-lattices}
At the used tuple, let \(\Lambda_\Delta\subseteq\Lambda_\nabla\)
be the smaller and ordinary section lattices of the lift in the
whole block. Transfer the two into the same coordinate space as in
Section~\ref{subsec:selected-period}, and put
\[
 \mathcal T=P_j(E^{-1}\Lambda_\nabla),\qquad
 \mathcal C=(VZ)^{-1}\Lambda_\nabla,\qquad
 \mathcal D=(VZ)^{-1}\Lambda_\Delta .
\]
These full lattices satisfy \(\mathcal D\subseteq\mathcal C\subseteq\mathcal T\). In the notation at the beginning of the section,
\(e_w\eta_j(w)=\operatorname{length}(\mathcal T_w/\mathcal C_w)\)
at the twisted place. Let \(J_w\) be the least nonnegative
integer such that \(\mathfrak p_w^{J_w}\mathcal C_w\subseteq
\mathcal D_w\). The congruence length uses \(\mathcal C_w\)
inside \(\mathcal T_w\). The exponent \(J_w\) measures
the further inclusion of \(\mathcal D_w\) into \(\mathcal C_w\).
\end{established}
\begin{proof}
The transfer identifies the generic spaces, and the block projection is surjective, so the three images span their prescribed spaces. The first inclusion follows by inverse image from \(\Lambda_\Delta\subseteq\Lambda_\nabla\). If \(v\in\mathcal C\), then \(EI_jv=VZv\in\Lambda_\nabla\). Thus \(I_jv\in E^{-1}\Lambda_\nabla\), and applying \(P_jI_j=1\) shows that \(v\in\mathcal T\). Localization preserves both inclusions. The quotient of the second pair has the stated normalized length by the definition of \(\eta_j(w)\), and a power of the maximal ideal kills the finite quotient of the first pair.
\end{proof}

\begin{proposition}[The local signed dual index on the common lift]
\label{prop:signed-lift-dual-index}
For the three lattices and \(J_w\) of
Lemma~\ref{hyp:signed-dual-lattices}, let \(D_G\) clear
the positive class and level denominators. Suppose \(s\) is a
positive integer satisfying the classwise value dual inclusion.
Then
\[
 [\mathcal T_w:(\mathcal T^\#)_w]_w
 \le e_w\eta_j(w)+D_0\bigl(J_w+v_w(sD_G)\bigr).
\]
If the same integer \(s\) works at every place in the sum and
\(D_*\) is a common multiple of the \(D_G\) at that point, divide
by \(D_0\) and sum with the normalized place weights. The
denominator correction is then at most \(\log D_*+\log s\).
The least exponent \(J_w\) is charged to the first error.
\end{proposition}

\begin{proof}
\proofstep{The class denominators and the global Hermitian dual.}
The Hermitian form on these coordinates is the restriction of the
level form. It is a sum of the class pairings divided by
the positive class weights, with the level normalization included.
Let the positive integer \(D_G\) be a common multiple of
these denominators, and suppose the positive integer \(s\)
satisfies the value-lattice
dual inclusion at every class. This means that the class pairing
of an ordinary value with \(s\) times any smaller-lattice
value is integral.

For class weight \(d_c\) and local pairing \(B_c\), the contribution of ordinary \(u\) and smaller-lattice \(v\) is integral after scaling:
\[
 sD_G\,\frac{B_c(u,v)}{d_c}
     =\frac{D_G}{d_c}B_c(u,sv).
\]
Here \(d_c\mid D_G\) and the value-lattice hypothesis makes \(B_c(u,sv)\) integral. Summing puts \(sD_G\) times the smaller section lattice in the ordinary dual. Restriction to the level subspace
uses the same class pairing and its normalizing denominator.
For \(z\in\mathcal D\) and \(y\in E^{-1}\Lambda_\nabla\)
orthogonality gives
\[
 \langle z,P_jy\rangle=\langle I_jz,y\rangle
\]
for these restricted forms. It follows that
\(sD_G\mathcal D\subseteq\mathcal T^\#\), where
\(\mathcal T^\#\) is the Hermitian dual over \(\mathcal O_K\).
Therefore
\[
 \mathfrak p_w^{J_w+v_w(sD_G)}\mathcal C_w \subseteq(\mathcal T^\#)_w.
\]
This inclusion concerns the
localization of the global dual and remains valid at places
interchanged by the coefficient involution.

\proofstep{The signed determinant index at the chosen place.}
For full local lattices take the signed index \([L:L']_w\)
to be the valuation of the determinant of a basis of \(L'\)
in a basis of \(L\). These indices add in triples of lattices
and increase as \(L'\) decreases. With \(d_j(w)=D_0\) the last
inclusion gives
\[
 [\mathcal T_w:(\mathcal T^\#)_w]_w
 \le
  e_w\eta_j(w)+D_0\bigl(J_w+v_w(sD_G)\bigr).
\]
\proofstep{Normalization of the place sum.}
Put \(\omega_w=\log\mathrm Nw/[K:\Q]\). Dividing the
local inequality by \(D_0\) and summing at the twisted places
gives
\[
 D_0^{-1}\sum_{w\ {\rm twisted}}\omega_w
        [\mathcal T_w:(\mathcal T^\#)_w]_w
 \le H_j/D_0
     +\sum_{w\ {\rm twisted}}\omega_w J_w
     +\sum_{w\ {\rm twisted}}\omega_w v_w(sD_G).
\]
Here \(H_j=\sum_{w\ {\rm twisted}}\omega_w e_w\eta_j(w)\).
Choose \(D_*\) divisible by every denominator \(D_G\) at
this point, and use the same \(s\) at every place. Then
\(0\le v_w(sD_G)\le v_w(sD_*)\). For a positive integer \(a\),
\[
 \sum_{w\mid p}\omega_w v_w(a)
   =\frac{\sum_{w\mid p}e_w f_w}{[K:\Q]}v_p(a)\log p
   =v_p(a)\log p,
\]
where \(f_w\) is the residue degree over \(p\).
This gives the upper bound \(\log D_*+\log s\) on the
last sum, also for a partial set of places. Thus the weighted signed index per lift,
divided by the positive row rank of the selected bundle,
is at most
\[
        (|C|+1)z(B_0+N_c)+\log D_*+\log s
        +\sum_{w\ {\rm twisted}}\frac{\log\mathrm Nw}{[K:\Q]}J_w .
\]

In the
global slope comparison the term \(J_w\) is charged in
the same per-lift allowance as the row comparison error.
Write \(J_{j,w}\) for this least exponent and \(e_w\) for
the ramification index. The chosen nonnegative errors
\(r_{j,w}\) bound the paired row-index remainder after
summing over lifts and including the declared \(2/p^2\)
term. The per-lift allowance \(E_{1,j}\)
must satisfy
\[
  \sum_{w\ {\rm twisted}}\frac{\log\mathrm Nw}{[K:\Q]}
                  (e_w r_{j,w}+J_{j,w})\le E_{1,j}.
\]
In the numerical scale the normalized exponent
\(J_0(j,w)\) is \(J_{j,w}/e_w\). Separately, the
allowance from the numerator divided by \(D_0\), plus
\(\log D_*+\log s\), is the quantity denoted \(E_{2,j}\)
in the division estimate. Thus the preceding signed-index
estimate gives the two terms with their specified
normalizations.

\proofstep{The rational-prime majorant for the least exponent.}
For the integral inclusion of lift lattices, let \(J(j,p)\) be the rational-prime exponent that bounds the local exponents:
\[
                     J_{j,w}\le e_w J(j,p)\qquad(w\mid p).
\] Each such inclusion concerns the smaller
lattice in the same generic lift space. A bound for
\(\sum_p J(j,p)\log p\) transfers to the weighted coefficient
places by the local degree identity
\[
 \sum_{w\mid p} e_w\,\frac{\log\mathrm Nw}{[K:\Q]}=\log p .
\] It remains to sum these lift exponents. Their contribution comes from lattice inclusion and is accounted for separately from the cohomology numerator.
\label{cor:division-length-6}
\label{cor:division-length-7}

\end{proof}

\paragraph{Aggregation of the least exponents.}
For the selected lift \(j\) let \(\mathcal C_w,\mathcal D_w\)
be the two local section lattices defining \(J_{j,w}\).
Both use the selected level, good subspace and tuple
\(\mu(w)\). The weight record of \(j\) uses the actual
weight at each embedding, up to an additive shift.
Let \(B_j(p)\) be the nonnegative integer obtained by the local
exponent rule below. The inclusion to be proved is
\[
                     p^{B_j(p)}\mathcal C_w\subseteq\mathcal D_w .
\]
Outside the active range the comparison uses the sum
of the full Weyl exponents over the \(f\) labels, with
the tail formula at the upper end. Inside the active
range the simultaneous local-model inclusions imply this
inclusion at the same rational-prime power. In those
models the indices with larger weights use the smaller of the full Weyl exponent
and the floor of the coarse bound for the surviving depths, and the indices with smaller weights
keep the elementary exponent.

\begin{lemma}[Summing the least local exponents]
In the specified size, jump and exceptional-prime range,
the least local exponents satisfy, uniformly in the lift and the coefficient field,
\[
 \sum_{w\ {\rm twisted}}\omega_wJ_{j,w}\le32\mathcal E,\qquad\omega_w=\frac{\log\mathrm Nw}{[K:\mathbb Q]}.
\]
\end{lemma}
\begin{proof}
Using \(p\) as a scalar costs \(e_w\) per power, so
\(J_{j,w}\le e_w B_j(p)\). Put \(\omega_w=\log\mathrm Nw/[K:\Q]\)
and let \(P_{\rm tw}\) be the rational primes below the
twisted places. Then
\begin{align}
 \sum_{w\ {\rm twisted}}\omega_w J_{j,w}
 &\le\sum_{p\in P_{\rm tw}}B_j(p)
            \sum_{\substack{w\ {\rm twisted}\\ w\mid p}}e_w\omega_w
 \le\sum_{p\in P_{\rm tw}}B_j(p)\log p \notag\\
 &\le 30|T_2|mbz^2+2mb
 \le32\mathcal E . \label{eq:division-least-sum}
\end{align} For the third inequality, use the size, jump range and exceptional-set conditions of the weight record in the prime-sum bound. Its
main and exceptional contributions together are at most
\(30|T_2|mbz^2\), and its tail contributes at most \(2mb\).
Here \(T_2\) bounds the number of indices with larger weights and has
size at most \(f\).

The last step uses \(z\ge1\) and
\(B_0=fmbz\). The estimate has the same constant
for every lift and is independent of the
coefficient degree. The calculation just given transports the prime sum.
Sections~\ref{subsec:lift-exponent-models} and
\ref{subsec:lift-exponent-primes} give the cover, weight
comparison and prime sum for this record, with the
representation-theoretic hypotheses used for the local linkage
and the integral graded bound. The paired row-index remainder
enters separately in Section~\ref{sec:word-division}.
\label{cor:division-length-8}

\end{proof}
\subsubsection{Lattice inclusions on the selected lift}\label{subsec:lift-exponent-models}
\hypertarget{source-lift-exponent-models}{}

The least exponents entering the signed index measure a further lattice inclusion beside the congruence length. To bound them, put the two coefficient lattices in the same space of forms at the selected lift.

Fix a selected lift, a tuple and a coefficient place \(w\mid p\).
Write \(Z\) for the selected space of forms at its row level.
In the coordinates of \(Z\), \(\mathcal C\) and \(\mathcal D\)
are the inverse images of the value lattices for the dual Weyl
and Weyl coefficients, respectively. The coordinate maps on
the finite class set identify both lattices simultaneously.
\begin{proposition}[Weyl inclusion on the original section lattices]
\label{prop:local-weyl-value-inclusion}
On the selected tuple and place \(w\mid p\), the sum of the exact Weyl exponents at all labels gives
\[
 p^{\sum_\tau e_p(\lambda_\tau)}\mathcal C_w\subseteq\mathcal D_w.
\]
For \(y<p\le h\) as in Lemma~\ref{hyp:weyl-prime-range},
exclude \(q,d,p_{\rm de},p'_{\rm de}\) and primes dividing an
integer of \(\mathcal I\). A power that gives the value inclusion at every local model under
an occurring integral Hecke character gives the same inclusion on
\(\mathcal C_w,\mathcal D_w\) at the original level.
\end{proposition}

\begin{proof}
\proofstep{Tensoring the exact inclusions and localizing the value map.}
At each coefficient label \(\tau\) choose a local basis of the
underlying lattice. The two coefficient models in that basis
use the dual Weyl and Weyl lattices \(W_\tau^\vee,W_\tau\) in
the prescribed rational realizations. The inclusion of
Section~\ref{subsec:lift-exponent-primes} tensors to
\[
 p^{\sum_\tau e_p(\lambda_\tau)}
          \bigotimes_\tau W_\tau^\vee
       \ \subseteq\  \bigotimes_\tau W_\tau .
\]
An additive weight translation twists both sides by the same
determinant line and does not change this exponent. If
\(\operatorname{ev}\) is the finite map of value coordinates,
the inverse image of the Weyl lattice is the kernel of its
composition with reduction modulo that lattice.
Writing \(q\) for reduction modulo the Weyl value lattice, exactness of localization gives
\[
 (\ker(q\circ\operatorname{ev}))_w=\ker(q_w\circ\operatorname{ev}_w).
\]
Equivalently one can
multiply the finitely many denominators outside \(w\) occurring
in the individual coordinates. The displayed inclusion therefore
gives the asserted exponent on each tuple.
\label{cor:lift-exponent-models-1}

\proofstep{Comparing the two lattices in the same operator space.}

We describe the sharper bound at primes in the weight range
\(y<p\le h\) outside the exceptional interval and specified
set, using the tuple assigned to \(w\). The base level and selected block hypotheses give finite free tensor value lattices over the same local coefficient ring. Choose
an embedding of the coefficient field into \(\overline{\Q}_p\)
inducing \(w\), and adjoin the roots of unity for the finite
level character to the induced local field. Its integer ring
\(O\) is a DVR. The integral frame data are changes of basis
on the two local value lattices, including the level group
matrices and character.

In the common interpolation frame a Hecke correspondence is represented on both lattices by the same operator
\[
 T=\sum_j a_j R(g_j),\qquad a_j\in O,
\]
with \(R(g_j)\) the specified right translations. Thus on the auxiliary space \(M\) of
integral forms the operators define an action of the commutative
polynomial algebra \(\mathcal T\) over \(O\) in these operators.
This action preserves the forms with values in the Weyl lattice,
whose submodule is \(M_{\rm W}\).

\proofstep{Detecting a failed inclusion at an occurring residual character.}
Take the integral form \(x_0\) from the dual Weyl reading of a
vector of \(\mathcal C\). Its value coordinates in the auxiliary
model satisfy covariance: it holds after scalar extension, and
extension is injective on the coordinates. The cyclic module
\[
 C=\mathcal T x_0\subset M
\]
is finite free over \(O\), and its
generic extension lies in the selected Hecke-stable span.
Write \(z=[p^a x_0]\) in \(M/M_{\rm W}\). If \(z\ne0\),
choose a maximal ideal \(\mathfrak m\) containing its annihilator.
Then \((\mathcal T z)_{\mathfrak m}\ne0\).
The kernel of \(\mathcal T\to\operatorname{End}_O C\)
kills \(x_0\) and hence \(z\). Thus the inclusions of ideals
\[
     \ker(\mathcal T\to\operatorname{End}_O C)
                   \subseteq\operatorname{Ann}_{\mathcal T}(z)
                   \subseteq\mathfrak m
\]
identify \(\mathfrak m\) with a maximal ideal of the
finite image algebra \(A\).

In a splitting field simultaneously triangularize the action on
\(C\). The diagonal characters \(\chi_j\) are integral,
and the matrices of \(J=\bigcap_j\ker\chi_j\) are strictly upper
triangular. In a triangular basis, with \(C_i\) the span of the first
\(i\) vectors, this gives
\[
 JC_i\subseteq C_{i-1}\quad(1\le i\le\dim C),
 \qquad J^{\dim C}C=0 .
\]
The action of the image algebra is faithful, so \(J\) is nilpotent. The product of integer rings
of the \(\chi_j\) is integral over \(A/J\). Lying over
gives a coordinate character with reduction \(\mathfrak m\).
Its generalized eigenspace is nonzero. On the generalized eigenspace \(V\), let \(a\) be the last exponent with \(\mathfrak n^aV\ne0\), for the maximal ideal \(\mathfrak n\) of the character. Then
\[
 0\ne\mathfrak n^aV\subseteq\ker\mathfrak n,\qquad \mathfrak n^{a+1}V=0.
\]
Any nonzero vector in this term is an eigenvector. It is an eigencharacter in the selected span:
extension of the interpolation frame preserves independence
because its nonzero full-rank minors remain nonzero.

\proofstep{Descent of the Weyl value test to the given row.}
Use the local models at all maximal ideals arising in this
way, equivalently those under integral eigencharacters occurring
in the selected span. If \(p^a\) gives the Weyl inclusion in
every such model, a maximal ideal detecting a nonzero class of
\(p^a x_0\) would contradict that inclusion. Thus \(p^a x_0\) has
Weyl values. Descent of the value tests gives the inclusion on
\(\mathcal C,\mathcal D\). In a localized coefficient basis the
condition is precisely nonnegative valuation on each coordinate.
It follows at the original level by the inverse-image description
above. This argument requires the same value and operator readings on both lattices.
\label{cor:lift-exponent-models-2}

\end{proof}

\paragraph{The two weights seen at a maximal ideal.}
At the chosen label put \(r=k_c+j+t+q\) and number the
coordinates from zero. Let \(A\) be the coefficient algebra
generated by the level matrices and the individual row-weight
projectors in the tensor. For a simple factor \(L\) of the
reduction, take decreasing real numbers \(\xi_{\tau,a}\) at
each label \(\tau\) and choose a joint weight \(\mu\) maximizing
\(\sum_{\tau,a}\xi_{\tau,a}\mu_{\tau,a}\) on \(L\).

If
\(\mu_{\tau,a}<\mu_{\tau,b}\) for \(a<b\), the transposition
of rows changes this sum by
\[
        (\xi_{\tau,a}-\xi_{\tau,b})
                     (\mu_{\tau,b}-\mu_{\tau,a})>0 .
\]
The row permutation preserves the weight set, so the
maximizer is dominant. The highest-weight theorem for this
coefficient algebra embeds \(L\) into the reduced dual Weyl
tensor of weight \(\mu\) and preserves the projectors. At one
label the theorem and the weight bounds give
\[
 \sum_a\mu_a=\sum_a\lambda_a,\qquad
 \lambda_{r-1}\le\mu_a\le\lambda_0,\qquad
 U_a=\lambda_a+r-1-a,\quad U_a'=\mu_a+r-1-a .
\]
In particular \(U,U'\) are strictly decreasing, of equal sum,
with entries in \([0,h_0]\), using \(\lambda_0+r-1\le h_0\).
\label{cor:lift-crystalline-depth-1}

Use the central Capelli generating polynomial
\(C_\nu(X)=\prod_{a=0}^{r-1}(X+\nu_a-a)\).
Its polarizations are coefficients of elementary row additions
followed by projection to the changed weights. The common degree
bound \(|\lambda|+|\mu|\) lets one use the same coefficient
operators on the reduced source and image. On a nonzero image
of \(L\) the two Capelli actions therefore have the same scalar
at every coefficient. The scalars are integral, so this gives
\[
       C_\lambda(X+r-1)\equiv C_\mu(X+r-1)\pmod p,
       \qquad \prod_a(X+U_a)=\prod_a(X+U_a')\ \text{in }k[X].
\]
Use the residue field \(k\) of the coefficient DVR here. Equality
of the two products gives the same multiset of residues, including
when several labels have the same residue.
\label{cor:lift-crystalline-depth-2}

Suppose the forms with values in \(L\) survive at the residual
ideal \(\mathfrak m\). At a chosen double coset \(c\), lift a
reduced value to the integral \(\mu\)-lattice. Its stabilizer
\(H_c\) has order invertible in \(O\). Averaging with the
level action \(R_c\) gives the lift
\[
        v_c^{\rm inv}=|H_c|^{-1}\sum_{h\in H_c} R_c(h)v_c .
\]
The reduction of \(v_c^{\rm inv}\) is the required value, since
the reduced value satisfies the stabilizer equations. Extend these
lifts by covariance. Let \(M_\mu\) be the integral weight-\(\mu\) forms and \(\bar v\) the surviving residual eigenform. Commutation of Hecke translations with reduction gives
\[
 \operatorname{Ann}(M_\mu)\subseteq\operatorname{Ann}(\bar v)=\mathfrak m.
\]
The finite torsion-free Hecke algebra argument
above now gives an integral characteristic-zero eigensystem
over \(\mathfrak m\), realized by a weight-\(\mu\) form since the
coordinate maps preserve independence. A residue-field embedding
aligning the two residual systems lifts to a local Galois
automorphism. This gives the congruence on every generator.
\label{cor:lift-crystalline-depth-3}

\input{lift-frame-galois}

\begin{lemma}[Identifying the marked factors in the companion reduction]
\label{lem:companion-residual-factors}
On the common coefficient extension of the two normalized
Hecke systems, the residual semisimplification is
\[
 \bar r_\mu^{\rm ss}\simeq \bar\rho^{\rm ss}\oplus\bigoplus_{s=0}^{t-1}\bar\chi_s.
\]
One may keep the marked occurrence of each scalar line in a
composition series of its actual reduced lattice.
\end{lemma}

\begin{proof}
At the good degree-one Frobenius classes, congruence of
the two Hecke systems gives
\[
 \det(X-r_\mu(g))\equiv
 \det(X-\rho(g))\prod_{s=0}^{t-1}(X-\chi_s(g))
                 \pmod {\mathfrak m_{\overline{\Q}_p}} .
\]
Indeed the reversed polynomial has Hecke coefficients times
signed integral norm powers and powers of the common twisting
character, which has unit values by compactness. For degree \(r\), reversing by
\[
 P(X)\mapsto X^rP(X^{-1})
\]
is a permutation of the coefficient list. It therefore
preserves the congruence after the common unit normalizations. The congruence
locus is closed because all coefficients are continuous and
the open unit ball in an ultrametric field is closed.
Degree-one Chebotarev density after finitely many omissions
extends it to every group element, also for the reciprocal
Frobenius convention. Choose a common finite coefficient field
and stable lattices.

For the two residual representations \(\bar R_1,\bar R_2\), Brauer--Nesbitt gives
\[
 \det(1-X\bar R_1(g))=\det(1-X\bar R_2(g))\ (\forall g)\ \Longrightarrow\ \bar R_1^{\rm ss}\simeq\bar R_2^{\rm ss}.
\]
A composition series therefore has the
prescribed line characters among its factors with their
multiplicities. Keep their positions in the filtration.
The block comparison takes these lattices over a common
finite coefficient extension with the marked characters
defined there. Compactness gives an invariant lattice
by taking the span of translates of one integral basis.
The coefficient field can be enlarged to contain the
labeled fields and the primitive \(p\)-th roots used in
the cyclotomic frame. This computation uses the original
residual filtration on the \(r_\mu\) side, with its
marked factors identified by the semisimplification.
\label{cor:lift-crystalline-depth-4}
\end{proof}

\input{lift-weight-exchange}

\paragraph{The localized tensor inclusion.}
For each embedding \(\tau\), let \(V_\tau=\nabla(\lambda_\tau)\)
be the dual Weyl module of highest weight \(\lambda_\tau\) over \(O\).
Write \(B_\tau=\Delta(\lambda_\tau)\subset V_\tau\) for its Weyl lattice.
Write \(V=\bigotimes_\tau V_\tau\).

\input{lift-pieri-cost}

\begin{lemma}[Descent from projected Weyl values]
\label{lem:localized-projected-weyl-values}
With the integral weights of the tensor \(\bigotimes_\tau
\nabla(\lambda_\tau)\), use the projected inclusions of
Lemmas~\ref{lem:rectangular-pieri-denominator}
and~\ref{lem:selected-pieri-depth}. The Weyl inclusion on the actual forms at
a selected residual ideal is obtained by clearing each factor
by its full exponent, or by the smaller allowable depth cost
at a larger-weight factor. Write \(d_\tau\) for the charge so chosen at each label, \(V\) for the unrestricted tensor lattice and \(D_\varnothing\) for the lattice using the Weyl factor at every label. After passing to the forms at the selected residual ideal \(\mathfrak m\), the inclusion is
\[
 p^{\sum_\tau d_\tau}\,\mathcal F(V)_{\mathfrak m}
 \subseteq \mathcal F(D_\varnothing)_{\mathfrak m},
\]
where \(\mathcal F(L)\) denotes the module of forms with values in \(L\).
\end{lemma}

\begin{proof}
For a set \(S\) of
unrestricted labels, let \(D_S\) be the image of the tensor
using \(V_\tau\) on \(S\) and \(B_\tau\) elsewhere. Each
\(D_S\) is stable under the coefficient algebra and its forms
are stable under the Hecke translations. The exact exponents
\(e_\tau\) give \(p^{\sum e_\tau}V\subset D_\varnothing\)
by multiplication on pure tensors. Thus \(V/D_\varnothing\)
has finite length. Every simple factor of a subquotient of
this quotient occurs in the reduction of \(V\) modulo the
uniformizer \(\varpi\). Indeed, multiplication by a power of \(\varpi\) gives the surjection onto each successive quotient:
\[
 V/\varpi V\ \twoheadrightarrow\ \varpi^a V/\varpi^{a+1}V,\qquad [v]\longmapsto[\varpi^a v]\quad(a\ge0).
\]

Write \(P_{\tau,w}\) for the integral weight projector,
\(i_\tau\) for the depth allowance and \(c_\tau\) for the
integral bounded-depth cost just obtained. For
\(a\notin S\), set \(D=D_S\) and \(D'=D_{S\cup\{a\}}\).
If the full exponent is used at \(a\), then
\(p^{e_a}D'\subset D\) on pure tensors. At an index with a larger weight
where \(c_a<e_a\), form the actual intermediate lattice
\[
                          E=D+p^{c_a}D'.
\]
For any weight \(w\) at depth at most \(i_a\) we have
\(p^{c_a}P_{a,w}V_a\subset B_a\). Applying that projector
in just factor \(a\) shows \(P_{a,w} E\subset D\).

In a simple factor \(L\) of \(E/D\), the projector at its
highest weight \(\mu_a\) has nonzero image. This follows
from the highest-weight line in the coefficient classification.
If the highest weight had depth at most \(i_a\), the integral inclusion would give
\[
 P_{a,\mu_a}E\subseteq D\quad\Longrightarrow\quad P_{a,\mu_a}(E/D)=0,
\]
contradicting the nonzero image on \(L\). Thus all simple factors here have
excess depth. By the preceding crystalline estimate their
localized forms vanish. This argument requires stability of
the intermediate lattices and the integral action of the
projector on the full tensor.

Take a composition series for \(E/D\). At each successive
quotient map a form to the form obtained by projecting its
values to the simple factor. This map is Hecke-equivariant.
Filter \(E\) by its series \(D=D_0\subset\cdots\subset D_b=E\), and denote the projection at the last step by \(\pi_b\). The vanishing for the simple factor reads
\[
 (\pi_b\phi_b)_{\mathfrak m}=0,\qquad \pi_b:D_b\to D_b/D_{b-1}.
\]
There is \(s_b\notin\mathfrak m\) killing that image before
localization, so \(\phi_{b-1}=s_b\phi_b\) takes values in
\(D_{b-1}\). Composing the successive denominator clearings gives
\[
 (s_1\cdots s_b)\phi_b\ \text{with values in }D_0=D .
\]
Multiplying these denominators down the series proves the
localized inclusion of the forms with values in \(E\) into
those with values in \(D\).

In particular \(p^{c_a}\) moves
the forms on \(D'\) into the forms on \(D\) after localization.
Inducting over the unrestricted labels gives
\[
 p^{\sum_{a\ {\rm thick}}\min(e_a,c_a)+
       \sum_{a\ {\rm other}}e_a}
            M_{\mathfrak m}\ \subset\ (M_{\rm W})_{\mathfrak m}.
\]
All denominators used are units in this localization.
The simple-factor argument works separately at each integral
joint character over the ideal. The surviving bound is the
same for every such character in the prime range under
consideration. This proves the passage from the projected
coefficient inequalities to the lattice of actual forms.
\label{cor:lift-graded-inclusion-1}
\end{proof}



\begin{lemma}[Least exponent on the original Weyl value pair]
\label{lem:weyl-least-exponent-local}
For the coefficient models of Proposition~\ref{prop:local-weyl-value-inclusion},
let \(B_j(p)\) be the bound from Section~\ref{subsec:lift-exponent-primes}
for the weight record of lift \(j\). For the ramification index \(e_w\) and the least local exponent,
\[
 J_{j,w}/e_w\le B_j(p)\qquad(w\mid p).
\]
This normalized upper bound descends through the coefficient extension.
\end{lemma}

\begin{proof}
The exponent in these models is computed with the actual
Frobenius and conjugation labels of the frame. In the indicated
range the two window correction counts vanish by the shape
calculation of Section~\ref{subsec:lift-exponent-primes}.
The depth and all other terms are therefore independent of the Frobenius and conjugation labels. One power applies at every maximal ideal in the cover. Outside the weight range use the sum of the full Weyl exponents, with exponent zero above height. The small and excluded primes likewise use the full Weyl sum. At a remaining prime the contribution of each larger-weight index is the minimum of its full exponent and its integral coarse correction. These comparisons fix the
tuple at \(w\) before taking the least exponent and keep the
original weight record for the prime sum. In compatible Smith bases the two lattices to be compared are
\[
 \mathcal C_w=\bigoplus_i O_we_i,\qquad \mathcal D_w=\bigoplus_i O_w\varpi_w^{d_i}e_i\subseteq\mathcal C_w.
\]

Then
\[
     J_{j,w}=\max_i d_i,\qquad
     p^{B_j(p)}\mathcal C_w\subseteq\mathcal D_w
             \ \Longrightarrow\ d_i\le e_wB_j(p)\quad\text{for all }i.
\]
The same comparison applies if the surviving eigensystem
requires an extension of the coefficient field. Over the extension the exponents and
the absolute ramification index both multiply by its relative
ramification index.
\label{cor:lift-exponent-models-3}
\end{proof}


\subsubsection{Summing the weight exponents}\label{subsec:lift-exponent-primes}
\hypertarget{source-lift-exponent-primes}{}

The integral inclusion gives local exponents indexed by highest weights. The sum over primes depends on which differences between shifted weights can be divisible by the prime. We keep the head weights and the consecutive block distinct in making this count.

\begin{established}[Heights, cutoff and exceptional interval]\label{hyp:weyl-prime-range}
Use the weight record for one lift. In this subsection \(r\) is
its coefficient rank, \(k_c=k-1\) is the head size, \(t\) is its
consecutive block length, and \(S,\eta\) bound the widths of
the head weights and the displaced weights, respectively.
The shifted weights in that
record are the highest weights plus the usual staircase. Their
height is bounded by \(h_0=\lfloor h\rfloor\), where
\(h=b+S+\eta+t+1\). Write \(X'\) for the auxiliary cutoff,
\(b'=xb/X'\), \(f=2n\), and \(w''=b/(nX')\).

Let \(N=k+m(t+1)\) be the final rank, and let \(\ell_*\)
be the auxiliary threshold
\[
 \ell_*=\frac{mb}{nt^2(N+t)z^2+kn^2(S+t+z)} .
\]
Put \(K_*=\lfloor 2b/\ell_*\rfloor\). Define the small-prime
threshold \(y\) and the exceptional interval \(\mathcal I\) by
\[
 y=\frac{mbz}{k(m+t)\log(3h)},\qquad
 \mathcal I=\left[b-\frac{t-1}{2}-K_*\,,\,
       b+S+\eta+1+\frac{t-1}{2}+K_*\right].
\]
Uniformly beyond a fixed threshold one has \(z\ge4\),
\(y>8\), \(y\le b\), \(r+1<p\) and \(S+k<p\) for \(p>y\),
\(h\le2b\), \(h+r\le3b<y^2\), and \(\mathcal I\subset(0,y^2)\).
Furthermore \(k(S+1)\le2kS\le2w''\) since \(S\ge1\). \label{cor:lift-exponent-primes-1}
\end{established}
\begin{proof}
Here \(b=\exp(x^\epsilon)\), \(z=\log b\). The ranks \(k,m,t,r\) are polynomially bounded in \(x\), and \(S,\eta\le x^{O(1)}\exp(x^{\epsilon/4})=b^{o(1)}\) by the weights of Lemma~\ref{lem:window-parameter-choice}. Hence \(h=(1+o(1))b\), \(K_*=b^{o(1)}\) by its displayed denominator, and \(\log(3h)\ge z\). As \(m\le k(m+t)\), one has \(b^{1-o(1)}\le y\le b\). In particular \(y^2>3b\), \(r+1,S+k<y\), and both endpoints of \(\mathcal I\) equal \(b+o(b)\). Finally \(kS=b^{o(1)}\) and \(w''=b^{1-o(1)}\), which proves the last comparison.
\end{proof}

For an integral highest weight \(\lambda\) in rank \(d\), let
\(e_p(\lambda)\) be the least exponent carrying the dual Weyl
lattice into the Weyl lattice over \(\Z_{(p)}\). Determinant
translation of both lattices leaves it fixed.

\begin{lemma}[Cutting a Weyl lattice at a consecutive block]
\label{lem:consecutive-weyl-cut}
The exponent satisfies the following estimate at a consecutive cut into ranks
\(a\) and \(d-a\). If the full shifted weights lie in an interval
of integer width \(W\), then
\begin{equation}\label{eq:lift-weyl-cut}
 e_p(\lambda)\le e_p(\lambda_{<a})+e_p(\lambda_{\ge a})
                   +a(d-a)\lfloor\log_p W\rfloor .
\end{equation}
Use zero for the logarithmic term at \(W=0\). Here the
restricted weights use the staircase in their own ranks.
Thus restriction changes the shift by a constant on each block
and preserves its width. \end{lemma}

\begin{proof}
We require a rational operator for the Levi blocks that realizes the two integral inclusions with scalars \(c_1,c_2\). In the adjoint action on endomorphisms,
apply the successive rank-one projections for the cross-block
pairs. Grid invariance gives an endomorphism of the full
representation. Its scalar on the highest weight vector is the
inverse of the product of the cross-block shifted differences.
Schur's lemma gives that same scalar on the module. In the
integral grid comparison the beta denominator for each pair
is bounded by the denominators of
\(i!j!/(i+j+1)!\) for \(i+j+1\le W\). Put \(m=p^a\).

At this power the factorial valuation contributes at most one, since
\[
 \left\lfloor\frac{i+j+1}{m}\right\rfloor-\left\lfloor\frac{i}{m}\right\rfloor-\left\lfloor\frac{j}{m}\right\rfloor=\left\lfloor\frac{(i\bmod m)+(j\bmod m)+1}{m}\right\rfloor\le1 .
\] Hence a factor
\(p^{\lfloor\log_p W\rfloor}\) per pair, times the two block
inclusion scalars, clears the full inclusion once the grid is
multiplied by its highest-weight scalar denominator. This gives the grid identities for the rank-one line, its scalar and the integral denominator. We use those grid identities
on the Weyl lattices with the stated width condition.
\label{cor:lift-exponent-primes-2}

\end{proof}

\begin{proposition}[Sum of least integral exponents on the selected record]
\label{prop:selected-weyl-prime-sum}
With the heights, cutoff and exceptional interval of
Lemma~\ref{hyp:weyl-prime-range}, for
each selected lift the original rational-prime contributions
\(B_j(p)\) in \eqref{eq:division-least-sum} have
\[
 \sum_p B_j(p)\log p\le30|T_2|mbz^2+2mb.
\]
The estimate also holds for a restricted set of rational primes
satisfying Lemma~\ref{hyp:weyl-prime-range}. The constants and threshold are
uniform over the selected lifts and coefficient fields.
\end{proposition}

\begin{proof}
\proofstep{The width estimate and the small primes.}
Rank zero and rank one have exponent zero. Iterating
\eqref{eq:lift-weyl-cut} with one rank-one block gives
\(e_p(\lambda)\le\binom d2\lfloor\log_p W\rfloor\), and hence
zero for \(p>W\). At an index with a smaller weight, the width satisfies \(W_\tau\le2b'\).
If \(W_\tau>1\), the bound
\(\pi(Y)<1.26Y/\log Y\) for \(Y>1\), used here with
\(Y=W_\tau\), follows from Rosser and Schoenfeld's explicit
estimate \cite{rosser-schoenfeld-explicit}. This gives
\[
 \sum_p e_p(\lambda_\tau)\log p
 \le \binom r2\log W_\tau\,\pi(W_\tau)
 \le 1.3r^2 b'.
\]

If \(W_\tau\le1\), this sum is zero. At an index with a larger weight the
crude pointwise bound is \(\binom r2\log h\) for the
exponent times \(\log p\). For \(p\ge2\eta+t\), cut off the
head of size \(k_c\). Its weights lie in an interval of width
at most \(S\). The
remaining weights have width strictly less than \(p\) under
the weight shape conditions. The lower block therefore has
exponent zero, and \eqref{eq:lift-weyl-cut} gives the bound
\[
 e_p(\lambda_\tau)\log p\le
 k(m+t)\log(3h)+k^2\log(S+k).
\]

In particular, the head by itself has zero exponent when
\(p>S+k\). At \(p\le y\), the displayed bound times \(p\)
is at most \(2mbz\) since \(k\le m+t\), \(S+k\le3h\),
and \(y\,k(m+t)\log(3h)=mbz\). In the other case
\(p<2\eta+t\), the size conditions \(t\le\eta\), \(r\le2t\)
and \(\eta\le b/(t^3 f k^3)\) give the same bound from
\(e_p(\lambda_\tau)\log p\le4t^2 z\) and
\(m f t k^3\ge6\). For an index with a larger weight the small-prime contribution is at most \(2mbz^2\), since
\[
 \sum_{p\le y}\frac1p\le\sum_{j=2}^{\lfloor y\rfloor}\frac1j\le \log y\le z .
\]

\proofstep{The exceptional primes above the cutoff.}
Above \(y\), an exceptional prime divides an integer of
\(\mathcal I\). Each such integer has at most one distinct
prime factor above \(y\), as all these integers are positive
and less than \(y^2\). Hence the number of excluded primes
above \(y\) is at most \(S+\eta+t+2K_*+1\), plus the
number of specified primes. The conductor of the auxiliary field factors that
are at most \(y\) are already charged in the small-prime range.
Indeed the additional small factors divide an integer at most
\(n^2\le kt\), and \(kt\le b/(4kt)\le y\) by the large
parameter conditions.

The four listed primes \(q,d,p_{\rm de},
p'_{\rm de}\) suffice in the specified set after this removal.
Allowing five gives an exceptional cost at most
\begin{equation}\label{eq:exponent-exceptional-prime-cost}
 (S+\eta+t+2K_*+6) f r^2 z,
\end{equation}
because each of the \(f\) label contributions satisfies
\[
 e_p(\lambda_\tau)\log p
 \le \binom{r}{2}\log h\le r^2 z
\]
whenever its width exceeds one. For a smaller weight use its
width, at most \(2b'\) and at most \(h\), in place of \(h\);
widths at most one have zero exponent.
The vanishing above height holds on the full record because the
smaller widths are at most \(h\) as well.
\label{cor:lift-exponent-primes-3}

\proofstep{The restricted exponent at the remaining primes.}
At primes \(y<p\le h\) outside the excluded interval and the
specified set, use the local weight-elimination formula with
the depth estimate of Lemma~\ref{lem:selected-pieri-depth}. Let \(j_\tau,q_\tau\ge0\), \(j_\tau+q_\tau\le m\),
be its two counts outside the head and consecutive block. At an index with a larger weight, the head interval \([x_-,x_+]\) and the top weight \(B_\tau\) of the consecutive block satisfy
\[
 0\le B_\tau<x_-\le x_+\le h,\qquad x_+-x_-\le S.
\]
Put \(d_\tau=x_--B_\tau\).

The head exponent in this formula is zero. The correction counts from the windows
around the consecutive blocks of radius \(\lfloor h/(p-1)\rfloor\)
are zero for the original and the Frobenius-labelled weights.
In the signed coordinates of the shape condition, center the
comparison at the midpoint of the consecutive block. Write \(c\) for this midpoint, \(z_0\) for a point of the
enlarged consecutive window, and \(H_p=\lfloor h/(p-1)\rfloor\).
The window and the signed displacement bounds give, for a displaced
weight \(u\),
\begin{equation}\label{eq:weyl-short-window-distance}
 |z_0-c|\le (t-1)/2+H_p,\qquad
 0<|u-z_0|\le\eta+t+H_p<p .
\end{equation}
The strict lower bound uses the lower displacement in the weight
record. Thus no displaced weight is congruent to \(z_0\) modulo
\(p\).

For the head weights there are two further tests. At an
index with a larger weight the signed difference, with the sign
prescribed by the shape, belongs to \(\mathcal I\).
At a top or an index with a smaller weight the absolute
difference lies in \((0,p)\). Summarizing the exclusions,
\[
 \begin{aligned}
 d\in\mathcal I\cap\mathbb Z&\quad\Longrightarrow\quad p\nmid d
       &&\text{by the choice of }p,\\
 0<|d|<p&\quad\Longrightarrow\quad p\nmid d
       &&\text{by size}.
 \end{aligned}
\]
Reflection preserves the absolute differences in \eqref{eq:weyl-short-window-distance} and the nondivisibility tests, so the reflected shape has the same vanishing correction counts.
Thus the remaining depth in the real bound for
this formula is \(q_\tau h\). Let \(I_p\) indicate whether
\(p\mid u-B_\tau\) for an integer \(u\in[x_-,x_+]\).

Write \(r_p\) for the least positive residue of \(d_\tau\) modulo
\(p\). The term \(G_p\) is \(\min(k,t,q_\tau h/r_p)\) if
\(r_p<d_\tau\), and zero otherwise. The larger-weight
part of the rational-prime summand, before multiplying by
\(\log p\), is bounded by
\[
 (q_\tau+2j_\tau)\left(\frac{h+r}{p-1}+6\right)
   +\frac{kS}{p}+k I_p+G_p
   +\frac{q_\tau h}{p}(1+\log(h/p)).
\]
Two elementary sums give
\[
 \sum_{y<p\le h} I_p\log p\le(S+1)\log h,\qquad
 \sum_{y<p\le h} G_p\log p\le q_\tau h(1+\log h)\log h.
\]

\proofstep{Summing the two residue terms.}
For the first, map each prime counted by the indicator to one
integer \(u\in[x_-,x_+]\) such that \(p\mid u-B_\tau\). Two distinct primes cannot give the same \(u\). Otherwise
\[
 y^2<pp'\le u-B_\tau\le h<y^2 .
\] The interval contains at most \(S+1\) such integers.
For the second, every prime with \(G_p\ne0\) divides the positive
integer \(d_\tau-r_p\le h\). Thus \(p\mapsto r_p\) is injective
by the same product inequality. Sum \(q_\tau h/r\) over
\(1\le r\le\lfloor h\rfloor\), and bound \(\log p\) by \(\log h\).
\label{cor:lift-exponent-primes-4}
\proofstep{Elementary weighted prime sums.}
The other sums use \(\sum_{p\le h+r}\log p\le(h+r)\log4\)
and \(\sum_{p\le h+r}(\log p)/p\le\log(h+r)+\log4\).
At any integer \(Y\ge2\), sum \(Y/p\le\lfloor Y/p\rfloor+1\) with weights \(\log p\). The contributions of the two sums give
\[
 Y\sum_{p\le Y}\frac{\log p}{p}\le \sum_{p\le Y}\lfloor Y/p\rfloor\log p+\sum_{p\le Y}\log p\le\log(Y!)+Y\log4 .
\]
The floor terms are bounded by \(\log(Y!)\) and the unit terms
by \(Y\log4\). Division by \(Y\) and taking floors in the
cutoff give the bound for real cutoffs as well.

For \(p>y>8\) one can replace \(1/(p-1)\) by \(9/(8p)\).
Here \(h+r\le3b\), \(\log h\le z+0.7\), and
\(\log(h+r)+\log4\le z+2.49\). The two terms with factor
\(k\) contribute at most \(2k(S+1)(z+2)\) since
\(S(z+2.49)+(S+1)(z+0.7)\le2(S+1)(z+2)\).
The remaining four terms together cost at most \(mb\) times
\[
 \tfrac{27}{4}(z+2.49)+50.04
   +2(z+1.7)(z+0.7)+2(z+1.7)(z+2.49)\le18z^2.
\]
For the last inequality the difference has value and derivative
positive at \(4\) and has positive quadratic coefficient.
These bounds apply to every subset of the indicated primes.

\proofstep{Absorption of the three tails.}
After summing the residue terms, the remaining tails involve different sources of divisibility. Their bounds must all fit the exponent allowance for this single lift.

Let \(T_2'\) be the number of indices with larger weights in the record.
Adding the three prime ranges gives the term
\(30T_2' mbz^2\), with room from \(2+18\le30\). The remaining
terms are
\begin{equation}\label{eq:exponent-prime-tails}
 2T_2' k(S+1)(z+2)+1.3 f r^2 b'
       +(S+\eta+t+2K_*+6) f r^2 z .
\end{equation}

To bound them let \(A\ge1\) be the polynomial size majorant
for the selected parameters. The active size rows give
\[
 \begin{gathered}
 f r^2\le4A^3,\quad kS\le w'',\quad S+\eta\le b/X',\quad
 K_*+t\le A^4x^2(z^2+S),\\
 4(z+2)\le X',\quad 5.2A^3x\le X',\quad
 40 A^7x^2 z\le X',\quad 40 A^7x^2 z^3\le b .
 \end{gathered}
\]
The bound for \(K_*+t\) follows by expanding \(2b/\ell_*\)
and using \(n,k,t\le A\), \(N\le Ax\) and \(m\ge2\).
The last four inequalities hold eventually for the common
polynomial majorant at \(X'=\exp(x^{\epsilon/3})\),
\(z=x^\epsilon\), \(b=\exp z\). The first term of
\eqref{eq:exponent-prime-tails} is at most
\(8(z+2)b/X'\le2b\), using \(T_2'\le f=2n\).
The second is at most \(5.2A^3xb/X'\le b\).

For the third use
\(S+\eta+t+2K_*+6\le5A^4x^2(z^2+b/X')\).
After multiplication by \(fr^2z\) the two terms are each
at most \(b/2\). Adding these three remainders and using \(m\ge2\) gives
\[
 2b+b+2(b/2)=4b\le 2mb .
\] This proves the rational-prime bound in
\eqref{eq:division-least-sum} with \(T_2'\le|T_2|\).
It is uniform over the chosen weights, lifts, coefficient fields
and subsets of rational primes in those size rows.
\label{cor:lift-exponent-primes-5}
\end{proof}


\paragraph{A common class denominator.}
For each lift \(j\) and tuple \(t\), the integer \(D_{G,j,t}\)
used above is the order of the row level subgroup times the
least common multiple of its class stabilizer orders. Take
\(D_*=\operatorname{lcm}_{j,t}D_{G,j,t}\).
The selected row bundle identifies the level as a subgroup of
the common finite group \(G\). Its stabilizer orders divide
the positive integer \(D_{\Gamma,u_0}\) of the distinguished
word \(u_0\). These divisibilities hold for every tuple in the
bundle, including all \(\mu(w)\). It follows that
\[
                   D_*\mid |G|D_{\Gamma,u_0}.
\]
This argument uses a single common multiple even when the
number of tuples grows. In the active fixed-scale comparison
the value-lattice multiplier is \(s=1\).

The row finite group has
\(\log|G|\le 2fr^2\log(2g_0)\).
Indeed each of its residue factors is bounded by the order of
the corresponding unitary or general linear matrix group.
In the unitary case its order is
\[
 q^{r^2}\prod_{i=1}^r(1-(-q)^{-i})\le (2q)^{r^2}.
\]

The local degrees over \(g_0\) and the selected level size give
the stated group bound. The size conditions give \(2r^2\le b\)
and \(\log(2g_0)\le2z^2\), so \(\log|G|\le2\mathcal E\).
The stabilizer comparison has the bound
\[
 \log D_{\Gamma,u}\le C_\Gamma fN_m\log(fN_m)
\]
for every
word \(u\), where \(N_m\) is the final unitary rank. Here the
stabilizers act faithfully on rational spaces of dimension
at most \(2fN_m\) by restriction of scalars. Minkowski's
finite-subgroup bound gives a common multiple independent
of the lattice.

For rational dimension \(n\), its exponent
at a prime \(p\) is
\(\sum_{a\ge0}\lfloor n/(p^a(p-1))\rfloor\).
The logarithm of that common multiple is \(O(n\log(2n))\):
only \(p\le n+1\) occur, and summing the exponents is
bounded by \(n\sum_{p\le n+1} p\log(p)/(p-1)^2\).
Chebyshev's bound and partial summation give
\[
 \sum_{p\le u}\frac{\log p}{p}=\frac{1}{u}\sum_{p\le u}\log p+\int_2^u\frac{\sum_{p\le t}\log p}{t^2}\,dt=O(\log(2u))\quad(u=n+1).
\]
\cite{serre-finite-orders}
The conductor and rank size condition gives
\(1\le fN_m\le3\sqrt b\). In particular
\(\log(fN_m)\le3z\) and
\(fN_m\log(fN_m)\le9bz\le9\mathcal E\).
Thus on the selected row bundle
\begin{equation}\label{eq:division-class-bound}
                \log D_*\le(2+9C_\Gamma)\mathcal E .
\end{equation}
The coefficient field and the number of words, lifts and
tuples cause no additional multiplier in this bound.
\label{cor:division-length-9}


\subsection{Real scales and compatible translates}
\label{subsec:window-points}

Let the primitive target and its weights be fixed as in
Proposition~\ref{prop:arithmetic}. At a real scale \(x\), the
base data consist of the chosen auxiliary field, the ranks,
and the two weight functions. The pointwise estimates can enter the Fourier average only if every required translate has local lattices and period maps over these base data. Construction~\ref{hyp:family-prescribed-realization} selects the characters at this translate, retains the reference tuple, and forms the member used for both index estimates.

\begin{definition}[Completed member]\label{def:completed-member}
A \emph{completed member} at scale \(x\) and translate \(u_0\)
of the base family with host \(E\) from Definition~\ref{def:family-original-host} consists of the realized system of theta spaces, row and
column lattices and period maps at that translate, with the
level exponents, local divisors and colengths read from this
system. Here \(u_0\in A\) translates the weights by the group law of
Definition~\ref{def:translate-scale}. Its Fourier average uses
\(u_0\in A_0\). The positive constant relating the original and rescaled real scales is \(\kappa_{\rm sc}\) from Lemma~\ref{lem:rescale-host-budget}. Compatibility means that the field and weight
identifications return the base data at \(x\), that each
place is assigned its prescribed tuple and level, and that
the base and refined period readings use the rational
isomorphisms of Section~\ref{subsec:completed-division}.
\end{definition}
Completion is needed because the two division bounds
use both the analytic period coordinates and the integral
indices, and these must come from the same system.

\paragraph{Successive bounds on the constants.}\label{par:window-successive-bounds}
Let \(\epsilon_0>0\) be a common upper bound for the small
parameter in the finitely many estimates. Choose
\(0<\epsilon<\min(1/53,\epsilon_0)\). Write \(v\) for the
list of first parameters, including a preliminary choice,
and \(w\) for the list of middle parameters, including the
proposed height exponent. A system of sufficient bounds has
a preliminary bound \(p_0\), a bound \(A(v)\) for the middle
parameters, and a bound \(B(v,w)\) for the final parameters.
Bounds from several estimates can be combined by
pointwise maximum at each stage. This is also how one keeps
the thresholds for the two kinds of arithmetic errors and the
height limit of Proposition~\ref{prop:marked-frame-height} in the
same system.

\begin{lemma}[Simultaneous cutoffs and indexing of the columns]\label{lem:window-parameter-choice}
For the selected base family with \(p_0,A(v),B(v,w)\) as in
Paragraph~\ref{par:window-successive-bounds},
one can choose the preliminary, middle and final parameters in order
so that both division comparisons and the height bound use the
sequence-prime exponent \(C_{Tu}=c_T(\epsilon,9)\).
After the choice, for every sufficiently late scale in \(S\) there is a single column index map
\[
 \{1,\ldots,N\}\lhook\joinrel\longrightarrow\{0,1\}^{k-2},
\]
fixed before choosing any translate.
\end{lemma}

\begin{proof}
\proofstep{Evaluate the bounds at the actual sequence-prime exponent.}
The base data fix the parameter \(C_{T2}\) for the second
index set and a common upper bound \(C_8\ge1\) for the
conductor parameters and the interval exponent \(J\).
The standard list uses \(C_t=401\) for the original cohomology
cutoff and \(c_{TO}=377\) in the rank formula of
Lemma~\ref{lem:rank-window}. It also uses \(C_{YL}=9\),
\(C_p=1\) and \(e_P=12\).
The exponent \(C_{Tu}\) for the prime assigned to a sequence
will have the value \(t_*=c_T(\epsilon,9)\) supplied by the
construction.

To impose the final system on a
list \(v\) temporarily made with the standard value
\(C_{Tu}=25\), substitute \(t_*\) in the entry for that exponent
of \(v\) in both \(A\) and \(B\). In notation that exhibits the only substitution, this list is
\[
 v_* = v[C_{Tu}\leftarrow t_*],\qquad t_*=c_T(\epsilon,9).
\] Leave the separate sufficient bounds for
constructing the point on the standard list. If their three
functions are \(p_1,A_1,B_1\), denote the chosen first, middle
and final bounds by \(p,h,l\), respectively. Set
\[
 p=\max(1000,p_0,p_1),\qquad
 h=\max\{1,A(v_*),A_1(v)\},\qquad
 l=\max\{1,B(v_*,w),B_1(v,w)\}.
\]
Here \(v\) is evaluated after setting \(p\), and \(w\) after
setting all free middle entries equal to \(h\). Set the free
last entries equal to \(l\).

Entries encoding inverse lower
constants in the standard list use \(1/l\).
The middle list reading the frame limit has both its prescribed
height parameter and the code exponent equal to \(h\). The
outer entry that reads \(\max(C_8,J)\) is \(C_8\) on the
constructed base data. Restoring \(C_{Tu}=t_*\) now gives precisely
the arguments used in the final bounds. After setting the free middle entries to \(h\), the same modified list occurs in both bounds:
\[
 h\ge A(v_*),\qquad l\ge B(v_*,w).
\] The point construction meanwhile uses
\(h\ge A_1(v)\), \(l\ge B_1(v,w)\) with its own standard list.
Neither pair requires evaluating a bound at an entry chosen later.
\label{cor:window-points-1}

\proofstep{Choose the numerical cutoffs in order.}
There remain two numerical cutoffs \(c_D,c_h\).
If \(d_i\) bound the first cutoff and \(h_i(c_D)\) bound the
second, set
\[
 c_D=\max(0,\max_i d_i),\qquad
 c_h=\max(402,\max_i h_i(c_D)).
\]
In particular a dependence of \(h_i\) on \(c_D\) creates no
dependence in the reverse direction. The scale threshold may now depend on all these
numbers. The construction of the base data gives a set
\(S\) containing arbitrarily late intervals
\([Y,Y^{3/2}]\). At any sufficiently large real \(x\in S\)
take its ranks, field, distinguished auxiliary prime, and
two weight functions from Section~\ref{subsec:frobenius-selection}.
Denote the weight functions by \(a(\cdot),b(\cdot)\), to
distinguish them from the numerical scales \(a,b\). They
satisfy the window conditions before selecting a translate.
\proofstep{Encode the columns before choosing the translate.}
The two estimates must also use identical column labels at every translate.

The two auxiliary prime signs vary independently by the Chinese
remainder theorem. Therefore the degree entering the number of columns is
\[
 [E:\mathbb Q]=4n,\qquad f=[E^+:\mathbb Q]=2n,\qquad n=|A|.
\]
Put \(L=x^\epsilon\). Let \(C_\Gamma>0\) be the fixed size
constant. Define the size parameter \(\Gamma\) and the
number \(N\) of sequence indices by
\[
 \Gamma=\exp(C_\Gamma m^5fL),\qquad
 N=\lceil 3\Gamma^2x^h\rceil .
\]

The rank bounds give \(8C_\Gamma m^5 f L\le k\). Since
\(h\) is fixed and \(k\ge x^{9/2+10\epsilon}\),
\[
 \log N\le\log4+2\log\Gamma+h\log x
       \le \log4+k/4+h\log x < (k-2)\log 2
\]
eventually, uniformly in the base data. Choose an injection
of the \(N\) indices into \(\{0,1\}^{k-2}\) and use it for
all translates. Call the sequence of these prescribed steps for an index a
\emph{word}. The bits prescribe the shifts of the interior
upper coordinates from the preceding lower coordinates. The top and bottom are
specified separately by the two weight functions.
\end{proof}

\paragraph{The single-translate realization used here.}
Fix a translate \(u_0\in A\), keeping the
field, ranks, weight functions and binary sequences fixed.
For each sequence index \(u\), write \(T_u\) for its assigned
split prime. Before choosing it, the construction needs two
positive dimension factors \(d_{u,*},D_*\) and a constant
\(c_s>0\), uniform in the sequence, such that
\[
  \log\!\left[\left(D_*/(d_{u,*}\Gamma^2)\right)^{1/(fk)}\right]
      \ge c_s mL
\]

The ranks determine the allowed level interval at each assigned prime.
Fix levels in those intervals. The following construction gives the
spaces and comparison maps on which the local divisors and heights
will be measured.

\begin{rowconstruction}[The signed rows at a prescribed list of levels]
\label{hyp:family-prescribed-realization}
Fix the numerical choices of Lemma~\ref{lem:window-parameter-choice}.
At a scale in its range fix a translate and any admissible list
\((T_u)_u\) of split word primes. Take the marked Hermitian lattices,
two split type markings and deletion characters of
Definition~\ref{hyp:initial-adelic-frame}, with the regular weights
of Proposition~\ref{prop:linked-type-multiplicity}, even column
count \(n_V\ge8fr+4\), and the word levels of
Section~\ref{subsec:selected-word-level}. The dimensions in this stable range satisfy
\[
        [E:\mathbb Q]=2f,\qquad n_V\in2\mathbb Z_{>0},
                     \qquad n_V\ge 4[E:\mathbb Q]r+4 .
\]
The finite class tests, their harmonic weights and the splitting characters are those of
Proposition~\ref{prop:signed-theta-data}.
At a served index use the signatures occurring on the initial lower row as in Lemma~\ref{lem:initial-row-signature-tests}. It gives integral conjugate duality and the residual tests of Lemma~\ref{hyp:simultaneous-retained-twist}. Each indicated jump \(a\) satisfies
\[
       a<-1\ \text{ or }\ a>t,\qquad r\ge100,\qquad
                                     \min(g_0,g_1)\ge2r .
\]
For these data the order of construction is
\[
 (1,\delta,\delta')\ \longmapsto\
   \text{initial signed spaces}\ \longmapsto\ (\varphi_i)_i
     \ \longmapsto\ \text{final member }P'.
\]

The initial sections and their arithmetic coordinates have the same
norm on every prescribed integral test. In the reconstruction the
embedded fields \(K_s,K_0\) and the reference tuple with its periods
are retained. Each selected nonreference character occurs on the
corresponding restricted Hecke space. The bounds for the two
division errors both use \(P'\).
\end{rowconstruction}

\begin{proof}
\proofstep{Construct the initial signed spaces.}
Use Lemmas~\ref{lem:original-marked-pel-field}
and~\ref{lem:marked-pel-base-levels} to put the marked moduli spaces
over the original embedded field. A row argument determines a
polarized point, its Hodge frame and the principal theta line
\((1,\lambda w_0)^*\mathcal P\).
Lemmas~\ref{lem:marked-period-domain} and~\ref{lem:placed-row-theta}
give these objects also at each refinement. The finite translations
and harmonic derivatives on its class representatives define the
sections of Lemma~\ref{lem:marked-class-section-descent}. At a
positive-dimensional stage the constant terms vanish by
Lemma~\ref{lem:marked-class-section-cuspidal}.
The full cusp space has the lowest-type reading of
Lemma~\ref{lem:marked-full-holomorphic-row}. The algebraic and
analytic class maps on its signed subspace have the same kernel
by Lemma~\ref{lem:marked-class-kernel-quotient}.

Write \(b_1,\ldots,b_d\) for a basis of the rational cusp space over
its field of definition. For the framed point \(y\) used in an
integrality test, the norm of the coordinate row is
\[
 \left\|(z_j)_j\right\|_y=\left|\sum_j z_j b_j(y)\right|_y .
\]
Lemma~\ref{lem:constructed-signed-coefficient-row} gives its full
integral ball, its intersections with the type subspaces and the
measured multiplicity lattice. This test includes semistable
reduction. Lemma~\ref{lem:marked-semistable-local-ball} proves
Hecke stability on it outside the residue characteristic.

The two signs use the same finite class inputs with their respective
coefficient maps. The common comparison of the vacuum and the
determinant normalization is made in
Lemma~\ref{lem:placed-row-tuple-phase}. Lemma~\ref{lem:placed-row-reading-phase} installs the resulting unit factors in the complex reading, keeping the algebraic sections over their original field. For each fixed stage,
sign, coefficient place and test, Lemmas~\ref{lem:marked-bounded-denominator}
and~\ref{lem:marked-prime-to-level-denominator} give
\[
  \exists b_v\ge0\quad\forall Q,M,\delta\quad
            v'(e_h(x_{Q,M,\delta}))\ge-b_v\quad(1\le h\le d),
                 \qquad v'(\pi_v)=1
\]
in the admitted ranges of Proposition~\ref{prop:signed-theta-data}.
If the torsion labels require a finite local extension \(K''/K'\),
then \(K'\cap\pi_v^{-b_v}\mathcal O_{K''}
=\pi_v^{-b_v}\mathcal O_{K'}\). The separating matrix and its
denominator are fixed before \(Q,M,\delta\).

\proofstep{Select the characters and reconstruct the row.}
Start with the nonzero abelian theta value of \ref{cor:toric-twist-5} and the common denominator of Lemma~\ref{lem:toric-common-denominator}. At the nonreference indices Lemma~\ref{lem:initial-row-signature-tests} gives the jumps, integral duality and local conditions on each occurring lower signature. Proposition~\ref{prop:signed-theta-data} clears the universal coefficients on the original lattice. Proposition~\ref{prop:retwisting-row-basis} selects independent projected vectors spanning the two signed images there. The simultaneous specialization of Proposition~\ref{prop:one-retained-character} therefore applies to both signs and every indicated signature at this index. If \(m_i\) is the order of the character, the resulting tests have
\[
          b\in\{+,-\},\qquad \varphi=\varphi_i^{\,a},
                                     \quad a\in(\mathbb Z/m_i)^\times .
\]

 Lemmas~\ref{lem:reference-field-transport}
and~\ref{lem:second-core-reference} extend the large coefficient
field, keeping the smaller two embedded fields fixed. On the
reference tuple the oscillator, both lift maps and the source
measure can be used again. Lemma~\ref{lem:copied-reference-period}
therefore identifies the actual period. For a copied eigencharacter
and a tested Hecke operator away from the excluded primes,
\[
              \lambda'(T)=j(\lambda(T)),\qquad j:K\hookrightarrow K'.
\]

For a new tuple the source classes and its level may change.
On the actual span of the new class inputs, Lemma~\ref{cor:marked-full-class-ball} also gives a full lattice. Its fixed finite list of jet denominators may depend on this tuple. Lemma~\ref{lem:new-ray-character-line} computes the character on
its new local type. Lemma~\ref{lem:new-ray-type-occurrence} puts the
required restricted eigencharacter in the new space by the
adjoint of the surjective specialized lift. Reduce these actual
Hecke characters at the coefficient place as in
Proposition~\ref{prop:marked-hecke-localization}. A full eigensystem
occurring with a given restriction contributes to the corresponding
block. More explicitly, for the tested algebra \(\mathcal H_0\subset\mathcal H\), the relevant dimension is
\[
       \dim X[\lambda_0]=
           \sum_{\lambda|_{\mathcal H_0}=\lambda_0}\dim X[\lambda].
\]
To read a single summand apply its separating projector.

\proofstep{Use both auxiliary covers on the resulting member.}
At the word prime the lattice with its trace conditions and its
Petersson norm is the one carried in Lemma~\ref{lem:word-prescribed-class-lattice}.
Assign its exceptional local sublattices as in
Lemma~\ref{lem:word-lattice-gluing}. Their valuations are now recorded
on this word. On the resulting member take the two covers for the
distinct auxiliary primes \(p_0,p_1\), retaining its level action.
At a coefficient place \(w\), the cover to use is indexed by
\[
 p(w)=\begin{cases}
           p_1,&\operatorname{char} k(w)=p_0,\\
           p_0,&\operatorname{char} k(w)\ne p_0.
       \end{cases}
\]
Its extra marking is prime to the residue characteristic, so the
level and integral evaluation pull back as in
Lemma~\ref{lem:copied-ray-level-correspondence}. The two calculations
restrict to the same original data of \(P'\). The actual values of the
two index bounds on this member are the subject of
Section~\ref{sec:word-division}.
\end{proof}

The construction is pointwise in the indicated range. The zero
criterion needs the uniform numerical inequality of
Theorem~\ref{hyp:uniform-arithmetic-data}. For the square
variant, Lemmas~\ref{lem:maximizing-translate}
and~\ref{lem:square-window-transfer} specify which translates
and scales suffice. \label{cor:window-points-2}

The signed coefficients use a fixed bound that may depend on the
stage, side and coefficient place. The following two tests
distinguish bounded denominators from an exact determination of
the integral lattice by point values.

\begin{auxlemma}[Clearing denominators on a fixed row]\label{lem:marked-bounded-denominator}
Let \(R\) be a discrete valuation ring with uniformizer \(\pi\)
and fraction field \(K\). Fix a full free lattice \(M\) of rank
\(d>0\) in a \(K\)-space \(V\), and a valuation ring \(R'\)
of a finite field extension \(K'/K\) over \(R\).
Suppose \(d\) scalar readings give an injective map
\(e:V\otimes_K K'\longrightarrow(K')^d\).
There is \(c\in\mathbb Z_{\ge0}\) such that for every integer \(b\ge0\) and \(x\in V\),
\[
           e(x)\in\pi^{-b}(R')^d
             \quad\Longrightarrow\quad \pi^{b+c}x\in M .
\]
\end{auxlemma}
\begin{proof}
Write \(A\) for the reading matrix in a basis of \(M\) and
\(v'(\pi)>0\) for the extended valuation. Select \(c\ge0\) with
\[
 c\,v'(\pi)\ \ge\
   \max_{(A^{-1})_{ij}\ne0}\{-v'((A^{-1})_{ij})\}.
\]
Since \(\pi^c A^{-1}\) is integral, the coordinates of
\(\pi^{b+c}x=(\pi^c A^{-1})(\pi^b e(x))\) lie in \(R'\cap K=R\).
The same \(c\) works for every vector in the fixed space and after
further extension of \(R'\).
\end{proof}

\begin{auxlemma}[The moving part of a signed theta denominator]\label{lem:marked-prime-to-level-denominator}
Fix the coefficient place \(v\mid\ell\), the stage, the sign,
the base level, the source classes and the harmonic polynomial.
Use the rigidified class formula of Lemmas~\ref{lem:marked-theta-fibre-comparison} and~\ref{lem:marked-class-section-descent} at the fixed good CM separating points. Write \(v'\) for the extended valuation
normalized by \(v'(\pi_v)=1\). For admissible moving levels prime
to \(\ell\) the descended coefficients satisfy
\[
             \inf_{Q,M,\delta,\,1\le h\le d}
                   v'\bigl(e_h(x_{Q,M,\delta})\bigr)>-\infty
                   \qquad(d>0).
\]
\end{auxlemma}
\begin{proof}
Work over a valuation ring containing a good model of the fixed
CM variety. The prime-to-\(\ell\) isogenies pull back sections of
the rigidified line integrally. The moving torsion labels are
finite \'etale. A theta-group lift can be normalized by adjoining
a prime-to-\(\ell\) root of a unit, which again has valuation
zero. On the chosen split half, with residue cardinality \(q\),
the denominator of the rank-\(r\) complement sum in dimension
\(s\) is a unit:
\[
                q^{r(s-r)}\in(\mathcal O'_v)^\times .
\]
Finite additive Fourier inversion on a moving half divides by
the order of an additive quotient of residue characteristic
different from \(\ell\). Split these operations from the finitely
many base half changes. Their clearing factors, the fixed class
masses, and the matrix of the fixed harmonic differential in the
integral tangent frames have one common denominator.

Finally, descent attaches a determinant class \(\alpha\) to
each open and closed component. If \(e_\alpha\) is its
idempotent and \(a_\beta\) the coefficients before that shift,
the coefficient at \(\delta\) is
\[
                      \sum_\alpha e_\alpha\,a_{\delta-\alpha}.
\]
No division by the ray-group order occurs, even when
\(\ell\mid q-1\). The valuation of a sum is bounded below by
the smallest valuation of a term. Thus the fixed factors suffice
after every admissible moving-level operation and after descent.
\end{proof}

In particular, once the framed readings at the base level agree,
one separating matrix works for every auxiliary ray level. A bound
\(b_v\) before the matrix is inverted costs at most \(b_v+c_v\)
afterwards. At \(v\mid\ell\) this bound is absorbed into the
allowed power of \(\ell\) in Proposition~\ref{prop:signed-theta-data}.
For either sign the basis may have its own matrix; the maximum of
the finitely many exponents serves both signs and the fixed test
generators. A unit minor gives the stronger conclusion below.

\begin{auxlemma}[Integral detection by marked evaluations]\label{lem:marked-evaluation-criterion}
Let \(R\) be a discrete valuation ring with fraction field \(K\),
uniformizer \(\pi\) and residue field \(k\). Let \(M\) be a finite free
\(R\)-lattice, and let \(F_i\) be finite free \(R\)-modules.
For any family of \(K\)-linear evaluations
\(e_i:M_K\to (F_i)_K\) with \(e_i(M)\subset F_i\), write
\(\bar e_i:M/\pi M\to F_i/\pi F_i\) for the reduction. Then
\[
 M=\bigcap_i e_i^{-1}(F_i)
 \quad\Longleftrightarrow\quad
       \bigcap_i\ker\bar e_i=0 .
\]
If \(d=\operatorname{rank}M>0\), this is equivalent to the existence of
\(d\) scalar coordinates of finitely many evaluations whose matrix on a basis
of \(M\) has determinant in \(R^\times\).
\end{auxlemma}
\begin{proof}
If a nonzero residue class is killed by every evaluation, lift it to
\(v\in M\setminus\pi M\). Then \(v/\pi\notin M\), while
\(e_i(v/\pi)\in F_i\) for every \(i\). Conversely, for
\(z\in M_K\setminus M\) let \(h\ge1\) be the least integer with
\(v=\pi^h z\in M\). If all the \(e_i(z)\) were integral we would have
\[
 v\notin\pi M,\qquad e_i(v)=\pi^h e_i(z)\in\pi F_i
                       \quad\text{for every }i .
\]
This contradicts the residue condition. Over \(k\), the scalar coordinate
forms of all the \(\bar e_i\) span the annihilator of their common kernel.
If this kernel is zero, choose \(d\) coordinates forming a basis of
\((M/\pi M)^\vee\). The resulting minor is a unit. A unit minor
conversely gives injectivity on reduction, also after any extension of
discrete valuation rings. For \(M=0\) both conditions hold directly.
\end{proof}

For the definite class functions the evaluations and the Hecke action
can be constructed at once. Let \(G\) be the definite group, \(K_f\)
its level and \(g_i\) representatives for
\(G(F)\backslash G(\mathbb A_{F,f})/K_f\). Set
\(\Gamma_i=G(F)\cap g_iK_fg_i^{-1}\). A form of coefficient
representation \(V\) has the transformation law
\[
 f(\gamma g k)=\gamma f(g)\qquad(\gamma\in G(F),\ k\in K_f).
\]
The map \(f\mapsto(f(g_i))_i\) has inverse
\((v_i)_i\mapsto[\gamma g_i k\mapsto\gamma v_i]\) on
\(\bigoplus_i V^{\Gamma_i}\). Two presentations of the same argument
differ by an element of \(\Gamma_i\), which proves the inverse is well
defined. This is the definite algebraic modular form description of
Gross. \cite[Sections~4 and~6]{gross-algebraic-forms}

Choose the transported \(\Gamma_i\)-stable coefficient lattices \(W_i\)
in the fibres and require integral values there. The coordinate lattice of
all classes is \(\bigoplus_i W_i^{\Gamma_i}\). For a rational subspace
\(V'\) of the class space its saturated lattice is
\[
  M'=V'\cap\bigoplus_i W_i^{\Gamma_i}.
\]
If \(v\in M'\) has every coordinate in \(\pi W_i\), then
\(v/\pi\in V'\) and all its coordinates lie in \(W_i^{\Gamma_i}\).
Thus \(v\in\pi M'\). The lemma gives integral detection for this
lattice, even at primes dividing a class stabilizer order.

For the double coset \(K_f aK_f=\coprod_b a_bK_f\) and
\(g_i a_b=\gamma_{ib}g_{j(i,b)}k_{ib}\), the counting normalization
gives the matrix formula
\[
        (T_a f)(g_i)=\sum_b\gamma_{ib}f(g_{j(i,b)}) .
\]
It follows by substituting the transformation law in right convolution.
Thus the definite class lattice is determined integrally by all its
class values. On the geometric row Lemma~\ref{lem:marked-bounded-denominator}
uses only rational separation for the coefficient bound. The
Hecke action on the geometric supremum lattice is a separate
assertion. An integral finite correspondence with unit counting
normalization preserves the condition of integral values on its
integral frames, provided its branches transport those frames
integrally. The finite list chosen to clear denominators need
only separate sections over the fraction field.

\begin{auxlemma}[Detection on the special fibre]\label{lem:marked-fibre-density}
Let \(X/R\) be flat with reduced special fibre over the residue field
\(k\) of the discrete valuation ring \(R\). Let \(\mathcal V\) be a
vector bundle and let \(M\subset H^0(X,\mathcal V)\) be a saturated
finite free sublattice. Suppose the marked sections of \(X/R\) reduce
to a dense set in each component of \(X_k\). Evaluate in integral
frames of their pullbacks of \(\mathcal V\). Then
\[
            \bigcap_i\ker(\bar e_i:M/\pi M\to F_i/\pi F_i)=0.
\]
\end{auxlemma}
\begin{proof}
Flatness and local freeness give the sequence
\[
 0\longrightarrow \mathcal V\xrightarrow{\ \pi\ }
 \mathcal V\longrightarrow\mathcal V|_{X_k}\longrightarrow0 .
\]
By left exactness of global sections and saturation, the induced map
\(M/\pi M\to H^0(X_k,\mathcal V_k)\) is injective. A residue class
killed by each evaluation maps to a section whose zero locus contains
the indicated dense set. In each trivialization its coordinates vanish
on that set and hence on the reduced fibre. The image section, and
therefore the residue class, is zero.
\end{proof}

The same proof applies over a common finite extension on whose integral
model these hypotheses hold. For the signed sections at the prescribed
levels, one must still identify the lattice and its evaluations, and
prove this density assertion or exhibit the unit minor directly.

The definite moduli points admit an integral construction. Let \(E/F\)
be a CM extension with CM type \(\Phi\) and reflex field \(E^*\).
For the signature \((d,0)\) at \(\Phi\), define the ideal
\[
 J_\Phi=\ker\!\left(\mathcal O_E\otimes_{\mathbb Z}\mathcal O_{E^*}
             \longrightarrow\prod_{\tau\in\Phi}\mathbb C\right),
 \qquad a\otimes b\longmapsto(\tau(a)b)_\tau .
\]

\begin{auxlemma}[Definite CM points and their prime-to-characteristic levels]\label{lem:marked-definite-cm-points}
Let \(A\) be a principally polarized CM abelian scheme of type \(\Phi\)
and CM rank one. Let \((L,h)\) be a finite projective
\(\mathcal O_E\)-module of rank \(d\) with a positive definite
unimodular Hermitian form. There is a principally polarized
\(\mathcal O_E\)-abelian scheme \(A_L\) of CM rank \(d\), with
\[
 A_L=L\otimes_{\mathcal O_E}A,\qquad
 \operatorname{Lie}(A_L)=L\otimes_{\mathcal O_E}\operatorname{Lie}(A),
 \qquad J_\Phi\operatorname{Lie}(A_L)=0 .
\]
On the definite moduli stack over the reflex integers, the last
condition and the principal polarization give a proper \'etale model.
At residue characteristic \(p\), fixing a compatible level structure of
order prime to \(p\) gives a finite \'etale cover of this model wherever
that level is defined.
\end{auxlemma}
\begin{proof}
Choose an idempotent with image \(L\) in a finite free module.
Apply it to the same power of \(A\). Its image is Serre's tensor
construction, and taking tangent spaces gives the displayed Lie
identity. The dual abelian scheme is obtained from \(L^\vee\) and
\(A^\vee\). The polarization is the isomorphism
\[
 h\otimes\lambda_A:L\otimes_{\mathcal O_E}A\longrightarrow
 L^\vee\otimes_{\mathcal O_E}A^\vee .
\]
Positivity gives a polarization and unimodularity makes it principal.
\cite[Theorem~17]{amir-khosravi-serre}

Amir-Khosravi's Theorem~35 proves the proper \'etale assertion for the
stack defined by \(J_\Phi\operatorname{Lie}=0\). On a flat integral
model a general-fibre CM point of this type satisfies the ideal
condition throughout: each section of \(J_\Phi\) acts on the locally
free Lie module by a matrix vanishing in the fraction field. For a
level of order \(N\) invertible on the base, \(A_L[N]\) is finite
\'etale. The full level frames form a finite \'etale sheaf; the
congruence orbit giving the specified level is a finite \'etale sheaf
as well. \cite[Theorems~35 and~52]{amir-khosravi-serre}
\end{proof}

When a principally polarized rank-one CM point exists for \(\Phi\),
the tensor construction exhausts the definite moduli objects \'etale
locally. Mere choice of a CM type is insufficient for this rank-one
assertion. If \(E/F\) is unramified at every finite prime, precisely
half of the CM types admit a principally polarized rank-one point:
\[
 \#\{\Phi:\mathcal M_\Phi^1(\mathbb C)\ne\varnothing\}
 =\tfrac12\#\{\Phi:\Phi\text{ is a CM type of }E\}.
\]
The proper \'etale assertion applies to the rank-\(d\) stack without
this extra existence assumption. For a level divisible by \(p\), even
the distinct primitive \(p\)-th roots of unity in a splitting field
have the same reduction. Thus one needs the actual moduli definition
at such a place to identify its integral level structures.
\cite[Theorems~46 and~52]{amir-khosravi-serre}

For a positive-dimensional unitary moduli space, the level and the
integral readings are different parts of the construction. Let
\(\mathcal M_K\) denote the characteristic-zero moduli space with
polarization, endomorphisms and level \(K\), using the associated
similitude group. At a rational prime \(p\), normalize a product
of auxiliary good-reduction moduli stacks in \(\mathcal M_K\), with
the auxiliary lattice data of Lan, Section~2. Thus, if \(\mathcal M_{j,\rm aux}\)
are the auxiliary models, the map into their product factors as
\[
 \mathcal M_K\longrightarrow\widetilde{\mathcal M}_K
 \longrightarrow \prod_j\mathcal M_{j,\rm aux},
 \qquad \widetilde{\mathcal M}_K\otimes\mathbb Q=\mathcal M_K .
\]
Proposition~6.1 of
that work gives a normal flat model over the localized reflex
integers. Its polarized abelian schemes with endomorphisms extend
on the open model. At a level involving \(p\), the original level
structure in this result is defined on its generic fibre.
\cite[Propositions~6.1 and~6.4]{lan-ramified-models}

For a Hecke element \(g\) use the level
\(K'=K\cap gKg^{-1}\), so that both \(K'\) and \(g^{-1}K'g\) lie
in \(K\). Keeping the original auxiliary lattices and their \(g\)-translates
gives a diagram extending the generic Hecke correspondence,
\[
   \widetilde{\mathcal M}_K\xleftarrow{\ [1]\ }
   \widetilde{\mathcal M}_{K'}\xrightarrow{\ [g]\ }
   \widetilde{\mathcal M}_K .
\]

The maps on these models are projective. For prime-to-\(p\) Hecke
change they are finite when the \(p\)-component preserves every auxiliary lattice. Writing \(G_j\) for its integral stabilizer, Lan, Proposition~13.19 assumes
\[
             g_p\in\bigcap_j G_j(\mathbb Z_p).
\] 

On the generic
fibre, for neat level, the first map is finite \'etale. If a section
\(s\) is pulled back on the right and the coefficient fibres are
identified by the universal quasi-isogeny, its trace on the left is
\[
      (T_gs)(x)=\sum_{x'\in [1]^{-1}(x)}\iota_{x'}s([g]x')
\]
at a geometric point \(x\), with counting measure. This follows by
splitting the \'etale fibre into its points, where trace on a product
is addition. These morphisms give the geometric correspondences;
integral detection requires the evaluations of the specified sections
on the particular marked points. \cite[Section~13]{lan-ramified-models}

\begin{auxlemma}[Finite detection by Serre--Tate coefficients at ordinary level]\label{lem:ordinary-finite-jet-detection}
Use the split, unramified unitary setting of Caraiani, Eischen,
Fintzen, Mantovan and Varma, with hyperspecial level at \(p\), neat
tame level and \(p\) split completely in the CM and reflex fields.
On an ordinary connected component choose a point \(x\) of its
Igusa tower over \(W=W(\overline{\mathbb F}_p)\). Take a fixed
weight and a finite free \(W\)-sublattice \(M\) of the corresponding
integral \(p\)-adic automorphic forms, saturated in that space.
For \(s\in M\) let \(a_\nu(s)\) be the coefficient of \(t^\nu\)
in its Serre--Tate expansion at \(x\) in the chosen coordinates.
There are finitely many indices \(\nu_1,\ldots,\nu_d\), with
\(d=\operatorname{rank}_WM\), such that
\[
  s\in pM\quad\Longleftrightarrow\quad
      a_{\nu_1}(s),\ldots,a_{\nu_d}(s)\in pW \qquad(s\in M).
\]
\end{auxlemma}
\begin{proof}
In fixed weight the character factors on the translates of the point
are units. Corollary~5.16 of the cited expansion principle, applied
to \(s\) and zero, gives divisibility by \(p\) in the full space if
all the coefficients are divisible by \(p\). Saturation gives the
same divisibility in \(M\). Thus the coefficient functionals on
\(M/pM\) have zero common kernel. Choose a basis of the dual from
these functionals, as in Lemma~\ref{lem:marked-evaluation-criterion}.
This proves the assertion. \cite[Theorem~5.14 and Corollary~5.16]{cef-mv-expansions}
\end{proof}

These are coefficients of a power series on the deformation space of
an ordinary point. Its value is the constant coefficient. On an
ordinary Hecke orbit actual point values can also detect the
integral lattice in the following good-reduction setting.

\begin{auxlemma}[Ordinary CM point evaluations at good level]\label{lem:ordinary-hecke-evaluations}
Let the positive-dimensional PEL datum in
Lemma~\ref{lem:marked-fibre-density} be of Hodge type. Fix a prime
\(p>2\) of good reduction, hyperspecial level at \(p\) and
sufficiently small tame level. Suppose the local reflex field is
\(\mathbb Q_p\). Fix an ordinary reduction of an integral CM point
\(z\) on the canonical model, and use its entire prime-to-\(p\)
Hecke orbit, allowing a finite extension for each evaluating point.
For a saturated finite free lattice of sections \(M\) as in that
lemma, of rank \(d>0\), there exist \(d\) scalar coordinates of
actual CM point values on this orbit with a unit determinant:
\[
     z_1,\ldots,z_d\in\mathcal H^{(p)}z,\qquad
       \det\big(e_{z_i}(m_j)\big)_{1\le i,j\le d}\in (R')^\times .
\]
Here \(m_j\) are an integral basis, the evaluations are in integral
frames, and \(R'\) is a common extension of their valuation rings.
Several coordinates may come from one point.
\end{auxlemma}
\begin{proof}
Over the strict henselization the prime-to-\(p\) torsion of the
abelian scheme is finite \'etale at every fixed order. Thus the
subgroup, its PEL tensors and the transported tame marking in a
prime-to-\(p\) Hecke correspondence lift along with \(z\). For
a general prime-to-\(p\) quasi-isogeny first clear its denominator
by multiplication prime to \(p\). The quotient of the lift by the
lifted finite subgroup gives a point on the generic Hecke orbit
reducing to the specified image. These quotients are again CM,
since isogeny preserves the rational endomorphism algebra. In
particular, with reduction taken over a geometric residue field,
\[
      \mathcal H^{(p)}\bar z\ \subseteq               \overline{\mathcal H^{(p)}z}^{\,\mathrm{red}},
\]
where the right side denotes the set of reductions of the orbit points.
Van Hoften's Theorem I gives Zariski density of the orbit on the
left in the special fibre. \cite[Theorem I]{van-hoften-ordinary}

The special fibre is smooth and reduced. The injection
\(M/\pi M\hookrightarrow H^0(X_k,\mathcal V_k)\) in
Lemma~\ref{lem:marked-fibre-density} stays injective after
extending the residue field. A section killing every orbit value
would vanish on a dense subset and hence everywhere. Pick \(d\)
independent geometric residue coordinates. Their finitely many
characteristic-zero points and frames can be taken over a common
finite extension; reduction of the determinant there is nonzero.
\end{proof}

For fixed marked tests there is a different way to identify exactly
where the same conclusion survives a change of coefficients. Suppose
the two signed evaluation matrices for one fixed realization have
already been constructed over a number field and have full column
ranks over that field. Spread their two lattices and entries over a
localization \(A\) of its integers. Denote their full-rank minor ideals
by \(J_+\) and \(J_-\). Then the bad primes for simultaneous
detection on these very tests in \(\operatorname{Spec}A\) are precisely
\[
          V(J_+J_-)=V(J_+)\cup V(J_-).
\]
Indeed a minor is a unit in the local ring if and only if its residue
is nonzero. Both ideals are nonzero, so this set is finite for the
fixed matrices.

For \(w\mid v\), extending the coefficient field gives
\[
     v_w(J_+J_-)=e(w/v)\,v_v(J_+J_-).
\]
A positive valuation persists under coefficient extension. It merely
increases the fixed clearing exponent in
Lemma~\ref{lem:marked-bounded-denominator}. The height and
divisibility comparisons still use their original lattices in
Proposition~\ref{prop:all-window-applications}.

\begin{auxlemma}[A CM point with the two prescribed types]\label{lem:marked-two-type-seed}
Let \(E/F\) be the CM extension of the row and fix its CM type
\(\Psi\). Suppose the Hermitian space \((V,H)\) has signature
\((p_+,p_-)\) at every embedding in \(\Psi\). Use the convention
that \(H\) is linear in the first argument. Write the alternating
PEL pairing on the marked lattice \(\Lambda\) as
\[
    e(v,w)=\operatorname{Tr}_{E/\mathbb Q}(\delta H(v,w)),\qquad
    \bar\delta=-\delta,\qquad \operatorname{Im}\tau(\delta)>0
             \quad(\tau\in\Psi).
\]
Assume \(\Lambda\) is integral and self-dual for \(e\). Its complex
unitary component contains a principally polarized CM point with
that lattice and its given finite marking. There is an \(E\)-linear
isogeny to it from a product of \(p_+\) factors of type \(\Psi\)
and \(p_-\) factors of type \(\bar\Psi\), each of rank one over \(E\).
\end{auxlemma}
\begin{proof}
Choose an orthogonal signature basis at each archimedean place.
The conditions on the signs of the consecutive principal Gram
determinants are open. Weak approximation for the entries of a
matrix over \(E\) therefore gives an \(E\)-basis satisfying them
simultaneously. Apply Gram--Schmidt over \(E\); its pivots
\(a_j\in F^\times\) are the successive determinant ratios. Thus
\[
 V=\mathop{\perp}_{j}L_j,\qquad
 \tau(a_j)>0\ (j\le p_+),\quad
 \tau(a_j)<0\ (j>p_+),\qquad \tau\in\Psi .
\]

On the \(\tau\)-component of the line \(L_j\) define
\(J_j=i\operatorname{sgn}(\tau(a_j))\), extending by conjugation on
the opposite embedding. For \(J=\bigoplus_j J_j\), the pairing is
\(J\)-invariant and in these coordinates
\[
  e(v,Jv)=2\sum_{\tau\in\Psi,j}
           \operatorname{Im}\tau(\delta)\,|\tau(a_j)|\,|v_{\tau j}|^2>0
                             \quad(v\ne0).
\]
The Riemann bilinear relations give a polarized complex abelian
variety \(A_J=(V_{\mathbb R},J)/\Lambda\); self-duality makes its
polarization principal. The rank-one factors have the claimed CM
types and their Hodge tori lie in the product of the corresponding
restriction-of-scalars tori. The chosen marking stays on the
original lattice. Over a finite extension one may also choose a
compatible point of the finite characteristic-zero theta-level cover
when such a cover is part of the level.

Set \(\Lambda_j=L_j\cap\Lambda\). The sum of these lattices has
full rank and is contained in \(\Lambda\). Addition of the
complex subspaces therefore gives an isogeny with explicit kernel
\[
  0\longrightarrow\Lambda/\!\bigoplus_j\Lambda_j
    \longrightarrow\prod_j (L_{j,\mathbb R},J_j)/\Lambda_j
    \longrightarrow A_J\longrightarrow0 .
\]
The polarizations restricted to the factors may have nontrivial
degree. Each fractional ideal lattice of a given type is
commensurable with a fixed rank-one lattice of that type. Clearing
these indices gives homomorphisms from two fixed CM factors, with
orthogonal images. The CM point and a fixed finite collection of
these homomorphisms and torsion markings descend to a finite
extension by the theory of complex multiplication.
\end{proof}

\begin{auxlemma}[Principal polarizations on the two CM source types]\label{lem:marked-principal-cm-sources}
Use the trace form \(e=\operatorname{Tr}_{E/\mathbb Q}(\delta H)\)
and the type \(\Psi\) of Lemma~\ref{lem:marked-two-type-seed}.
Let \(\mathfrak D_E\) be the absolute different. Each of the
following data gives a principally polarized rank-one
\(\mathcal O_E\)-abelian variety of each of the types
\(\Psi,\bar\Psi\):
\[
 \begin{array}{ll}
 \text{(i)} &\Lambda=\Lambda^{\vee_e},\quad
             \operatorname{rk}_{\mathcal O_E}\Lambda=2k+1,\quad
             p_-(H_\tau)=p_-\quad(\tau\in\Psi),\\
 \text{(ii)}& I\text{ a fractional ideal of }E,\quad
              (\delta)I\bar I=\mathfrak D_E^{-1}.
 \end{array}
\]
For the field \(E\) in \eqref{eq:host-character-compositum},
(ii) holds for the type extending either embedding of
\(k_0=\mathbb Q(\sqrt{-d})\), with opposite signs of \(\delta\).
\end{auxlemma}
\begin{proof}
For (i) put \(J=(\delta\mathfrak D_E)^{-1}\), so \(\bar J=J\).
Since \(\Lambda\) is an \(\mathcal O_E\)-module, testing its
trace pairing on every scalar multiple of each argument gives
\[
 \Lambda=\{x\in V:H(x,\Lambda)\subseteq J\}.
\]
Thus \(H\) is a perfect \(J\)-valued Hermitian pairing.

Taking determinants in a local \(\mathcal O_E\)-basis gives a
perfect \(J^{2k+1}\)-valued pairing on \(\det\Lambda\).
Tensor that line with \(J^{-k}\) and multiply its Hermitian
pairing by \((-1)^{p_-}\), where \(p_-\) is the common negative
index. The resulting line \(P\) and form \(h_P\) satisfy
\[
 P=(\det\Lambda)\otimes J^{-k},\qquad
 h_P(P,P)=J^{2k+1}J^{-k}\bar J^{-k}=J.
\]
Its real scalar at every member of \(\Psi\) is positive, since
the sign of the Hermitian determinant is \((-1)^{p_-}\).
The alternating form \(\operatorname{Tr}_{E/\mathbb Q}(\delta h_P)\)
is perfect on \(P\). With the complex structure of type \(\Psi\)
it satisfies the Riemann positivity condition. Negate both this
complex structure and this alternating form to obtain type
\(\bar\Psi\) with a principal polarization.

For (ii) take \(P=I\) with \(H(x,y)=x\bar y\). Its trace dual is
\[
   \{x:\operatorname{Tr}_{E/\mathbb Q}(\delta x\bar I)
                       \subseteq\mathbb Z\}
           =(\delta)^{-1}\mathfrak D_E^{-1}\bar I^{-1}=I .
\]
Choose the sign of \(\delta\) as above, and reverse it for the other
type. In the host construction put \(k_0=\mathbb Q(\sqrt{-d})\)
and take \(\delta=\pm1/\sqrt{-d}\) with positive imaginary part
at the chosen embedding of \(k_0\). We now construct \(I\).

Write \(\mathfrak D=\mathfrak D_{E/k_0}\). At a rational prime
\(p\nmid 2q f_\psi\), the extension over \(k_0\) is obtained
inside \(k_0(\mu_{Q_Sq})\). For any completion of \(k_0\) above
\(p\), the order \(Q_Sq\) is prime to \(p\), so
Corollary~10.16 of \cite{sutherland-local-fields} gives an
unramified extension there. This includes \(p=d\), where the
base field itself ramifies. Since the residue fields are finite,
hence separable, Theorem~12.19 of the same reference gives
\[
  \mathfrak P\mid\mathfrak D\ \Longrightarrow\ 
             p\mid 2qf_\psi\quad(\mathfrak P\mid p) .
\]

Every rational prime on this list splits in \(k_0\).
At two, \(-d\equiv1\pmod8\) is a dyadic unit square.
All odd divisors of the fixed \(f_\psi\) are at most \(w\) on
sufficiently late intervals. For these primes and for \(q\),
Proposition~\ref{prop:quadratic-prime-pair} gives
\[
        d\equiv-1\pmod p,\qquad -d\in\mathbb Q_p^{\times2} .
\]

Indeed the simple root \(1\) of \(X^2+d\) modulo \(p\) lifts.
Thus complex conjugation exchanges the two primes of \(k_0\)
under every possible factor of \(\mathfrak D\). In particular
\(\bar{\mathfrak P}\ne\mathfrak P\) above each such prime.
Conjugation preserves the relative different since it carries the
trace dual of \(\mathcal O_E\) over \(\mathcal O_{k_0}\) to itself.
Choose one prime from each conjugate pair in its support and put
\begin{equation}\label{eq:host-different-factors}
 a_{\mathfrak P}=v_{\mathfrak P}(\mathfrak D)
                  =v_{\bar{\mathfrak P}}(\mathfrak D),\qquad
 \mathfrak b=\prod_{\{\mathfrak P,\bar{\mathfrak P}\}}
                   \mathfrak P^{a_{\mathfrak P}},\qquad
 \mathfrak D=\mathfrak b\bar{\mathfrak b},\quad I=\mathfrak b^{-1}.
\end{equation}

The rings of integers in the tower \(E/k_0/\mathbb Q\) are
Dedekind domains and both extensions are separable.
Transitivity of the different \cite[Proposition~12.28]{sutherland-local-fields}
now gives
\[
   (\delta)=\mathfrak D_{k_0/\mathbb Q}^{-1}\mathcal O_E,\qquad
   I\bar I=\mathfrak D_{E/k_0}^{-1},\qquad
   (\delta)I\bar I=\mathfrak D_E^{-1}.
\]
The multiplication action is by \(\mathcal O_E\), with Rosati
involution complex conjugation. The Riemann form thus gives the
required polarized CM varieties over \(\mathbb C\); CM descent
defines them and their polarizations over a finite extension.
\cite{cm-reciprocity}
\end{proof}

The same half different converts an integral Hermitian model into
a trace-self-dual period model. If \(L=L^{\sharp_H}\), with
\(L^{\sharp_H}=\{v:H(v,L)\subseteq\mathcal O_E\}\), then for
\(\Lambda=I L\) one has
\[
 (IL)^{\vee_e}
     =(\delta\mathfrak D_E)^{-1}\bar I^{-1} L^{\sharp_H}
     =I L .
\]
For a trivial target the only possible rational support of the
relative different is \(\{2,q\}\). For a general target the
factorization \eqref{eq:host-different-factors} includes every
conductor prime as well.

The two principal sources may now be used in the product isogeny
of Lemma~\ref{lem:marked-two-type-seed}. Each of its rank-one
lattices is commensurable with the chosen ideal of the same type.

This supplies the initial CM lift for
Lemma~\ref{lem:marked-generic-cm-separation} on a constant-signature
component with the indicated trace form. The indices in the kernel
above record exactly where a bound on the polarizations of the
individual homomorphisms would be needed when the integral markings
and the weight estimates vary.

\begin{auxlemma}[Separation by CM points on a fixed complex component]\label{lem:marked-generic-cm-separation}
Let \(\Gamma\backslash D\) be a Hermitian domain quotient in a
unitary PEL variety, at a fixed finite level. For a CM lift
\(z\in D\) and a representative of the finite level component, the
resulting CM evaluations have zero common kernel on each space of
algebraic sections on this component:
\[
   \bigcap_{\gamma\in U(F)^+}\ker(\operatorname{ev}_{[\gamma z]})=0 .
\]
Here \(F\) is the totally real ground field and \(U(F)^+\) denotes
the rational unitary elements preserving the real component.
\end{auxlemma}
\begin{proof}
For the Lie space of skew-adjoint endomorphisms use the Cayley map
\[
                 X\longmapsto (1-X)(1+X)^{-1}.
\]
Weak approximation in this vector space gives rational unitary
elements dense in a neighborhood of the identity at infinity. The
closure of the rational subgroup therefore contains the real
identity component, which acts transitively on \(D\). Rational
conjugation carries the CM Hodge torus to a rational torus. Thus
the images of \(\gamma z\) are CM points, dense on this component.
A section vanishing there vanishes identically. In a finite
dimensional space this gives finitely many scalar fibre readings
whose matrix has full column rank.
\end{proof}

When the signed evaluations include these points and frames on
every contributing component, they supply the matrix of
Lemma~\ref{lem:marked-bounded-denominator}. For the finite list,
CM abelian varieties acquire good reduction over a common finite
extension at any fixed coefficient prime. Their additional
prime-to-characteristic torsion markings then extend over finite
unramified extensions. \cite[Theorem~6(a)]{serre-tate}

This separation argument allows CM points to move. A bounded
geometric height test can require two fixed complementary CM types,
integral \(\mathcal O_E\)-linear homomorphisms from those factors
whose sum on the product is an isogeny onto the tested polarized
variety, and the transported eigenframes.
If the polarization on the \(i\)-th factor pulls back to its own
polarization multiplied by \(h_i\), the marked bound requires
\[
         \max_{i,\tau:E\hookrightarrow\mathbb C}|\tau(h_i)|
                 \le L_a^2 .
\]
Transport by a rational unitary element preserves the isogeny
class. It does not bound the denominators of the individual
homomorphisms or these polarization factors. The quantitative height test uses the CM approximations of Lemma~\ref{lem:deep-cover-detecting-frames} and the estimates of Section~\ref{subsec:deep-cm-interior}. Generic separation alone supplies a fixed coordinate bound.

\begin{auxlemma}[A coefficient field fixed before the ray characters]\label{lem:marked-original-geometric-field}
Let \(E/\mathbb Q\) be the abelian CM host, \([E:\mathbb Q]=2f\),
with class number \(h_E\) and embedded imaginary quadratic
subfield \(k\). Let \(M_0,M_1\) be the two fixed CM torsion
moduli for the auxiliary levels. Let \(\mathcal N\) be the finite
list of root-of-unity orders for the original level, including
\(h_E\). Write \(r\) for the rank and \(m\) for the number
of stage signatures. Put
\[
   s(K)=\max\{\log[K:\mathbb Q],\log\operatorname{rd}(K)\},
   \qquad \operatorname{rd}(K)=|D_K|^{1/[K:\mathbb Q]} .
\]
Fix an algebraic closure of \(\mathbb Q\) and the original embedding
of \(E\). The two sets of Weber values, these roots of unity, the
generators principalizing the ideals of \(E\), and the generators for
the \(m\) stage type norms all lie in one field \(K_{\rm geom}\) with
\[
 s(K_{\rm geom})\le s(E)+\sum_{a\in\mathcal N}\log a
       +s(k_{M_0})+s(k_{M_1})+(m+1)\log h_E .
\]
Here \(k_M\) is the ray field of \(k\) of modulus \(M\mathcal O_k\).
This construction uses neither the subsequently chosen ray characters
nor their fields of values.
\end{auxlemma}
\begin{proof}
Adjoin the two ray fields and the indicated roots. The main theorem
of complex multiplication places the Weber values in the ray fields
as in Section~\ref{subsec:joint-family-estimate}. Compositum does
not increase either logarithmic cost beyond the sum of the costs:
for the discriminant use the tower formula and divisibility of the
relative different after base change. The cyclotomic cost is at most
\(\log a\). The ray exact sequence and conductor--discriminant formula give
\[
  [k_M:\mathbb Q]\le2h_k M^2,\qquad
   \operatorname{rd}(k_M)\le |D_k|^{1/2} M .
\]
Indeed every character of the ray class group has conductor dividing
\(M\mathcal O_k\), whose absolute norm is \(M^2\).

Let an ideal \(\mathfrak a\) represent a class of order \(b\)
modulo the subgroup already principalized in the growing field \(K\).
Then \(\mathfrak a^b\mathcal O_K=(u)\). Adjoin \(t\) with
\(t^b=u\); equality of ideal valuations shows that
\(\mathfrak a\mathcal O_{K(t)}=(t)\).

At a finite place \(v\) divide \(t\) by a local generator of
\(\mathfrak a\mathcal O_{K_v}\), obtaining a unit \(t_v\).
Because \(t_v^b\) is a base unit and \(\mu_b\subset K_v\), the
extension is Galois. The map
\[
 \operatorname{Gal}(K_v(t_v)/K_v)\hookrightarrow\mu_b,\qquad
          \sigma\longmapsto \sigma(t_v)/t_v
\]
has image \(\mu_{e_0}\), where \(e_0=[K_v(t_v):K_v]\).
Thus \(e_0\mid b\), \(u_0=t_v^{e_0}\in\mathcal O_{K_v}^{\times}\),
and \(X^{e_0}-u_0\) is the minimal polynomial by degree.
The derivative bounds the different of the maximal order:
\[
   \mathfrak D_{K(t)/K}\mid b\mathcal O_{K(t)},\qquad
   \log\operatorname{rd}(K(t))-\log\operatorname{rd}(K)\le\log b .
\]
The degree cost has the same bound. At each step the subgroup of
classes killed grows by at least the chosen order \(b\). The product
of these orders in one pass is at most \(h_E\).

For one signature form the ideal
\[
 N_j(\mathfrak a)=
   \prod_{\tau\in\Psi}\tau(\mathfrak a)^{p_{j,\tau}}
             \bar\tau(\mathfrak a)^{r-p_{j,\tau}} .
\]
For a principal ideal \((u)\) its natural generator has squared
absolute value \(|N_{E/\mathbb Q}u|^r\) at every embedding.
Indeed, if \(v_j(u)\) denotes that generator, the product uses
exponent \(r\) at each embedding:
\[
       v_j(u)\overline{v_j(u)}
               =\prod_{\tau:E\hookrightarrow\mathbb C}\tau(u)^r
               =|N_{E/\mathbb Q}u|^r .
\]
For fractional ideals \(\mathfrak a\), require a generator
of \(N_j(\mathfrak a)\) of squared absolute value
\((N\mathfrak a)^r\) at all embeddings. Those \(\mathfrak a\)
form a subgroup containing the principal ideals. If a
power \(N_j(\mathfrak a)^b\) has such a generator, a \(b\)-th
root restores the required size as well as principalizing the ideal.
The same local derivative bound applies. Repeating once for each
signature adds at most \(m\log h_E\) to each logarithmic cost.
\end{proof}

All larger fields needed for finite tests may now be taken in the
fixed algebraic closure, extending this embedding. The arithmetic
cost of the geometric field is independent of those later choices.
What the family construction needs in addition is for the two actual
section systems at every selected translate to be defined with these
same geometric coefficients and signed weight lists.


\begin{auxlemma}[The embedded field of the marked PEL datum]\label{lem:original-marked-pel-field}
Let \(K_s\subset K\) be the original small coefficient field inside the
large coefficient field, and let \(\iota:K\hookrightarrow\mathbb C\)
be the prescribed embedding. Record the integral coefficients by the
map \(\phi:\mathcal O_{K_s}\to\mathbb C\), \(a\mapsto\iota(a)\).
The field of these coefficients and its structural embedding in \(K\)
are
\[
       F=\mathbb Q(\phi(\mathcal O_{K_s}))=\iota(K_s),
       \qquad \eta:F\longrightarrow K,\quad
               \eta(\phi(a)/\phi(b))=a/b\quad(b\ne0).
\]
Suppose a second presentation uses a subfield \(F'\subset\mathbb C\)
and an isomorphism \(h:\mathcal O_{F'}\simeq\mathcal O_{K_s}\)
with \(x=\phi(h(x))\) for every \(x\in\mathcal O_{F'}\).
Then \(F'=F\) inside \(\mathbb C\) and the structural maps into any
fixed extension \(K'/K\) agree on \(F\).
\end{auxlemma}
\begin{proof}
Every \(z\in K_s\) is a fraction of two elements of
\(\mathcal O_{K_s}\). Its prescribed complex image is the quotient of
their recorded images. Conversely all such quotients belong to
\(\iota(K_s)\). The integral comparison gives the equality of subsets
\[
          \mathcal O_{F'}=\phi(\mathcal O_{K_s})\subset\mathbb C .
\]
Taking fractions proves the second assertion. An extension of the
integral comparison must send a fraction to the quotient of the two
integral images, proving uniqueness. In particular a fixed row
coefficient already belonging to \(F\) has the same two signed
readings whenever the second embedding restricts to the first and
its conjugate on the two prescribed sides.
\end{proof}

\begin{auxlemma}[The generic schemes with the two fixed markings]\label{lem:marked-pel-base-levels}
Fix a self-dual polarized \(\mathcal O_E\)-lattice
\((\Lambda,e)\), its Hodge multiplicities \(d_\tau\) and the
associated unitary PEL datum. Let the embedded field \(F\) contain
its reflex field. Take two distinct odd marking primes \(g_0,g_1\)
with a chosen half \(\mathfrak b_i\) at each, as in the prescribed
split levels. Thus
\[
     M=g_0g_1,\qquad \mathfrak b=\mathfrak b_0\mathfrak b_1,\qquad
                  \mathfrak b\bar{\mathfrak b}=M\mathcal O_E .
\]
Use the compact open fixing the given paired torsion markings, and
use its neat marking. The PEL moduli problem of this level and the
problems with any additional prime-to-\(M\) torsion labels at
smaller neat compact opens have smooth quasi-projective representing
schemes over \(F\). The maps forgetting the extra labels are finite
\'etale. Every algebraic highest weight defined over the coefficient field gives
a bundle on this same generic scheme.
\end{auxlemma}
\begin{proof}
The objects are polarized abelian schemes \(A\), with endomorphism
action satisfying the Rosati condition and the determinant equation
\[
       \det(T-a\mid\operatorname{Lie}A)
                 =\prod_\tau (T-\tau(a))^{d_\tau}
                   \qquad(a\in\mathcal O_E),
\]
and the chosen orbit of level isomorphisms on their Tate modules.
The paired principal markings at the two distinct primes are neat. Indeed for an element of an arithmetic stabilizer the eigenvalues reduce to one at each marked prime. This remains true for their products and inverses. If \(\zeta\) is such a root of unity, its order satisfies
\[
        \operatorname{ord}(\zeta)\mid g_0^{a_0},\qquad
        \operatorname{ord}(\zeta)\mid g_1^{a_1}
            \quad\Longrightarrow\quad \zeta=1 .
\] Paired half-markings give the same test because the other half is their dual. Lan's representability theorem gives the generic scheme over the
reflex field. Base change gives the stated scheme over \(F\).
\cite[Section~2]{lan-ramified-models}

For \(U'\subset U\), choose \(n\) with the principal congruence
subgroup of level \(n\) inside \(U'\). In characteristic zero the
paired level-\(n\) isomorphisms form a finite \'etale local system.
Over a geometric point with one chosen \(U\)-marking its lifts form
the finite set \(U/U'\). Descent along a trivializing \'etale cover
gives the finite \'etale forgetful map.
Over a field splitting the embedding labels, the Hodge bundle has
the eigensummands
\[
          \omega_\tau=\bigcap_{a\in\mathcal O_E}
                       \ker(a-\tau(a):\omega\longrightarrow\omega).
\]
The two weight representations are applied to the frame torsors of
these same summands. This construction commutes with field
extension. Over \(F\) use the weight representations defined there;
an individual non-invariant labelled weight requires its field of
definition. Prescribing an eigensection and its period in a bundle
is a further condition on that section.
\end{proof}

At a positive-dimensional vertex the signed class formula is a
geometric section, whereas the period is evaluated on a vector of
a real representation. We describe the map between these two spaces.
The choice of the archimedean vector matters even when the finite
marking is unchanged.

\begin{auxlemma}[A common period domain for the marked level and its refinements]
\label{lem:marked-period-domain}
Use the polarized Hermitian lattice, linear-first convention and
CM type of Lemma~\ref{lem:marked-two-type-seed}. Let
\(G=\operatorname{Res}_{F/\mathbb Q}U(V,H)\). Fix a compact
open \(U\) preserving the finite level structure and its lattice.
Consider the complex PEL locus with this rational Hermitian space,
polarization and finite genus. With exact isometries as markings,
its complex uniformization is
\[
 G(\mathbb Q)\backslash
    \bigl(\mathcal D\times G(\mathbb A_f)/U\bigr),\qquad
 \mathcal D=\prod_{\tau\in\Psi}
      \{W\subset V_\tau:\dim W=p_-,\ H_\tau|_W<0\}.
\]
For \(U'\subset U\) forgetting part of the marking sends
\([W,aU']\) to \([W,aU]\). A prescribed unitary basis at infinity
can be used for every such refinement.
\end{auxlemma}
\begin{proof}
Write \(\widehat\Lambda=\Lambda\otimes\widehat{\mathbb Z}\).
The finite component \(a\) gives
\(\Lambda_a=V\cap a\widehat\Lambda\), whose completion is
\(a\widehat\Lambda\). This follows prime by prime by clearing a
common denominator and using the Chinese remainder theorem. The
lattice is self-dual since \(a\) is an isometry. Put
\(J_W=i\) on \(W_\tau^\perp\) and \(J_W=-i\) on \(W_\tau\).
Extend by conjugation on the opposite embeddings. If
\(v_\tau=v_\tau^++v_\tau^-\) in this decomposition, then
\[
 e(v,J_Wv)=2\sum_{\tau\in\Psi}\operatorname{Im}\tau(\delta)
    \{H_\tau(v_\tau^+,v_\tau^+)-H_\tau(v_\tau^-,v_\tau^-)\}>0
                 \qquad(v\ne0).
\]

The marked torus \((V_{\mathbb R},J_W)/\Lambda_a\) is therefore
principally polarized. Conversely the complex structure on a marked
point commutes with \(E\). The Riemann relations make its two
eigenspaces orthogonal, and the displayed positivity identifies
the negative one. The finite marking gives \(aU\) in this genus.
Changing the rational isometry multiplies both entries on the left
by the same element of \(G(\mathbb Q)\).

Fix an orthonormal positive and negative basis at a base point.
Every negative plane is the graph of a matrix, since its
projection to the negative factor is injective and both have
dimension \(p_-\). The domain in this basis is
\[
                  W=\operatorname{graph}(Z),\qquad I-Z^*Z>0 .
\]
The graph coordinates are holomorphic Grassmannian coordinates.
Neither the splitting used to write them nor the polarizing form
changes under \(U'\subset U\). The construction also applies if
one signature multiplicity is zero, when this factor is a point.
This is the fixed-isometry part of the PEL uniformization.
\cite[Section~8]{milne-pel-period-domain}
\end{proof}

\begin{auxlemma}[Reading a cusp section on the lowest compact type]
\label{lem:marked-lowest-reading}
On the domain of Lemma~\ref{lem:marked-period-domain}, fix a
stage at a holomorphic coefficient embedding. Use the regular
holomorphic discrete series \(D_\infty\) for the prescribed
weight, its irreducible lowest compact type \(W\), and a unit
extremal vector \(v_0\in W\). Write \(\mathcal H\) for the
cuspidal \(L^2\)-space with the specified finite level and
character. Evaluation on \(v_0\) is injective:
\[
 \operatorname{Hom}_{G(\mathbb R)}(D_\infty,\mathcal H)
       \longrightarrow\mathcal H,\qquad A\longmapsto A(v_0).
\]
Let \(s\) be a smooth holomorphic cusp section of the associated
weight bundle in the dual compact-type convention. Assume its
framed evaluations are bounded on the finite-volume quotient and
have zero constant Fourier coefficient along every rational
isotropic-line translation. The scalar evaluation of \(s\) on
\(v_0\) lies in the displayed image. This construction intertwines
the finite correspondences preserving its type.
\end{auxlemma}
\begin{proof}
\proofstep{Pass from a section to a function.}
In a unitary frame over the component represented by \(a\), a
section is a function with values in \(W^\vee\). If \(K_\infty\)
is the compact stabilizer, its evaluations \(F_w\), linear in
\(w\in W\), obey the frame transformation rule
\[
     F_w(gk)=F_{kw}(g)\quad(k\in K_\infty),\qquad
     F_w(\gamma g)=F_w(g)\quad(\gamma\in G(\mathbb Q)).
\]
The second identity uses simultaneous transport of the point and
its frame. The first is the defining equivariance of the associated
bundle. Holomorphy says that the antiholomorphic tangent
algebra \(\mathfrak p^-\) annihilates this compact type in the lifted
functions. Boundedness and finite volume give \(F_w\in L^2\).

For the unipotent radical \(N\) of a proper rational parabolic
choose a rational isotropic vector in the first space of its flag.
The radical fixes that first isotropic space pointwise. Its transvections along the chosen vector therefore form a central additive subgroup \(Z\subset N\).
The Fourier coefficient of frequency zero is the mean in these
coordinates. At fixed finite level, integration over
\(Z(\mathbb Q)\backslash Z(\mathbb A)\) reduces to integration over
real period parallelepipeds: the fixed open subgroup at the
finite places and additive approximation give the period
lattice at infinity. Thus the constant term satisfies
\[
 \int_{[N]}F_w(ng)\,dn=
    \int_{[N/Z]}\left(\int_{[Z]}F_w(zng)\,dz\right)d\bar n=0 .
\]
Here \([U]=U(\mathbb Q)\backslash U(\mathbb A)\) in the integration
formula, with compatible quotient measures. If there is no proper
rational parabolic the cusp condition is empty.

\proofstep{Identify the representation generated by the reading.}
If the map \(w\mapsto F_w\) is zero, take the zero intertwiner.
Otherwise its pullback inner product on \(W\) is a positive scalar
multiple of the original one, by compact irreducibility. On
compact-finite vectors use the triangular decomposition
\[
 \mathfrak g_{\mathbb C}=
      \mathfrak p^+\oplus\mathfrak k_{\mathbb C}\oplus\mathfrak p^- .
\]

The invariant inner product of any two vectors generated by \(W\)
is determined by its restriction to \(W\). Indeed, move derivatives
from the left entry to the right by their unitary adjoints. Order
the resulting word with \(\mathfrak p^-\) on the right. A word
ending in a lowering operator kills \(W\); an unpaired raising
factor has a different compact central weight and pairs to zero
with \(W\). The remaining words act by
\(U(\mathfrak k_{\mathbb C})\) on \(W\).
Equivalently, for words \(P,Q\) in raising operators this gives
\[
 \langle P F_w,Q F_{w'}\rangle
     =c\,\langle P w,Q w'\rangle_{D_\infty},
       \qquad c>0 .
\]

For the use of the unitary adjoints, fix concentric coordinate
polydiscs at the origin of the period domain and a holomorphic
frame on the larger disc. Translate the discs and frames by
\(g_\infty\). The change from unitary frames to the translated
holomorphic frame has the same bound on each smaller closed disc,
by invariance. Cauchy's inequalities on finitely many coordinates
then give, for each order \(a\),
\[
    \sup_g |D^a F_w(g)|\le C_{a,W}\!
           \max_{w'\text{ in a fixed basis of }W}\sup_g|F_{w'}(g)|
               \qquad(w\text{ in that basis}).
\]
The compact directions act in the finite-dimensional representation
\(W\). Thus the smooth derivatives used in the adjoint calculation
are square integrable by finite volume.

The Gram equality makes \(P w\mapsto P F_w\) well-defined,
including every linear relation among such vectors. It extends
to a scalar multiple of an isometry on the Hilbert completion.
Compact-finite vectors of the holomorphic discrete series are analytic vectors. The Gram equality gives the same norm-convergent exponential series for their images near the identity. Writing \(I\) for that extension and \(R\) for right translation, analytic continuation from a neighborhood of the identity in the connected group gives
\[
       I D_\infty(g)=R(g)I\quad(g\in G(\mathbb R)),\qquad I(v_0)=F_{v_0}.
\]
Uniqueness follows because the orbit of \(v_0\) spans a dense
subspace of the irreducible \(D_\infty\). At a finite Hecke
correspondence the evaluation of each branch commutes with
the real action. Apply uniqueness after summing the branches.
At a conjugate embedding complex conjugation gives the
corresponding antiholomorphic construction.
\end{proof}

\begin{auxlemma}[Theta readings on a common marked fibre]\label{lem:marked-theta-fibre-comparison}
Fix a principally polarized complex abelian variety with a
rigidified theta line, an \(\mathcal O_E\)-action and eigenframes in its
tangent space. Require the alternating form \(e\) of the polarization
to satisfy the PEL adjoint identity
\[
          e(ax,y)=e(x,\bar a y)\qquad (a\in\mathcal O_E).
\]

Give two presentations of the same fibre the same
finite theta structures, transported to a common refinement of
their levels. Suppose their principal theta generators correspond
by \(\vartheta_2=t\vartheta_1\), \(t\in\mathbb C^\times\).
Fix a class input using \(n\) theta factors and a homogeneous joint
harmonic polynomial \(P\) of degree \(d\). Use corresponding
torsion labels, translation lifts, Fourier weights and class
measures in the two presentations. Then their algebraic harmonic
readings \(\mathcal T_i(P)\) satisfy
\[
              \mathcal T_2(P)=t^n\mathcal T_1(P).
\]
The same relation holds under a change of the polarized splitting
of the formal coordinates whose logarithmic gauge is a linear
combination of the polarization pairings between opposite
eigenspaces. If the common line, markings and both generators are
defined over \(K\subset\mathbb C\), then \(t\in K^\times\).
\end{auxlemma}
\begin{proof}
\proofstep{Reading the principal line at the marked division points.}
Use the lattice \(\Lambda\) of the principally polarized fibre.
For the integral marked half ideal
\(\mathfrak a\bar{\mathfrak a}=(M)\),
\(M\ge2\), put \(L'=\mathfrak a\Lambda\). For the alternating
trace form \(e\), the dual of \(L'\) with respect to \(e/M\) is
\[
                     M\bar{\mathfrak a}^{-1}\Lambda=L'.
\]
Thus on the torus with lattice \(L'\) the pulled-back theta line
has the form \(B^{\otimes M}\), with \(B\) principal: any
algebraically trivial line discrepancy has an \(M\)-th root
in the complex dual torus. Write \(g\) for its dimension.
Riemann--Roch gives \(h^0(B^{\otimes M})=M^g\).

In a marked symplectic decomposition of \(M\)-torsion, one
isotropic half acts on sections by commuting operators.
If \(b\) acts on \(v\ne0\) by the character \(\lambda(b)\), the
commutator pairing \(e_M\) gives, for \(a\) in the other half,
\[
 b\cdot(a\cdot v)=e_M(b,a)\lambda(b)\,(a\cdot v).
\]
These \(M^g\) characters are distinct by perfectness, so their
eigenvectors form a basis. The same argument in any nonzero invariant
subspace proves irreducibility, the concrete form of Mumford's
uniqueness result for the given central character.
\cite[Section~1, Proposition~3]{mumford-theta-isogeny}

The origin is not a base point of \(B^{\otimes M}\). Indeed,
translate a principal section by \(b_1,\ldots,b_M\) with
\(\sum b_i=0\). The product belongs to \(B^{\otimes M}\) by the
theorem of the square. An open dense set of choices of
\(b_1,\ldots,b_{M-1}\) avoids the theta divisor in every factor.
Hence the kernel of all evaluations at division points is a
proper invariant subspace:
\[
 \ker\!\left(H^0(B^{\otimes M})\longrightarrow
       \bigoplus_{u\in A'[M]} B^{\otimes M}|_u\right)=0,
       \qquad A'=\mathbb C^g/L'.
\]
These evaluations compare the two principal generators even at a
fibre where the value at zero vanishes. A detecting division
point can depend on the fibre.

At identity level test equality of sections by the finite
surjective map induced by inclusion of lattices
\[
       \mathbb C^g/s\Lambda\longrightarrow\mathbb C^g/\Lambda,
       \qquad \mathfrak a=(s),\quad M=s^2,\quad s\ge2.
\]
The class evaluation keeps its original level and masses.
If the principal generators and
the identification are over \(K\), the two elements of the
one-dimensional \(K\)-space \(H^0(A,\mathcal L)\) have ratio in
\(K^\times\). Individual division coordinates may require a
further extension. The compatible lifts of translations and
isogeny pullbacks use the theta structure of
\cite[Section~1, Theorem~4]{mumford-theta-isogeny}.

\proofstep{A single class in the two readings.}
Let \(c\) be a definite source class, \(\Gamma_c\) its stabilizer,
\(v'\) its level volume, \(e_c\) its determinant coefficient,
and \(N_c\) its odd denominator prime to the marked level. Write
\(Q_2\) for the finite Fourier quotient, \(w_c(q)\) for its
coefficients and \(Z_c\) for the representatives of the finite
translation group projecting to the placed vacuum. If \(m_c\)
is the prescribed short-theta index, the class term is
\[
 \frac{v'e_c}{|\Gamma_c|m_cN_c^{d}|Q_2|}
 \sum_{q\in Q_2}w_c(q)\sum_{u\in Z_c}
                \bigl[P(\partial)F_{c,q,u}\bigr](0).
\]
Here \(F_{c,q,u}\) is the product of \(n\) theta factors
with the marked lifts of their translations and pullback by
\([N_c]\). Every translation is linear on the theta section,
so replacing its generator scales the product by \(t^n\).

On tangent coordinates \([N_c]\) acts by \(N_c\). If \(D_P\)
denotes \(F\mapsto[P(\partial)F](0)\), homogeneity and character
orthogonality give, respectively,
\[
 D_P(F(N_c\,\cdot))=N_c^dD_PF,\qquad
 \sum_{u\in Z_c}U(u)=|Z_c|\,\operatorname{pr}_{\rm vac}.
\]
Here \(U(u)\) are the translations with their prescribed inverse
character in the finite projector. Thus the derivative removes
\(N_c^{-d}\); the class index left after projection is
\(|Z_c|/m_c\). All other factors are common at this marked point.

In logarithmic coordinates, write \(Q(z)\) for the quadratic
change of polarization gauge on the product. The adjoint of
multiplication by a coordinate \(z_a\) in the jet pairing is
\(\partial_a\) acting on the symbol \(P\). Expanding the
exponential gives the finite identity
\[
 D_P(e^Q F)=
    \sum_{h=0}^{\lfloor d/2\rfloor}
       \frac{1}{h!}D_{Q(\partial)^hP}F .
\]
Each term of \(Q(\partial)\) is a scalar multiple of a mixed
contraction \(\sum_{b=1}^n\partial_{\tau,i,b}
\partial_{\bar\tau,j,b}\), which annihilates a joint harmonic.
Only the term \(h=0\) survives. Summing the class identity
proves the formula on the common fibre.
\end{proof}

\begin{auxlemma}[Cancellation of a change of theta generator]\label{lem:marked-vacuum-scalar}
At a common point as in Lemma~\ref{lem:marked-theta-fibre-comparison}
fix the same adelic oscillator vector \(f_0\), determinant
splitting character and finite argument \(h\) after forgetting
the additional torsion labels. Let the rational unitary elements
\(\gamma_1,\gamma_2\) place this lattice in the two presentations.
Write \(V_i\) for its analytic vacuum in the corresponding frame
and \(c_i\) for the scalar comparing it, in holomorphic coordinates,
with the \(n\)-fold theta section:
\[
                  V_i=c_i\vartheta_i^{\boxtimes n}.
\]
For matching class inputs one then has
\[
     c_2t^n=c_1,\qquad c_2\mathcal T_2(P)=c_1\mathcal T_1(P).
\]
\end{auxlemma}
\begin{proof}
Set \(\alpha=\gamma_1\gamma_2^{-1}\). A rational unitary isometry
preserves the rational Heisenberg subgroup and its additive
character. On functions on the adelic Heisenberg quotient it
acts by pullback; the determinant splitting character is trivial
on a rational determinant. Thus in the placed coordinates
\[
  \bigl(\omega(1,\gamma_1 h)f_0\bigr)(\alpha w)
     =\bigl(\omega(1,\gamma_2 h)f_0\bigr)(w).
\]
Here the two real Gaussian factors agree under \(\alpha\) by
isometry of the polarized period coordinates. Removing them and
transporting the finite marking gives \(V_1(\alpha w)=V_2(w)\).

At a marked point of the half-level torus where \(\vartheta_1\)
does not vanish, evaluate the tensor powers on the diagonal
\(n\)-tuple of that point. Lemma~\ref{lem:marked-theta-fibre-comparison}
supplies such a label \(u\) on each fibre. Represent it by a point of \(M^{-1}L'\) at infinity and zero translation at the finite places, keeping \(h\) fixed. In the common frame the equality of the vacua at that point reads
\[
 c_2 t^n\vartheta_1(u)^{\otimes n}=c_1\vartheta_1(u)^{\otimes n},
 \qquad \vartheta_1(u)^{\otimes n}\ne0 .
\]
Cancel this tensor and apply the scaling formula for the jet.
The proof allows \(t\) to vary with the framed point.

On a connected parameter domain, suppose the two vacua are
holomorphic in a common theta trivialization, with the same Hodge
weight and their Gaussian-weighted probability norms in constant
ratio. Near a nonvanishing torsion value their ratio is then
holomorphic with constant modulus. If \(R\) is the holomorphic ratio, then
\[
             0=\partial |R|^2=\bar R\,\partial R.
\]
Since \(R\ne0\), it is constant on the component. This last assertion requires the
comparison of metrics throughout the component.
\end{proof}

For the complex Fock reading, fix the sign \(\delta\in\{1,-1\}\)
of the additive variable and the positive real Gaussian scale
\(\xi_{\mathbb R}\). The substitution in the homogeneous
polynomial contributes the separate factor
\[
 P\!\left(-\delta(2\pi\xi_{\mathbb R})^{-1/2}z\right)
      =(-\delta)^d(2\pi\xi_{\mathbb R})^{-d/2}P(z).
\]

Two algebraic generators of unit norm for the same
Gaussian-weighted probability Haar metric have \(|t|=1\).
The scalar relating the Fock and harmonic maps also has modulus
one when their extremal vectors have equal norm and are
identified on the same irreducible source type.
Section~\ref{subsec:joint-family-estimate} recovers the resulting
common comparison scalar in the coefficient field from a nonzero
coordinate on the reference tuple. This argument assumes the
specified reference comparison and the two actual rational
readings. On each common fibre the lemma identifies the
individual class terms, hence also their ray-indexed sums.
Lemma~\ref{lem:placed-row-theta} realizes this common fibre on the covers of a fixed marked point. The following calculation tracks its phase when the point varies.


The covers also carry a finite calculation. Before choosing a character, keep each coefficient of the class sum in the regular basis of the ray group.

\begin{auxlemma}[The principal theta line at a placed row]
\label{lem:placed-row-theta}
At a stage of Lemma~\ref{lem:marked-period-domain}, let
\(g=(g_\infty,a)\) be a row argument and put
\(x=g_\infty x_0\), \(\Lambda_a=V\cap a\widehat\Lambda\).
Let \(A_g=(V_{\mathbb R},J_x)/\Lambda_a\), with polarization
\(\lambda:A_g\to A_g^\vee\), and let \(w_0\in\mathcal O_E\)
have trace one to \(F\) as in the half-trace convention.
Let \(\mathcal P\) be the rigidified Poincar\'e bundle and put
\[
  \mathcal L_g=(1,\lambda w_0)^*\mathcal P,\qquad
  \dim H^0(A_g,\mathcal L_g)=1 .
\]
Both vacua lie in this space. Each level refinement
\([x,aU_t]\mapsto[x,aU]\) retains the polarized abelian variety,
the line and the Hodge frame.
\end{auxlemma}
\begin{proof}
The homomorphism to the dual abelian variety associated with
the pullback is
\[
 \phi_{\mathcal L_g}=\lambda w_0+(\lambda w_0)^\vee
                   =\lambda(w_0+\bar w_0)=\lambda .
\]
Indeed the two mixed factors under translation are the pairings
with \(\lambda w_0\) and its dual. Simultaneous negation on
the arguments preserves the rigidified biextension. Thus the
line is symmetric and principal.

Write \(e\) for the alternating form of \(\lambda\). On
\(\Lambda_a\oplus\Lambda_a^\vee\) the Poincar\'e semicharacter
at \((u,\ell)\) equals \((-1)^{\ell(u)}\). Its polarization
modulo two is the alternating difference of the two evaluation
pairings, and the rigidifications make its restrictions to both
axes trivial. These conditions characterize it. Its diagonal
pullback consequently has sign
\[
 \alpha_a(u)=(-1)^{e(u,w_0u)},\qquad
 \frac{\alpha_a(u+v)}{\alpha_a(u)\alpha_a(v)}
   =(-1)^{e(u,(2w_0-1)v)}=(-1)^{e(u,v)} .
\]
In Weyl coordinates the central correction of the marked
lattice splitting is the half of the integer \(e(u,w_0u)\),
read by the additive character. This gives the same sign at
each lattice vector. The two Appell--Humbert factors have the
same Riemann form and semicharacter, and rigidification identifies
the lines. In particular both generators lie in the same
one-dimensional space \(H^0(A_g,\mathcal L_g)\).

For a refined level \(U_t\subset U\) use \([x,aU_t]\)
above \([x,aU]\). Its period lattice, complex structure and
diagonal Poincar\'e pullback are given by identical formulas.
The finite component \(a\) transports the standard residue
marking at each level and commutes with reduction. A change
of representative by \(\gamma\in G(\mathbb Q)\) transports
all the data by the same rational isometry.
\end{proof}

\begin{auxlemma}[One phase on the tuple covers of a row]
\label{lem:placed-row-tuple-phase}
In Lemma~\ref{lem:placed-row-theta} fix the even number \(n\)
of source columns and the degree \(d\) of the stage's joint
harmonic coefficient map. For a base component \(C\) there is
a \(u_C\in\mathbb C\), independent of the tuple \(t\), such that
its algebraic and adelic readings in matched frames satisfy
\[
 \operatorname{adelic}_t
   =(2\pi\xi_{\mathbb R})^{-d/2}u_C\,
     \operatorname{algebraic}_t,\qquad |u_C|=1 .
\]
Here the determinant normalization of the stage is fixed and
its lattice splitting factor has been extracted.
\end{auxlemma}
\begin{proof}
\proofstep{Compare the vacuum on the base component.}
In a symplectic basis of \(\Lambda_a\) let
\(\Omega=X+iY\) be the period matrix,
\(z=u+\Omega v\), and \(s=\dim A_g\). The semicharacter fixes
a half characteristic \((\epsilon,\eta)\). The principal
series and its squared probability norm satisfy
\begin{align*}
 \theta(z,\Omega)&=\sum_{m\in\mathbb Z^s}
 e^{\pi i(m+\epsilon)^t\Omega(m+\epsilon)+
                 2\pi i(m+\epsilon)^t(z+\eta)},\\
 \int_{[0,1]^{2s}}|\theta(u+\Omega v,\Omega)|^2
     e^{-2\pi v^tYv}\,du\,dv
 &=\int_{\mathbb R^s}e^{-2\pi t^tYt}\,dt
    =(\det(2Y))^{-1/2}.
\end{align*}
Integration in \(u\) kills distinct Fourier indices. Translating
the remaining cubes by \(m+\epsilon\) fills \(\mathbb R^s\).
Thus the probability-normalized generator in this unit theta
frame is \((\det(2Y))^{1/4}\theta\). In the self-dual adelic
lattice model the finite lattice has volume one; unfolding the
norm on the Heisenberg quotient leaves this same real integral,
once for each of the \(n\) factors. In particular the geometric
and adelic vacuum comparison factors \(C_C,\kappa_C\) against
the classical vacuum have equal absolute value.

To couple the theta frame to the varying Hodge frame, arrange
the first and second halves of the periods as the columns of
\(A(x)\) and \(B(x)\) in the chosen Hodge coordinates \(z\). Put
\[
 h(x)=A(x)^{-1},\qquad \Omega(x)=h(x)B(x),\qquad
 c_*=\left(\frac{2^s}{\prod_{j=1}^s2|\tau_j(\xi_0)|}\right)^{1/2}.
\]
Here \(\tau_j\) labels the \(j\)-th tangent direction and
\(\xi_0\) is the scalar in the alternating trace form.
In the rigidified holomorphic theta gauge choose the generator
\[
       S_x(z)=\beta_x\theta(h(x)z,\Omega(x)),\qquad
       \beta_x^2=c_*\det h(x).
\]
The periods are holomorphic in the frame-domain coordinates,
and the positive Riemann relations give \(\det A(x)\ne0\).
A local holomorphic square root exists. Its sign disappears in
the required even power,
\[
       \beta_x^n=(c_*\det h(x))^{n/2}.
\]

This prescribes the relative normalization between frames.
For example if all period coordinates are changed by
\(T\in\mathrm{GL}_s(\mathbb C)\), the new values are
\(A'=TA\), \(h'=hT^{-1}\), \(z'=Tz\). Thus
\[
         h'z'=hz,\qquad
         \frac{(\beta'_x)^n}{\beta_x^n}=(\det T)^{-n/2}.
\]
For separate Hodge eigenspace frame changes \(T\) is block
diagonal, so the determinant in this formula is the product
of their determinants. A holomorphic change of theta
trivialization multiplies both fibre expressions by the same
transition scalar.

The normalized Gaussian in the oscillator model has the same
Fourier series in symplectic coordinates. In a local choice of
square root its half-density acquires \((\det T)^{-1/2}\) per
column. Thus both \(n\)-fold expressions transform by the
single-valued character \(T\mapsto(\det T)^{-n/2}\).

In a unit Hodge frame the alternating area element has coefficient
\(2|\tau_j(\xi_0)|\) in the \(j\)-th complex coordinate.
Unimodularity equates the symplectic covolume to one, so
\[
  \frac{\det\Im\Omega}{|\det h|^2}
     =\frac{1}{\prod_{j=1}^s2|\tau_j(\xi_0)|},\qquad
  |\beta_x|^2=c_*|\det h|=(\det(2\Im\Omega))^{1/2}.
\]
The integral just computed now gives norm one. This calculation
uses the unit frame solely to measure the holomorphically
coupled generator and its even powers.

For a detecting division label, take the geometric translated
evaluation on \(n\) copies and the corresponding translated
periodized series, with polynomial \(P=1\). In holomorphic frames
both are holomorphic of the same determinant weight above, so
their quotient extends across each change of Hodge frame on its
nonvanishing set. In unit Hodge frames its modulus is constant
by the covolume and Gaussian identities. More explicitly, on
the nonvanishing set for a label \(t\), write the quotient as
\(q_t\). Its weight is zero. The two sections are proportional
on each fibre by Lemma~\ref{lem:placed-row-theta}, so
\[
      |q_t|=\mathrm{const},\qquad dq_t=0,\qquad
      q_t=q_{t'}\quad\text{where both readings are nonzero}.
\]
Lemma~\ref{lem:marked-theta-fibre-comparison} supplies a
detecting label at every point. The quotients therefore give
one constant on the connected component. The underlying
algebraic rule is homogeneous of degree \(n\) in the generator
of the direct-image theta line.

For descent, write a stabilizing isometry in the symplectic
basis adapted to the dyadic halves. It commutes with \(w_0^4\),
which reduces modulo \(8\) to the projector onto the first half.
Consequently
\[
 \gamma=\begin{pmatrix}A&B\\ C&D\end{pmatrix},\qquad
 \gamma\begin{pmatrix}1&0\\0&0\end{pmatrix}
   \equiv\begin{pmatrix}1&0\\0&0\end{pmatrix}\gamma\pmod8
   \quad\Longrightarrow\quad B\equiv C\equiv0\pmod8 .
\]
Every conjugate of the integral unit \(\det_E\gamma\) has
absolute value one, so it is a root of unity. The principal
congruence at level at least three forces it to be one. The
determinant on either dyadic half is the norm of the local
determinant. Applying this identity modulo \(8\) gives
\[
       \det A\equiv\det D\equiv1\pmod8,\qquad \chi_4(\det D)=1 .
\]
In the squared theta transformation the rigidified gauge
absorbs the Gaussian and the determinant cancels with the
change of \(\beta_x^2\), as in
\ref{cor:slice-marked-theta-cover-1}. Its \(n/2\)-th power
therefore descends to the marked quotient.

Changing the finite representative of a lattice multiplies the
adelic vacuum by the determinant splitting factor \(\mu\):
before splitting, its stabilizer acts by pullback and preserves
the lattice and half-trace sign. The infinite component transports
the Hodge frame, and rational transport preserves the Heisenberg
quotient and additive character. Thus after extracting \(\mu(g)\)
the adelic comparison factor is constant on the component. A change
of symplectic basis acts on both readings of the same rigidified
line, so the theta multiplier cancels in their quotient. The
resulting unit is \(\kappa_C/C_C\).

\proofstep{Match the character on the refined component.}
For the finite-order relative ray character \(\chi_t\) put
\[
       \xi_t(b/\bar b)=\chi_t(b).
\]
Hilbert's theorem 90 supplies \(b\), and two choices differ by a
base idele on which \(\chi_t\) is trivial. At a rational unitary determinant a rational lift has
character value one. To check the refined level, write \(z\)
for its local determinant. The following two lifts suffice at
the marked places:
\[
 \begin{array}{ll}
 z=(z_1,z_1^{-1})\text{ at a split place}:&
 b=(z_1,1),\\
 z\equiv1\text{ at an odd place fixed by conjugation}:&
 b=1+z,\quad b/\bar b=z,\quad b\equiv2 .
 \end{array}
\]
The marked half kills the split residue character, and \(2\)
belongs to the base field at the fixed place. Away from the
conductor a local unit has trivial character. In the unramified
quadratic host multiply any lift by a base power of the common
uniformizer to obtain a unit lift. These computations give the
level invariance also used in \ref{cor:completed-division-7}.

After an extension splitting the refined components and the
character values, write \(e_{C'}\) for their idempotents and
choose a representative \(a_{C'}\) on each. The determinant
section and its inverse are
\[
 D_t=\sum_{C'}\xi_t(\det a_{C'})e_{C'},\qquad
 D_t^{-1}=\sum_{C'}\xi_t(\det a_{C'})^{-1}e_{C'} .
\]

Changing the splitting character in the Weil representation
multiplies the row action by this same \(\xi_t(\det a)\).
It therefore cancels against the determinant section. On each
definite source class apply the common finite theta translations
and the harmonic derivative on the fibre of
Lemma~\ref{lem:placed-row-theta}. The finite Fourier projection,
derivative scaling and class mass agree by
Lemma~\ref{lem:marked-theta-fibre-comparison}. The two normalized
coefficient maps differ by a unit \(k\) on the multiplicity-one
joint harmonic type with the specified extremal product of minors.
A degree \(d\) Gaussian derivative contributes the additive
sign \((-\delta)^d\), \(\delta\in\{1,-1\}\), and the positive
factor stated in the lemma. For reference determinant one the
ratio on the joint harmonic type is therefore
\[
     (2\pi\xi_{\mathbb R})^{-d/2}\,
          k(-\delta)^d\,\kappa_C/C_C .
\]
It is independent of the class input and the tuple at this stage.

If an already chosen reference determinant has value \(a_0\)
and the new finite family has reference value \(b_0\), form the
disjoint index set consisting of a copy of the old reference and
all new labels, including any new copy of the reference tuple.
The values \(a_0,b_0\) are roots of unity. Over the finite
compositum containing them, the prescription is
\[
 D^{\rm merged}_{\rm old}=D_{\rm old},\qquad
 D^{\rm merged}_t=(a_0/b_0)D_t,\qquad (a_0/b_0)b_0=a_0 .
\]
The finite tests, Fourier weights and class masses are transported
by the coefficient inclusion. The common reference value is \(a_0\).
Its reciprocal enters the component factor. Restriction to the
distinguished old label retains its original coefficient map,
finite level action and framed reading. Each tuple may have its
own source classes, stabilizers and detecting torsion labels.
\end{proof}

\begin{auxlemma}[The common phase in the complex reading of a rational row]
\label{lem:placed-row-reading-phase}
Let \(S_K\) be the rational cusp space of the marked row in
Lemma~\ref{lem:placed-row-tuple-phase}. At its distinguished
embedding \(\iota\), write \(r_\iota(s;g)\) for the evaluation
of \(s\) in the prescribed frame at \(g\), \(C(g)\) for its
component and \(\mu(g)\) for its splitting factor. There is a
choice of the complex comparison map \(J_\iota\) on \(S_K\) for
which the phase outside the entire comparison is one. With
\(L_i=(2\pi\xi_{\mathbb R})^{-d/2}\), its lowest-type reading is
\[
 (J_\iota s)(g;v_0)
      =L_i\,c_{i,\iota}(C(g))\,\mu(g)\,r_\iota(s;g),\qquad
 c_{i,\iota}(C)=\frac{k_i(-\delta_i)^d}{a_{0,i}}\,
                       \frac{\kappa_C}{C_C}.
\]
This same function of the component is used for the reference
class sum and for every refined tuple after descent to the
base. The \(K\)-linear algebraic section attached to any class
input is the section before the displayed complex evaluation.
\end{auxlemma}
\begin{proof}
The finite class sum gives \(s\in S_K\) in the frames of the
algebraic moduli functor. Lemma~\ref{lem:placed-row-tuple-phase}
compares the vacuum, harmonic derivatives and finite translations
one term at a time. It gives the displayed function on the right
of the evaluation map. The scalar \(k_i\) belongs to the one
coefficient intertwiner for the entire stage, and \(a_{0,i}\)
to its common determinant normalization. For the class section \(s_t\) attached to tuple \(t\), and
its analytic theta lift \(\Theta_t\), the identities are
\[
        \Theta_t=J_\iota(s_t),\qquad
        J_\iota(bs)=\iota(b)J_\iota(s)\quad(b\in K).
\]
Thus \(J_\iota(S_K)\) supplies the rational structure on its
complex span. The finite-place seminorms are evaluated on \(s\)
in the algebraic frames.

There are two consistency checks on the component factor.
Let \(z\) be a level move and \(\tilde z\) an adelic lift.
The same vacuum evaluated in the moved marking and in the
argument \(g\tilde z\) gives the equation
\[
 c(C(g))\,\mu(g)\,r_\iota(zs;g)=
 \iota(\chi(z))\,
 c(C(g\tilde z))\,\mu(g\tilde z)\,r_\iota(s;g\tilde z).
\]
Indeed the level move transports the geometric theta translations
and the finite Fourier weights contragrediently. The lattice
stabilizer acts on the reference adelic vacuum by its determinant
splitting character, as in Lemma~\ref{lem:placed-row-tuple-phase}.
The residual character is the prescribed \(\chi\) of the marked
level group. Thus the factors attached to different components
already satisfy the required finite action, including when a move
permutes components. For the Petersson integral with its prescribed component measures,
\[
    \sum_C\int_C|c(C)F|^2=\sum_C\int_C|F|^2
                  \qquad (|c(C)|=1).
\]
The positive Lie factor and measure normalization are those of
the stage.

For another holomorphic coefficient embedding the equation is
imposed with that embedding's character values and frames.
To see the compatibility of these choices, divide its bare
frame-change factor by the required character and splitting
factor. For right moves, denote the resulting transition from
\(C\) to \(Cz\) by \(m(C,z)\), oriented so that the required
equation reads \(c(Cz)=m(C,z)c(C)\). On an orbit supporting
a nonzero reading the group actions and the vacuum law on its
stabilizer give
\[
 |m(C,z)|=1,\qquad
 m(C,z_1z_2)=m(Cz_1,z_2)m(C,z_1),\qquad
 m(C,h)=1\quad(Ch=C).
\]
Choosing a unit at one component and applying the transition
formula now prescribes \(c\) on the orbit, independently of the
move used. On an orbit of zero readings every unit satisfies
the equation. Complex conjugation supplies the conjugate reading.

If a later comparison retains the reference row, its component
factors have already been assigned. For the same normalized
oscillator vector and source frame the reference factor is
then recovered by evaluating at a division point at which the
principal section is nonzero, separately on each component.
More generally let \(b\) be one fixed rational unitary source
basis change preserving the prescribed source lattice and its
standard real Gaussian. It preserves the same one-dimensional
line of standard oscillator vacua. For its unit vector
\(\varphi_0\), the commuting source and row actions give
\[
 \omega(b)\varphi_0=u_b\varphi_0,\qquad |u_b|=1,\qquad
 \omega(b)\omega(g)\varphi_0=u_b\,\omega(g)\varphi_0 .
\]

On each component evaluate at a detecting division point. Its
nonzero theta tensor and the splitting and Gaussian factors
cancel. With the prescribed reference normalization on both
sides the two reference functions therefore satisfy
\[
     c_{\rm transported}(C)/c_{\rm old}(C)=u_b
                                  \quad\text{for every }C .
\]
Transporting the source frame itself gives \(b=1\) and \(u_b=1\).
This identifies the retained reference across components in the
second construction.
\end{proof}

\begin{auxlemma}[Descent of the theta class sum on a marked PEL model]
\label{lem:marked-class-section-descent}
Fix one of the holomorphic row stages, its polarized Hermitian
lattice, and its principal level. Pass to a neat refinement to
carry out the construction. Let \(K\subset\mathbb C\) contain the
reflex field and the fields of definition of the weight, the finite
test and the definite coefficient space. Use the rigidified
polarization line, its theta structures on the prescribed finite
cover and the \(E\)-linear Hodge frames. Use rational level volume
in the class sum, placing any archimedean scalar in the subsequent
Petersson reading. For a \(K\)-rational class input \(x\), the finite
harmonic expression of Lemma~\ref{lem:marked-theta-fibre-comparison}
defines an algebraic section \(s_x\) of its weight bundle on the
corresponding marked PEL space \(X/K\). Writing \(X_{\rm cl}\)
for the space of class inputs and \(\mathcal W\) for the weight bundle,
one obtains a \(K\)-linear map compatible with field extension:
\[
 X_{\rm cl}\longrightarrow H^0(X,\mathcal W),\quad x\longmapsto s_x,
 \qquad s_{1\otimes x}=1\otimes s_x\quad(K_1/K).
\]
\end{auxlemma}
\begin{proof}
Choose local Hodge frames and generators of the principal theta line
and split the finite torsion labels after a finite \'etale cover.
For one class \(c\) the operations on these charts are pullback by
\([N_c]\), finite theta-group translations and an algebraic fibre
jet. In characteristic zero the logarithm of the commutative
formal group, with linear term the chosen tangent frame, is defined
through the finite order \(d=\deg P_c\) by integrating the invariant
differentials term by term. For a multi-index \(\alpha\),
\[
    \partial^\alpha F(0)=\alpha!\,[z^\alpha]F,\qquad |\alpha|=d .
\]
This construction uses finitely many coefficients over the chart.

Under a change of polarized Hodge splitting the additional jet
terms are mixed contractions of \(P_c\); their vanishing is the
last identity of Lemma~\ref{lem:marked-theta-fibre-comparison}.
Indeed the difference of two splittings is linear from the tangent
space to its dual, symmetric for the polarization. Its quadratic
primitive in logarithmic coordinates pairs only opposite
\(E\)-eigenspaces, by the Rosati identity. On the \(n\) columns
it acts by the sum of these pairings.

For clarity, write \(t\) for a finite label, \(U_t\) for its chosen
theta translation and \(w(t)\) for the coefficient of the test in
the dual Fourier basis. Under a change of marking \(g\), these
bases transform contragrediently. In a new labelling with lifted
phase \(b_g(t)\) the paired terms transform as
\[
      U_t\longmapsto b_g(t)U_{g t},\qquad
      w(t)\longmapsto b_g(t)^{-1}w(t)
                 \quad\hbox{at the new label }g t .
\]
For a change of Fourier polarization the same assertion is the
invariance of the pairing between a finite vector space and its
dual. A Hodge frame change acts by the weight representation on the
harmonic symbol. Rescaling the generator by \(a\) multiplies the
expression by \(a^n\), its prescribed theta weight. Changes of
definite class representative permute the terms; the level action
on its coefficient is the inverse of that on its test. Conjugation
of a stabilizer preserves its order.

Let \(p_1,p_2:T\times_X T\to T\) be the two projections for the
resulting faithfully flat cover by frames and markings. Under the
transition isomorphism for the weight bundle \(\mathcal W\), the
chart formulas therefore satisfy
\[
           p_1^*s_T=p_2^*s_T
                 \quad\hbox{in }\Gamma(T\times_X T,\mathcal W).
\]
This proves descent. If individual torsion labels require a splitting
field \(K'/K\), the same identity for the Galois change of labels
descends the sum to \(K\), provided the class input and its finite
test are defined there. For a character valued only in \(K'\), the
individual character section is defined over the compositum of its
field of values with the fields of the model, coefficient and other
fixed data. Conjugating it conjugates these data and its character. The descent equality
and the sums are preserved by every further coefficient extension.
\end{proof}

\begin{auxlemma}[Cuspidality and the algebraic coefficient row]
\label{lem:marked-class-section-cuspidal}
Use the definite source of Lemma~\ref{lem:marked-class-section-descent}.
For an indefinite row write \(p_\tau,q_\tau\) for its signed
multiplicities at the real place \(\tau\). Suppose at a place with
\(p_\tau,q_\tau>0\) its positive-block harmonic partition has
last part \(\lambda^+_{\tau,p_\tau}\ge1\). In the prescribed
holomorphic stages, \(p_\tau=k+j\) and the last part is
\[
     \lambda^+_{\tau,p_\tau}
       =\left\lfloor e_{\tau,j}-b_{\tau,j}-(t+1)/2\right\rfloor
       \ge 1
\]
by the gap \(e_{\tau,j}-b_{\tau,j}\ge(t+3)/2\).
For an anisotropic row no weight inequality is needed. Let
\(\bar X\) be a smooth proper toroidal compactification at the
neat level, \(D\) its reduced boundary, and \(\mathcal W^{\rm can}\)
the canonical extension. The class input defines
\[
  s_x\in S_K=H^0(\bar X,\mathcal W^{\rm can}(-D)),\qquad
  S_K\otimes_K K_1\simeq S_{K_1}\quad(K_1/K\text{ a field extension}).
\]
These are finite-dimensional spaces. The character and good
idempotents defined over \(K\) commute with formation of their images
after base change.
\end{auxlemma}
\begin{proof}
Take a rational isotropic line in the row and split off its hyperbolic
plane. In the mixed Schr\"odinger model a theta summand has one
coordinate \(y\) in the definite source, including its prescribed
torsion shift. The centre of the line unipotent acts on this
coordinate through the additive character of its Hermitian norm.
Thus its integral on the compact central quotient is
\[
 \int_{[Z]}\psi\bigl(z\,H_{\rm def}(y,y)\bigr)\,dz
  =\begin{cases}1,&y=0,\\0,&y\ne0,\end{cases}
\]
with probability measure and the usual trace pairing.

Definiteness
equates vanishing of the frequency to \(y=0\). The finite level
can change which \(y\)'s occur. Every surviving archimedean
test is evaluation at \(y=0\) in the smaller oscillator with row
signature \((p_\tau-1,q_\tau-1)\). This evaluation commutes with
the definite unitary action. Writing \(\omega_{p,q}\) for the
oscillator at \(\tau\), the change of variables in \(y\) gives
\[
  \operatorname{ev}_0\omega_{p,q}(k)
       =\omega_{p-1,q-1}(k)\operatorname{ev}_0
          \quad(k\in U(n)),\qquad
          (p-1)-(q-1)=p-q .
\]
Thus the determinant shift of the vacuum is unchanged.

To see why that evaluation kills the prescribed type, remove the
common vacuum shift. The polynomial model on the remaining positive
and negative rows is contained in the tensor product of the
polynomial algebras on \(p_\tau-1\) copies of the standard definite
module and \(q_\tau-1\) copies of its dual, in the positive-row
convention. Write \(n\) for the definite rank and \(h=p_\tau-1\). The
Schur decomposition and the dual Pieri rule give
\begin{align*}
 \operatorname{Sym}\big((\mathbb C^n)^{\oplus h}\big)
   &=\bigoplus_{\ell(\nu)\le\min(n,h)}
           \mathbb S_\nu(\mathbb C^n)\otimes\mathbb S_\nu(\mathbb C^h),\\
 V_\gamma\otimes (\mathbb C^n)^\vee
   &=\bigoplus_{i:\,\gamma-e_i\ {\rm dominant}} V_{\gamma-e_i}.
\end{align*}
The first line bounds the number of positive highest-weight
entries by \(h\). Each dual tensor factor in the second line
can only lower a coordinate, so the bound persists.

The Fock polynomials are dense in their unitary completion, and
projection onto a compact type by its character is continuous.
Hence the same exclusion holds for all smooth vectors there. If \(W_\lambda\) is the source type, its \(p_\tau\) positive
entries imply
\[
 \operatorname{Hom}_{U(n)}
      (W_\lambda,\omega_{p_\tau-1,q_\tau-1})=0
\]
after removing the common shift on both sides. The zero-frequency
term therefore has zero projection on this type at every rational
isotropic line.

The same calculation applies to translated tests. For a rational
parabolic choose a line in the smallest nonzero space of its
isotropic flag. The Hermitian translations supported on this
line form a central subgroup of its unipotent radical. Integrating
first over this subgroup therefore annihilates the constant term.

The Gaussians on the definite lattice times the fixed polynomials
give normally convergent theta sums on compact subsets of the
domain. At a boundary chart their Fourier indices are positive
semidefinite Hermitian Gram matrices. Let \(H\ne0\) generate a rational semidefinite toroidal ray.
Vanishing order zero for a semidefinite index \(T\) forces
\[
  0=\operatorname{tr}(TH)
    =\|T^{1/2}H^{1/2}\|_{\rm HS}^2
    \quad\Longrightarrow\quad \operatorname{im}H\subseteq\ker T .
\]
At a multivariable cusp, sum these nonnegative traces over all real
embeddings. The rational image contains a line, already excluded by the
zero-frequency calculation. The Fourier--Jacobi expansion has nonnegative
exponents and vanishes at every boundary divisor.

In terms of the toroidal comparison this gives the extension to
\(\mathcal W^{\rm can}_{\mathbb C}(-D)\).
\cite[Section~1, p.~1788]{lan-stroh-subcanonical}
There is no boundary on an anisotropic component. Uniqueness of the
extension and faithful flatness descend it to \(K\): a section on
\(X\) belongs to the subspace of extendible sections if its scalar
extension does, since field extension is exact.

Properness gives finite dimension. For the base-change assertion,
compute sections as the equalizer on a finite affine cover of the
separated compactification and its pairwise intersections. Flatness
of \(K_1/K\) preserves that equalizer. For any idempotent \(e\)
defined over \(K\), taking its image commutes with tensor product:
\[
    S_K=eS_K\oplus(1-e)S_K,\qquad
    (eS_K)\otimes_K K_1=e(S_K\otimes_K K_1).
\]
At the original level this gives the
prescribed invariant or character subspace of the neat refinement.
\end{proof}

\begin{auxlemma}[The full holomorphic multiplicity space at a fixed level]
\label{lem:marked-full-holomorphic-row}
With \(S_K\) as in Lemma~\ref{lem:marked-class-section-cuspidal},
fix the holomorphic embedding, the regular holomorphic discrete
series \(D_\infty\) and its lowest type \(W\) as in
Lemma~\ref{lem:marked-lowest-reading}. Let \(\mathcal H_{\rm cusp}\)
be the full cuspidal Hilbert space in the same PEL genus, with the
same finite level and character. Unitary-frame evaluation, in the
dual weight convention, gives a complex-linear isomorphism
\[
 S_K\otimes_{K}\mathbb C
       \ \simeq\
 \operatorname{Hom}_{G(\mathbb R)}
                 (D_\infty,\mathcal H_{\rm cusp}) .
\]
It restricts to the images of corresponding finite idempotents.
The conjugate embedding gives the conjugate isomorphism on the
antiholomorphic weight and conjugate finite character.
\end{auxlemma}
\begin{proof}
The characteristic-zero comparison for the subcanonical extension
identifies its global sections over \(\mathbb C\) with the holomorphic
cuspidal automorphic forms of the bundle weight.
\cite[Section~1, p.~1788]{lan-stroh-subcanonical}
In the unitary frame such a section is a copy of \(W\) of smooth
cuspidal functions annihilated by \(\mathfrak p^-\). Its derivatives
are square integrable, by cuspidal rapid decay at this fixed finite
level and infinitesimal character. The lowest-type construction of
Lemma~\ref{lem:marked-lowest-reading} therefore produces a bounded
intertwiner of \(D_\infty\). It is unique by cyclicity.
At a compact real factor this is the elementary identification
\[
          D_\tau=W_\tau,\qquad \mathfrak p^-_\tau=0 .
\]

Conversely suppose \(A:D_\infty\to\mathcal H_{\rm cusp}\) is an
intertwiner and put \(F_w=A(w)\) for \(w\in W\). At the fixed finite
level these are smooth automorphic functions: smoothness of \(w\)
gives all local \(L^2\) derivatives of \(F_w\), and local Sobolev
regularity gives smooth representatives. They satisfy
\[
 F_w(gk)=F_{kw}(g),\qquad
          R(Z)F_w=A(Zw)=0\quad(Z\in\mathfrak p^-).
\]
The first formula is exactly the descent rule for the associated
bundle. The second is its holomorphy equation in unitary frames.
The functions are cuspidal by their ambient representation.
The same characteristic-zero comparison identifies the resulting
holomorphic cusp form with a unique subcanonical section. The two
assignments recover the same functions \(F_w\) on all frames and
hence are inverse. Each finite correspondence preserving the
bundle and level commutes with the right real action. This gives
the assertion on corresponding finite idempotents.
\end{proof}

\begin{auxlemma}[The rational quotient generated by the same theta classes]
\label{lem:marked-class-kernel-quotient}
Fix the definite \(K\)-space \(X_{\rm cl}\) of class inputs at one
level and weight of Lemma~\ref{lem:marked-class-section-descent}.
Use the PEL dictionary of Lemma~\ref{lem:marked-period-domain} for
these classes. Use its polarization, marked finite tests, class masses
and the common equivariant harmonic polynomial model in both
the algebraic formula \(A:x\mapsto s_x\) and its analytic theta
operator \(B\), with the same finite level projection. At the chosen
embedding of \(K\) into \(\mathbb C\), write \(A_{\mathbb C}\)
for the scalar extension. Then
\[
  \ker B=\ker A_{\mathbb C}
       =\mathbb C\otimes_K\ker A,\qquad
  (X_{\rm cl}/\ker A)\otimes_K\mathbb C\simeq\operatorname{im}B .
\]
\end{auxlemma}
\begin{proof}
For a dictionary point \(z_g\) choose the unitary frame used by the
lowest-type evaluation. Its analytic vacuum and its geometric
principal theta generator span the same one-dimensional space on
the marked fibre. The nonvanishing test at division points of
Lemma~\ref{lem:marked-theta-fibre-comparison} recovers the nonzero
ratio there. The pullback, finite translations and harmonic jets in
each class term are linear in the same tensor of \(n\) vacua.
The weight is fixed throughout \(X_{\rm cl}\). Its real harmonic
coefficient and Fock coefficient maps differ by a fixed nonzero
scalar on this irreducible type, since their extremal vectors give
the same flag-minor polynomial up to normalization. Hence the factor
relating the two evaluations at this frame is independent of \(x\):
\[
       (Bx)(g;v_0)=m(g)\,\operatorname{rd}_{g}(s_x),
                \qquad m(g)\in\mathbb C^\times .
\]
This calculation uses the actual marked determinant and Fourier
weights in both class sums. A level projection can equivalently be
taken on the source characters before applying either map.

As \(g\) varies, the dictionary covers the complex components in
this genus and all the unitary frames. The compact-group orbit of
\(v_0\) spans the lowest type. Vanishing of these evaluations is
therefore equivalent both to vanishing of \(s_x\) and to vanishing
of \(Bx\), by Lemma~\ref{lem:marked-lowest-reading}. This proves
the first equality without requiring the factors \(m(g)\) to be
algebraic. Tensoring the \(K\)-linear class map with the field
\(\mathbb C\) preserves its kernel, which proves the second.
\end{proof}

Thus the generic spectral span on these matched classes has the
rational form \(X_{\rm cl}/\ker A\). Every subspace cut out by
\(K\)-rational type projectors intertwining the two class maps
inherits the same description. The coordinates and the integral norms of this rational form on
the marked row are constructed in Lemma~\ref{lem:constructed-signed-coefficient-row}.

\begin{auxlemma}[Full span of the integral cusp lattice of a row]
\label{lem:placed-row-full-cusp-ball}
At any of the selected row stages let \(r\) be the row rank,
\(N=g_0g_1\) its squarefree odd principal level, \(F\) its placed
coefficient field and \(E\) its CM field. Give the algebraic cusp
space \(S_F(a,n)\) its dominant harmonic exponents \(a\) and its
positive even theta column count \(n\). For each \(f\in S_F(a,n)\)
there is a \(c_f\in\mathcal O_F\setminus\{0\}\) such that
\[
      v_R\bigl(c_ff(y)\bigr)\ge0
\]
at every framed semistable test \(y\) over a mixed-characteristic
complete discrete valuation ring with finite residue field, over
every finite place of \(F\). In particular the full geometric
local unit ball of the entire cusp space has rational span
\(S_F(a,n)\). The assertion holds on restriction to a rational
type subspace of this space.
\end{auxlemma}
\begin{proof}
\proofstep{Increase the number of theta columns on the same row.}
A cusp section is a scalar rule on the Hodge and theta frame
bundle, homogeneous in its two kinds of frames. Therefore the
pointwise power belongs to the weight and column count
\[
       f^k\in S_F(ka,kn),\qquad k\in\mathbb Z_{\ge1}.
\]
Dominance is preserved. On a toroidal chart a cusp section belongs
to its boundary ideal times the canonical extension; its \(k\)-th
power lies in the \(k\)-th boundary ideal, hence still vanishes
on the boundary. The implication also follows from multiplying
the nonnegative Fourier--Jacobi expansions. On a compact
component the cusp condition is empty. Choose, for example,
\[
        k=(4[E:\mathbb Q]r+4)|\mu(E)|.
\]
Write \(\epsilon_{a,n}\) for the action on the scalar rule
of a central scalar \(u\) stabilizing the lattice. Such a
scalar lies in \(\mu(E)\), and the even column count makes the
half-column exponent integral. Thus
\[
 kn\ge 4[E:\mathbb Q]r+4,\qquad 2\mid kn,\qquad
 \epsilon_{ka,kn}(u)=\epsilon_{a,n}(u)^k=1 .
\]

\proofstep{Clear the high-column power on complementary sets of primes.}
Keep the original marked row. In its theta comparison use the
identity level cover, the unit ideal for the added half and the
constant section \(1\) for its relative determinant. These choices
give the trivial relative character in every component. The
rational Hermitian space and half-trace polarization are those
of Lemma~\ref{lem:marked-period-domain}. The stabilizer calculation
at \ref{cor:slice-marked-theta-cover-1} identifies the vacuum
multipliers. These are the geometric inputs of the stable-range
span in Lemma~\ref{lem:theta-overlap-invertible} and
Proposition~\ref{prop:span-original-theta}, before applying a
good-type projection.

For every \(b_0\ge1\) its short theta expansion can be taken
with auxiliary unimodular lattices locally standard at
the primes of \(b_0N\), as in the two expansions
at \ref{cor:deep-cover-denominators-9}. Let \(W\) be their
\(F\)-span. Faithful scalar extension gives the descent
\[
  [f^k]\otimes1=0\text{ in }
  \bigl(S_F(ka,kn)/W\bigr)\otimes_{F,\iota}\mathbb C
                \quad\Longrightarrow\quad f^k\in W .
\]
For each of its finitely many short forms \(g\), let \(n_g\)
be its positive odd denominator. The local short-rule calculation
at \ref{cor:deep-cover-denominators-10} gives one nonzero clearing
scalar for all its integral framed evaluations at the primes
\[
                  p\mid b_0\quad\text{or}\quad p\nmid n_g .
\]
It is a calculation on the local lattice and the semistable
theta lift. It includes \(p\mid N\), since the auxiliary lattice
is locally standard there and the fixed \(p\)-power part of
the marked level is included in the clearing scalar.

Use first \(b_0=1\), clear the coefficients of this finite
expansion, and take the product \(B\) of the integers \(n_g\)
in that expansion. This gives \(c_1\in\mathcal O_F\setminus\{0\}\)
clearing \(f^k\) outside \(B\). Use next \(b_0=B\) on the same
\(f^k\) to obtain \(c_2\in\mathcal O_F\setminus\{0\}\) clearing
it at all primes dividing \(B\). For \(c_f=c_1c_2\) and a
nonzero reading,
\[
       k v_R(c_ff(y))
          =v_R(c_ff(y)^k)+(k-1)v_R(c_f)\ge0 .
\]
A zero reading is integral. Choosing such scalars on a basis
proves the span assertion, also for a rational subspace.
At a fixed coefficient place, finitely many independent
integral-frame evaluations bound the other inclusion by the inverse
of their evaluation matrix, as in
Lemma~\ref{lem:marked-bounded-denominator}. A module between two
full lattices over a discrete valuation ring is finite free of the
same rank.
\end{proof}

The positive column count in this lemma is the actual even rank
of the definite source at the selected lift. The argument applies
also when a signature multiplicity is zero. Proposition
\ref{prop:deep-cover-full-lattice} treats a different initial
weight, with zero central exponent and a positive gap between
two nonempty signature blocks. There the determinant and flat
multipliers provide columns which a power of a zero-column
weight could never supply.

The pullback of a base section to a refinement admits a second
description of its unit ball. Let \(U'\subset U\) be a finite
torsion-level refinement, \(y\) an integral test on the base,
and \(\pi:X_{U'}\to X_U\) the forgetful map. Over the fraction
field the torsion is finite \'etale. Its compatible full
markings form a nonempty finite scheme there. A closed point
therefore supplies a finite extension and a lift
\[
      K=\operatorname{Frac}(R)\subset K',\qquad
      y'_{K'}\in X_{U'}(K'),\qquad
      \pi(y'_{K'})=y_K\otimes_K K'.
\]

The integral test asks for the semistable identity model, its
saturated eigentangent frames and a generator of its direct-image
theta line, with the additional marking on the generic fibre.
Thus the same semistable model after extension is a test
upstairs. For an eigenlattice defined by the simultaneous
kernel of endomorphisms \(D_1,\ldots,D_m\) on \(M\), flatness
of \(R'\) over \(R\) gives
\[
 \left(\bigcap_{j=1}^m\ker(D_j:M\to M)\right)\otimes_R R'
       =\bigcap_{j=1}^m\ker(D_j\otimes1:M_{R'}\to M_{R'}) .
\]

The cubical line comes from the same rigidified polarization.
Its generator stays primitive under extension by
Lemma~\ref{lem:marked-hecke-semistable-unit}.
Normalized by \(v_{R'}(\pi_v)=v_R(\pi_v)=1\), the two evaluations
satisfy
\[
              v_{R'}\bigl((\pi^*f)(y')\bigr)=v_R(f(y)).
\]
Conversely every upstairs test forgets to a base test. This
proves that integral readings pull back and are detected by this
refinement, including when the level is divisible by the residue
characteristic. It compares the old subspace with its pullback.
To apply the spanning argument to all cusp forms at a new
principal level, the separate stable-range and short-rule
conditions at that level must hold.

\begin{auxlemma}[Constructing the signed coefficient row and its lattice]
\label{lem:constructed-signed-coefficient-row}
On the marked PEL row of Lemma~\ref{lem:marked-full-holomorphic-row}
use the prescribed finite character and the following evaluations
at the coefficient field \(K\). For \(v\) finite, let
\(\mathcal P_v\) consist of the valuation profiles of the admissible
integral frames over complete discrete valuation rings with finite
residue field and coefficient embedding inducing \(v\). This includes
the semistable frames admitted by the moduli problem. Write \(|s(y)|_v\)
for the norm of the framed scalar, using the valuation extended
with its given normalization, and set
\[
 B_v=\{s\in S_K: |s(y)|_v\le1\ \text{for every }y\in\mathcal P_v\},
 \qquad B=\bigcap_{v\ {\rm finite}} B_v .
\]

An arbitrary \(K\)-basis of \(S_K\) identifies the initial rational
row with a coordinate space and carries its integral tests to these
very balls. At a valuation ring \(O'\) over \(v\) use
\[
 B_{O'}=O'\cdot(1\otimes B_v),\qquad
 B_{O',V}=B_{O'}\cap V_{K'}
\]
for the generic type or good subspace \(V_{K'}\)
over \(K'=\operatorname{Frac}(O')\). This defines scalar extension from
the original ball. Evaluating every new \(K'\)-section anew is a
separate operation.

At the base level of odd
squarefree modulus with the prescribed common components and
stable-range short theta tests of positive column degree, \(B\) is a full lattice over
\(\mathcal O_K\). The corresponding lattice in the complex
multiplicity model is obtained by inverse image under its
Petersson coordinate map.
\end{auxlemma}
\begin{proof}
Let \(b_\alpha\) be a basis of \(S_K\). For an auxiliary
coefficient field \(L\supset K\) take the coordinate map
\[
 \iota:\sum_\alpha x_\alpha b_\alpha\longmapsto
                  (x_\alpha)_\alpha\in L^{\dim S_K}.
\]
Its inverse on its \(K\)-image reads these same coefficients as a
section. For an admitted embedding \(\sigma\) and the framed
norm one has
\[
  J_\sigma(\iota(s))=\sum_\alpha\sigma(x_\alpha)J_\sigma(b_\alpha),
  \qquad \|\iota(s)\|_y=|s(y)|_v .
\]
The conjugate construction uses
the conjugate lowest weight and finite character. A level
operator preserving the bundle acts through its matrix on the
\(b_\alpha\). Changes of an allowed integral level frame preserve
the collection of all valuations, by applying the inverse move.
Thus their geometric actions preserve \(B\) and each \(B_v\).

For fullness, take a finite set of detecting CM evaluations as in
Lemmas~\ref{lem:marked-generic-cm-separation}
and~\ref{lem:deep-cover-detecting-frames}. Put their fields of
definition in a single number field \(K_1\). At each such point
the CM abelian variety acquires good reduction after a fixed
finite extension. Its chosen algebraic frame differs from an
integral frame at only finitely many places; on the fixed weight
the changes have a common denominator. Thus for some nonzero
integer \(c\),
\[
       B\ \hookrightarrow\ c^{-1}\mathcal O_{K_1}^{\,h}
\]
by these evaluations. The right side is a finite
\(\mathcal O_K\)-module, and consequently so is \(B\).

For rational spanning, Lemma~\ref{lem:placed-row-full-cusp-ball}
puts an integral multiple of each basis section in \(B\). Thus
\(KB=S_K\). At a fixed finite place \(v\), localization
inside the original rational space satisfies
\[
                     B_v=\mathcal O_{K,(v)} B .
\]
Indeed, given \(s\in B_v\), choose a nonzero \(c\in\mathcal O_K\)
with \(cs\in B\). The Chinese remainder theorem gives
\[
 h\equiv1\pmod v,\qquad
 h\in\prod_{w\ne v,\ w\mid(c)}w^{\operatorname{ord}_w(c)} .
\]
Then \(hs\) is integral at \(v\) by assumption and at all other
places by construction. Hence \(s=h^{-1}(hs)\) lies in the
localization. Conversely, a denominator prime to \(v\) acts by a
unit on the integral evaluations there.

Let \(V\) be the generic good character subspace and let
\(J\) denote its complex lowest-type reading. The measured
multiplicity coordinates may be chosen by an isomorphism
\[
 U:\mathbb C^{\dim V}\overset{\sim}{\longrightarrow}J(V_{\mathbb C}),
       \qquad U^*U=n_P I ,
\]
where \(n_P>0\) is the prescribed level index in the Petersson
measure. Take an orthonormal basis of the image and scale its
vectors by \(\sqrt{n_P}\). Define the integral subgroup in these
coordinates by
\[
            \Lambda_{\rm mult}=U^{-1}J(B\cap V).
\]
The intersection is full in \(V\): clear the finitely many
denominators of a rational basis of \(V\) in \(B\), and use
finite generation of submodules of \(B\). This proves both
the norm factor and the membership equivalence in the measured
model. Localized good lattices use the displayed extension and
intersection operations. In dimension zero these rows and all their lattices are zero.
\end{proof}

The initial rational row, its ball, and the lattice in this
measured image can therefore all be chosen from the same framed
cusp sections. The geometric inputs to this construction concern
the allowed marked moduli charts, their complete sets of
valuation tests and the realization of the theta class data and
Hecke correspondences on them at the prescribed levels.

\begin{auxlemma}[Integral comparison on a prime-to-characteristic refinement]
\label{lem:copied-ray-level-correspondence}
Let \(R\) be a strictly henselian discrete valuation ring of
residue characteristic \(\ell\), and let \(x\) be a marked
principally polarized \(\mathcal O_E\)-abelian scheme over \(R\).
Suppose the added level of one test has order \(N\) prime to
\(\ell\), and that its frame orbit is nonempty on the geometric
generic fibre. If \(\mathcal F_x\) is its finite \'etale scheme of additional markings, then
\[
                 \mathcal F_x(R)\longrightarrow \mathcal F_x(\bar k)
                            \quad\text{is a bijection},
                 \qquad \mathcal F_x(R)\ne\varnothing ,
\]
where \(\bar k\) is an algebraic closure of the residue field.

Suppose its Hecke correspondence at a prime \(p\)
is given by the admissible \(p\)-power subgroups of \(x\), where
\(p\nmid N\), and has integral coefficient transport in the
old frames. For \(p\ne\ell\) its branches lift with the old
coset multiplicities. For \(p=\ell\), assume that the old
integral correspondence is finite locally free over its source
and acts by a trace with the specified integral normalization.
In both cases the induced correspondence acts integrally in
the pulled-back frames.
\end{auxlemma}
\begin{proof}
The polarization pairing and the \(\mathcal O_E\)-module
structure give locally constant data on \(x[N]\), a finite
\'etale group scheme over \(R\). The prescribed orbit of
compatible level frames is consequently finite \'etale.
Its coordinate algebra is
\[
                    R^d,\qquad d>0 ,
\]
since \(R\) is strictly henselian and the generic fibre is
nonempty. Every chosen frame over the geometric generic fibre
comes from a section over \(R\).

For a Hecke kernel \(C\subset x[p^s]\) the quotient isogeny
\(q_C:x\to x/C\) induces an isomorphism
\[
                x[N]\xrightarrow[\;\sim\;]{q_C}(x/C)[N].
\]
Hence the frame at each target is determined by the source
frame. Its pairing is taken with the prescribed polarization
similitude of the branch. Forgetting the \(N\)-frame recovers
the original kernel \(C\). With a fixed source frame count the
kernels as labels of the original double coset, retaining the
stack stabilizers in any class integral. For \(p\ne\ell\) each
such branch occurs with the old coefficient matrix.

For \(p=\ell\), the construction applies to finite flat kernels
over every test algebra. It identifies the refined source
correspondence with the base change of the old one along the
frame cover. For its finite locally free coordinate algebra
\(B/A\) and an \(A\)-algebra \(A'\) one has
\[
 \operatorname{Tr}_{B\otimes_A A'/A'}(b\otimes a')
       =a'\operatorname{Tr}_{B/A}(b).
\]
Indeed the multiplication matrix of \(b\otimes1\) is obtained
from that of \(b\) by change of coefficients in any local
basis. The same trace and coefficient transport thus give the
integral operator on the refined row, including its original
normalization.
\end{proof}

For a finite ray group \(\Delta\), the universal coefficient
module may be read in its regular basis. If its translation by
\(t\in\Delta\) sends \(a_\delta\) to \(a_{\delta-t}\), then
\[
 \sum_{\delta\in\Delta}\varphi(\delta)a_{\delta-t}
    =\varphi(t)\sum_{\delta\in\Delta}\varphi(\delta)a_\delta .
\]
This proves the asserted character covariance in the integral
coefficient ring, even when \(\ell\mid|\Delta|\). The inverse
character convention results from translating by \(-t\).
These level and covariance computations apply wherever the
given integral PEL model, source classes and old framed
correspondence have been realized. At a residue characteristic
dividing the added level the required local calculation is the
one for that characteristic.

\begin{auxlemma}[Specialization at the same ray characters]\label{lem:marked-ray-specialization}
Let \(R\to R'\) be a finite extension of discrete valuation rings,
\(\Delta\) a finite abelian ray quotient, and \(j:M\otimes_R R'\to M'\)
an \(R'\)-linear map of the marked coefficient lattices. For a sign and a
fixed test write the two universal series of
Lemma~\ref{lem:universal-theta-type} as \(\Theta=\sum_\delta a_\delta[\delta]\)
and \(\Theta'=\sum_\delta a'_\delta[\delta]\). If the class sums in the
matched frames give \(a'_\delta=j(a_\delta\otimes1)\) for every
\(\delta\), then for every finite extension \(S/R'\) of valuation rings
and every allowed character \(\xi:\Delta\to S^\times\),
\[
 (j\otimes_{R'}S)(\Theta(\xi))=\Theta'(\xi),\qquad
              \Theta(\xi)=\sum_{\delta\in\Delta}\xi(\delta)a_\delta .
\]
For corresponding tests, this identifies their generated images
under the comparison. A common integral denominator before
specialization also clears each specialized coefficient.
\end{auxlemma}
\begin{proof}
The map \(S[\Delta]\to S\) given by \([\delta]\mapsto\xi(\delta)\)
is an algebra homomorphism. Applying it to the equality of universal
series gives the formula. Taking the \(S\)-span over the test list
commutes with applying \(j\otimes S\). Each scalar \(\xi(\delta)\)
is a unit in \(S\), so a denominator clearing the coefficients of
the finite sum also clears its value. This argument applies even at
primes dividing \(|\Delta|\).
\end{proof}

Thus a comparison of the actual class sums for each sign covers even
the finite characters chosen after the first evaluation. It requires
the marked specialization map and the compatible test frames; its
conclusion concerns the generated image, with any saturation or index
at the rebuilt level still measured on that level.

Section~\ref{subsec:selected-word-level} computes the two
masses and the discriminant and Weyl factors, and in its stated weight ranges gives the explicit mass estimate 
\[
 \log\!\left[\left(\frac{D_*}{d_{u,*}\Gamma^2}\right)^{1/(fk)}\right]\ge \tfrac12 mL.
\]

The factor \(d_{u,*}\) is a positive integer even before imposing
the level at \(T_u\). Write \(g_0\) for the auxiliary prime
at the common type level used to impose the marking. Reduction of the marked lattice at the common type prime sends the class stabilizer \(A\) into a unitriangular group. If the local rank is \(n\), the target has order
\[
 |U_n(\mathbb F_{g_0})|=g_0^{n(n-1)/2}.
\] For the finite stabilizer \(A\), let \(A_1\) be the kernel of reduction. It is a finite subgroup of the pro-\(g_0\) kernel, while \(A/A_1\) embeds into the unitriangular group. Consequently
\[
 |A|=|A_1|\,|A/A_1|\quad\text{is a power of }g_0 .
\]
If a nonidentity \(a\in A\) has an eigenvalue of order \(g_0^e>1\),
its characteristic polynomial at the completely split place contains
the corresponding local cyclotomic factor, so
\[
           k\ge [\mathbb Q_{g_0}(\zeta_{g_0^e}):\mathbb Q_{g_0}]
                     =g_0^{e-1}(g_0-1)>k .
\]
Thus every eigenvalue of \(a\) is one.

Since \(a\) has finite order
in characteristic zero, it is semisimple and equals the identity. Write \(M_k\) for the genus
mass and \(I_k\) for its level index. The count of these free
classes is \(M_k I_k\). The Weyl dimensions multiplying
it are positive integers. The positive real number \(D_*\)
uses the common row formal degrees and the level index.
The ray class count supplies split generators in the specified
intervals with norm \(T_u\) and the congruence at the common
level. The lower bound for \(m\) places these levels beyond \(X'\), uniformly once \(c_s x^{2\epsilon/3}>1\):
\[
 \log X'=x^{\epsilon/3}<c_s mL\le\log T_u.
\]
\paragraph{Integral lattices at the chosen levels.}
Take the numerical size data from the compatible system at
these levels. Complete them by the prime exponents
of its level and period lattices. Its resulting weight
functions agree with the base data, its field has the
prescribed identification, and its binary sequences give
the prescribed shifts at every interior coordinate.
The local prime assignments are obtained from the valuations
on that same system.

Fix a sequence and write \(M\) for the full ramified column
lattice in its rational span. There are two sorts of exceptional
local lattices. At a place over the prime assigned to the sequence, let \(c_w\) be the column map and \(L_{\rm marked}\) the level lattice. The prescribed inverse image is
\[
 c_w^{-1}(L_{\rm marked})=\{z:c_w(z)\in L_{\rm marked}\}.
\] At a
comparison place selected for an orbit condition, take the
inverse image of the specified integral lattice from
Section~\ref{subsec:column-orbit-length}. These prescriptions
are disjoint. The latter places have rational prime
\(p\le X'\), and the coefficient field has finite degree.
For a finite place \(w\), write \(M_w\) for the localization
of \(M\). At each exceptional place the specified lattice
and \(M_w\) are commensurable by the local comparisons.
At the places selected for orbit conditions, the specified lattice lies in \(M_w\).
\paragraph{Gluing the local lattices.}
\begin{lemma}[Lattices with prescribed localizations]\label{lem:word-lattice-gluing}
Let \(M\) be a full lattice over the ring of integers of a number
field. Prescribe full lattices commensurable with \(M_w\) at a finite
set of places \(w\), in the localizations of the same rational
space. If \(A\) is that set and \(L_w\) is the prescribed lattice at \(w\in A\), there is a unique full lattice \(L\) with
\[
 L\otimes O_w=\begin{cases}L_w,&w\in A,\\ M_w,&w\notin A.\end{cases}
\]
\end{lemma}

\begin{proof}
Choose one nonzero rational integer \(d\) clearing the
finitely many local inclusions in both directions. Each
lattice then lies between \(dM_w\) and \(d^{-1}M_w\).
The finite module \(d^{-1}M/dM\) is the direct sum of its
primary components over the ring of integers of the
coefficient field. In those components take the images of
the specified lattices and elsewhere the images of \(M_w\).
The inverse image of their sum in \(d^{-1}M\) has exactly
the desired localizations. A lattice is the intersection of its
localizations in the rational space, which also gives uniqueness.
\end{proof}

Apply Lemma~\ref{lem:word-lattice-gluing} to the two disjoint kinds of
places just described. The order used in the size data
away from the prime assigned to the sequence is its local colength in \(M_w\)
divided by the ramification index \(e_w\). Multiplication
of this glued lattice by the fractional ideal of
Section~\ref{subsec:column-orbit-length} gives the shifted
column. The height on it is the numerical height of the
same sequence, so the subsequent rank mean in
\eqref{eq:marked-frame-height} refers to these columns.
These equalities are needed even for
arithmetic bounds invariant under relabeling, since the
Frobenius, deleted character and valuation labels all occur
in the same sums. \label{cor:window-points-3}

\subsubsection*{Starting the automorphic construction at the assigned levels}
\label{subsec:automorphic-realization}
\hypertarget{source-automorphic-realization}{}

The field, levels and weight functions have been chosen at a translate. We now give the Hecke and character calculations used in Construction~\ref{hyp:family-prescribed-realization}. The first set of signed tests selects the characters. The second calculation determines their actual spaces and local orders on the rebuilt member.

\begin{paramdefinition}[Fields, frames and levels at a translate]\label{hyp:initial-adelic-frame}
Fix the auxiliary CM field \(E\) of Definition~\ref{def:family-original-host},
one translate, its two weight functions, and all the
primes assigned to the sequences. Denote the small field and the two coefficient
fields by \(K_s\subset K_0\subset K\subset\mathbb C\), in
their distinguished embeddings. The underlying Hermitian spaces,
the marked split identifications, the two deletion characters
\(\delta,\delta'\) of Lemma~\ref{lem:split-deletion-characters}, and the levels at the primes assigned to the sequences belong to the initial data.
Let \(\Xi\) denote the finite set of tuple characters. The successive field requests are finite sets of algebraic numbers with
\[
 S_1(K_s)\subseteq K_0,\qquad S_2(K_s,K_0,\Xi)\subseteq K .
\] Combine requests at each stage by taking their union with the preceding arguments fixed. The prescribed rational prime set includes the two type primes, \(d\), and every prime at most \(X''\).
\end{paramdefinition}

To meet these finite requests, first fix the embedded small field, then adjoin the union of the first sets, and only then evaluate the second sets in that field and take their compositum. Every extension in these two steps is finite. The field cost for the geometric requests on the initial frames is proved in Lemma~\ref{lem:marked-original-geometric-field}.
\begin{lemma}[Initial characters and support over the smaller field]
For Definition~\ref{hyp:initial-adelic-frame} with trivial auxiliary characters,
\(\Xi=\{1,\delta,\delta'\}\). For every \(a\in\mathcal O_{K_0}\), the support test factors through contraction:
\[
 v(a)>0\quad\Longleftrightarrow\quad a\in v\cap\mathcal O_{K_0}.
\]
Thus any support selected by these tests is a union of fibres of the contraction map.
\end{lemma}
\begin{proof}
The reference contributes \(1\). At a coefficient place the
deletion rule uses \(\delta'\) over the first deletion prime
and \(\delta\) otherwise. The three distinguished tuples realize
the three listed values. Thus the final finite field request
\(S_2(K_s,K_0,\Xi)\) is determined before localization in \(K\);
adjoining it requires no new word level. For \(x\in\mathcal O_{K_0}\),
\[
 v(x)>0\quad\Longleftrightarrow\quad x\in v\cap\mathcal O_{K_0}.
\]
The positivity tests coincide at two places exactly when their
contractions coincide. Apply this to the prospective nonbase support.
\end{proof}

Begin with the untwisted saturated form of the marked adelic datum.
Its localization set is empty and its prescribed orders are
\[
                 \eta_j(v)=0\qquad(1\le j\le m,\ \text{all selected }v).
\] The two
weight comparisons used in making the choice are the lower
unitary comparison and the comparison at the residue characteristic
of the lift. The first division lattice at this step uses the zero
orders, the unit coordinate frame and the full word-column lattices.
It gives a starting row before the local arithmetic orders are
inserted. The row stage comes from the lowest discrete series at
the specified archimedean weights and the finite types on that
adelic datum. \label{cor:automorphic-realization-1}

\paragraph{Forming the adelic row on the prescribed group.}
Retain the original Hermitian frame, the finite level group and
the marked split half. Let \(P\) be the finite rational prime set in
Definition~\ref{hyp:initial-adelic-frame}. Before restricting the
support, allow the entire set of coefficient places
\[
             \{v:v\mid p\text{ for some }p\in P\},
\]
which is finite and stable under conjugation. The split-half condition is used even
when two different tuple conductors contain places at the same
rational prime. Both must use the same chosen half there. The
place of each real subfield is the one below the specified
column place, so the pinned choice gives this condition for
all the conductors at once.

Use the active tuple characters and their conductor sets in the
level datum. The distinguished tuple has character one. At the
two deletion characteristics use the prescribed split side, and
carry the original ramification test at every other place. If \(S_{\rm type}\) is the set of chosen type places and \(\bar G_v\) the matrix group given by reduction there, the quotient is identified place by place:
\[
 G_{\rm type}\simeq\prod_{v\in S_{\rm type}}\bar G_v .
\] Use the diagonal self-dual
frames when they have been supplied by the Hermitian frame.

In
the standard lattice case the coordinate frames with unit scales
are self-dual. These give the frames for the row and each of the
lifts. The finite group and character comparison then supplies
the type projectors on this product, together with the inert place
and tuple assignments. The identity map of the finite group can
serve as the preliminary label map. The Fock and lowest-type
models are to be built for these very groups and characters.
\label{cor:automorphic-realization-2}


\begin{proposition}[Adelic realization on the prescribed frame]
Use the fields and levels of Definition~\ref{hyp:initial-adelic-frame}
and the marked self-dual frames at the chosen split halves. Take the
level volumes in their common Haar measure, the regular discrete-series
parameters and their specified lowest compact types, and the definite
standard lattices with \(n_V\ge 8fr+4\) in the matched geometric model.
Then the initial row admits the prescribed real
intertwiners and theta lifts on the chosen finite types. At the distinguished tuple, if \(\theta\) is the theta series of a rational input and \(f\) its algebraic class sum, the normalized comparison is
\[
 f=\theta .
\]
The phase is one for this normalization.
\end{proposition}
\begin{proof}
The finite-place characters take values in the fixed coefficient
subfield on ideles with trivial archimedean component. The analytic construction
uses these character values. The proof uses the volume tie on this
same group to supply the core, its coefficient algebra and these
finite character values. The real models at every place use
the regular parameter of the prescribed weight and its
specified lowest compact type. The discrete series of the
determinant cover give these models. We use Harish--Chandra's
discrete-series theorem for the regular parameters and the
Casselman--Wallach theorem to extend the admissible
Harish--Chandra-module intertwiners to the smooth
moderate-growth globalizations.

For the Harish--Chandra modules \(V,W\) in this comparison and their
smooth moderate-growth globalizations \(V^\infty,W^\infty\), the
Casselman--Wallach equivalence identifies
\[
 \operatorname{Hom}_{G}(V^\infty,W^\infty)
       \simeq\operatorname{Hom}_{(\mathfrak g,K)}(V,W).
\]
Thus the extension preserves the intertwining
identities on the dense compact-finite vectors. With the
chain and the finite-action identities on these models, take
the theta map from the lowest-type Fock coefficient map.
Rapid decay of the cusp input gives the convergent integrals
for the intermediate lifts. The final projected lift uses the
good projector built on the chosen finite quotient. This is
the role of the intertwiner and lift existence assertions in
the fine datum. \cite{harish-chandra-degree,casselman-wallach}\label{cor:automorphic-realization-3}

For the geometric realization use the standard definite lattices.
The dimension inequality required there is
\[
                          n_V\ge8fr+4 .
\]
Use the same frame for the volume normalization and the algebraic
identifications. At each lift the finite comparison has one scalar for all allowed tuples. Determine it at the distinguished tuple in this common frame.

To determine the scalar relating the two forms, take the
theta series \(\theta\) of a rational input there, and let \(f\)
be its algebraic class sum. For a fixed class \(x\), the definite theta integral is a finite class sum. 

Let \(K_f\) be its finite level, \(h\) range over the double cosets, and \(\Gamma_h\) be the stabilizer in that class (with the same central quotient in the measure and the class sum). Write \(k(x,h)\) for the already normalized kernel times the input, paired at infinity. Splitting the integral into the individual level orbits gives
\[
 \theta(x)=\sum_{[h]}\frac{\operatorname{vol}(K_f)}{|\Gamma_h|}\,k(x,h)=f(x).
\]
Here each orbit has the indicated mass by the finite stabilizer formula; the second equality uses the same character, kernel and measure in the algebraic class sum. Evaluation over all classes therefore gives
\[
  \bigl({\rm ev}_x(f-\theta)\bigr)_x=0,\qquad
  \bigcap_x\ker({\rm ev}_x)=0 .
\]
Thus \(f=\theta\) with phase one. The weight scalars and geometric
indices transport this equality on the same levels, finite
group, tuple characters and selected places. The construction
keeps the intertwiners, group action and lattices together
in the normalized adelic realization. \label{cor:automorphic-realization-4}
\end{proof}

\begin{proposition}[The same multiplicity on the linked real and finite types]\label{prop:linked-type-multiplicity}
For the linked unitary groups with the same finite level and
the real and split actions in the frames of
Definition~\ref{hyp:initial-adelic-frame}, write \(m_H(\sigma)\) and \(m_{H'}(\sigma')\) for the multiplicities of matching types in the two groups. At either initial weight the trace comparison gives
\[
 m_H(\sigma)=m_{H'}(\sigma').
\]
\end{proposition}

\begin{proof}
The multiplicity assertion compares the ranks of spaces of
finite-adelic intertwiners into cuspidal forms of fixed real
discrete-series type. It compares the row group with the group
of the indicated lift. Their odd ranks agree, their real
parameters agree at every real place with adjacent gaps at
least \(100r^4\), and the finite group isomorphism identifies
the lattice frames modulo the type places. The selected type
has the no-fixed-vector condition along the prescribed proper
parabolic radicals. The representations in the comparison are
irreducible admissible and have a level-fixed vector surviving
the type projector. Denote the two groups by \(H,H'\) and the
transport of a finite representation \(\sigma\) by \(\sigma'\).
With the real models included in the multiplicities, the required identity is
\[
                         m_H(\sigma)=m_{H'}(\sigma').
\]
\label{cor:automorphic-realization-5}

\proofstep{Match the projected trace at the prescribed weight.}
At each real place use the pseudo-coefficient normalized to
select the specified discrete series for the trace. At the
finite places use the type tests and an auxiliary split place
outside their support. The constant terms along proper
rational parabolics vanish for these tests. The prescribed
split test allows only regular semisimple elements on the
diagonal kernel, and the nonfield terms vanish in the class
integrals. The trace expansion is then a finite sum over
the regular polynomial classes with field centralizer. If \(T\) is the common torus identifying the two centralizers, use a single measure \(dt\) on it. For the corresponding identifications \(i,i'\) the measures on the two sides are
\[
 d\mu_{\mathrm{cent}}=i_*dt,\qquad d\mu'_{\mathrm{cent}}=i'_*dt .
\]
Transport the finite Haar measures with the frame isometries.
This identifies the finite orbital factors, including the
normalizations of the lattice stabilizers.
\label{cor:automorphic-realization-6}

At a real place where the torus is compact, write its distinct
eigenvalues as \(z_1,\ldots,z_r\). The real conjugacy classes
are the choices of which eigenlines have positive sign, with
the number of signs determined by the real form. Write \(W_K\subset S_r\) for the stabilizer of the sign pattern.
The classes provide the disjoint decomposition
\[
                         S_r=\coprod_y y^{-1}W_K .
\] For ordered integral
Harish--Chandra parameters \(\lambda_j\), put
\(a_j=\lambda_j+(r-1)/2\). The discrete-series character
on the class indexed by \(y\) is the alternating sum over
\(W_K\) on the reordered eigenvalues divided by their
Vandermonde. Its overall factor \((-1)^{pq}\) is one
because \(p+q=r\) is odd.

Moving each denominator to the
same ordering contributes the sign of \(y\), so the sum of the
quotients is
\[
   \frac{\det(z_i^{a_j})_{1\leq i,j\leq r}}
        {\det(z_i^{r-j})_{1\leq i,j\leq r}}.
\]
Indeed, expanding the compact alternating sums on the cosets gives
\[
 \sum_y\sum_{w\in y^{-1}W_K}\operatorname{sgn}(w)
         \prod_{i=1}^r z_i^{a_{w(i)}}
 =\sum_{w\in S_r}\operatorname{sgn}(w)\prod_{i=1}^r z_i^{a_{w(i)}}
 =\det(z_i^{a_j}).
\] These are the Weyl character
coordinates of highest weight \(a_j-r+j\). In the orbital
integral the formula is conjugated and divided by the same
torus volume. At a real place with noncompact centralizer
the orbital integrals of the chosen pseudo-coefficients
vanish. This computes the two real stable factors in terms
of the same eigenvalues and weights even when the signatures
are different. Thus, writing \(h'\) for the transport of the finite test \(h\), the equality assembled from the orbital integrals is
\[
 \operatorname{tr}_{\rm cusp}(h)=\operatorname{tr}'_{\rm cusp}(h').
\] \label{cor:automorphic-realization-7}

Choose a common compact open subgroup on the finite side
for the tests, and transport it to \(H'\). The fixed-vector
spaces in the corresponding discrete-series multiplicity
spaces have finite dimension by cuspidal spectral finiteness.
Make one finite list \(\sigma_1,\ldots,\sigma_d\) of distinct
irreducible admissible representations covering both spectra
seen by the tests and include the given \(\sigma\). The
spectral trace and the preceding equality give, for every
allowed finite test \(h\),
\[
   \sum_{i=1}^d a_i\operatorname{tr}\sigma_i(h)=0,
   \qquad a_i=m_H(\sigma_i)-m_{H'}(\sigma_i').
\]
Let \(e_v\) be the local convolution idempotent for the chosen
type at a type place, the normalized Iwahori indicator at
the auxiliary split place \(u\), and the normalized maximal
compact indicator at the remaining places. Their pure tensor
acts on \(\sigma_i\) by a finite-rank idempotent \(P_i\).

The resulting projector is nonzero on the stipulated \(\sigma\). For a finite-adelic
element \(g=(g_v)_v\), set
\[
 b_v(x)=e_v(g_v^{-1}x),\qquad f_v=e_v*b_v.
\]
Since \(b_v*e_v=b_v\), the test at a type place satisfies
\(e_v*f_v*e_v=f_v\). Away from the type places and \(u\)
it is bi-invariant under the maximal compact. The family
agrees with the unit Hecke family almost everywhere. At \(u\)
it is left Iwahori-invariant. Replacing only \(f_u\) by
\(1_{ItI}\), for each \(t\) in the strict dominant torus
cone, preserves all the conditions for the trace comparison.
The original pure tensor acts by \(P_i\sigma_i(g)P_i\).
\label{cor:automorphic-realization-8}

For a local isomorphism class \(\lambda\) at \(u\), collect all
the factors away from \(u\) in the coefficient
\[
 b_\lambda=\sum_{i:\,\sigma_{i,u}\simeq\lambda}
             a_i\prod_{v\ne u}{\rm tr}\,\sigma_{i,v}(f_v).
\]
Classes without Iwahori invariants contribute zero. For \(t\) in
the strict dominant cone the trace relation therefore reads
\[
 0=\sum_\lambda b_\lambda{\rm tr}\,\lambda(1_{ItI})
   ={\rm vol}(I)\delta_B(t)^{-1}
       \sum_\lambda b_\lambda\sum_{\xi\in J(\lambda)}\xi(t).
\]
Here \(J(\lambda)\) lists the unnormalized Jacquet exponents with
multiplicities (normalized exponents give the common factor
\({\rm vol}(I)\delta_B(t)^{-1/2}\)). The cited exponent separation
on the local classes in this comparison forces \(b_\lambda=0\).
Reinsert \(f_u\). Cyclicity of finite-rank trace gives
\({\rm tr}(P_i\sigma_i(g)P_i)={\rm tr}(\sigma_i(g)P_i)\),
so \(\sum_i a_i{\rm tr}(\sigma_i(g)P_i)=0\) for all \(g\).
\cite{zelevinsky-induction}\label{cor:automorphic-realization-9}
The commutant of \(\sigma_i\) is \(\mathbb C\): an endomorphism
has an eigenvector on a nonzero finite-dimensional fixed space,
and that vector generates the irreducible module.

For a prescribed
index \(k\), apply Jacobson density to bases of the finite spaces
\({\rm im}\,P_i\) in the distinct modules. There is a group
algebra element \(a\) such that
\[
 \sigma_i(a)|_{{\rm im}\,P_i}=\delta_{ik}\,{\rm id},\qquad
 0=\sum_i a_i{\rm tr}(\sigma_i(a)P_i)=a_k\operatorname{rk}P_k .
\]
Since \(P_k\ne0\) at the selected type, its multiplicities agree.
The identical finite and real actions transport this equality at
both initial weights. \label{cor:automorphic-realization-10}
\end{proof}

\begin{lemma}[Rational separation by common row tests]\label{lem:rational-row-separation}
Use the prescribed tests of Definition~\ref{hyp:fitted-pointwise-nonvanishing} and their nonvanishing in Lemma~\ref{lem:fitted-completed-nonvanishing}. For the rational row maps \(B_t:A_t\to E_{\mathbb C}\) at the allowed tuples, one finite coefficient request supplies complex-linear functionals \(\lambda_a\) with
\[
 \bigcap_a\ker\lambda_a=0,\qquad\lambda_a\bigl(B_t(A_t)\bigr)\subseteq K\quad\text{for every }a,t,
\]
where \(K\) denotes its specified embedded copy.
The fixed-field and support comparisons used at rank zero
preserve these maps under restriction of the coefficient
places.
\end{lemma}

\begin{proof}
Write \(K\) for the large coefficient field with its fixed
embedding, \(A_t\) for a rational row of tuple \(t\), and
\(B_t(x)=\Theta_t(j_t x)\) for its image in the
common complex row space \(E_\mathbb C\). We need complex-linear
functionals \(\psi_u\) with zero common kernel on
\(E_\mathbb C\) and every \(\psi_u(B_t(x))\) in the
embedded \(K\). This is the field supplied by the finite
coefficient request on the actual tuple set.

For each allowed finite test type \(\alpha\), let
\(C_\alpha\) be the column space and \(c_{\alpha,u}\) its fitted
spanning columns. The common theta integral identifies the row
tests with the Bessel pairings for the specified real vector and
covector. For a row \(y\) and \(c=\sum_u a_uc_{\alpha,u}\), convergence
and linearity give
\[
 {\cal B}_\alpha(c,y)=\sum_u a_u\,{\cal B}_\alpha(c_{\alpha,u},y).
\]
If all the column tests vanish on \(y\), linearity kills every pairing on every \(C_\alpha\). On the given base multiplicity space Proposition~\ref{prop:fitted-grid-base-separation} chooses one real test and a basis of the dual from these allowed tests. Evaluation against this basis gives \(y=0\).
A finite family suffices: if the current common kernel has positive
dimension, choose one further test nonzero there. Each choice
strictly lowers that dimension. \label{cor:automorphic-realization-11}

For the field assertion first apply the rational good
projector to the input \(x\), giving \(x'\). If \(P\) denotes the good projector, equivariance and the good target give
\[
 B_t(x')=B_t(Px)=P B_t(x)=B_t(x).
\] Write \(C_t=\prod_{i,j}c_{ij,t}\) for
the product of the nonzero fitted square constants and
\(0\) for the distinguished tuple. The finite Fourier and
class sum reads \(v_{t,u}(x')\) in the embedded coefficient
field. Write \(\iota\) for the coefficient embedding, \(\zeta_t\)
for the coefficient-field root of unity, and \(\kappa\) for
the single complex scalar for the chosen test family.
On the good input the formula reads
\[
   \iota(\zeta_t) C_t\,\psi_u(B_t(x))
     = \kappa v_{t,u}(x').
\]
The bottom toric comparison in its nonvanishing integer range gives
\(C_t/C_0=\prod_{i,j}(c_{ij,t}/c_{ij,0})\in\iota(K)^\times\).
If \(\kappa\ne0\), set \(\psi'_u=(C_0/\kappa)\psi_u\). Then
\[
 \psi'_u(B_t x)=\frac{v_{t,u}(x')}{\iota(\zeta_t)(C_t/C_0)}
             \in\iota(K),\qquad \bigcap_u\ker\psi'_u=0.
\]

If \(\kappa=0\), the displayed Fourier identity gives
\(\psi_u(B_tx)=0\), already in the field. Make the request on the whole
tuple family so that these comparisons use the same field.
For a complex row with conjugate scalar action, precompose
with its conjugate-linear identification and conjugate the
output. The resulting functionals are complex-linear, and
their values are in the conjugate coefficient embedding.
If \(X_1,X_2\) are the cutoffs for the test and field assertions, impose them simultaneously by taking
\[
 x\ge \max\{X_1,X_2\}
\]
and increasing the scale further whenever the application requires it. The Hida size setting is supplied
on a bare state with the same group, so the tests pull back
to the original state along its complex row isomorphism.
Their values and field are preserved. This is the second
assertion at the residue-characteristic weights. It is used
jointly with multiplicity on the restricted datum.
\label{cor:automorphic-realization-12}

With respect to a \(K\)-basis \((a_i)\) of \(A_t\), the normalized
tests have coefficient matrix
\[
               M_t=\big(\iota^{-1}\psi_u(B_ta_i)\big)_{u,i}.
\]
The finite test map \(\Psi=(\psi_u)_u\) is injective, hence
\[
 \ker\Theta_t=\ker(\Psi\Theta_t)=j_t(\mathbb C\otimes_K\ker M_t).
\]
The last equality follows by row reduction over \(K\). Equivariance
of \(\Theta_t\) makes this descended kernel stable under the level group.

To pass from the full allowed place set to the actual nonbase
support, restrict to the places detected by the index tests,
together with their conjugates. The other places have trivial
auxiliary character. Restrict the tuple-indexed data along
the inclusion of the smaller set, then update the weight
scalars on the restrictions of the same group and geometric
maps. The two index tests commute with that restriction.
They depend on the place readings retained on the restriction
and are unchanged by the accompanying scalar update.

Denote restriction to the detected places and their conjugates by \(R\), and the support recovered from the two tests by \(S\). The preceding compatibility gives
\[
 S(RD)=S(D)
\]
for the datum \(D\) under consideration. Thus a second restriction discards no further places. The residue-characteristic weight comparison is
applied \emph{after} this step. Its hypotheses include equality
with the recovered support, the definite weight condition,
the core and group conditions and membership of all residue
characteristics in the allowed prime set. It returns the
comparison used to construct the first row-lift family.
\label{cor:automorphic-realization-13}

The fixed-field part of the support argument uses an earlier fine
adelic datum on the same pre-level frame. For the fixed small field
one requests a finite list of algebraic coefficients in \(K_0\)
and a finite list of residue characteristics. The uniformity here
is over the generator choices and auxiliary numerical parameters
at the prescribed levels. Let \(p\) be a type prime or the prime
assigned to a sequence. The norm relation gives the same bound at every embedding:
\[
              \pi\bar\pi=p\quad\Longrightarrow\quad
               |\sigma(\pi)|=\sqrt p\qquad(\sigma:E\hookrightarrow\mathbb C).
\] In the
Minkowski embedding the possible \(\pi\)'s lie in the finite set
\[
 \mathcal O_E\cap\prod_{\sigma:E\hookrightarrow\mathbb C}
                          \{z:|z|\le\sqrt p\}.
\]

For each occurring tuple of generators at the finitely many word
and type primes, choose one fine reference realization with its
lower-unitary comparison. Its two zero-content index tests compare
full lattices in a common space. For each pair choose \(N>0\) clearing both lattice inclusions.
Localizing gives
\[
 NL\subseteq L',\quad NL'\subseteq L
        \quad\Longrightarrow\quad L_v=L'_v\quad(v\nmid N).
\] Request the finitely many resulting
residue primes and their conjugates, and an embedded rational basis
of the reference coefficient field. The comparison at fixed generators
preserves the lower-unitary maps and these zero-content tests when
the other numerical parameters vary. Thus the union over reference
tuples is finite before \(K_0\) is chosen.

The chosen reference coefficient field embeds into \(K_0\).
Transport its support with zero added content to \(K\).
In the enlarged datum this inherited support is supplemented by
places with nonzero numerical content at the prescribed scale and
their conjugates. The latter residue characteristics are among the
type primes, the auxiliary prime and primes at most \(X''\), so they are
already allowed. The inherited residue characteristics were included
with the finite list. For either part of the support let \(S\) denote its set of places. Saturation over \(K_0\) means precisely
\[
 w\cap K_0=w'\cap K_0\quad\Longrightarrow\quad(w\in S\ \Longleftrightarrow\ w'\in S).
\]

For the inherited part, two primes over the same \(K_0\)-prime
restrict to the same prime of the earlier coefficient field, and
the comparison identifies membership with the test at that prime.
For the numerical content part the test on a lift is a function
of the rational residue characteristic. It is invariant under
the same contraction. This gives the saturation on the
actual datum used to initialize the zero row orders. The fixed-field
comparison uses an enlargement of the fine datum carrying the
lower-unitary maps, with the same finite group and rational readings.
\label{cor:automorphic-realization-14}
\end{proof}

\begin{proposition}[Ray splitting and Galois representations on the initial row]\label{prop:initial-row-arithmetic-readings}
For the fixed initial row and its marked finite types, the ray
splitting test identifies the pinned ray field over the marked
quadratic base. The local type and the Galois representation attached to the same
row agree on the prescribed inertia data. For a continuous de Rham representation of rank \(a\) at a coefficient place, let \(h_{\tau,i}\) be the multiplicity of jump \(i\) at an embedding \(\tau\) inducing that place. Then
\[
 \sum_i h_{\tau,i}=a\qquad\text{for each such }\tau.
\]
\end{proposition}

\begin{proof}
\proofstep{Identify the two principal prime ideals and the pinned ray field.}
Keep the split places, the old row representation for each
finite type \(\sigma\), and the two central actions on the
ambient and multiplicity spaces. Let \(E\) be the auxiliary field,
and let \(g_0,g_1\) be the two primes used for the finite
types. Proposition~\ref{prop:two-type-primes} supplies generators
\(\pi_k\) in the integers of its distinguished quadratic subfield.
With \(Q_S=\operatorname{lcm}(2^{e+2},f_\psi)\) as in
\eqref{eq:host-character-compositum}, they satisfy
\[
   \pi_k\overline{\pi_k}=g_k,
   \qquad \pi_k\equiv 1\pmod{Q_Sqd}.
\]

Under the fixed embedding \(j_E\) of \(E\), both
\(j_E(\pi_k)\) and \((j_E(\pi_k)-1)/(Q_Sqd)\)
are integral over \(\mathbb Z\). Extend complex conjugation on \(E\) to an automorphism of the algebraic closure. The
product of \(j_E(\pi_k)\) and its image under this
automorphism is \(g_k\). This is the ray-prime test used
for splitting in the ray field of the marked base. To
see the splitting, put \(c=Q_Sqd\).

In the imaginary
quadratic base the two conjugate factors are integral
and each is one modulo \(c\). Thus \(g_k\equiv1\pmod c\)
and \(g_k\nmid c\). For the pinned ray field \(M\), the
ray-field ramification comparison gives
\[
              g\mid\operatorname{disc}(M/\mathbb Q)
                  \ \Longrightarrow\ g\mid c
              \qquad(g\text{ a rational prime}).
\]
The quadratic base has discriminant \(-d\), and \(d\)
divides \(c\). At \(g\nmid c\) the base is unramified and
the ray extension is unramified at every prime above \(g\)
by reciprocity for the modulus \((c)\). Multiplicativity of
ramification indices in the tower gives the comparison for
the absolute discriminant. Apply it with \(g=g_k\).
The two principal ideals have
norm \(g_k\) and multiply to \((g_k)\). They are distinct
because their product is \((g_k)\) and \(g_k\) is unramified.
Each has trivial ray class
and splits in the ray class field by reciprocity.
\label{cor:automorphic-realization-15}

Conversely let \(\mathfrak p\) be a degree-one prime away from \(c\)
splitting in the ray field. For the modulus \((c)\), reciprocity gives
\[
 {\rm Frob}_{\mathfrak p}=1
 \ \Longleftrightarrow\ [\mathfrak p]_{(c)}=1
 \ \Longleftrightarrow\
 \mathfrak p=(\alpha),\quad \alpha\equiv1\bmod c.
\]
In the imaginary quadratic base the norm of \(\alpha\) is the
positive rational prime below \(\mathfrak p\). These are the
generator and congruence conditions under the fixed embedding of \(E\). The ray
field, its marked conjugation and the old finite type actions are
transported with that embedding. \label{cor:automorphic-realization-16}
Indeed, if every marked ray-test prime splits completely in a finite
extension of the base, so does every such degree-one prime.

Degree-one primes suffice here: the primes of the imaginary quadratic
base of higher residue degree have Dirichlet density zero. Chebotarev's
split-prime inclusion criterion puts the extension inside the ray
field, which is therefore maximal for this test. For the embedded field \(K\) the corresponding ring is recovered from integrality alone:
\[
 \mathcal O_K=\{z\in K:f(z)=0\text{ for some monic } f\in\mathbb Z[X]\}.
\] Transport also preserves the integral row
exponents at the type primes and the valuation zero in the
lower-left comparison. \label{cor:automorphic-realization-17}
\proofstep{Use the fixed-vector comparison on the actual split local component.}
At a marked place over \(g_k\), let \(K_1\) be the first
principal congruence subgroup and \(\pi_v\) the split component
of the old row of type \(\sigma\). The split comparison embeds
\(gK_1g^{-1}\) into both bottom levels, where both level
characters are trivial. A nonzero bottom-level vector therefore
gives, by local factorization,
\[
 \pi_v^{gK_1g^{-1}}\ne0,\qquad
       \pi_v(g^{-1})\pi_v^{gK_1g^{-1}}=\pi_v^{K_1}.
\]
This uses the two level maps on the actual bottom frame.
It supplies the same depth condition for all finite types and
marked places in that row. At these places the completion
is \(\mathbb Q_{g_k}\). Depth-zero compatibility of local Langlands gives trivial wild inertia for the Weil--Deligne representation \(r_v\):
\[
 r_v(P_v)=\{1\},\qquad P_v\subset I_v\ \text{the wild inertia subgroup}.
\] Keep the local reciprocity parameter
at every split place and the Satake multiset at the spherical
inert places. \label{cor:automorphic-realization-18}

Use a tame inertia generator coupled to the fixed residue
generator at each type place. The Green labelling gives the
series idempotents in these coordinates. A local unit
representing this inertia generator can be compared with
its rational residue representative by reciprocity. Testing
the character on that representative gives the Teichm\"uller value
of the inverse finite character in the allowed orientation.

If
\(\alpha_1,\ldots,\alpha_r\) are the tame inertia eigenvalues
listed with multiplicity and \([\cdot]\) denotes Teichm\"uller lift,
the residual support polynomial is
\[
 P_\sigma(T)=\prod_{h=1}^r(T-\overline{\alpha_h}),\qquad
 [\overline{\alpha_h}]=\alpha_h .
\]
The reductions here are at the marked residue valuation, with the
chosen tame eigenvalues in the prime-to-residue-characteristic
roots of unity. Its degree is \(r\) because the parameter is tame. Transport the
place on the closure and its generator by the marked conjugation;
the residue generator law transports the polynomial with them, on the same Galois place frame and with the two original finite type actions.
\label{cor:automorphic-realization-19}
On this parameter family the character and cuspidal
comparison realizes the local parameters on the automorphic
side of the bare level.

This comparison uses local reciprocity, the
finite-dimensional fixed-level automorphic comparisons and
the character lifting at the selected ray field. In the ray
stage replace the auxiliary parameter records by those of
the split ray class field. For every chosen residue label \(u\), Teichm\"uller lifting followed by reduction at the specified ray place returns that very label:
\[
 [u]\longmapsto\overline{[u]}=u .
\]
Distinct roots in the reduction field stay distinct because
their orders are prime to the residue characteristic.
The ray modification preserves the two central actions,
the series idempotents and the marked places. Throughout
these operations the localization index set is empty and
the prescribed order data are carried along. The arithmetic
enrichment is applied before the later specialization of
the orders. \label{cor:automorphic-realization-20}

Let \(\iota_G\) identify the old and new finite groups using the
indexed type places and their matrix coordinates. The series
idempotents have matching coefficients. For either row comparison
\(J\) and central element \(z=\sum_g z_g[g]\), equivariance gives
\[
 J\rho(z)=\sum_g z_g J\rho(g)
         =\sum_g z_g\rho'(\iota_G g)J=\rho'(\iota_G z)J .
\]
The old row fits and geometric and finite isotypic readings
therefore use the same central projectors. The underlying host,
rank, empty index set, algebraic closure with its host embedding,
prime valuations and type-place indices are fixed. The places
above the real subfield still project to the specified list.
\label{cor:automorphic-realization-21}
The algebraic-closure frame extends the chosen embedding
of the host.

On the same Weil--Deligne realization impose
continuity of each \(\ell\)-adic matrix representation in
the de Rham condition. Write \(f:A\to B\) for either coefficient map. Continuity of the transported representation follows from entrywise composition:
\[
 G\xrightarrow{\ \rho\ }\mathrm{Mat}_a(A)\xrightarrow{(a_{ij})\mapsto(f(a_{ij}))}\mathrm{Mat}_a(B).
\] Thus the old
continuous local lifts still satisfy the strengthened
condition. The unramified and monodromy tests use the
original de Rham condition. \proofstep{Read the dimension of the de Rham filtration at the chosen embedding.}
Apply the Hodge--Tate
comparison to a continuous rank-\(a\) representation
that satisfies the local de Rham condition at \(w\).
For an embedding \(\tau\) into the \(\ell\)-adic
coefficient closure inducing \(w\), extend it to the
algebraic closure with the prescribed valuation ring.
One obtains the extension by conjugating a field
embedding over the host, since the two valuation
rings restrict to the same ring on the host. The continuous
representation descends to a finite extension of \(\mathbb Q_\ell\)
containing the local field. In the labelled de Rham space \(D_\tau\) this gives
\[
 \sum_i\dim\operatorname{gr}^i D_\tau
 =\sum_i(\dim\operatorname{Fil}^i D_\tau-
               \dim\operatorname{Fil}^{i+1}D_\tau)=\dim D_\tau=a .
\]
The jumps are counted with multiplicity at each \(\tau\) inducing
\(w\), equivalently \(x\in\mathfrak m_w\) if and only if
\(|\tau(x)|_\ell<1\) on host integers. \label{cor:automorphic-realization-22}
\end{proof}

\begin{proposition}[Isotypic span in the rational row]\label{prop:rational-row-isotypic-span}
At a lift \(1\le j\le m\), write \(V_{j,\lambda}\) for the space with prescribed label \(\lambda\). Its internal isotypic decomposition is
\[
 V_{j,\lambda}=\bigoplus_{\sigma\ {\rm occurring},\ {\rm label}(\sigma)=\lambda}I_{j,\sigma}.
\] In the rational
model, the allowed split Hecke operators preserve the
indicated coefficient lattices and furnish eigenvalues
on the full good simultaneous eigenspaces.
\end{proposition}

\begin{proof}
The lower-unitary multiplicity equality applies on the
same frame at every state with those pins. To see which
finite types it compares, use the good projector for the
actual two type primes. Its coefficients at \(g\) and
\(g^{-1}\) agree by the finite type inversion identity.
Inversion takes an irreducible character to the character
of its dual. In the expansion of the good projector in primitive central idempotents, inversion interchanges the labels \(\sigma\) and \(\sigma^\vee\). Independence of irreducible characters now gives
\[
 \sigma\in\operatorname{supp}(e_{\rm good})\quad\Longleftrightarrow\quad\sigma^\vee\in\operatorname{supp}(e_{\rm good}).
\] The
character comparison descends through the injective
coefficient embedding. This is the duality stability
condition in the multiplicity comparison.
\label{cor:automorphic-realization-23}

\proofstep{Compare the dimension of each label with the sum of its cuspidal isotypic dimensions.}
Let \(I_{j,\sigma}\) be the row's complex isotypic space
at a lift \(1\le j\le m\), and let \(n(\sigma)\) be
the dimension of the good finite-type space of \(\sigma\)
on the compact base row. Transport through the lift
pairing and the tensor-product identification reads
\(I_{j,\sigma}\) as the intersection of the global
\(\sigma\)-part with the fixed level, level character and
good projector of the lift. In the finite cuspidal
decomposition, the multiplicity equality identifies its
dimension with the compact counterpart. The compact
comparison identifies that dimension with the base count
\(n(\sigma)\). Hence
\[
                      \dim I_{j,\sigma}=n(\sigma).
\]
Only the fixed locus, the level character and the same
good projector enter this dimension comparison.
\label{cor:automorphic-realization-24}

The rational reading available at this point consists of
the rational coordinate isomorphisms for the sharp and
base forms, the two complex readings, their metrics at
the other embeddings, and the lift maps from rational
coordinates to the complex rows. No order comparison
is needed for these maps. The bottom Hecke action assigns
each \(\sigma\) a constituent label \(c(\sigma)\).
The theta adjoint identity on the rational generators
extends to arbitrary test vectors.

For a theta map
\(\Theta\), project a source vector to
\((\ker\Theta)^\perp\). Both sides of the adjoint identity vanish
on the discarded kernel; linearity extends it from the generators.
On a labelled image the adjoint identity gives
\[
 \Theta\Theta^*y=\lambda y,\qquad
 \lambda\|y\|^2=\|\Theta^*y\|^2>0\quad(y\ne0).
\] For the common Hecke label \(c\)
and its lift space \(V_{j,c}\) the dimension comparison gives
\[
 \bigoplus_{\sigma:\,c(\sigma)=c}I_{j,\sigma}\subset V_{j,c},
 \qquad \dim V_{j,c}=\sum_{\sigma:\,c(\sigma)=c}n(\sigma).
\]
Distinct cuspidal projectors give the direct sum, and
\(\dim I_{j,\sigma}=n(\sigma)\). Thus the inclusion is an equality
at every lift. \label{cor:automorphic-realization-25}

\proofstep{Read the operators and eigenvalues on the same integral coefficient model.}
At the arithmetic level these decompositions give the
span test for the good simultaneous Hecke eigenspaces.
The rational forms come from the full geometric lattices
and are read through \(K_s\hookrightarrow K_0\).
At each type prime \(g_k\), the inverse of the
chosen primitive \(g_k\)-th root of unity defines
a nontrivial additive character of \(\mathbb F_{g_k}\).
Compose it with the marked linear combination of the
superdiagonal entries of the upper unipotent matrices.

Write \(L(u)\) for the marked linear combination of the superdiagonal entries and \(\zeta\) for the chosen primitive root. The level character is the inverse unipotent character because
\[
 (\zeta^{-1})^{L(u)}=(\zeta^{L(u)})^{-1}.
\] The split
Hecke transform and its rational realization give the operators
on the same rational forms, with the index-zero operator
equal to the identity. The integral comparison gives
stability for the indicated coefficient lattices.
Use the isotypic span and the good finite-type condition
to choose their eigenvalue arrays. The finite set of
excluded places contains the indicated ramified and level
places. These rational operators act on the same integral lattices
whose characters and Hecke span were just determined. \label{cor:automorphic-realization-26}
\end{proof}

\paragraph{Theta series and marked points on the rational row.}
Proposition~\ref{prop:rational-row-isotypic-span} supplies the
integral operators and their simultaneous eigenvalues.
The cuspidal parameter comparison uses the same good lower
types that occur in the rational row. At split primes outside
its finite exclusions, the rational arrays read the elementary
symmetric polynomials in the normalized Satake parameters.
The index-zero identity fixes the normalization of the
Hecke operators. The matching of Frobenius polynomials on
these split primes, with semisimplicity and Chebotarev, gives
the coefficient-field condition for the cuspidal polynomials
used in the local multiplicity models. The comparison also
uses the fixed lower-left valuation and the reciprocal
character twist in passing back to the original row.
\label{cor:automorphic-realization-27}


\begin{proposition}[A common grid for the coefficient and geometric models]
At the prescribed grid size, depth and finite list of levels, the cuspidal forms, toric coefficients, and marked-point evaluations can be chosen with common grid characters, excluded places and signed indices. If \(J\) identifies the original row with its twisted copy in the identity multiplicity frame, every Hecke operator \(T\) outside the common exclusions satisfies
\[
 T^{\rm new}J=JT^{\rm old}.
\]
Use the multiplicity comparison of Proposition~\ref{prop:linked-type-multiplicity}
and the coefficient-ring and marked-point identifications for this
grid. \label{cor:automorphic-realization-28}
\end{proposition}
\begin{proof}
Take a grid on this row at the prescribed size and depth.
The cusp multiplicity construction and the class-group
coefficient-ring construction each specify a finite
coefficient extension at every indicated residue prime and
lift. Use their compositum. Choose one finite set \(S_x\)
of places of \(E\) containing the exclusions of the rational
Hecke action and the allowed conductors of the grid.
The twisted coefficient rings can then be chosen with
excluded place sets containing \(S_x\). The two coefficient
extensions compare in these very rings. In particular the
coefficient matrix of the cuspidal row and the finite-order
toric character can both be read there. This common choice
is needed because the toric construction compares at the
same grid entries and away from the same exclusions.
\label{cor:automorphic-realization-29}

For the cyclotomic coefficient valuation ring \(O_v\), the
integrality test in the signed comparison is
\[
 a/b\in O_v\quad\Longleftrightarrow\quad v(a)\ge v(b)
       \qquad(a,b\in O_v,\ b\ne0).
\]

The toric-character construction takes the
grid characters with their Teichm\"uller lifts and
their values in the fixed field. It uses the lower-left
unit and the inclusion of their conductor in the excluded
set. The toric ratio comparison identifies the coefficient
reading of these characters. Apply the multiplicity
comparison both at the trivial character and at the
selected toric characters, using the same full
coefficient-module basis. The settings obtained this way
carry the rational Hecke reading and the complete grid.
\label{cor:automorphic-realization-30}

Let \(J\) identify the old row with its twisted copy in the
identity multiplicity frame. The exclusion sets at corresponding
places agree, so for an allowed Hecke operator \(H\) one has
\[
 H^{\rm new}J=JH^{\rm old},\qquad
 H^{\rm old}v=\lambda v\ \Longrightarrow\
             H^{\rm new}(Jv)=\lambda Jv .
\]
The excluded residue characteristics obey their previous bound.
The level identification transports the unipotent character and
the grid with the coefficient forms. The series comparison on
these same forms supplies the designated lift support and the
signed level condition. \label{cor:automorphic-realization-31}

For the geometric step from the series, the marked
points detect the coefficient algebra in the framed
space. The signed comparison descends the support
condition to the geometric model. Let \(k_H(x)\) count the allowed kernels of a split Hecke correspondence \(H\) at a framed point \(x\). Outside the excluded characteristics the local level gives a single count: 
\[
 k_H(x)=k_H^{\rm level}=k_H(x')\qquad(x,x'\text{ in the prescribed frame}).
\]
 Its quotient read over the
complete field is integral in the chosen geometric
coordinates. The integral correspondence and the series support outside the selected range give the marked moduli-point model. Transport of level preserves the grid characters, their field values and their roots. \label{cor:automorphic-realization-32}

In particular the character assignment to coefficient places uses an index map \(v\mapsto i(v)\) satisfying
\[
 i(\bar v)=i(v)\qquad\text{at each assigned place }v .
\] Apply reciprocity to its finite-order Hecke
character, valued in the roots of unity of the specified
prime-power order. Let \(C_v\) be the set of conductor places of the character chosen
at \(v\), and \(B_{i(v)}\) the exclusions for its index. The two
support conditions are
\[
 C_v\cap B_{i(v)}=\varnothing,\qquad
 C_v=\varnothing\quad(v\text{ outside the twist support}).
\] Use the
integral ring-class descent at the two readings to put
the twist into the coefficient model. The rank test
for each reading is the level comparison at its actual
signed place. Equality of the characters gives equality
of their ramification supports, so the two conjugate
choices agree after the signed transport. The construction
at a fixed finite set of levels makes these choices for
the whole list of prescribed indices. Its spectral
specialization must keep the chosen twist labels while
inserting the actual row orders.
\label{cor:automorphic-realization-33}

\end{proof}

\begin{proposition}[Localization of the integral Hecke algebra]\label{prop:marked-hecke-localization}
At a prescribed coefficient place, tuple, and lift, a finite
extension of the local coefficient field splits the occurring
full Hecke eigensystems and contains any prescribed finite list
of scalar values. Over its integers \(O\), the Hecke order \(\mathbb T\) decomposes into residual-character blocks. If \(V_\sigma\) is a full eigenspace with character \(\lambda_\sigma\), the summand with residual character \(\bar\lambda\) is described by
\[
 \mathbb T=\prod_{\bar\lambda}\mathbb T_{\bar\lambda},\qquad V_{\bar\lambda}=\bigoplus_{\sigma:\,\overline{\lambda_\sigma}=\bar\lambda} V_\sigma.
\] The construction preserves the original
integral action under restriction to the block.
\end{proposition}

\begin{proof}
\proofstep{Split the commuting spherical operators over one finite extension.}
Let \(v\) be a marked coefficient place of residue characteristic
\(\ell\), and keep its comparison \(\iota_v\colon
\overline{\mathbb Q}_\ell\simeq\mathbb C\). The strict-henselian
character choice fixes the tuple, level and conductor on the definite
multiplicity spaces. For each lift \(j\), use the existing rational
model of the corresponding space and its integral lattice. At the
allowed split primes outside the conductor, \(\ell\) and the other
exceptional places, the normalized spherical double cosets act by
rational operators preserving that lattice. On each complex isotypic summand \(V_\sigma\) of the definite decomposition, the hyperspecial fixed line of the local spherical representation has a Hecke character \(\lambda_\sigma\). The spherical operator \(T\) acts on each copy through this same character, so
\[
 T|_{V_\sigma}=\lambda_\sigma(T)\,\mathrm{id}_{V_\sigma}.
\] Equivalently, the convolution adjoint is obtained
by inversion and conjugation, so the commuting spherical operators
are simultaneously diagonalizable for the invariant pairing.

A basis for the rational Hecke algebra involves only finitely many
of the allowed double cosets. Let \(T_1,\ldots,T_s\) be generators
so obtained, and choose \(F_{j,v}/\mathbb Q_\ell\) finite containing
the marked coefficient field and their eigenvalues. The simultaneous
eigenspace for a character \(\lambda\) is
\[
                 \bigcap_{a=1}^s\ker(T_a-\lambda(T_a)I) .
\]
The equations are defined over \(F_{j,v}\), so they split the entire
action there. Take the compositum
with the fields for the local parameter comparisons and the finite
character values. Adjoin the requested complex coefficients in the specified embedding.

For a finite set \(S_v\subset\mathbb C\) this amounts to the extension
\[
 F_{j,v}\ \subseteq\ F_{j,v}(\iota_v^{-1}(S_v)).
\] All local factors over \((j,v)\)
thus have the required coefficient inclusion with the original
embedding. The chosen finite character takes values in units, and
its auxiliary factor is trivial on inertia at the conductor primes
where ramification of the marked factor is required. The product
retains those ramified primes. Its conductor is away from \(\ell\)
and the excluded level primes by the character choice.
\label{cor:automorphic-realization-34}

\proofstep{Identify maximal ideals with reductions of full characters.}
Put \(O=\mathcal O_{F_{j,v}}\). On the base-changed lattice the integral
algebra \(A\) acting faithfully is finite over \(O\), as a submodule
of the integral endomorphism ring. If the space is zero, the index
list is empty. Otherwise
\[
 A\otimes_O F_{j,v}=\prod_{\alpha}F_{j,v},
 \qquad A\subset\prod_{\alpha} O ,
\]
where \(\alpha\) runs through the occurring full characters.
Each eigenvalue belongs to \(O\) by the monic characteristic
polynomial of its integral operator. Write \(k=O/\mathfrak m_O\).
Since \(\prod_\alpha O\) is finite over \(A\), lying over gives
\[
 \operatorname{Max}A
   =\{\ker(A\xrightarrow{\alpha}O\longrightarrow k):\alpha\}.
\]
Two residual characters with this kernel induce the same map
to \(k\), since both fix \(k\) on diagonal scalars. Completeness
lifts the residual factors to orthogonal idempotents \(e_{\mathfrak m}\).
Thus for the original lattice \(L\),
\[
 (e_{\mathfrak m}L)\otimes_O F_{j,v}
   =\bigoplus_{\ker\bar\alpha=\mathfrak m} (L\otimes_O F_{j,v})[\alpha].
\]
The decomposition intertwines every integral Hecke operator.

\proofstep{Recover the complex isotypic summands from the same character.}
Write \(v=\sum_\sigma v_\sigma\) in the complex isotypic
decomposition. For an allowed generator \(H\),
\[
 (H-\alpha(H))v=0
 \quad\Longleftrightarrow\quad
   (\alpha_\sigma(H)-\alpha(H))v_\sigma=0\quad\text{for every }\sigma .
\]
Intersecting these conditions over \(H\) gives precisely the
summands with full character \(\alpha\). The hyperspecial factor
has trivial level character at an allowed split prime. A nonzero
summand there is spherical, so its Hecke polynomial reads the
unramified parameter with the rank-\(a\) twist
\(\lvert\det\rvert^{(1-a)/2}\). The tuple, lift and multiplicities
are common across the coefficient factors. At the two marked
type primes away from \(\ell\), only the prescribed residue
character acts on inertia, with its inverse at the conjugate
place. The Weil--Deligne and infinity comparisons use the
designated embeddings. \label{cor:automorphic-realization-35}

The block lists are determined by the normalized split polynomials
on the constituents used here. For a block of length \(d\) centred on a conjugate self-dual cusp
\(\pi\) of rank \(a\), its terms, with multiplicity, are
\[
           \pi |\det|^{(d-1)/2-k},\qquad 0\le k<d .
\] In total rank \(r\) the common
Galois normalization adds \((1-r)/2\). Relative to the normalization
in rank \(a\), the requisite algebraic shift is \(u+(a-r)/2\).
The allowed infinity parameters are regular with the block's
parity, so the shifted cuspidal Galois realizations in this comparison
are supplied by the regular algebraic cuspidal input. The norm twist
changes the polarizing character by a norm power. For a regular
algebraic polarized cusp over a CM field the attachment is
continuous, with the normalized local polynomial at each unramified
place away from the residue characteristic. This is the part of
Barnet-Lamb, Gee, Geraghty and Taylor's attachment theorem used here.

To pass from the split Frobenius elements to the Galois group, compare each coefficient of the two direct-sum polynomials. The difference is continuous, so
\[
 \{g:\det(1-X\rho_1(g))=\det(1-X\rho_2(g))\}
\]
is closed and contains a dense set. Here \(\rho_1,\rho_2\) denote the two direct sums in the comparison. In particular there is equality of the normalized
polynomials at all but finitely many primes. Cancel the common
normalization and apply Jacquet and Shalika's uniqueness for the
isobaric sums with their real norm twists. Unitarity and the
central characters distinguish the real exponent of each
cuspidal term. Fix one cusp; let \(e\) be the largest exponent
left and \(m(e)\) its multiplicity. Its longest remaining blocks
have length \(2e+1\), with one occurrence of \(e\) per block.
Update all exponent multiplicities by
\[
 m_{\rm new}(s)=m(s)-m(e)\,\mathbf 1_{\{e,e-1,\ldots,-e\}}(s).
\]
Induction on the number of terms recovers even the blocks with
overlapping strings. \cite{blggt-attachment,jacquet-shalika-ii}\label{cor:automorphic-realization-36}

For each row type and residue place which has an integral Galois
realization, reduce one such realization and take the direct sum of
its composition factors. The kernel is open since it contains
the open kernel of the reduction. Any two admissible realizations
of this type at the same place read the same split polynomials
through the common coefficient embedding. Reduce the coefficient
identities in that marked valuation ring outside the finite union
of their bad primes.

For two such residual representations \(\bar\rho_1,\bar\rho_2\),
the coefficient functions are locally constant. Equality on the dense
split Frobenius conjugates therefore gives, for every Galois element \(g\),
\[
 \det(X-\bar\rho_1(g))=\det(X-\bar\rho_2(g)).
\] The polynomial is unchanged by
semisimplification, by the block triangular determinant formula.
If there is no such realization the comparison has an empty source
and the auxiliary residual representation may be taken trivial.
This supplies the semisimple data for arbitrary prescribed
\(S_v\). The comparison is coefficientwise also at the small residue
characteristics. Its assertion about the polynomials is independent
of the degree of the coefficient enlargement.
\label{cor:automorphic-realization-37}
\end{proof}

\begin{proposition}[Simultaneous coefficient requests and local row orders]\label{prop:joint-row-orders}
At each prescribed tuple and lift, finitely many conditions
on the compatible semisimple readings may be imposed at once
by adjoining the union of their coefficient requests. Write \(l_\alpha\) for the ordinary \(O\)-lengths of the indicated localizations. On these specializations
\[
 \eta_j(v)=\frac{\sum_\alpha l_\alpha}{\operatorname{length}_O(O/\ell O)};
\]
in the untwisted case all the lengths and the order vanish. The comparison uses the same good
projector and is unchanged by enlarging the coefficient DVR
at the given place.
\end{proposition}

\begin{proof}
\proofstep{Meet the field requests by union before applying specialization.}
Use the localized order of Proposition~\ref{prop:marked-hecke-localization}.
Every coefficient condition here has the form of a finite set
\(S_v\) at each original place \(v\), followed by an assertion
about every compatible semisimple realization whose marked embeddings
contain \(S_v\). For the two conditions with requests \(S_v^{(1)}\) and \(S_v^{(2)}\), impose the single request
\[
 S_v=S_v^{(1)}\cup S_v^{(2)}\qquad\text{at each original place }v .
\]

The coefficient construction just given works for that union and
keeps the conductor, level and tuple unchanged. For a compatible specialization the two implications can therefore be applied at the same homomorphism. For an occurring eigencharacter, choose a nonzero eigenvector \(u\) in the geometric row. The resulting specialization is the homomorphism
\[
    \mathbb T\longrightarrow O,\qquad T\longmapsto\lambda_u(T),
                         \qquad Tu=\lambda_u(T)u .
\]
Here the row is the one constructed in Lemma~\ref{lem:constructed-signed-coefficient-row}. Its integrality follows from the stable ball of Lemma~\ref{lem:marked-semistable-local-ball}. At a new ray type, occurrence on the tested Hecke algebra is proved in Lemma~\ref{lem:new-ray-type-occurrence}. In particular the order identity,
which is proved on specializations satisfying the full coefficient
ring conditions, can be imposed together with those conditions.

At the fixed tuple and weight only finitely many systems of split
Hecke polynomials occur. Enlarge their exceptional set to contain
the conductors and excluded level places. For each realized system
fix one continuous integral representative
\(\rho:G_E\to{\rm GL}_a(\overline{\mathbb Q}_\ell)\).
Writing the algebraic closure as a countable union of finite
extensions \(F_b\), one has
\[
 G_E=\bigcup_b\rho^{-1}({\rm GL}_a(F_b)).
\]
These are closed subgroups of the compact group \(G_E\); by Baire
one is open. Adjoin the matrix entries of finitely many coset
representatives to obtain a finite field of definition \(F_\rho\).
There are finitely many extensions of \(F_\rho\) of degree
\(\le a\) in the chosen closure. Their compositum contains every
root of \(\det(T-\rho(g))\). By split Frobenius density any other
realization of the system has the same characteristic polynomials.
Request generators of these finitely many composita. An unrealized
system imposes no condition.

To impose the integral block condition, first specialize with the
rational eigenvalues, split polynomials and decomposition of the
good locus. Every occurring good block lies over a good residual
block here by descent of residual irreducibility. For each such
block use its separating Galois element and the matrices at its
finitely many characteristic-zero points. At each original \(v\)
make the \emph{single} request
\[
 S_v=S_v^{\rm split}\ \cup\
    \bigcup_{\text{these blocks and points at }v}
       \{\text{roots of the chosen characteristic polynomials}\}.
\]
Integral base change carries those matrices into the final blocks.
Thus all the root tests and order comparisons hold on the one
specialization obtained after the union. \label{cor:automorphic-realization-38}

For the row orders, fix a lift \(j\) and a place \(v\) with
nonreference tuple. At this tuple the rational definite space
carries an integral lattice \(L\). Let \(p_G\) denote its rational good projector.
Put \(G=\operatorname{im}p_G\) for the good space and \(B=\ker p_G\)
for its complement in the ambient row. Use the matched theta reading,
so \(B\) is the kernel of its adjoint. The tuple identification,
twisted vacuum identity and rational class-sum comparison give the
following pair of geometric lattices in this row:
\[
 L_{\rm upper}=p_G L=G\cap(L+B),\qquad L_{\rm lower}=L\cap G .
\]
For the middle equality write \(a=p_Ga+(1-p_G)a\) for \(a\in L\).
Conversely a vector of \(G\) in \(L+B\) projects to itself.
The lower identity follows by pulling \(L\) back along the injective
good-space embedding. These identifications give the same lattice
quotient in the two realizations. \label{cor:automorphic-realization-39}

The good projector in the localized definite model is the sum of
the good point projectors. For the finite type check on the
original row the supported lifts at a marked split factor are
\((g,(g^{-1})^{\mathsf t})\) in the paired integral coordinates.
The principal congruence kernels fix the forms at both marked
row levels. If \(e=\sum_g c_g[g]\) is the finite idempotent, choose a lift \(\tilde g\) of each element by surjectivity.

On the representation factoring through reduction one obtains
\[
 \sum_g c_g\,\rho(\tilde g)=\sum_g c_g\,\bar\rho(g)=\bar\rho(e).
\] A nonzero vector in the good part of an isotypic
summand retains that idempotent after projection onto an
irreducible copy in which it has a nonzero component. One full
good product of finite series projectors therefore acts
nontrivially in its local components. The tame type dictionary
identifies the degree-\(r\) support polynomials with their
inertia polynomials including multiplicity. Conjugation gives
the dual series, and the good set is dual stable. This gives
the original good-label test carried into the localization.
\label{cor:automorphic-realization-40}

For the matched theta range, fix a localized index and denote the
upward theta map by \(T\). On every nonzero lower isotypic summand the adjoint identity reads
\[
             T^*T=\lambda_\sigma I,\qquad \lambda_\sigma>0.
\]
At a split prime of norm \(q\), remove the common normalization.
The lifted root multiset is the smaller one together with the
specified string of length \(t\). Its largest modulus is
\(q^{(t-1)/2}\), attained exactly once. For a block of length \(d<t\)
with cuspidal root \(\alpha\),
\[
 |\alpha|q^{(d-1)/2}<q^{d/2}\le q^{(t-1)/2}.
\]

If a block of length \(d\ge t\) has centre of rank \(a\), its rank
alone consumes \(ad\) dimensions. Thus
\[
                ad\le r+t<2t\le2d,\qquad a=1 .
\] That center is unitary, so its largest modulus forces
\(d=t\), and uniqueness of the top root forces a single such
block. Dividing that root by \(q^{(t-1)/2}\) recovers its center.
Reciprocity and split Frobenius density identify the global
character. Remove one copy of this block, with multiplicities.
At a good point \(x\),
let \(U_x\) be the sum of the lower summands with the corresponding
parameters, and \(V_x\) the point eigenspace. The parameter
comparison gives \(TU_x\subset V_x\). Any other lower summand lifts orthogonally to \(V_x\).

Indeed, writing \(C\) for the common tuple block and \(L,L'\) for the lower parameter multisets, cancellation gives
\[
 L\sqcup C=L'\sqcup C\quad\Longrightarrow\quad L=L' .
\] Consequently \(T^*V_x\)
is contained in \(U_x\). The downward nonvanishing condition
is injectivity of \(T^*\) on the whole \(V_x\) for a good point.
Section~\ref{subsec:selected-point-adjoint} gives injectivity of
\(T^*|_{V_x}\) before the row orders are inserted, for the specified
rank-one string \(t\ge r+5\), good finite type and marked moduli
model. For \(v\in V_x\) one has
\[
 v\perp TU_x\ \Longrightarrow\
 T^*v\in U_x\cap U_x^\perp=0\ \Longrightarrow\ v=0.
\]
The inclusion \(TU_x\subset V_x\) is therefore an equality.
The class-sum identity and the good-type incidence rule now
identify the sum of these ranges with the range of the good
projector. The remaining simultaneous eigenspaces give its
complement. \label{cor:automorphic-realization-41}

Localize and extend scalars to the marked coefficient DVR.
The integral double cosets preserve \(L\), with the same conductor
in both its rational and spectral realizations. It decomposes
integrally as \(\bigoplus_\alpha L_\alpha\) over the occurring
residual blocks, and \(G=\bigoplus_\alpha G_\alpha\),
\(B=\bigoplus_\alpha B_\alpha\) after scalar extension.
The two lattices in the quotient
\[
          \frac{G\cap(L+B)}{G\cap L}
\]
therefore split by the same residual idempotents. Write \(Q_\alpha\) for the resulting localized quotient, and
\(O_v\) for the coefficient DVR at \(v\mid\ell\). The normalized total
order of this quotient is
\[
              \frac{\sum_\alpha\operatorname{length}_{O_v} Q_\alpha}
                   {\operatorname{length}_{O_v}(O_v/\ell O_v)} .
\]
\proofstep{Normalize length by the original coefficient valuation.} This convention is invariant under extension of DVRs
at the same place. Flatness carries the intersections and images
to the extended comparison, and a composition factor over the
old ring becomes a module of length the ramification index.
For an extension \(O_v\to O_w\) of ramification index \(e\), this
is the cancellation of Lemma~\ref{lem:retained-dvr-normalized-length}:
\[
 \frac{e\operatorname{len}_{O_v}(M)}{e v_v(\ell)}
          =\frac{\operatorname{len}_{O_v}(M)}{v_v(\ell)}
\]
for the index quotient \(M\). At an untwisted place \(v\) there are no localization terms, and for every lift one obtains
\[
 \sum_{\text{localizations at }v}\operatorname{ord}(M)=0=\eta_j(v).
\] This proves
the identity with the prescribed \(\eta_j(v)\) at every place
and every \(1\le j\le m\), with no factor from the local degrees
or the number of localizations. The equality uses the good
projector and integral comparison on the same chosen tuple.
\label{cor:automorphic-realization-42}

\proofstep{Carry the order to the level-three row specialization.}
Apply the archimedean and type comparison to the localization that
meets the united request. It uses the Lie and Galois comparisons
on its compatible local parameter data and obtains the normalized
character, local type and level data on that specialization.
The row-lift family there is read in the same finite group and
marked adelic frame. Write \(o_j(v)\) for the integral localization order on this row. The input to the normalized family is the pointwise equality
\[
 o_j(v)=\eta_j(v)\qquad(1\le j\le m,\ v\text{ a coefficient place}).
\] Transporting the row to the normalized theta family preserves
the original weight functions and character. The grid list transports by the
preserved marked readings at each allowed depth. The initial
pass and the rebuilt pass apply this argument to their respective
strict-henselian inputs. Their resulting row specializations are identified by the same
row and character comparisons.
\label{cor:automorphic-realization-43}
\end{proof}

\begin{lemma}[Type of the universal theta series]\label{lem:universal-theta-type}
Fix an index, a lift, and a sign in the prescribed row.
Let \(M_0\subset M\) be the submodule with the marked level
type. If \(\Theta=\sum_{g\in\Delta}a_g[g]\) is the universal theta series in \(M\otimes_O O[\Delta]\), then
\[
 a_g\in M_0\quad(g\in\Delta),\qquad \Theta_{\rm ordinary}\in M_0.
\] After an admitted specialization the
series lies in the scalar-extended type submodule.
\end{lemma}

\begin{proof}
Fix an index, a lift \(1\le j\le m\), and a sign.
Write \(M\) for the coefficient module of this row
and \(M_0\) for the submodule satisfying the marked level
type. Writing \(T_h\) for a level operator and \(A_h\) for its prescribed
type action in these coordinates, the defining equations are
\[
                    M_0=\bigcap_h\ker(T_h-A_h)\subseteq M.
\] The injection into the complex model intertwines these
operators, so their kernel conditions may be tested there. There the theta transformation law
gives the marked type. The same law holds for each admitted
twisted generator and for every admitted universal coefficient
viewed in the original row. The real comparison uses the
uniform lift reading and the identical complex models
of the actual row. \label{cor:automorphic-realization-44}

\proofstep{Extract the coefficient for each element of the finite twist group.}
The universal series lies in \(M\otimes_O O[\Delta]\)
for the finite twist group \(\Delta\) and the coefficient
ring \(O\) at this index. On the group-algebra basis the extraction of one coefficient is
\[
  (1\otimes[\delta]^*)\!
    \left(\sum_{d\in\Delta}m_d\otimes[d]\right)=m_\delta .
\] In the admitted case its
class-sum formula is a finite sum over terms of that label.
Each term belongs to the prescribed type submodule \(M_0\) of \(M\), so their sum is in
\(M_0\) as well. Off the specified range the universal
series is zero.

Since the coefficient maps recover
every coefficient of any expansion in the group basis,
this proves the type assertion for every expansion
and every label. The ordinary theta series itself lies
in \(M_0\) by the same faithful type test. Thus every
datum in the designated lift range may be used for the
base equivariance of the theta data. For an admitted
specialization the corresponding type assertion holds
in the scalar-extended type submodule.
\label{cor:automorphic-realization-45}
\end{proof}

\begin{established}[The real row and measured columns at the prescribed scale]\label{hyp:real-row-measured-models}
Fix an arithmetic row specialization of Proposition~\ref{prop:joint-row-orders}. Let \(F_u^0\) be
its degree-zero column form model for the word \(u\), including
the chosen column values, the level at \(T_u\) and the two
type markings. If \(V_{\rm good}\subset V_{\rm full}\) are the two rows recorded here, the dimension conditions used in the replacements are
\[
 \dim F_u^0>0\quad\text{for each word }u,\qquad\dim V_{\rm good}\le\dim V_{\rm full}.
\]
The local matrix models, Petersson measure and quotient measures on the flag fibres can all be computed in the same Hermitian coordinates. In these coordinates the unit lowest vectors satisfy the absolute Bessel bounds used for the finite term expansion.
\end{established}
\begin{proof}
The dimension formula of Proposition~\ref{prop:marked-column-dimension} is a product of positive volume, weight and finite marking factors. The other dimension inequality follows from inclusion of vector spaces. Keep this row and these columns while making the following two
replacements of the spectral data. In the first replacement use
the scale of the actual Haar data of the form, the discrete-series
parameters obtained from the labelled Harish--Chandra weights
and signatures, and the numerical discriminant of the fixed
scale context.

For each finite place the local matrix model is
obtained from the completion of the marked quadratic algebra
with its involution. The adjoint form in the local matrix model comes from scalar extension of the same Hermitian form. For a Gram matrix \(H\) in the row or column coordinates, the adjoint is read locally by
\[
 A\longmapsto H^{-1}\bar A^{\,t}H ,
\] This gives the placewise
isomorphisms for the replaced spectral data, with the same
exponent and size context. Division lattices, their height
functions, the chosen character, and the localization grid
are transported by the row identifications. The transport for
the initial state carries the origin comparison, and that for
the rebuilt state carries the origin and geometric norm
comparisons.
\label{cor:real-row-realization-1}

Let \(X\) be the adelic unitary model of this Hermitian space.
The measures in \(X\) are chosen together with the standard
additive character. Start with the self-dual lattice from
the reference form \(F_{u_*}^0\). At the marked split primes keep the fixed frames. For each finite place \(v\), take the self-dual Haar measure on the additive unipotent coordinates for the standard character \(\psi_v\). In the Fourier convention this means
\[
 \widehat{\widehat f}(x)=f(-x)\qquad(f\text{ a Schwartz--Bruhat function}).
\]

At the real places use the prescribed unit
lowest vectors and the Bessel functions of the discrete
series of the selected weight. The local choice is made
for the weight gaps of the scale context. The stabilizer
quotient measures come from the Gram fibres. Use the
convergent unitary Siegel--Weil identity for the two
stabilizers and the rational quotient density to normalize
their covolumes simultaneously. In this normalization the
Petersson measure and the product measures on the flag
fibres refer to the same ambient Haar, as in the fibre
comparison of Section~\ref{sec:absolute}. The gap bound
gives the weighted absolute integrability of the Bessel
functions through the matrix-coefficient estimate of Proposition~\ref{prop:bessel-sign-gap-decay} and the local Bessel projection in Section~\ref{sec:absolute}. These are the measure and decay
conditions used in the finite term expansion.
\label{cor:real-row-realization-2}
\end{proof}

\begin{proposition}[Identification of the measured columns and their fibres]\label{prop:measured-columns}
For the Hermitian column models and measure of Lemma~\ref{hyp:real-row-measured-models}, match the ordered
frames, the distinguished line of norm \(-T_u\), and the characters
at both type primes. The resulting finite adelic transport induces an
isometry of the degree-zero column spaces. For a nonempty term fibre \(\mathcal F\), write \(\mathcal F/U\) for its lift classes modulo flag unipotents, \(S\) for the marked target stabilizer and \(X_0\) for the target isotropic space. There is an injection
\[
 \mathcal F/U\lhook\joinrel\longrightarrow\operatorname{im}\bigl(S\to\operatorname{Aut}(X_0^\perp/X_0)\bigr),
\]
with finite target. Corank and the existence of a scalar lift are
independent of the representatives of a pair with equal row reading.
\end{proposition}

\begin{proof}
\proofstep{Matching the local lattices.}
For each sequence, two ordered frames have the same Gram matrix.
Their hyperbolic pairs have pairing one, and the distinguished
line has norm \(-T_u\). If \(A\) is the matching between their spans and \(G,G'\)
are their Gram matrices in any bases, the Gram equalities give
\(A^*G'A=G\). This extends by Witt's theorem. In the
orthogonal complement the two degree-zero lattices are
self-dual off \(T_u\). At an unramified inert place split off a unit-norm line at each step. If a chosen vector has unit norm \(a\), surjectivity of the norm on units gives a unit \(c\) with
\[
 N(c)=a^{-1},\qquad \langle cv,cv\rangle=N(c)a=1 .
\] At a split place the two halves of a self-dual
lattice are paired duals. At \(T_u\) the dual quotient on
each split factor has length one. Indeed the primitive
distinguished line splits rationally from its orthogonal
complement in a self-dual middle lattice.

Projection gives
the two dual quotients the same length, and the norm
\(-T_u\) gives that length at the split factors. In split
coordinates the resulting lattice chain is a full lattice
and the inverse image of a hyperplane modulo a uniformizer.
A change of basis \(A\) on one factor forces the inverse transpose on the dual factor, since for paired coordinate vectors
\[
 (Ax)^t(A^{-t}y)=x^ty .
\] At the two type
primes a diagonal unit correction in the common finite
reduction matches the superdiagonal characters in the
chosen frames. Conjugation by the inverse of
\(\operatorname{diag}(1,b,b^2,\ldots)\) multiplies all
first-superdiagonal entries by \(b\). In reversed dual
coordinates inverse transpose negates their sum. The
primitive residue characters are related by a nonzero
residue scalar, so one chooses \(b\) to match the inverse
character in the sharp reading and lifts its entries to units.

\proofstep{Transport of the class sets.}
These local comparisons are identity at almost all places.
They give a finite adelic transport identifying the
column levels and compact models. For the induced bijection
\(a\mapsto a'\) on double quotients, let \(\mathcal I_a\) be the
unitary value comparison in the compact model at \(a\). It sends a
class function \(F\) to \(F'\) with \(F'(a')=\mathcal I_a F(a)\).
Covariance extends this identity to all representatives, and
\[
          \sum_{a'}\|F'(a')\|^2=\sum_a\|F(a)\|^2
\]
uses the bijection and the finer-level trivial stabilizers. The chosen unit vectors in the compact
highest-weight representations identify the invariant
norms. At the finer level the stabilizers are trivial by
the split type-prime argument of
Section~\ref{subsec:selected-word-level}. Indeed the class bijection preserves the values in the identified compact representations. With trivial stabilizers the norm comparison is simply reindexing: for a column function \(f\) and class bijection \(b\),
\[
 \sum_C\|f(C)\|^2=\sum_{C'}\|f(b^{-1}C')\|^2 .
\] It identifies the actual column forms
entering the Bessel and Gram computations.

The finite
term fibres use these very classes and their shorter
and complete flag representatives. Fix one lift \(s_0\) of a term in a nonempty fibre. Write
\(X_0\) for the isotropic space of its marked target flag and
\(\Lambda'\) for its target lattice. Modulo the two flag unipotents the transition for any other lift is \(s_0 c\). If \(b_i\) are the marked flag basis vectors and \(v\) the marked vector, then \(c\) satisfies
\[
 c\Lambda'=\Lambda',\qquad c b_i=b_i\ (\text{each }i),\qquad c v=v .
\] Choose one such \(c\) per lift. Its
action on \(X_0^\perp/X_0\) lies in the finite image of the marked
lattice stabilizer. Here the definite lattice automorphism bound
applies with the signature at stage zero and the positive norm
parameter \(T_{u'}\) of the target word. Write \(\bar c\) for the induced endomorphism of
\(X_0^\perp/X_0\). If \(\bar c=\bar c'\), put \(u=c^{-1}c'\).

The three conditions defining the flag-unipotent equivalence are
\[
 u|_{X_0}=1,\qquad u\Lambda'=\Lambda',\qquad
 (u-1)X_0^\perp\subseteq X_0 .
\]
The last inclusion is just \(\bar c^{-1}\bar c'=1\); the first
two follow from the marking and the integral lattice. Such \(c\) also normalizes
the target flag unipotents, so the two lifts represent the same
class. The chosen map into the finite image is injective.

\proofstep{Independence of representatives.}
Consider also a pair of lifts having the same row reading, and a
second pair with the same respective terms and again equal row
readings. In the transport just used, equality of the two marked
embedding readings identifies the induced Levi actions. The neat
level makes their rational stabilizer representatives equal. In particular the relative target transporter induces the identity on the flag quotient:
\[
 (c_2c_1^{-1})v-v\in X_0\qquad(v\in X_0^\perp).
\] Conjugate it by the first lift. The new relative
matrix differs from the old one by this conjugate and the
lattice-preserving source and frame unipotents. They carry the
space defining the scalar \(s_0(\cdot,\cdot)\) isomorphically
onto the new space, preserving its codimension, and carry any
scalar lift to a scalar lift. Interchange the two pairs to get
the reverse implication. This proves representative independence
before using absolute convergence to reindex the expansion.
\cite{companion-row-count}\label{cor:real-row-realization-3}
\end{proof}

The product discrete series exists with the prescribed
labelled real parameters. Its formal degrees are computed
in the actual product Haar. In the real tie their product
is the one at the true scale, and the finite factor is the
value of the actual counting projector at the identity.
The comparison of quotient volume with mass unfolds the
same product measures and uses the definite genus
decomposition. Hence the identity contribution in the
measured model is exactly the predicted \(D_*\) of
Section~\ref{sec:count}. In particular it uses the
finite group quotient and level index of this row. The
finite good operator is fixed under the conjugate reading
in the matched frames. The local series projector has
its coefficients paired by dualizing the supported tame
labels and reversing their order. This is the coefficient
identity which compares the two readings of the operator.
\label{cor:real-row-realization-4}

\begin{proposition}[The isometry of the refined row]\label{prop:refined-real-row}
Let \(J\) be the sharp geometric level, \(n_P\) its level index
in the normalization of the Petersson form, and \(M(J)\) the full
multiplicity space. Evaluation at the chosen unit lowest vectors
identifies \(M(J)\) with the conjugate measured row, with one positive
squared-norm ratio \(c\). It respects the good projector. If \(j\) is the normalized composite on the arithmetic good row and \(V_{\rm full}\) the row before projection, then
\[
 \|jv\|^2=\|v\|^2\quad(v\text{ in the arithmetic good row}),\qquad \dim V_{\rm full}=D_{\mathrm{all}}.
\]
\end{proposition}

\begin{proof}
\proofstep{The two finite characters.}
Let \(J\) be the geometric level defining the refined row space and \(K_2\) the
corresponding subgroup in \(X\). Use the simultaneous
finite reduction of the type frames to first match
principal kernels. In the common finite group the two
marked sharp levels are conjugate. Lift the finite
conjugator inside the compact group and correct the
finite translate. If \(\iota\) is the group identification,
the corrected translate \(t\) satisfies
\[
 t^{-1}\iota^{-1}(K_2)t=J,\qquad
 \chi_J(t^{-1}\iota^{-1}(k)t)=\overline{\chi_X(k)}
 \quad(k\in K_2).
\]
\proofstep{Inverting lowest-vector evaluation.}
The central idempotents are unchanged by this conjugation
of frames. The real representation comparison uses the
contragredient discrete series at the same infinitesimal
character in the conjugate Hilbert space. Choose the
geometric unit lowest vectors corresponding to those
of \(X\). For an intertwiner in the full sharp multiplicity
space \(M(J)\), evaluation at these lowest vectors after
the translate \(t\) defines a map \(L\) onto the conjugate
row space of \(X\). Automorphy and the displayed level
and character conditions give the prescribed finite
equivariance. At infinity, for a unit lowest vector \(w\)
in a discrete series, let \(P_{v,w}\) be the operator
of its normalized matrix coefficient taking the column
\(w\) to \(v\). Schur orthogonality gives the rank-one matrix-unit identities
\[
 P_{w,w}^2=P_{w,w},\qquad P_{v,w}P_{w,w}=P_{v,w}.
\]
A row function \(f\) in the image of \(P_{w,w}\) gives the
intertwiner \(v\mapsto P_{v,w}f\) by covariance under
real translation of the matrix coefficients.

Evaluation
at \(w\) recovers \(f\). Apply the product of these
operators, using the contragredient identification and
the inverse finite translation. This gives surjectivity
and the inverse to evaluation. The construction goes
through the smooth discrete-series vectors, using the
matrix-coefficient projection with the actual formal
degrees. Rapid decay of cusp forms makes the evaluations
square integrable.
Petersson integration on the two identified quotients has
a single positive ratio \(c\). With the real vectors of
norm one this gives
\[
 \|LS\|^2=c\|S\|^2\qquad(S\in M(J)).
\]
The map is complex linear into the space of conjugated
functions with its complex structure there. The
finite action on both sides now intertwines with
conjugation of its coefficients. Thus \(L\) carries the
good part of \(M(J)\) onto the conjugate good subspace,
as seen by applying the central good projector inside
the full evaluation isomorphism.
\label{cor:real-row-realization-5}

\proofstep{Normalization on the good subspace.}
The geometric realization of the arithmetic row
\(E_{\rm row}\) is an isomorphism onto the linearized
good multiplicity space at \(J\). In the multiplicity evaluation, write \(d\mu\) for the measure before sharp-level normalization. The Petersson measure on the row is
\[
 d\mu_{\rm sharp}=n_P^{-1}d\mu .
\]
The good space at \(J\) is contained in the principal
model. Invert the faithful principal realization on
this subspace to obtain the geometric row embedding.
Its composition \(e_E\) with multiplicity realization
has image the whole good space of \(M(J)\), by surjectivity
of this geometric realization. Using the same unit tuple
in both evaluations gives \(\|e_E x\|^2=n_P\|x\|^2\).

It follows that
\[
             j=(c n_P)^{-1/2}L e_E
\]
is the isometry onto the conjugate good row space of the
measured model. The full evaluation bijection identifies
its full dimension with the recorded \(D_{\mathrm{all}}\).
This dimension is from the full space at the refined level before applying
the good projector. The two positive factors in \(j\) cancel
exactly in the pairing for every host and every allowed
word level.
\label{cor:real-row-realization-6}
\end{proof}

For the second replacement choose the spectral datum
using the row with its normalized Petersson measure and the comparison map \(j\). This choice fixes
\(t,c,j\) before any replacement of the division lattices
and heights on the word spaces. If a replacement retains
the degree-zero levels and readings, the term fibres are
carried by the two induced bijections of class labels.
The fibre of a term is the image of its original finite
fibre under those bijections, with equal function values.
Write \(b\) for the isometry induced by reindexing the good space. The transported row map is the composite, so for every row vector \(v\),
\[
 j_{\rm new}=b\circ j,\qquad\|j_{\rm new}v\|=\|jv\|=\|v\|.
\]

To apply
this comparison here take the original division lattices
and their height function. Transport first through the
symmetrized column form model, which preserves the
degree-zero values, then through the equivalence to the
actual rebuilt carriers of the same columns. It follows
that the new comparison of the row spaces at the refined level uses the original
column levels and actual column forms. The full
dictionary, the real tie, the discriminant and the
transported order and character data constitute the
level-four output for this pass.
\label{cor:real-row-realization-7}

\begin{lemma}[The ambient lattices in the used tuples]\label{lem:real-ambient-tuples}
At the sharp and bottom levels separately, inclusion of the used
tuples in the full tuple list identifies their coefficient spaces,
integral ambient lattices and positive Hermitian forms. Writing \(b\) for this identification at one of the two levels, \(\Lambda\) for the lattice, and \(H\) for the form, one has at each used tuple \(\mu\)
\[
 b\Lambda_{\mu,{\rm used}}=\Lambda_{\mu,{\rm full}},\qquad b^*H_{\mu,{\rm full}}=H_{\mu,{\rm used}}
\]
at every embedding where the form is positive. These identifications use the
same degree-zero columns as Proposition~\ref{prop:refined-real-row}.
\end{lemma}

\begin{proof}
The ambient lattices at the sharp and bottom levels use
the rational coordinates of this row specialization. Fix one of these levels and a tuple \(\mu\). Let
\(V_\mu\) be its coefficient-field coordinate space, with
lattice \(\Lambda_\mu\) and Hermitian matrix \(H_\mu\).
Scalar extension at the marked embedding sends its coordinate vectors to a
basis at the marked embedding, and carries \(\Lambda_\mu\)
onto the marked ambient lattice. At a positive embedding
\(\sigma\), choose an orthonormal frame for
\(\sigma(H_\mu)\).

Let \(F\) send the marked basis coordinates to coordinates in this orthonormal frame at \(\sigma\). Thus \(F^*F=\sigma(H_\mu)\), and
\[
 \|F\sigma(x)\|^2=\sigma(x)^*F^*F\sigma(x)=\sigma(x)^*\sigma(H_\mu)\sigma(x).
\] This is the positive form
in the comparison of the ambient norms. At the marked
embedding the original Hermitian and operator identities
give the designated complex form and adjoint under the
reading. Finite generation and spanning of the ambient
lattice follow from those of \(\Lambda_\mu\) and the two
identities for the basis and lattice.

The used tuple list maps into the full tuple list by the
same tuple of local character labels. The prescribed tuple
map agrees with this inclusion because the row indices are
defined by those labels. At either level the coordinate map on the selected tuples and the restriction of the full coordinate map agree on every basis vector. If these maps are \(f,g\),
\[
 f\!\left(\sum_i a_i e_i\right)=\sum_i a_i f(e_i)=\sum_i a_i g(e_i)=g\!\left(\sum_i a_i e_i\right).
\] The identity on coordinates, with the
dimension identified through that equality of indices,
therefore identifies the rational readings, integral
lattices and positive forms at every embedding where the
form is positive. This constructs the two ambient
identifications consumed by the division construction.
The sharp and bottom matrices and lattices are taken from
their respective levels.
\label{cor:real-row-realization-8}
\end{proof}

Both passes have this construction at their own
level-three states. The first level-four state is
completed to a built member before taking its signed
export. The rebuilt pass makes the corresponding choices
on the carried specialization. Their division and theta
constructions consume the corresponding real isometry and
ambient identifications on that pass.
\label{cor:real-row-realization-9}

\paragraph{The theta maps on this row.}
Fix one specialization and its levels at the primes assigned to the sequences. Denote the
unitary spaces in its chain by \(Y_0,\ldots,Y_m\), identifying
the definite top with \(Y_m\) by its given isometry. For a used tuple \(\mu\), take the multiplicity space \(M_j\) in the row at the refined level on \(Y_j\). The downward maps form the chain
\[
 M_m\xrightarrow{D_{m-1}}M_{m-1}\longrightarrow\cdots\longrightarrow M_1\xrightarrow{D_0}M_0 .
\] We need a simultaneous realization of the \(D_j\) by theta
integrals with the Petersson pairings on these very spaces.
This is stronger than existence of each local intertwiner.

\begin{proposition}[Isometric realization of the theta chain]\label{prop:isometric-refined-theta-chain}
At the fixed specialization and assigned levels, choose
isometric readings \(R_j\) of every multiplicity space on
a used tuple. With the stated Petersson measures and chosen
Weil and level characters, write \(\Theta_j^{\downarrow}\) for the downward theta integral transported back to the multiplicity rows by these isometries. On every edge of the chain, including the last,
\[
 \Theta_j^{\downarrow}=D_j.
\] On unused tuples the zero row realizes the
prescribed zero maps.
\end{proposition}

\begin{proof}
\proofstep{Fix the inner product multiplier at each stage.}
Use the identifications of the row spaces at the refined level for both the multiplicity
pairing and the downward integral. At zero these include \(M_0\)
in the base row, at intermediate indices in the cuspidal
multiplicity spaces, and at the last index in the definite one.
For the sharp index \(n_P\), the stage-zero measure is
\(n_P^{-1}\nu_0\) with pairing multiplier \(c_0=1\).
At the intervening stages the original quotient measures
have \(c_j=n_P^{-1}\). At the top the local definite factors \(M_{X,\tau}\) from the row enter with the common level divisor:
\[
 c_m=\frac{\prod_\tau \dim M_{X,\tau}}{n_P}.
\] The latter
identity is Schur orthogonality at the definite factors.
The same lower multiplier \(c_j\) occurs when the downward
model integral is paired against a vector of \(M_j\).
\label{cor:level-four-core-1}

\proofstep{Choose a common evaluation tensor serving both edges at each interior stage.}
For each edge its nonzero real intertwiner has a nonzero
matrix coefficient on a polynomial test and a pair of pure
smooth tensors, since such tensors span the algebraic tensor
product and separate its elements by the inner product. Fix
one such polynomial and a pair of tensors giving a nonzero coefficient
on each edge, including the last. Vary each tensor factor in the
finite real span of the corresponding chosen factors.

Let \(P_j(\mathbf t)\) be the resulting coefficient on edge \(j\) in the
common real coordinate space. At the pair chosen for edge \(j\) its coefficient is nonzero, so
\(P_j\not\equiv0\), even with the unused coordinates adjoined.
The polynomial ring has no zero divisors, so
\[
 \prod_j P_j(\mathbf t)\overline{P_j(\mathbf t)}\not\equiv0 .
\] At a point where the product is nonzero,
normalize each factor to unit norm and scale the test on each
edge to make its coefficient one. Thus both occurrences of a
stage use one and the same evaluation tensor. The resulting
polynomial test has nonzero Schwartz function by the model
intertwining identity.

\proofstep{Scale each row isometrically and then normalize its edge integral.}
Write \(\kappa_j>0\) for the required quotient multiplier on
\(Y_j\), including its finite and real Petersson factors.
Let \(r_j>0\) be the quotient-measure ratio. In the convention of this comparison the two pairings are related by
\[
 \langle\ ,\ \rangle_{Y_j,\kappa_j}=(\kappa_j/r_j)\,\langle\ ,\ \rangle_{\rm transported\ model}.
\] The Haar ratios include transport from the definite
top. Evaluate the multiplicity map on the chosen tensor and
multiply the transported cusp function by
\(s_j=(c_jr_j/\kappa_j)^{1/2}\). Schur orthogonality and rapid
decay give the original \(M_j\) inner product, because the multiplier for two readings is
\[
 \frac{\kappa_j}{r_j}s_j^2
 =\frac{\kappa_j}{r_j}\frac{c_jr_j}{\kappa_j}=c_j .
\] Denote this isometric reading by \(R_j\).

On the edge with upper space \(W=Y_{j+1}\), let \(r_W^I\)
be the transported quotient ratio for its integration variable.
Use the oriented Weil kernel from \(W\) to \(Y_j\), with its
recorded splitting characters. Multiply its test by
\[
 a_j=\frac{c_jr_j r_W^I}{\kappa_j\kappa_{j+1}s_js_{j+1}} .
\]
The integral against \(R_{j+1}y\), followed by the lower
Petersson pairing against \(R_jx\), gives the model pairing
with scalar
\[
 (\kappa_{j+1}/r_W^I)\, a_js_{j+1}s_j\,(\kappa_j/r_j)=c_j .
\]
The real intertwiner coefficient in this pairing is one.
The model chain identity now gives \(\langle D_jy,x\rangle\)
for all \(x,y\). This argument uses the convergent upper
theta integral of the marked cuspidal chain. It includes the
last edge by transporting its Weil action as well as its Haar.

Use the unique marked infinity data, the standard character of
sign \(-1\), and a unitary Weil realization for that character.
For a finite tuple character \(\chi\) there is an integer \(n\ge1\) such that, writing \(F\) for the totally real subfield of the auxiliary CM field,
\[
 \chi^n=1,\qquad \chi|_{\mathbb A_F^\times}=1 .
\] Their splitting characters
agree with the marked ones by the character identities on
ideles. On unused tuples the specified \(D_j\) are zero, so zero
row functions and the original nonzero tests give both required
edge integrals. The row isometry is required on used tuples.
In the theta realization the quotients and their Petersson
multipliers are exactly the ones just used. Equality of these
integrals therefore gives the required downward maps under
pairing, and the construction has the prescribed Weil
equivariance and finite characters on the same realization.
\label{cor:level-four-core-2}
\end{proof}

\begin{proposition}[Divisor with the marked row and finite support]\label{prop:core-divisor-support}
Use the refined theta chain of
Proposition~\ref{prop:isometric-refined-theta-chain} with
its sharp and bottom row comparisons at the same specialization.
The finite prime support of the divisor, denoted \(S_D\), is
\[
 S_D=\{p:\text{some component of }\xi_\mu\text{ above }p\text{ is ramified}\},
\]
where the character is the actual tuple character of the divisor. It has the same marked
coefficient-place valuations and row norm at the
initial flag datum.
\end{proposition}

\begin{proof}
Use the simultaneous theta kernels of
Proposition~\ref{prop:isometric-refined-theta-chain}.
Let \(E/F\) be the marked CM extension, \(P(v)\) the chosen prime
of \(E\) over each \(v\) of \(F\), and \(\xi_\mu\) the
original character of the column tuple. Use the rational lift
spaces, bases, integral modules and Gram matrices of the fitted
row. The sharp scalar is fixed before the finite type is chosen. For each type \(\sigma\) used here it is
\[
 \kappa(\sigma)=\kappa_{\rm sharp}.
\]

For each spectral label choose a reference finite type
above it if one occurs. A fallback type exists by positivity
of the reference primitive block. At the distinguished embedding
the logarithmic value is \(\log\lambda_{\mu_0,j,\sigma}\) from
the row. At another embedding the value at a tested lift is zero
on a zero isotypic component and is the logarithm of the chosen
type's reference reading otherwise. One choice of type for each
label works at every lift.

\proofstep{Identify finite support by contracting the same ramification places.}
Put \(S_\mu=\{v:\xi_\mu\text{ ramifies at some }P\mid v\}\).
The character pins give ramification at the chosen \(P(v)\) if
\(v\in S_\mu\), so the rational set is
\[
 Q_\mu=\{v\cap\mathbb Z:v\in S_\mu\}
      =\{P(v)\cap\mathbb Z:v\in S_\mu\}.
\]
For the reference \(\mu\) it is empty. For every tuple it
excludes the two type primes, and for \(\mu(w)\) it excludes
the residue characteristic at \(w\), by the local twist choice
on the support and the trivial tuple elsewhere.

If \(p=P(v)\cap\mathbb Z\) is a ramified character place,
the pin there gives \({\rm N}P(v)=p\) and
\(p\in P(v)\setminus P(v)^2\). The norm determines the residue degree, while \(p\notin P(v)^2\) determines the ramification index:
\[
 f(P(v)/p)=1,\qquad e(P(v)/p)=v_{P(v)}(p)=1 .
\]
Both indices are multiplicative in the tower, so the norm
downstairs is the same rational prime. This gives the
degree-one local field used for the tame determinant character
\(\delta\) and the unramified unitary quotient characters \(\nu_i\).

\label{cor:level-four-core-3}

Here is the local normalization. Write \(q\) for the residue
cardinality, \(A\) for the local ring of integers, and
\(\varpi\) for its uniformizer. On the edge out of stage \(j\)
use the Levi with blocks \(r,jt\). Write \(P\) for the parabolic in the second block with \(j\) blocks of size \(t\). Its spherical representation is obtained from normalized induction of the ordered determinant characters, whose inducing representation is
\[
 (g_1,\ldots,g_j)\longmapsto\prod_{i=1}^j\nu_i(\det g_i).
\] Take the algebraic cyclic span of its smooth spherical unit vector
\(v_0\). The maximal compact \(K\) has a one-dimensional fixed space
\(\mathbb C v_0\) by Iwasawa decomposition.

If a nonzero invariant
subspace contains \(u\ne0\), there is a translate \(gu\) with
\(\langle gu,v_0\rangle\ne0\), since the translates of \(v_0\)
span. With probability measure on \(K\),
\[
  \int_K \pi(k)\pi(g)u\,dk=\langle \pi(g)u,v_0\rangle v_0 ,
\]
using the linear-first Hermitian convention and \(\|v_0\|=1\).
Smoothness makes the average a finite linear combination, so
the invariant subspace contains \(v_0\) and equals its cyclic span.
For a diagonal entry \(z\) in block \(a=0,\ldots,j-1\) the
Borel law has factor
\(\nu_{a+1}(z)|z|^{(jt-t)/2-at}\) by the half modulus of
this parabolic. The first Levi block acts trivially in this
spherical model. For \(j=0\) the second block is trivial.
\label{cor:level-four-core-4}

Let \(J\) be the parahoric reducing to the stated upper block
parabolic \(P\) modulo \(\varpi\). For \(g\in J\) let \(d(g)\) be its lower diagonal block. Multiplication respects this block modulo \(\varpi\), so for the character of \(J\) one obtains
\[
 \chi_J(gg')=\delta(\det d(g))\delta(\det d(g'))=\chi_J(g)\chi_J(g').
\]
On \(P\cap J\) it agrees on the fixed unit vector with the
inducing law, since the diagonal blocks have unit determinant
norm. Extend this vector to a section on \(PJ\) by the two
transformation laws and set it to zero elsewhere. On the
maximal compact \(K\) its squared norm is the indicator of
\(K\cap PJ\). Give \(J\) volume one and scale the section by
\(\operatorname{vol}(K\cap PJ)^{-1/2}\). The matrix coefficient
at the identity is then one.

Use the dimension of this type
space and the doubling integral of the section in the divisor
data. The first-edge value can be read directly. The two-block test in these coordinates is supported on integral invertible \(X\) and \(\varpi\)-divisible \(Y\). There it has the value
\[
 (X,Y)\longmapsto\delta(\det X)^{-1},\qquad X\in\mathrm{GL}_r(A),\quad Y\in\varpi M_{t,r}(A).
\] Testing also \((Xg,Yg)\) forces
\(g\in{\rm GL}_r(A)\). On this compact group the ratio of the
tests cancels \(\delta(\det g)\) in the Weil action, and the
norm and unramified factors equal one. With the two full
integral matrix boxes of additive volume one, the first
value is therefore
\[
                q^{-tr}\prod_{i=1}^{r}(1-q^{-i}).
\]
If the residue field has cardinality \(q\) and the block has size \(d\), this product is the invertible-matrix density. Counting linearly independent columns gives
\[
 \frac{|\mathrm{GL}_d(\mathbb F_q)|}{q^{d^2}}=\prod_{a=0}^{d-1}\frac{q^d-q^a}{q^d}=\prod_{j=1}^d(1-q^{-j}).
\] All of these values refer to the primes of
\(Q_\mu\) just selected. \label{cor:level-four-core-5}

\proofstep{Realize the original coefficient valuations at the constant-field flag.}
For the initial flag datum use the constant function field equal
to the row field, with full section space there, one auxiliary
level supported on the two type primes, zero curve degrees and
zero maps from the row to sections. At each marked place it is the actual coefficient valuation that has to be recovered. For an embedding \(e\) into a valued field with valuation ring \(O\), the test is
\[
 a\in\mathcal O_{\rm coeff,\,w}\quad\Longleftrightarrow\quad e(a)\in O
\]
on the coefficient field. Embed the Witt fraction
field for each place in \(\C\) using a transcendence basis and
extend its valuation ring. The Witt field has at most continuum
cardinality. Conjugacy of the extensions of the \(p\)-adic place
to \(\overline{\Q}\) permits an embedding of the coefficient
number field inducing the prescribed place on restriction.
A complex field automorphism taking that embedding to the
distinguished one transports the Witt embedding and its
extended valuation ring together.

If \(e\) denotes one of the transported embeddings and \(O\) its chosen extended valuation ring, the local reading on the coefficient field is recovered by inverse image:
\[
 e^{-1}(O)=\{a:e(a)\in O\}.
\] Adjoin
this flag datum to the divisor, retaining its original row
basis, lattice, norm, logarithmic values, sharp factor and
quotient-prime readings. This is the base case of the geometric
condition closed under changing division lattices and heights.
The later geometric degree estimates use the parameter models
constructed on a completed point.
\label{cor:level-four-core-6}
\end{proof}

\begin{proposition}[The integral kernel on the selected divisor]\label{prop:level-four-integral-kernel}
On the selected divisor of Proposition~\ref{prop:core-divisor-support}
and the prescribed stage lattices,
let \(C_j\) be evaluation on the class charts, \(\Lambda_j\) the integral row, \(O\) the upper valuation ring and \(M_j^{\rm int}\) the module of equivariant functions with integral tensor Weyl values. At each valid positive stage, \eqref{eq:core-upper-image} reads
\[
 O\,C_j(\Lambda_j)=M_j^{\rm int}.
\] At each selected pair of
class representatives the transported Taylor kernel is
the prescribed entry series in the adapted bases.
The first maps and all these entries use the same divisor
and the ramified word lattice.
\end{proposition}

\begin{proof}
Fix a word \(u\) and a coefficient place \(w\) satisfying the
local orbit condition of Section~\ref{subsec:column-orbit-length}.
Use its marked column and the tuple \(\mu(w)\). Write \(k\) for
the starting column rank, \(t\) for the block size added on each
edge, and \(\Delta=\lceil x^{c_D}\rceil\). The bound
\(c=16k\Delta\) works for the column entries and the first
filling weight simultaneously on all sequences. The interior Harish--Chandra entries \(u\) and the two extreme entries off the top have the respective bounds and values
\[
 3\Delta-\tfrac12\le u\le4k\Delta,\qquad u_{\max}=4(k-2)\Delta+4k\Delta+\tfrac12,\qquad u_{\min}=-\tfrac12 .
\]
The rank shift in passing to column entries has absolute value
at most \(k\).

At the first edge the gap above the lower
endpoint \(b_1\) is at least \((t+3)/2\). The prescribed
top floor, lower bound for the larger weights \(b_1\ge\lceil b\rceil\), and thin
choice \(b_1=\lceil w''\rceil\) each give \(b_1+1\ge w''\),
where \(w''=b'/(xn)\). Uniformly in \(k\le x^5\), \(n\le x^3\),
we have eventually
\[
 c+k+2\le35x^{c_D+5}\le x^{c_D+6}\le w''\le b.
\]
Thus even after subtracting the shift \((k+t+1)/2\) the first
edge weight is at least \(c\). The size conditions give
\(c\le c_G(L)\) on every compatible size record \(L\).
This is a single bound chosen before those records.

In each orbit of the local step \(h\), take the residue of
least numerical label. If \(b\) is the initial label in an orbit and \(d\) the order of the step, the arithmetic and geometric rules both take the minimum of the same finite set:
\[
 \min_{0\le s<d}\,\operatorname{label}(b+sh),\qquad d=\operatorname{ord}(h).
\] Send its translate by \(sh\) to the
corresponding marked embedding \(\tau\) at slot \(s\).

At top slots the endpoints are the first and the negative of
the last column entry. Put zero in these two positions off
the top. After subtracting the rank shift \(a_j\), the added weight at lift \(j\) and embedding \(\tau\) is
\[
 w_j(\tau)-a_j .
\] These are the
negatives of the first entries of the successive added blocks
of a dominant column. They increase weakly, and the first
dominates the negative last entry of the original block.
For every compatible canonical completion the orbit, the
positions, the top endpoints and the shifted weights just
computed are therefore exactly its prescribed filling. This
comparison concerns the columns and weights alone, so it
persists for each choice of the subsequent level and maps.
\label{cor:level-four-core-7}

Take the standard primitive lift chain at the \emph{actual}
\(Q_\mu\). Use its nonzero unfolding factors, integral jet
identities and adapted bases. In the basis of an entry place one
rescales the stage \(j\) frame at \(\tau\) by the prescribed unit
\(a_\tau^j\). A product of minors of shape \(\lambda\) has total degree \(|\lambda|\). Thus stage-\(j\) scaling by the unit \(a_\tau^j\) gives the multiplier
\[
 (a_\tau^j)^{|\lambda|}=a_\tau^{j|\lambda|}.
\] The scalar substitution in the kernel polynomial contributes a further
\(a_\tau^{|\lambda|}\). The homogeneous factors on the edge are therefore
\[
 a_\tau^{j|\lambda|}a_\tau^{|\lambda|}
        =a_\tau^{(j+1)|\lambda|},
\]
the factor for the target frame. The source and target determinant
normalizations occur in the unfolding scalar with the
corresponding powers. Dividing by this nonzero scalar cancels
them. The unit preserves both integral tests and the local norms.
Thus the entry coefficient in the adapted bases is the Taylor
coefficient of the kernel for these lift objects.
\label{cor:level-four-core-8}

For each marked local component \(v\) let \(A_v\) be its ring
of integers and \(L_v\) the prescribed full lattice at one
stage. The first split embedding \(\iota_v:E\to F_v\) on \(F\)
is its completion at \(v\), and these places of \(F\) are
distinct. Put
\[
 \Lambda=\{z\in E^d:\ \iota_v(z)\in L_v\text{ for all marked }v\}.
\]
It is an \(\mathcal O_E\)-module. Its image spans \(L_v\) over
\(A_v\): approximate a basis matrix for \(L_v\) within its
open integral basis coset at \(v\), and zero within the other
lattices at all the other marked places, by a matrix over \(F\).
Its columns belong to \(\Lambda\) and are a local basis at \(v\).
This weak approximation argument asserts a span at the specified
places. The global finite lattice with all its other localizations
prescribed was constructed separately above.

For the local class map first fix an away-place representative
\(a\), then a local matrix carrying the reference lattice to
a tuple \(L'\). Every admissible adelic class is obtained in
this manner. For a fixed away class \(a\), let \(\Gamma_a\) denote its rational stabilizer. Two presentations with local tuples \(L'_1,L'_2\) give the same class exactly when
\[
 L'_2\in\Gamma_a\cdot L'_1 .
\] This is the double
coset equivalence after fixing \(a\). Stage zero uses the given
word-column identification. At a positive stage the reference
tuple uses its member charts, with just the end chart at stage
\(m\).

For a nonreference tuple use the symmetric chain of its
own character. For the frame matrix \(J=C^*C\), the reference basis \(B=C^{-1}\) therefore satisfies
\[
 B^*JB=(CB)^*(CB)=1.
\] Its lattice and the
level lattice agree upon completion at \(v\), since the latter
is between the lattice and its multiple by the word modulus.
In the nonreference frame \(J=B=1\) the last-coordinate
condition disappears there since \(\pi_T\bar\pi_T=T\) is a
unit. Moreover the quotient-prime set for \(\mu(w)\)
excludes its own characteristic. Thus on both kinds of charts
the local level is the stabilizer of the indicated lattice.
\label{cor:level-four-core-9}

Here is the integral image assertion needed for the upper maps.
Let \(e:H\to K_w\) be the coefficient embedding, \(O\) the upper valuation ring and \(A\) the valuation ring of \(H\) at \(w\). The compatibility of the two integrality tests is
\[
 e^{-1}(O)=A,\qquad h\in A\ \Longleftrightarrow\ e(h)\in O .
\]
Let \(V_j,\Lambda_j\) be the rational row and its integral
lattice at a valid positive stage. Let \(M_j^{\rm int}\) be
the local module of equivariant functions on the class
presentations, taking values in the tensor Weyl lattices of
the filled weights. Then the chart reading \(C_j\) satisfies
\[
                O\,C_j(\Lambda_j)=M_j^{\rm int}.
                \tag{*}\label{eq:core-upper-image}
\]

To see the two integral tests involved, write
\(d=\lfloor j(t+1)/2\rfloor\) for the determinant shift. At
a class \(y\), the normalized evaluation multiplies the
original frame evaluation by
\[
             \chi_D(\det y)^d\prod_\tau
                   \det\bigl(\tau(B)^{-1}z_\tau\bigr)^d .
\]
In a rational basis \(B'\) of the local dual lattice the change
of determinant weight multiplies the coefficient polynomial by
\[
             s=\chi_D(\det y)^d
                \prod_\tau\tau(\det B'/\det B)^d.
\]
The finite determinant-character values are rational over the
coefficient field by the marked character reading and its
values on uniformizers. The projections of the bases are
the marked embeddings of the host. To check the valuation of
\(s\), use finiteness of the class set to write
\(y^{b+l}=\gamma y^b k'\) with \(l>0\), \(\gamma\) rational
unitary, and \(k'\) in the level.

On each first split component,
both \(B'^{-1}yB\) and \(B'^{-1}yk'B\) have unit determinant
by the local basis equality and self-duality. Taking determinants
in the power relation shows that
\(z=(\det B'/\det B)^l/\det\gamma\) is a unit there.
Also \(z\bar z=\det(B'^*JB')^l\) is a unit by self-duality,
so the second split component gives a unit. Applying the character law at \(\gamma\) and \(k'\) now yields
\[
 s^l=\prod_\tau\tau(z)^d .
\] Hence \(s\) and its inverse are
integral. This calculation uses both projections even when the
slot reads the conjugate embedding. \label{cor:level-four-core-10}

The coefficient change at a translated lattice uses
\((B^{-1}\gamma B')^{-1}B^{-1}=B'^{-1}\gamma^{-1}\).
On an ordinary block this substitutes the translated bases
in products of minors. On a dual block the complementary-minor
identity for the inverse transpose, with the column order
reversed, gives integral signs and the determinant correction.
The inverse substitutions give both directions. Tensor over
the slots. The level character multiplies by a root of unity,
and descent along \(e^{-1}O=A\) tests coefficients on a basis
of the rational Weyl lattice. Thus for \(x\in V_j\)
\[
             C_jx\in M_j^{\rm int}\quad\Longleftrightarrow\quad
             x\in A\Lambda_j .
\]
For surjectivity after taking the span, read a local function
at representatives of the finitely many class orbits.

At each
representative its values obey the stabilizer equations with
the inverse character of the glued away-level element. These
are linear equations on a finite-dimensional rational Weyl
space, hence finitely many equations suffice over every scalar
extension. Conversely each rational solution gives a complex
form on that adelic class. Indeed if \(y=\gamma yk'\) the
determinant shift contributes one to the stabilizer equation
by the character law. At the chosen orbit representative the stabilizer is conjugate at level to \(k'^{-1}\). Since its character is constant on conjugacy classes, inverting the glued character recovers
\[
 \chi(k'^{-1})^{-1}=\chi(k').
\] Ordinary positions use the
marked entry of \(\gamma\), dual positions its conjugate in
the reversed dual frame. The resulting frame polynomials are
independent whenever their restrictions to invertible blocks
are independent, since these blocks are Zariski dense and the
frame changes are invertible there. Extend by equivariance and
by zero on other classes.

Independence after restriction to the open set of invertible blocks therefore bounds the integral rank by
\[
 \operatorname{rank} M_j^{\rm int}\le \dim_{\mathbb C}(\text{complex row at }j).
\] In the other direction, the original
frame evaluations separate \(V_j\). A nonzero minor of finitely
many rational coefficient readings remains nonzero over \(K_w\),
giving the opposite bound. Equal rank and that same minor give
a rational preimage for every rational target in the span.
\proofstep{Obtain the integral generators after imposing the stabilizer equations.}
Finally choose bases of all marked local lattices by simultaneous
weak approximation at the distinct places of \(F\). Their local
Weyl lattices extend rational \(A\)-lattices. Intersecting with
the stabilizer equations commutes with extension to \(O\) by
flatness over the DVR \(A\). The resulting rational integral
generators extend by equivariance to the finitely many orbits.
Their rational preimages pass the two-way integral test. This
proves \eqref{eq:core-upper-image} with the original lattice at
each stage. \cite{companion-selected-core}\label{cor:level-four-core-11}

\proofstep{Pull the entry back by the inverse of the exact coordinate substitution.}
For a pair of adjacent chosen bases with image matrices \(A_j\)
and \(A_{j+1}\), the forward substitution of a kernel block is
\[
 X\longmapsto A_{j+1}XA_j^{-1}
 \quad\hbox{or}\quad
 X\longmapsto A_{j+1}^{-\mathsf T}X A_j^{\mathsf T}
\]
according to the marked ordinary or dual position. Pull the
fixed entry series for the pair of global classes back by the
inverse linear substitution. Composing with the forward map
returns it coefficient by coefficient, since the matrices
cancel and all substitutions preserve the degree filtration.
The split idempotents identify the embeddings with those of
the chosen column, and the orbit-position equality identifies
the variables. Attach these series, the local class maps and
the image lattices \eqref{eq:core-upper-image} to the fixed
filling. Use their first comparison maps and the ramified
submodule specified by the word column. This gives the kernel
at every selected pair with the adapted entry identity and the
same divisor throughout. \label{cor:level-four-core-12}
\end{proof}

\paragraph{Assembling the theta maps at the specified level.}
Use the section and kernel of Proposition~\ref{prop:level-four-integral-kernel}
with the theta maps of Proposition~\ref{prop:isometric-refined-theta-chain}.
The two ambient identifications, the quotient primes and the
theta realization were obtained for the same level-four data.
Form the object whose sharp and bottom ambient modules are those
just identified, whose divisor has those primes, and whose
kernel and lattice comparisons are those just constructed for
that divisor. Include its original degree-zero column and
level, normalized character, signed row, and a CM gauge for the
host. The divisor comparison applies to every object with
this divisor, base row and row reading. All three equalities
hold here by construction. Thus the logarithmic type choice,
the rational roots on the base row, the stable multiplicity
reading and the sharp normalization belong to the assembled
object.

In particular, writing \(S\) for the selected twist support, nontriviality on \(S\) gives
\[
 \mu(w)=\mu_0\quad\Longrightarrow\quad w\notin S .
\] At its conjugate the reference tuple is again
assigned. The untwisted identities therefore give zero for the
signed twisting order on this tuple. For any admissible size record of the same scale, a nonzero forced content would itself put the place in the support. Thus for each such record
\[
 v\notin\operatorname{supp}\quad\Longrightarrow\quad\text{forced content at }v=0 .
\] The comparison of
the local first maps and the orbit weights above holds for
each such record with its canonical coordinates. This gives
the required row and column data at this level.
\label{cor:level-four-core-13}

Perform this construction first on the initial member. Its
signed export below uses the rational and integral theta
readings of this member. After transporting its auxiliary
character, index and series to the new field, construct the
second member using its own measured row and columns by this
same argument. On that member the transported norm choice
gives both the original norm condition and the condition on
all twist places. The signed tests descend along the kept
identifications. The second member consequently supplies a
single row, kernel and divisor over the chosen levels.
\label{cor:level-four-core-14}


\begin{auxlemma}[The theta generator on a semistable split Hecke branch]\label{lem:marked-hecke-semistable-unit}
Let \(R\) be a complete discrete valuation ring of mixed
characteristic \((0,\ell)\) with finite residue field. Use a
principally polarized semistable PEL fibre \((A,\lambda)\),
with \(\mathcal O_E\)-action of Rosati involution \(a\mapsto\bar a\).
Write \(\mathcal P_A\) for the rigidified Poincar\'e bundle on
\(A\times A^\vee\). The marked theta line on this fibre is
\[
 L_A=(1,\lambda w_0)^*\mathcal P_A,\qquad
            w_0\in\mathcal O_E,\quad w_0+\bar w_0=1 .
\]

Extend \(L_A\) cubically to the identity component of the
semistable model. For a split Hecke prime \(p\ne\ell\) away
from the marking, the \(\mathcal O_E\)-linear branch and its
target line satisfy
\[
       f:A\longrightarrow B,\qquad f^\vee\lambda_B f=p^2\lambda,
              \qquad L_B=(1,\lambda_Bw_0)^*\mathcal P_B .
\] The theta trace on this branch carries an integral
generator to an integral generator after the finite extension
defining the branch. This holds also for \(p=2\) when this split
prime is admitted.
\end{auxlemma}
\begin{proof}
\proofstep{Identify the model norm for this theta line.}
After finite extension take split Raynaud uniformization with period
lattice \(Y\), character lattice \(X\) and principal polarization
\(\Phi:Y\overset{\sim}{\to}X\). Write \(t(y,x)\) for the
biextension trivialization on \(Y\times X\) and set
\[
     b(y,z)=-\log\|t(y,\Phi z)\| .
\]
This is the positive definite monodromy pairing.

The Poincar\'e trivializations on the two period axes are rigidified.
Translating both arguments multiplies the axis factors by the
biextension value at the pair of periods. At the rigidified origin
the axis factors are one, so the graph trivialization is
\[
                     c(y)=t(y,\Phi w_0y)\qquad(y\in Y).
\] The resulting line on the abelian part is the same graph
pullback of its Poincar\'e bundle. Bilinearity identifies the
rigidified biextension with its pullback by \((-1,-1)\), so both
graph lines are symmetric. Functoriality with respect to
\(w_0\) and the adjoint identity give
\[
          -\log\|c(y)\|=b(y,w_0y)=b(y,\bar w_0 y)
                      =\tfrac12 b(y,y).
\]
The graph description thus specifies the characteristic of the
actual line. \cite[preprint, Section~3.4 and Theorems~3.10, 4.10]{frss-nonarch-theta}

For a nonzero principal theta section \(s\), the coefficient at
\(\Phi y\) in the Raynaud expansion is its coefficient at zero
translated on the good abelian part and multiplied by \(c(y)\)
in the corresponding Poincar\'e fibres. Translation preserves the
model norm there. Write \(\gamma\) for the negative logarithm of
the norm of the zero coefficient at the generic point of the good
special fibre. In skeleton coordinates \(z\in Y_{\mathbb R}\),
the Gauss formula and the quadratic correction for the cubical
metric give
\[
 \begin{aligned}
   -\log\|s(\xi_z)\|_{\rm can}
    &=\gamma+\tfrac12 b(z,z)
       +\min_{y\in Y}\{\tfrac12 b(y,y)+b(y,z)\}\\
    &=\gamma+\tfrac12\min_{y\in Y}b(y+z,y+z).
 \end{aligned}
\]
Indeed the correction \(\tfrac12 b(z,z)\) cancels the change
\(b(y,z)+b(y,y)/2\) in the period functional equation. It gives the
rigidified canonical metric by compatibility with pullback by
\([2]\), which multiplies the quadratic form by four.
\cite[Proposition~9.1 and its proof]{dejong-shokrieh-theta}

The displayed minimum is zero at the origin of the skeleton.
The supremum on each retraction fibre is attained at its Gauss point.
At the identity component the negative logarithm there is the
vertical order of the section in the cubical extension, using the
valuation with \(-\log|\pi|=1\). Hence
\[
       -\log\sup_{x\in A^{\rm an}}\|s(x)\|_{\rm can}
                    =\operatorname{ord}_{0}(s).
\]
A section on the generic fibre extends to the identity model
precisely when this order is nonnegative: the model is smooth,
the special fibre is integral, and all its other divisorial points
lie on the generic fibre. Thus a generator has supremum norm one.
The same calculation applies to \(B\). \cite[Proposition~10.3 and
Corollary~10.4]{dejong-shokrieh-theta}

\proofstep{Compare the two kernel actions at the actual level.}
Let \(C=\ker f\), \(C_0=C\cap A[p]\) and \(\mathcal L=f^*L_B\).
Write \(\mathcal B_A=(1\times\lambda)^*\mathcal P_A\) on
\(A\times A\), with \(\mathcal B_B\) defined similarly.
Functoriality and the polarization relation give rigidified
isomorphisms of biextensions and of their restrictions to the graph:
\[
 (f,f)^*\mathcal B_B\simeq([p],[p])^*\mathcal B_A,\qquad
             \mathcal L\simeq[p]^*L_A.
\]
For \(d\in C_0\), the two lifts of translation by \((d,0)\)
differ by an invertible function on the proper connected
\(A\times A\), hence by a scalar. Its restriction to
\(A\times\{0\}\) is one by rigidification. Apply the same
argument on \(\{0\}\times A\) to \((0,w_0d)\), using
\(w_0d\in C_0\) by \(\mathcal O_E\)-linearity. Composing along
the graph identifies the splittings \(U\) from \(f\) and \(V\)
from \([p]\):
\[
                    U_d=V_d\qquad(d\in C\cap A[p]).
\]

Put \(S=[p]^*s\). The trace is \(C\)-invariant and therefore
defines a section \(s'\) on \(B\) by
\[
                       f^*s'=\sum_{c\in C/C_0}U_cS .
\]
The rigidified pullback isomorphism preserves the canonical metrics
by their multiplication law. Thus \(U_c,V_a\) are isometries for
the supremum norm: each is a deck action on its own pullback line.
The summands have distinct \(A[p]\)-characters by the perfect
commutator pairing on \(A[p^2]\), since \(A[p]^\perp=A[p]\) there.
Over a finite extension splitting the torsion, the zero-character
projector therefore gives
\[
        |A[p]|^{-1}\sum_{a\in A[p]}V_a(f^*s')=S .
\]
Its denominator is an \(\ell\)-adic unit. Apply the ultrametric
inequality first to the trace and then to the projector:
\[
 \|f^*s'\|_{\sup}\le\|S\|_{\sup}\le\|f^*s'\|_{\sup}.
\]
Surjectivity of \(f\) and \([p]\) now gives
\(\|s'\|_{\sup}=\|s\|_{\sup}\). The first part of the proof now identifies the
integral generator on the target.

The complementary isogeny
\(f'=\lambda^{-1}f^\vee\lambda_B\) satisfies \(f'f=ff'=[p^2]\).
After extension both maps act on the identity components of the
N\'eron models. Since \(p\) is a unit, \(p^{-1}df\) and
\(p^{-1}df'\) are mutually inverse integral maps on each saturated
eigensummand of the Lie lattice. Thus the tangent frame and theta
generator remain integral on every such branch, including traits
with toric reduction. Counting the branches gives an integral sum
on these frames. Identifying this operator with a spectral Hecke
operator on a specified rational subspace further requires the
equivariance of its realization in the marked sections.
\end{proof}

\begin{auxlemma}[The geometric local unit ball and its scalar extensions]
\label{lem:marked-semistable-local-ball}
Let \(v\mid\ell\) be a coefficient place of the embedded small field
\(F\). At a fixed marked level and weight, take \(V\) in the space of
\(F\)-rational cusp sections.
For every complete discrete valuation ring \(R\) with finite residue
field and fraction field \(K\), take all embeddings
\(F\to K\) over \(v\) and all framed semistable tests \(y\)
over \(R\). The tangent vectors are bases of the saturated
\(\mathcal O_E\)-eigenlattices and the theta section is a generator
on the identity component. Use the extended valuation with
\(v_R(\pi_v)=1\), where \(\pi_v\) is a uniformizer at \(v\). Set
\[
 M_v(V)=\{s\in V:\ v_R(s(y))\ge0\quad
                \text{for every }(R,F\to K,y)\}.
\]

Let \(P\) be an admitted split place away from the level and
splitting-character conductor, of rational characteristic \(p\ne\ell\).
Write \(q_P=|\mathcal O_E/P|=\mathrm NP\). Fix a Hecke degree
\(0\le a\le r\) and let \(n_V\) be the even theta weight.
The geometric branch sum with scalar
\[
 c=q_P^{-a n_V/2},\qquad u_{\rm geom}=c(\mathrm NP)^{-a n_V/2}
\]
preserves this ball if it preserves \(V\). For a field extension
\(F'/F\) and any valuation ring \(O\) of \(F'\) inducing \(v\), it
preserves
\[
  \langle1\otimes M_v(V)\rangle_O
       \ \cap\ (F'\otimes_F V_0)
\]
for every invariant subspace \(V_0\subseteq V\). The same conclusion
holds for an invariant subspace over a larger coefficient field in
the corresponding scalar-extended ambient space.
\end{auxlemma}
\begin{proof}
On a test \(y\) each geometric kernel and its quotient may be defined
over a finite extension. By Lemma~\ref{lem:marked-hecke-semistable-unit}
the transported frame is integral there. Replacing its eigenbasis by
an integral basis of the same flag and multiplying its theta generator
by a unit change the weight evaluation by a unit. With the valuation
still normalized on \(F\), each term in the branch sum satisfies
\[
    v_{R'}(s(y_C))\ge0,\qquad
    v_{R'}\!\left(u_{\rm geom}\sum_C s(y_C)\right)\ge0 .
\]
Both integers \(q_P\) and \(\mathrm NP\) are powers of \(p\).
The estimate descends to \(R\) on the original evaluation.
Preservation of the \(O\)-span follows by linearity. Intersecting
two invariant modules gives the stated restriction for valuation
rings of arbitrary rank.
\end{proof}

In geometric coordinates the normalized trace contributes
\((\mathrm NP)^{-a n_V/2}\). Writing its adelic evaluation as
a sum of right translations introduces the splitting character
of the reading. With \(c=q_P^{-a n_V/2}\) the respective
coefficients are
\[
 u_{\rm geom}=c(\mathrm NP)^{-a n_V/2},\qquad
 u_{\rm ad}=c\,\mu_0(\varpi_P)^{-a}.
\]
To compute the valuation of the character, choose \(h>0\) with
\(P^h=(\alpha)\) and \(\alpha\equiv1\) modulo its conductor.
The algebraic infinity type expresses a fixed positive power of
\(\mu_0(\varpi_P)^h\) as a root of unity times a monomial in the
conjugates of \(\alpha\). Each conjugate has all its finite prime
divisors over \(p\), so
\[
                    v(\mu_0(\varpi_P))=0\qquad(v\nmid p).
\]

A finite family of CM evaluations can be used to bound a prescribed
collection of vectors inside a full \(M_v(V)\). Its fixed inverse
matrix costs the power of \(\pi_v\) in
Lemma~\ref{lem:marked-bounded-denominator}. If the realized
spectral span is the quotient of the same class-input space by the
kernel of the geometric evaluation, this construction gives a
specific integral module in that quotient, by inverse image under
the injection into cusp sections. Using its ball for an independently
specified rational row requires the specified reading to identify
the two evaluations and their normalizations.

\begin{auxlemma}[Finite harmonic jets in a semistable theta frame]
\label{lem:marked-bounded-jet-frame}
Fix the embedded CM field \(E\), the dimensions, an order \(d\)
of differentiation and a coefficient place \(v\mid\ell\). Use the
principal graph line of Lemma~\ref{lem:marked-hecke-semistable-unit}.
On each of the semistable tests in
Lemma~\ref{lem:marked-semistable-local-ball}, evaluate the harmonic
jet of order \(d\) in any integral eigenframe. Start with sections
of canonical supremum norm at most one, and form a fixed finite
linear combination of their rigidified isogeny pullbacks and
finite-order theta translations, of the marked tensor weight.
Writing \(J_P\) for this contracted jet, there is a constant \(b\)
depending only on the order, the polynomial, the tensor weight and
the finite operations such that
\[
             v_R(J_P(y))\ge-b\quad\text{for every test }y,
                         \qquad v_R(\pi_v)=1 .
\]
\end{auxlemma}
\begin{proof}
\proofstep{Trivialize the biextension to finite order.}
Let \(G\) be the identity model. The biextension
\(\mathcal B=(1,\lambda)^*\mathcal P\) on the generic fibre extends
to \(G\times G\) as the polar biextension of its cubical line.
In particular the two rigidified restrictions have the generic identification
\[
            \mathcal B|_{(1,w_0)}\simeq L_A .
\]
It extends to the identity model: the
generic isomorphism can have vertical divisor only a multiple
of the integral special fibre, and comparison at the identity
makes its order zero. The same argument extends the rigidified
adjoint isomorphisms for the \(\mathcal O_E\)-action.

Use logarithmic coordinates \(z,w\) on the two formal groups over
the fraction field. In smooth integral parameters the logarithm
is obtained by integrating the invariant differentials with identity
linear term. Its coefficients through degree \(D=\max(2,d)\) and
those of its inverse have bounded denominators depending only on
\(D\): integration divides by integers at most \(D\), and inversion
with identity linear term uses finitely many additions and products.
Thus for the logarithm \(l\) and exponential \(e\) of the formal group there is an integer \(c_D>0\), independent of the test, with
\[
    c_D [z^\alpha]l,\quad c_D[z^\alpha]e\ \in R^{\dim G}
                   \qquad(1\le|\alpha|\le D).
\]
Substitution along integral tangent eigenvectors preserves this bound.

Pick an integral formal frame for \(\mathcal B\) trivialized on
both axes. Denote by \(h_1,h_2\) the logarithms of its discrepancies
under doubling the first or second variable, using the biextension
law to identify the resulting line with its tensor square.
The discrepancy of addition in the first variable is one if either
summand is zero. On putting the summands equal, its logarithm
therefore has first-variable degree at least two. Likewise \(h_2\)
has second-variable degree at least two. Commutativity of the two
doublings gives in bidegree \((a,b)\)
\[
 (2^b-2)(h_1)_{ab}=(2^a-2)(h_2)_{ab}\qquad(a,b\ge1).
\]

Multiplying the frame by \(\exp H\) changes the two logarithms by
\((2^a-2)H_{ab}\) and \((2^b-2)H_{ab}\), respectively. Thus one
can make both vanish by setting
\[
 H_{ab}=\begin{cases}
 -(h_1)_{ab}/(2^a-2),&a\ge2,\\
 -(h_2)_{ab}/(2^b-2),&a=1,\ b\ge2,\\
 0,&a=b=1.
 \end{cases}
\]
For the addition discrepancy of the new frame, simultaneous
doubling of the two summands in the first variable doubles its
logarithm. It has total degree at least two in those summands
and consequently vanishes. The second variable gives the same
conclusion. This constructs a frame compatible with both addition
laws; at order \(D\) only finitely many of the indicated integers
and exponential factorials occur in denominators.

\proofstep{Impose the adjoint identity.}
For \(t\in\mathcal O_E\), the discrepancy between the frame at
\((tz,w)\) and at \((z,\bar t w)\) is the exponential of a bilinear
form \(B_t(z,w)\), since both frames respect both additions.
For a pair of eigenlabels \(\tau,\sigma\), put
\(\sigma^*(t)=\sigma(\bar t)\). On this pair, addition and composition
of endomorphisms give
\[
 B_{t+s}=B_t+B_s,\qquad
 B_{ts}=\tau(s)B_t+\sigma^*(t)B_s,\qquad B_m=0\quad(m\in\mathbb Z).
\]

If \(\tau=\sigma^*\), these identities give a derivation of the
separable number field \(E/\mathbb Q\) with values in the coefficient
space. It vanishes, as follows by differentiating a separable
minimal polynomial. Otherwise choose one primitive integral element
\(t_0\) of \(E\). Replacing the frame by the exponential of a
bilinear form changes \(B_t\) by
\((\tau(t)-\sigma^*(t))\) times that form. Choose its coefficient
to kill \(B_{t_0}\). Commuting \(t\) and \(t_0\) in the adjusted frame gives
\[
       (\tau(t_0)-\sigma^*(t_0))B_t
           =(\tau(t)-\sigma^*(t))B_{t_0}=0 .
\]
The first factor is nonzero, so every \(B_t\) vanishes.

The extra denominator is one of the finitely many nonzero
differences of conjugates of \(t_0\). Each coefficient of
\(B_{t_0}\) is bounded in the actual eigenvectors, since it
uses only the integral adjoint isomorphism and the preceding
finite-order frame. The adjusted frame has the required uniform
denominator when evaluated in those vectors.

The graph of \(w_0\) now gives a formal frame of the theta line.
Two such adjoint-compatible frames differ on the graph by the
exponential of a quadratic form coupling only opposite eigenlabels.
Over all the columns its coefficients give the contractions
\[
       \sum_h\partial_{\tau,i,h}\partial_{\bar\tau,j,h} P=0
\]
of the specified harmonic symbol \(P\). Multiplication by this
exponential therefore leaves its contracted jet unchanged. This
identifies the computed harmonic jet with the one in the marked
formula, including the one read in complex coordinates.

\proofstep{Bound the translated germs.}
A rigidified pullback between the corresponding symmetric lines
is isometric for their canonical metrics by the multiplication
law. A finite-order theta lift above torsion of order \(h\) also
preserves the metric. Indeed take \(m>1\) with \(m\equiv1\bmod h\).
The discrepancy of the two metrics, as a real logarithm \(b\) on
the compact analytic abelian variety, obeys
\[
                  b([m]x)=m^2 b(x)+c
\]
for a constant \(c\), since the two rigid line isomorphisms differ
by a scalar. Surjectivity of \([m]\) forces the oscillation of \(b\)
to equal \(m^2\) times itself. Thus \(b\) is constant, and a power
of the lift is the identity, so \(b=0\).

By the identification at the identity Gauss point in
Lemma~\ref{lem:marked-hecke-semistable-unit}, each resulting
norm-bounded section has vertical order at least zero there.
For exterior products and symmetric powers use the tensor powers
of the same cubical metric at this point. The section consequently
has an integral germ in any integral local frame.
Converting its \(d\)-jet to the graph frame, applying the harmonic
differential and summing the fixed coefficients cost only the
finitely many denominators above. This proves the bound.
\end{proof}

\begin{auxlemma}[Fullness on the rational span of theta classes]
\label{cor:marked-full-class-ball}
Let \(V\) be the image of a finite-dimensional space of fixed
definite class inputs under the marked algebraic theta map. At
every coefficient place \(v\), its local unit ball \(M_v(V)\)
satisfies
\[
      M_v(V)\simeq\mathcal O_{F,v}^{\dim_F V},\qquad
                  F\otimes_{\mathcal O_{F,v}}M_v(V)=V .
\]
Its inverse
image on the class space modulo the kernel gives the same lattice
in that rational realization of the spectral span.
\end{auxlemma}
\begin{proof}
Choose a basis \(s_1,\ldots,s_m\) of \(V\) among the class images.
The principal theta generator has norm one by
Lemma~\ref{lem:marked-hecke-semistable-unit}. Its fixed class
formula consists of the finite operations of
Lemma~\ref{lem:marked-bounded-jet-frame}. Thus a single \(b\ge0\)
makes \(\pi_v^b s_i\) integral on every test. Use the separating
marked CM evaluations on \(V\), at integral frames after one finite
extension of the local field. Their fixed matrix is invertible on
these basis sections. Every vector in the unit ball has integral
readings there, so the inverse matrix bounds its coordinates:
\[
  \pi_v^b\sum_i\mathcal O_{F,v}s_i\ \subseteq\ M_v(V)
       \ \subseteq\ \pi_v^{-c}\sum_i\mathcal O_{F,v}s_i
                         \qquad(c\ge0).
\]
Over the discrete valuation ring this submodule is finite free of
rank \(m\). The assertion includes \(m=0\). Pulling back to the
quotient by the kernel preserves both inclusions. For a different
prescribed realization one needs its identification with this
particular quotient before transporting the lattice.
\end{proof}

\begin{proposition}[Signed theta coefficients on the original specialization]\label{prop:signed-theta-data}
On the initial member with the prescribed row, tuple, and
coefficient embeddings, the normalized signed theta coefficients
are rational in the original coefficient coordinates. Fix an admitted
index \(i\), lift, sign and theta datum \((s,h,P)\) consisting of support,
adelic argument and harmonic polynomial. Let \(L_i^{\rm sup}\) be the
supremum lattice over the coefficient DVR \(O_i\), and \(R\) the signed
series map on the rational row. Write \(x_{Q,M,\delta}\) for the
normalized row vector of the coefficient at \([\delta]\) of the universal
series with auxiliary level \(Q\) and depth \(M\). There is one integer
\(N\ge0\), depending on the fixed data, such that, with
\(x=x_{Q,M,\delta}\),
\[
 \ell_i^Nx\in L_i^{\rm sup},\qquad
 \ell_i^N R(x)\in R(L_i^{\rm sup}).
\]

This holds for every allowed \(Q\) away from the norm product of
\(s\), every admitted depth \(M\), and each \(\delta\) in the resulting
finite ray quotient. The same coefficient
choices give the prescribed residual Galois representations on the
averaged row.
\end{proposition}

\begin{proof}
\proofstep{Compare the real and theta-series realizations of the rational row.}
The real and series realizations of the initial member use the same marked lattice, point evaluations and finite Hecke action. At index \(i\), lift \(j\) and sign \(b\), let
\(O_i\) be the coefficient DVR, \(F_i\) its fraction field
and \(\ell_i\) its residue characteristic. Denote the signed
rational series map by \(R\), with values in the complex
coefficient functions. The matched rational ray-level test
applies to the base, twisted and universal series over the
split deletion places.

The ray-level test uses an orthonormal frame on the
chosen test space and a common nonzero normalization scalar.
One can move the frame to obtain this test. The moved and
original readings agree after the phase and coordinate changes,
so rationality of the three series transfers to the original
reading. In the zero cusp space every test reads the zero vector. Thus if \(A,B\) denote the two readings,
\[
 A=B=0\quad\Longrightarrow\quad A=cB\quad\text{for every }c\ne0,
\]
so the normalization can still be chosen nonzero.
\label{cor:automorphic-realization-46}

\proofstep{Clear the universal coefficients before changing the level or depth.}
A single denominator must work across the allowed levels and depths.

Fix a nonzero index, a lift, a sign and one admitted theta datum
\((s,h,P)\), consisting of its support, adelic argument and harmonic
polynomial. The coefficient estimate requires an integer \(c\)
for this datum. It must apply to every auxiliary level \(Q\) in the
allowed range whose places lie away from the integer norm product
of \(s\), every auxiliary exponent \(M\), and every group-algebra
coefficient.

For a harmonic
polynomial whose evaluation is a generator of the small-field
coefficient lattice, the integral core calculation supplies a
nonzero factor \(u\) and a power of \(\ell_i\). Multiplying
the coefficient by their product gives a sum of evaluations
of integral core forms in the supremum lattice over \(O_i\).
The two constants have different dependencies:
\[
 u=u(i,j,b,s,h)\ne0,\qquad
 c_\gamma=c_\gamma(i,j,b,s,h,P_\gamma).
\]
They are uniform in the admissible auxiliary level, the auxiliary
exponent and the coefficient of the group algebra.
Lift the finite integral expansion of the evaluation of \(P\).
Injectivity of harmonic evaluation gives
\[
 P=\sum_\gamma a_\gamma P_\gamma,\qquad
 c=\max_\gamma c_\gamma,\qquad
 \ell_i^c u\,x(P)=\sum_\gamma a_\gamma\ell_i^{c-c_\gamma}
                              \bigl(\ell_i^{c_\gamma}u\,x(P_\gamma)\bigr).
\]
Here \(x\) denotes the coefficient evaluation before normalization.
The right side is integral for each allowed \(Q,M\) and coefficient.
For \(P=0\) no clearing is necessary. \label{cor:automorphic-realization-47}

Dividing by the common normalization \(\nu\ne0\) gives
\[
 \lambda R(x)=R(x'),\qquad \lambda=\ell_i^c u/\nu,\qquad
 x'\in L_{i}^{\rm sup}.
\]
The \(R\)-images of an \(F_i\)-basis are complex linearly
independent. Hence \(x'_k=\lambda x_k\) in each coordinate.
For any nonzero coordinate this gives
\[
                 \lambda=x'_k/x_k\in F_i
\]
in the chosen embedding. In this case choose \(N\) with
\(\ell_i^N/\lambda\in O_i\); then \(\ell_i^Nx=(\ell_i^N/\lambda)x'\)
is integral for every admissible coefficient vector \(x\). If \(\lambda\notin F_i\), every such vector vanishes. Thus the exponent may be chosen as
\[
 N=0\qquad(\lambda\notin F_i).
\] The exponent depends
on the datum, index, lift and sign and is common to all \(Q,M\)
and coefficients. \label{cor:automorphic-realization-48}

For the full Hecke preservation, take an allowed split prime \(P\)
away from the two type primes and its normalized finite double coset
on the lower row. Use a coefficient valuation of rational
characteristic different from that of \(P\). The framed evaluation
identifies the double-coset action with the corresponding transform
of cusp sections, since their evaluations at the test points agree and
the evaluation map is injective.

For the degree \(a\) double coset, let \(c=q_P^{-a n_V/2}\).
In the geometric frame the scalar and its adelic reading are
\[
  u_{\rm geom}=c(\mathrm NP)^{-a n_V/2},\qquad
  u_{\rm ad}=c\,\mu_0(\varpi_P)^{-a}.
\]
Lemma~\ref{lem:marked-semistable-local-ball} proves the pointwise bound in the geometric frame, using the transported integral generator on each branch. The finite action commutes with the Hecke sum, since it changes the markings on its source and target simultaneously. Intersecting the extended ball with the type or good space gives the asserted preservation.

\label{cor:automorphic-realization-49}

The Galois input begins with the regular algebraic polarized
attachment just used in block uniqueness. For a regular
algebraic real norm shift \(\nu\) of a conjugate self-dual
cusp \(\pi\), use the representation of \(\pi\nu\).
Its labelled de Rham weights are those of the shifted infinity
type with the rank normalization. At a place \(w\mid\ell\) with spherical local component, the unramified case of the stated attachment identifies its local Galois representation as crystalline:
\[
 \dim D_{\rm cris}(\rho|_{G_w})=\dim\rho ,
\]
with the dimensions taken over the respective base coefficient fields. For the potential crystallinity
test, suppose that the local parameter has zero monodromy
and the stipulated local type. First take a conjugate
self-dual algebraic shift of \(\pi\). For an integral norm
exponent take the trivial character.

For a half-integral norm
exponent use the following construction on the marked CM
extension \(E/F\) with the chosen CM type \(\Phi\).
This extension is unramified at every finite place. Let
\(\eta\) be its quadratic idele class character on \(F\). Its local values in the unramified finite normalization are
\[
 \eta_v(\varpi_v)=
 \begin{cases}1,&v\text{ split},\\-1,&v\text{ inert} .\end{cases}
\]
The real factors are sign characters.
The CM type and field supply, for every root of unity \(\zeta\) in
\(E\), an element \(\delta\in E^\times\) for which
\[
 \frac{\delta}{\bar\delta}=\zeta,\qquad
 v_P(\delta)\text{ is even for }P\text{ inert over }F,\qquad
 \prod_{\phi\in\Phi}\frac{\phi(\delta)}{|\phi(\delta)|}=1 .
\]
For \(\zeta\ne-1\), take \(\delta=1+\zeta\); explicitly
\[
                   \frac{1+\zeta}{1+\zeta^{-1}}=\zeta .
\] The marked CM type has a fixed-point-free involution
induced by a nonidentity Galois involution in its indexing
subgroup that inverts the roots of unity. The phases in each
pair are complex conjugates, so their product is one.

For \(\zeta=-1\) take the marked square root of \(-d\).
It has phase \(i\) at each member of the CM type, whose
cardinality is a multiple of four. Both choices have
valuation zero away from the rational conductor of the host.
For the square root this follows from \(d\) dividing the
conductor. For \(\zeta\ne-1\), a prime away from the conductor cannot divide \(1+\zeta\). Otherwise injectivity of reduction on the relevant roots of unity would give
\[
 \bar\zeta=\overline{-1}\quad\Longrightarrow\quad\zeta=-1 .
\] The prime divisors of the
conductor split in \(E/F\), giving the inert valuation condition.
\label{cor:automorphic-realization-50}

On ideles with unit finite components prescribe the product of
the phases at \(\Phi\). On principal ideles prescribe one, and on
finite lifts of uniformizers of \(F\) prescribe the values of
\(\eta\). To verify compatibility, let \(k\in E^\times\) and let
\(m_v\) be the exponents of the chosen uniformizers. Since all ramification indices here are one, the valuation of the product at \(P\mid v\) is the sum of the scalar valuation and the prescribed exponent. The unit condition is
\[
 v_P(k)+m_v=0\quad\Longleftrightarrow\quad v_P(k)=-m_v\qquad(P\mid v).
\] Thus \(k/\bar k\) is a unit of absolute value
one at every embedding. Kronecker's theorem and the preceding
choice of \(\delta\) give
\[
 k/\bar k=\delta/\bar\delta,\qquad a=k/\delta\in F^\times,\qquad
 m_v\equiv-v(a)\pmod 2\quad(v\text{ inert}).
\]
The split uniformizers have assigned value one.

Write
\(\operatorname{ph}_\Phi(x)=\prod_{\phi\in\Phi}\phi(x)/|\phi(x)|\).
The intersection test becomes the quadratic product formula
\[
 \prod_v\eta_v(\varpi_v)^{m_v}
 =\eta_{\rm fin}(a)^{-1}=\eta_\infty(a)
 =\operatorname{ph}_\Phi(a)=\operatorname{ph}_\Phi(k).
\]
The character so defined on the generated open subgroup extends
to the idele group by divisibility of the circle group. It is
continuous because its restriction to the open subgroup is
continuous. It is trivial on principal ideles, unramified at
finite places and restricts to \(\eta\) on the ideles of \(F\).
On an idele \(x\) the product evaluates the character on a norm; triviality of \(\eta\) there gives
\[
 (\xi\xi^c)(x)=\xi(x)\xi(x^c)=\xi(xx^c)=1 .
\]
The factor \(z/|z|\) at a member of \(\Phi\) has exponent
\(1/2\) there and \(-1/2\) at its conjugate, in the convention
of labelled exponents.
\label{cor:automorphic-realization-51}

Call the chosen character in either parity \(\xi\).
Its exponents differ from those of \(\nu\) by integers at each
label. It follows that \(\pi\xi\) is regular algebraic, since
the change is a common integral shift within each infinity
parameter. The quotient \(\nu/\xi\) is an algebraic Hecke
character unramified at finite places. Its rank-one Galois
character is crystalline above \(\ell\).
At the specified local type in rank at least two, local-global
compatibility gives zero monodromy on the de Rham attachment
to \(\pi\xi\). It is potentially semistable and hence potentially
crystalline when the monodromy vanishes. The twist back to the
original attachment satisfies
\[
 r_\ell(\pi\nu)\simeq r_\ell(\pi\xi)\otimes r_\ell(\nu/\xi).
\]
Indeed the two semisimple sides have equal characteristic
polynomials at dense Frobenius classes. The second factor on the
right is crystalline, so the tensor product is potentially crystalline.
\cite{blggt-attachment,caraiani-monodromy,berger-differential}\label{cor:automorphic-realization-52}

The remaining Galois inputs are read at the actual indices.
The full degree label gives the trivial wild inertia
required for the level type. At conjugate places the local
conjugate self-dual comparison transports the dual parameter.
Inverting the eigenvalues at a transported tame generator
preserves separation. In rank one the split Frobenius values identify the row character with the reciprocity character. Tameness makes its restriction to the principal units trivial:
\[
 \chi_w|_{1+\mathfrak p_w}=1 .
\] The marked signature comparison and the
toric datum join these properties in the cohomological and
toric hypotheses of the export. Here the original rank and
the two type primes satisfy \(r\ge100\) and \(g_k\ge2r\).

The toric denominator uses
integral normalized representatives at the tested prime and
the finite norm calculation of
Section~\ref{subsec:fitted-square-ratios}. The same class-zeta
bound applies to each normalized class term and each class
sum twisted by a ring-class character. The added level consists
of split places prime to the base conductor and away from the
tested characteristic. This gives the threshold at the
given index before choosing the added level. The required roots of unity
come from the joint coefficient choice for the two type
primes. Thus the signed conditions use the row, series and
character identifications of the averaged bases.
\label{cor:automorphic-realization-53}
\end{proof}

\subsubsection*{Coefficient fields and character data for the second twist}

The averaged theta bases already realize the assigned characters for the initial rational rows. We now compare the coefficient fields and character data required by the second twist, using those rows and their period values.
For an index \(i\) over a prime \(\ell_i\), a lift
\(1\le j\le m\) and a sign \(b\), the relevant theta vectors
are the class averages of the signed theta series in the
type-selected Hecke space. The rank used to choose the finite
set of vectors is the rank of this span. The twist theorem uses
the full-rank assertion to identify the span with the rational
reading of the type-selected subspace.

\begin{proposition}[Finite basis for the retained retwisting row]\label{prop:retwisting-row-basis}
For an index \(i\) over \(\ell_i\), a lift \(1\le j\le m\) and
a sign \(b\) in Proposition~\ref{prop:signed-theta-data}, the class averages
of the signed theta series generate a finite free lattice
\(L_{i,j,b}\) over \(O_i\). One can select class averages \(a_1,\ldots,a_d\) that give an integral basis:
\[
 L_{i,j,b}=\bigoplus_{k=1}^{d} O_i a_k.
\] For the nonbase tuple, the tied
range and faithfulness hypotheses identify its rational
span with the good subspace used by the retwisting theorem.
\end{proposition}

\begin{proof}
\proofstep{Use the integral comparison before extending the fraction field.}
The
finite-class construction gives the integral comparison of the
class sums with the Hecke coefficient module. The roots of
unity of orders \(g_0(g_0-1)\) and \(g_1(g_1-1)\) are in the
embedded field \(K_0\) for the coefficient values at the type
primes. Write \(O_i\) for the local coefficient DVR. Take the
integral comparison at trivial twist and empty auxiliary level.
The \(O_i\)-span of the raw series injects via \(u\mapsto 1\otimes u\)
into the image of the finitely generated multiplicity lattice under
the theta map over the trivial-twist extension ring.

For the finite faithfully flat extension ring \(B/O_i\), the descent map is injective:
\[
 U\hookrightarrow B\otimes_{O_i}U,\qquad u\longmapsto1\otimes u .
\]
The multiplicity image in which the raw series are read is finitely generated over \(O_i\). Since \(O_i\) is
Noetherian, the raw span is finitely generated. Its image under
the integral averaging operator is consequently a finitely generated
torsion-free \(O_i\)-module \(L_{i,j,b}\).

\proofstep{Choose a basis among the reductions, then lift by Nakayama.}
Write \(L=L_{i,j,b}\), \(k_i=O_i/\mathfrak m_i\) and
\(d=\dim_{k_i} L/\mathfrak m_iL\). Choose \(d\) averaged vectors
whose reductions form a basis. Their column map satisfies
\[
 O_i^d\xrightarrow{\,C\,}L,\qquad
 k_i^d\xrightarrow[\sim]{\,\bar C\,}L/\mathfrak m_iL .
\]
Nakayama gives surjectivity of \(C\). Since \(L\) is finite and torsion-free over the DVR, it is free of
rank \(d\). The surjection \(C\) fits into the exact sequence
\[
       0\longrightarrow \ker C\longrightarrow O_i^d
              \xrightarrow{C} L\longrightarrow0 .
\]
If \(K\) denotes the kernel over the coefficient DVR \(O\), its rank-zero and torsion-free conditions together give
\[
 K\hookrightarrow K\otimes_O\operatorname{Frac}(O)=0 .
\]
This also gives the empty basis when \(d=0\).

\proofstep{Identify the rational span using the actual signed theta range.}
The second argument determines the rational rank of \(L_{i,j,b}\)
inside the signed Hecke space. The real-side map from each signed
finite model is injective under the rational row comparison.
The nonzero scalar multiplying its theta reading and the chosen
real vector give the same coefficient reading on the rational
form.

At a nonbase index, let \(V_{\rm good}\) denote the good complex space at the indicated base reading. The range hypothesis and the finite class sum give 
\[
 V_{\rm good}=\operatorname{span}_{\mathbb C}(\text{projected lifts})\subseteq\operatorname{span}_{\mathbb C}(\text{averaged series}).
\]
 For the base reading the projected image is the good row by Proposition~\ref{prop:span-original-theta} in its stated stable range. 

At a nonbase reading the marked comparison must identify the projected lifts and their signed class sums on the same space. On a tied good eigenspace \(V\), with lower summand \(U\) and lift \(T:U\to V\), Proposition~\ref{prop:downward-on-good-space} gives
\[
 V\ominus T(U)=\ker(T^*|_V)=0
\]
in the rank range of Lemma~\ref{hyp:tied-adjoint-level}, using \(T^*V\subseteq U\). Thus no choice of a basis of the integral lattice by itself provides the rational spanning assertion. 

Over the rational form giving
the complex reading, the dimension of this span is the largest
size of a nonzero minor of a coefficient matrix. For the coefficient matrix \(C\) and local embedding \(\tau\), apply \(\tau\) to each minor:
\[
 \det\tau(C_{I,J})=\tau(\det C_{I,J})\quad\Longrightarrow\quad\operatorname{rank}\tau(C)=\operatorname{rank}C .
\] This proves full rank
of \(L_{i,j,b}\). The integral index in the Hecke lattice
continues to enter the twist estimates.
\label{cor:automorphic-realization-54}

Use these bases in the all-coprime-powers choice of
Section~\ref{subsec:retained-twist}. The type comparison
above supplies the universal coefficient membership. Use the
clearing powers and full Hecke preservation with the Galois
and toric hypotheses on this same initial system. For a
nonzero index \(i\) served by a nonbase tuple, the output is
\(N_0(i)\), a level \(Q_i\) and a character with values in
\(\mu_{\ell_i^{k_i}}\), \(k_i\ge1\). The quantifier on the signed
character and valuation tests is
\[
 \exists\, Q_i,k_i,\varphi_i\quad
 \forall a\in\mathbb Z,\quad (a,\ell_i)=1
       \ \Longrightarrow\ \varphi_i^a\ \text{satisfies both tests}.
\]
Before each choice put all host primes over the exclusions into
the avoidance set, including the two deletion primes, the sequence
primes, the residual test prime and the original level primes.

For successive indices also exclude all rational primes used
earlier. The auxiliary places have degree one; this gives disjoint
levels even after taking their rational norms.
To choose the index at a coefficient place, group the contracted
primes of \(K_0\) into orbits under conjugation. The marked
two-place covering provides the original index and signed
valuation on the fixed coefficient subfield at each place. Let \(\mathcal O\) be an orbit of contracted primes meeting the nonbase support. Choose one supported place \(r\) over the orbit, then use the single index
\[
 i(w)=i_{\rm orig}(r)\qquad(w\text{ supported},\ w|_{K_0}\in\mathcal O).
\] The
export gives the full rank and the basis on that very index,
since \(r\) is assigned a nonbase tuple.

At places contracting
to the same prime as \(r\), its original side recovers the same
valuation tests on \(\mathcal O_{K_0}\). At the conjugate prime
use the conjugate side. Its embedding is the first embedding
composed with conjugation. If the contracted prime is fixed,
the representative's side already gives the valuation tests.
At each place off the support keep its original index and
signed reading, with trivial auxiliary character. Thus there
is a reading at every place. The index is constant under
conjugation and on each contraction fibre on the support.
\end{proof}

\begin{auxlemma}[Transport of the reference coefficients]\label{lem:reference-field-transport}
Let \(K_s\subset K_0\subset K\) be the three coefficient fields of
the reference calculation, embedded by \(\iota:K\hookrightarrow\mathbb C\).
Let \(\Xi\) be its finite set of selected Hecke characters. Suppose a
finite set \(S(\Xi)\subset\overline{\mathbb Q}\subset\mathbb C\) of
additional character values, roots of unity and ray-field generators
is required. There is a finite extension \(K'/K\), with embedding
\(\iota'\) extending \(\iota\), containing these prescribed complex
values. The embedded copies of \(K_s,K_0\) may be kept equal. If
conjugation stability is required, assume that \(\iota K\) is stable
under complex conjugation. In this case a possible image is
\[
    \iota'K'=\iota K\bigl(S(\Xi),\overline{S(\Xi)}\bigr)
                           \ \subset\mathbb C .
\]

Let \(\mathcal H\) be the Hecke algebra away from the excluded
places, let \(V\) be the reference space over \(K\), and give it and
its two marked Hecke actions their original embeddings. If
\(e\in\mathcal H\otimes K\) is either prescribed reference idempotent
and \(\lambda\) a character acting on a reference vector, their copies
on \(V\otimes_K K'\) satisfy
\[
  (eV)\otimes_K K'=(e\otimes1)(V\otimes_KK'),\qquad
  \lambda'(T)=\iota'^{-1}\iota(\lambda(T))
                  \quad(T\in\mathcal H).
\]
These identities apply to the reference summands carried to the new
construction.
\end{auxlemma}
\begin{proof}
Inside \(\mathbb C\) take the compositum of \(\iota K\) and the fields
generated by the finitely many specified elements, and also by their
conjugates when required. Give the two smaller fields their original
embeddings. The required inclusion is the unique one satisfying
\[
    j=(\iota')^{-1}\circ\iota:K\longrightarrow K'.
\]
Tensor the direct sum \(V=eV\oplus(1-e)V\) with \(K'\). On a
reference eigenvector \(v\) one has
\((T\otimes1)(v\otimes1)=(v\otimes1)j(\lambda(T))\).
A reference evaluation expressed in a fixed cusp basis has the same
complex coordinates after replacing each coefficient \(a\) by \(j(a)\).
When \(w\) lies above \(v\), contraction of prime ideals gives
\[
    j^{-1}(\mathfrak p_w)=\mathfrak p_v,\qquad
    \operatorname{ord}_w(j(a))=e(w/v)\operatorname{ord}_v(a).
\]
Consequently an original vanishing condition modulo \(v\), or a
nonzero residue, pulls back to the corresponding condition modulo
every \(w\) above it. For a marked set specified by these original
conditions, the resulting set of places is the entire inverse image.
\end{proof}

\begin{auxlemma}[The reference datum after adjoining the ray characters]
\label{lem:second-core-reference}
Fix the Hermitian spaces, lattices, levels at the word primes, split
halves and deletion characters at one translate as in
Definition~\ref{hyp:initial-adelic-frame}. Suppose the preliminary
realization at this translate has been constructed. Write
\[
 K_s\subset K_0\subset K\subset\mathbb C,\qquad
 \iota:K\hookrightarrow K'\subset\mathbb C
\]
for the distinguished fields and the finite enlargement of
Proposition~\ref{prop:second-specialization-row}. Keep the images
of \(K_s,K_0\) fixed in \(\mathbb C\). The geometric and adelic
data of the reference tuple admit a presentation over this
enlargement with the original small field, marked points, class
representatives, oscillator, measures and analytic theta maps.

For a coefficient-valued reference class form \(\phi\), denote the
resulting form by \(\phi^\iota\). For its Hecke operators away
from the levels and conductors one has
\[
 \phi^\iota(h)=\iota_*\phi(h),\qquad
 T\phi^\iota=(T\phi)^\iota .
\]
Here \(\iota_*\) applies the inclusion to the coefficients of the
value polynomial. These assertions hold on the reference row and
on each copy of a common test at its own level.
\end{auxlemma}
\begin{proof}
\proofstep{Retain the geometric model.}
The PEL equations, integral structures, marked base point, component
maps and the correspondence with the analytic quotient are already
defined in the preliminary presentation. If \(F\subset K_0\)
denotes its embedded geometric coefficient field, use the map
\[
                 F\longrightarrow K\xrightarrow{\iota}K'.
\]
Both embeddings of \(F\) into \(\mathbb C\) agree. Keep the
model over \(F\), its ring of integers, its polarization, frames,
and all the maps at its original levels. Its presentation over
\(K'\) is obtained by pullback. The Hermitian groups and their
adelic quotients depend on the original host and Hermitian
spaces. We can use the original fundamental sets and Haar measures.
Thus both their orbit masses and the Petersson integrals on the
old complex forms have the original values.

\proofstep{Identify the reference level and the copied forms.}
Let \(S_{\rm type}\) be the prescribed set of type places and
\(Q(\mu)\) the conductor support used for a tuple \(\mu\).
For the reference tuple \(\mu_0\),
\[
              Q(\mu_0)=\varnothing,\qquad
              S_{\rm type}\cup Q(\mu_0)=S_{\rm type}.
\]
At these places a pair consisting of a level subgroup \(J\) and
its character \(\lambda\) becomes \((J,\iota\lambda)\).
For each definite class represented by \(h\), apply \(\iota\)
to the coefficients of \(\phi(h)\). The stabilizer and
right-\(J\) equations commute with \(\iota_*\), since the
weight action and the level character extend coefficientwise.
Thus these polynomials define the form \(\phi^\iota\). On
evaluation in \(\mathbb C\) at the distinguished embeddings,
\[
              (\iota_*\phi(h))_{\mathbb C}=\phi(h)_{\mathbb C}.
\]
For any fixed integral class lattice given by these finite
stabilizer and level equations in the same value lattices, scalar
extension to a finite extension of discrete valuation rings is
flat. It commutes with their common kernel. This observation
applies to the copied equations and their coefficient lattices.

\proofstep{Transport the operators.}
With the same oscillator vector and integration domain, the
reference theta operator is given by the same integral on the
same smooth inputs. The normalized algebraic reading agrees by
Lemmas~\ref{lem:marked-theta-fibre-comparison}
and~\ref{lem:marked-vacuum-scalar}. At a good prime a Hecke
double coset has the same right-coset representatives on the
copied level. If \(A_d\) is the coefficient action in the
chosen convention, then at each representative \(h\)
\[
 (T\phi^\iota)(h)
   =\sum_d \iota_*(A_d\phi(hd))
   =\iota_*\!\left((T\phi)(h)\right).
\]
In particular a copied nonzero eigenvector of eigenvalue \(a_T\)
has eigenvalue \(\iota(a_T)\). If a supported tuple requires a
finer level, use its forgetful map to the reference level for
the pullback of the old form. At a good prime the local factors
of the two level subgroups coincide, so the same computation
intertwines pullback and \(T\).

\proofstep{Insert a moving character at its actual level.}
For a carried tuple \(\mu\), the prescribed local level uses
\(S_{\rm type}\cup Q(\mu)\), its split halves, and the finite
level character specified there. Realize a level character on
the cover by its kernel, where it acts by deck transformations.
For an allowed fixed ray test with universal polynomial
\(\sum_{\delta\in\Delta} a_\delta[\delta]\), choose the same
test on this cover and pull back its value frames. Then for any
allowed character \(\varphi\), with values read in a common
extension,
\[
       \sum_{\delta\in\Delta}\iota_*(a_\delta)
                   (\iota\varphi)(\delta)
        =\iota_*\!\left(\sum_{\delta\in\Delta}
                         a_\delta\varphi(\delta)\right).
\]

Here \(\iota\) also denotes the chosen extension to the field
of the character values. Using a different primitive-root
identification at a place above it replaces \(\varphi\) by
\(\varphi^a\) for one \(a\) prime to its order.
Proposition~\ref{prop:retwisting-row-basis} supplies the
selected tests for all such powers. The equality compares the
same fixed test and character in two coefficient presentations.
The reference eigenvalue identity concerns the copied
eigenvector. A newly twisted eigenvector is tested with its
own finite character and level.
\end{proof}

The construction therefore starts with one preliminary realization.
It exports the chosen finite indices, ray supports, characters and
their all-powers tests. After adjoining their values the reference
data can be carried as in Lemma~\ref{lem:second-core-reference}.
At a place \(w\) of the larger field the index is obtained by
contraction to the preceding field. The prescribed character
and all-powers comparison then apply at \(w\). The final row
must realize the new local types, the corresponding integral
and real comparisons, and the estimates at its actual levels.
Once it has been so constructed, write \(D'\) for that final
member, with \(M_{w,D'}\) its column lattices.
Both division estimates use the lattices of \(D'\); the
preliminary member supplies the selection of characters. This
order uses no comparison between the division errors of two
different members.

\begin{auxlemma}[A single period on the copied rational row]
\label{lem:copied-reference-period}
Fix a served coefficient index \(i\), a lift \(1\le j\le m\)
and a sign \(b\). Use the reference realization of
Lemma~\ref{lem:second-core-reference} and its fixed extremal
vector \(v_0\). Suppose the first signed reading on its rational
row is given by a period \(\nu_{ijb}\ne0\). On the copy one may
take the very same \(\nu_{ijb}\), including for the common ray
tests. Under the fixed-field identification the resulting
reading at \(K'\) restricts to the old reading.
Writing \(R,R'\) for the two signed readings and \(j\) for transport of coefficients, this gives on every common test
\[
                         R'(jz)(g)=R(z)(g).
\]
\end{auxlemma}
\begin{proof}
Write \(\sigma\) and \(\sigma'\) for the signed coefficient
embeddings into the common complex realization, with
\(\sigma'\iota=\sigma\). On a pure tensor \(z\otimes f\) and a
marked representative \(g\), the normalized value on the copy is
\[
   \frac{\sigma'(\iota z)}{\nu_{ijb}}\,
            \operatorname{ev}_{v_0}(f^\iota)(g)
       = \frac{\sigma(z)}{\nu_{ijb}}\,
            \operatorname{ev}_{v_0}(f)(g).
\]

The evaluation vectors and the analytic kernels are the copied
ones. The identity on tensors gives equality on the original
rational span and determines its scalar extension. To compare a
common test supported at \(Q\), use the local section of each
ray label, the original Fourier weights, and the determinant
translation at \(g\). For its coordinate \(a\in\Delta_Q\), the
translated local section is indexed by
\[
                      a-d_Q(g),\qquad d_Q(g)\in\Delta_Q .
\]
This label, the input polynomial and the placed finite
representative coincide in the copied formulas. The stabilizer
and Haar factor at each matching definite source class coincide
as well. Their weighted sums give the same value for that
coordinate. Finite specialization and the normalized value
therefore commute at every eligible common test in this copy.
\end{proof}

\begin{auxlemma}[The character line at an added split prime]\label{lem:new-ray-character-line}
At an added split prime \(v\in Q\) of a retained tuple, let
\(F_v\) have ring of integers \(A\) and residue field \(k\).
Write \(\xi:k^\times\to\mathbb C^\times\) for the nontrivial tame
restriction of the chosen ray character, in the selected split half.
Use a row of rank \(r+d\), where \(d=jt\), and the inverse image \(J\)
in \(K=\mathrm{GL}_{r+d}(A)\) of the upper \((r,d)\) block parabolic.
On the conjugate automorphic input the suffix is twisted by \(\xi\).
After putting the determinant multiplier of the Weil action into the
group coefficient, the normalized induced representation and the
character to be tested there are
\[
 \pi'=\operatorname{Ind}_{P_{r,d}}^{\mathrm{GL}_{r+d}(F_v)}
       \big((\tau_0\otimes\widetilde\xi^{-1}\circ\det)
                           \boxtimes\sigma_0\big),
 \qquad \eta(j)=\xi^{-1}(\det\bar j_{00}).
\]

Here \(\widetilde\xi\) extends \(\xi\) to \(F_v^\times\), with its
unramified part absorbed in the spherical factors \(\tau_0,\sigma_0\).
The representation on the rank-\(r\) block is spherical before the
displayed twist. The rank-\(d\) factor is the spherical representation
of the preceding \(j\) determinant strings. Then
\[
               \dim_{\mathbb C}\operatorname{Hom}_{J}(\eta,\pi')=1.
\]
The formula includes \(d=0\) and every power of the retained character
prime to its order.
\end{auxlemma}
\begin{proof}
Write \(X=\left(\begin{smallmatrix}A_0&B_0\\C_0&D_0\end{smallmatrix}\right)\)
in the split oscillator model, with \(A_0\) square of size \(r\).
On integral matrices take the test
\[
 f(X)={\bf1}_{\{A_0\in\mathrm{GL}_r(A),\,C_0\equiv0\bmod\varpi\}}
                   \xi(\det\bar A_0).
\]
Right multiplication by \(j\in J\) multiplies this test by
\(\xi(\det\bar j_{00})\). The scalar of the oscillator action on units
is \(\xi(\det\bar j)^{-1}\). Their product is
\(\xi(\det\bar j_{11})^{-1}\), the suffix character in the level.
After conjugating the automorphic coefficient and moving the
oscillator scalar to that coefficient the character is
\[
     \xi(\det\bar j_{11})\xi(\det\bar j)^{-1}=\eta(j).
\]
The same computation with left inverse multiplication gives the
inverse level law on the other leg.

Take unit spherical vectors in \(\tau_0,\sigma_0\), with product \(u\).
The normalized induction on \(P_{r,d}\cap J\) acts on \(u\) by
\(\eta\), since the modulus character is one there. Thus
\[
 v(pj)=\delta_{P_{r,d}}(p)^{1/2}
             (\tau\boxtimes\sigma_0)(p)\eta(j)u,\qquad
 \tau=\tau_0\otimes(\widetilde\xi^{-1}\circ\det),
\]
with value zero off \(P_{r,d}J\), is well defined and \(v(1)=u\).

To bound the dimension, take unramified principal-series
presentations surjecting onto the two spherical factors and induce
in stages. Taking a \(J\)-character is exact by compact averaging.
The cells \(B\backslash \mathrm{GL}_{r+d}(F_v)/J\), for the upper Borel
\(B\), are indexed by \(r\)-element subsets of the diagonal positions.
On diagonal units the characters of the principal series are
\[
        \underbrace{\xi^{-1},\ldots,\xi^{-1}}_{r},
        \underbrace{1,\ldots,1}_{d}.
\]
Only the subset of the first \(r\) positions can match \(\eta\),
because \(\xi\ne1\). On that cell a section is determined by a single
scalar. This also proves the assertion when the suffix is empty.
A power prime to the order leaves the inertia character nontrivial.
\end{proof}

\begin{auxlemma}[Occurrence on the signed row after specializing the ray character]\label{lem:new-ray-type-occurrence}
Use the signed row of Proposition~\ref{prop:retwisting-row-basis} and the
coefficient tests of Lemma~\ref{hyp:simultaneous-retained-twist}.
Fix a served index, a lift \(1\le j\le m\) and either side. Let
\(\varphi\) be its retained character and \(a\) prime to \(\ell\).
Let \(V\) be the good rational span in that row. On the prescribed
finite-class realization at the tuple \(\varphi^a\), denote by \(X_a\)
the character space of source forms and by \(L_a:X_a\to V\) the theta
class sum followed by the good projection, extending scalars as needed.

Then \(L_a\) is onto. In its complex Petersson realization, for
each joint bottom eigencharacter \(\lambda\) of the spherical operators
outside the exceptional places where the attached local theta
transfer is specified,
\[
 L_a^*(V[\lambda])\ \subseteq\ X_a[\lambda_a],\qquad
 \dim X_a[\lambda_a]\ \geq\ \dim V[\lambda],
\]
where \(\lambda_a\) is obtained by that transfer. The assertion also
holds after imposing the common finite type idempotent intertwined by
the two theta actions. Thus every occurring bottom eigensystem has a
nonzero vector on the prescribed upper character space.
\end{auxlemma}
\begin{proof}
Put \(n=\dim V\). Proposition~\ref{prop:row-pair-removal} supplies a basis
of the row lattice and a common clearing scalar \(x=\ell^c\).
For the \(n\) selected source class generators \(z_i(a)\) at this
tuple the specialized matrix in that basis satisfies
\[
       C_Q(u^a-1)=
         [\,xL_a z_1(a)\ \cdots\ xL_a z_n(a)\,],\qquad
       \det C_Q(u^a-1)=f_M(u^a-1)\ne0 .
\]
Use the inverse tuple on the side defined with the inverse character.
The nonvanishing for both sides and all these powers is the simultaneous
choice of Proposition~\ref{prop:one-retained-character}.
Lemma~\ref{lem:universal-theta-type} gives the signed type on each
specialized column. For the source character law, translate the
class representative in its finite sum by a level element: the
change in its determinant index multiplies its coefficient by the
prescribed character, and its polynomial is transported by the
coefficient representation. These are precisely the transformation
laws defining \(X_a\). The columns span \(V\), including the empty
case \(n=0\).

For a tied arrow let \(h\) be an upper spherical operator and
\(\mathcal T(h)\) its lower transferred operator. The good projection
is orthogonal and commutes with this algebra. The downward
intertwining law on the kernel and the Petersson adjoint give
\[
 L_a h^*=\mathcal T(h^*)L_a,\qquad
 h L_a^*=L_a^*\mathcal T(h).
\]

Indeed the local transfer respects the involution. In normalized
Satake coordinates the involution conjugates and inverts the
parameters. The splitting characters are unitary and the appended
string is fixed under this operation, by \(j'\leftrightarrow t-1-j'\).
At an unramified place the local transfer on normalized parameters of
one upward arrow, with its oriented splitting characters
\(\chi_S,\chi_T\), is the multiset map
\[
 B_w\longmapsto
     (\chi_S/\chi_T)(\varpi_w)B_w
       \ \sqcup\
     \{\chi_S(\varpi_w)({\rm N}w)^{\,j'-(t-1)/2}:0\le j'<t\}.
\]
It is injective by multiset cancellation. At inert hyperspecial
places this formula reads the base-change parameter over \(E_w\);
the Hecke operator is the relative spherical operator over \(F_v\).

Iterate the formula along the tied arrows. At the residue-type
places simultaneous translation of the two copies of the first
block in the Weil kernel intertwines the common finite projector.
The spherical operators commute with this projector by disjoint
support. Equivalently the latter statements can be imposed first
in forming \(X_a,V\).

If \(v\in V[\lambda]\), the intertwining identity gives
\(hL_a^*v=\lambda(\mathcal T(h))L_a^*v\). Surjectivity gives
\[
          \ker L_a^*=(\operatorname{im}L_a)^\perp=0 .
\]
This proves the asserted lower bound with the same finite test
for each vector of the bottom eigenspace. When several full labels
have the same restriction to the tested operators, the argument
applies to their direct sum. It supplies a bound on the whole
restricted eigenspace. It gives no numerical multiplicity for a
particular representation inside that eigenspace without a projector
separating that representation. In the comparison of the two signed
coefficient models alone, the tested operators have split primes
\(P\notin S_0(i)\) and degrees \(1\le k\le r\); the signed
normalization multiplies their degree-\(k\) eigenvalue by
\(c_{i,b,P}^{\,k}\).
\end{proof}

\begin{proposition}[Preservation of the row under the second specialization]\label{prop:second-specialization-row}
In the two-pass construction, retain the embedded fields
\(K_s,K_0\), the selected character data and the assigned
word levels. Extension of the coefficient field and the second
specialization preserve the transported row and the tuple
character at each marked place and the grid tests at every
admitted depth. If \(w\) lies over an original marked place \(v\), its tuple character is inherited from \(v\). In particular
\[
 \{\chi_{v(w)}:w\text{ over the marked places}\}=\{\chi_v:v\text{ marked}\},
\]
and the same tuple list includes the distinguished characters \(1,\delta,\delta'\).
\end{proposition}

\begin{proof}
Carry the averaged bases of Proposition~\ref{prop:retwisting-row-basis}
into this comparison, with their level type from
Lemma~\ref{lem:universal-theta-type}.

Keep the embedded fields \(K_s,K_0\) and the character frame.
For each selected index choose a complex identification of its
cyclic \(\ell_i^{k_i}\)-th roots and use class field theory for
its ray quotient to form the complex finite-order Hecke
character. At a prime outside \(Q_i\), let \(\gamma\) be
the image of its Frobenius in the ray quotient, and let \(j_i\)
be the chosen embedding of its character values. In the given
Artin convention,
\[
 \chi_i(\mathfrak p)=j_i(\varphi_i(\gamma))^{-1},\qquad
 \chi_i^c=\chi_i^{-1}.
\]
The second equality follows from conjugation on the ray character.
The kernel is open because the quotient is finite and continuous.
Injectivity of \(j_i\) preserves the nontrivial inertia at each
place of \(Q_i\); the character is unramified elsewhere, and
\(Q_i\) avoids the primes over \(\ell_i\).

For every original
coefficient place, use the character for its assigned index
on the prescribed nonbase support, and use the trivial
character on the complement. A new place of any eventual
coefficient extension then inherits the complex character,
level and root order from the place below it. If \(w\) is an upper prime and \(K\) the lower field in the contraction, compatibility of the field inclusion with conjugation gives
\[
 (\bar w)\cap K=\overline{w\cap K}.
\]
Under contraction,
places with the same prime of \(K_0\) have contractions with the
same prime of \(K_0\). The saturation of the initial support
and the constancy of the index on this support show that
these assignments agree there, including on the complement.
The conjugate pair of local readings inherits the same index.
The rational prime set of the marked data is unchanged.

Let \(S\) be the set of the three distinguished characters and the
deletion-character products from the assigned places. Contraction
\(\pi\) is onto, so for the assignment \(v\mapsto\chi_v\) one has
\[
 \{\chi_{\pi(w)}:w\text{ in the extension}\}
                 =\{\chi_v:v\text{ in the original field}\}.
\]
Thus \(\Xi'=S\) is known before adjoining roots. Adjoin to \(K\)
the roots of orders \(\ell_i^{k_i}\), the values required by
\(\Xi'\), and \(S_2(K_s,K_0,\Xi')\). Only finitely many places
lie over the prescribed rational primes; at the other places
the assigned root order is one. Choose finitely many generators
of \(K/\mathbb Q\).

Let \(A\) consist of these generators, the requested entries and
roots, and \(i\). The required enlargement is
\[
 K'=\operatorname{Spl}_{\mathbb Q}
          \left(\prod_{a\in A}\operatorname{minpol}_{\mathbb Q}(a)\right)
           \subset\mathbb C .
\]
It is finite and normal, and complex conjugation permutes the roots
of every factor. Hence it acts on \(K'\).
The inclusions retain both coefficient embeddings: if \(i'\) denotes the inclusion of the enlarged field, then
\[
 i'|_{K_s}=i_s,\qquad i'|_{K_0}=i_0,
\]
where \(i_s,i_0\) are the original maps into \(\mathbb C\). If an additional field requirement concerns the
reconstruction of the adelic datum or finite coefficient matrices,
its finite list must be included in this step, with its
additional avoidance set imposed before selecting the ray
characters. The interval caller itself uses empty additional
avoidance sets. \label{cor:automorphic-realization-55}

There is a root identification to check when comparing the
old character with the rebuilt one at \(w\) above \(v\).
Write \(i=i(v)\), \(n=\ell_i^{k_i}\) and use the transported ray
quotient with both marked maps on \(K_0\). In the common local
fraction field let \(j_1,j_2\) read the old and new roots. Both embeddings account for all \(n\) distinct roots of \(X^n-1\) in the common local fraction field \(L\). Thus
\[
 j_1(\mu_n)=\mu_n(L)=j_2(\mu_n).
\] For one
primitive \(\zeta\) there is therefore a unit \(a\bmod n\) with
\[
 j_2(\zeta)=j_1(\zeta)^a,\qquad
 j_2(\zeta^k)=j_1(\zeta^k)^a\quad(k\in\mathbb Z).
\]
The same exponent works on the whole ray-character image.

Both the new arithmetic character and the old ray character
read the same complex Hecke character at Frobenius outside
\(Q_i\) and the residue characteristic. For an arbitrary element
\(\gamma\) of the finite ray quotient, surjectivity of the Artin
map gives a preimage in the Galois group of the ray extension.
The abelian Chebotarev theorem realizes \(\gamma\) by a
Frobenius at a place outside those exclusions.

At that place
each reading uses the inverse of the complex Hecke value, so
the common exponent gives
\[
             \varphi^{\mathrm{new}}_{w}(\gamma)
                         =\varphi_i(\gamma)^a
\]
after identifying the common coefficient rings. Equality in
the fraction field gives equality in the root ring by
injectivity. Hence the rebuilt tuple of index, level, root order
and raw character is the original tuple with its character
raised to a power prime to \(\ell_i\). The exponent may vary
with \(w\), and the all-powers twist bound covers each new place.
For the base support use the trivial auxiliary character.
\par
For the scalar extension of a fixed coefficient lattice, the
length comparisons use normalized valuations. If a finite
length module over the first local integer ring has elementary
divisors \(\pi_v^{a_1},\ldots,\pi_v^{a_h}\), their orders after
extension are \(e(w/v)a_1,\ldots,e(w/v)a_h\). Transitivity of ramification indices gives the equality of
normalized lengths explicitly:
\[
 \frac{\sum_j e(w/v)a_j}{e(w/\ell_i)}
       =\frac{\sum_j a_j}{e(v/\ell_i)} .
\] The same calculation applies to the
valuation of a scalar. Applying this observation to an
automorphic comparison requires identifying the actual
transported lattice and period on both sides.
\label{cor:automorphic-realization-56}

The rebuilt adelic datum uses the same column frame,
choice of split places and character coordinates and the
extended finite readings. Its bare row is built on the
extended fine datum. The same-Galois construction and the
rational and cuspidal readings give the intermediate state
there. The first pass fixes the moduli point, its signed integral series and its eigensystem. For each good Hecke operator \(T\) at a split place away from both levels, write \(H_T\) for the finite level correspondence there. For a form \(a\) the test on a marked point \(z\) has the double-coset description
\[
             (Ta)(z)=\sum_{h\in H_T}a(zh)
\]
with the same normalization (and coefficient action) on both levels. Thus pullback of the values along the marked correspondence intertwines the two Hecke operators. If the marked value tests separate the form, equality of these evaluations gives equality of the carried eigensystems. 

Lemma~\ref{lem:placed-row-theta} gives the marked points. The tests and their integral evaluations on the reference space are transported as in Lemma~\ref{lem:constructed-signed-coefficient-row}. At the chosen new type, Lemma~\ref{lem:new-ray-type-occurrence} supplies the eigencharacter on the tested algebra. The two bare Galois
frames agree by transitivity, through the comparison
with the bare source in each construction.
\label{cor:automorphic-realization-57}

Write \(i',i\) for the two fixed-field embeddings. Equality of their images defines the identification by
\[
 K'_0\xrightarrow{\ i^{-1}\circ i'\ }K_0 .
\] This identification respects
conjugation and integer rings. For the place map \(\pi\) and an
integer \(a\in\mathcal O_{K_0}\) the local comparison reads
\[
 a\in\mathfrak p_{\pi(v')}\ \Longleftrightarrow\
   a\in\mathfrak p_{v'}\cap\mathcal O_{K_0},\qquad
 S_{\rm new}=\pi^{-1}(S_{\rm old}) .
\]
If \(S_i\) denotes the old support at index \(i\), the transported
support is \(\pi^{-1}(S_i)\). Since contraction is surjective,
\[
                \pi^{-1}(S_i)\ne\varnothing
                     \quad\Longleftrightarrow\quad S_i\ne\varnothing .
\] \label{cor:automorphic-realization-58}

These comparisons give a copy of the original
eigensystem in the new intermediate state. It transports
the twisting, scalar and theta data after the finite
group and trivial-tuple projector have been identified.
The series on the copy uses the comparison of the
Hecke readings after restriction at the new types.
The positive multiplicity and signed level comparisons
are made on this series and give the new second
intermediate level. They refer to the copied
eigenvalues, the same range of lifts and the copied
unipotent character. They also supply the conductor
and level data for the two signed readings of the
appropriate local twist.
\label{cor:automorphic-realization-59}

\proofstep{Keep the selected place and transport the old level tests using the same support.}
The field comparison must preserve the selected place at which the row tests are imposed.

Let \(p_*\) be the least split auxiliary prime specified for the grid and \(v_*\) its prescribed base place. This choice meets both exclusions simultaneously:
\[
 v_*\notin S_{\rm grid}\cup \bigcup_\mu\operatorname{supp}(\operatorname{cond}\chi_\mu),
\]
where \(S_{\rm grid}\) denotes the bad places of the grid.
Use the compatible
place on the enlarged frame, with the same selected
host place above it. The field values and required
roots lift through the inverse of the fixed-field
identification. Since the distinguished and deleted
character coordinates agree, the formula for each
grid character agrees. The old level test at a supported
index applies at its new presentation by the equality
of support just proved. This gives transport of the list
at every depth for which the old list is given.
\label{cor:automorphic-realization-60}

Let \(\pi\) denote contraction on coefficient places and \(i(v)\)
the exported index. At the geometric point of the carried series
the character choice is
\[
                 i_{\rm new}(w)=i(\pi w),\qquad
                 i_{\rm new}(\bar w)=i_{\rm new}(w).
\]
The first equality uses the pulled-back support and its allowed
reading; the second uses conjugation on the original assignments.
The geometric comparison carries the grid, and spectral
specialization preserves this index. Hence its signed tests use
the exported signature and basis. The first specialization supplies
these data from the original real and series readings and the
second receives them on the transported series.
\label{cor:automorphic-realization-61}

For the tuple set \(U\) on this system, injectivity of the tuple-to-character map and the description of its image give the bijection
\[
 U\xrightarrow{\ \sim\ }\Xi',\qquad \mu\longmapsto\chi_\mu .
\] Its distinguished tuples read
\(1,\delta,\delta'\), and its assigned nonbase tuples read the
products just constructed. Off the selected support the
auxiliary character is trivial. On the base tuple the two
conjugate readings have twisting order zero. The mass
formula of Section~\ref{subsec:selected-word-level} returns
the original \(d_{u,*}\) and \(D_*\) on this assembled member.
The volume-to-mass identity uses its measured row, and its
signatures and Weyl weights are the prescribed ones.

For the
stabilizer divisor use its single value across words and
tuples, supplied by the column comparison. The congruence
orders off the twist support vanish as above. These assignments exhaust the tuple range. On the twist support this follows from the tuple character formula; elsewhere its auxiliary factor is one. Since the tuple-to-character map is injective, equality of the two readings also identifies their labels:
\[
 \chi_{\mu_1}=\chi_{\mu_2}\quad\Longrightarrow\quad \mu_1=\mu_2 .
\] These identities extend the member
by its numerical conditions. Forgetting those conditions
returns exactly its constructed object. In particular the
pre-level masses are valid for the member obtained after both
passes at any permitted levels at the primes assigned to the sequences.
\label{cor:automorphic-realization-62}
\end{proof}

\paragraph{Substitution of the division lattices.}

\begin{lemma}[Simultaneous substitution of the division lattices]
At a fixed good member let \(M_u\) denote its rational ramified word columns. The divisor, arithmetic and period comparisons may be transported simultaneously for any assignment
\[
 u\longmapsto(N_u,h_u),\qquad N_u\subseteq M_u\text{ an integral submodule},\quad h_u\in\mathbb R.
\] For a compatible completion these are its local
lattices and heights.
\end{lemma}
\begin{proof}
Write \(M_u\) for the rational space of the ramified word
column. The substitution is made independently in every word column. Its entire input is
\[
 (N_u,h_u)_u,\qquad N_u\subset M_u\ \text{a coefficient-integer submodule},\quad h_u\in\mathbb R .
\] Substitute \(N_u,h_u\) in the division
coordinates of the member. Its Hermitian forms, local
comparison bases, positions, split levels and degree-zero column
are carried in the same coordinates. In particular their joint
reconstruction can use the old matrices, local coefficient
embeddings, bases and weight arrays with the new division
coordinates inserted in the original column datum. Its kernel
family on each selected pair uses the same class, series and
substitution matrices. The expression of its entries from
the chosen places and adapted bases is carried with the series.
The downward maps and their paired integral identities persist
with the same test, function, character and Haar measure.

Let \(\pi\) pull an index of the substituted flag datum back to
its original flag. At a flag \(a\) the quotient-prime set and
model degree satisfy
\[
 Q^{(N,h)}(a)=Q(\pi a),\qquad
 \deg^{(N,h)}(a)=\deg(\pi a).
\]
The marked chart, weights and model fibre are read at the same
\(\pi a\). Closure of the geometric condition under a single
\((N,h)\) substitution thus applies both to the initial flags
and to flags resulting from earlier substitutions. \label{cor:automorphic-realization-63}

For the arithmetic comparison use the two row realizations with their
rational coordinates, lattices, finite types and signatures.
In the revised norm choice retain the normalized character and
signed basis. The same-row comparison transfers its grid list.
If \(a\) identifies the earlier field and series with the present origin and \(b\) identifies that origin with the new one, the induced identification is
\[
 x\longmapsto b(a(x)).
\]

In the measured comparison use the saturation for the
degree-zero column just transported and its original scale and
evaluation. Its spectral data therefore transport on these
columns with the prescribed \(N,h\). The sharp normalization
and the type choices defining the logarithmic values transfer
on the carried row. The two finite sets in the twist test and the
same signed theta data give its transported test. The error
model is the one read from the carried Hermitian column.
For a fixed \((N,h)\) the transported column, comparison map, twist form and kernel are identified together. Thus the spectral and error estimates concern the same lattices and heights. \label{cor:automorphic-realization-64}

Now fix a compatible completion \(p\). Take the actual local
prime exponents prescribed by \(p\) and its heights \(h_u(p)\).
By the primary-component construction the submodule \(N_u\)
realizing these conditions satisfies, at each marked place \(v\),
\[
              (N_u)_v=N_{u,v}(p).
\]
Insert this joint choice in the preceding transport. Every
coordinate, including the ramified column in the adapted entry,
is the coordinate of \(p\). The transported column thus has the prescribed
localizations and heights \(h_u(p)\). The
twist and tuple laws give its local divisor and character tests.
The scale, mass and size comparisons fixed the threshold before
\(p\) was chosen. \label{cor:automorphic-realization-65}

\end{proof}

\subsection{The period operator at the selected subspace}
\label{subsec:selected-completion-period}

The localizations and heights now come from a single completed point \(q\). We must identify the operator that acts on its selected subspace. We first obtain the
reciprocity data there from rational coordinates and the
rational-pairing law of the chosen norm extension. We then
identify the period of Section~\ref{sec:selected-row-gram}
in the ramified coordinates of that same point. The measured
row, the word cosets and the evaluation translate are those
already chosen on \(q\).


\begin{lemma}[The period pairing in rational coordinates]\label{lem:completed-rational-period-reading}
Let \(D_c(Z)\) be the definite column data at an allowed subspace
of the completion, with its chosen levels, jets, tuple characters
and integral kernel. For a word \(u\) and tuple \(a\) in use, write
\(\mathcal P^{\rm col}_{u,a}\) for the complex period pairing between
its column forms and the row of that tuple. For a rational column
coordinate \(x\) and face coordinate \(F\), the embedded scalar of \(D_c(Z)\) equals
\[
 \mathcal P^{\rm col}_{u,a}
       \big(j_{N,u}(e_{\nu,u}x),\,j_{A,a}(r_{V,0}F)\big),
\]
where \(r_{V,0}\) restricts the original row at the refined level.
\end{lemma}
\begin{proof}
The rational coordinate comparison accepts any Hilbert space,
period operator and distinguished subspace on each \(D_c(Z)\).
Use \((\mathbb C,0,0)\) provisionally. Its word index is the
identity. The fitted rational column isomorphism satisfies
\[
       e_{\nu,u}\big(e_{\nu,u}^{-1}\Lambda_u\big)=\Lambda_u .
\]
For the other argument, the form recovered from \(r_{V,0}F\)
and the restriction of the form defining the face have the same
evaluation. Evaluation on restricted forms is injective, so their
embeddings in the tuple row agree. Applying the definite column
pairing gives the stated scalar, using the twisted form at a
nonreference tuple. \label{cor:automorphic-realization-66}
\end{proof}

\begin{proposition}[Rationality of a nonzero period factor]\label{prop:rational-period-factor}
In the coordinates of Lemma~\ref{lem:completed-rational-period-reading}, the reference sharp
period descends to the rational good quotient. At a used tuple write \(P_\mu\) for the period and \(B_\mu\) for the indicated quotient pairing. For the coefficient embedding \(e:H\hookrightarrow\mathbb C\),
\[
 0\ne P_\mu=c B_\mu\quad\Longrightarrow\quad c\in e(H).
\]
\end{proposition}

\begin{proof}
\proofstep{Keep the provisional periods while forming the rational quotient.}
Put \(H\) for the coefficient field and \(e:H\to\C\) for its
given embedding. For this calculation keep the provisional
periods and the original size coordinates. The new real errors
are those read on the resulting face. The original Petersson
scalar and row scalars still give its primitive, separation
and character identities, since its \(l_{6m,j}\) have their
original values. The tuple labels, sharp and bottom embeddings
and all the coordinate and lattice readings in the displayed
period identity are therefore still available.

On the reference tuple let \(T=j_A r_{V,0}\) be the sharp
coordinate map and \(P_q\) the rational matrix to the good row.
Its range is the entire good coordinate space by the primitive
comparison. On the sharp coordinates the reference row identification makes the two composites agree:
\[
 \iota T=j_{EE}e_RP_q,\qquad T=j_A r_{V,0}.
\]
The complex span of the image of \(T\) is the degree-zero
space \(V_0\) by equality of the two sharp projector ranges.
The subspace removed by the orthogonal projection is
\[
 U=V_0\cap\ker\iota=\operatorname{span}_{\mathbb C}T(\ker P_q).
\] Indeed the images of a rational row basis
under \(j_{EE}\) are complex linearly independent, so ordinary
elimination in the rational matrix \(P_q\) describes its
kernel after extension of scalars. The same computation
applied to \(T\) gives the asserted image in \(V_0\).

Choose a rational \(G\) with \(e_RP_qG=v\). Pick a basis
\(u_a=TG_a\) of \(U\) with \(G_a\in\ker P_q\). The projection
coefficients of \(TG\) satisfy the Gram system
\[
 \left\langle u_a,TG-\sum_b c_b u_b\right\rangle=0
 \quad(\text{each }a).
\]
The Gram matrix is invertible and its coefficients and the
right sides lie in \(e(H)\), a field stable under conjugation.
Hence all \(c_b\) lie in \(e(H)\). Their rational preimages
give \(G'=\sum_b e^{-1}(c_b)G_b\in\ker P_q\). Now
\[
 T(G-G')\in V_0\cap U^\perp,\qquad \iota T(G-G')=j_{EE}v .
\]
This is the representative used by the quotient pairing.
For every rational \(\phi,v\) it follows that
\[
   B_u^{\C}(j_N\phi,j_{EE}v)
    = \mathcal P^{\rm col}_{u,0}
       \big(j_N\phi,T(G-G')\big)\ \in e(H),
\]
by the rational scalar of the definite column with
\(x=e_{\nu,u}^{-1}\phi\). This holds even when the orthogonal
projection of \(TG\) is nonzero.

The division construction requires a family of forms over \(H\)
before rationality has been proved. Put
\(b_u(\phi,v)=B_u^{\C}(j_N\phi,j_{EE}v)\) on rational coordinates.
Use the single test \(b_{u'}(\phi',v')\in e(H)\) for all words and
all rational pairs. Define, simultaneously for every \(u\),
\[
 B_{k,u}(\phi,v)=
 \begin{cases}
  e^{-1}(b_u(\phi,v)),&\text{if the test holds for every pair and word},\\
  0,&\text{otherwise}.
 \end{cases}
\]
In the first case the two coordinate maps are semilinear over \(e\).

For \(a,b\in H\), complex bilinearity gives
\[
 b_u(a\phi,bv)=e(a)e(b)b_u(\phi,v).
\]
Applying \(e^{-1}\) on the image of \(H\) gives scalar bilinearity
there, and additivity follows in each argument in the same way.
The other case gives the zero forms. This defines the full family
before the rationality proof. Localization carries the
resulting conditional identity from the initial state, and
changing \(N,h\) preserves the underlying rational spaces,
maps and this identity. The calculation above satisfies its
hypothesis on the completed data. Thus for all these pairs
the displayed value is \(e(B_{k,u}(\phi,v))\).
\label{cor:automorphic-realization-67}

\proofstep{Choose a nonzero pair of rational coordinates before taking the ratio.}
A rational pairing determines the comparison of two period normalizations from any pair of coordinates on which the denominator is nonzero. The choice below avoids dividing by an individual period that could vanish.

Fix a word \(u\) and tuple \(\mu\). On \(V_0(\mu)\) suppose the period is given by the scalar \(d\) and the projected reading
\[
 v\longmapsto d\,B_u^{\C}\bigl(\iota(p_\perp v)\bigr),
\]
where \(p_\perp\) projects orthogonally onto \(V_0(\mu)\cap U(\mu)^\perp\). Suppose some period there is
nonzero. Write \(p_\perp\) for that orthogonal projection. The kernel calculation identifies the two inputs to the period:
\[
 v-p_\perp v\in U(\mu)\quad\Longrightarrow\quad \iota(v)=\iota(p_\perp v)\qquad(v\in V_0(\mu)).
\]
An unused tuple has zero row image and could not have such a
period. For a used tuple the sharp coordinate vectors project
to a spanning set in \(V_0(\mu)\), and the rational column
vectors span the complex column.

If the right-hand pairing
vanished on every pair of those vectors, complex linearity in
each variable would give vanishing throughout. Thus there
are rational coordinates \(x,F\) with nonzero denominator.
The quotient image of the row vector is \(j_{EE}v\) for a
rational \(v\) by the identification of the row space at the refined level at that tuple.
Consequently
\[
        d=e\!\left(
          \frac{D_c(Z).\operatorname{Pt}_{u,\mu}(x,F)}
               {B_{k,u}(e_{\nu,u}x,v)}\right)
\]
at the selected row subspace \(Z\) of the rational reading.
The original row chooses its scalar to agree with any such
nonzero factor lying in \(e(H)\). This proves that agreement
for every complex factor meeting the stated hypotheses. Here
the factor multiplies the quotient pairing to give the column
period. On a nonzero reading this is the reciprocal convention
to the factor multiplying the column period in
Section~\ref{subsec:selected-period}. \label{cor:automorphic-realization-68}
\end{proof}

\paragraph{The compatible reciprocity data.}

\begin{proposition}[Compatibility of the marked reciprocity data]
Let \(E\) be the original marked CM host. The spectral comparison and standard kernel at the adapted entry admit compatible Galois and automorphic embeddings with
\[
 \iota_{\rm Galois}|_E=\iota_{\rm aut}|_E .
\]
They retain the prescribed ray field, auxiliary extension and local choices.
\end{proposition}
\begin{proof}
The host level in the completed point is supported on the two
type primes, the two distinguished primes and the specified
base auxiliary prime. Its level realization follows from the
word system over this marked host. Retain the existing full and
degree-zero word forms and their local comparison maps, the
normalized character, the divisor geometry and the theta action.
Use the twist supplied by the good member in the local Galois
and norm data. Its choice agrees with the retained selection,
so its character frame and tuple characters, local types,
cover and Galois frame are the prescribed ones.

Thus, writing \(E\) for the CM host, the two comparison embeddings restrict to the same map on \(E\):
\[
 \iota_{\rm Galois}|_E=\iota_{\rm aut}|_E .
\] Artin reciprocity and the Weber
ray-field comparison on the marked imaginary quadratic field
identify the two prescribed field pieces (the ray field and
the chosen auxiliary extension) in this frame. The row
coefficient comparison just proved supplies the scalar
condition on the same rows. These identities assemble the
spectral data with the original geometry, tuple choices and
word columns. \label{cor:automorphic-realization-69}

Use the transported adapted entry to choose the standard
kernel with its original class, frame, characteristic and
support readings. In the standard coordinates the two Hermitian forms have the same
Gram matrix and the two lattices impose the same local membership
test. At each word, stage and admissible frame the identity matrix
therefore satisfies both comparisons:
\[
           I^*H_{\rm std}I=H_{\rm std},\qquad
           I\Lambda_{\rm std}=\Lambda_{\rm std}.
\] The
frame equality uses the ordinary comparison at an ordinary
slot and the conjugate comparison at a dual slot. The
original and twisted lower maps give the two row comparisons
in these coordinates. Use the compatibility of the divisor
character, the additive character of sign \(-1\), and the
common Weil action to retain the joint comparisons.

Let \(S\) contain the assigned sequence primes and all contractions
of nonreference quotient primes. For a base ideal class represented
by \(\mathfrak a\), choose \(b\in\mathfrak a^{-1}\) by approximation
with \(v_P(b)=-v_P(\mathfrak a)\) for \(P\in S\). Then
\[
 b\mathfrak a\subset\mathcal O_E,\qquad
 v_P(b\mathfrak a)=0\ (P\in S),\qquad [b\mathfrak a]=[\mathfrak a].
\]
Over the totally complex field \(E\), Artin reciprocity
identifies the given Hilbert class field with the maximal unramified
abelian extension. Its integers lie in the prescribed field by the
marked integral containment. Transport through the specified
embedding preserves the Galois property of the coefficient extension
over \(\mathbb Q\). These choices give the required realization on
the original host. \label{cor:automorphic-realization-70}
\end{proof}

Let \(Z_*\) be the subspace determined by the actual real row and
its face. Use the positive definite counting Hilbert space
\(\mathcal H_X\) of that measured comparison. Its character
level and its good projector \(G\) were fixed on \(q\).
Transport of the degree-zero word forms identifies them with the
automorphic column functions at this level. Write \(J_u\) for
the resulting bijection from the automorphic column to the
abstract full column. On the former, let \(\mathcal P_u\) be
the Poincar\'e sum using the specified coset representatives for the sequence.
Its membership in \(\mathcal H_X\) and its linearity are the
period identities for the sequence in the measured model.

Write \(e_i\) for the rational coordinate vector and
\(e_{\mu,u}\) for the complex column evaluation.
In the ramified
subspace of the full column put
\[
       b_{u,i}=e_{\mu,u}\,j_{N,u}(e_{\nu,u}e_i).
\]
Identify the word
at a column label by the inverse of the bijection between the chosen sequence indices.
Let \(\sigma\) be identity or complex conjugation according to
the comparison of the coefficient embedding with the selected
archimedean place. Use the Hilbert space with scalar action
twisted by \(\sigma\). On its underlying additive space the
operator at \(Z_*\) is
\[
   \operatorname{out}\,P_{u,*}(z)
     =\mathcal P_u\,J_u^{-1}
        \left(\sum_i\sigma(z_i)b_{u,i}\right).
\]
\begin{lemma}[The operator on the selected row]
With the \(\sigma\)-twisted target action, the underlying operator satisfies
\[
 P_{u,*}(av+bv')=\sigma(a)P_{u,*}(v)+\sigma(b)P_{u,*}(v')\qquad(a,b\in\mathbb C).
\]
Thus it is complex linear into that target. On the good row its restriction is the ramified column reading at \(Z_*\).
\end{lemma}
\begin{proof}
On the underlying Hilbert space, substitution of \(az+bz'\) in
the displayed formula gives \(\sigma(a)\operatorname{out}P_{u,*}(z)
+\sigma(b)\operatorname{out}P_{u,*}(z')\), which is the required
linearity in the twisted space. Membership in the distinguished
row is tested by \(Gv=v\) after inclusion into \(L^2\).
Consequently unfolding \(\mathcal P_u\) on these word functions
restricts that very operator to the ramified columns.
\label{cor:automorphic-realization-71}
\end{proof}

Replace the provisional period record at \(Z_*\) by this one.
Keep the fitted words and their maps \(e_{\nu,u}\). The selected
map into the row space at the refined level \(s_E\) is the measured evaluation at the
original translate and Haar scale times the positive
normalization chosen for the placement. On the sharp subspace
it satisfies \(n_P\|s_Ev\|^2=\|v\|^2\).
Read the size coordinates with the actual \(\Delta,h_0\) and
\(N_c\) from the completion, where \(\eta_{\rm sz}=mh_0\).
For the given real coefficient \(c'\), the update acts simultaneously
on the size coordinates:
\[
 \widetilde l_{6m,j}=l_{6m,j}-c'l_{6u},\qquad \widetilde l_{6u}=0 .
\]
Every subtraction uses the original \(l_{6u}\).

\paragraph{Placement on the updated face.}
The row and divisor comparisons identify the two ambient realizations,
the intermediate modules and the finite group identification
after this replacement. For each nonzero base isotypic row labelled by \(\sigma\), the occurrence condition on the real row has two parts. At each required tuple \(\mu\), and at each other infinite place \(v\), it requires
\[
 V_{\mu,\sigma}\ne0,\qquad \pi_{\sigma,v}\simeq\pi_{\sigma,v}^{\rm prescribed},
\]
where \(V_{\mu,\sigma}\) is the corresponding tuple summand. This is the multiplicity comparison
on the real row, applied with the generic series character and
the chosen type data. With these inputs the column construction
permits the new period and scale at fixed rational readings.

Apply the definite-column comparison to the resulting face.
It uses the equality of kernels and definite columns, the
fitted period map for the sequence and the original rational coordinates. It then
identifies its entire family of period spaces with the
\emph{updated} family, including \(P_{u,*}\) at its selected
subspace. The same comparison gives the actual weights, finite
character labels and compatible column characters.
\label{cor:automorphic-realization-72}

\subsection{The adelic division pair on the completed point}
\label{subsec:completed-division}

The preceding period comparison identifies the operator in the ramified coordinates. For the division argument this identification must also respect the degree-zero forms, their metrics and the refinements of level at every word.

Fix a completed member \(q\) of Definition~\ref{def:completed-member}
and its compatible reciprocity data \(R\). For each word \(u\) the
degree-zero column forms of \(R\) agree with those of \(q\). Apply this equality to the measured
saturation and its remaining spectral data. This gives the
measured saturation for \(R\) with the same group, measure,
evaluation translate and comparison of the full and good rows.
The finite part of the saturation puts each period indexed by the sequence in the
level space of its indicated character. Thus the period
construction above uses the prescribed level projector on
every word.

\begin{proposition}[The completed division pair]\label{prop:completed-division-pair}
On a completed member in the sense of Definition~\ref{def:completed-member}, the real, Hecke, and lower
comparisons identify the square-integrable and geometric
realizations of the definite columns, periods, and size
functions at all subspaces. The identified point carries
the original tuple characters and reciprocity data,
including the prescribed side at each selected prime.
In particular, writing \(q,R\) for the two identified realizations and \(F_u^0\) for the degree-zero forms of word \(u\),
\[
 F_u^0(q)=F_u^0(R)\qquad\text{for every }u.
\]
Both parameter runs are evaluated on this same completed point.
\end{proposition}

\begin{proof}
Use the two ambient identifications on \(q\), the intermediate
modules and the finite group comparison of the preceding
construction. Denote the resulting adelic row data and face by
\(x\). The adelic division pair is formed from the base and
sharp form readings of \(x\) and its word columns. Its
intersection comparison over the coefficient field is the
base change of the one on the original row. The rational word
isomorphisms and the sharp and base ambient identifications
are those in the construction of \(x\).

Let \(n_{\rm sh}\) be the sharp-level index. Write \(E\) for
sharp evaluation with the measured translate and Petersson
comparison, using the conjugate scalar action at a conjugate
reading. The placement chosen in the preceding construction
is \(s_{\rm sh}=\kappa E\) with \(\kappa>0\) and
\[
                 n_{\rm sh}\|s_{\rm sh}v\|^2=\|v\|^2
\]
on the refined row. In particular its phase is \(1\). Writing \(F_u^0\) for the degree-zero form space, the identification on the two realizations is wordwise:
\[
 F_u^0(q)=F_u^0(R)\qquad\text{for each word }u .
\]
The saturation therefore uses the same level projector, character,
evaluation translate and Haar measure in the norm identity.

Let \(e_F,\mathrm{fi}\) be the two coefficient maps in the normal
expansion. A substitution of face sizes or error changes the
arguments of the expansion. Its form and polynomial coefficients
are read in the original rational frame, so the coefficient maps
are still \(e_F,\mathrm{fi}\). The diagonal forms, positive
terms, standard Hermitian frames and their lattices are those
of \(q\); uniqueness gives the standard additive character of
sign \(-1\) on both sides. Use the bounded real comparison on
this core with its frame weights, extremal vectors, level action
and local uniformization maps. All its phase factors are the
one just computed. This identifies the intersection, norm and
period comparisons on a single pair. \label{cor:completed-division-1}

\proofstep{Identify the actual Hecke eigenvalues and tuple characters.}

On the base row write the decomposition by complete Hecke systems as
\[
 V_{\rm base}=\bigoplus_{\lambda}V_{\rm base}[\lambda].
\] Fix one
occurring eigensystem \(\lambda\) and a required tuple \(\mu\).
Multiplicity gives its nonzero constituent at \(\mu\) with
the prescribed infinite components off the distinguished place.
The generic series choice makes the product of the finite
character and level projectors nonzero on the required labels.
At each tuple in use the adelic and marked characters satisfy
\[
 \chi^{\rm adelic}_\mu=\chi^{\rm marked}_\mu=\chi^R_\mu .
\]
The first equality uses the tuple map and the second uses
compatibility with \(R\). In particular the character depends only on the tuple label. For two positions \(a,b\) the implication is
\[
 \mu(a)=\mu(b)\quad\Longrightarrow\quad \chi_{\mu(a)}=\chi_{\mu(b)} .
\] Discrete-series reciprocity at the actual
weights gives the relative parameter as the logarithm of the
positive lift parameter; intersection at the refined level gives
both rational realizations of this constituent. \label{cor:completed-division-2}

\proofstep{Identify the lower maps in both ambient realizations.}

At the base tuple use the reference lower map of the
standard columns. At a nonreference tuple the lower map uses
the integral columns of the actual twist. To see the
polynomial comparison there, let \(P\) be the highest-weight
polynomial in the span for the prescribed top weight. Take
the standard original and twisted columns and a last-stage
frame whose determinant is a unit.
Suppose the evaluation of the specified twisted vector on the
given class agrees at every argument with the function of \(P\)
in that frame, with its prescribed coordinate shift. In the standard last-stage coordinates these evaluations recover the polynomial
\[
 P_{\rm recovered}=\rho_\lambda(A)P,
\]
where \(A\) is the matrix of the last-stage frame and \(\rho_\lambda\) its polynomial action. Applying the lower
map means taking its branching projection. The two true
weights interlace.

In those weights the restriction to the
last block commutes with the projection in the sense
\[
                 \operatorname{br}\big(\rho_{\lambda}(A)P\big)
                  =\operatorname{br}_{\rm tens}(P)
\]
with the lower map expressed in its standard coordinates.
Substitution in the top formula followed by extraction in the standard twisted coordinates recovers \(\rho_\lambda(A)P\). The block branching formula, with the upper and
lower weight shifts, then gives the displayed equality and
the determinant factor of the lower coordinate is \(1\).
This comparison uses the equality of the twist datum, the
standard twisted column maps, and the standard original
columns on \(q\). The unit condition restricts the frame in this statement.

The sharp ambient map is the map obtained from the sharp
identification at \(q\) and the two forms of the tuple
reading. The splitting of the lower modules is the splitting
of the adelic pair. If \(s_1,s_2\) are the stabilizer orders at the two levels and \(C\) the common stabilizer multiple, the original class and level conditions give
\[
 s_1\mid C,\quad s_2\mid C,\qquad \operatorname{lcm}(s_1,s_2)\mid C .
\] These identities supply
both lower maps as well as the two ambient realizations on
\(x\). \label{cor:completed-division-3}

\proofstep{Identify the two sets of marked primes under the prescribed field isomorphism.}

Let \(\theta\) be the auxiliary field isomorphism from the
standard kernel to the column model. Write \(Q_{\ker}(\mu)\)
for its finite set of quotient primes. The two prime comparisons
in the construction give
\[
 \theta_*Q_{\ker}(\mu)=Q(\mu),\qquad
 \theta_*(\operatorname{im}\mathfrak P_{\ker})=
                         \operatorname{im}\mathfrak P .
\]
Here \(\mathfrak P(v)\) is the selected prime of \(E\) over
\(v\) in \(E^+\). A prime contracting under \(\theta\) to
\(P\in Q_{\ker}(\mu)\) is necessarily \(\theta_*P\). The
level comparison at a used tuple consequently puts that very
prime in the adelic quotient-prime set. For
\(\ell\in\{p_d,p_d'\}\), if \(\ell\in\mathfrak P(v)\), then
\(v\mid\ell\). The marked side law at \(v\) identifies the
contraction of \(\mathfrak P(v)\) to the imaginary quadratic
field with the \(\ell\)-prime specified by \(R\).
\label{cor:completed-division-4}

\proofstep{Substitute the common columns before comparing the periods and sizes.}
Compare periods and heights on these identified columns.


For any subspace \(U\), identify its square-integrable column
with its geometric column by the definite comparison. The
resulting period spaces on this one domain agree, including the
Poincar\'e period in the updated component. In symbols, writing
\(C,\mathcal P,H\) for the column, its period space and its size,
the successive identifications give
\[
 C_{L^2}(U)=C_{\rm geom}(U),\qquad
 \mathcal P_{L^2}(U)=\mathcal P_{\rm geom}(U),\qquad
 H_{L^2}(U)=H_{\rm geom}(U).
\]
For the last equality both sides use the record of \(q\), the
nonnegative integer values assigned to \(\Delta,\eta_{\rm sz}/m\) and the same
replacement of \(l_{6m,j},l_{6u}\) at the given real coefficient.
Transport of the bounded real datum through these identifications
preserves the weights, word placement, side and lower maps, on
the point with record \(R\). \label{cor:completed-division-5}
\end{proof}

\begin{construction}
\begin{tabular}{@{}l@{\hspace{1.2em}}l@{}}
\toprule
\textit{Datum of the completed point} & \textit{Identification used in both runs} \\
\midrule
Word and tuple character & Same rational column and level projector \\
Measured sharp evaluation \(E\) & \(s_{\rm sh}=\kappa E,\ \kappa>0\) \\
Sharp-level index \(n_{\rm sh}\) & \(n_{\rm sh}\|s_{\rm sh}v\|^2=\|v\|^2\) \\
Base and sharp forms & One adelic division pair at the chosen translate \\
\bottomrule
\end{tabular}
\caption{The comparison data carried into both parameter runs by Proposition~\ref{prop:completed-division-pair}.}\label{fig:analytic-completed-pair}
\end{construction}

\subsection{Two parameter runs on one division pair}
\label{subsec:completed-prime-runs}

Choose the point of Proposition~\ref{prop:completed-division-pair}.
The depth estimates and the exceptional-prime estimates call for two choices of numerical parameters. Their comparison uses the same row, level and period identification at this point.

\paragraph{Row normalization and the numerical point.}
The numerical point used for the runs has the face, word
lattices, period spaces and scale of the bundle. To form it,
take a recursive realization with the base and refined level
characters, a normalization of the real rows, and the local
row comparison between that normalization and the recursive
realization. This is the choice of comparison denoted by a
summand of the bundle. Denote the row spaces before and after this comparison by superscripts \(1,2\). The same identifications apply simultaneously to both rows and the word columns:
\[
 V_{\rm full}^{(1)}=V_{\rm full}^{(2)},\qquad V_{\rm good}^{(1)}=V_{\rm good}^{(2)},\qquad M_u^{(1)}=M_u^{(2)}\quad(\text{each word }u).
\]

The row normalization uses the equality of the character
frames, the identification of the type primes with the
kernel's type places and the conjugate-level and character
identities of the finite translate. It keeps the standard
frames and the coefficient values on the standard word
vectors. For the recursive comparison use the standard
additive character on these levels and the local row
reading of the marked Galois level. The degree-three
normalization uses the equality of the zero section with
the one in the original scale datum.

For \(z=v\) or \(\bar v\), write \(I_j(z),J_j(z),Z_j(z)\) for the
base-row index, lower-map index and value of the zero-section test at lift \(j\);
nonzero means that the corresponding test detects the place. The twist
support has the following intrinsic description:
\[
 \mathcal S_{\rm tw}
 =\{v:\ (I_j(z),J_j(z),Z_j(z))\ne(0,0,0)
            \text{ for some }j\text{ and }z\in\{v,\bar v\}\}.
\]
All three tests precede the choice of the height and error constants.
Let \(S\) be the set selected by the three tests and \(\mu_0\) the base tuple. The tests identify the tuple map on conjugate places by
\[
 S=\bar S,\qquad v\in S\ \Longleftrightarrow\ \mu(v)\ne\mu_0,\qquad \mu(\bar v)=\mu(v).
\] In particular each
single-index place is twisted.
These are comparisons on \(x\). Substitution of the three
family equalities of Section~\ref{subsec:completed-division}
transports them to the bundle with \(R\) fixed.
\label{cor:completed-division-6}

\input{linked-family-rows}

The two runs below use the rows of Proposition~\ref{prop:linked-comparison-rows}
for the selected compatible comparison.

\begin{lemma}[The two prime orientations on the common core]
For \(\ell=p_d,p_d'\) there is a full short-section extension of
the same core \(J\) with distinguished prime \(\ell\) at every
short weight. For a twisted place \(v\), the determinant character and short cover use the marked half prime of the assigned orientation, where
\[
 p(v)=\begin{cases}p_d',&v\mid p_d,\\ p_d,&v\nmid p_d.\end{cases}
\]
\end{lemma}
\begin{proof}
Write \(J\) for the common geometric core at \(P\). It
has its canonical frames, positive ranks and half-level
divisibility of the error data. The two field conditions and
the values of both distinguished characters hold in the
prescribed fields. The half-level ideal \(\mathfrak m_b\)
of the common core is the product of the marked type prime
ideals of the row. In fact the full level is squarefree
and factors into its two indicated halves. The divisor
comparisons and the equality of the core and row half
primes give the same prime factors on the indicated
side. Equality follows from the squarefree factorization.

Fix \(\ell\in\{p_d,p_d'\}\). For this orientation choose
the distinguished character \(\delta_\ell\) of the frame at
\(\ell\), and use the field condition at \(\ell\). Define
the own prime of a coefficient place \(v\) by
\[
  p(v)=
  \begin{cases}
    p_d' & \text{if the residue characteristic of \(v\) is \(p_d\)},\\
    p_d & \text{otherwise}.
  \end{cases}
\]
For twisted places with \(p(v)=\ell\), the marked side
comparison supplies one prime of the imaginary quadratic
field over \(\ell\) for all the quotient primes over
\(\ell\) on the tuples of those places. Extension of that
prime is contained in every one of these quotient primes,
by contraction. If the class of such places is empty, any
prime above \(\ell\) suffices for this orientation.

On the finite ideles the character construction gives
characters \(\xi_\ell,d_\ell\) with
\[
                        \xi_\ell(a/\bar a)=d_\ell(a).
\]
Here \(d_\ell\) is the prescribed finite-idelic restriction of
\(\delta_\ell\). A character of finite order on the ideles
is trivial on the archimedean subgroup of the CM field,
since its archimedean part is a product of copies of \(\C^\times\).
For a character \(\xi\) of order \(n\) and any archimedean idele
\(z\), take an archimedean \(w\) with \(w^n=z\). Then
\[
 \xi(z)=\xi(w^n)=\xi(w)^n=1 .
\] Thus
its restriction to the finite ideles inherits triviality on principal ideles,
triviality on real ideles, and the specified coefficient
values.

The finite character is trivial on integral units congruent to \(1\)
at the prescribed level. At a prime over \(\ell\), choose
a unit with a nontrivial character value, using the
distinguished character's prescribed tame ramification.
An idele equal to this unit at that prime and to \(1\)
elsewhere is congruent to \(1\) at the two type primes
because their rational primes differ from \(\ell\).
This gives the required nontrivial test there.

On finite ideles put \(h(b)=b/\bar b\). Finite idele Hilbert
90 gives \(\operatorname{im}h=\ker N_{E/E^+}\). Since
\(\ker h\) consists of real ideles and \(d_\ell\) is trivial
there, the assignment
\[
 \xi_\ell(h(b))=d_\ell(b)
\]
is well defined. Extend this character to the finite idele
group; its required congruence and finite-order properties
concern the norm-one subgroup. For a principal norm-one
element \(a\), explicit preimages are
\[
 b=1+a\ (a\ne-1),\qquad
 b\in E^\times,\ \bar b=-b\ (a=-1).
\]
Thus \(\xi_\ell(a)=1\) on principal norm-one elements.
The congruence form of Hilbert 90 transports the tame
congruences at the two type primes and at \(\ell\), giving
the stipulated ramification and coefficient values.
\label{cor:completed-division-7}

Extend the core to the upper geometric model, then take the
full short-section extension with the prime and the
character just chosen. These constructions use the laws
of \(J\), the canonical frames, positive ranks and the
host hypotheses at \(\ell\). The full extension has the
same core \(J\), and its short cover at each lift has
distinguished prime \(\ell\) at every short weight.
For a determinant root \(b\) of an adelic element \(g\), let \(b_f\)
be its finite part. The finite determinants and character values satisfy
\[
 \det(g)_f=b_f/\bar b_f,\qquad
 \xi_\ell(\det(g)_f)=d_\ell(b_f).
\] Rational determinant roots give
value \(1\). At a twisted \(v\) with \(p(v)=\ell\), the
same formula gives the specified determinant character and
the chosen short cover above the marked half prime. This
verifies the own-prime and character pinning at that place.
Taking \(\ell=p_d\) and \(\ell=p'_d\) gives the two extensions.
\label{cor:completed-division-8}
\end{proof}

\input{split-level-sections}

\begin{lemma}[One level action for both orientations]
Let \(\bar A_g^{(1)},\bar A_g^{(2)}\) be the two induced operators of the finite row group on the core half-level quotient \(\Lambda/\mathfrak m_b\Lambda\) at lift \(j\). With the common level-form identification,
\[
 \bar A_g^{(1)}=\bar A_g^{(2)}\quad\text{on }\Lambda/\mathfrak m_b\Lambda\qquad\text{for every }g.
\]
This action intertwines that on the level forms.
\end{lemma}
\begin{proof}
Write \(\Lambda\) for this lattice and \(A_g\) for the core level
map at \(g\). The congruences of the level comparison give
\[
 \bar A_1=1,\qquad \bar A_g\bar A_h=\bar A_{gh}
       \quad\text{on }\Lambda/\mathfrak m_b\Lambda .
\]
In particular \(\bar A_{g^{-1}}\) is the two-sided inverse
of \(\bar A_g\). This action is defined on \(J\) before either
extension. Two integral lifts of \(\bar A_g\) differ by a map
into \(\mathfrak m_b\Lambda\); on the complementary level
sublattice their difference lands in the full level multiple.
Thus they give the same allowed change of level vector. Pullback
on forms depends only on the quotient map, and the geometric
action in the common frame intertwines it with the row action.
\label{cor:completed-division-9}
\end{proof}

At \(v\) use the extension at \(p(v)\). Its own prime
differs from the residue characteristic, including when
the residue characteristic is \(p_d'\). The place
\(\bar v\) uses the same orientation as \(v\).
Consequently the short-cover laws, the
determinant comparisons and the one level action apply
on the indicated parts of this single numerical point.
No scale cutoff is used in this selection. The linked row
hypotheses and the numerical inequalities for section counts
have their own thresholds.
\label{cor:completed-division-10}


\begin{lemma}[Common form coordinates on the two runs]
At each lift \(j\), write \(I_j\) for the common arithmetic frame identification. The same inverse coordinate map can be used for both orientations:
\[
 z_j^{(1)}=z_j^{(2)}=I_j^{-1}=:z_j.
\] It reads the
same finite action, good projector and depth scalar at \((v,j)\).
The short sections at \(v\) come from its own prime \(p(v)\).
\end{lemma}
\begin{proof}
Use the phase-one base row of \(P\). Its weights are the positive
entries \(a^+_{j,\tau}\) and the reversed negatives of
\(a^-_{j,\tau}\), of signature \((k+j-1,m-j)\) and central
weight \(n_V\). These are the product order and the light and
short weights of the numerical record.

Let \(I_j\) denote the injective comparison of core cusp forms
with the arithmetic frame. The level quotient determines the
pullback \(A_g\) for each element of the mirabolic group.
It suffices to read the representation \(\rho\) at one coefficient
place, because \(I_j A_g=\rho(g)I_j\) in the common frame.
In particular
\[
 I_jA_gA_h=\rho(g)\rho(h)I_j=I_jA_{gh}.
\]
Injectivity gives composition, and gives linearity by the same
argument. The basis comparison makes the scalar extension of
\(I_j\) invertible.

Take \(z_j=I_j^{-1}\) on this extension, as
in Paragraph~\ref{par:selected-form-coordinates}. For the
arithmetic projector \(e_{\rm good}\) and a rational class-sum
map \(T_\mu\) in these coordinates, the corresponding form
readings are \(z_je_{\rm good}I_j\) and \(z_jT_\mu\).
Thus
\[
 I_j(z_je_{\rm good}I_j)=e_{\rm good}I_j,\qquad
 I_j(z_jT_\mu)=T_\mu .
\]
Use this one isomorphism for the ambient lattice and the scaled
and unscaled short sections at \(p(v)\). The section lattice in
the good row is the projector image of the theta lattice at
\(\mu(v)\).

The embedding of the host integers into the point field factors
through the integral coefficient ring of the core, and agrees
with complex conjugation on the host involution. It therefore
reads the same half-level and error generators and the same
scalar in the form and depth comparisons. Its integral exponent
is the depth in the local section estimate. The remaining
parameters \(f,m,r,k,x,\epsilon,b',X''\), the type primes and
the determinant and cutoff data are those of the common numerical
point. Only the local short-section spans depend on the
distinguished prime. \label{cor:completed-division-11}
\end{proof}


\begin{proposition}[One remainder before choosing the cut]
At the common point of Proposition~\ref{prop:completed-division-pair},
assume the linked rows of Proposition~\ref{prop:linked-comparison-rows}
and the section counts of Proposition~\ref{prop:slice-rank-shares}.
Use the simultaneous short basis of
Lemma~\ref{lem:slice-short-basis-before-localization} on the cut by
the selected light columns, with \(C_{n_W}=10\), \(\log g_i\le9L+1\), \(C_P\ge1000\)
and \(C_P^{20}\le frL\). At sufficiently large scale the section
remainder on each lift of either orientation satisfies
\[
                 \langle r\rangle\le R_j fmb\delta_2 .
\]
The scale threshold is common to the orientations and lifts.
At each place the selected remainder precedes the choice of the
cut and its integer depth.
\end{proposition}
\begin{proof}
Use the core coefficient bound \(C_K\) in \(S_\kappa\le C_K kfb\)
in the scale inequalities \eqref{eq:slice-remainder-scales}.
Write \(X_{\rm net}(20,C_K)\) for a scale beyond which the size
inequalities of Proposition~\ref{prop:slice-evaluation-net} hold
with exponent \(20\) and this core bound. Increase the common
scale threshold \(x_1\) to satisfy
\[
 x_1\ge\max\{10^{60},X_{\rm net}(20,C_K)\}.
\]
Read the half-error divisibility on \(J\), and the nonzero short
sections on the selected full extension. The full, short and
light weights and the index depth are those of
Section~\ref{subsec:paired-row-indices}.

At \(p=p(v)\) fix the split marking, once for all three weight spaces. The input level and its special character are
\[
 g_0g_1p,\qquad g\longmapsto\xi_p(\det g),\qquad \xi_p(y/\bar y)=d_p(y).
\] Read them in the core weight spaces at
\(\kappa_j,w_{F,j},w_{L,j}\) by the common form isomorphism.
The incidence chart is the one for the light columns of this
run on the core. Its full level in the light-weight inequalities
is \(g_0g_1\). Select the \(d_s\) scaled short sections before localization by
Lemma~\ref{lem:slice-short-basis-before-localization}, using independent
evaluation on the geometric generic intersection. The shifted
discriminant of this basis is computed in \eqref{eq:slice-discriminant-order}.

If \(A\) is the selected square evaluation minor over the generic
field \(K\), then for each branch field \(K_\zeta/K\) one has
\[
                    0\ne\det A\in K\hookrightarrow K_\zeta .
\]
Thus the same sections give full rank on every branch. Their local
integral index still contributes to \(I_\lambda\). The choice by rank in the
boundary and complementary spaces preserves the light columns
and, before the final lift, the same auxiliary section
\eqref{eq:selected-eta-product}. Hence the finite valuation
test, the two sections for the height calculation and the
determinants concern this one family with the new choice of sections. The evaluation
net gives the second norm on the same rational spaces and
integral section lattices. Changing this norm in the
determinant proof preserves the index on those modules.


Let \(e\) be the branch ramification index. At its reference point the weight-\(w_F\) section \(g_0\) fixes the valuation offset
\[
 c_\zeta=v(g_0(y_\zeta)).
\]
It cancels against the denominator in the quotient reading.
The Gauss value \(\nu_\zeta\) of the generic light row is the
minimum over its light columns, giving the shifted order
\[
              e\bigl(v_{y_\zeta}(F)-T\nu_\zeta\bigr).
\]
This uses the normalized point valuation of the depth estimate.
The vertical, horizontal and archimedean comparisons in
\eqref{eq:slice-height-disc} bound the discriminant and auxiliary
costs on this slice. The remaining two costs are its boundary
determinant and complementary CM determinant heights. Their
ranks give the \(s^2/T\) savings in \eqref{eq:geometric-four-terms}.
At the definite lift the auxiliary cost and complementary rank
are zero. The four bounds give the claimed inequality with the
chosen finite net and the same polynomial bounds at \(frL\)
for every \(j\) and both \(p\).
\label{cor:completed-division-12}
\end{proof}

Extending the form frame to the point field now puts this remainder
estimate in the coordinates of the paired index bound. The
restriction and depth argument uses that bound on the same
lattice and level action. \label{cor:completed-division-13}

\input{joint-family-estimate}



\begin{auxlemma}[The integral word column in two presentations]
\label{lem:word-prescribed-class-lattice}
Fix a word of Lemma~\ref{lem:window-parameter-choice}, its split prime
\(T=T_u\) and the degree-zero column of
Proposition~\ref{prop:measured-columns}. Write \(U\) for its marked
level and \(D_U=H(F)\backslash H(\mathbb A_{F,f})/U\) for its
definite class set, where \(H\) is the column group over the totally
real field \(F\). Use the transported integral tensor Weyl values of
Proposition~\ref{prop:level-four-integral-kernel} at each representative.
Let \(L_U\) be the module of class forms with these integral values.
For \(v\mid T\) let \(\mathrm{Tr}_v\) sum the translates by the
residue diamond group at \(v\), using the prescribed level
character. In a second presentation obtained by the marked class
transport the original and trace-zero lattices and their metrics
are carried into one another:
\[
 L_U\simeq L_{U'},\qquad
 \bigcap_{v\mid T}\ker(\mathrm{Tr}_v|_{L_U})\simeq
 \bigcap_{v\mid T}\ker(\mathrm{Tr}'_v|_{L_{U'}}).
\]
Here \(U'\) denotes the transported level. On an auxiliary refinement
the equalities concern its subspace of pullbacks.
\end{auxlemma}
\begin{proof}
Retain the finite representative, split halves and ordered flag when
changing coefficient coordinates. For a class \(a\) with stabilizer
\(\Gamma_a\), let \(B_a\) denote its integral value lattice.
A class form is determined by one value per class. Changing a lift of
that class applies the stabilizer with its specified character. Thus
the lattice in its rational function space is
\[
 L_U=\bigoplus_{a\in D_U}
  \{z\in B_a:\ \rho_a(\gamma)z=\chi_a(\gamma)z
                              \quad(\gamma\in\Gamma_a)\}.
\]

At the word prime write the chain with reduction \(\pi_T\) as
\[
   L_1=\pi_T^{-1}(H)\subset L_0,\qquad
       H\subset L_0/T L_0,\qquad [L_0:L_1]=T .
\] In an adapted basis, a basis change on
one half acts by its inverse transpose on the paired half. Both are
integral and invertible in the respective bases. Tensor powers,
their divided-power lattices and the integral maps defining the Weyl
values are functorial for such isomorphisms. One can also keep exactly
the same adapted basis in the two presentations. This identifies
\(B_a\), its stabilizer equations and the residue diamond translations.

For the local diamond action \(R_h\) the operator used on integral values is
\[
              \mathrm{Tr}_v=\sum_{h\in\mathbb F_T^\times} R_h ,
\]
in the marked character convention. Its coefficients are integral even at residue primes dividing \(T-1\). Over an extension of coefficient rings of
integers, flatness preserves the kernels of the stabilizer equations
and of the finite list of traces. This proves the lattice assertions,
including at places dividing the order of a diamond group.

The compact real value space and its Hermitian form are transported
with each representative. At the same level this only reindexes the
Petersson sum. At a refinement \(U_1\subset U\), the classes over
\(a\) are the \(\Gamma_a\)-orbits in \(U/U_1\). Orbit--stabilizer gives
\[
 \sum_{a_1\mapsto a}\frac{1}{|\Gamma_{a_1}|}
             =\frac{[U:U_1]}{|\Gamma_a|},\qquad
 \operatorname{vol}(U_1)[U:U_1]=\operatorname{vol}(U).
\]
The Petersson norm of a pullback is therefore unchanged for the
original Haar measure. Integrality of its old values can be tested on
the pullback since the projection of class sets is onto. For an
auxiliary prime different from \(T\), the diamond action above \(T\)
commutes with the forgetful map. Hence the trace-zero condition on
old forms is the same in the refinement. The mixed coefficient
lattice and the image of the higher-level Hecke map are subsequently
computed on this column in Propositions~\ref{prop:ramified-mixed-column}
and~\ref{prop:ramified-selected-hecke-colength}.
\end{proof}

\paragraph{Simultaneous choices at every translate.}
\begin{proposition}[Completion at every translate]\label{prop:translate-completion}
For the family of Lemma~\ref{lem:window-parameter-choice} with the marked data of Construction~\ref{hyp:family-prescribed-realization} at its admissible word levels, above a common threshold the compatible
realizations can be chosen at all translates \(u_0\in A\) of the
same scale. Their field, ranks, top labels and larger-weight label
sets agree with the base data through the specified identifications.
At every translate the strict scalar lengths on that same field and rank satisfy
\[
 \sum_\ell(\log\ell)\sum_{1\le a<t}V(E,\ell,a)\le\sqrt b .
\]
Each member also inherits the same logarithmic derivative bound with constant \(4\).
\end{proposition}

\begin{proof}
Fix the scale and translate before choosing any realization. By Lemma~\ref{lem:window-parameter-choice} the two weight lists and the labels of their columns are already fixed. Write \(I_u\) for the interval determined by the pre-level mass. The ray class count chooses primes subject to the prescribed ray and splitting conditions with
\[
                   (T_u)_u\in\prod_u I_u .
\] Apply the initial construction in Construction~\ref{hyp:family-prescribed-realization} to this list of primes. The signed integral bases of Proposition~\ref{prop:retwisting-row-basis} supply finitely many character requirements. For every served index the character and its level are selected once for both signs and all coprime powers.

Reconstruct the row at these characters using Lemmas~\ref{lem:second-core-reference} and~\ref{lem:new-ray-type-occurrence}. Proposition~\ref{prop:second-specialization-row} identifies the pulled-back marked places and the spectra on the rebuilt row. The numerical masses are read there before prescribing the local divisor lattices. For each word, Lemma~\ref{lem:word-lattice-gluing} constructs the lattice with the specified localizations. Its fractional-ideal shift and its height can now be computed on that same word. Making the simultaneous substitution of Section~\ref{subsec:automorphic-realization} gives an element of the set \(\mathcal P(x,u_0)\) of completions over this translate.

For fixed \(x\), the preceding argument gives each factor of the finite product below. Choosing one element of each factor therefore gives
\[
 (P_{u_0})_{u_0\in A}\in\prod_{u_0\in A}\mathcal P(x,u_0).
\]
The threshold was fixed before the translate and the list of primes. No supremum over point-dependent cutoffs is used. Each completed member has the prescribed
ranks, field, torus, Artin map, two auxiliary primes, type
primes and weight functions under its identifications with
the prescribed data. On the set of indices with larger weights this follows by
transport of the actual label set; the index map on labels
is injective, so taking its image preserves the set and its
cardinality. Writing \(a\mapsto a'\) for the index map and \(s,s'\) for the corresponding binary sequences, the coordinate comparison is
\[
 s'_{a'}=s_a\quad\text{at every corresponding index}.
\]

In particular every translate of
the original weights occurs, and the logarithmic derivative
bound with constant \(4\) and the signed scalar length bound
\(\sqrt b\) pass from the base data to each such point. The latter
bound is for the scalar cohomology sum over the same field. It enters the
cohomological contribution to the divisibility bound after the length comparison of
Section~\ref{subsec:division-length}.
\label{cor:window-points-4}
\end{proof}

\begin{auxlemma}[A common prime pair on each interval of one family]
\label{lem:disjoint-common-pair-windows}
Fix the primitive target \(\psi\), its weights from
Definition~\ref{def:family-weights}, and \(0<\epsilon<1/52\).
There are intervals \(I_\nu=[Y_\nu,Y_\nu^{3/2}]\), \(\nu\ge1\), and
prime pairs \((q_\nu,d_\nu)\) such that \(Y_1\ge 2^{100}\),
\(Y_{\nu+1}>Y_\nu^{3/2}+1\), and the same pair serves every
\(x\in I_\nu\). More precisely, write \(\kappa=\kappa_{\rm sc}\)
as in Lemma~\ref{lem:rescale-host-budget}, and let \(e\) be the
dyadic index function of Proposition~\ref{prop:original-window-budget}.
On \(I_\nu\) we may use the field
\[
 E(x)=E_{\psi,e(x/\kappa);q_\nu,d_\nu}
\]
of Definition~\ref{def:family-original-host}. At \(x\) the prime-pair
bounds of Proposition~\ref{prop:quadratic-prime-pair} and the
strict-length bound at the rank \(t(x)\) of
Proposition~\ref{prop:actual-rank-budget} both hold.

Put \(S=\bigcup_{\nu\ge1} I_\nu\). In the small-\(\epsilon\) range and with the marked data of
Proposition~\ref{prop:translate-completion}, the assignment on \(S\)
can be retained at every completed translate above its threshold.
For any finite collection of estimates valid eventually on the resulting
family, with thresholds uniform over its scales and translates, there
are arbitrarily late entire intervals \(I_\nu\) on which all these
estimates hold together.
\end{auxlemma}

\begin{proof}
Fix \(e\), and take \(J\ge1\) and its threshold \(X_0\) from
Proposition~\ref{prop:original-window-budget}. With this \(J\), let
\(U_0\) be the threshold in Proposition~\ref{prop:quadratic-prime-pair}
with \(M=f_\psi\) and \(K=J\). Its character list contains the
primitive inducers at level \(8f_\psi\), and its pair meets the
conditions of Proposition~\ref{prop:common-host-prime-pair} with
\(\tau=\epsilon/2\). Given \(z\ge1\), put
\[
 M=\max\{z+1,2^{100},U_0,\kappa^3,\kappa^{400}\},
 \qquad X=\max\{X_0,M^J\},\qquad U=X^{1/J}.
\]
Choose \(q,d\) for the whole interval \([U,U^J]=[U,X]\).
For this pair take \([W,W^2]\subseteq[U,X]\) from
Proposition~\ref{prop:original-window-budget} and set \(Y=\kappa W\).

For every \(Y\le x\le Y^{3/2}\), the inequalities
\(\kappa^3\le W\) and \(\kappa^{400}\le x/\kappa\) give
\[
 U\le W\le x/\kappa\le x\le(\kappa W)^{3/2}\le W^2\le X,
 \qquad
 t(x)\le\lfloor x^{400}\rfloor\le\lfloor(x/\kappa)^{401}\rfloor .
\]
The rank inequality is Lemma~\ref{lem:rank-window}.
Thus the same pair is available at \(x\) and at the argument
\(x/\kappa\) of the strict-length estimate. Its field may vary
through \(e(x/\kappa)\). Restricting \eqref{eq:cohomwin} to
\(1\le j<t(x)\) gives, for every finite set of primes \(P\),
\[
 \sum_{\ell\in P}\log\ell\sum_{1\le j<t(x)}
             V(E(x),\ell,j)\le \exp((x/\kappa)^\epsilon/2)\le\sqrt b .
\]

Start with \(z=1\). For the next interval take \(z=Y^{3/2}+1\). Iteration gives the claimed
separation, so every \(x\in S\) has a unique assigned pair. Specify
this pair as part of the original field data at \(x\).
Proposition~\ref{prop:translate-completion} retains that field and its
marked prime data at each translate. In particular the construction
can use the assignment throughout \(S\). Fix the family, including
its successive parameter choices of Lemma~\ref{lem:window-parameter-choice}.
If the finitely many estimates on this family have thresholds
\(X_1,\dots,X_h\), then for any \(Z\ge1\) take \(\nu\) with
\(Y_\nu>\max\{Z,X_1,\dots,X_h\}\). Their hypotheses on the scale
then hold at every real \(x\in I_\nu\), for all the translates
covered by the respective estimates.
\end{proof}

\begin{construction}
\begin{tabular}{@{}ll@{}}
\toprule
\textit{Range} & \textit{Choice on that range}\\
\midrule
\([U,U^J]\) & One pair \((q,d)\) for the quadratic twists\\
\([W,W^2]\subset[U,U^J]\) & Strict lengths on the original field at \(x'=x/\kappa\)\\
\([Y,Y^{3/2}],\ Y=\kappa W\) & Actual rank \(t(x)\) and all translated rows at \(x\)\\
\(Y_\nu>\max(Z,X_0)\) & All estimates on one of the intervals already chosen\\
\bottomrule
\end{tabular}
\caption{Selection of an interval in Lemma~\ref{lem:disjoint-common-pair-windows}. The first three choices give a sequence of disjoint intervals. The last line is used after bounding the two division errors and the Gram error on the resulting family.}\label{fig:common-window-choice}
\end{construction}

\begin{proposition}[A common application range]\label{prop:all-window-applications}
Fix a primitive target and its weights. In the range of
Lemma~\ref{lem:window-parameter-choice}, let \(S\) be the union of
intervals in Lemma~\ref{lem:disjoint-common-pair-windows}. One can
choose the base data at the scales in \(S\) and a family of their
completed translates with the following property. There is \(X_0\)
such that
\[
 \forall x\in S\cap(X_0,\infty)\quad
 \forall u_0\in A(x)\quad
 \exists\,P(x,u_0)\text{ over the chosen base at }x .
\]
On these members the two division estimates, the projected Gram
estimate and the height cutoff hold simultaneously with family
constants. The weight functions at a scale are fixed before its
translate. On each interval the prime pair is fixed before the scale.
\end{proposition}

\begin{proof}
\proofstep{Fix the family.}
Use the outer, middle and final parameters of
Lemma~\ref{lem:window-parameter-choice}, including the prescribed
exponent for the split word primes. Use the intervals with their
prime pairs as in Lemma~\ref{lem:disjoint-common-pair-windows}.
At each sufficiently large \(x\in S\), choose the top set and the
two weight functions by Section~\ref{subsec:frobenius-selection}
and Lemma~\ref{lem:one-share-vector}. The same choice has the
required coset counts and mean loss at every translate.
Proposition~\ref{prop:two-type-primes} supplies the two split types
on this base. The rank bounds in Lemma~\ref{lem:rank-window} give
\(r=k+m-1\le2x^{9/2+10\epsilon}+2\sqrt{x}\le x^5\)
eventually. The chosen primes satisfy
\[
  (kb)^8<g_i\le 2(kb)^8,\qquad
  B_r(g_0)\cap B_r(g_1)=\varnothing .
\]

Thus for each residual prime one of the two types is available,
and \(\min_i g_i\ge2r\) eventually. On the standard Hermitian
space use the diagonal form with \(p_+\) positive and \(p_-\)
negative unit entries. For its polarized homology lattice take a
copy of the fractional ideal \(I\) of
\eqref{eq:host-different-factors} on each coordinate, as in
Lemma~\ref{lem:marked-principal-cm-sources}(ii).
Writing \(\beta\) for the imaginary scalar in the trace pairing,
the dual on each coordinate is
\[
 (\beta)^{-1}\mathfrak D_E^{-1}\bar I^{-1}=I .
\]

At a split place choose local generators on one half and the dual
basis on the other. This supplies the self-dual marking used in
Section~\ref{subsec:automorphic-realization}. Take the two deletion
characters from Lemma~\ref{lem:split-deletion-characters}. Their
ramification sets are disjoint. The
finite field requests of Definition~\ref{hyp:initial-adelic-frame}
are met by adjoining their entries in the specified order. On the
resulting rank interval Proposition~\ref{prop:assigned-prime}
gives the levels for Construction~\ref{hyp:family-prescribed-realization}.
Proposition~\ref{prop:translate-completion} constructs the compatible
completions at those levels. In each fibre use the choice from
Proposition~\ref{prop:estimates-each-translate} for the weight pair
just fixed. It gives a member with the stated count and prime losses.
Choose one such member for each index in
\[
     \{(x,u_0):x\in S,\ x>x_1,\ u_0\in A(x)\},
\]
where \(x_1\) is the common threshold for constructing points.
If an index below \(x_1\) is needed, use a single already
constructed fibre there. All sufficiently late indices still read
their actual scale.

\proofstep{Identify the arithmetic estimates on this family.}
The two extensions at distinct auxiliary primes are made on the
reconstructed member. Proposition~\ref{prop:linked-comparison-rows}
identifies its reference pairing through the projection to the good
row. Lemma~\ref{lem:parent-slice-lattices} gives the integral
branch and product lattices there. In the notation of
\eqref{eq:run-cut}, restriction to a subspace over \(v\mid\lambda\)
uses
\[
     E=e(v/\lambda)h,\qquad r_v=e(v/\lambda)r_\lambda.
\]
Proposition~\ref{prop:uniform-parent-remainders} gives common
thresholds for the two covers. For the numerator, use
Proposition~\ref{prop:actual-rank-budget} on the field of the
member and Lemma~\ref{lem:strict-budget-after-cut} after restriction
to its good characters. The weighted return in
Proposition~\ref{prop:weighted-dual-cost} uses their original
multiplicities. Lemma~\ref{lem:joint-family} and
Proposition~\ref{prop:parent-dual-cost} therefore bound the two
indices at this same member.

\proofstep{Take the threshold after choosing the family.}
The measures for its period maps are the quotient measures in
Proposition~\ref{prop:gram-unfolded-measures}. Proposition~\ref{prop:pretransport-ramified-period}
uses the dimension at each assigned word prime before transporting
the period. This gives the column dimensions for the Gram comparison
with the same identity contribution \(D_*\). With the height
exponent chosen by Proposition~\ref{prop:marked-frame-height},
Proposition~\ref{prop:selected-rank-comparison} estimates the
normalized Gram operator and the dimension of the projected space
at that exponent. Proposition~\ref{prop:estimates-each-translate}
gives the count, prime losses and division bounds with
\[
      N=\lceil2a\rceil,\qquad (e_D,e_L)=(2,0).
\]

These bounds apply at the members selected in the first step.
Restricting a bound valid throughout the constructed family to
these members preserves its constant. The substitution in
Lemma~\ref{hyp:absorption-inputs} now involves finitely many
uniform constants \(K_i\) and scale thresholds \(X_i\).
Take \(K=\max(6,\max_i K_i)\) before applying
Lemma~\ref{lem:uniform-numerical-application}. With its threshold
\(X_{\rm num}\), put
\[
             X_0=\max\{x_1,X_{\rm num},\max_i X_i\}.
\]
For every prescribed lower bound there is an entire interval in
\(S\) above \(X_0\) and that bound. This proves the simultaneous
application with the stated order of choices.
\end{proof}

Apply the arithmetic comparison in Section~\ref{subsec:fixed-weight-absorption}
on the family of Proposition~\ref{prop:all-window-applications}.
It bounds the total
difference at every translate by \(x^{5\epsilon}\mathcal B\).
The normalized average and Lemma~\ref{lem:comparison} give
\[
       G_{\mathcal K}(x)\le x^{1/2+9\epsilon}
          \qquad(Y\le x\le Y^{3/2}).
\]
The character family and weights are the ones fixed before
choosing \(\epsilon\). Since the lower bound \(Z\) is free, the order of choices in Proposition~\ref{prop:arithmetic} is
\[
 \forall Z\ \exists\,Y\ge Z\quad\text{with the bound throughout }[Y,Y^{3/2}].
\]
\label{cor:window-points-5}


For the interval construction, Section~\ref{sec:absolute} normalizes the row amplitudes and Section~\ref{sec:transport} identifies the integral periods under the level and type conditions. The jet estimates give a fractional ideal of local divisibilities (Section~\ref{sec:finite-valuations}). The index and slopes of the trace-zero column at its assigned prime and the other places are calculated in Sections~\ref{sec:ramified-hecke-index} and~\ref{subsec:column-orbit-length}.

At the assigned prime, Sections~\ref{subsec:hecke-period} and~\ref{subsec:hecke-entries} bound the split kernel in an incident frame. Sections~\ref{subsec:incident-sections} and~\ref{subsec:second-period} identify the series and sum the complementary cycles on the same pairs. Section~\ref{subsec:word-period} transports both bounds along the column, using the measure and projections of Sections~\ref{subsec:word-moments}--\ref{subsec:word-projection}. The entry and period errors for the reference and twisted tuples (Sections~\ref{subsec:baseline-jets}--\ref{subsec:division-shift}) then enter the upper comparison together with the cohomology, lift, row-index and geometric bounds at the same scale. \label{cor:family-window-4}
